% ============ 《表示论》中文翻译 LaTeX 重排版 ============
% 原书: Sophie Morel, MAT 449: Representation Theory (April 26, 2019)
% 模板: AJbook.cls (李文威\emph{代数学方法}书系模板)
%       ★ 定理编号规则已按原书修改(见 AJbook.cls): 定理/推论/引理/命题/定义/
%         例/注记/记号/约定/习题/注意 共享一个计数器, 按节编号
% 编号体系 (原书约定):
%   - 章 = 原书部分: I-VI 罗马数字 + 附录 A、B
%   - 节 = 原书节 (如 I.1, II.3), 编号 \thechapter.\arabic{section}
%   - 子节 = 原书子节 (如 I.3.1, IV.9.2), 编号 \thesection.\arabic{subsection}
%   - 定理类: 共享计数器, 编号 N.M 或 N.M.K 与原书一致:
%       无子节的节: I.1.1  (Definition I.1.1, Remark I.1.13)
%       有子节的节: I.3.1.1 (Definition I.3.1.1, Exercise I.5.1.1)
%   - 公式: 原书基本不编号, 少数引用处用显式 \tag{...}
\documentclass[geometry=a4, fontsetup=font-setup-open.tex, titlesetup=titles-setup, coverpage=cover, CJKthechapter=false]{AJbook}

% ============ 宏包 ============
\usepackage{mycommand}      % 常用数学命令(\dd \R \Z \Q \CC \e \ii \Hom \End 等)
\usepackage{myarrows}       % 可伸缩箭头
\usepackage{amssymb}
\usepackage{chngcntr}       % \counterwithout
\usepackage{etoolbox}       % \pretocmd 等
\usepackage{tikz-cd}        % 交换图
\usetikzlibrary{calc}       % 封面坐标计算
\usepackage{imakeidx}
\usepackage[bookmarksnumbered=true]{hyperref}

% ============ 编号体系 ============
% 章号: 正文用罗马数字 (I-VI), 附录用字母 (A, B)
\renewcommand{\thechapter}{\Roman{chapter}}
\renewcommand{\AJchapterttl}[1]{#1}             % 章号直接显示编号 (I, II, ...)

% 天眉: ``I 拓扑群的表示''. AJbook.cls 在 \AtBeginDocument 里定义 \chaptermark, 故此处覆盖须放在其之后
\AtBeginDocument{%
  \renewcommand{\chaptermark}[1]{\markboth{\thechapter\quad #1}{}}%
}

% 定理编号格式切换 (原书: 每个最底层编号单元内从 1 重新计数):
%   无子节的节: 编号 I.1.1  (\thetheorem = \thesection.\arabic{theorem})
%   有子节的节: 编号 I.3.1.1 (\thetheorem = \thesubsection.\arabic{theorem})
% 通过 titlesetup.tex 中 titleformat 的 after-code 钩子在 \section/\subsection 后自动切换
% (不能用 \apptocmd: titlesec 的 \section/\subsection 含 \new@ifnextchar, 会破坏参数扫描)
% 注意: after-code 在 group 内执行, 须用 \gdef 使重定义全局生效
\newcommand{\AJseccount}{%
  \gdef\thetheorem{\thesection.\arabic{theorem}}%
  \gdef\theHtheorem{\thesection.\arabic{theorem}}%
}
\newcommand{\AJsubseccount}{%
  \setcounter{theorem}{0}%
  \gdef\thetheorem{\thesubsection.\arabic{theorem}}%
  \gdef\theHtheorem{\thesubsection.\arabic{theorem}}%
}

% 超链接
\hypersetup{
    pdfauthor={Sophie Morel},
    pdftitle={表示论 (Representation Theory)},
    pdfsubject={表示论, 拓扑群, 调和分析}
}

% 索引 (intoc=false: 目录条目由 titlesec 无号章 after-code 添加, 避免重复)
\makeindex[columns=2, intoc=false, title={索引}]

% 参考文献 (biblatex 由 AJbook.cls 引入, sorting=nyt 与原书字母序一致)
\addbibresource{refs.bib}

% 全局断行弹性: 消除边界性 Overfull (与 Gelbart/古典调和分析项目同款)
\emergencystretch=3em

% ============ 文档开始 ============
\begin{document}
\frontmatter        % 封面(cover.tex) + 目录(模板自动处理, 罗马页码)

\mainmatter         % 正文(阿拉伯页码从 1 开始)

% 原书部分 I-VI
\chapter[拓扑群的表示]{拓扑群的表示}
% ============ 第 I 部分: \index{拓扑群}拓扑群的\index{表示}表示 (Representations of Topological Groups) ============
% 原书 I.1 - I.5 (原书 pp.7-118)
% 编号: 无子节的节用 \section 后定理编号 I.1.x; 有子节的节用 \subsection 后编号 I.3.1.x

% 本部分文件结构 (与原书一致):
%   I.1 Topological groups (pp.7-10)
%   I.2 Haar measures (pp.11-18)
%   I.3 Representations (pp.19-29): I.3.1-I.3.5
%   I.4 The convolution product and the group algebra (pp.30-44): I.4.1-I.4.3
%   I.5 Exercises (pp.44-118): I.5.1-I.5.6

\section{拓扑群}
% 原书 I.1 (pp.7-10)

\begin{definition}[拓扑群 (topological group)]
\label{def:I.1.1}
一个拓扑群 (topological group) 是一个带群结构的拓扑空间 $G$, 使得乘法映射
$G \times G \to G$, $(x,y) \longmapsto xy$ 与取逆映射 $G \to G$, $x \longmapsto x^{-1}$ 都是连续的.
\end{definition}

我们通常用 $1$ 或 $e$ 表示 $G$ 的单位元.

\begin{example}[拓扑群的例子]
\label{exa:I.1.2}
\begin{itemize}
\item 任何赋予离散拓扑的群都是拓扑群. 常用的例子包括有限群、自由群 (自由交换群与自由非交换群) 以及``算术''矩阵群, 如 $\GL_n(\Z)$ 与 $\SL_n(\Z)$.
\item $\R$ 和 $\CC$ 的加法群是拓扑群.
\item 群 $\GL_n(\CC)$ 赋以 $\CC$-向量空间 $M_n(\CC)$ 上任意范数给出的拓扑, 是一个拓扑群 (见习题 \ref{ex:I.5.1.1}), 因此其一切子群在诱导拓扑下也是拓扑群. 例如 $S^1 := \{z \in \CC \mid |z| = 1\}$, $\GL_n(\R)$, $\SU(n)$, $\SO(n)$ 等等.
\item $\Q_p$ 的加法群与群 $\GL_n(\Q_p)$ 是拓扑群. (见习题 \ref{ex:I.5.1.4}.)
\end{itemize}
\end{example}

\begin{definition}[\index{局部紧}局部紧 (locally compact)]
\label{def:I.1.3}
设 $X$ 是一个拓扑空间. 若 $X$ 的每一点都有一个紧邻域, 则称 $X$ 是局部紧 (locally compact) 的.
\end{definition}

\begin{remark}[局部紧与紧邻域基]
\label{rem:I.1.4}
若 $X$ 是 Hausdorff 的, 则这等价于 $X$ 的每一点都有一个由紧邻域组成的邻域基.\footnote{参考文献 ?}

注意我们并不假定拓扑空间中点的邻域是开的.
\end{remark}

\begin{notation}[子集的运算]
\label{not:I.1.5}
设 $G$ 是一个群, $A$, $B \subset G$, $x \in G$, $n \geq 1$. 我们使用以下记号:
\[
xA = \{xy \mid y \in A\} \quad \text{与} \quad Ax = \{yx \mid y \in A\},
\]
\[
AB = \{yz \mid y \in A,\, z \in B\},
\]
\[
A^n = AA \cdots A \quad (\text{$n$ 个因子}),
\]
\[
A^{-1} = \{y^{-1} \mid y \in A\}.
\]
\end{notation}

\begin{definition}[\index{对称子集}对称子集 (symmetric subset)]
\label{def:I.1.6}
设 $G$ 是一个群. 若 $G$ 的子集 $A$ 满足 $A = A^{-1}$, 则称 $A$ 是对称 (symmetric) 的.
\end{definition}

\begin{proposition}[拓扑群的基本性质]
\label{prop:I.1.7}
设 $G$ 是一个拓扑群.
\begin{enumerate}
\item[(1)] 若 $U$ 是 $G$ 的开子集而 $A$ 是 $G$ 的任意子集, 则集合 $UA$, $AU$ 与 $U^{-1}$ 都是开的.
\item[(2)] 若 $U$ 是 $1$ 在 $G$ 中的一个邻域, 则存在 $1$ 的开对称邻域 $V$ 使得 $V^2 \subset U$.
\item[(3)] 若 $H$ 是 $G$ 的子群, 则其闭包 $\overline{H}$ 也是 $G$ 的子群.
\item[(4)] 若 $H$ 是 $G$ 的开子群, 则它也是闭的.
\item[(5)] 若 $A$ 和 $B$ 是 $G$ 的紧子集, 则集合 $AB$ 也是紧的.
\item[(6)] 设 $H$ 是 $G$ 的子群. 则商 $G/H$ (赋予商拓扑) 具有以下性质:
  \begin{enumerate}
  \item[(a)] 若 $H$ 是闭的, 则 $G/H$ 是 Hausdorff 的;
  \item[(b)] 若 $G$ 是局部紧的, 则 $G/H$ 是局部紧的;
  \item[(c)] 若 $H$ 是正规的, 则 $G/H$ 是一个拓扑群.
  \end{enumerate}
\end{enumerate}
\end{proposition}

\begin{proof}
(1). 对 $x \in G$, 记 $l_x \colon G \to G$ (相应地 $r_x \colon G \to G$) 为左 (相应地右) 乘以 $x$. 又记 $\iota \colon G \to G$ 为映射 $x \longmapsto x^{-1}$. 由拓扑群的公理, 所有这些映射都是连续的.

注意到 $U^{-1} = \iota^{-1}(U)$, $AU = \bigcup_{x \in A} l_{x^{-1}}^{-1}(U)$, $UA = \bigcup_{x \in A} r_{x^{-1}}^{-1}(U)$. 因此 $U^{-1}$, $AU$ 与 $UA$ 都是开的.

(2). 我们可以假定 $U$ 是开的. 设 $m \colon G \times G \to G$, $(x,y) \longmapsto xy$. 则 $m$ 连续, 所以 $W := m^{-1}(U)$ 是开的. 因为 $1^2 = 1 \in U$, 我们有 $(1,1) \in W$. 由 $G \times G$ 上乘积拓扑的定义, 存在 $G$ 的开子集 $\Omega \ni 1$ 使得 $\Omega \times \Omega \subset W$. 由 $W$ 的定义有 $\Omega^2 \subset U$. 令 $V = \Omega \cap \Omega^{-1}$. 由 (1) 知 $\Omega^{-1}$ 是开的, 所以 $V$ 是开的, 且由定义知它是对称的. 显然 $1 \in V$ 且 $V^2 \subset \Omega^2 \subset U$.

(3). 考虑映射 $u \colon G \times G \to G$, $(x,y) \longmapsto xy^{-1}$; 则 $G$ 的非空子集 $A$ 是子群当且仅当 $u(A \times A) \subset A$. 又由拓扑群的公理, 映射 $u$ 连续. 因此对每个 $Z \subset G \times G$, 有 $u(\overline{Z}) \subset \overline{u(Z)}$. 将其应用于 $H \times H$ (其闭包为 $\overline{H} \times \overline{H}$), 我们得到 $\overline{H}$ 是 $G$ 的子群.

(4). 我们有 $G = H \sqcup ((G - H)H)$. 若 $H$ 是开的, 则由 (1) 知 $(G - H)H$ 也是开的, 因此 $H$ 是闭的.

(5). 乘法映射 $m \colon G \times G \to G$ 由假设连续. 由于 $AB = m(A \times B)$ 且 $A \times B$ 紧, 集合 $AB$ 也是紧的.

(6). (a) 设 $x, y \in G$ 满足 $xH \neq yH$. 由 (1), $x(G - H)y^{-1}$ 是开的, 所以其余集 $xHy^{-1}$ 是闭的. 又由假设 $xH \neq yH$, 单位元 $1$ 不在 $xHy^{-1}$ 中. 由 (2), 存在对称开集 $1 \in U$ 使得 $U^2 \subset G - xHy^{-1}$. 我们来证明 $UxH \cap UyH = \emptyset$, 这将证明结论, 因为 $UxH$ (相应地 $UyH$) 是 $xH$ (相应地 $yH$) 在 $G/H$ 中的开邻域. 若 $UxH \cap UyH \neq \emptyset$, 则存在 $u_1, u_2 \in U$ 与 $h_1, h_2 \in H$ 使得 $u_1xh_1 = u_2yh_2$. 但此时 $xh_1h_2^{-1}y^{-1} = u_1^{-1}u_2 \in xHy^{-1} \cap U^2$, 这是不可能的.

(b) 设 $xH \in G/H$. 若 $K$ 是 $x$ 在 $G$ 中的紧邻域, 则其在 $G/H$ 中的像是 $xH$ 在 $G/H$ 中的紧邻域.

(c) 若 $H$ 正规, 则 $G/H$ 是一个群. 我们来证明其乘法连续. 设 $x, y \in G$. $xyH$ 在 $G/H$ 中的任何开邻域都具有 $UxyH$ 的形式, 其中 $U$ 是 $xy$ 在 $G$ 中的开邻域. 由 $G$ 上乘法的连续性, 存在 $x$ 与 $y$ 在 $G$ 中的开邻域 $V$ 与 $W$ 使得 $VW \subset U$. 则 $VH$ 与 $WH$ 是 $xH$ 与 $yH$ 在 $G/H$ 中的开邻域, 且有 $(VH)(WH) \subset UH$. (记住, 由于 $H$ 正规, 对 $G$ 的任何子集 $A$ 有 $AH = HA$.) 我们来证明 $G/H$ 上的取逆连续. 设 $x \in G$. $x^{-1}H$ 在 $G/H$ 中的任何开邻域都具有 $UH$ 的形式, 其中 $U$ 是 $x^{-1}$ 在 $G$ 中的开邻域. 由 (1), 集合 $U^{-1}$ 是开的, 所以 $U^{-1}H$ 是 $xH$ 在 $G/H$ 中的开邻域, 且有 $(U^{-1}H)^{-1} = HU = UH$.
\end{proof}

\begin{remark}[Hausdorff 化]
\label{rem:I.1.8}
特别地, 若 $G$ 是拓扑群, 则 $G/\overline{\{1\}}$ 是 Hausdorff 拓扑群. 我们感兴趣的是 $G$ 在向量空间上的连续群作用, 因此可以把 $G$ 换成 $G/\overline{\{1\}}$ 来研究它们. 因此在以下的叙述中, 我们只考虑 Hausdorff 拓扑群 (除非另有说明).
\end{remark}

\begin{definition}[\index{紧群}紧群 (compact group)]
\label{def:I.1.9}
紧群 (compact group) (相应地局部紧群 (locally compact group)) 是 Hausdorff 且紧 (相应地局部紧) 的拓扑群.
\end{definition}

\begin{example}[紧群与局部紧群的例子]
\label{exa:I.1.10}
在例 \ref{exa:I.1.2} 的群中, 有限离散群与群 $S^1$, $\SU(n)$, $\SO(n)$ 是紧的. 其余所有群都是局部紧的. 若考虑一个无穷维巴拿赫空间的可逆有界线性自同态群, 则得到一个非局部紧的群 (见习题 \ref{ex:I.5.1.1}).
\end{example}

\textbf{平移算子 (Translation operators).} 设 $G$ 是一个群, $x \in G$, $f \colon G \to \CC$ 是一个函数. 我们定义两个函数 $L_xf, R_xf \colon G \to \CC$ 如下:
\[
L_xf(y) = f(x^{-1}y) \quad \text{与} \quad R_xf(y) = f(yx).
\]
我们这样选择约定是为了 $L_{xy} = L_x \circ L_y$ 且 $R_{xy} = R_x \circ R_y$. 注意, 若 $G$ 是拓扑群且 $f$ 连续, 则 $L_xf$ 与 $R_xf$ 也连续.

\textbf{函数空间 (Function spaces).} 设 $X$ 是一个拓扑空间. 若 $f \colon X \to \CC$ 是一个函数, 我们记
\[
\|f\|_{\infty} = \sup_{x \in X} |f(x)| \in [0, +\infty].
\]
我们还使用以下记号:
\begin{itemize}
\item $\mathcal{C}(X)$: 从 $X$ 到 $\CC$ 的连续函数 $f \colon X \to \CC$ 的集合;
\item $\mathcal{C}_b(X)$: 有界连续函数 $f \colon X \to \CC$ 的集合 (即满足 $\|f\|_{\infty} < +\infty$ 的 $\mathcal{C}(X)$ 的元素);
\item $\mathcal{C}_0(X)$: 在无穷远处消失 (vanish at infinity) 的连续函数 $X \to \CC$ 的集合 (即对每个 $\varepsilon > 0$, 存在 $X$ 的紧子集 $K$ 使得对所有 $x \notin K$ 有 $|f(x)| < \varepsilon$);
\item $\mathcal{C}_c(X)$: 具有紧支撑 (compact support) 的连续函数 $X \to \CC$ 的集合.
\end{itemize}
注意我们有 $\mathcal{C}(X) \supset \mathcal{C}_b(X) \supset \mathcal{C}_0(X) \supset \mathcal{C}_c(X)$, 且若 $X$ 紧, 则这些包含都是等式. 函数 $\|\cdot\|_{\infty}$ 是 $\mathcal{C}_b(X)$ 及其子空间上的范数, 且 $\mathcal{C}_b(X)$ 与 $\mathcal{C}_0(X)$ 对此范数完备 (但 $\mathcal{C}_c(X)$ 不完备, 除非 $X$ 紧).\footnote{参考文献 ?}

\begin{definition}[\index{一致连续}一致连续 (uniformly continuous)]
\label{def:I.1.11}
设 $G$ 是一个拓扑群. 若当 $x \to 1$ 时 $\|L_xf - f\|_{\infty} \to 0$ (相应地当 $x \to 1$ 时 $\|R_xf - f\|_{\infty} \to 0$), 则称函数 $f \colon G \to \CC$ 是左 (相应地右) 一致连续 (left (resp.\ right) uniformly continuous) 的.
\end{definition}

\begin{proposition}[紧支撑函数的均匀连续性]
\label{prop:I.1.12}
若 $f \in \mathcal{C}_c(G)$, 则 $f$ 既是左一致连续又是右一致连续的.
\end{proposition}

\begin{proof}
我们证明 $f$ 是右一致连续的 (左一致连续的证明类似). 设 $K$ 是 $f$ 的支撑. 设 $\varepsilon > 0$. 对每个 $x \in K$, 我们选取 $1$ 的一个邻域 $U_x$ 使得对所有 $y \in U_x$ 有 $|f(xy) - f(x)| < \frac{\varepsilon}{2}$; 由命题 \ref{prop:I.1.7}, 可以找到 $1$ 的对称开邻域 $V_x$ 使得 $V_x^2 \subset U_x$. 我们有 $K \subset \bigcup_{x \in K} xV_x$. 由于 $K$ 紧, 可以找到 $x_1, \ldots, x_n \in K$ 使得 $K \subset \bigcup_{i=1}^n x_iV_{x_i}$. 令 $V = \bigcup_{i=1}^n V_{x_i}$, 这是 $1$ 的一个对称开邻域.

我们断言: 若 $y \in V$, 则 $\|R_yf - f\|_{\infty} < \varepsilon$. 事实上, 设 $y \in V$, $x \in G$. 先假定 $x \in K$. 则存在 $i \in \{1, \ldots, n\}$ 使得 $x \in x_iV_{x_i}$. 于是 $xy \in x_iV_{x_i}V_{x_i} \subset x_iU_{x_i}$, 从而
\[
|f(xy) - f(x)| \leq |f(xy) - f(x_i)| + |f(x_i) - f(x)| < \frac{\varepsilon}{2} + \frac{\varepsilon}{2} = \varepsilon.
\]
再假定 $xy \in K$. 则存在 $i \in \{1, \ldots, n\}$ 使得 $xy \in x_iV_{x_i}$, 且有 $x = xyy^{-1} \in x_iV_{x_i}V_{x_i} \subset x_iU_{x_i}$. 因此
\[
|f(xy) - f(x)| \leq |f(xy) - f(x_i)| + |f(x_i) - f(x)| < \frac{\varepsilon}{2} + \frac{\varepsilon}{2} = \varepsilon.
\]
最后, 若 $x, xy \notin K$, 则 $f(x) = f(xy) = 0$, 当然有 $|f(xy) - f(x)| < \varepsilon$.
\end{proof}

\begin{remark}[一致连续与连续性的等价刻画]
\label{rem:I.1.13}
我们在 $\mathcal{C}_b(G)$ 上赋予由 $\|\cdot\|_{\infty}$ 给出的拓扑. 则函数 $f \in \mathcal{C}_b(G)$ 是左 (相应地右) 一致连续的, 当且仅当映射 $G \to \mathcal{C}_b(G)$, $x \longmapsto L_xf$ (相应地 $x \longmapsto R_xf$) 在 $G$ 的单位元处连续.

利用 $L_{xy} = L_x \circ L_y$ 与 $R_{xy} = R_x \circ R_y$, 以及算子 $L_x$ 与 $R_x$ 保持 $\mathcal{C}_c(G)$, 我们看到上述命题蕴含: 若 $f \in \mathcal{C}_c(G)$, 则把 $x \in G$ 分别映到 $L_xf$ 与 $R_xf$ 的两个映射 $G \to \mathcal{C}_c(G)$ 都是连续的.
\end{remark}

% ============ I.2 \index{Haar 测度}Haar 测度 (原书 pp.11-18) ============
\section{Haar 测度}

\begin{definition}[\index{Borel 测度}Borel 测度与\index{正则 Borel 测度}正则 Borel 测度]
\label{def:I.2.1}
设 $X$ 是一个拓扑空间.
\begin{enumerate}
\item[(1)] $X$ 上 Borel 集的 $\sigma$-代数 (the $\sigma$-algebra of Borel sets) 是由 $X$ 的开子集生成的 $\sigma$-代数. $X$ 上的 Borel 测度 (Borel measure) 是该 $\sigma$-代数上的测度.
\item[(2)] $X$ 上的正则 Borel 测度 (regular Borel measure) 是 $X$ 的 Borel 集的 $\sigma$-代数上满足下列性质的测度 $\mu$:
  \begin{enumerate}
  \item[(a)] 对 $X$ 的每个紧子集 $K$, $\mu(K) < +\infty$;
  \item[(b)] $\mu$ 是外正则 (outer regular) 的: 对 $X$ 的每个 Borel 子集 $E$, 有 $\mu(E) = \inf\{\mu(U) \mid U \supset E \text{ 开}\}$;
  \item[(c)] $\mu$ 是内正则 (inner regular) 的: 对每个 $E \subset X$, 只要 $E$ 是有限测度的 Borel 集或是开集, 就有 $\mu(E) = \sup\{\mu(K) \mid K \subset E \text{ 紧}\}$.
  \end{enumerate}
\end{enumerate}
\end{definition}

\begin{notation}[非负紧支撑函数]
\label{not:I.2.2}
我们记 $\mathcal{C}_c^+(X)$ 为所有非零的 $f \in \mathcal{C}_c(X)$ 且 $f(X) \subset \R_{\geq 0}$ 组成的集合.
\end{notation}

\begin{theorem}[Riesz 表示定理 (Riesz representation theorem)]
\label{thm:I.2.3}
设 $X$ 是局部紧 Hausdorff 空间, 设 $\Lambda \colon \mathcal{C}_c(X) \to \CC$ 是满足对每个 $f \in \mathcal{C}_c^+(X)$ 都有 $\Lambda(f) \geq 0$ 的线性泛函.\footnote{这样的线性泛函称为正的 (positive).} 则存在 $X$ 上唯一的正则 Borel 测度 $\mu$ 使得对每个 $f \in \mathcal{C}_c(X)$ 有
\[
\Lambda(f) = \int_X f \,\dd\mu.
\]
\end{theorem}

\begin{definition}[左/右 Haar 测度]
\label{def:I.2.4}
设 $G$ 是局部紧群. $G$ 上的左 (相应地右) Haar 测度 (left (resp.\ right) Haar measure) 是非零正则 Borel 测度 $\mu$, 使得对 $G$ 的每个 Borel 集 $E$ 和每个 $x \in G$, 有 $\mu(xE) = \mu(E)$ (相应地 $\mu(Ex) = \mu(E)$).
\end{definition}

\begin{example}[Haar 测度的例子]
\label{exa:I.2.5}
\begin{enumerate}
\item[(1)] 若 $G$ 是离散群, 则计数测度 (counting measure) 是 $G$ 上的左和右 Haar 测度.
\item[(2)] Lebesgue 测度是 $\R$ 的加法群上的左和右 Haar 测度.
\end{enumerate}
\end{example}

\begin{proposition}[Haar 测度的刻画]
\label{prop:I.2.6}
设 $G$ 是局部紧群, $\mu$ 是 $G$ 上的正则 Borel 测度.
\begin{enumerate}
\item[(1)] 设 $\tilde{\mu}$ 是由 $\tilde{\mu}(E) = \mu(E^{-1})$ 定义的 Borel 测度. 则 $\mu$ 是左 Haar 测度当且仅当 $\tilde{\mu}$ 是右 Haar 测度.
\item[(2)] $\mu$ 是 $G$ 上的左 Haar 测度当且仅当: 对每个 $f \in \mathcal{C}_c(G)$, 对每个 $y \in G$, 有 $\int_G L_yf \,\dd\mu = \int_G f \,\dd\mu$.
\item[(3)] 若 $\mu$ 是 $G$ 上的左 Haar 测度, 则对 $G$ 的每个非空开子集有 $\mu(U) > 0$, 且对每个 $f \in \mathcal{C}_c^+(G)$ 有 $\int_G f \,\dd\mu > 0$.
\end{enumerate}
\end{proposition}

\begin{proof}
(1). 首先注意 $\tilde{\mu}$ 是 $G$ 上的正则 Borel 测度, 因为 $x \longmapsto x^{-1}$ 是 $G$ 到自身的同胚.

若 $E \subset G$ 是 Borel 集且 $x \in G$, 则 $\tilde{\mu}(Ex) = \mu(x^{-1}E^{-1})$. 这就蕴含了结论.

(2). 设 $x \in G$, 令 $\mu_x$ 是由 $\mu_x(E) = \mu(xE)$ 定义的 Borel 测度. (这确实是 $G$ 上的正则 Borel 测度, 因为 $y \longmapsto xy$ 是 $G$ 到自身的同胚.) 则对每个可测函数 $f \colon G \to \CC$, 有 $\int_G f \,\dd\mu_x = \int_G L_xf \,\dd\mu$. (这对 Borel 子集的特征函数是显然的, 一般情形通过对 $f$ 用特征函数的线性组合逼近得到.)

一方面, 测度 $\mu$ 是左 Haar 测度当且仅当对每个 $x \in G$ 有 $\mu = \mu_x$. 另一方面, 由 Riesz 表示定理中的唯一性 (以及上一段), 对 $x \in G$, $\mu = \mu_x$ 当且仅当对每个 $f \in \mathcal{C}_c(G)$ 有 $\int_G f \,\dd\mu = \int_G L_xf \,\dd\mu$. 结论由此得出.

(3). 假设存在 $G$ 的非空开子集 $U$ 使得 $\mu(U) = 0$. 则对每个 $x \in G$ 有 $\mu(xU) = 0$, 所以我们可以假设 $1 \in U$. 设 $K$ 是 $G$ 的紧子集. 则 $K \subset \bigcup_{x \in K} xU$, 所以存在 $x_1, \ldots, x_n \in K$ 使得 $K \subset \bigcup_{i=1}^n x_iU$. 由于对每个 $i$ 有 $\mu(x_iU) = 0$, 这蕴含 $\mu(K) = 0$. 但由 $\mu$ 的内正则性, 我们得到 $\mu(G) = 0$, 这与 $\mu$ 非零矛盾.

设 $f \in \mathcal{C}_c^+(G)$. 则 $U := \{x \in G \mid f(x) > \frac{1}{2}\|f\|_{\infty}\}$ 是 $G$ 的非空开子集, 所以 $\mu(U) > 0$. 但 $f \geq \frac{1}{2}\|f\|_{\infty} \mathbb{1}_U$, 因此 $\int_G f \,\dd\mu \geq \frac{1}{2}\|f\|_{\infty}\mu(U) > 0$.
\end{proof}

\begin{theorem}[左 Haar 测度的存在性与唯一性]
\label{thm:I.2.7}
设 $G$ 是局部紧群. 则:
\begin{enumerate}
\item[(1)] $G$ 上存在左 Haar 测度.
\item[(2)] 若 $\mu_1$ 和 $\mu_2$ 是 $G$ 上的两个左 Haar 测度, 则存在 $c \in \R_{>0}$ 使得 $\mu_2 = c\mu_1$.
\end{enumerate}
\end{theorem}

由命题 \ref{prop:I.2.6}, 此定理蕴含右 Haar 测度的类似结论.

\begin{proof}
我们先证明存在性. 其想法与 $\R$ 上 Lebesgue 测度的构造非常相似. 假设 $c > 0$, 且 $\phi \in \mathcal{C}_c^+(\R)$ 被 $1$ 控制且非常接近区间 $[0,c]$ 的特征函数. 若 $f \in \mathcal{C}_c(\R)$ 在长度为 $c$ 的区间上变化不太快, 则我们可以用 $\phi$ 的左平移的线性组合逼近 $f$: $f \simeq \sum c_iL_{x_i}\phi$, 然后 $\int f \,\dd\mu \simeq \sum c_i \int L_{x_i}\phi \,\dd\mu$. 当 $c \to 0$ 时, 我们将能逼近每个 $f \in \mathcal{C}_c(\R)$ (因为我们知道这些函数是一致连续的), 并且能通过取极限来定义 $\int f \,\dd\mu$. 在一般的局部紧群上, 我们把区间换成 $1$ 的越来越小的紧邻域.

下面是严格的证明. 设 $f, \phi \in \mathcal{C}_c^+(G)$. 则 $U := \{x \in G \mid \phi(x) > \frac{1}{2}\|\phi\|_{\infty}\}$ 是 $G$ 的非空开子集, 且 $\phi \geq \frac{1}{2}\|\phi\|_{\infty}\mathbb{1}_U$. 由于 $f$ 的支撑紧, 它可被 $U$ 的有限多个平移覆盖, 所以存在 $x_1, \ldots, x_n \in G$ 和 $c_1, \ldots, c_n \in \R_{\geq 0}$ 使得 $f \leq \sum_{i=1}^n c_iL_{x_i}\phi$. 因此, 若定义 $(f : \phi)$ 为所有有限和 $\sum_{i=1}^n c_i$ (其中 $c_1, \ldots, c_n \in \R_{\geq 0}$ 且存在 $x_1, \ldots, x_n \in G$ 使得 $f \leq \sum_{i=1}^n c_iL_{x_i}\phi$) 的下确界, 则 $(f : \phi) < +\infty$. 我们断言:
\begin{align}
\label{eq:I.2.1}
(f : \phi) &= (L_xf : \phi) \quad \forall x \in G, \\
\label{eq:I.2.2}
(f_1 + f_2 : \phi) &\leq (f_1 : \phi) + (f_2 : \phi), \\
\label{eq:I.2.3}
(cf : \phi) &= c(f : \phi) \quad \forall c \geq 0, \\
\label{eq:I.2.4}
(f_1 : \phi) &\leq (f_2 : \phi) \quad \text{若 } f_1 \leq f_2, \\
\label{eq:I.2.5}
(f : \phi) &\geq \frac{\|f\|_{\infty}}{\|\phi\|_{\infty}}, \\
\label{eq:I.2.6}
(f : \phi) &\leq (f : \psi)(\psi : \phi) \quad \forall \psi \in \mathcal{C}_c^+(G) - \{0\}.
\end{align}

前四条性质是容易的. 对第五条性质, 注意若 $f \leq \sum_{i=1}^n c_iL_{x_i}\phi$, 则
\[
\|f\|_{\infty} \leq \sum_{i=1}^n c_i\|L_{x_i}\phi\|_{\infty} = \sum_{i=1}^n c_i\|\phi\|_{\infty}.
\]

最后, 最后一条性质是以下事实的推论: 设 $\psi \in \mathcal{C}_c^+(G)$. 若 $f \leq \sum_{i=1}^n c_iL_{x_i}\psi$ 且 $\psi \leq \sum_{j=1}^m d_jL_{y_j}\phi$, 则 $f \leq \sum_{i=1}^n \sum_{j=1}^m c_id_jL_{x_iy_j}\phi$.

现在我们固定 $f_0 \in \mathcal{C}_c^+(G)$. 由 \eqref{eq:I.2.5}, 我们知道 $(f_0 : \phi) > 0$. 我们定义 $I_{\phi} \colon \mathcal{C}_c^+(G) \to \R_{\geq 0}$ 为
\[
I_{\phi}(f) = \frac{(f : \phi)}{(f_0 : \phi)}.
\]
由 \eqref{eq:I.2.1}-\eqref{eq:I.2.4}, 我们有
\begin{align*}
I_{\phi}(f) &= I_{\phi}(L_xf) \quad \forall x \in G, \\
I_{\phi}(f_1 + f_2) &\leq I_{\phi}(f_1) + I_{\phi}(f_2), \\
I_{\phi}(cf) &= cI_{\phi}(f) \quad \forall c \geq 0, \\
I_{\phi}(f_1) &\leq I_{\phi}(f_2) \quad \text{若 } f_1 \leq f_2.
\end{align*}

若第二个不等式是等式 (即若 $I_{\phi}$ 可加), 我们就能把 $I_{\phi}$ 延拓成 $\mathcal{C}_c(G)$ 上的\index{正线性泛函}正线性泛函并应用 Riesz 表示定理. 这并不完全成立, 但我们有以下结果:

\textbf{断言 (Claim):} 对所有 $f_1, f_2 \in \mathcal{C}_c^+(G)$ 和 $\varepsilon > 0$, 存在 $1$ 在 $G$ 中的邻域 $V$ 使得只要 $\operatorname{supp}(\phi) \subset V$, 就有 $I_{\phi}(f_1) + I_{\phi}(f_2) \leq I_{\phi}(f_1 + f_2) + \varepsilon$.

我们首先证明这个断言. 选取函数 $g \in \mathcal{C}_c^+(G)$ 使得对每个 $x \in \operatorname{supp}(f_1 + f_2)$ 有 $g(x) = 1$, 并设 $\delta$ 是一个正实数. 令 $h = f_1 + f_2 + \delta g$. 我们定义函数 $h_1, h_2 \colon G \to \R_{\geq 0}$ 为
\[
h_i(x) = \begin{cases} \frac{f_i(x)}{h(x)} & \text{若 } f_i(x) \neq 0, \\ 0 & \text{若 } f_i(x) = 0. \end{cases}
\]
注意 $h_i$ 等于 $\frac{f_i}{h}$, 因此在开子集 $\{x \in G \mid h(x) \neq 0\}$ 上连续. 由于 $G$ 是这个开子集与开子集 $G - \operatorname{supp}(f_i)$ (在其上 $h_i$ 也连续) 的并, 这证明 $h_i$ 连续, 因此 $h_i \in \mathcal{C}_c^+(G)$. 又注意 $f_i = h_ih$.

由命题 \ref{prop:I.1.12}, 存在 $1$ 的邻域 $V$ 使得对 $i \in \{1, 2\}$ 及满足 $y^{-1}x \in V$ 的 $x, y \in G$, 有 $|h_i(x) - h_i(y)| < \delta$. 设 $\phi \in \mathcal{C}_c^+(G)$ 满足 $\operatorname{supp}(\phi) \subset V$. 若 $c_1, \ldots, c_n \in \R_{\geq 0}$ 与 $x_1, \ldots, x_n \in G$ 满足 $h \leq \sum_{j=1}^n c_jL_{x_j}\phi$, 则对每个 $x \in G$ 和 $i \in \{1, 2\}$,
\[
f_i(x) = h(x)h_i(x) \leq \sum_{j=1}^n c_j\phi(x_j^{-1}x)h_i(x) \leq \sum_{j=1}^n c_j\phi(x_j^{-1}x)(h_i(x_j) + \delta),
\]
因为 $\phi(x_j^{-1}x) = 0$ 除非 $x_j^{-1}x \in V$. 因此
\[
(f_1 : \phi) + (f_2 : \phi) \leq \sum_{j=1}^n c_j(h_1(x_j) + h_2(x_j) + 2\delta).
\]
由于 $h_1 + h_2 \leq 1$, 我们得到
\[
(f_1 : \phi) + (f_2 : \phi) \leq (1 + 2\delta)\sum_{j=1}^n c_j,
\]
从而, 对族 $(c_1, \ldots, c_n)$ 取下确界并除以 $(f_0 : \phi)$, 得到
\[
I_{\phi}(f_1) + I_{\phi}(f_2) \leq (1 + 2\delta)I_{\phi}(h) \leq (1 + 2\delta)(I_{\phi}(f_1 + f_2) + \delta I_{\phi}(g)).
\]
当 $\delta$ 趋于 $0$ 时, 右端趋于 $I_{\phi}(f_1 + f_2)$, 所以取足够小的 $\delta$ 就得到所要的不等式. 这就完成了断言的证明.

我们回到左 Haar 测度的构造. 对每个 $f \in \mathcal{C}_c^+(G)$, 设 $X_f = [(f_0 : f)^{-1}, (f : f_0)] \subset \R$. 令 $X = \prod_{f \in \mathcal{C}_c^+(G)} X_f$, 赋予乘积拓扑. 由 Tychonoff 定理,\footnote{参考文献 ?} $X$ 是紧 Hausdorff 空间. 它正是函数 $I \colon \mathcal{C}_c^+(G) \to \R$ (满足对每个 $f$ 有 $I(f) \in X_f$) 的空间 (赋予逐点收敛拓扑). 又由 \eqref{eq:I.2.6}, 对每个 $\phi \in \mathcal{C}_c^+(G)$ 有 $I_{\phi} \in X$. 对 $1$ 在 $G$ 中的每个邻域 $V$, 设 $K(V)$ 为 $\{I_{\phi} \mid \operatorname{supp}(\phi) \subset V\}$ 在 $X$ 中的闭包. 对每个 $V$ 有 $K(V) \neq \emptyset$, 所以对 $1$ 在 $G$ 中的任何有限族 $V_1, \ldots, V_n$ 个邻域, $K(V_1) \cap \cdots \cap K(V_n) \supset K(\bigcap_{i=1}^n V_i) \neq \emptyset$. 由于 $X$ 紧, 这蕴含所有集合 $K(V)$ 的交非空. 我们选取这个交中的一个元素 $I$.

我们来证明 $I$ 在左平移下不变、可加且是一次齐次的. (也就是说, 它和 $I_{\phi}$ 有相同的性质, 但它是可加的面不是次可加的.) 设 $f_1, f_2 \in \mathcal{C}_c^+(G)$, $c \in \R_{\geq 0}$, $x \in G$, $\varepsilon > 0$. 选取 $1$ 在 $G$ 中的邻域 $V$ 使得只要 $\operatorname{supp}(\phi) \subset V$ 就有 $I_{\phi}(f_1) + I_{\phi}(f_2) \leq I_{\phi}(f_1 + f_2) + \varepsilon$; 这样的 $V$ 由断言保证存在. 由 $I$ 的定义, 它属于 $\{I_{\phi} \mid \operatorname{supp}(\phi) \subset V\}$ 的闭包, 这意味着存在 $\phi \in \mathcal{C}_c^+(G)$ 使得 $\operatorname{supp}(\phi) \subset V$ 且对 $g \in \{f_1, f_2, f_1 + f_2\}$, $y \in \{1, x\}$, $a \in \{1, c\}$ 有 $|I(aL_yg) - I_{\phi}(aL_yg)| < \varepsilon$. 于是我们得到:
\begin{align*}
|I(L_xf_1) - I(f_1)| &\leq |I(L_xf_1) - I_{\phi}(L_xf_1)| + |I_{\phi}(L_xf_1) - I_{\phi}(f_1)| + |I_{\phi}(f_1) - I(f_1)| < 2\varepsilon, \\
|I(cf_1) - cI(f_1)| &\leq |I(cf_1) - I_{\phi}(cf_1)| + |I_{\phi}(cf_1) - cI_{\phi}(f_1)| + |cI_{\phi}(f_1) - cI(f_1)| < \varepsilon(1 + c),
\end{align*}
以及
\begin{multline*}
|I(f_1 + f_2) - I(f_1) - I(f_2)| \leq |I_{\phi}(f_1 + f_2) - I_{\phi}(f_1) - I_{\phi}(f_2)| \\
+ |I(f_1 + f_2) - I_{\phi}(f_1 + f_2)| + |I(f_1) - I_{\phi}(f_1)| + |I(f_2) - I_{\phi}(f_2)| < 4\varepsilon.
\end{multline*}
由于 $\varepsilon$ 任意, 这蕴含 $I(L_xf_1) = I(f_1)$, $I(cf_1) = cI(f_1)$ 且 $I(f_1 + f_2) = I(f_1) + I(f_2)$.

现在我们把它延拓成线性泛函 $\mathcal{C}_c(G) \to \CC$, 仍记作 $I$. 设 $f \in \mathcal{C}_c(G)$. 则我们可以把 $f$ 写成 $f = (f_1 - f_2) + i(g_1 - g_2)$, 其中 $f_1, f_2, g_1, g_2 \in \mathcal{C}_c^+(G) \cup \{0\}$ (例如取 $f_1 = \max(0, \operatorname{Re}(f))$, $f_2 = \max(0, -\operatorname{Re}(f))$, $g_1 = \max(0, \operatorname{Im}(f))$, $g_2 = \max(0, \operatorname{Im}(f))$). 我们令 $I(f) = I(f_1) - I(f_2) + i(I(g_1) - I(g_2))$ (约定 $I(0) = 0$). 若 $f = (F_1 - F_2) + i(G_1 - G_2)$, 其中 $F_1, F_2, G_1, G_2 \in \mathcal{C}_c^+(G) \cup \{0\}$, 则 $F_1 + f_2 = F_2 + f_1$ 且 $G_1 + g_2 = G_2 + g_1$, 所以 $I(f)$ 得到相同的结果. 容易验证 $I$ 是从 $\mathcal{C}_c(G)$ 到 $\CC$ 的线性泛函, 且由构造它是正的. 由 Riesz 表示定理, 存在 $G$ 上的正则 Borel 测度 $\mu$ 使得 $I(f) = \int_G f \,\dd\mu$. 由命题 \ref{prop:I.2.6}, 这个测度是左 Haar 测度.

我们现在证明定理的第二部分 (左 Haar 测度在相差一个常数意义下唯一). 设 $\mu_1$, $\mu_2$ 是 $G$ 上的两个左 Haar 测度. 由 Riesz 表示定理中的唯一性 (以及 $\mathcal{C}_c^+(G)$ 生成 $\mathcal{C}_c(G)$ 这一事实), 只需找到一个正实数 $c$ 使得对每个 $f \in \mathcal{C}_c^+(G)$ 有 $\int_G f \,\dd\mu_1 = c\int_G f \,\dd\mu_2$. 由命题 \ref{prop:I.2.6}, 对每个 $f \in \mathcal{C}_c^+(G)$ 有 $\int_G f \,\dd\mu_2 > 0$. 所以只需证明: 若 $f, g \in \mathcal{C}_c^+(G)$, 则
\begin{equation}
\label{eq:I.2.star}
\frac{\int_G f \,\dd\mu_1}{\int_G f \,\dd\mu_2} = \frac{\int_G g \,\dd\mu_1}{\int_G g \,\dd\mu_2} \quad (*).
\end{equation}

设 $f, g \in \mathcal{C}_c^+(G)$. 设 $V_0$ 是 $1$ 的对称紧邻域, 并令
\[
A = (\operatorname{supp}(f))V_0 \cup V_0(\operatorname{supp}(f))
\]
及
\[
B = (\operatorname{supp}(g))V_0 \cup V_0(\operatorname{supp}(g)).
\]
由命题 \ref{prop:I.1.7}, $A$ 和 $B$ 紧. 若 $y \in V_0$, 则函数 $x \longmapsto f(xy) - f(yx)$ 与 $x \longmapsto g(xy) - g(yx)$ 分别支撑在 $A$ 和 $B$ 上.

设 $\varepsilon > 0$. 由命题 \ref{prop:I.1.12}, 存在 $1$ 的对称邻域 $V \subset V_0$ 使得对每个 $x \in G$ 和每个 $y \in V$, 有 $|f(xy) - f(yx)| < \varepsilon$ 且 $|g(xy) - g(yx)| < \varepsilon$. 设 $h \in \mathcal{C}_c^+(G)$ 满足 $\operatorname{supp}(h) \subset V$ 且对每个 $x \in G$ 有 $h(x) = h(x^{-1})$. 则
\begin{align*}
\Bigl(\int_G h \,\dd\mu_2\Bigr)\Bigl(\int_G f \,\dd\mu_1\Bigr) &= \int_{G \times G} h(y)f(x) \,\dd\mu_1(x)\,\dd\mu_2(y) \\
&= \int_{G \times G} h(y)f(yx) \,\dd\mu_1(x)\,\dd\mu_2(y).
\end{align*}
(我们用了 $\mu_1$ 的左不变性. 又, 我们可以应用 Fubini 定理, 因为所有函数都支撑在紧集上, 而紧集有有限测度.) 类似地,
\begin{align*}
\Bigl(\int_G h \,\dd\mu_1\Bigr)\Bigl(\int_G f \,\dd\mu_2\Bigr) &= \int_{G \times G} h(x)f(y) \,\dd\mu_1(x)\,\dd\mu_2(y) \\
&= \int_{G \times G} h(y^{-1}x)f(y) \,\dd\mu_1(x)\,\dd\mu_2(y) \\
&= \int_{G \times G} h(x^{-1}y)f(y) \,\dd\mu_1(x)\,\dd\mu_2(y) \\
&= \int_{G \times G} h(y)f(xy) \,\dd\mu_1(x)\,\dd\mu_2(y).
\end{align*}

因此
\begin{align*}
\Bigl(\int_G h \,\dd\mu_1\Bigr)\Bigl(\int_G f \,\dd\mu_2\Bigr) - \Bigl(\int_G h \,\dd\mu_2\Bigr)\Bigl(\int_G f \,\dd\mu_1\Bigr) &= \int_{G \times G} h(y)(f(xy) - f(yx)) \,\dd\mu_1(x)\,\dd\mu_2(y) \\
&\leq \varepsilon\mu_1(A)\int_G h \,\dd\mu_2,
\end{align*}
因为 $\operatorname{supp}(h) \subset V$. 除以 $\bigl(\int_G f \,\dd\mu_2\bigr)\bigl(\int_G h \,\dd\mu_2\bigr)$, 得到
\[
\Bigl(\int_G h \,\dd\mu_1\Bigr)\Bigl(\int_G h \,\dd\mu_2\Bigr)^{-1} - \Bigl(\int_G f \,\dd\mu_1\Bigr)\Bigl(\int_G f \,\dd\mu_2\Bigr)^{-1} \leq \varepsilon\mu_1(A)\Bigl(\int_G f \,\dd\mu_2\Bigr)^{-1}.
\]
类似地,
\[
\Bigl(\int_G h \,\dd\mu_1\Bigr)\Bigl(\int_G h \,\dd\mu_2\Bigr)^{-1} - \Bigl(\int_G g \,\dd\mu_1\Bigr)\Bigl(\int_G g \,\dd\mu_2\Bigr)^{-1} \leq \varepsilon\mu_1(B)\Bigl(\int_G g \,\dd\mu_2\Bigr)^{-1}.
\]
两式相加给出
\[
\left|\frac{\int_G f \,\dd\mu_1}{\int_G f \,\dd\mu_2} - \frac{\int_G g \,\dd\mu_1}{\int_G g \,\dd\mu_2}\right| \leq \varepsilon\left(\frac{\mu_1(A)}{\int_G f \,\dd\mu_2} + \frac{\mu_1(B)}{\int_G g \,\dd\mu_2}\right).
\]
由于 $\varepsilon$ 任意, 这给出所要的等式 $(*)$.
\end{proof}

我们现在来比较左、右 Haar 测度.

\begin{proposition}[\index{模函数}模函数 (modular function)]
\label{prop:I.2.8}
设 $G$ 是局部紧群, $x \in G$. 则存在 $\Delta(x) \in \R_{>0}$ 使得对 $G$ 的每个左 Haar 测度 $\mu$, 有 $\mu(Ex) = \Delta(x)\mu(E)$. 而且, $\Delta \colon G \to \R_{>0}$ 是连续群同态 (其中 $\R_{>0}$ 上的群结构由乘法给出), 且对 $G$ 的每个左 Haar 测度 $\mu$, 每个 $x \in G$ 和每个 $f \in L^1(G, \mu)$, 有
\[
\int_G R_xf \,\dd\mu = \Delta(x^{-1})\int_G f \,\dd\mu.
\]
\end{proposition}

\begin{proof}
设 $x \in G$, $\mu$ 是 $G$ 上的左 Haar 测度. 则由 $\mu_x(E) = \mu(Ex)$ 定义的测度 $\mu_x$ 也是 $G$ 上的左 Haar 测度, 所以由定理 \ref{thm:I.2.7} 中的唯一性, 存在 $\Delta(x) \in \R_{>0}$ 使得 $\mu_x = \Delta(x)\mu$, 即对 $G$ 的每个 Borel 子集 $E$ 有 $\mu(Ex) = \Delta(x)\mu(E)$. 假设 $\lambda$ 是 $G$ 上的另一个左 Haar 测度. 则再次由定理 \ref{thm:I.2.7}, 存在 $c > 0$ 使得 $\lambda = c\mu$, 于是对 $G$ 的每个 Borel 子集 $E$,
\[
\lambda(Ex) = c\mu(Ex) = c\Delta(x)\mu(E) = \Delta(x)\lambda(E).
\]
这就证明了第一个结论.

我们证明 $\Delta$ 是群同态. 设 $x, y \in G$, $E$ 是 $G$ 的满足 $\mu(E) \neq 0$ 的 Borel 子集. 则
\[
\Delta(xy)\mu(E) = \mu(Exy) = \Delta(y)\mu(Ex) = \Delta(y)\Delta(x)\mu(E),
\]
因此 $\Delta(xy) = \Delta(x)\Delta(y)$.

我们证明最后一个结论. 若 $E$ 是 $G$ 的 Borel 子集且 $x \in G$, 则 $R_x\mathbb{1}_E = \mathbb{1}_{Ex^{-1}}$, 所以由 $\Delta$ 的定义,
\[
\int_G R_x\mathbb{1}_E \,\dd\mu = \mu(Ex^{-1}) = \Delta(x^{-1})\mu(E) = \Delta(x)^{-1}\int_G \chi_E \,\dd\mu.
\]
这就对 $f = \chi_E$ 证明了结论. 一般情形通过用函数 $\mathbb{1}_E$ 的线性组合逼近 $f$ 得到.

最后, 我们证明 $\Delta$ 连续. 设 $f \in \mathcal{C}_c^+(G)$. 我们知道函数 $G \to \mathcal{C}_c(G)$, $x \longmapsto R_{x^{-1}}f$ 连续 (见注记 \ref{rem:I.1.13}), 所以函数 $G \to \CC$, $x \longmapsto \int_G R_{x^{-1}}f \,\dd\mu$ 也连续. 但刚才已见 $\int_G R_{x^{-1}}f \,\dd\mu = \Delta(x)\int_G f \,\dd\mu$, 且由命题 \ref{prop:I.2.6} 知 $\int_G f \,\dd\mu > 0$. 因此 $\Delta$ 连续.
\end{proof}

\begin{definition}[模函数与\index{幺模群}幺模群 (unimodular group)]
\label{def:I.2.9}
前一命题中的函数 $\Delta$ 称为 $G$ 的模函数 (modular function). 若 $\Delta = 1$ (即 $G$ 的某个 (或任何) 左 Haar 测度也是右 Haar 测度), 则称群 $G$ 是幺模 (unimodular) 的.
\end{definition}

\begin{remark}[推广的模因子]
\label{rem:I.2.10}
假设 $\alpha \colon G \to G$ 是同胚, 使得对每个 $x \in G$, 存在 $\beta(x) \in G$ 满足: 对每个 $y \in G$, $\alpha(xy) = \beta(x)\alpha(y)$. (例如, $\alpha$ 可以是 $G$ 中固定元素的右平移, 或具有连续逆的连续群同构.) 则我们可以推广命题 \ref{prop:I.2.8} 的构造, 得到 $\Delta(\alpha) \in \R_{>0}$ 满足: 对每个 $f \in \mathcal{C}_c(G)$, 对 $G$ 的每个左 Haar 测度 $\mu$,
\[
\Delta(\alpha)\int_G f(\alpha(x)) \,\dd\mu(x) = \int_G f(x) \,\dd\mu(x)
\]
(或等价地, 对 $G$ 的每个 Borel 子集 $E$, $\mu(\alpha(E)) = \Delta(\alpha)\mu(E)$). 而且, 若 $\beta \colon G \to G$ 满足与 $\alpha$ 相同的条件, 则 $\alpha \circ \beta$ 也满足, 且 $\Delta(\alpha \circ \beta) = \Delta(\alpha)\Delta(\beta)$.
\end{remark}

\begin{example}[幺模群的例子]
\label{exa:I.2.11}
\begin{enumerate}
\item[(1)] 任何紧群都是幺模的. 事实上, 若 $G$ 紧, 则 $\Delta(G)$ 是 $\R_{>0}$ 的紧子群, 但 $\R_{>0}$ 唯一的紧子群是 $\{1\}$. 特别地, 紧群 $G$ 有唯一的满足 $\mu(G) = 1$ 的左、右 Haar 测度 $\mu$; 我们称这个测度为 $G$ 的\index{规范化 Haar 测度}规范化 Haar 测度 (normalized Haar measure).
\item[(2)] 任何离散群都是幺模的. 事实上, 我们有 $G$ 上的一个同时是右 Haar 测度的左 Haar 测度: 计数测度.
\item[(3)] 若 $G$ 交换, 则 $G$ 上的左平移与右平移相等, 所以 $G$ 是幺模的.
\item[(4)] 群 $\GL_n(\R)$ 和 $\GL_n(\CC)$ 是幺模的. (对 $\GL_n(\R)$ 在习题 \ref{ex:I.5.3.2}(c) 中证明, 同样的证明对 $\GL_n(\CC)$ 也成立.)
\item[(5)] $M_2(\R)$ 中可逆上三角矩阵的群不是幺模的 (见习题 \ref{ex:I.5.3.2}(d)). 事实上, 其模函数为
  \[
  \Delta \colon \begin{pmatrix} a & b \\ 0 & c \end{pmatrix} \longmapsto |ac^{-1}|.
  \]
\item[(6)] 回忆换位子群 (commutator subgroup) $[G, G]$ 是由所有 $xyx^{-1}y^{-1}$ ($x, y \in G$) 生成的子群. 它是 $G$ 的正规子群, 且从 $G$ 到交换群的任何群同态在 $[G, G]$ 上都平凡. 特别地, 模函数 $\Delta$ 在 $[G, G]$ 上平凡, 所以若 $G = [G, G]$, 则 $G$ 是幺模的. 更一般地, 利用第一个例子, 我们看到若商群 $G/[G, G]$ 紧, 则 $G$ 是幺模的.
\end{enumerate}
\end{example}

\begin{proposition}[$f(x^{-1})$ 的积分公式]
\label{prop:I.2.12}
设 $G$ 是局部紧群, $\mu$ 是 $G$ 上的左 Haar 测度. 我们用 $\nu(E) = \mu(E^{-1})$ 定义 $G$ 上的右 Haar 测度 $\nu$ (见命题 \ref{prop:I.2.6}).
则对每个 $f \in \mathcal{C}_c(G)$, 有
\[
\int_G f(x^{-1}) \,\dd\mu_G(x) = \int_G f(x) \,\dd\nu(x) = \int_G \Delta(x^{-1})f(x) \,\dd\mu(x).
\]
\end{proposition}

我们也将此性质记为 $d\nu(x) = \Delta(x^{-1})d\mu(x)$, 或 $d\mu(x^{-1}) = \Delta(x^{-1})d\mu(x)$.

% ============ I.3 表示 (原书 pp.19-30) ============
\section{表示}
在本节中, $G$ 是一个拓扑群.

\subsection{连续表示}

\begin{definition}[$\Hom(V,W)$ 与 $\End(V)$ 上的拓扑]
\label{def:I.3.1.1}
若 $V$ 和 $W$ 是赋范 $\CC$-向量空间, 我们记 $\Hom(V, W)$ 为从 $V$ 到 $W$ 的有界线性算子组成的 $\CC$-向量空间, 并在其上赋予由算子范数 (operator norm) $\|\cdot\|_{\mathrm{op}}$ 给出的拓扑. 又记 $\End(V) = \Hom(V, V)$, $\GL(V) = \End(V)^{\times}$, 赋予 $\End(V)$ 的拓扑所诱导的拓扑.
\end{definition}

\begin{definition}[连续表示 (continuous representation)]
\label{def:I.3.1.2}
设 $V$ 是赋范 $\CC$-向量空间. 则 $G$ 在 $V$ 上的 (连续) 表示是群同态 $\rho \colon G \to \Aut_{\CC}(V)$ 使得作用映射 $G \times V \to V$, $(g, v) \longmapsto \rho(g)(v)$ 连续.
\end{definition}

我们用 $(\rho, V)$, $\rho$ 或常常简单地用 $V$ 来指代表示. 有时我们不显式给出映射 $\rho$ 的名字, 而把 $G$ 在 $V$ 上的作用写成 $(g, v) \longmapsto gv$.

\begin{remark}[定义的一般性]
\label{rem:I.3.1.3}
\begin{itemize}
\item 若 $V$ 是任意拓扑向量空间 (在拓扑域上), 此定义仍有意义.
\item 若 $(\rho, V)$ 是 $G$ 的连续表示, 则每个 $g \in G$ 在 $V$ 上的作用是 $V$ 的连续自同态, 所以我们得到一个群同态 $\rho \colon G \to \GL(V)$. 但除非 $V$ 有限维, 这个同态不一定连续 (见命题 \ref{prop:I.3.5.1}). 例 \ref{exa:I.3.1.11} 中定义的 $G$ 在 $L^p(G)$ 上的正则表示就是这样的例子.
\item 若 $\rho \colon G \to \GL(V)$ 是关于 $\End(V)$ 上的弱*拓扑连续的群同态, 则它不一定是连续表示. (例如, 取 $G = \GL(V)$, 赋予 $\End(V)$ 上弱*拓扑诱导的拓扑, $\rho = \mathrm{id}$. 若 $V$ 无穷维, 则这不是 $G$ 在 $V$ 上的连续表示.)
\end{itemize}
\end{remark}

\begin{example}[连续表示的例子]
\label{exa:I.3.1.4}
\begin{itemize}
\item $G$ 在 $V$ 上的\index{平凡表示}平凡表示 (trivial representation) 是由 $\rho(x) = \mathrm{id}_V$ (对每个 $x \in G$) 给出的表示. (它是连续表示.)
\item 若 $V$ 有限维, 则 $\GL(V)$ 的恒等映射是 $\GL(V)$ 在 $V$ 上的连续表示.
\item 若 $G = S^1$ 且 $n \in \Z$, 则映射 $G \to \CC$, $z \longmapsto z^n$ 是 $G$ 在 $\CC$ 上的连续表示.
\item 映射 $\rho \colon \R \to \GL_2(\CC)$, $x \longmapsto \begin{pmatrix} 1 & x \\ 0 & 1 \end{pmatrix}$ 是 $\R$ 在 $\CC^2$ 上的连续表示.
\item 关于 $G$ 在其函数空间上的表示, 见例 \ref{exa:I.3.1.11}.
\end{itemize}
\end{example}

\begin{definition}[\index{交错算子}交错算子 (intertwining operator)]
\label{def:I.3.1.5}
设 $(\rho_1, V_1)$ 和 $(\rho_2, V_2)$ 是 $G$ 的两个表示. 从 $V_1$ 到 $V_2$ 的交错算子 (或 $G$-\index{等变映射}等变映射 (G-equivariant map)) 是从 $V_1$ 到 $V_2$ 的有界 $\CC$-线性映射 $T$ 使得对每个 $g \in G$ 和每个 $v \in V_1$, 有 $T(\rho_1(g)v) = \rho_2(g)T(v)$.

我们记 $\Hom_G(V_1, V_2)$ 为从 $V_1$ 到 $V_2$ 的交错算子空间, $\End_G(V_1)$ 为从 $V_1$ 到自身的交错算子空间.

若存在交错算子 $T \colon V_1 \to V_2$ 与 $T' \colon V_2 \to V_1$ 使得 $T' \circ T = \mathrm{id}_{V_1}$ 且 $T \circ T' = \mathrm{id}_{V_2}$, 则称表示 $(\rho_1, V_1)$ 与 $(\rho_2, V_2)$ \index{同构的表示}同构 (isomorphic) (或\index{等价的表示}等价 (equivalent)).
\end{definition}

\begin{definition}[\index{子表示}子表示与不可约等概念]
\label{def:I.3.1.6}
设 $(\rho, V)$ 是 $G$ 的一个表示.
\begin{enumerate}
\item[(1)] $V$ 的子表示 (subrepresentation) (或 $G$-不变子空间 (G-invariant subspace)) 是满足对每个 $g \in G$ 都有 $\rho(g)(W) \subset W$ 的线性子空间 $W$.
\item[(2)] 若 $V \neq 0$ 且其唯一的闭 $G$-不变子空间是 $0$ 和 $V$, 则称表示 $(\rho, V)$ 不可约 (irreducible). 否则称其为可约 (reducible) 的.
\item[(3)] 若每当 $V = W_1 \oplus W_2$ (其中 $W_1$, $W_2$ 是 $V$ 的两个闭 $G$-不变子空间) 就有 $W_1 = 0$ 或 $W_2 = 0$, 则称表示 $(\rho, V)$ \index{不可分解表示}不可分解 (indecomposable).
\item[(4)] 若存在一族 $V$ 的闭 $G$-不变子空间 $(W_i)_{i \in I}$ 处于直和中且 $\bigoplus_{i \in I} W_i$ 在 $V$ 中稠密, 则称表示 $(\rho, V)$ \index{半单表示}半单 (semisimple). (若 $I$ 有限, 则直和也在 $V$ 中闭, 所以这蕴含 $V = \bigoplus_{i \in I} W_i$.)
\end{enumerate}
\end{definition}

\begin{remark}[$G$-稳定子空间的闭包]
\label{rem:I.3.1.7}
若 $(\rho, V)$ 是 $G$ 的表示且 $W \subset V$ 是 $G$-稳定子空间, 则其闭包 $\overline{W}$ 也被 $G$ 稳定.
\end{remark}

\begin{example}[不可分解但不不可约的表示]
\label{exa:I.3.1.8}
$\R$ 在 $\CC^2$ 上由 $\rho(x) = \begin{pmatrix} 1 & x \\ 0 & 1 \end{pmatrix}$ 给出的表示 $\rho$ 是不可分解的但不是不可约的.
\end{example}

\begin{lemma}[交错算子的核与像]
\label{lem:I.3.1.9}
设 $(\rho_1, V_1)$ 和 $(\rho_2, V_2)$ 是 $G$ 的两个表示, $T \colon V_1 \to V_2$ 是一个交错算子. 则 $\Ker(T)$ 是 $V_1$ 的子表示, $\Im(T)$ 是 $V_2$ 的子表示.
\end{lemma}

\begin{proof}
设 $v \in \Ker(T)$, $g \in G$. 则 $T(\rho_1(g)(v)) = \rho_2(g)(T(v)) = 0$, 所以 $\rho_1(g)(v) \in \Ker(T)$.

现在设 $w \in \Im(T)$, 并选取 $v \in V_1$ 使得 $w = T(v)$. 则 $\rho_2(g)(w) = T(\rho_1(g)(v)) \in \Im(T)$.
\end{proof}

\begin{proposition}[连续表示的一个判定]
\label{prop:I.3.1.10}
设 $V$ 是赋范向量空间, $\rho \colon G \to \End(V)$ 是乘法映射. 我们记 $\|\cdot\|_{\mathrm{op}}$ 为 $\End(V)$ 上的算子范数. 假设:
\begin{enumerate}
\item[(a)] 对每个 $g \in G$, 有 $\|\rho(g)\|_{\mathrm{op}} \leq 1$;
\item[(b)] 对每个 $v \in V$, 映射 $G \to V$, $g \longmapsto \rho(g)(v)$ 连续.
\end{enumerate}
则 $(\rho, V)$ 是 $G$ 的连续表示.
\end{proposition}

\begin{proof}
设 $g_0 \in G$, $v_0 \in V$, $\varepsilon > 0$. 我们要找 $g$ 在 $G$ 中的邻域 $U$ 和 $\delta > 0$ 使得: $g \in U$ 且 $\|v - v_0\| < \delta \Rightarrow \|\rho(g)(v) - \rho(g_0)(v_0)\| < \varepsilon$.

选取 $g$ 在 $G$ 中的邻域 $U$ 使得: $g \in U \Rightarrow \|\rho(g)(v_0) - \rho(g_0)(v_0)\| < \varepsilon/2$, 并取 $\delta = \varepsilon/2$. 则若 $g \in U$ 且 $\|v - v_0\| < \delta$, 有
\begin{multline*}
\|\rho(g)(v) - \rho(g_0)(v_0)\| \leq \|\rho(g)(v) - \rho(g)(v_0)\| + \|\rho(g)(v_0) - \rho(g_0)(v_0)\| \\
< \|\rho(g)\|_{\mathrm{op}}\|v - v_0\| + \varepsilon/2 < \varepsilon/2 + \varepsilon/2 = \varepsilon,
\end{multline*}
因为 $\|\rho(g)\|_{\mathrm{op}} \leq 1$.
\end{proof}

\begin{example}[正则表示]
\label{exa:I.3.1.11}
\begin{enumerate}
\item[(1)] 我们已对每个 $x \in G$ 定义了函数空间上的两个自同态 $L_x$ 和 $R_x$, 且这些自同态保持 $\|\cdot\|_{\infty}$. 所以由命题 \ref{prop:I.3.1.10} 和注记 \ref{rem:I.1.13}, 它们定义了 $G$ 在 $\mathcal{C}_c(G)$ 上的两个表示.
\item[(2)] 假设 $G$ 是局部紧 Hausdorff 的. 我们固定 $G$ 上的左 Haar 测度 $dx$, 并对 $1 \leq p \leq \infty$ 记 $L^p(G)$ 为此测度对应的 $L^p$ 空间. 测度的左不变性蕴含算子 $L_x$ 保持 $L^p$ 范数, 所以我们得到 $G$ 在 $L^p(G)$ 上的 $\CC$-线性左作用; 由命题 \ref{prop:I.3.1.10}, 要证明它是一个表示, 只需证明: 若 $f \in L^p(G)$, 则映射 $G \to L^p(G)$, $x \longmapsto L_xf$ 连续. 当 $p = \infty$ 时这不一定成立, 但对 $1 \leq p < \infty$ 成立, 见下面的命题 \ref{prop:I.3.1.13}. 所以对 $1 \leq p < \infty$ 我们得到 $G$ 在 $L^p(G)$ 上的表示.
  若我们改取 $G$ 上的右 Haar 测度, 则算子 $R_x$ 将对 $1 \leq p < \infty$ 定义 $G$ 在 $L^p(G)$ 上的表示. 所以若 $G$ 是幺模的, 我们得到 $G$ 在 $L^p(G)$ 上的两个可交换的表示.
\end{enumerate}
\end{example}

\begin{definition}[左/\index{右正则表示}右正则表示 (left/right regular representation)]
\label{def:I.3.1.12}
设 $G$ 是带有左 (相应地右) Haar 测度 $dx$ 的局部紧群, $L^2(G)$ 是对应的 $L^2$ 空间. 由算子 $L_x$ (相应地 $R_x$) 给出的 $G$ 在 $L^2(G)$ 上的表示称为 $G$ 的左 (相应地右) 正则表示.
\end{definition}

\begin{proposition}[$L^p$ 中平移的连续性]
\label{prop:I.3.1.13}
设 $G$ 是局部紧群, $\mu$ 是 $G$ 上的左 Haar 测度, $L^p(G)$ 是对应的 $L^p$ 空间. 假设 $1 \leq p < \infty$.
则对每个 $f \in L^p(G)$, 当 $x \to 1$ 时有 $\|L_xf - f\|_p \to 0$ 且 $\|R_xf - f\|_p \to 0$.
\end{proposition}

\begin{proof}
先假设 $f \in \mathcal{C}_c(G)$, 并固定 $1$ 的紧邻域 $V$. 则 $K := V(\operatorname{supp} f) \cup (\operatorname{supp} f)V$ 由命题 \ref{prop:I.1.7} 是紧的, 所以 $\mu(K) < +\infty$. 对每个 $x \in V$, 有 $\operatorname{supp}(f), \operatorname{supp}(L_xf), \operatorname{supp}(R_xf) \subset K$, 所以 $\|L_xf - f\|_p \leq \mu(K)^{1/p}\|L_xf - f\|_{\infty}$ 且 $\|R_xf - f\|_p \leq \mu(K)^{1/p}\|R_xf - f\|_{\infty}$. 结论由命题 \ref{prop:I.1.12} 得出.

现在设 $f$ 是 $L^p(G)$ 的任意元素. 我们仍固定 $1$ 的紧邻域 $V$, 并令 $C = \sup_{x \in V}\Delta(x)^{-1/p}$. 设 $\varepsilon > 0$. 存在 $g \in \mathcal{C}_c(G)$ 使得 $\|f - g\|_p < \varepsilon$. 则对 $x \in V$,
\[
\|L_xf - f\|_p \leq \|L_x(f - g)\|_p + \|L_xg - g\|_p + \|g - f\|_p \leq 2\varepsilon + \|L_xg - g\|_p
\]
(因为 $\|L_x(f - g)\|_p = \|f - g\|_p$) 且
\[
\|R_xf - f\|_p \leq \|R_x(f - g)\|_p + \|R_xg - g\|_p + \|g - f\|_p \leq (1 + C)\varepsilon + \|R_xg - g\|_p
\]
(因为 $\|R_x(f - g)\|_p = \Delta(x)^{-1/p}\|f - g\|_p$). 由证明的第一部分, 当 $x$ 趋于 $1$ 时 $\|L_xg - g\|_p$ 和 $\|R_xg - g\|_p$ 趋于 $0$, 所以我们可以找到 $1$ 的邻域 $U \subset V$ 使得对 $x \in U$ 有 $\|L_xf - f\|_p \leq 3\varepsilon$ 且 $\|R_xf - f\|_p \leq (2 + C)\varepsilon$.
\end{proof}

\subsection{\index{酉表示}酉表示}

回忆 (复) \index{希尔伯特空间}希尔伯特空间是带 Hermitian 内积\footnote{我们总假定 Hermitian 内积对第一个变量是 $\CC$-线性的.} 的 $\CC$-向量空间 $V$, 使得 $V$ 对相应范数完备. 若 $V$ 是带 Hermitian 内积的有限维 $\CC$-向量空间, 则它自动完备, 因此是希尔伯特空间. 我们通常用 $\langle \cdot, \cdot \rangle$ 表示所有 Hermitian 内积空间上的内积 (除非另有说明).

\begin{notation}[伴随算子与正交补]
\label{not:I.3.2.1}
设 $V$ 和 $W$ 是 Hermitian 内积空间. 对每个连续 $\CC$-线性映射 $T \colon V \to W$, 若 $T$ 的伴随 (adjoint) 存在, 我们记 $T^* \colon W \to V$. 回忆对每个 $v \in V$ 和 $w \in W$ 有 $\langle T(v), w \rangle = \langle v, T^*(w) \rangle$, 且若 $V$ 和 $W$ 是希尔伯特空间则 $T^*$ 总存在.
若 $V'$ 是 $V$ 的子空间, 我们记 $V'^{\perp}$ 为 $V'$ 的正交补 (orthogonal complement); 它由
\[
(V')^{\perp} = \{v \in V \mid \forall v' \in V',\, \langle v, v' \rangle = 0\}
\]
定义.

最后, 我们记 $\UU(V)$ 为 $V$ 的酉自同构 (unitary endomorphism) 群, 即保持内积的 $V$ 的自同构 $T$ (对所有 $v, w \in V$ 有 $\langle T(v), T(w) \rangle = \langle v, w \rangle$) 的群. 酉自同构 $T$ 自动有界且可逆 (逆等于 $T^*$).
\end{notation}

下面的结果是命题 \ref{prop:I.3.1.10} (以及酉算子范数为 $1$ 这一事实) 的直接推论.

\begin{corollary}[酉表示的等价条件]
\label{cor:I.3.2.2}
若 $V$ 是希尔伯特空间且 $\rho \colon G \to \UU(V)$ 是群同态, 则以下条件等价:
\begin{enumerate}
\item[(1)] 映射 $G \times V \to V$, $(g, v) \longmapsto \rho(g)(v)$ 连续.
\item[(2)] 对每个 $v \in V$, 映射 $G \to V$, $g \longmapsto \rho(g)(v)$ 连续.
\end{enumerate}
\end{corollary}

\begin{definition}[酉表示 (unitary representation)]
\label{def:I.3.2.3}
若 $V$ 是希尔伯特空间, 则 $G$ 在 $V$ 上的酉表示是满足推论 \ref{cor:I.3.2.2} 的条件的群同态 $\rho \colon G \to \UU(V)$.
\end{definition}

这些表示是我们研究的主要对象.

\begin{example}[$L^2$ 空间上的酉表示]
\label{exa:I.3.2.4}
若 $(X, \mu)$ 是任意测度空间, 则 $L^2(X)$ 是希尔伯特空间, 内积为
\[
\langle f, g \rangle = \int_X f(x)\overline{g(x)} \,\dd\mu(x).
\]
所以若 $G$ 是局部紧群, 则 $G$ 的\index{左正则表示}左正则表示和右正则表示是 $G$ 的酉表示 (若 $G$ 幺模则在同一个空间上).
\end{example}

\begin{remark}[$\rho$ 不一定连续]
\label{rem:I.3.2.5}
注意 $\rho$ 一般来说仍不一定是连续映射. (除非 $\dim_{\CC}V < +\infty$.) 例如, 对 $S^1$ 的左正则表示它不连续.
又注意, 证明中不需要 $V$ 的完备性, 所以推论 \ref{cor:I.3.2.2} 对任意 Hermitian 内积空间都成立.
\end{remark}

\begin{lemma}[不变子空间的正交补]
\label{lem:I.3.2.6}
设 $(\rho, V)$ 是 $G$ 的酉表示. 则对 $V$ 的每个 $G$-不变子空间 $W$, 子空间 $W^{\perp}$ 也被 $G$ 不变.

特别地, 若 $W$ 是 $V$ 的闭 $G$-不变子空间, 则 $V = W \oplus W^{\perp}$, 其中 $W^{\perp}$ 是闭 $G$-不变子空间.
\end{lemma}

\begin{proof}
设 $v \in W^{\perp}$, $g \in G$. 则对每个 $w \in W$,
\[
\langle \rho(g)(v), w \rangle = \langle v, \rho(g)^{-1}w \rangle = 0
\]
(最后一个等式来自 $\rho(g)^{-1}w \in W$), 因此 $\rho(g)(v) \in W^{\perp}$.
\end{proof}

\begin{lemma}[交错算子的伴随]
\label{lem:I.3.2.7}
设 $(\rho_1, V_1)$ 和 $(\rho_2, V_2)$ 是 $G$ 的两个酉表示, $T \colon V_1 \to V_2$ 是交错算子. 则 $T^* \colon V_2 \to V_1$ 也是交错算子.
\end{lemma}

\begin{proof}
设 $w \in V_2$, $g \in G$. 则对每个 $v \in V_1$,
\begin{multline*}
\langle v, T^*(\rho_2(g)(w)) \rangle = \langle T(v), \rho_2(g)(w) \rangle = \langle \rho_2(g)^{-1}T(v), w \rangle = \langle T(\rho_1(g)^{-1}(v)), w \rangle \\
= \langle \rho_1(g)^{-1}(v), T^*(w) \rangle = \langle v, \rho_1(g)T^*(w) \rangle.
\end{multline*}
所以 $T^*(\rho_2(g)(w)) = \rho_1(g)(T^*(w))$.
\end{proof}

\begin{theorem}[平均化内积]
\label{thm:I.3.2.8}
假设群 $G$ 是紧 Hausdorff 的. 设 $(V, \langle \cdot, \cdot \rangle_0)$ 是希尔伯特空间, $\rho \colon G \to \GL(V)$ 是 $G$ 在 $V$ 上的连续表示. 则存在 $V$ 上的 Hermitian 内积 $\langle \cdot, \cdot \rangle$ 满足以下性质:
\begin{enumerate}
\item[(1)] 存在实数 $c, C > 0$ 使得对每个 $v \in V$ 有 $c|hv, vi_0| \leq |\langle v, v \rangle| \leq C|\langle v, v \rangle_0|$. 换句话说, 两个内积给出的范数等价, 所以 $V$ 对内积 $\langle \cdot, \cdot \rangle$ 仍是希尔伯特空间.
\item[(2)] 表示 $\rho$ 对内积 $\langle \cdot, \cdot \rangle$ 是酉的.
\end{enumerate}
\end{theorem}

\begin{remarkstar}[I.3.1]
(a) 若 $V$ 不可约, 则由 \index{Schur 引理}Schur 引理 (见定理 \ref{thm:I.3.4.1}) 知, 这个内积在相差一个常数意义下唯一.

(b) 对非紧群这是错的. 例如, 考虑 $\R$ 在 $\CC^2$ 上由 $\rho(t) = \begin{pmatrix} 1 & t \\ 0 & 1 \end{pmatrix}$ 给出的表示 $\rho$. 不存在使该表示为酉的 $\CC^2$ 上的内积 (否则 $\rho(\R)$ 将是该内积的酉群的闭子群, 从而是紧的, 但这是不可能的, 因为 $\rho(\R) \simeq \R$).
\end{remarkstar}

\begin{proof}[定理的证明]
我们用下面的公式定义 $\langle \cdot, \cdot \rangle \colon V \times V \to \CC$: 对所有 $v, w \in V$,
\[
\langle v, w \rangle = \int_G \langle \rho(g)v, \rho(g)w \rangle_0 \,\dd g,
\]
其中 $dg$ 是 $G$ 上的规范化 Haar 测度. 这定义了 $V$ 上的 Hermitian 形式, 且由测度的左不变性, 对每个 $v, w \in V$ 和 $g \in G$ 有 $\langle \rho(g)v, \rho(g)w \rangle = \langle v, w \rangle$.

若我们证明了性质 (1), 它也将蕴含 $\langle \cdot, \cdot \rangle$ 正定 (从而是内积), 于是大功告成. 设 $v \in V$. 则把 $v$ 分别映到 $\rho(g)(v)$ 和 $\rho(g)^{-1}(v)$ 的两个映射 $G \to V$ 连续. 由于 $G$ 紧, 它们都有界. 由共鸣定理 (定理 \ref{thm:I.3.2.11}), 存在 $A, B \in \R_{>0}$ 使得对每个 $g \in G$ 有 $\|\rho(g)^{-1}\|_{\mathrm{op}} \leq A$ 且 $\|\rho(g)\|_{\mathrm{op}} \leq B$. 由算子范数的次乘性, 第一个不等式蕴含对每个 $g \in G$ 有 $\|\rho(g)\|_{\mathrm{op}} \leq A^{-1}$. 所以 $\langle \cdot, \cdot \rangle$ 的定义 (以及 $G$ 体积为 $1$ 的事实) 给出性质 (1), 其中 $c = A^{-2}$, $C = B^2$.
\end{proof}

\begin{corollary}[紧群的有限维表示半单]
\label{cor:I.3.2.9}
若 $G$ 是紧 Hausdorff 的, 则 $G$ 的每个非零有限维连续表示都是半单的.
\end{corollary}

\begin{proof}
由定理, 我们可以假设表示是酉的. 我们对 $\dim V$ 归纳证明此推论. 若 $\dim V \leq 1$ 结论显然, 所以假设 $\dim V \geq 2$ 且已知结论对所有维数更小的空间成立. 若 $V$ 不可约, 则已完成. 否则, 存在 $G$-不变子空间 $W \subsetneq V$ 且 $W \neq 0$. 该子空间因有限维而闭, 且由引理 \ref{lem:I.3.2.6} 有 $V = W \oplus W^{\perp}$, 其中 $W^{\perp}$ 不变. 由于 $\dim(W), \dim(W^{\perp}) < \dim(V)$, 可以对 $W$ 和 $W^{\perp}$ 应用归纳假设并得出它们半单. 但此时它们的直和 $V$ 也半单.
\end{proof}

\begin{remark}[半单性的进一步讨论]
\label{rem:I.3.2.10}
这对紧群的无穷维酉表示仍然成立 (但更难证明, 见定理 \ref{thm:IV.2.1}), 但对非紧群的无穷维酉表示是错的 (例如, 若 $G$ 交换且不紧, 则由推论 \ref{cor:I.3.4.4} 和习题 \ref{ex:III.6.1.2}(c), 其正则表示不半单), 对非紧群的有限维 (非酉) 表示也是错的 (见例 \ref{exa:I.3.1.8}).
\end{remark}

\begin{theorem}[共鸣定理 (Uniform Boundedness Principle)]
\label{thm:I.3.2.11}
设 $V$ 和 $W$ 是赋范向量空间, 并假设 $V$ 是巴拿赫空间 (即它对范数完备). 设 $(T_i)_{i \in I}$ 是从 $V$ 到 $W$ 的有界线性算子族.
若族 $(T_i)_{i \in I}$ 逐点有界 (即对每个 $v \in V$ 有 $\sup_{i \in I}\|T_i(v)\| < +\infty$), 则有界 (即 $\sup_{i \in I}\|T_i\|_{\mathrm{op}} < +\infty$).
\end{theorem}

\begin{proof}
假设 $\sup_{i \in I}\|T_i\|_{\mathrm{op}} = +\infty$, 并选取 $I$ 的元素序列 $(i_n)_{n \geq 0}$ 使得 $\|T_{i_n}\|_{\mathrm{op}} \geq 4^n$. 我们用如下方式定义 $V$ 的元素序列 $(v_n)_{n \geq 0}$:
\begin{itemize}
\item $v_0 = 0$;
\item 对 $n \geq 1$, 由下面的引理 \ref{lem:I.3.2.12}, 可以找到 $V$ 的元素 $v_n$ 使得 $\|v_n - v_{n-1}\| \leq 3^{-n}$ 且 $\|T_{i_n}(v_n)\| \geq \frac{2}{3} \cdot 3^{-n}\|T_{i_n}\|_{\mathrm{op}}$.
\end{itemize}
对 $m \geq n$ 有 $\|v_n - v_m\| \leq \frac{1}{2} \cdot 3^{-n}$, 所以序列 $(v_n)_{n \geq 0}$ 是 Cauchy 列; 由于 $V$ 完备, 它有极限 $v$, 且对每个 $n \geq 0$ 有 $\|v_n - v\| \leq \frac{1}{2} \cdot 3^{-n}$. 不等式 $\|T_{i_n}(v_n)\| \geq \frac{2}{3} \cdot 3^{-n}\|T_{i_n}\|_{\mathrm{op}}$ 和三角不等式现在蕴含 $\|T_{i_n}(x)\| \geq \frac{1}{6} \cdot 3^{-n}\|T_{i_n}\|_{\mathrm{op}} \geq \frac{1}{6}\bigl(\frac{4}{3}\bigr)^n$, 所以序列 $(\|T_{i_n}(x)\|)_{n \geq 0}$ 无界, 这与假设矛盾.
\end{proof}

\begin{lemma}[范数的局部下界]
\label{lem:I.3.2.12}
设 $V$ 和 $W$ 是赋范向量空间, $T \colon V \to W$ 是有界线性算子. 则对任意 $v \in V$ 和 $r > 0$, 有
\[
\sup_{v' \in B(v,r)} \|T(v')\| \geq r\|T\|_{\mathrm{op}},
\]
其中 $B(v, r) = \{v' \in V \mid \|v - v'\| < r\}$.
\end{lemma}

\begin{proof}
对每个 $x \in V$, 有
\[
\|T(x)\| \leq \frac{1}{2}(\|T(v + x)\| + \|T(v - x)\|) \leq \max(\|T(v + x)\|, \|T(v - x)\|).
\]
对 $x \in B(0, r)$ 取上确界就得到引理中的不等式.
\end{proof}

最后, 我们有下面的结果, 其证明用到 Schur 引理 (定理 \ref{thm:I.3.4.1}), 并见习题 \ref{ex:I.5.5.9}.

\begin{theorem}[紧群的不可约酉表示有限维]
\label{thm:I.3.2.13}
若 $G$ 是紧群, 则 $G$ 的每个不可约酉表示都是有限维的.\footnote{取自 Alan Sokal 的一篇论文.}\footnote{精确的参考文献.}
\end{theorem}

\subsection{\index{循环表示}循环表示}

\begin{definition}[循环表示 (cyclic representation)]
\label{def:I.3.3.1}
设 $(\rho, V)$ 是 $G$ 的连续表示, $v \in V$. 则 $\operatorname{Span}\{\rho(g)(v) \mid g \in G\}$ 的闭包 $W$ 是 $V$ 的子表示, 称为由 $v$ 生成的\index{循环子空间}循环子空间 (cyclic subspace).
若 $V = W$, 则称 $V$ 是循环表示, $v$ 是 $V$ 的\index{循环向量}循环向量 (cyclic vector).
\end{definition}

\begin{example}[循环表示的例]
\label{exa:I.3.3.2}
\index{不\index{可约表示}可约表示}不可约表示是循环的, 且其每个非零向量都是循环向量.
其逆不真. 例如, 考虑置换群 $S_n$ 在 $\CC^n$ 上由 $\rho(\sigma)(x_1, \ldots, x_n) = (x_{\sigma^{-1}(1)}, \ldots, x_{\sigma^{-1}(n)})$ 定义的表示 $\rho$, 并取 $v = (1, 0, \ldots, 0) \in \CC^n$. 则集合 $\rho(S_n)(v)$ 是 $\CC^n$ 的典范基, 因此生成 $\CC^n$, 所以 $v$ 是 $\rho$ 的循环向量. 但 $\rho$ 不可约吗? 不是, 因为 $\CC(1, 1, \ldots, 1)$ 是一个子表示.
\end{example}

\begin{proposition}[循环分解]
\label{prop:I.3.3.3}
$G$ 的每个酉表示都是循环表示的直和.
\end{proposition}

若指标集无穷, 我们理解的直和是闭直和 (即代数直和的闭包).

\begin{proof}
设 $(\pi, V)$ 是 $G$ 的酉表示. 由 Zorn 引理, 可以找到 $V$ 的一族两两正交的循环子空间 $(W_i)_{i \in I}$ 的极大族. 假设 $V$ 不是 $W_i$ 的直和, 则存在非零向量 $v \in \bigl(\bigoplus_{i \in I} W_i\bigr)^{\perp}$. 由引理 \ref{lem:I.3.2.6}, 由 $v$ 生成的循环子空间包含于 $\bigl(\bigoplus_{i \in I} W_i\bigr)^{\perp}$, 这与族 $(W_i)_{i \in I}$ 的极大性矛盾. 因此 $V = \bigoplus_{i \in I} W_i$.
\end{proof}

\subsection{Schur 引理}

下面的定理是基本的. 在我们得到希尔伯特空间正规自同态的\index{谱}谱定理 (定理 \ref{thm:II.4.1}) 之前, 我们无法完全证明它.

\begin{theorem}[Schur 引理 (Schur's lemma)]
\label{thm:I.3.4.1}
设 $(\rho_1, V_1)$ 和 $(\rho_2, V_2)$ 是 $G$ 的两个表示, $T \colon V_1 \to V_2$ 是交错算子.

\begin{enumerate}
\item[(1)] 若 $V_1$ 不可约, 则 $T$ 要么是零要么是单射.
\item[(2)] 若 $V_2$ 不可约, 则 $T$ 是零或有稠密像.
\item[(3)] 假设 $V_1$ 是酉的. 则 $V_1$ 不可约当且仅当 $\End_G(V_1) = \CC \cdot \mathrm{id}_{V_1}$.
\item[(4)] 假设 $V_1$ 和 $V_2$ 是酉且不可约的. 则 $\Hom_G(V_1, V_2)$ 的维数是 $0$ (若 $V_1$ 和 $V_2$ 不同构) 或 $1$ (若 $V_1$ 和 $V_2$ 同构).
\end{enumerate}
\end{theorem}

\begin{proof}
我们证明前两点. 由引理 \ref{lem:I.3.1.9}, $\Ker(T)$ 和 $\Im(T)$ 是 $V_1$ 和 $V_2$ 的 $G$-不变子空间. 而且 $\Ker(T)$ 是 $V_1$ 的闭子空间. 若 $V_1$ 不可约, 则其唯一的闭不变子空间是 $0$ 和 $V_1$; 这给出第一点. 若 $V_2$ 不可约, 则其唯一的闭不变子空间是 $0$ 和 $V_2$; 这给出第二点.

我们证明第三点. 先假设 $V_1$ 不可约. 则它有闭不变子空间 $W$ 满足 $0 \neq W \neq V_1$, 且到 $W$ 上的正交投影是 $G$-等变自同态 (由引理 \ref{lem:I.3.4.3}). 所以 $\End_G(V_1)$ 严格包含 $\CC \cdot \mathrm{id}_{V_1}$.

现在假设 $V_1$ 不可约, 设 $T \in \End_G(V_1)$. 我们要证明 $T \in \CC\mathrm{id}_{V_1}$. 若 $V_1$ 有限维, 则 $T$ 有特征值 $\lambda$, 于是 $\Ker(T - \lambda\mathrm{id}_{V_1})$ 是 $V_1$ 的非零 $G$-不变子空间, 从而是 $V_1$, 得到 $T = \lambda\mathrm{id}_{V_1}$. 一般情形下, 我们仍知道每个 $T \in \End(V)$ 都有非空谱 (由定理 \ref{thm:II.1.1.3}), 但若 $\lambda$ 在 $T$ 的谱中, 我们只知道 $T - \lambda\mathrm{id}_V$ 不可逆, 不知道 $\Ker(T - \lambda\mathrm{id}_V) \neq 0$. 所以不能应用同样的策略. 我们改用谱定理 (定理 \ref{thm:II.4.1}) 的一个推论. 注意 $\End(V_1)$ 的子群 $\rho_1(G)$ 满足推论 \ref{cor:II.4.4} 的假设, 因为 $V_1$ 不可约, 所以它在 $\End(V_1)$ 中的中心化子等于 $\CC\mathrm{id}_{V_1}$; 但这个中心化子正是 $\End_G(V_1)$, 于是得证.

我们证明第四点. 设 $T \colon V_1 \to V_2$ 是交错算子. 则由引理 \ref{lem:I.3.2.7}, $T^* \colon V_2 \to V_1$ 也是交错算子, 所以 $T^*T \in \End_G(V_1)$ 且 $TT^* \in \End_G(V_2)$. 由第三点, 存在 $c \in \CC$ 使得 $T^*T = c\mathrm{id}_{V_1}$. 若 $c \neq 0$, 则 $T$ 单射且 $\Im(T)$ 闭 (因为对每个 $v \in V_1$ 有 $\|T(v)\| \geq \frac{|c|}{\|T\|_{\mathrm{op}}}\|v\|$, 见引理 \ref{lem:I.3.4.2}), 所以由第二点 $T$ 是同构, 其逆为 $c^{-1}T^*$; 因此 $V_1$ 和 $V_2$ 同构, 且 $\Hom_G(V_1, V_2) \simeq \End_G(V_1)$ 是 $1$ 维的. 假设 $c = 0$. 若 $T \neq 0$, 则由第二点它有稠密像, 但此时由第一点 $T^* = 0$, 从而 $T = (T^*)^* = 0$, 这荒谬; 所以 $T = 0$. 因此我们证明了: 若 $\Hom_G(V_1, V_2) \neq 0$, 则 $V_1$ 和 $V_2$ 必同构; 这就完成了第四点的证明.
\end{proof}

\begin{lemma}[压缩映射的闭像]
\label{lem:I.3.4.2}
设 $V$, $W$ 是赋范向量空间, $T \colon V \to W$ 是有界线性算子. 假设 $V$ 完备. 若存在 $c > 0$ 使得对每个 $v \in V$ 有 $\|T(v)\| \geq c\|v\|$, 则 $\Im(T)$ 闭.
\end{lemma}

\begin{proof}
设 $(v_n)_{n \in \N}$ 是 $V$ 的元素序列, 使得序列 $(T(v_n))_{n \in \N}$ 收敛到 $w \in W$. 我们要证明 $w \in \Im(T)$. 注意对所有 $n, m \in \N$, 有 $\|v_n - v_m\| \leq c^{-1}\|T(v_n) - T(v_m)\|$. 这蕴含 $(v_n)_{n \in \N}$ 是 Cauchy 列, 所以它因 $V$ 完备而有极限 $v \in V$. 由于 $T$ 连续, $w = \lim_{n \to +\infty} T(v_n) = T(v)$, 所以 $w \in \Im(T)$.
\end{proof}

\begin{lemma}[正交投影的等变性]
\label{lem:I.3.4.3}
设 $(\rho, V)$ 是 $G$ 的酉表示, $W$ 是 $V$ 的闭子空间, $\pi$ 是到 $W$ 上的正交投影, 看作 $V$ 的线性自同态.
则 $W$ 被 $G$ 不变当且仅当 $\pi$ 是 $G$-等变的.
\end{lemma}

\begin{proof}
假设 $\pi$ 是 $G$-等变的. 设 $w \in W$, $g \in G$. 则 $\rho(g)(w) = \rho(g)(\pi(w)) = \pi(\rho(g)(w)) \in W$. 所以 $W$ 被 $G$ 不变.

反之, 假设 $W$ 被 $G$ 不变. 由引理 \ref{lem:I.3.2.6}, 其正交补 $W^{\perp}$ 也被 $G$ 不变. 设 $v \in V$, $g \in G$. 记 $w = \pi(g)$ 且 $w' = g - \pi(g)$. 则 $\rho(g)(v) = \rho(g)(w) + \rho(g)(w')$, 其中 $\rho(g)(w) \in W$ 且 $\rho(g)(w') \in W^{\perp}$, 所以 $\pi(\rho(g)(v)) = \rho(g)(w)$.
\end{proof}

\begin{corollary}[交换群的不可约酉表示]
\label{cor:I.3.4.4}
若 $G$ 交换, 则 $G$ 的每个不可约酉表示都是 $1$ 维的.
\end{corollary}

因此 $G$ 的每个不可约酉表示都等价于唯一一个连续群同态 $G \to S^1$.

\begin{proof}
设 $(\rho, V)$ 是不可约酉表示. 由于 $G$ 交换, 对所有 $x, y \in G$, 算子 $\rho(x)$ 与 $\rho(y)$ 交换, 所以对每个 $x \in G$ 有 $\rho(x) \in \End_G(V)$. 由 Schur 引理, 这蕴含对每个 $x \in G$ 有 $\rho(x) \in \CC \cdot \mathrm{id}_V$. 特别地, $V$ 的每个线性子空间都被 $G$ 不变. 由于 $V$ 不可约, 它没有非平凡的闭不变子空间, 所以它必是 $1$ 维的.
\end{proof}

\begin{example}[$\R$ 的不可约酉表示]
\label{exa:I.3.4.5}
设 $G = \R$. 则 $G$ 的每个不可约酉表示都具有 $\rho_y \colon x \longmapsto e^{ixy}$ 的形式, 其中 $y \in \R$. 表示 $\rho_y$ 通过 $S^1 \simeq \R/\Z$ 分解当且仅当 $y \in 2\pi\Z$. (见习题 \ref{ex:I.5.4.1}(c) 和 (d).)
\end{example}

\subsection{有限维表示}

回忆若 $V$ 是有限维 $\CC$-向量空间, 则 $V$ 上的所有范数等价.\footnote{参考文献 ?} 所以 $V$ 有典范拓扑, $\End(V)$ 也有 (因为它也是有限维向量空间).

\begin{proposition}[有限维表示的条件]
\label{prop:I.3.5.1}
设 $V$ 是赋范 $\CC$-向量空间, $\rho \colon G \to \GL(V)$ 是群同态. 考虑以下条件.
\begin{enumerate}
\item[(i)] 映射 $G \times V \to V$, $(g, v) \longmapsto \rho(g)(v)$ 连续 (即 $\rho$ 是 $G$ 在 $V$ 上的连续表示).
\item[(ii)] 对每个 $v \in V$, 映射 $G \to V$, $g \longmapsto \rho(g)(v)$ 连续.
\item[(iii)] 映射 $\rho \colon G \to \GL(V)$ 连续.
\end{enumerate}
则 $(iii) \Rightarrow (i) \Rightarrow (ii)$. 若此外 $V$ 有限维, 则三个条件等价.
\end{proposition}

\begin{proof}
(i)$\Rightarrow$(ii) 显然.
(ii)$\Rightarrow$(iii): 假设 $V$ 有限维, 设 $(e_1, \ldots, e_n)$ 是 $V$ 的基, 设 $\|\cdot\|$ 是由 $\Bigl\|\sum_{i=1}^n x_ie_i\Bigr\| = \sup_{1 \leq i \leq n}|x_i|$ 定义的 $V$ 上的范数. 我们在 $\End(V)$ 上使用对应的算子范数并仍记之为 $\|\cdot\|$. 设 $g_0 \in G$, $\varepsilon > 0$; 我们要找 $g_0 \in G$ 的邻域 $U$ 使得: $g \in U \Rightarrow \|\rho(g) - \rho(g_0)\| \leq \varepsilon$.
对每个 $i \in \{1, \ldots, n\}$, 由假设, 函数 $G \to V$, $g \longmapsto \rho(g)(e_i)$ 连续, 所以存在 $g_0$ 在 $G$ 中的邻域 $U_i$ 使得: $g \in U_i \Rightarrow \|\rho(g)(e_i) - \rho(g_0)(e_i)\| \leq \varepsilon/n$. 令 $U = \bigcap_{i=1}^n U_i$. 则若 $g \in U$, 对每个 $v = \sum_{i=1}^n x_ie_i \in V$,
\[
\|\rho(g)(v) - \rho(g_0)(v)\| \leq \sum_{i=1}^n |x_i|\|\rho(g)(e_i) - \rho(g_0)(e_i)\| < \sum_{i=1}^n |x_i|\varepsilon/n \leq \varepsilon\|v\|,
\]
这意味着 $\|\rho(g) - \rho(g_0)\| \leq \varepsilon$.
(iii)$\Rightarrow$(i): 设 $g_0 \in G$, $v_0 \in V$, $\varepsilon > 0$. 我们要找 $g$ 在 $G$ 中的邻域 $U$ 和 $\delta > 0$ 使得: $g \in U$ 且 $\|v - v_0\| < \delta \Rightarrow \|\rho(g)(v) - \rho(g_0)(v_0)\| < \varepsilon$.
选取 $\delta$ 使得 $0 < \delta \leq \frac{\varepsilon}{2\|\rho(g_0)\|}$, 设 $U$ 是 $g_0$ 在 $G$ 中的邻域使得: $g \in U \Rightarrow \|\rho(g) - \rho(g_0)\| < \frac{\varepsilon}{2(\|v_0\| + \delta)}$. 则若 $g \in U$ 且 $\|v - v_0\| < \delta$, 有 $\|v\| \leq \|v_0\| + \delta$, 从而
\begin{align*}
\|\rho(g)(v) - \rho(g_0)(v_0)\| &\leq \|\rho(g)(v) - \rho(g_0)(v)\| + \|\rho(g_0)(v) - \rho(g_0)(v_0)\| \\
&\leq \|\rho(g) - \rho(g_0)\|\|v\| + \|\rho(g_0)\|\|v - v_0\| \\
&< \frac{\varepsilon}{2(\|v_0\| + \delta)}(\|v_0\| + \delta) + \|\rho(g_0)\|\delta \\
&\leq \varepsilon/2 + \varepsilon/2 = \varepsilon.
\end{align*}
\end{proof}

% ============ I.4 \index{卷积}卷积与\index{群代数}群代数 (原书 pp.30-44) ============
\section{卷积与群代数}

设 $G$ 是局部紧群, $dx$ 是 $G$ 上的左 Haar 测度. 我们记 $L^p(G)$ 为此测度对应的 $L^p$ 空间. 又记 $\Delta$ 为 $G$ 的模函数.

\subsection{\texorpdfstring{$L^1(G)$ 上的卷积与 $G$ 的群代数}{L^1(G) 上的卷积与 G 的群代数}}

\begin{definition}[卷积 (convolution)]
\label{def:I.4.1.1}
设 $f$ 和 $g$ 是从 $G$ 到 $\CC$ 的函数. $f$ 与 $g$ 的卷积 (convolution), 记为 $f * g$, 是函数 $x \longmapsto \int_G f(y)g(y^{-1}x)\,\dd y$ (如果它有定义).
\end{definition}

\begin{proposition}[卷积的 $L^1$ 范数估计]
\label{prop:I.4.1.2}
设 $f, g \in L^1(G)$. 则积分 $\int_G f(y)g(y^{-1}x)\,\dd y$ 对 $G$ 中几乎每个 $x$ 绝对收敛, 所以 $f * g$ 几乎处处有定义, 且 $f * g \in L^1(G)$, 并有
\[
\|f * g\|_1 \leq \|f\|_1\|g\|_1.
\]
\end{proposition}

\begin{proof}
由 Fubini-Tonelli 定理和 $G$ 上测度的左不变性, 函数 $G \times G \to \CC$, $(x, y) \longmapsto f(y)g(y^{-1}x)$ 可积, 且
\[
\int_{G \times G}|f(y)g(y^{-1}x)|\,\dd x\,\dd y = \int_{G \times G}|f(y)||g(x)|\,\dd x\,\dd y = \|f\|_1\|g\|_1.
\]
所以第一个结论也由 Fubini 定理得出, 而第二个结论显然.
\end{proof}

注意卷积对两个变量都显然是线性的.

\begin{proposition}[卷积的性质]
\label{prop:I.4.1.3}
设 $f, g \in L^1(G)$.
\begin{enumerate}
\item[(1)] 对 $G$ 中几乎每个 $x$,
\begin{align*}
f * g(x) &= \int_G f(y)g(y^{-1}x)\,\dd y \\
&= \int_G f(xy)g(y^{-1})\,\dd y \\
&= \int_G f(y^{-1})g(yx)\Delta(y^{-1})\,\dd y \\
&= \int_G f(xy^{-1})g(y)\Delta(y^{-1})\,\dd y \\
&= \int_G f(y)L_yg(x)\,\dd y \\
&= \int_G g(y^{-1})R_yf(x)\,\dd y.
\end{align*}
\item[(2)] 对每个 $h \in L^1(G)$, 有
\[
(f * g) * h = f * (g * h).
\]
(换句话说, 卷积乘积是结合的.)
\item[(3)] 对每个 $x \in G$, 有
\[
L_x(f * g) = (L_xf) * g
\]
且
\[
R_x(f * g) = f * (R_xg).
\]
\item[(4)] 若 $G$ 交换, 则 $f * g = g * f$.
\end{enumerate}
\end{proposition}

\begin{proof}
(1). 前四行的等式由代换 $y \longmapsto xy$ 和 $y \longmapsto y^{-1}$, $dy$ 的左不变性及命题 \ref{prop:I.2.12} 得到. 最后两行只是前两行的改写.

(2). 对 $G$ 中几乎每个 $x$,
\begin{align*}
((f * g) * h)(x) &= \int_G (f * g)(y)h(y^{-1}x)\,\dd y \\
&= \int_{G \times G} f(z)g(z^{-1}y)h(y^{-1}x)\,\dd z\,\dd y \\
&= \int_G f(z)\Bigl(\int_G g(z^{-1}y)h(y^{-1}x)\,\dd y\Bigr)\,\dd z \\
&= \int_G f(z)\Bigl(\int_G g(y)h(y^{-1}z^{-1}x)\,\dd y\Bigr)\,\dd z \\
&= \int_G f(z)(g * h)(z^{-1}x)\,\dd z \\
&= (f * (g * h))(x).
\end{align*}

(3). 这由定义和 (1) 中前两行的等式立即得出.

(4). 这由 (1) 和 $\Delta = 1$ 的事实得出.
\end{proof}

\begin{definition}[\index{Banach 代数}Banach 代数 (Banach algebra)]
\label{def:I.4.1.4}
($\CC$ 上的) Banach 代数是结合的 $\CC$-代数 $A$, 带有范数 $\|\cdot\|$ 使 $A$ 成为巴拿赫空间 (即完备赋范向量空间), 且对每个 $x, y \in A$ 有 $\|xy\| \leq \|x\|\|y\|$ (即范数是次乘性的). 若 $A$ 有单位元 $e$, 我们还要求 $\|e\| = 1$.
\end{definition}

注意我们不假定 $A$ 有单位元. 若它有, 则称 $A$ 是\index{含幺}含幺 (unital) 的.

\begin{example}[Banach 代数的例子]
\label{exa:I.4.1.5}
(a) 若 $V$ 是巴拿赫空间, 则 $\End(V)$ 是含幺 Banach 代数.
(b) 由命题 \ref{prop:I.4.1.2} 和 \ref{prop:I.4.1.3}, 空间 $L^1(G)$ 带卷积乘积是 Banach 代数. 我们称它为 $G$ 的 ($L^1$) 群代数 (group algebra).
\end{example}

\begin{remark}[\index{测度代数}测度代数 $\mathcal{M}(G)$]
\label{rem:I.4.1.6}
若群 $G$ 离散且 $dx$ 是计数测度, 则 $\delta_1 := \mathbb{1}_{\{1\}}$ 是卷积乘积的单位元. 一般地, $L^1(G)$ 不一定有单位元. (当且仅当 $G$ 离散时它有.)\footnote{参考文献 ?} 实际上我们可以把它看作一个确实有单位元的更大的 Banach 代数的子代数, 即 $G$ 的测度代数 (measure algebra) $\mathcal{M}(G)$ (例如见 {[}11{]} 的第 2.5 节):
回忆 $G$ 上的 (复) Radon 测度是 $\mathcal{C}_0(G)$ (带范数 $\|\cdot\|_{\infty}$) 上的有界线性泛函. 我们记 $\mathcal{M}(G)$ 为 Radon 测度的空间, 记 $\|\cdot\|$ 为其范数 (即算子范数); 这是巴拿赫空间. 若 $\mu$ 是 Radon 测度, 我们记 $f \longmapsto \int_G f(x)\,\dd\mu(x)$ 为 $\mathcal{C}_0(G)$ 上对应的线性泛函. 我们定义两个 Radon 测度 $\mu$ 与 $\nu$ 的卷积乘积 $\mu * \nu$ 为线性泛函
\[
f \longmapsto \int_{G \times G} f(xy)\,\dd\mu(x)\,\dd\nu(y).
\]
则不难验证 $\|\mu * \nu\| \leq \|\mu\|\|\nu\|$ 且 $\mathcal{M}(G)$ 上的卷积乘积结合. 这使得 $\mathcal{M}(G)$ 成为 Banach 代数, 且 $1$ 处的 Dirac 测度是 $\mathcal{M}(G)$ 的单位元.
又注意 $\mathcal{M}(G)$ 交换当且仅当 $G$ 交换. 事实上, 由 $*$ 的定义显然 $\mathcal{M}(G)$ 在 $G$ 交换时交换. 为证其逆, 我们记 $\delta_x$ 为 $x$ 处的 Dirac 测度 (于是 $\int_G f\,\dd\delta_x = f(x)$). 则显然对每个 $x, y \in G$ 有 $\delta_x * \delta_y = \delta_{xy}$. 所以若 $\mathcal{M}(G)$ 交换, 则对每个 $x, y \in G$ 有 $\delta_{xy} = \delta_{yx}$, 这蕴含 $G$ 交换.
尽管 $L^1(G)$ 不含 $\mathcal{M}(G)$ 的单位元, 我们有一族称为``\index{近似单位元}近似单位元 (approximate identities)''的函数, 它们在实际上几乎和 $\delta_1$ 一样好. 特别地, 我们将能证明 $L^1(G)$ 交换当且仅当 $G$ 交换 (见推论 \ref{cor:I.4.1.10}).
\end{remark}

\begin{definition}[近似单位元 (approximate identity)]
\label{def:I.4.1.7}
带支撑在 $1$ 在 $G$ 中的邻域基 $\mathcal{U}$ 中的 (对称、连续的) 近似单位元是 $\mathcal{C}_c^+(G)$ 中的函数族 $(\psi_U)_{U \in \mathcal{U}}$ 使得对每个 $U \in \mathcal{U}$:
\begin{itemize}
\item $\operatorname{supp}(\psi_U) \subset U$;
\item $\psi_U(x^{-1}) = \psi_U(x)$, $\forall x \in G$;
\item $\int_G \psi_U(x)\,\dd x = 1$.
\end{itemize}
\end{definition}

对某些结果, 我们不需要 $\psi_U$ 的连续性或 $\psi_U(x^{-1}) = \psi_U(x)$ 这一事实.

\begin{proposition}[近似单位元的存在性]
\label{prop:I.4.1.8}
对 $1$ 在 $G$ 中的每个邻域基 $\mathcal{U}$, 存在支撑在 $\mathcal{U}$ 中的近似单位元.
\end{proposition}

\begin{proof}
设 $U \in \mathcal{U}$. 则 $U$ 包含 $1$ 的对称邻域 $V \subset U$ 和 $1$ 的紧邻域 $K \subset V$, 且由推论 \ref{cor:A.3.1}, 存在支撑包含于 $V$ 的连续函数 $f \colon X \to [0, 1]$ 使得 $f|_K = 1$. 特别地, $f \neq 0$, 所以 $f \in \mathcal{C}_c^+(X)$. 定义 $g \colon X \to [0, 2]$ 为 $g(x) = f(x) + f(x^{-1})$. 则 $g \in \mathcal{C}_c^+(X)$ (因为 $g|_K = 2$) 且 $\operatorname{supp}(g) \subset V \subset U$. 现在取 $\psi_U = \frac{g}{\int_G g(x)\,\dd x}$.
\end{proof}

\begin{proposition}[近似单位元的逼近性质]
\label{prop:I.4.1.9}
设 $\mathcal{U}$ 是 $1$ 在 $G$ 中的邻域基, $(\psi_U)_{U \in \mathcal{U}}$ 是支撑在 $\mathcal{U}$ 中的近似单位元.
\begin{enumerate}
\item[(1)] 对每个 $f \in L^1(G)$, 当 $U \to \{1\}$ 时 $\|\psi_U * f - f\|_1 \to 0$ 且 $\|f * \psi_U - f\|_1 \to 0$. 事实上, 我们有:
\[
\|\psi_U * f - f\|_1 \leq \sup_{y \in U}\|L_yf - f\|_1
\]
且
\[
\|f * \psi_U - f\|_1 \leq \sup_{y \in U}\|R_yf - f\|_1.
\]
\item[(2)] 若 $f \in L^{\infty}(G)$ 且 $f$ 左 (相应地右) 一致连续, 则当 $U \to \{1\}$ 时 $\|\psi_U * f - f\|_{\infty} \to 0$ (相应地 $\|f * \psi_U - f\|_{\infty} \to 0$). 事实上, 我们有:
\[
\|\psi_U * f - f\|_{\infty} \leq \sup_{y \in U}\|L_yf - f\|_{\infty}
\]
且
\[
\|f * \psi_U - f\|_{\infty} \leq \sup_{y \in U}\|R_yf - f\|_{\infty}.
\]
\end{enumerate}
\end{proposition}

在 (2) 中, 注意若 $f \colon G \to \CC$ 有界且 $g \in \mathcal{C}_c(G)$, 则对每个 $x \in G$, 定义 $(f * g)(x)$ 的积分绝对收敛.

\begin{proof}
(1). 设 $U \in \mathcal{U}$. 对每个 $x \in G$,
\[
(\psi_U * f)(x) - f(x) = \int_G \psi_U(y)(L_yf(x) - f(x))\,\dd y
\]
(因为 $\int_G \psi_U(y)\,\dd y = 1$). 所以
\begin{align*}
\|\psi_U * f - f\|_1 &= \int_G \Bigl|\int_G \psi_U(y)(L_yf(x) - f(x))\,\dd y\Bigr|\,\dd x \\
&\leq \int_{G \times G} \psi_U(y)|L_yf(x) - f(x)|\,\dd y\,\dd x \\
&\leq \int_G \psi_U(y)\|L_yf - f\|_1\,\dd y \\
&\leq \sup_{y \in U}\|L_yf - f\|_1.
\end{align*}
第一个收敛结果由当 $y \to 1$ 时 $\|L_yf - f\|_1 \to 0$ 得出, 即命题 \ref{prop:I.3.1.13}.
第二个收敛结果的证明类似 (我们得到 $\|f * \psi_U - f\|_1 \leq \sup_{y \in U}\|R_yf - f\|_1$ 并应用命题 \ref{prop:I.3.1.13}).

(2). 设 $U \in \mathcal{U}$. 则对每个 $x \in G$,
\[
|(\psi_U * f)(x) - f(x)| \leq \int_G \psi_U(y)|L_yf(x) - f(x)|\,\dd y.
\]
由于对 $y \notin U$ 有 $\psi_U(y) = 0$, 这蕴含
\[
|(\psi_U * f)(x) - f(x)| \leq \bigl(\sup_{y \in U}|L_yf(x) - f(x)|\bigr)\Bigl(\int_G \psi_U(y)\,\dd y\Bigr) = \sup_{y \in U}|L_yf(x) - f(x)|.
\]
对 $x \in G$ 取上确界给出
\[
\|\psi_U * f - f\|_{\infty} \leq \sup_{y \in U}\|L_yf - f\|_{\infty}.
\]
所以第一个结论由左一致连续的定义立即得出. 第二个结论的证明类似.
\end{proof}

\begin{corollary}[$L^1(G)$ 的交换性与\index{理想}理想]
\label{cor:I.4.1.10}
\begin{enumerate}
\item[(1)] Banach 代数 $L^1(G)$ 交换当且仅当群 $G$ 交换.
\item[(2)] 设 $I$ 是 $L^1(G)$ 的闭线性子空间. 则 $I$ 是左 (相应地右) 理想当且仅当它在算子 $L_x$ (相应地 $R_x$) ($x \in G$) 下稳定.
\end{enumerate}
\end{corollary}

\begin{proof}
(1). 若 $G$ 交换, 则我们已见 $L^1(G)$ 交换. 反之, 假设 $L^1(G)$ 交换. 设 $x, y \in G$. 设 $f \in \mathcal{C}_c(G)$, 并选取近似单位元 $(\psi_U)_{U \in \mathcal{U}}$. 由命题 \ref{prop:I.4.1.3}, 对每个 $U \in \mathcal{U}$,
\[
(R_xf) * (R_y\psi_U) = R_y((R_xf) * \psi_U) = R_y(\psi_U * (R_xf)) = R_yR_x(\psi_U * f) = R_{yx}(f * \psi_U)
\]
且
\[
(R_xf) * (R_y\psi_U) = (R_y\psi_U) * (R_xf) = R_x(f * (R_y\psi_U)) = R_xR_y(f * \psi_U) = R_{xy}(f * \psi_U).
\]
在 $1$ 处取值给出 $(f * \psi_U)(xy) = (f * \psi_U)(yx)$. 但命题 \ref{prop:I.4.1.9} (以及命题 \ref{prop:I.1.12}) 蕴含当 $U \to \{1\}$ 时 $\|f * \psi_U - f\|_{\infty} \to 0$, 所以我们得到
\[
f(xy) = \lim_{U \to \{1\}}(f * \psi_U)(xy) = \lim_{U \to \{1\}}(f * \psi_U)(yx) = f(yx).
\]
由于这对每个 $f \in \mathcal{C}_c(G)$ 成立, 必有 $xy = yx$ (这由局部紧性和 Urysohn 引理得出).\footnote{参考文献 ?}

(2). 我们证明左理想的情形 (右理想的证明类似). 假设 $I$ 是左理想, 设 $x \in G$. 选取近似单位元 $(\psi_U)_{U \in \mathcal{U}}$. 我们知道当 $U \to \{1\}$ 时 $\psi_U * f \to f$ ($L^1(G)$ 中), 所以当 $U \to \{1\}$ 时 $L_x(\psi_U * f) \to L_xf$ (因为 $L_x$ 保持 $L^1$ 范数). 但由命题 \ref{prop:I.4.1.3}, $L_x(\psi_U * f) = (L_x\psi_U) * f$; 由于 $I$ 是左理想, 对每个 $U \in \mathcal{U}$ 有 $(L_x\psi_U) * f \in I$, 又因 $I$ 闭, 这最终蕴含 $L_xf \in I$.
反之, 假设 $I$ 被所有算子 $L_x$ ($x \in G$) 稳定. 设 $f \in L^1(G)$, $g \in I$. 由命题 \ref{prop:I.4.1.3}, $f * g = \int_G f(y)L_yg\,\dd y$. 由积分的定义, 函数 $f * g$ 在 $L_yg$ ($y \in G$) 的线性张成的空间的闭包中, 所以由假设 (以及 $I$ 闭) 它在 $I$ 中.
\end{proof}

\subsection{\texorpdfstring{$G$ 的表示与 $L^1(G)$ 的表示}{G 的表示与 L^1(G) 的表示}}

\begin{definition}[\index{Banach $*$-代数}Banach $*$-代数与\index{对合}对合]
\label{def:I.4.2.1}
Banach $*$-代数 (Banach $\ast$-algebra) 是带对合反自同构 $*$ 的 Banach 代数 $A$. (即对每个 $x, y \in A$ 和 $\lambda \in \CC$, 有 $(x + y)^* = x^* + y^*$, $(\lambda x)^* = \bar{\lambda}x^*$, $(xy)^* = y^*x^*$ 且 $(x^*)^* = x$.)
反自同构 $*$ 称为 Banach 代数 $A$ 上的对合 (involution).
\end{definition}

\begin{example}[Banach $*$-代数的例子]
\label{exa:I.4.2.2}
(a) $\CC$, 带对合 $z^* = \bar{z}$.
(b) 若 $G$ 是带左 Haar 测度的局部紧群, 则 $L^1(G)$ 带卷积乘积和由 $f^*(x) = \Delta(x)^{-1}\overline{f(x^{-1})}$ 定义的对合 $*$ 是 Banach $*$-代数 (注意 $f^*$ 属于 $L^1(G)$, 且由命题 \ref{prop:I.2.12} 有 $\int_G f^*(x)\,\dd x = \int_G \overline{f(x)}\,\dd x$ 和 $\int_G |f^*(x)|\,\dd x = \int_G |f(x)|\,\dd x$; 所以 $\|f^*\|_1 = \|f\|_1$). 它交换当且仅当 $G$ 交换, 它有单位元当且仅当 $G$ 离散.
(c) 若 $X$ 是局部紧 Hausdorff 空间, 空间 $\mathcal{C}_0(X)$ 带范数 $\|\cdot\|_{\infty}$, 通常的 (逐点) 乘法和由 $f^*(x) = \overline{f(x)}$ 定义的对合 $*$ 是交换 Banach $*$-代数. 它有单位元当且仅当 $X$ 紧 (单位元是常值函数 $1$).
(d) 设 $H$ 是希尔伯特空间. 则 $\End(H)$ 带算子范数和对合 $T \longmapsto T^*$ (其中 $T^*$ 是如上 $T$ 的伴随) 是含幺 Banach $*$-代数. 它交换当且仅当 $\dim_{\CC}(H) = 1$.
\end{example}

\begin{definition}[\index{$*$-同态}$*$-同态与表示]
\label{def:I.4.2.3}
(i) 若 $A$ 和 $B$ 是两个 Banach $*$-代数, 从 $A$ 到 $B$ 的 $*$-同态 ($*$-homomorphism) 是 $\CC$-代数同态 $u \colon A \to B$, 它作为线性算子有界, 且对每个 $x \in A$ 满足 $u(x^*) = u(x)^*$.
(ii) Banach $*$-代数 $A$ 在希尔伯特空间 $H$ 上的表示是从 $A$ 到 $\End(H)$ 的 $*$-同态 $\pi$. 若对每个 $v \in H - \{0\}$, 存在 $x \in A$ 使得 $\pi(x)(v) \neq 0$, 则称该表示非退化 (nondegenerate).
\end{definition}

我们将需要下面的结果, 它将在下一节中证明. (见推论 \ref{cor:II.3.9}.)

\begin{proposition}[正规算子的\index{谱半径}谱半径公式]
\label{prop:I.4.2.4}
设 $V$ 是希尔伯特空间. 则对每个满足 $TT^* = T^*T$ 的 $T \in \End(H)$, 有
\[
\|T\|_{\mathrm{op}} = \lim_{n \to \infty}\|T^n\|_{\mathrm{op}}^{1/n}.
\]
\end{proposition}

\begin{corollary}[表示自动收缩]
\label{cor:I.4.2.5}
设 $A$ 是对每个 $x \in A$ 满足 $\|x^*\| = \|x\|$ 的 Banach $*$-代数, $\pi$ 是 $A$ 在希尔伯特空间 $V$ 上的表示. 则 $\|\pi\|_{\mathrm{op}} \leq 1$.
\end{corollary}

\begin{proof}
由定义, 算子 $\pi$ 有界; 设 $C = \|\pi\|_{\mathrm{op}}$. 设 $x \in A$, $T = \pi(x^*x) \in \End(H)$. 注意 $T = T^*$. 对每个 $n \geq 1$,
\[
\|T^n\| \leq C\|(x^*x)^n\|_1 \leq C\|x\|^{2n}
\]
(因为 $\|x^*\| = \|x\|$). 另一方面, 由命题 \ref{prop:I.4.2.4},
\[
\|T\|_{\mathrm{op}} = \lim_{n \to +\infty}\|T^n\|_{\mathrm{op}}^{1/n},
\]
从而
\[
\|\pi(x)\|_{\mathrm{op}} = \|\pi(x)^*\pi(x)\|_{\mathrm{op}}^{1/2} = \|T\|_{\mathrm{op}}^{1/2} \leq \bigl(\lim_{n \to +\infty}C^{1/n}\|x\|^{2n/n}\bigr)^{1/2} = \|x\|.
\]
换句话说, $\|\pi\|_{\mathrm{op}} \leq 1$.
\end{proof}

我们现在如前固定一个局部紧群 $G$. 我们将使用习题 \ref{ex:I.5.6.1} 中定义的向量值积分, 以及习题 \ref{ex:I.5.6.2} 和 \ref{ex:I.5.6.3} 中证明的性质.

\begin{theorem}[$G$ 的酉表示与 $L^1(G)$ 的表示]
\label{thm:I.4.2.6}
\begin{enumerate}
\item[(1)] 设 $(\pi, V)$ 是 $G$ 的酉表示. 我们按下述方式定义从 $L^1(G)$ 到 $V$ 的线性自同态空间的映射, 仍记为 $\pi$: 若 $f \in L^1(G)$, 我们令
\[
\pi(f) = \int_G f(x)\pi(x)\,\dd x,
\]
其含义是对每个 $v \in V$,
\[
\pi(f)(v) = \int_G f(x)\pi(x)(v)\,\dd x
\]
(积分由习题 \ref{ex:I.5.6.3} 收敛). 则这是 Banach $*$-代数 $L^1(G)$ 在 $V$ 上的\index{非退化表示}非退化表示, 而且对每个 $x \in G$ 和每个 $f \in L^1(G)$,
\[
\pi(L_xf) = \pi(x)\pi(f) \quad \text{且} \quad \pi(R_xf) = \Delta(x)^{-1}\pi(f)\pi(x)^{-1}.
\]
\item[(2)] Banach $*$-代数 $L^1(G)$ 在希尔伯特空间 $V$ 上的每个非退化表示 $\pi$ 都由 (1) 中那样的群 $G$ 的酉表示 $\pi$ 得到.
而且, 若 $(\psi_U)_{U \in \mathcal{U}}$ 是近似单位元, 则对每个 $x \in G$ 和每个 $v \in V$,
\[
\pi(x)(v) = \lim_{U \to \{1\}}\pi(L_x\psi_U)(v).
\]
\item[(3)] 设 $(\pi, V)$ 是 $G$ 的酉表示, $\pi \colon L^1(G) \to \End(V)$ 是伴随的 $*$-同态. 则 $V$ 的闭子空间 $W$ 被 $G$ 不变当且仅当对每个 $f \in L^1(G)$ 有 $\pi(f)(W) \subset W$.
\item[(4)] 设 $(\pi_1, V_1)$ 和 $(\pi_2, V_2)$ 是 $G$ 的酉表示, $\pi_i \colon L^1(G) \to \End(V_i)$, $i = 1, 2$, 是伴随的 $*$-同态. 则有界线性映射 $T \colon V_1 \to V_2$ 是 $G$-等变的当且仅当对每个 $f \in L^1(G)$ 有 $T \circ \pi_1(f) = \pi_2(f) \circ T$.
\end{enumerate}
\end{theorem}

\begin{proof}
(1). 若 $f \in L^1(G)$, 则映射 $\pi(f) \colon V \to V$ 显然是 $\CC$-线性的, 且对每个 $v \in V$:
\[
\|\pi(f)(v)\| = \Bigl\|\int_G f(x)\pi(x)(v)\,\dd x\Bigr\| \leq \int_G |f(x)|\|v\|\,\dd x \leq \|v\|\|f\|_1,
\]
所以 $V$ 的自同态 $\pi(f)$ 有界且 $\|\pi(f)\|_{\mathrm{op}} \leq \|f\|_1$. 又容易看出把 $f$ 映到 $\pi(f)$ 的映射 $\pi \colon L^1(G) \to \End(H)$ 是 $\CC$-线性的, 而不等式 $\|\pi(f)\|_{\mathrm{op}} \leq \|f\|_1$ 蕴含它也有界 (我们也看到 $\|\pi\|_{\mathrm{op}}$ 被 $1$ 控制, 正如推论 \ref{cor:I.4.2.5} 所应有).
设 $f, g \in L^1(G)$. 则对每个 $v \in V$,
\begin{align*}
\pi(f * g)(v) &= \int_{G \times G} f(y)g(y^{-1}x)\pi(x)(v)\,\dd x\,\dd y \\
&= \int_G f(y)\Bigl(\int_G g(y^{-1}x)\pi(x)(v)\,\dd x\Bigr)\,\dd y \\
&= \int_G f(y)\Bigl(\int_G g(x)\pi(yx)(v)\,\dd x\Bigr)\,\dd y \\
&= \int_G f(y)\pi(y)(\pi(g)(v))\,\dd y \\
&= \pi(f)(\pi(g)(v)).
\end{align*}
所以 $\pi(f * g) = \pi(f) \circ \pi(g)$. 又,
\begin{align*}
\pi(f^*)(v) &= \int_G \Delta(x)^{-1}\overline{f(x^{-1})}\pi(x)(v)\,\dd x \\
&= \int_G \overline{f(x)}\pi(x^{-1})(v)\,\dd x \quad \text{(由命题 \ref{prop:I.2.12})} \\
&= \int_G \overline{f(x)}\pi(x)^*(v)\,\dd x,
\end{align*}
所以对 $w \in V$,
\[
\langle \pi(f^*)(v), w \rangle = \int_G \langle \overline{f(x)}\pi(x)^*(v), w \rangle\,\dd x = \int_G \langle v, \overline{f(x)}\pi(x)(w) \rangle\,\dd x = \langle v, \pi(f)(w) \rangle.
\]
这意味着 $\pi(f^*) = \pi(f)^*$. 所以我们已经证明了 $\pi$ 是 $*$-同态.
设 $f \in L^1(G)$, $x \in G$. 则对每个 $v \in V$,
\[
\pi(x)(\pi(f)(v)) = \int_G f(y)\pi(x)(\pi(y)(v))\,\dd y = \int_G f(x^{-1}y)\pi(y)(v)\,\dd y = \pi(L_xf)(v)
\]
且
\[
\pi(f)(\pi(x)^{-1}(v)) = \int_G f(y)\pi(y)(\pi(x^{-1})(v))\,\dd y = \Delta(x)\int_G f(yx)\pi(y)(v)\,\dd y = \Delta(x)\pi(R_xf)(v).
\]

最后, 我们证明表示 $\pi \colon L^1(G) \to \End(V)$ 非退化. 设 $v \in V - \{0\}$, 选取 $1$ 在 $G$ 中的紧邻域 $K$ 使得对每个 $x \in K$ 有 $\|\pi(x)(v) - v\| \leq \frac{1}{2}\|v\|$. 令 $f = \operatorname{vol}(K)^{-1}\mathbb{1}_K$. 则
\[
\|\pi(f)(v) - v\| = \frac{1}{\operatorname{vol}(K)}\Bigl\|\int_K (\pi(x)(v) - v)\,\dd x\Bigr\| \leq \frac{1}{2}\|v\|,
\]
特别地 $\pi(f)(v) \neq 0$.

(2). 设 $\pi$ 是 Banach $*$-代数 $L^1(G)$ 在希尔伯特空间 $V$ 上的非退化表示. 选取 $G$ 的近似单位元 $(\psi_U)_{U \in \mathcal{U}}$. 证明的想法是 $\pi(x)$ 应当是当 $U$ 趋于 $\{1\}$ 时 $\pi(L_x\psi_U)$ 的极限.

我们现在把上述证明想法严格化. 注意由推论 \ref{cor:I.4.2.5}, $\|\pi\|_{\mathrm{op}} \leq 1$. 设 $W$ 是 $\pi(f)(v)$ ($f \in L^1(G)$, $v \in V$) 的线性张成. 我断言 $W$ 在 $V$ 中稠密. 事实上, 设 $v \in W^{\perp}$. 则对每个 $f \in L^1(G)$, 对所有 $v' \in V$ 有 $\langle \pi(f)(v), v' \rangle = \langle v, \pi(f^*)(v') \rangle = 0$, 因此 $\pi(f)(v) = 0$. 由于 $\pi$ 非退化, 只有 $v = 0$ 才可能. 因此 $W^{\perp} = 0$, 这意味着 $W$ 在 $V$ 中稠密.
设 $x \in G$. 我们要定义元素 $\tilde{\pi}(x) \in \End(V)$ 使得对每个 $f \in L^1(G)$ 有 $\tilde{\pi}(x)\pi(f) = \pi(L_xf)$. 这迫使我们把 $W$ 的元素 $w = \sum_{j=1}^n \pi(f_j)(v_j)$ ($f_j \in L^1(G)$, $v_j \in V$) 上的 $\tilde{\pi}(x)$ 定义为
\[
\tilde{\pi}(x)(w) = \sum_{j=1}^n \pi(L_xf_j)(v_j).
\]
这是良定义的, 因为对每个 $n \geq 1$ 和所有 $f_1, \ldots, f_n \in L^1(G)$, $v_1, \ldots, v_n \in V$,
\begin{align*}
\sum_{j=1}^n \pi(L_xf_j)(v_j) &= \lim_{U \to \{1\}}\sum_{j=1}^n \pi(L_x(\psi_U * f_j))(v_j) \\
&= \lim_{U \to \{1\}}\sum_{j=1}^n \pi((L_x\psi_U) * f_j)(v_j) \\
&= \lim_{U \to \{1\}}\pi(L_x\psi_U)\Bigl(\sum_{j=1}^n \pi(f_j)(v_j)\Bigr),
\end{align*}
所以若 $\sum_{j=1}^n \pi(f_j)(v_j) = 0$, 则 $\sum_{j=1}^n \pi(L_xf_j)(v_j) = 0$.

而且, 由于对每个 $U \in \mathcal{U}$ 有 $\|\pi(L_x\psi_U)\|_{\mathrm{op}} \leq \|\pi\|_{\mathrm{op}}\|\psi_U\|_1 \leq 1$, 对每个 $w \in W$ 有 $\|\tilde{\pi}(x)(w)\| \leq \|w\|$, 所以 $\tilde{\pi}(x)$ 是 $W$ 上范数 $\leq 1$ 的有界线性算子, 因此由连续性延拓成范数 $\leq 1$ 的有界线性算子 $\tilde{\pi}(x) \in \End(V)$.
其次, 利用 $L_{xy} = L_x \circ L_y$, 我们看到对所有 $x, y \in G$, 在 $W$ 上有 $\tilde{\pi}(xy) = \tilde{\pi}(x)\tilde{\pi}(y)$, 从而在 $V$ 上成立. 类似地, $L_1 = \mathrm{id}_{L^1(G)}$ 这一事实蕴含 $\tilde{\pi}(1) = \mathrm{id}_V$. 又对每个 $x \in G$, 若 $v \in V$, 则
\[
\|v\| = \|\tilde{\pi}(x^{-1})\tilde{\pi}(x)(v)\| \leq \|\tilde{\pi}(x^{-1})\|_{\mathrm{op}}\|\tilde{\pi}(x)(v)\| \leq \|\tilde{\pi}(x)(v)\| \leq \|v\|,
\]
所以 $\|\tilde{\pi}(x)(v)\| = \|v\|$, 即 $\tilde{\pi}(x)$ 是酉算子.
设 $v \in V$. 我们要证明映射 $G \to V$, $x \longmapsto \tilde{\pi}(x)(v)$ 连续. 由命题 \ref{prop:I.3.1.10}, 这将蕴含 $\tilde{\pi} \colon G \to \End(V)$ 是 $G$ 在 $V$ 上的酉表示. 我们先假设 $v = \pi(f)(v')$, 其中 $f \in L^1(G)$, $v' \in V$. 则 $\tilde{\pi}(x)(v) = \pi(L_xf)(v')$, 所以结论由映射 $G \to L^1(G)$, $x \longmapsto L_xf$ 的连续性 (见命题 \ref{prop:I.3.1.13})、$\pi$ 的连续性及求值映射 $\End(V) \to V$, $T \longmapsto T(v')$ 的连续性得出. 由于连续函数 $G \to V$ 的有限和仍连续, 我们对每个 $v \in W$ 得到结论. 现在处理一般情形. 设 $x \in G$, $\varepsilon > 0$. 我们必须找到 $x$ 在 $G$ 中的邻域 $U$ 使得对每个 $y \in U$ 有 $\|\tilde{\pi}(y)(v) - \tilde{\pi}(x)(v)\| < \varepsilon$. 选取 $w \in W$ 使得 $\|v - w\| < \varepsilon/3$, 以及 $x$ 在 $G$ 中的邻域 $U$ 使得对每个 $y \in U$ 有 $\|\tilde{\pi}(y)(w) - \tilde{\pi}(x)(w)\| < \varepsilon/3$ (由本段的第一部分这是可能的). 则对每个 $y \in U$,
\begin{multline*}
\|\tilde{\pi}(y)(v) - \tilde{\pi}(x)(v)\| \leq \|\tilde{\pi}(y)(v) - \tilde{\pi}(y)(w)\| + \|\tilde{\pi}(y)(w) - \tilde{\pi}(x)(w)\| + \|\tilde{\pi}(x)(w) - \tilde{\pi}(x)(v)\| \\
< \|v - w\| + \varepsilon/3 + \|v - w\| \quad \text{(因为 $\tilde{\pi}(x)$ 和 $\tilde{\pi}(y)$ 酉)} < \varepsilon,
\end{multline*}
如所愿.

我们证明由 $\tilde{\pi}$ 诱导的 $L^1(G)$ 的表示 $\tilde{\pi}$ 就是我们一开始出发的表示 $\pi$. 设 $f, g \in L^1(G)$. 则对每个 $v \in V$,
\begin{align*}
\tilde{\pi}(f)\pi(g)(v) &= \int_G f(x)\tilde{\pi}(x)\pi(g)(v)\,\dd x \\
&= \int_G f(x)\pi(L_xg)(v)\,\dd x \\
&= \int_G \pi(f(x)L_xg)(v)\,\dd x \\
&= \pi\Bigl(\int_G f(x)L_xg\,\dd x\Bigr)(v) \\
&= \pi(f * g)(v) \\
&= \pi(f)\pi(g)(v).
\end{align*}
所以, 若 $f \in L^1(G)$, 则 $\tilde{\pi}(f)$ 和 $\pi(f)$ 在 $W$ 上相等. 由于 $W$ 在 $V$ 中稠密, 这蕴含 $\tilde{\pi}(f) = \pi(f)$.

最后, 我们证明最后一条结论. 设 $(\psi_U)_{U \in \mathcal{U}}$ 是如上近似单位元. 设 $x \in G$. 我们已见对每个 $v \in W$,
\[
\tilde{\pi}(x)(v) = \lim_{U \to \{1\}}\pi(L_x\psi_U)(v).
\]
由于两边都是 $v \in V$ 的连续函数 (对右端, 我们用 $\|\pi(L_x\psi_U)\|_{\mathrm{op}} = 1$ 这一事实), 此恒等式延拓到所有 $v \in V$.

(3). 假设 $W$ 被 $G$ 不变. 设 $f \in L^1(G)$, $w \in W$. 由于 $\pi(f)(w) = \int_G f(x)\pi(x)(w)\,\dd x$ 是形如 $\pi(x)(w)$ ($x \in G$) 的元素的线性组合的极限, 它仍在 $W$ 中.
反之, 假设对每个 $f \in L^1(G)$ 有 $\pi(f)(W) \subset W$. 设 $x \in G$, $(\psi_U)_{U \in \mathcal{U}}$ 是近似单位元. 则由 (2) 的最后一条结论, 对每个 $w \in W$,
\[
\pi(x)(w) = \lim_{U \to \{1\}}\pi(L_x\psi_U)(w) \in W.
\]
所以 $W$ 被 $G$ 不变.

(4). 设 $T \colon V_1 \to V_2$ 是有界线性映射, $W \subset V_1 \times V_2$ 是 $T$ 的图 (graph); 这是 $V_1 \times V_2$ 的闭线性子空间. 则 $T$ 是 $G$-等变的当且仅当 $W$ 被 $G$ 不变, 而 $T$ 是 $L^1(G)$-等变的当且仅当 $W$ 被所有 $\pi_1(f) \times \pi_2(f)$ ($f \in L^1(G)$) 稳定. 所以结论由 (3) 得出.
\end{proof}

\begin{example}[左正则表示与卷积]
\label{exa:I.4.2.7}
设 $\pi$ 是 $G$ 的由 $\pi(x)(f) = L_xf$ 给出的表示 (见例 \ref{exa:I.3.1.11}). 则对每个 $f, g \in L^1(G)$, 有 $\pi(f)(g) = f * g$. 事实上, 由 $\pi(f)$ 的定义有 $\pi(f)(g) = \int_G f(x)L_xg\,\dd x$, 所以结论由习题 \ref{ex:I.5.6.4} 得出.
\end{example}

\subsection{\texorpdfstring{其他 $L^p$ 空间上的卷积}{其他 L^p 空间上的卷积}}

我们只介绍以后证明紧群 Peter-Weyl 定理所需的一些结果. 最重要的情形是 $L^2(G)$.
大多数结果基于 Minkowski 不等式, 它在习题 \ref{ex:I.5.6.7} 中证明. 这里我们只对 $G$ 上的函数陈述它.

\begin{proposition}[Minkowski 不等式 (Minkowski's inequality)]
\label{prop:I.4.3.1}
设 $p \in [1, +\infty)$, $\phi$ 是从 $G \times G$ 到 $\CC$ 的函数. 则
\[
\Bigl(\int_G \Bigl(\int_G \phi(x, y)\,\dd\mu(y)\Bigr)^p \,\dd\mu(x)\Bigr)^{1/p} \leq \int_G \Bigl(\int_G |\phi(x, y)|^p \,\dd\mu(x)\Bigr)^{1/p} \,\dd\mu(y),
\]
其含义是: 若右端有限, 则 $\int_G \phi(x, y)\,\dd\mu(y)$ 对几乎所有 $x \in G$ 绝对收敛, 左端有限且不等式成立.
\end{proposition}

\begin{corollary}[$L^1 * L^p \subset L^p$]
\label{cor:I.4.3.2}
设 $p \in [1, +\infty)$, $f \in L^1(G)$, $g \in L^p(G)$.
\begin{enumerate}
\item[(1)] 定义 $f * g(x)$ 的积分对几乎所有 $x \in G$ 绝对收敛, 且 $f * g \in L^p(G)$, $\|f * g\|_p \leq \|f\|_1\|g\|_p$.
\item[(2)] 若 $G$ 幺模, 则把 $f * g$ 换成 $g * f$ 后同样的结论成立.
\end{enumerate}
\end{corollary}

\begin{proof}
(1). 我们对函数 $\phi(x, y) = f(y)g(y^{-1}x)$ 应用 Minkowski 不等式. 对每个 $y \in G$,
\[
\int_G |\phi(x, y)|^p \,\dd\mu(x) = |f(y)|^p \int_G |g(x)|^p \,\dd\mu(x) = |f(y)|^p\|g\|_p^p
\]
由 $\mu$ 的左不变性, 所以
\[
\int_G \Bigl(\int_G |\phi(x, y)|^p \,\dd\mu(x)\Bigr)^{1/p} \,\dd\mu(y) = \|g\|_p \int_G |f(y)|\,\dd\mu(y) = \|f\|_1\|g\|_p.
\]
Minkowski 不等式首先说明 $\int_G \phi(x, y)\,\dd\mu(y) = f * g(x)$ 对几乎所有 $x \in G$ 绝对收敛, 这是第一个结论. Minkowski 不等式的其余部分正是 $\|f * g\|_p \leq \|f\|_1\|g\|_p$.

(2). 假设 $G$ 幺模. 则
\[
g * f(x) = \int_G g(y)f(y^{-1}x)\,\dd\mu(x) = \int_G g(xy^{-1})f(y)\,\dd\mu(y).
\]
所以证明与 (1) 相同, 对函数 $\phi(x, y) = g(xy^{-1})f(y)$ 应用 Minkowski 不等式.
\end{proof}

现在我们把命题 \ref{prop:I.4.1.9} 推广到其他 $L^p$ 空间.

\begin{corollary}[$L^p$ 中的近似单位元]
\label{cor:I.4.3.3}
设 $\mathcal{U}$ 是 $1$ 在 $G$ 中的邻域基, $(\psi_U)_{U \in \mathcal{U}}$ 是支撑在 $\mathcal{U}$ 中的近似单位元. 则对每个 $1 \leq p < +\infty$, 若 $f \in L^p(G)$, 则当 $U \to \{1\}$ 时 $\|\psi_U * f - f\|_p \to 0$ 且 $\|f * \psi_U - f\|_p \to 0$.
\end{corollary}

\begin{proof}
设 $U \in \mathcal{U}$, $f \in L^p(G)$. 则对每个 $x \in G$,
\[
\psi_U * f(x) - f(x) = \int_G \psi_U(y)(L_yf(x) - f(x))\,\dd\mu(y)
\]
(因为 $\int_G \psi_U\,\dd\mu = 1$). 对函数 $\phi(x, y) = \psi_U(y)(L_yf(x) - f(x))$ 应用 Minkowski 不等式, 得到
\[
\|\psi_U * f - f\|_p \leq \int_G \|L_yf - f\|_p\psi_U(y)\,\dd\mu(y) \leq \sup_{y \in U}\|L_yf - f\|_p.
\]
类似地,
\begin{align*}
f * \psi_U(x) - f(x) &= \int_G f(xy)\psi_U(y^{-1})\,\dd\mu(y) - f(x)\int_G \psi_U(y)\,\dd\mu(y) \\
&= \int_G (R_yf(x) - f(x))\psi_U(y)\,\dd\mu(y).
\end{align*}
所以对函数 $\phi(x, y) = (R_yf(x) - f(x))\psi_U(y)$ 应用 Minkowski 不等式给出
\[
\|f * \psi_U - f\|_p \leq \int_G \|R_yf - f\|_p\psi_U(y)\,\dd\mu(y) \leq \sup_{y \in U}\|R_yf - f\|_p.
\]
因此两个结论都由命题 \ref{prop:I.3.1.13} 得出.
\end{proof}

最后, 我们证明在某些情形下卷积乘积使函数更光滑. 对我们最重要的情形是 $G$ 紧且 $p = q = 2$.

\begin{proposition}[卷积使函数光滑]
\label{prop:I.4.3.4}
假设 $G$ 幺模. 设 $p, q \in (1, +\infty)$ 满足 $p^{-1} + q^{-1} = 1$, $f \in L^p(G)$, $g \in L^q(G)$.
则 $f * g$ 存在, $f * g \in \mathcal{C}_0(G)$ 且 $\|f * g\|_{\infty} \leq \|f\|_p\|g\|_q$.
\end{proposition}

\begin{proof}
设 $x \in G$. 我们有
\[
f * g(x) = \int_G f(y)R_xg(y^{-1})\,\dd\mu(y).
\]
由于 $G$ 幺模, 函数 $y \longmapsto R_xg(y^{-1})$ 仍在 $L^q(G)$ 中且与 $g$ 有相同的 $L^q$ 范数. 所以由 H\"older 不等式,\footnote{当 $p = q = 2$ 时它化为 Cauchy-Schwarz 不等式.} 上面的积分绝对收敛, 且 $|f * g(x)| \leq \|f\|_p\|g\|_q$. 这证明了 $f * g$ 的存在性和关于其范数的结论. 它也说明双线性映射 $L^p(G) \times L^q(G) \to L^{\infty}(G)$, $(f, g) \longmapsto f * g$ 连续. 由于 $\mathcal{C}_0(G)$ 在 $L^{\infty}(G)$ 中闭且 $\mathcal{C}_c(G)$ 在 $L^p(G)$ 和 $L^q(G)$ 中都稠密, 只需证明 $f, g \in \mathcal{C}_c(G)$ 时 $f * g \in \mathcal{C}_0(G)$.
所以设 $f, g \in \mathcal{C}_c(G)$. 设 $x \in G$, $\varepsilon > 0$, 选取 $x$ 的邻域 $U$ 使得对每个 $y \in G$ 和 $x' \in U$, 有 $|g(yx) - g(yx')| \leq \varepsilon$. 则若 $x' \in U$,
\[
|f * g(x) - f * g(x')| \leq \int_G |f(y)||g(y^{-1}x) - g(y^{-1}x')|\,\dd\mu(y) \leq \varepsilon \int_G |f(y)|\,\dd\mu(y).
\]
这证明 $f * g$ 连续. 设 $K = (\operatorname{supp} g)(\operatorname{supp} f)$; 这是 $G$ 的紧子集. 我们要证明 $\operatorname{supp}(f * g) \subset K$, 这将完成证明. 设 $x \in G$, 假设 $f * g(x) \neq 0$. 则存在 $y \in G$ 使得 $f(y)g(y^{-1}x) \neq 0$. 必有 $y \in \operatorname{supp} f$ 且 $y^{-1}x \in \operatorname{supp} g$, 所以 $x \in y(\operatorname{supp} g) \subset K$.
\end{proof}

% ============ I.5 习题 (原书 pp.44-118) ============
\section{习题}

\subsection{拓扑群的例子}

\begin{exercise}[$GL(V)$ 的拓扑]
\label{ex:I.5.1.1}
设 $V$ 是 $\CC$ 上的巴拿赫空间. (即 $V$ 是赋范 $\CC$-向量空间, 对由范数给出的度量完备.) 我们记 $\End(V)$ 为从 $V$ 到自身的有界线性算子空间, 赋予算子范数. 回忆若 $\|\cdot\|$ 是 $V$ 上的范数, 则算子范数 $\|\cdot\|_{\mathrm{op}}$ 定义为: 对每个 $f \in \End(V)$,
\[
\|f\|_{\mathrm{op}} = \inf\{c \in \R_{\geq 0} \mid \forall v \in V,\, \|f(v)\| \leq c\|v\|\} = \sup_{v \in V, \|v\| = 1}\|f(v)\|.
\]
设 $\GL(V)$ 为 $\End(V)$ 中可逆元素的群, 赋予 $\End(V)$ 的拓扑所诱导的拓扑.
\begin{enumerate}
\item[(a)] 证明 $\GL(V)$ 是 $\End(V)$ 的开子集.
\item[(b)] 证明 $\GL(V)$ 是拓扑群.
\item[(c)] 证明 $\GL(V)$ 局部紧当且仅当 $V$ 有限维.
\end{enumerate}
\end{exercise}

\begin{solution}
(a). 注意由算子范数的定义, 对所有 $x, y \in \End(V)$ 有 $\|xy\|_{\mathrm{op}} \leq \|x\|_{\mathrm{op}}\|y\|_{\mathrm{op}}$. (此性质称为``次乘性 (submultiplicativity)''.) 所以若 $x \in \End(V)$ 满足 $\|x\|_{\mathrm{op}} < 1$, 则级数 $\sum_{n \geq 0}x^n$ 收敛 (按约定取 $x^0 = \mathrm{id}_V$), 且 $(\mathrm{id}_V - x)\bigl(\sum_{n \geq 0}x^n\bigr) = \bigl(\sum_{n \geq 0}x^n\bigr)(\mathrm{id}_V - x) = \mathrm{id}_V$. 因此若 $\|x\|_{\mathrm{op}} < 1$, 则 $\mathrm{id}_V - x \in \GL(V)$.
现在设 $x \in \GL(V)$. 我们要证明 $\GL(V)$ 包含 $x$ 在 $\End(V)$ 中的邻域. 设 $y \in \End(V)$ 满足 $\|y\|_{\mathrm{op}} < \|x^{-1}\|_{\mathrm{op}}$. 则 $\|x^{-1}y\|_{\mathrm{op}} < 1$, 所以 $\mathrm{id}_V - x^{-1}y$ 可逆, 从而 $x - y = x(\mathrm{id}_V - x^{-1}y)$ 可逆. 所以 $\End(V)$ 中满足 $\|x - x'\|_{\mathrm{op}} < \|x^{-1}\|_{\mathrm{op}}$ 的每个元素 $x'$ 都在 $\GL(V)$ 中, 这证明了结论.
若 $V$ 有限维, 我们也可以用行列式是连续映射 $\det \colon \End(V) \to \CC$ 且 $\GL(V)$ 是 $\CC$ 的开子集 $\CC^{\times}$ 的原像这一事实.

(b). 我们来证明乘法是从 $\End(V) \times \End(V)$ 到 $\End(V)$ 的连续映射. (这立即蕴含乘法在 $\GL(V)$ 上连续.) 这由算子范数的次乘性立即得出. 事实上, 若 $x, x', y, y' \in \End(V)$, 则
\[
\|xy - x'y'\|_{\mathrm{op}} = \|x(y - y') + (x - x')y'\|_{\mathrm{op}} \leq \|x\|_{\mathrm{op}}\|y - y'\|_{\mathrm{op}} + \|x - x'\|_{\mathrm{op}}\|y'\|_{\mathrm{op}}.
\]
利用
\[
\|y'\|_{\mathrm{op}} = \|y + (y' - y)\|_{\mathrm{op}} \leq \|y\|_{\mathrm{op}} + \|y - y'\|_{\mathrm{op}},
\]
我们看到, 若固定 $x$ 和 $y$, 则当 $(\|x - x'\|_{\mathrm{op}}, \|y - y'\|_{\mathrm{op}})$ 趋于 $(0,0)$ 时 $\|xy - x'y'\|_{\mathrm{op}}$ 趋于 $0$.

我们来证明取逆在 $\GL(V)$ 上连续. 设 $x \in \GL(V)$. 设 $y \in \End(V)$, 记 $h = x - y$, $c = \|h\|_{\mathrm{op}}\|x^{-1}\|_{\mathrm{op}}$. 则 $y = x - h = x(\mathrm{id}_V - x^{-1}h)$. 我们在 (a) 的解答中已见, 若 $c < 1$, 则 $y$ 可逆且 $y^{-1} = \bigl(\sum_{n \geq 0}(x^{-1}h)^n\bigr)x^{-1} = x^{-1} + \sum_{n \geq 1}(x^{-1}h)^nx^{-1}$; 特别地,
\[
\|y^{-1} - x^{-1}\|_{\mathrm{op}} \leq \sum_{n \geq 1}\|(x^{-1}h)^nx^{-1}\|_{\mathrm{op}} = \|x^{-1}\|_{\mathrm{op}}\sum_{n \geq 1}c^n = \frac{c}{1-c}\|x^{-1}\|_{\mathrm{op}}.
\]
这证明: 若 $x$ 固定, 则当 $\|x - y\|_{\mathrm{op}}$ 趋于 $0$ 时 $\|x^{-1} - y^{-1}\|_{\mathrm{op}}$ 趋于 $0$, 这蕴含结论.
若 $V$ 有限维, 还有另一种证明第二点的方法. 事实上, 此时我们可以假设 $V = \CC^n$ (对某个 $n \in \N$), 于是 $\GL(V) = \GL_n(\CC)$. 然后利用: 若 $x \in \GL_n(\CC)$, 则 $x$ 的逆等于 $(\det x)^{-1}y^T$, 其中 $y$ 是 $x$ 的余因子矩阵. 由于 $y$ 的系数是 $x$ 的连续函数 (因为它们是 $x$ 的子矩阵的行列式的 $\pm 1$ 倍), 这证明 $x^{-1}$ 的系数是 $x$ 的连续函数.

(c). 由 I.5.1.1(a) 的解答, 拓扑群局部紧当且仅当其单位元有紧邻域. 由于由 (a) $\GL(V)$ 在 $\End(V)$ 中开, 这蕴含 $\GL(V)$ 局部紧当且仅当 $e$ 在 $\End(V)$ 中有紧邻域. 由于 $\End(V)$ 的拓扑由范数定义, 这等价于 $\End(V)$ 中的闭球紧. 由 Riesz 引理, 这等价于 $\End(V)$ 有限维. 若 $V$ 有限维, 则 $\End(V)$ 也有限维. 若 $V$ 无穷维, 则由 Hahn-Banach 定理可知 $\End(V)$ 也无穷维.
\end{solution}

\begin{exercise}[拓扑群的乘积]
\label{ex:I.5.1.2}
设 $(G_i)_{i \in I}$ 是一族拓扑群.
\begin{enumerate}
\item[(a)] 证明 $\prod_{i \in I} G_i$ (赋予乘积拓扑) 是拓扑群.
\item[(b)] 若所有 $G_i$ 都局部紧, $\prod_{i \in I} G_i$ 是否总局部紧? (给出证明或反例.)
\end{enumerate}
\end{exercise}

\begin{solution}
(a). 我们来证明乘法连续. 设 $(x_i), (y_i) \in \prod_{i \in I} G_i$. 设 $U$ 是 $(x_iy_i)$ 在 $\prod_{i \in I} G_i$ 中的邻域. 由乘积拓扑的定义, 存在 $I$ 的有限子集 $J$ 和 $x_iy_i$ 在 $G_i$ 中的开邻域 $U_i$ ($i \in J$) 使得 $U \supset \bigl(\prod_{i \in J} U_i\bigr) \times \bigl(\prod_{i \in I-J} G_i\bigr)$. 由 $G_i$ ($i \in J$) 上乘法的连续性, 对每个 $i \in J$ 可以找到 $x_i$ 和 $y_i$ 的开邻域 $V_i$ 和 $W_i$ 使得 $V_iW_i \subset U_i$. 设 $V = \bigl(\prod_{i \in J} V_i\bigr) \times \bigl(\prod_{i \in I-J} G_i\bigr)$, $W = \bigl(\prod_{i \in J} W_i\bigr) \times \bigl(\prod_{i \in I-J} G_i\bigr)$. 则 $V$ 和 $W$ 是 $(x_i)$ 和 $(y_i)$ 在 $\prod_{i \in I} G_i$ 中的开邻域, 且有 $VW \subset U$.

我们来证明取逆连续. (证明类似.) 设 $(x_i) \in \prod_{i \in I} G_i$. 设 $U$ 是 $(x_i^{-1})$ 在 $\prod_{i \in I} G_i$ 中的邻域. 由乘积拓扑的定义, 存在 $I$ 的有限子集 $J$ 和 $x_i$ 在 $G_i$ 中的开邻域 $U_i$ ($i \in J$) 使得 $U \supset \bigl(\prod_{i \in J} U_i\bigr) \times \bigl(\prod_{i \in I-J} G_i\bigr)$. 由 $G_i$ ($i \in J$) 上取逆的连续性, 对每个 $i \in J$ 可以找到 $x_i$ 的开邻域 $V_i$ 使得 $V_i^{-1} \subset U_i$. 设 $V = \bigl(\prod_{i \in J} V_i\bigr) \times \bigl(\prod_{i \in I-J} G_i\bigr)$. 则 $V$ 是 $(x_i)$ 在 $\prod_{i \in I} G_i$ 中的开邻域, 且有 $V^{-1} \subset U$.

(b). 只要有无穷多个 $G_i$ 不紧, 答案就是``否''. 事实上, 记 $p_j \colon \prod_{i \in I} G_i \to G_j$ 为投影映射. 这些映射连续, 所以把紧集映到紧集. 现在假设使得 $G_i$ 不紧的 $i \in I$ 的集合无穷. 若 $\prod_{i \in I} G_i$ 局部紧, 则其单位元有紧邻域 $K$. 由乘积拓扑的定义, $K$ 必包含形如 $\bigl(\prod_{i \in J} U_i\bigr) \times \bigl(\prod_{i \in I-J} G_i\bigr)$ 的集合 $U$, 其中 $J$ 是 $I$ 的有限子集, 且对每个 $i \in J$, $U_i$ 是 $e$ 在 $G_i$ 中的邻域. 由假设, 存在 $i \in I - J$ 使得 $G_i$ 不紧. 但 $G_i \supset p_i(K) \supset p_i(U) = G_i$, 所以 $G_i = p_i(K)$ 紧, 这荒谬.
反之, 假设存在 $I$ 的有限子集 $J$ 使得对每个 $i \in I - J$, $G_i$ 紧. 则 $\prod_{i \in I} G_i$ 局部紧. 事实上, 有 $\prod_{i \in I} G_i = \bigl(\prod_{i \in J} G_i\bigr) \times \bigl(\prod_{i \in I-J} G_i\bigr)$, 且由 Tychonoff 定理 $\prod_{i \in I-J} G_i$ 紧, 所以只需证明 $\prod_{i \in J} G_i$ 局部紧. 换句话说, 我们可以假设 $I$ 有限. 但此时, 若 $(x_i) \in \prod_{i \in I} G_i$, 且对每个 $i \in I$, $K_i$ 是 $x_i$ 的紧邻域, 则乘积 $\prod_{i \in I} K_i$ 是 $(x_i)$ 的紧邻域.
\end{solution}

\begin{exercise}[投射极限 (projective limit)]
\label{ex:I.5.1.3}
设 $(I, \leq)$ 是一个有序集. 考虑一族集合 $(X_i)_{i \in I}$ 和一族映射 $(u_{ij} \colon X_i \to X_j)_{i \geq j}$ 满足:
\begin{itemize}
\item 对每个 $i \in I$, $u_{ii} = \mathrm{id}_{X_i}$;
\item 对所有 $i \geq j \geq k$, $u_{ik} = u_{ij} \circ u_{jk}$.
\end{itemize}
这称为由有序集 $I$ 指标的集合的投射系统 (projective system). 此投射系统的投射极限 (projective limit) 是 $\prod_{i \in I} X_i$ 的如下子集:
\[
\varprojlim_{i \in I} X_i = \Bigl\{(x_i)_{i \in I} \in \prod_{i \in I} X_i \;\Big|\; \forall i, j \in I \text{ 且 } i \geq j,\, u_{ij}(x_i) = x_j\Bigr\}.
\]
\begin{enumerate}
\item[(a)] 若所有 $X_i$ 都是 Hausdorff 拓扑空间且所有 $u_{ij}$ 都连续, 证明 $\varprojlim_{i \in I} X_i$ 是 $\prod_{i \in I} X_i$ 的闭子集. 从现在起, 我们总在 $\varprojlim_{i \in I} X_i$ 上赋予诱导拓扑.
\item[(b)] 若所有 $X_i$ 都是紧 Hausdorff 拓扑空间且所有 $u_{ij}$ 都连续, 证明 $\varprojlim_{i \in I} X_i$ 也紧 Hausdorff. (提示: Tychonoff 定理.)
\item[(c)] 若所有 $X_i$ 都是群 (相应地环) 且所有 $u_{ij}$ 都是群 (相应地环) 同态, 证明 $\varprojlim_{i \in I} X_i$ 是 $\prod_{i \in I} X_i$ 的子群 (相应地子环).
\item[(d)] 若所有 $X_i$ 都是拓扑群且所有 $u_{ij}$ 都是连续群同态, 证明 $\varprojlim_{i \in I} X_i$ 是拓扑群.
\item[(e)] 设 $p$ 是一个素数. 取 $I = \N$ (通常序), $X_n = \Z/p^n\Z$, $u_{nm} \colon \Z/p^n\Z \to \Z/p^m\Z$ 是模 $p^m$ 的约化映射. 证明 $\Z_p := \varprojlim_{i \in I} X_i$ 是环, 且对加法是紧拓扑群.
\end{enumerate}
\end{exercise}

\begin{solution}
(a). 记 $X = \prod_{i \in I} X_i$, $X' = \varprojlim_{i \in I} X_i$. 对每个 $i \in I$, 设 $p_i \colon X \to X_i$ 是投影; 这是连续映射. 因此, 若 $i, j \in I$ 满足 $i \geq j$, 则 $X$ 的子集 $\{x \in X \mid u_{ij} \circ p_i(x) = p_j(x)\}$ 闭 (因为它是连续映射 $(u_{ij} \circ p_i, p_j) \colon X \to X_j \times X_j$ 下对角线的原像, 而对角线因 $X_j$ Hausdorff 而闭). 但由投射极限的定义,
\[
X' = \bigcap_{i,j \in I,\, i \geq j}\{x \in X \mid u_{ij} \circ p_i(x) = p_j(x)\}.
\]
所以 $X'$ 也闭.

(b). 若所有 $X_i$ 都是紧 Hausdorff 拓扑空间, 则由 Tychonoff 定理 $X$ 紧 Hausdorff. 由 (a), 投射极限是 $X$ 的闭子空间, 所以它也紧 Hausdorff.

(c). 我们沿用 (a) 的记号. 则所有投影 $p_i$ 都是群 (相应地环) 同态, 所以对所有满足 $i \geq j$ 的 $i, j \in I$, $X$ 的子集 $\{x \in X \mid u_{ij} \circ p_i(x) = p_j(x)\}$ 是子群 (相应地子环). 由投射极限的定义,
\[
X' = \bigcap_{i,j \in I,\, i \geq j}\{x \in X \mid u_{ij} \circ p_i(x) = p_j(x)\}.
\]
所以 $X'$ 也是子群 (相应地子环).

(d). 由习题 \ref{ex:I.5.1.2}(a), 直积是拓扑群. 由 (c), 投射极限 $X'$ 是 $X$ 的子群. 因此 $X'$ 是拓扑群.

(e). 集合 $\Z_p$ 由 (c) 是环, 由 (d) 是拓扑群. 由 (b) 它紧 (注意有限集赋予离散拓扑是紧 Hausdorff 的).
\end{solution}

\begin{exercise}[$p$-进数域 $\Q_p$]
\label{ex:I.5.1.4}
设 $p$ 是一个素数. 我们在 $\Q$ 上如下定义 $p$-进范数 (p-adic norm) $|\cdot|_p$:
\begin{itemize}
\item $|0|_p = 0$;
\item 若 $x$ 是非零有理数, 把 $x$ 写成 $x = p^ny$, 其中 $y$ 的分子和分母都与 $p$ 互素的有理数, 并令 $|x|_p = p^{-n}$.
\end{itemize}
\begin{enumerate}
\item[(a)] 证明对每个 $x, y \in \Q$ 有:
  \begin{itemize}
  \item $|x + y|_p \leq \max(|x|_p, |y|_p)$, 且当 $|x|_p \neq |y|_p$ 时等号成立;
  \item $|xy|_p = |x|_p|y|_p$.
  \end{itemize}
特别地, $p$-进距离函数 $d(x, y) = |x - y|_p$ 是 $\Q$ 上的度量. 我们记 $\Q_p$ 为 $\Q$ 对此度量的完备化.
\item[(a)] 证明 $p$-进范数 $|\cdot|_p$、$\Q$ 的加法和乘法通过连续性延拓到 $\Q_p$, 且 $\Q_p$ 是域 (称为 $p$-进数域 (field of p-adic numbers)), (a) 的结论延拓到 $\Q_p$.
\item[(b)] 证明 $\Q_p$ 的加法群是拓扑群.
\item[(c)] 计算 $\R$ 的子集 $|\Q_p|_p$.
\item[(d)] 证明 $\Q_p$ 中的每个开球也是闭球, 且 $\Q_p$ 中每个正半径的闭球也是开球.
\item[(e)] 证明 $\Q_p$ 完全不连通 (totally disconnected) (即其唯一的非空连通子集是单点集) 但不离散.
\item[(f)] 证明级数 $\sum_{n \geq 0}x_n$ 收敛当且仅当 $\lim_{n \to +\infty}|x_n|_p = 0$.
\item[(g)] 若 $m \in \Z$ 且 $(c_n)_{n \geq m}$ 是一族整数, 证明级数 $\sum_{n \geq m}c_np^n$ 在 $\Q_p$ 中收敛, 且其 $p$-进绝对值 $\leq p^{-m}$, 当 $c_m$ 与 $p$ 互素时等号成立.
\item[(h)] 设 $x \in \Q_p - \{0\}$. 证明存在唯一的 $m \in \Z$ 和唯一的 $\{0, 1, \ldots, p-1\}$ 的元素族 $(c_n)_{n \geq m}$ 使得 $x_m \neq 0$ 且 $x = \sum_{n \geq m}c_np^n$, 且 $|x|_p = p^{-m}$.
\item[(i)] 设 $B = \{x \in \Q_p \mid \|x\|_p \leq 1\}$. 证明这是 $\Q_p$ 的子环, 且是 $\Z$ 在 $\Q_p$ 中的闭包.
\item[(j)] 我们如下定义从 $B$ 到 $\prod_{n \geq 0}\Z/p^n\Z$ 的映射 $u$: 若 $x \in B$, 则由 (e) 可以找到收敛到 $x$ 的 $\Z$ 的元素 Cauchy 序列 $(x_n)_{n \geq 0}$. 取子序列后, 可以假设对每个 $n$ 有 $|x - x_n|_p \leq p^{-n}$. 我们令 $u(x) = (x_n \bmod p^n\Z)_{n \geq 0}$.
  证明 $u$ 是良定义的, 是从 $B$ 到 $\Z_p$ 的同胚, 且也是环同态. 我们将用此来等同 $B$ 与 $\Z_p$.
\item[(k)] 我们把 $M_n(\Q_p)$ 等同于 $\Q_p^{n^2}$, 在其上赋予乘积拓扑, 并在 $\GL_n(\Q_p)$ 上使用诱导拓扑. 证明 $\GL_n(\Q_p)$ 是局部紧拓扑群.
\item[(l)] 证明 $\GL_n(\Z_p)$ 是 $\GL_n(\Q_p)$ 的开紧子群. (提示: 证明 $\Z_p^{\times}$ 在 $\Z_p$ 中闭.)
\end{enumerate}
\end{exercise}

\begin{solution}
(a). 我们首先注意, 若 $x \in \Z - \{0\}$, 则可以把 $x$ 写成 $x = p^mx'$, 其中 $m \geq 0$ 且 $x' \in \Z$ 与 $p$ 互素, 所以 $|x|_p = p^{-m} \leq 1$. 当然, 若 $x = 0$, 也有 $|x|_p \leq 1$.
又注意由 $|\cdot|_p$ 的定义立即得出: 若 $x \in \Q^{\times}$, 则 $|x^{-1}|_p = |x|_p^{-1}$.

设 $x, y \in \Q_p$. 若 $x = 0$, 则 $x + y = y$ 且 $xy = 0$, 所以两点都显然; $y = 0$ 的情形类似. 所以我们假设 $x$ 和 $y$ 都非零, 并把 $x$ 写成 $x = p^nx'$, $y = p^my'$, 其中 $x'$ 和 $y'$ 是分子分母都与 $p$ 互素的有理数. 则 $x'y'$ 的分子分母也都与 $p$ 互素, 且 $xy = p^{n+m}x'y'$, 所以
\[
|xy|_p = p^{-n-m} = p^{-n}p^{-m} = |x|_p|y|_p.
\]
为证明第一个恒等式, 注意由于该恒等式对 $x$ 和 $y$ 对称, 可以假设 $n \leq m$. (注意此时 $p^{-n} = |x|_p = \max(|x|_p, |y|_p)$.) 把 $x'$ 和 $y'$ 写成 $x' = \frac{a}{b}$, $y' = \frac{c}{d}$, 其中 $a, b, c, d \in \Z$ 都与 $p$ 互素. 则
\[
x + y = p^n(x' + p^{m-n}y) = p^n\frac{ad + p^{m-n}cb}{bd},
\]
因此, 由我们已经证明的结论,
\[
|x + y|_p = |p^n|_p|ad + p^{m-n}bc|_p|bd|_p^{-1}.
\]
由于 $bd$ 与 $p$ 互素, 有 $|bd|_p = 1$ (由 $|\cdot|_p$ 的定义). 由于 $ad + p^{m-n}bc \in \Z$, 由开头的注记有 $|ad + p^{m-n}cb|_p \leq 1$. 最后得到
\[
|x + y|_p \leq |p^n|_p = p^{-n} = \max(|x|_p, |y|_p).
\]
最后, 若 $|x|_p \neq |y|_p$, 则 $n < m$, 所以 $ad + p^{m-n}cb$ 与 $p$ 互素, 由 $|\cdot|_p$ 的定义得 $|x + y|_p = p^{-n}$.

(b). 由定义, $\Q_p$ 是 $\Q$ 的元素 Cauchy 序列 $(x_n)_{n \geq 0}$ (对 $p$-进距离给出的度量) 商去等价关系 $\sim$ 的商集, 其中 $\sim$ 定义为: $(x_n)_{n \geq 0} \sim (y_n)_{n \geq 0}$ 当且仅当 $n \to +\infty$ 时 $|x_n - y_n|_p \to 0$.
注意 (a) 的第二个恒等式蕴含三角不等式: 对所有 $x, y \in \Q$, 有 $|x + y|_p \leq |x|_p + |y|_p$.
设 $x \in \Q_p$, $(x_n)_{n \geq 0}$ 是表示 $x$ 的 Cauchy 序列. 由三角不等式, 对所有 $n, m \in \N$ 有 $||x_n|_p - |x_m|_p| \leq |x_n - x_m|_p$. 所以 $(|x_n|_p)_{n \geq 0}$ 是 $\R$ 中的 Cauchy 序列, 且因 $\R$ 完备它有极限. 设 $(y_n)_{n \geq 0}$ 是表示 $x$ 的另一个 Cauchy 序列. 由三角不等式, 对每个 $n \geq 0$ 有 $||x_n|_p - |y_n|_p| \leq |x_n - y_n|_p$, 所以 $(|x_n|_p)_{n \geq 0}$ 和 $(|y_n|_p)_{n \geq 0}$ 的极限相等. 因此我们可以定义 $|x|_p = \lim_{n \to +\infty}|x_n|_p$.
现在设 $x, y \in \Q_p$, 并选取表示 $x$ 和 $y$ 的 Cauchy 序列 $(x_n)_{n \geq 0}$ 和 $(y_n)_{n \geq 0}$. 首先注意序列 $(|x_n|_p)_{n \geq 0}$ 和 $(|y_n|_p)_{n \geq 0}$ 有界 (例如因为它们如上所述收敛). 现在利用 (a), 对所有 $n, m \in \N$:
\[
|(x_n + y_n) - (x_m + y_m)|_p \leq \max(|x_n - x_m|_p, |y_n - y_m|_p)
\]
且
\[
|x_ny_n - x_my_m|_p = |x_n(y_n - y_m) + (x_n - x_m)y_m|_p \leq \max(|x_n|_p|y_n - y_m|_p, |y_m|_p|x_n - x_m|_p).
\]
所以序列 $(x_n + y_n)_{n \geq 0}$ 和 $(x_ny_n)_{n \geq 0}$ 是 Cauchy 序列 (对第二个, 用 $(|x_n|_p)_{n \geq 0}$ 和 $(|y_n|_p)_{n \geq 0}$ 有界的事实), 因此它们表示 $\Q_p$ 的元素. 我们想把这些元素称为 $x + y$ 和 $xy$, 但首先必须验证它们与表示 $x$ 和 $y$ 的 Cauchy 序列的选取无关. 所以设 $(x'_n)_{n \geq 0}$ 和 $(y'_n)_{n \geq 0}$ 是分别表示 $x$ 和 $y$ 的另外两个 Cauchy 序列. 则对每个 $n \geq 0$,
\[
|(x_n + y_n) - (x'_n + y'_n)|_p \leq \max(|x_n - x'_n|_p, |y_n - y'_n|_p)
\]
且
\[
|x_ny_n - x'_ny'_n|_p = |x_n(y_n - y'_n) + (x_n - x'_n)y'_n|_p \leq \max(|x_n|_p|y_n - y'_n|_p, |y'_n|_p|x_n - x'_n|_p).
\]
所以两个序列 $((x_n + y_n) - (x'_n + y'_n))_{n \geq 0}$ 和 $(x_ny_n - x'_ny'_n)_{n \geq 0}$ 都收敛到 $0$, 这意味着序列 $(x_n + y_n)_{n \geq 0}$ 与 $(x'_n + y'_n)_{n \geq 0}$ (相应地 $(x_ny_n)_{n \geq 0}$ 与 $(x'_ny'_n)_{n \geq 0}$) 有相同的极限, 所以 $x + y$ 和 $xy$ 的定义有意义.
$\Q_p$ 的环公理由运算的定义立即得出. 我们来验证 $\Q_p$ 是域. 设 $x \in \Q_p - \{0\}$, 选取表示 $x$ 的 Cauchy 序列 $(x_n)_{n \geq 0}$. 由于 $x \neq 0$, 序列 $(|x_n|_p)_{n \geq 0}$ 不能收敛到 $0$, 所以其极限 (即 $|x|_p$) 非零. 所以对足够大的 $n, m$ 有 $|x_n - x_m|_p \leq |x|_p/2$, 且经过用等价的 Cauchy 序列替换 $(x_n)_{n \geq 0}$ 后, 可以假设它对所有 $n, m \geq 0$ 成立. 设 $n, m \geq 0$. 由 (a), 有 $|x_n|_p \leq \max(|x_m|_p, |x_n - x_m|_p)$. 令 $m \to +\infty$ 取极限, 得 $|x_n|_p \leq |x|_p$. 类似地, 令 $n \to +\infty$ 取极限, 并利用 $|x|_p > |x|_p/2 \geq \lim_{n \to +\infty}|x_n - x_m|_p$, 得 $|x|_p \leq |x_m|_p$. 这蕴含对每个 $n \geq 0$ 有 $|x|_p = |x_n|_p$. 现在, 若我们能证明序列 $(x_n^{-1})_{n \geq 0}$ 是 Cauchy 序列, 则它表示的 $\Q_p$ 的元素显然将是 $x$ 的逆. 但对所有 $n, m \geq 0$,
\[
|x_n^{-1} - x_m^{-1}|_p = |x_n^{-1}x_m^{-1}|_p|x_m - x_n|_p = |x|_p^{-2}|x_m - x_n|_p,
\]
所以 $|x_n^{-1} - x_m^{-1}|_p$ 在 $n, m \to +\infty$ 时确实收敛到 $0$.

我们最后证明 (a) 的恒等式在 $\Q_p$ 中仍成立. 若 $x, y \in \Q_p - \{0\}$, 则刚才已见可以找到收敛到 $x$ 和 $y$ 的 Cauchy 序列 $(x_n)_{n \geq 0}$ 和 $(y_n)_{n \geq 0}$ 使得对每个 $n \geq 0$, $|x|_p = |x_n|_p$ 且 $|y|_p = |y_n|_p$. 当然, 若 $x$ 或 $y$ 为 $0$ 这也成立. 则恒等式由 (a) 和 $\Q_p$ 上运算的定义立即得出.

(c). $\Q_p$ 上的加法由定义连续 (它定义为映射的连续性延拓). 取逆映射 $x \longmapsto -x$ 连续因为 $|x|_p = |-x|_p$.

(d). 我们已见: 若 $x \in \Q_p$, 则存在收敛到 $x$ 的 $\Q$ 的 Cauchy 序列且对每个 $n \geq 0$ 有 $|x|_p = |x_n|_p$. 所以 $|\Q_p|_p = |\Q|_p = \{0\} \cup p^{\Z}$.

(e). 若 $x \in \Q_p$, $r \in \R$, 记 $B(x, r)$ 和 $\overline{B}(x, r)$ 为以 $x$ 为心、$r$ 为半径的开球和闭球.
设 $x \in \Q_p$, $r \in \R$. 若 $r \leq 0$, 则 $B(x, r) = \emptyset = \overline{B}(x, -1)$. 假设 $r > 0$, 设 $n$ 是满足 $p^n < r \leq p^{n+1}$ 的唯一整数. 由前一问, 若 $a \in \Q_p$ 满足 $|a|_p < r$, 则 $|a|_p \leq p^n$, 且显然其逆也成立. 所以 $B(x, r) = \overline{B}(x, p^n)$. 现在设 $m$ 是满足 $p^m \leq r < p^{m+1}$ 的唯一整数. 则类似地有 $B(x, r) = \overline{B}(x, p^{m+1})$.

(f). 设 $x, y \in \Q_p$ 且 $x \neq y$. 则 $|x - y|_p > 0$, 所以可以选取 $r > 0$ 使得 $r < |x - y|_p$. 则 $B(x, r)$ 是 $\Q_p$ 的包含 $x$ 而不包含 $y$ 的开闭子集, 所以 $x$ 和 $y$ 不能在 $\Q_p$ 的同一连通子集中. 这证明 $\Q_p$ 完全不连通.
为证明 $\Q_p$ 不离散, 只需证明其子集 $\{0\}$ 不开. 这由序列 $(p^n)_{n \geq 0}$ 在 $\Q_p$ 中收敛到 $0$ 且对每个 $n \in \Z$ 有 $p^n \neq 0$ 的事实得出.

(g). 定义有理数序列 $(S_n)_{n \geq 0}$ 为 $S_n = \sum_{i=0}^n x_i$. 则级数 $\sum_{n \geq 0}x_n$ 收敛当且仅当序列 $(S_n)_{n \geq 0}$ 收敛. 特别地, 若级数收敛, 则 $|x_n|_p = |S_{n+1} - S_n|_p$ 在 $n \to +\infty$ 时趋于 $0$.
反之, 假设 $\lim_{n \to +\infty}|x_n|_p = 0$. 对所有 $n, n' \in \N$, 若 $n \leq n'$, 则 (利用 (b)):
\[
|S_{n'} - S_n|_p = \Bigl|\sum_{i=n+1}^{n'}x_i\Bigr|_p \leq \max_{n+1 \leq i \leq n'}|x_i|_p \leq \sup_{i \geq n+1}|x_i|_p.
\]
当 $n \to +\infty$ 时这趋于 $0$, 所以 $(S_n)_{n \geq 0}$ 是 Cauchy 序列, 因此它在 $\Q_p$ 中收敛, 级数 $\sum_{n \geq 0}x_n$ 也随之收敛.

(h). 收敛性由前一问以及 (b) 的证明中注意到的 $|c|_p \leq 1$ (对每个 $c \in \Z$) 得出. 设 $x = \sum_{n \geq m}c_np^n$. 由定义,
\[
x = \lim_{n \to +\infty}\sum_{i=m}^n c_ip^i.
\]
对每个 $n \geq m$,
\[
\Bigl|\sum_{i=m}^n c_ip^i\Bigr|_p \leq \max_{m \leq i \leq n}|c_i|_p|p^i|_p = p^{-m},
\]
所以 $|x|_p \leq p^{-m}$. 假设 $c_m$ 与 $p$ 互素; 则 $|c_m|_p = 1$. 因此 $|c_mp^m|_p = p^{-m} > |c_ip^i|_p$ (对所有 $i > m$), 所以再次利用 (b), 对每个 $n \geq m$,
\[
\Bigl|\sum_{i=m}^n c_ip^i\Bigr|_p = p^{-m}.
\]
这给出 $|x|_p = p^{-m}$.

(i). 让我们证明存在性. 可以假设 $x \neq 0$ (否则结论平凡). 我们知道对某个 $m \in \Z$ 有 $|x|_p = p^{-m}$. 选取收敛到 $x$ 的 Cauchy 序列 $(x_n)_{n \geq 0}$. 取子序列后, 可以假设对每个 $n \geq 0$ 有 $|x - x_n|_p < p^{-n}$.
设 $n \geq 0$. 把 $x_n$ 写成 $p$ 进制形式 $x_n = \sum_{i=a_n}^{b_n}c_{i,n}p^i$, 其中 $a_n, b_n \in \Z$, $c_{i,n} \in \{0, 1, \ldots, p-1\}$. 可以假设 $c_{a_n,n} \neq 0$. 则 $c_{a_n,n}$ 与 $p$ 互素, 所以 $|c_{a_n,n}p^{a_n}|_p = p^{-a_n} > |c_{i,n}p^i|_p$ (对所有 $i > a_n$), 于是由 (b),
\[
p^{-m} = |x_n|_p = p^{-a_n},
\]
最终得到 $m = a_n$.
又, 通过在 $i > b_n$ 时令 $c_{i,n} = 0$, 可以把 $b_n$ 换成 $+\infty$.
设 $n, n' \in \N$, 并假设 $n \geq n'$. 则 $|x_n - x_{n'}|_p \leq \max(|x_n - x|_p, |x - x_{n'}|_p) < p^{-n'}$. 另一方面, $x_n - x_{n'} = \sum_{i \geq m}(c_{i,n} - c_{i,n'})p^i$. 注意 $c_{i,n} - c_{i,n'}$ 属于 $\{1 - p, \ldots, p - 1\}$, 所以它们要么是 $0$ 要么与 $p$ 互素. 由 (h), 这蕴含 $|x_n - x_{n'}|_p = p^{-r}$, 其中 $r$ 是使得 $c_{r,n} - c_{r,n'} \neq 0$ 的最小整数. 这又蕴含 $n' < r$, 即对 $m \leq i \leq n'$ 有 $c_{i,n} = c_{i,n'}$.
我们现在对 $i \geq m$ 定义整数 $c_i$: 若 $i \leq 0$ 取 $c_i = c_{0,i}$, 若 $i \geq 0$ 取 $c_i = c_{i,i}$. 由上一段, 若 $0 \leq i \leq n$ 则 $c_i = c_{i,n}$. 对每个 $n \geq 0$, 令 $y_n = \sum_{i=m}^n c_ip^i$. 则
\[
x_n - y_n = \sum_{i \geq m}(c_{i,n} - c_i)p^i = \sum_{i \geq n+1}(c_{i,n} - c_i)p^i,
\]
所以由 (h) 有 $|x_n - y_n|_p \leq p^{-n-1}$. 因此序列 $(y_n)_{n \geq 0}$ 也收敛到 $x$, 这证明 $x = \sum_{i \geq m}c_ip^i$.
让我们证明唯一性. 假设有两个整数序列 $(c_n)_{n \geq m}$, $(d_n)_{n \geq m}$ 使得 $x = \sum_{n \geq m}c_np^n = \sum_{n \geq m}d_np^n$ 且对每个 $n$ 有 $c_n, d_n \in \{0, \ldots, p-1\}$. 则 $0 = \sum_{n \geq m}(c_n - d_n)p^n$. 又, 对每个 $n$, $c_n - d_n$ 属于 $\{1 - p, \ldots, p - 1\}$, 所以它是 $0$ 或与 $p$ 互素. 若存在 $n$ 使得 $c_n - d_n \neq 0$, 则由 (h) 这将蕴含 $|0|_p \neq 0$, 这不可能. 所以对每个 $n$ 有 $c_n = d_n$.

(j). $B$ 是子环这一事实由 (b) (以及 $|-x|_p = |x|_p$, 这在定义上显然) 得出, 且已在 (b) 的证明中看到 $\Z \subset B$. 又, $B$ 是闭球, 所以它在 $\Q_p$ 中闭, 因此包含 $\Z$ 的闭包.
设 $x \in B$. 由 (i), 可以把 $x$ 写成 $x = \sum_{n \geq 0}c_np^n$, 其中 $c_n \in \{0, \ldots, p-1\}$. 这意味着 $x$ 是整数序列 $\bigl(\sum_{i=0}^n c_ip^i\bigr)_{n \geq 0}$ 的极限, 所以 $x$ 在 $\Z$ 于 $\Q_p$ 中的闭包里.

(k). 我们证明 $u$ 良定义. 设 $x \in B$, $(x_n)_{n \geq 0}$, $(x'_n)_{n \geq 0}$ 是如题设的两个序列. 设 $n \geq 0$. 则 $|x_n - x'_n|_p \leq \max(|x_n - x|_p, |x - x'_n|_p) \leq p^{-n}$, 这意味着 $p^n$ 整除 $x_n - x'_n$, 所以 $x_n$ 和 $x'_n$ 在 $\Z/p^n\Z$ 中有相同的像. 这证明 $u(x)$ 良定义.
$u$ 是环同态这一事实由 $\Q_p$ 上环运算的定义以及对每个 $n$ 从 $\Z$ 到 $\Z/p^n\Z$ 的约化是环同态这一事实立即得出.
我们证明 $u$ 单射. 设 $x, y \in B$ 满足 $u(x) = u(y)$, 并选取满足题设条件的收敛到 $x$, $y$ 的序列 $(x_n)_{n \geq 0}$, $(y_n)_{n \geq 0}$. 则对每个 $n \geq 0$ 有 $x_n \equiv y_n \pmod{p^n}$, 所以 $p^n$ 整除 $x_n - y_n$, 即 $|x_n - y_n|_p \leq p^{-n}$. 令 $n \to +\infty$ 取极限, 得 $|x - y|_p = 0$. 但已在 (b) 中看到 $\Q_p$ 中 $p$-进范数为 $0$ 的唯一元素是 $0$, 所以 $x = y$.
我们证明 $u$ 满射. 设 $(x_n + p^n\Z)_{n \geq 0}$ 是 $\Z_p$ 的元素. 对每个 $n \geq 0$, 在 $\{0, \ldots, p^n - 1\}$ 中选取 $x_n + p^n\Z$ 的代表元, 记为 $x_n$. 又把 $x_n$ 写成 $p$ 进制形式 $x_n = \sum_{i=0}^{n-1}c_{i,n}p^i$, 其中 $0 \leq c_{i,n} \leq p - 1$. 设 $m \geq n$. 我们知道 $x_m \equiv x_n \pmod{p^n}$, 所以对 $0 \leq i \leq n$ 有 $c_{i,m} = c_{i,n}$. 我们定义序列 $(c_i)_{i \geq 0}$ 为 $c_i = c_{i,0} = c_{i,1} = \cdots = c_{i,i}$, 并令 $x = \sum_{i \geq 0}c_ip^i \in \Q_p$. 则由 (c), $x \in B$, 且容易验证 $(x_n)_{n \geq 0}$ 是收敛到 $x$ 且满足题设条件的整数 Cauchy 序列. 因此 $u(x) = (x_n + p^n\Z)_{n \geq 0}$.
我们证明 $u$ 连续. $\Z_p$ 中的每个开集都是形如 $\Z_p \cap \bigl((\prod_{n \geq m+1}\Z/p^n\Z) \times \{(x_m, \ldots, x_0)\}\bigr)$ 的开集的并, 其中 $m \geq 0$, 且对 $0 \leq i \leq m$ 有 $x_i \in \Z/p^i\Z$. 所以只需证明这种形式的集合的原像在 $B$ 中开. 记 $A = \bigl(\prod_{n \geq m+1}\Z/p^n\Z\bigr) \times \{(x_m, \ldots, x_0)\}$. 选取 $x \in \Z$ 使得 $x \equiv x_m \pmod{p^m}$. 则对 $0 \leq n \leq m$ 有 $x \equiv x_n \pmod{p^n}$ (因为 $x_n \equiv x_m \pmod{p^n}$). 我们将证明 $y \in B$ 属于 $u^{-1}(A)$ 当且仅当 $|x - y|_p < p^{-m+1}$, 这证明 $u^{-1}(A)$ 开. 首先注意, 由于 $|\cdot|_p$ 的值总是 $0$ 或 $p$ 的整数次幂, 条件 $|x - y|_p < p^{-m+1}$ 等价于 $|x - y|_p \leq p^{-m}$.\footnote{所以在 $\Q_p$ 中, 每个开球都是闭球, 反之亦然.} 设 $y \in B$, $(y_n)_{n \geq 0}$ 是如 $u$ 的定义中那样收敛到 $y$ 的 Cauchy 序列. 假设 $|x - y|_p \leq p^{-m}$. 则对 $n \in \{0, \ldots, m\}$, 有 $|y_n - x|_p \leq \max(|y_n - y|_p, |y - x|_p) \leq p^{-n}$, 所以 $y_n \equiv x \equiv x_n \pmod{p^n}$. 所以 $u(y) \in A$. 反之, 假设 $u(y) \in A$. 则 $y_m \equiv x_m \equiv x \pmod{p^m}$, 所以 $|y_m - x|_p \leq p^{-m}$. 由于 $|y - y_m|_p \leq p^{-m}$, 这蕴含 $|x - y|_p \leq p^{-m}$.
最后, 我们证明 $u$ 是开映射. 由于 $u$ 双射, 只需证明开球的像是开的. 我们几乎已经做了这件事: 设 $x \in B$, $r \in \R_{>0}$, $A'$ 是以 $x$ 为心、$r$ 为半径的开球. 若 $m$ 是满足 $p^{-m} < r$ 的最小整数, 则 $A'$ 也是以 $x$ 为心、$p^{-m}$ 为半径的闭球 (因为 $|\cdot|_p$ 取值于 $\{0\} \cup p^{\Z}$). 设 $y$ 是满足 $|x - y|_p < p^{-m}$ 的整数. 则 (a) 的第二个恒等式蕴含对 $z \in \Q_p$, $|x - z|_p \leq p^{-m}$ 当且仅当 $|y - z|_p \leq p^{-m}$, 这意味着在 $A'$ 的定义中可以用 $y$ 代替 $x$. 则上面已见 $u(A') = \bigl(\prod_{n \geq m+1}\Z/p^n\Z\bigr) \times \{(x_m, \ldots, x_0)\}$, 其中对 $0 \leq n \leq m$ 有 $x_n = y + p^n\Z$.

(l). $\GL_n(\Q_p)$ 是拓扑群的证明与习题 \ref{ex:I.5.1.2}(b) 相同 (有限维情形). 它也在 $M_n(\Q_p)$ 中开, 因为它是连续函数 $\det$ 下 $\Q_p$ 的开子集 $\Q_p^{\times}$ 的原像. 所以要证明 $\GL_n(\Q_p)$ 局部紧, 只需证明 $M_n(\Q_p)$ 局部紧, 而这将由 $\Q_p$ 局部紧得出. 但对每个 $x \in \Q_p$, 以 $x$ 为心、半径为 $1$ 的闭球 $x + \Z_p$ 是 $x$ 的紧邻域: 它紧因为 $\Z_p$ 紧且平移 $x$ 连续 (由度量的定义), 它开因为它等于以 $x$ 为心、$p$ 为半径的开球.

(m). 回忆 $\GL_n(\Z_p)$ 是 $M_n(\Z_p)$ 的可逆元群, 所以
\[
\GL_n(\Z_p) = \{A \in \GL_n(\Q_p) \mid A \text{ 与 } A^{-1} \in M_n(\Z_p)\}.
\]
换句话说, $\GL_n(\Z_p)$ 是连续映射 $\GL_n(\Q_p) \to M_n(\Q_p) \times M_n(\Q_p)$, $A \longmapsto (A, A^{-1})$ 下 $M_n(\Q_p) \times M_n(\Q_p)$ 的开子集 $M_n(\Z_p) \times M_n(\Z_p)$ 的原像, 所以它在 $\GL_n(\Q_p)$ 中开. 由于 $M_n(\Z_p)$ 紧 (因为 $\Z_p$ 紧), 为证明 $\GL_n(\Z_p)$ 紧, 只需证明它在 $M_n(\Z_p)$ 中闭. 由于它是连续映射 $\det \colon M_n(\Z_p) \to \Z_p$ 下 $\Z_p^{\times}$ 的原像, 只需证明 $\Z_p^{\times}$ 在 $\Z_p$ 中闭. 但 $\Z_p^{\times} = \{x \in \Z_p \mid |x|_p = 1\}$ (这蕴含它闭). 事实上, 设 $x \in \Z_p$. 若 $x$ 在 $\Z_p$ 中有逆, 则 $|x^{-1}|_p = |x|_p^{-1} \leq 1$, 所以 $|x|_p \geq 1$, 从而 $|x|_p = 1$. 反之, 若 $|x|_p = 1$, 则 $|x^{-1}|_p = 1$, 所以 $x^{-1} \in \Z_p$.
\end{solution}

\subsection{Van Dantzig 定理}
\label{sec:I.5.2}
本习题来自 Terry Tao 的博客.

\begin{exercise}[紧 Hausdorff 完全不连通空间]
\label{ex:I.5.2.1}
在此题中, $X$ 是紧 Hausdorff 完全不连通拓扑空间. (回忆``完全不连通 (totally disconnected)''意味着 $X$ 的唯一的非空连通子集是单点集.)
设 $x \in X$, 设 $A$ 是 $X$ 的所有包含 $x$ 的开闭子集的交集. 证明 $A = \{x\}$. (提示: 这等价于证明 $A$ 连通. 还要回忆在任何紧 Hausdorff 空间中, 不相交的闭集可被开集分离.)
\end{exercise}

\begin{solution}
注意 $A$ 在 $X$ 中闭. 假设 $A = A_1 \cup A_2$, 其中 $A_1$ 和 $A_2$ 在 $X$ 中闭且不相交, 且 $x \in A_1$. 由于任何紧 Hausdorff 空间正规, 可以找到 $X$ 的开子集 $U_1 \supset A_1$ 和 $U_2 \supset A_2$ 使得 $U_1 \cap U_2 = \emptyset$. 让我们找一个包含 $x$ 的开闭邻域 $V$ 使得 $V \cap \partial U_2 = \emptyset$. 对每个 $y \in \partial U_2$, 由于 $y \notin A$, 可以找到 $x$ 的开闭邻域 $V_y$ 使得 $y \notin V_y$. 注意 $X - V_y$ ($y \in \partial U_2$) 构成覆盖 $\partial U_2$ 的 $X$ 的开子集族; 由于 $\partial U_2$ 紧, 此族有仍然覆盖 $\partial U_2$ 的有限子族, 设为 $(X - V_{y_1}, \ldots, X - V_{y_n})$. 设 $V = V_{y_1} \cap \cdots \cap V_{y_n}$; 则 $V$ 仍开且闭, $x \in V$, 且 $V \cap \partial U_2 = \emptyset$. 最后这条性质蕴含 $B := V - U_2$ 也等于 $V - \overline{U_2}$, 所以它仍开且闭. 又, 有 $x \in B$ (因为 $x \notin U_2$), 且 $A_2 \cap B = \emptyset$. 但由定义 $A$ 必须包含于 $B$, 所以 $A_2 = \emptyset$. 这证明 $A$ 连通, 从而是单点集, 所以等于 $\{x\}$.
\end{solution}

\begin{exercise}[局部紧完全不连通群的紧开子群]
\label{ex:I.5.2.2}
在此题中, $G$ 是局部紧完全不连通拓扑群.
\begin{enumerate}
\item[(a)] 证明 $G$ 的单位元有紧开邻域 $K$.
\item[(b)] 证明存在包含于 $K$ 的开子群 $G_0$.
\item[(c)] 证明 $G$ 的紧开子群构成 $1$ 在 $G$ 中的邻域基.
\item[(d)] 设 $G$ 是习题 \ref{ex:I.5.1.4} 中的群 $\GL_n(\Q_p)$. 找出 $1$ 在 $G$ 中由紧开子群组成的邻域基.
\end{enumerate}
\end{exercise}

\begin{solution}
(a). 设 $V$ 是 $1$ 的紧邻域. 则 $\partial V$ 也紧且不含 $1$. 由习题 \ref{ex:I.5.2.1}, 对每个 $y \in \partial V$, 存在 $1$ 的开闭邻域 $B_y$ 使得 $B_y \cap \partial V = \emptyset$. 由于 $\partial V \subset \bigcup_{y \in \partial V}(X - B_y)$ 且 $X - B_y$ 开, 存在 $y_1, \ldots, y_n \in \partial V$ 使得 $\partial V \cap B = \emptyset$, 其中 $B = \bigcap_{i=1}^n B_{y_i}$. 注意 $B$ 仍开且闭, 且 $1 \in B$. 又, 由于 $\partial V \cap B = \emptyset$, 有 $B \cap V = B \cap \mathring{V}$, 所以 $K := B \cap V$ 开且紧 (因为它在 $V$ 中闭) 且包含 $1$.

(b). 设 $U$ 是 $1$ 的开对称邻域使得 $UK \subset K$, 设 $G_0$ 是由 $U$ 生成的 $G$ 的子群. 让我们证明 $G_0$ 是 $G$ 的开紧子群且 $G_0 \subset K$. 首先证明 $G_0$ 开. 设 $g \in G_0$; 则 $gU \subset G_0$ 且 $gU$ 在 $G$ 中开, 所以 $G_0$ 包含 $g$ 的邻域. 由于拓扑群的每个开子群也闭, 我们又得 $G_0$ 闭. 所以要证明它紧, 只需证明它包含于 $K$. 注意, 由于 $U$ 对称且含 $1$, 有 $G_0 = \bigcup_{n \geq 1}U^n$. 由于 $U \subset K$ (因为 $1 \in K$) 且 $UK \subset K$, 简单归纳证明对每个 $n \geq 1$ 有 $U^n \subset K$. 所以 $G_0 \subset K$.

(c). (a) 的解答中的论证实际上证明 $1$ 的每个紧邻域都包含 $1$ 的开紧邻域, 然后 (b) 蕴含它也包含 $G$ 的紧开子群. 因此, 由于 $G$ 局部紧, $1$ 在 $G$ 中的每个邻域都包含 $G$ 的紧开子群.

(d). 让我们在 $M_n(\Q_p)$ 上选取诱导乘积拓扑的范数. 例如, 由
\[
\|(a_{ij})_{1 \leq i,j \leq n} - (b_{ij})_{1 \leq i,j \leq n}\| = \sup_{1 \leq i,j \leq n}|a_{ij} - b_{ij}|_p
\]
定义的范数 $\|\cdot\|$ 即可. 对每个整数 $m \geq 1$, 设 $K_m = I_n + p^mM_n(\Z_p)$. 按我们选的范数, 这正是 $M_n(\Q_p)$ 中以 $I_n$ 为心、$p^{-m+1}$ 为半径的开球 (也是以 $I_n$ 为心、$p^{-m}$ 为半径的闭球). 特别地, 集合 $K_m$ ($m \geq 1$) 构成 $I_n$ 在 $M_n(\Q_p)$ 中的开邻域族, 从而 $K_m$ ($m \gg 0$) 构成 $I_n$ 在 $\GL_n(\Q_p)$ 中的开邻域族 (因为 $\GL_n(\Q_p)$ 在 $M_n(\Q_p)$ 中开, 它是连续映射 $\det$ 下 $\Q_p$ 的开子集 $\Q_p^{\times}$ 的原像).
又注意 $K_m$ 与 $M_n(\Z_p) \simeq \Z_p^{n^2}$ 同胚 (通过映射 $I_n + p^mX \longmapsto X$), 所以它也紧.
至此, 我们已经有了由紧开子群组成的邻域基. 实际上我们可以更精确, 证明对每个 $m \geq 1$ 有 $K_m \subset \GL_n(\Q_p)$ (而不只是对足够大的 $m$), 这只需 $K_1 \subset \GL_n(\Q_p)$. 事实上, 甚至有 $K_1 \subset \GL_n(\Z_p)$. 确实, 显然 $K_1 \subset M_n(\Z_p)$. 而且, 若 $X \in K_1$, 则容易看出 $\det(X) \in 1 + p\Z_p \subset \Q_p$, 由习题 \ref{ex:I.5.1.4}(a) 这蕴含 $|\det(X)|_p = 1$, 从而 $\det(X)^{-1}$ 也在 $\Z_p$ 中, 即 $\det(X) \in \Z_p^{\times}$.
\end{solution}

\subsection{Haar 测度的例子}

\begin{exercise}[仿射平移的群上的左 Haar 测度]
\label{ex:I.5.3.1}
设 $G$ 是拓扑群. 假设存在 $G$ 与某个 $\R^N$ 的开子集的同胚 (不要求与任何群结构相容), 使得 $G$ 上的左平移由仿射映射给出. 也就是说, 若把 $G$ 等同于它在 $\R^N$ 中的像 (只是作为拓扑空间!), 则对每个 $x \in G$, 存在 $N \times N$ 矩阵 $A(x) \in M_N(\R)$ 和元素 $b(x) \in \R^N$ 使得对每个 $y \in G$, $xy = A(x)y + b(x)$.
证明 $|\det A(x)|^{-1}dx$ 是 $G$ 上的左 Haar 测度, 其中 $dx$ 表示 $\R^N$ 上的 Lebesgue 测度. (提示: 换元公式. 另外, 首先证明 $x$ 唯一确定 $A(x)$ 和 $b(x)$, 且 $x \longmapsto A(x)$ 是从 $G$ 到 $\GL_N(\R)$ 的群同态.)
\end{exercise}

\begin{solution}
设 $x \in G$. 假设有 $A, A' \in M_N(\R)$ 和 $b, b' \in \R^N$ 使得对每个 $y \in G$ 有 $xy = Ay + b = A'y + b'$. 则对每个 $y \in G$ 有 $(A - A')y = b' - b$. 但方程 $(A - A')y = b' - b$ 的解集是 $\R^N$ 的仿射子空间 (即线性子空间的平移), 所以除非它等于 $\R^N$, 否则它有空的内部. 由于 $G$ 在 $\R^N$ 中开, 这意味着必有 $(A - A')y = b' - b$ 对所有 $y \in \R^N$ 成立. 这只有在 $\Ker(A - A') = \R^N$ 时才可能, 故 $A = A'$, 然后也有 $b = b'$. 所以 $x$ 确定 $A(x)$ 和 $b(x)$.
我们证明对每个 $x \in G$, $A(x)$ 可逆. 事实上, 若 $A(x)$ 不可逆, 则映射 $G \to G$, $y \longmapsto xy$ 的像包含于 $b(x) + \Im(A(x))$, 这是 $\R^N$ 的真线性子空间经 $b(x)$ 平移后的集, 所以它有空的内部. 但此像必须等于 $G$, 从而是 $\R^N$ 的开子集, 所以我们得到矛盾, 因此 $A(x)$ 必可逆.
我们证明 $x \longmapsto A(x)$ 是群同态. 设 $x_1, x_2 \in G$. 对每个 $y \in G$,
\[
A(x_1x_2)y + b(x_1x_2) = x_1x_2y = A(x_1)A(x_2)y + A(x_1)b(x_2) + b(x_1).
\]
由第一段, 这蕴含 $A(x_1x_2) = A(x_1)A(x_2)$ 且 $b(x_1x_2) = A(x_1)b(x_2) + b(x_1)$.

现在回忆换元公式蕴含: 若 $U$ 是 $\R^N$ 的可测子集, $A \in \GL_N(\R)$, $b \in \R^N$, 且 $V$ 是 $U$ 在变换 $y \longmapsto Ay + b$ 下的像, 则对每个 $f \in L^1(V)$,
\[
\int_V f(v)\,\dd v = |\det A|\int_U f(Ay + b)\,\dd y.
\]
将其应用于 $U = V = G$, $A = A(y)$, $b = b(y)$ (对某个 $y \in G$), 对每个 $f \in L^1(G)$ 得到
\[
\int_G f(x)\,\dd x = |\det A(y)|\int_G f(yx)\,\dd x.
\]
设 $f \in \mathcal{C}_c(G)$. 则函数 $x \longmapsto |\det A(x)|^{-1}f(x)$ 也在 $\mathcal{C}_c(G)$ 中, 所以可以把上一段应用于它. 对每个 $y \in G$ 得到
\[
\int_G f(yx)|\det A(yx)|^{-1}\,\dd x = |\det A(y)|\int_G f(x)|\det A(x)|^{-1}\,\dd x.
\]
利用 $\det(A(xy)) = \det(A(x))\det(A(y))$, 两边除以 $|\det(A(y))|^{-1}$ 得
\[
\int_G f(yx)|\det A(x)|^{-1}\,\dd x = \int_G f(x)|\det A(x)|^{-1}\,\dd x,
\]
这就是所要的结果.
\end{solution}

\begin{exercise}[经典李群的 Haar 测度]
\label{ex:I.5.3.2}
在此题中, $dx$ 总表示 $\R$ 上的 Lebesgue 测度.
\begin{enumerate}
\item[(a)] 证明 $\frac{dx}{|x|}$ 是乘法群 $\R^{\times}$ 上的 Haar 测度.
\item[(b)] 证明 $\frac{dxdy}{x^2 + y^2}$ 是乘法群 $\CC^{\times}$ (坐标为 $z = x + iy$) 上的 Haar 测度.
\item[(c)] 设 $dT$ 是 $M_n(\R)$ 上的 Lebesgue 测度. 证明 $|\det T|^{-n}dT$ 是 $\GL_n(\R)$ 上的左、右 Haar 测度.
\item[(d)] 设 $G = \Bigl\{\begin{pmatrix} x & y \\ 0 & z \end{pmatrix} \;\Big|\; x, z \in \R^{\times},\, y \in \R\Bigr\}$. 证明 $\frac{dxdydz}{x^2|z|}$ 是 $G$ 上的左 Haar 测度. 它是右 Haar 测度吗?
\end{enumerate}
\end{exercise}

\begin{solution}
当然, 我们将用前一题来解每一问.

(a). 显然的包含 $R^{\times} \subset \R$ 使 $\R^{\times}$ 成为 $\R$ 的开子集. 设 $x \in \R^{\times}$. 则对每个 $y \in \R^{\times}$, $xy = A(x)y + b(x)$, 其中 $A(x) = x \in \GL_1(\R)$, $b(x) = 0$. 所以结论由 $\det(A(x)) = x$ 得出.

(b). 通过映射 $z \longmapsto (\Re(z), \Im(z))$ 把 $\CC^{\times}$ 嵌入 $\R^2$. 这使 $\CC^{\times}$ 成为 $\R^2$ 的开子集. 设 $z = x + iy \in \CC^{\times}$, 其中 $x, y \in \R$. 则 $\CC^{\times}$ 上按 $z$ 的左平移由左乘 $2 \times 2$ 矩阵 $A(z) = \begin{pmatrix} x & -y \\ y & x \end{pmatrix}$ 给出 (这只是公式 $(x + iy)(a + ib) = (xa - yb) + i(ya + xb)$). 所以结论由 $\det(A(z)) = x^2 + y^2$ 得出.

(c). 群 $\GL_n(\R)$ 是 $M_n(\R) \simeq \R^{n^2}$ 的开子集 (因为它由方程 $\det(x) \neq 0$ 给出). 设 $x \in \GL_n(\R)$. 则 $M_n(\R)$ 上按 $x$ 的左平移是线性变换, 我们需要计算它的行列式. 注意 $M_n(\R) = \R^n \oplus \cdots \oplus \R^n$, 其中我们有 $n$ 个和项, 对应 $n \times n$ 矩阵的 $n$ 个列. 按 $x$ 的左乘保持这个分解, 且它在每个和项上的作用就是 $x$ 在 $\R^n$ 上的通常作用, 其行列式为 $\det(x)$. 所以 $M_n(\R)$ 上按 $x$ 的左平移的行列式为 $\det(x)^n$.
为看出 $|\det(T)|^{-n}dT$ 也是右 Haar 测度, 我们使用前一题中把左平移换成右平移的显然类似, 并如上看到 $x \in \GL_n(\R)$ 在 $M_n(\R)$ 上按右平移的作用的行列式为 $\det(x)^n$.

(d). 我们通过把 $\begin{pmatrix} x & y \\ 0 & z \end{pmatrix}$ 送到 $(x, y, z)$ 把 $G$ 嵌入为 $\R^3$ 的开子集. 设 $g = \begin{pmatrix} x & y \\ 0 & z \end{pmatrix} \in G$. 利用
\[
\begin{pmatrix} x & y \\ 0 & z \end{pmatrix}\begin{pmatrix} a & b \\ 0 & c \end{pmatrix} = \begin{pmatrix} xa & xb + yc \\ 0 & zc \end{pmatrix},
\]
我们看到我们处于前一题的情形, 其中
\[
A(g) = \begin{pmatrix} x & 0 & 0 \\ 0 & x & y \\ 0 & 0 & z \end{pmatrix}
\]
且 $b(g) = 0$. 所以 $\det(A(g)) = x^2z$, 我们得到结论.
同 (c), 使用前一题对右平移的类似, 我们看到 $g$ 在 $G$ 上按右平移的作用是线性的, 由矩阵
\[
\begin{pmatrix} x & 0 & 0 \\ y & z & 0 \\ 0 & 0 & z \end{pmatrix}
\]
给出, 其行列式为 $xz^2$. 所以 $\frac{dxdydz}{|x|z^2}$ 是 $G$ 上的右 Haar 测度. 它不是形如 $c\frac{dxdydz}{x^2|z|}$ 的 (其中 $c$ 是常数), 所以 $\frac{dxdydz}{x^2|z|}$ 不能是右 Haar 测度.
\end{solution}

\begin{exercise}[$(\Z/2\Z)^{\N}$ 的 Haar 测度]
\label{ex:I.5.3.3}
考虑群 $G = (\Z/2\Z)^{\N}$.
\begin{enumerate}
\item[(a)] 证明存在 $G$ 上的 Haar 测度 $\mu$ 使得 $\mu(G) = 1$.
\item[(b)] 证明 $G$ 的每个开子集都是形如 $U = V \times (\Z/2\Z)^{\N_{\geq n+1}}$ 的集合的可数并, 其中 $n \in \N$, $V \subset (\Z/2\Z)^{\{0,\ldots,n\}}$, 且 $\mu(U) = \frac{|V|}{2^{n+1}}$.
\item[(c)] 考虑映射 $u \colon G \to [0,1]$, 把 $(x_n)_{n \in \N} \in G$ 送到 $\sum_{n \geq 0}x_n2^{-n-1}$. (在 $u$ 的定义中把 $\Z/2\Z$ 等同于 $\{0,1\}$.) 证明 $u$ 可测, 并把 $\mu$ 送到 $[0,1]$ 上的 Lebesgue 测度 $\lambda$. 即证明: 若 $B \subset [0,1]$ 是 Borel 集, 则 $u^{-1}(B)$ 是 Borel 集且 $\lambda(B) = \mu(u^{-1}(B))$. (提示: 证明形如 $[j2^{-k}, (j+1)2^{-k}]$ 的半开区间生成 $[0,1]$ 上的 Borel $\sigma$-代数, 并计算它们在 $u$ 下的原像.)
\end{enumerate}
\end{exercise}

\begin{solution}
(a). 设 $\mu$ 是 $G$ 上的左 Haar 测度. 由于 $G$ 交换, $\mu$ 也是右 Haar 测度. 又, 由习题 \ref{ex:I.5.1.3}(b), 群 $G$ 紧, 所以 $\mu(G) < +\infty$, 且把 $\mu$ 乘以 $\mu(G)^{-1}$ 后可以假设 $\mu(G) = 1$.

(b). 由乘积拓扑的定义, $G$ 的每个开子集都是形如 $U = V \times (\Z/2\Z)^{\N - I}$ 的集合的并, 其中 $I \subset \N$ 有限, $V \subset (\Z/2\Z)^I$. 由于 $\N$ 的每个有限子集都包含于形如 $\{0,1,\ldots,n\}$ 的集合, 可以假设 $I = \{0,1,\ldots,n\}$ (对某个 $n \in \N$). 我们还需证明并可以是可数的. 假设有 $G$ 的开集 $\Omega$ 具有形式 $\bigcup_{i \in I}U_i$, 其中 $U_i = V_i \times (\Z/2\Z)^{\N_{\geq n_i+1}}$, $V_i \subset (\Z/2\Z)^{\{0,\ldots,n_i\}}$. 对每个 $n \in \N$, 设 $I_n = \{i \in I \mid n_i = n\}$. 则
\[
\bigcup_{i \in I_n}U_i = V_n \times (\Z/2\Z)^{\N_{\geq n+1}},
\]
其中
\[
V_n = \bigcup_{i \in I_n}V_i \subset (\Z/2\Z)^{\{0,\ldots,n\}}.
\]
所以 $\Omega = \bigcup_{n \in \N}V_n \times (\Z/2\Z)^{\N_{\geq n+1}}$.

现在计算 $\mu(U)$, 其中 $U = V \times (\Z/2\Z)^{\N_{\geq n+1}}$, $V \subset (\Z/2\Z)^{\{0,\ldots,n\}}$. 记 $W = (\Z/2\Z)^{\N_{\geq n+1}}$. 若 $v, v' \in (\Z/2\Z)^{\{0,\ldots,n\}}$, 则 $\{v\} \times W = (v' - v) + \{v'\} \times W$, 所以 $\mu(\{v\} \times W) = \mu(\{v'\} \times W)$. 由于 $G = \coprod_{v \in (\Z/2\Z)^{\{0,\ldots,n\}}}\{v\} \times W$, 这蕴含对每个 $v \in (\Z/2\Z)^{\{0,\ldots,n\}}$,
\[
1 = \mu(G) = |(\Z/2\Z)^{\{0,\ldots,n\}}|\mu(\{v\} \times W),
\]
所以 $\mu(\{v\} \times W) = \frac{1}{2^{n+1}}$. 另一方面, $U = \coprod_{v \in V}\{v\} \times W$, 所以我们得到 $\mu(U) = \frac{|V|}{2^{n+1}}$.

(c). 记 $I_{j,k} = [j2^{-k}, (j+1)2^{-k})$. 我们首先证明 $I_{j,k}$ ($k \geq 0$, $0 \leq j \leq 2^k - 1$) 生成 $[0,1]$ 上的 Borel $\sigma$-代数. $[0,1]$ 的每个开子集都是开区间 $(a, b)$ ($0 \leq a < b \leq 1$) 以及可能的一个或两个半开区间 $[0, b)$ ($0 < b \leq 1$) 和 $(a, 1]$ ($0 \leq a < 1$) 的可数并. 所以我们只需验证其中任何一个都可以写成 $I_{j,k}$ 的可数并.
假设 $0 \leq a < b \leq 1$. 若 $b - a > 2^{-k}$, 且 $i = 1 + \lfloor b2^k a \rfloor$, $i' = -1 + \lceil b2^k \rceil$ (其中 $\lfloor\cdot\rfloor$ 和 $\lceil\cdot\rceil$ 是取整函数), 则 $0 < i2^{-k} - a \leq 2^{-k}$ 且 $0 < b - i'2^{-k} \leq 2^{-k}$. 这蕴含
\[
(a, b) = \bigcup_{k > -\log_2(b-a)}\ \bigcup_{j=1+\lfloor b2^k a \rfloor}^{-2+\lceil b2^k \rceil} I_{j,k}
\]
(其中 $\log_2$ 是以 $2$ 为底的对数). 类似地, 若 $0 < b \leq 1$ 且 $0 \leq a < 1$, 则
\[
[0, b) = \bigcup_{k > -\log_2(b)}\ \bigcup_{j=0}^{-2+\lceil b2^k \rceil} I_{j,k}
\]
且
\[
(a, 1] = \bigcup_{k > -\log_2(1-a)}\ \bigcup_{j=1+\lfloor b2^k a \rfloor}^{2^k - 1} I_{j,k}.
\]
这证明了关于 $[0,1]$ 的 Borel $\sigma$-代数的结论. 为完成此题, 我们只需证明对所有 $k \geq 0$ 和所有 $j \in \{0,\ldots,2^k-1\}$, 原像 $u^{-1}(I_{j,k})$ 是 $G$ 中的 Borel 集且 $\mu(u^{-1}(I_{j,k})) = 2^{-k}$. 所以我们计算这些原像.
首先注意, 若 $x \in [0,1]$, 则 $u^{-1}(x)$ 是单点集, 除非 $x$ 具有 $j2^{-k}$ 的形式 ($0 < j < 2^k$); 在后一情形, $x$ 有第二个二进制展开, 其中所有系数在某一秩之后都是 $1$.
现在设 $k \geq 0$, $j \in \{0,1,\ldots,2^k-1\}$. 若 $k = 0$, 则 $j = 0$ 且 $I_{j,k} = [0,1]$, 所以 $u^{-1}(I_{j,k}) = G$ 且由 $\mu$ 的选取有 $\mu(G) = \lambda([0,1]) = 1$. 假设 $k \geq 1$. 由于 $0 \leq j \leq 2^k - 1$, 可以把 $j$ 写成二进制形式 $j = \sum_{i=0}^{k-1}a_{k-1-i}2^i$, 其中 $a_i \in \{0,1\}$. 若 $0 < j$, 还把 $j - 1$ 写成 $j - 1 = \sum_{i=0}^{k-1}b_{k-1-i}2^i$, 其中 $b_i \in \{0,1\}$. 若 $j + 1 < 2^k$, 还把 $j + 1$ 写成 $j + 1 = \sum_{i=0}^{k-1}c_{k-1-i}2^i$, 其中 $c_i \in \{0,1\}$. 则
\[
j2^{-k} = \sum_{i=0}^{k-1}a_i2^{-(i+1)},
\]
\[
j2^{-k} = \sum_{i=0}^{k-1}b_i2^{-(i+1)} + \sum_{i=k}^{+\infty}2^{-i} \quad \text{若 } j > 0
\]
且
\[
(j+1)2^{-k} = \sum_{i=0}^{k-1}c_i2^{-(i+1)} \quad \text{若 } j + 1 < 2^k,
\]
所以 $u^{-1}(I_{j,k}) = X \cup Y$, 其中
\[
Y = \{(a_0,\ldots,a_{k-1})\} \times (\Z/2\Z)^{\N_{\geq k}}
\]
且
\[
X = \begin{cases} \{(b_0,\ldots,b_{k-1},1,1,\ldots),(c_0,\ldots,c_{k-1})\} & \text{若 } 0 < j < 2^k - 1, \\ \{(b_0,\ldots,b_{k-1},1,1,\ldots)\} & \text{若 } j = 2^k - 1, \\ \{(c_0,\ldots,c_{k-1})\} & \text{若 } j = 0. \end{cases}
\]
由于 $X$ 闭且 $Y$ 开, 这是 $G$ 的 Borel 子集. 我们又由 (b) 知道 $\mu(Y) = 2^{-k} = \lambda(I_{j,k})$, 所以只需证明 $\mu(X) = 0$. 即要证明 $G$ 中的单点集体积为 $0$. 由于所有单点集互相是平移, 只需处理 $\{0\}$ 的情形. 这由对每个 $n \geq 0$,
\[
\{0\} \subset \{0\}^{\{0,1,\ldots,n\}} \times (\Z/2\Z)^{\N_{\geq n+1}}
\]
的事实得出, 因为由 (b), 右端体积为 $2^{-(n+1)}$.
\end{solution}

\begin{exercise}[$\Q_p$ 的 Haar 测度]
\label{ex:I.5.3.4}
对 $x \in \Q_p$ 和 $r \in \R$, 记 $B(x, r) = \{y \in \Q_p \mid |x - y|_p \leq r\}$ (以 $x$ 为心、$r$ 为半径的闭球). 设 $\lambda$ 是 $\Q_p$ 上满足 $\lambda(\Z_p) = 1$ 的 Haar 测度.
\begin{enumerate}
\item[(a)] 若 $x \in \Q_p$, $m \in \Z$, 证明 $\lambda(B(x, p^m)) = p^m$.
\item[(b)] 对每个 Borel 集 $X \subset \Q_p$, 证明
\[
\lambda(X) = \inf\Bigl\{\sum_{i \geq 0}p^{m_i} \;\Big|\; \exists x_0, x_1, \ldots \in \Q_p \text{ 使得 } X \subset \bigcup_{i \geq 0}B(x_i, p^{m_i})\Bigr\}.
\]
\end{enumerate}
\end{exercise}

\begin{solution}
(a). 首先注意, 若 $x, y \in \Q_p$, 则 $B(x, p^m) = B(y, p^m) + x - y$, 所以 $\lambda(B(x, p^m)) = \lambda(B(y, p^m))$. 又注意
\[
B(0, p^m) = \{x \in \Q_p \mid |x|_p \leq p^m\} = \{x \in \Q_p \mid |p^mx|_p \leq 1\} = p^{-m}\Z_p
\]
(对每个 $m \in \Z$). 所以对每个 $x \in \Q_p$ 和每个 $m \in \Z$,
\[
B(x, p^m) = x + p^{-m}\Z_p.
\]
又, 由习题 \ref{ex:I.5.1.4}(i),
\[
\Z_p = \coprod_{a=0}^{p-1}(a + p\Z_p).
\]
乘以 $p^{-m}$ 给出
\[
B(0, p^m) = \coprod_{i=0}^{p-1}B(p^{-m}i, p^{m-1}),
\]
所以 $\lambda(B(0, p^m)) = p\lambda(B(0, p^{m-1}))$. 由于由假设 $\lambda(B(0, 1)) = \lambda(\Z_p) = 1$, 结论对 $|m|$ 归纳得出.

(b). 首先, 由习题 \ref{ex:I.5.1.4}(f), 球 $B(x, r)$ 构成 $\Q_p$ 的拓扑的一组 (开!) 集基, 所以 $\Q_p$ 的每个开子集都是球 $B(x, r)$ 的并. 由于 $\Q$ 在 $\Q_p$ 中稠密且可数, $\Q_p$ 的每个开子集都是球 $B(x, r)$ 的可数并 (且可以把 $x$ 取在 $\Q$ 中, 但这无关紧要). 又注意, 由习题 \ref{ex:I.5.1.4}(b), 若 $y \in B(x, r)$, 则 $B(y, r) = B(x, r)$. 所以若 $\Q_p$ 的两个闭球相交, 则其中一个必包含另一个. 这蕴含 $\Q_p$ 的每个开子集都是球 $B(x, r)$ 的不相交可数并. 结论现在由 (a) 和 $\lambda$ 的外正则性立即得出.
\end{solution}

\begin{exercise}[商群上的 Haar 测度]
\label{ex:I.5.3.5}
设 $G$ 是局部紧群, $H$ 是 $G$ 的闭子群. 记 $\pi$ 为从 $G$ 到 $G/H$ 的商映射. 记 $\Delta_G$ (相应地 $\Delta_H$) 为 $G$ (相应地 $H$) 的模函数, 并假设 $\Delta_G|_H = \Delta_H$. 固定 $G$ 和 $H$ 上的左 Haar 测度 $\mu_G$ 和 $\mu_H$.
\begin{enumerate}
\item[(a)] 证明对 $G/H$ 的每个紧子集 $K'$, 存在 $G$ 的紧子集 $K$ 使得 $\pi(K) = K'$.
\item[(b)] 设 $f \in L^1(G)$. 证明函数 $G \to \CC$, $x \longmapsto \int_H f(xh)\,\dd\mu_H(h)$ 被 $H$ 的元素的右平移不变. 因此它定义了一个函数 $G/H \to \CC$, 记为 $f^H$.
\item[(c)] 若 $f \in \mathcal{C}_c(G)$, 证明 $f^H \in \mathcal{C}_c(G/H)$.
\item[(d)] 证明映射 $\mathcal{C}_c(G) \to \mathcal{C}_c(G/H)$, $f \longmapsto f^H$ 满射. (提示: 可以利用对 $G$ 的每个紧子集 $K$, 存在 $\phi \in \mathcal{C}_c^+(G)$ 使得对每个 $x \in K$ 有 $\phi(x) > 0$ 这一事实.)
\item[(e)] 若 $f \in \mathcal{C}_c(G)$ 满足 $f^H = 0$, 证明 $\int_G f(x)\,\dd\mu_G(x) = 0$. (提示: 用 $\mathcal{C}_c(G/H)$ 中在 $\pi(\operatorname{supp}(f))$ 上等于 $1$ 的函数, 以及命题 \ref{prop:I.2.12}.)
\item[(f)] 证明存在 $G/H$ 上唯一的正则 Borel 测度 $\mu_{G/H}$, 它被 $G$ 的元素的左平移不变, 且对每个 $f \in \mathcal{C}_c(G)$ 有 $\int_G f(x)\,\dd\mu_G(x) = \int_{G/H}f^H(y)\,\dd\mu_{G/H}(y)$.
\item[(g)] 若 $P$ 是 $G$ 的闭子群且 $\pi$ 诱导同胚 $P \xrightarrow{\sim} G/H$, 证明 $\mu_{G/H}$ 在此同胚下的原像是 $P$ 上的左 Haar 测度.
\item[(h)] 若 $P$ 是 $G$ 的闭子群, 映射 $P \times H \to G$, $(p, h) \longmapsto ph$ 是同胚, 且 $d\mu_P$ 是 $P$ 上的左 Haar 测度, 证明线性泛函 $\mathcal{C}_c(G) \to \CC$, $f \longmapsto \int_H\int_P f(ph)\,\dd\mu_P(p)\,\dd\mu_H(h)$ 定义 $G$ 上的左 Haar 测度.
\end{enumerate}
\end{exercise}

\begin{solution}
(a). 设 $V$ 是 $1$ 在 $G$ 中的紧邻域. 则 $\pi(V)$ 是 $\pi(1)$ 在 $G/H$ 中的紧邻域. 我们有 $K' \subset \bigcup_{x \in \pi^{-1}(K')}\pi(xV)$. 由于 $K'$ 紧, 可以找到 $x_1, \ldots, x_n$ 使得 $K' \subset \bigcup_{i=1}^n\pi(x_iV)$. 设 $K = \pi^{-1}(K') \cap \bigl(\bigcup_{i=1}^n x_iV\bigr)$. 则 $K$ 是紧集 $\bigcup_{i=1}^n x_iV$ 的闭子集, 所以它紧, 且有 $\pi(K) = K'$.

(b). 设 $x \in H$. 则对每个 $g \in G$,
\[
\int_H f(gxh)\,\dd\mu_H(h) = \int_H f(gh)\,\dd\mu_H(h)
\]
由 $\mu_H$ 的左不变性.

(c). 我们需要证明 $f^H$ 连续且有紧支撑.
固定 $1$ 的对称紧邻域 $V_0$, 注意 $A := \operatorname{supp} f \cup V_0(\operatorname{supp} f)$ 紧. 设 $\varepsilon > 0$. 由于 $f$ 左一致连续, 存在 $1$ 的邻域 $V \subset V_0$ 使得对每个 $x \in G$ 和每个 $y \in V$, 有 $|f(yx) - f(x)| \leq \varepsilon$. 则对每个 $x \in G$ 和每个 $y \in V$,
\[
|f^H(\pi(yx)) - f^H(x)| = \Bigl|\int_H (f(yxh) - f(xh))\,\dd\mu_H(h)\Bigr| \leq \varepsilon\mu_H(x^{-1}A \cap H),
\]
因为除非 $y \in (x^{-1}\operatorname{supp} f) \cup (x^{-1}y^{-1}\operatorname{supp} f) \subset x^{-1}A$, 有 $f(yxh) = f(xh) = 0$. 由于 $x^{-1}A \cap H$ 紧, 它有有限测度, 且上面的计算蕴含 $f^H$ 在点 $\pi(x)$ 处连续.
现在我们证明 $f^H$ 有紧支撑. 由 $f^H$ 的定义, 若 $x \notin KH$ 则 $f^H(\pi(x)) = 0$. 所以 $f^H$ 的支撑包含于 $\pi(KH) = \pi(K)$, 从而是紧的.

(d). 设 $g \in \mathcal{C}_c(G/H)$, $K'$ 是它的支撑. 由 (a), 存在 $G$ 的紧子集 $K$ 使得 $\pi(K) = K'$. 设 $\phi \in \mathcal{C}_c^+(G)$ 使得对每个 $x \in K$ 有 $\phi(x) > 0$. 我们证明对每个 $y \in K'$ 有 $\phi^H(y) > 0$. 设 $y \in K'$, 把 $y$ 写成 $y = \pi(x)$ ($x \in K$). 由于 $\phi(x) > 0$ 且 $\phi$ 连续, 可以找到 $1$ 在 $G$ 中的开邻域 $V$ 和 $c \in \R_{>0}$ 使得对每个 $x' \in xV$ 有 $\phi(x') \geq c$. 特别地,
\[
\phi^H(y) = \int_H \phi(xh)\,\dd\mu_H(h) \geq \int_{H \cap V}\phi(xh)\,\dd\mu_H(h) \geq c \cdot \mu_H(U \cap H) > 0
\]
(因为 $U \cap H$ 是 $H$ 的非空开子集, 有 $\mu_H(U \cap H) > 0$).
我们如下定义函数 $F \colon G \to \CC$:
\[
F(x) = \begin{cases} \frac{g(\pi(x))}{\phi^H(\pi(x))} & \text{若 } \phi^H(\pi(x)) > 0, \\ 0 & \text{否则}. \end{cases}
\]
注意 $F$ 在开子集 $U_1 = \{x \in G \mid \phi^H(\pi(x)) > 0\}$ 和 $U_2 = G - \operatorname{supp}(g \circ \pi)$ 上连续 (在第二个子集上它恒为零). 由于 $U_1 \supset \pi^{-1}(K')$ 且 $\pi^{-1}(K') = \operatorname{supp}(g \circ \pi)$, 有 $U_1 \cup U_2 = G$, 函数 $F$ 在 $G$ 上连续. 最后, 取 $f = F\phi$. 则 $f \in \mathcal{C}_c(G)$, 我们只需证明 $f^H = g$.
设 $x \in G$. 若 $\phi^H(\pi(x)) = 0$, 则对每个 $h \in H$ 有 $f(xh) = 0$, 所以 $f^H(\pi(x)) = 0$. 我们已见 $\phi^H$ 在 $K' = \operatorname{supp}(g)$ 上取正值, 所以我们也有 $x \notin \operatorname{supp}(g)$, 即 $g(x) = 0 = f^H(x)$. 现在假设 $\phi^H(\pi(x)) > 0$. 注意函数 $H \to \CC$, $h \longmapsto F(xh)$ 是常数. 所以
\[
f^H(\pi(x)) = F(x)\int_H \phi(xh)\,\dd\mu_H(h) = \frac{g(\pi(x))}{\phi^H(\pi(x))}\phi^H(\pi(x)) = g(\pi(x)).
\]
最后, 注意若 $g \in \mathcal{C}_c^+(G/H)$ 则 $f \in \mathcal{C}_c^+(G)$, 且我们顺带证明了若 $f \in \mathcal{C}_c^+(G)$ 则 $f^H \in \mathcal{C}_c^+(G/H)$ (我们是对 $\phi$ 证明的).

(e). 设 $\psi \in \mathcal{C}_c(G/H)$ 使得对每个 $y \in \pi(\operatorname{supp} f)$ 有 $\psi(y) = 1$. 由 (d), 存在 $\phi \in \mathcal{C}_c(G)$ 使得 $\phi^H = \psi$. 我们有
\begin{align*}
\int_G f(x)\,\dd\mu_G(x) &= \int_G f(x)\phi^H(\pi(x))\,\dd\mu_G(x) \\
&= \int_{G \times H} f(x)\phi(xh)\,\dd\mu_G(x)\,\dd\mu_H(h) \\
&= \int_H \Bigl(\int_G f(x)\phi(xh)\,\dd\mu_G(x)\Bigr)\,\dd\mu_H(h) \\
&= \int_H \Bigl(\Delta_G(h)^{-1}\int_G f(xh^{-1})\phi(x)\,\dd\mu_G(x)\Bigr)\,\dd\mu_H(h) \\
&= \int_G \phi(x)\Bigl(\int_H \Delta_H(h)^{-1}f(xh^{-1})\,\dd\mu_H(h)\Bigr)\,\dd\mu_G(x) \\
&= \int_G \phi(x)\Bigl(\int_H f(xh)\,\dd\mu_H(h)\Bigr)\,\dd\mu_G(x) \quad \text{(由命题 \ref{prop:I.2.12})} \\
&= 0 \quad \text{(因为 } f^H = 0\text{)}.
\end{align*}

(f). 由 (e), 正线性函数 $\mathcal{C}_c(G) \to \CC$, $f \longmapsto \int_G f\,\dd\mu_G$ 穿过线性映射 $\mathcal{C}_c(G) \to \mathcal{C}_c(G/H)$, $f \longmapsto f^H$. 由 (d) (及其解答末尾的注记), 它定义了 $\mathcal{C}_c(G/H)$ 上的正线性泛函. 由 Riesz 表示定理, 这来自 $G/H$ 上的正则 Borel 测度 $\mu_{G/H}$. 展开定义, 对每个 $f \in \mathcal{C}_c(G)$ 得到
\[
\int_G f\,\dd\mu_G = \int_{G/H}f^H\,\dd\mu_{G/H}.
\]
由 $\mu_G$ 的左不变性和 (d), 若 $f \in \mathcal{C}_c(G/H)$, $x \in G$, 则
\[
\int_{G/H}f(xy)\,\dd\mu_{G/H}(y) = \int_{G/H}f(y)\,\dd\mu_{G/H}.
\]
使用 Riesz 表示定理的唯一性部分 (如课堂上那样), 看到对 $G/H$ 的每个 Borel 子集 $E$ 有 $\mu_{G/H}(xE) = \mu_{G/H}(E)$.

(g). 设 $\nu$ 是同胚 $\alpha \colon P \xrightarrow{\sim} G/H$ 下 $\mu_{G/H}$ 的原像. 它是正则 Borel 测度, 因为 $\alpha$ 是同胚. 又注意对每个 $x \in P$ 有 $\alpha(xy) = x\alpha(y)$ (这在 $\alpha$ 的定义上显然). 由于 $\mu_{G/H}$ 被 $P$ 的元素的左平移不变, $\nu$ 也一样.

(h). 假设蕴含 $\pi$ 诱导同胚 $P \xrightarrow{\sim} G/H$, 因此如 (g) 得到 $P$ 上的左 Haar 测度 $\nu$. 由左 Haar 测度的唯一性, 有 $\mu_P = c\nu$ (对某个 $c \in \R_{>0}$). 所以对每个 $f \in \mathcal{C}_c(G)$,
\begin{align*}
\int_H\int_P f(ph)\,\dd\mu_P(p)\,\dd\mu_H(h) &= c\int_P\Bigl(\int_H f(ph)\,\dd\mu_H(h)\Bigr)\,\dd\nu(p) \\
&= c\int_{G/H}f^H(y)\,\dd\mu_{G/H}(y) \\
&= c\int_G f(x)\,\dd\mu_G(x).
\end{align*}
所以泛函 $f \longmapsto \int_H\int_P f(ph)\,\dd\mu_P(p)\,\dd\mu_H(h)$ 正, 且对应于左 Haar 测度 $c\mu_G$.
\end{solution}

\begin{exercise}[半直积的模函数]
\label{ex:I.5.3.6}
设 $G$ 是局部紧群. 设 $A$ 和 $N$ 是 $G$ 的两个闭子群, 使得 $A \times N \to G$, $(a, n) \longmapsto an$ 是同胚, 且 $A$ 正规化 $N$ (即对每个 $a \in A$ 和 $n \in N$, 有 $ana^{-1} \in N$).
\begin{enumerate}
\item[(a)] 若 $\mu_A$ 和 $\mu_N$ 是 $A$ 和 $N$ 上的左 Haar 测度, 证明线性泛函 $\mathcal{C}_c(G) \to \CC$, $f \longmapsto \int_A\int_N f(an)\,\dd\mu_A\,\dd\mu_N$ 定义 $G$ 上的左 Haar 测度.
\item[(b)] 设 $a \in A$. 证明存在 $\alpha(a) \in \R_{>0}$ 使得对每个 $f \in \mathcal{C}_c(N)$,
\[
\int_N f(ana^{-1})\,\dd\mu_N(n) = \alpha(a)\int_N f(n)\,\dd\mu_N(n).
\]
\item[(c)] 若 $\Delta_G$, $\Delta_A$ 和 $\Delta_N$ 分别是 $G$, $A$ 和 $N$ 的模函数, 证明当 $a \in A$, $n \in N$ 时 $\Delta_G(an) = \alpha(a)\Delta_A(a)\Delta_N(n)$.
\end{enumerate}
\end{exercise}

\begin{solution}
(a). 情形与习题 \ref{ex:I.5.3.5} 非常相似 (例如让 $N$ 扮演 $H$ 的角色), 区别是我们对模函数不做任何假设. 尽管如此, 习题 \ref{ex:I.5.3.5} 的 (a)-(d) 的结果仍然成立, 因为它们的证明不用模函数的假设. 特别地, 我们得到从 $\mathcal{C}_c(G)$ 到 $\mathcal{C}_c(G/N) \simeq \mathcal{C}_c(A)$ 的满线性变换 $f \longmapsto f^N$, 且它把 $\mathcal{C}_c^+(G)$ 送到 $\mathcal{C}_c^+(A)$. 题设的线性函数把 $f \in \mathcal{C}_c(G)$ 送到 $\int_A f^N(a)\,\dd\mu_A(a)$, 所以它正, 且 Riesz 表示定理说存在 $G$ 上唯一的正则 Borel 测度 $\mu_G$ 使得对每个 $f \in \mathcal{C}_c(G)$,
\[
\int_G f\,\dd\mu_G = \int_A\int_N f(an)\,\dd\mu_A(a)\,\dd\mu_N(n).
\]
由于 $\mu_A$ 是 $A$ 上的左 Haar 测度, 上面的公式蕴含对每个 $f \in \mathcal{C}_c(G)$ 和每个 $a \in A$ 有 $\int_G L_af\,\dd\mu_G = \int_G f\,\dd\mu_G$. 我们证明 $\mu_G$ 被 $N$ 左不变. 设 $x \in N$, $f \in \mathcal{C}_c(G)$. 则
\begin{align*}
\int_G L_xf\,\dd\mu_G &= \int_A\int_N f(xan)\,\dd\mu_A(a)\,\dd\mu_N(n) \\
&= \int_A\Bigl(\int_N f(a(a^{-1}xa)n)\,\dd\mu_N(n)\Bigr)\,\dd\mu_A(a) \\
&= \int_A\Bigl(\int_N f(an)\,\dd\mu_N(n)\Bigr)\,\dd\mu_A(a) \quad \text{因为 } a^{-1}xa \in N \text{ 且 } \mu_N \text{ 左不变} \\
&= \int_G f\,\dd\mu_G.
\end{align*}
由于 $G = AN$, 这蕴含对每个 $x \in G$ 和每个 $f \in \mathcal{C}_c(G)$ 有 $\int_G L_xg\,\dd\mu_G = \int_G f\,\dd\mu_G$. 由命题 \ref{prop:I.2.6}, $\mu_G$ 是 $G$ 上的左 Haar 测度.

(b). 注意映射 $N \to N$, $n \longmapsto a^{-1}na$ 是同胚. 所以公式 $E \longmapsto \mu_N(a^{-1}Ea)$ 定义 $N$ 上的正则 Borel 测度, 记为 $\nu$. 若 $E$ 是 Borel 子集且 $n \in N$, 则
\[
\nu(nE) = \mu(a^{-1}nEa) = \mu((a^{-1}na)a^{-1}Ea) = \mu(a^{-1}Ea) = \nu(E).
\]
所以 $\nu$ 是 $N$ 上的左 Haar 测度, 因此存在 $\alpha(a) \in \R_{>0}$ 使得 $\nu = \alpha(a)\mu_N$. 现在, 若 $E$ 是 $N$ 的 Borel 子集且 $f = \mathbb{1}_E$, 则函数 $n \longmapsto f(ana^{-1})$ 是 $a^{-1}Ea$ 的特征函数, 所以
\[
\int_N f(ana^{-1})\,\dd\mu_N(n) = \mu(a^{-1}Ea) = \alpha(a)\mu(E) = \alpha(a)\int_N f\,\dd\mu_N.
\]
这按通常方式延拓到所有函数 $f \in L^1(N)$, 特别地延拓到 $f \in \mathcal{C}_c(N)$.

(c). 设 $a \in A$, $n \in N$, 固定 $f \in \mathcal{C}_c^+(G)$. 则
\begin{align*}
\Delta_G(an)\int_G f\,\dd\mu_G &= \int_G R_{an}(f)\,\dd\mu_G \\
&= \int_A\int_N f(bman)\,\dd\mu_A(b)\,\dd\mu_N(m) \\
&= \int_A\Bigl(\int_N f(ba(a^{-1}ma)n)\,\dd\mu_N(m)\Bigr)\,\dd\mu_A(b) \\
&= \alpha(a)^{-1}\int_A\Bigl(\int_N f(bamn)\,\dd\mu_N(m)\Bigr)\,\dd\mu_A(b) \quad \text{由 (b)} \\
&= \alpha(a)^{-1}\Delta_N(n)^{-1}\int_A\Bigl(\int_N f(bam)\,\dd\mu_N(m)\Bigr)\,\dd\mu_A(b) \quad \text{由 } \Delta_N \text{ 的定义} \\
&= \alpha(a)^{-1}\Delta_N(n)^{-1}\int_N\Bigl(\int_A f(bam)\,\dd\mu_A(b)\Bigr)\,\dd\mu_N(m) \\
&= \alpha(a)^{-1}\Delta_N(n)^{-1}\Delta_A(a)^{-1}\int_N\Bigl(\int_A f(bm)\,\dd\mu_A(b)\Bigr)\,\dd\mu_N(m) \quad \text{由 } \Delta_A \text{ 的定义} \\
&= \alpha(a)^{-1}\Delta_N(n)^{-1}\Delta_A(a)^{-1}\int_G f\,\dd\mu_G.
\end{align*}
由于 $\int_G f\,\dd\mu_G > 0$, 这蕴含 $\Delta_G(an) = \alpha(a)\Delta_A(a)\Delta_N(n)$.
\end{solution}

\begin{exercise}[$\SL_n(\R)$ 的 Iwasawa 分解与模函数]
\label{ex:I.5.3.7}
设 $G = \SL_n(\R)$, $H = \SO(n)$, $P \subset G$ 是对角线元素为正的上三角矩阵 (行列式为 $1$) 的子群.
\begin{enumerate}
\item[(a)] 证明映射 $P \times H \to G$, $(p, h) \longmapsto ph$ 是同胚. (提示: Gram-Schmidt.)
\item[(b)] 给出与习题 \ref{ex:I.5.3.2}(d) 中公式类似的 $P$ 上左 Haar 测度的公式.
\item[(c)] 计算 $P$ 的模函数.
\item[(d)] 证明 $G$ 幺模. (有几种方法.)
\item[(e)] 若 $n = 2$, 证明 $\SO(n) \simeq S^1$ (圆群), 并给出 $G$ 上的左 Haar 测度.
\end{enumerate}
\end{exercise}

\begin{solution}
(a). 在此题中, 记 $\langle \cdot, \cdot \rangle$ 为 $\R^n$ 上的通常 Euclid 内积, $\|\cdot\|$ 为对应的范数.
记题设的映射 $P \times H \to G$ 为 $\alpha$. 此映射连续, 因为 $\SL_n(\R)$ 是拓扑群. 我们首先证明它单射. 假设有 $p, p' \in P$ 和 $h, h' \in H$ 使得 $ph = p'h'$. 则 $p^{-1}p' = h(h')^{-1} \in P \cap H$ 是对角线元素为正的上三角特殊正交矩阵. 这样的矩阵必为单位阵. 事实上, 设 $(v_1, \ldots, v_n)$ 是它的列, $(e_1, \ldots, e_n)$ 是 $\R^n$ 的典范基. 我们要证明 $(v_1, \ldots, v_n) = (e_1, \ldots, e_n)$. 由于 $v_1$ 是范数 $1$ 的向量且是 $e_1$ 的正倍数, 必有 $v_1 = e_1$. 由于向量 $v_2, \ldots, v_n$ 与 $v_1$ 正交, 它们的第一个分量全为 $0$. 所以 $v_2$ 是 $e_2$ 的正倍数; 由于 $v_2$ 范数 $1$, 必有 $v_2 = e_2$. 现在向量 $v_3, \ldots, v_n$ 与 $v_2$ 正交, 所以它们的第二个分量为零, 于是 $v_3$ 是 $e_3$ 的正倍数, 等等.
现在回忆 Gram-Schmidt 正交化过程. 若 $(v_1, \ldots, v_n)$ 是 $\R^n$ 的基, 它按下述方式产生正交基 $(w_1, \ldots, w_n)$ 和标准正交基 $(u_1, \ldots, u_n)$:
\begin{itemize}
\item $w_1 = v_1$ 且 $u_1 = \frac{1}{\|w_1\|}w_1$;
\item 对 $1 \leq k \leq n - 1$, $u_{k+1} = \frac{1}{\|w_{k+1}\|}w_{k+1}$, 其中 $w_{k+1} = v_{k+1} - \sum_{i=1}^k\frac{\langle w_i, v_{k+1} \rangle}{\langle w_i, w_i \rangle}w_i$.
\end{itemize}
特别地, 若 $A$ (相应地 $B$, 相应地 $C$) 是列向量为 $(v_1, \ldots, v_n)$ (相应地 $(w_1, \ldots, w_n)$, 相应地 $(u_1, \ldots, u_n)$) 的矩阵, 则 $B = AN$ 且 $C = AND$, 其中 $N$ 是对角线为 $1$ 的上三角矩阵, $D$ 是对角线为 $(\|w_1\|^{-1}, \ldots, \|w_n\|^{-1})$ 的对角矩阵. 又注意 $N$ 和 $D$ 的元素是 $v_1, \ldots, v_n$ 的连续函数, 所以 $B$ 和 $C$ 的元素也是, 且 $C$ 是正交矩阵. 若 $x \in \SL_n(\R)$, 把此过程应用于 $x$ 的列, 得到一个正交矩阵 $h$ 和一个矩阵 $p \in P$, 两者都连续依赖于 $x$, 使得 $h = xp$, 即 $x = hp^{-1}$. 又, $\det(h) = \det(xp) = \det(p) > 0$, 所以 $h$ 实际上在 $\SO(n)$ 中. 由于 $g \longmapsto g^{-1}$ 是 $\GL_n(\R)$ (从而是其子群 $P$) 上的连续函数, 我们构造了连续映射 $\beta \colon G \to P \times H$ 使得 $\alpha \circ \beta = \mathrm{id}_G$. 特别地, 映射 $\alpha$ 满射, 所以它双射. 则它的逆必为 $\beta$, 且我们知道 $\beta$ 连续. 所以 $\alpha$ 是同胚.

(b). 注意 $P$ 是 $M_n(\R)$ 中上三角矩阵构成的 $\R$-向量空间 $V$ 的开子集. 而且, 对每个 $p \in P$, $P$ 上按 $p$ 的左平移是线性自同态 $T_p \colon V \to V$, $x \longmapsto px$ 的限制. 所以我们可以应用习题 \ref{ex:I.5.3.1} 把 $P$ 上的 Haar 测度定义为 $|\det(T_p)|^{-1}dV(p)$, 其中 $dV$ 是 $V$ 上的 Lebesgue 测度.
我们还需计算 $p \in P$ 时的 $\det(T_p)$. 设 $p \in P$, $a_1, \ldots, a_n$ 是它的对角线元素. 设 $(e_1, \ldots, e_n)$ 如前为 $\R^n$ 的典范基, $V_i = \operatorname{Span}(e_1, \ldots, e_i) \subset \R^n$ ($1 \leq i \leq n$). 注意 $p \in \GL_n(\R)$ 的作用保持子空间 $V_1, \ldots, V_n$, 且 $p$ 诱导的 $V_i$ 的自同态的行列式为 $a_1 \cdots a_i$. 用矩阵的列分解 $V$ (如习题 \ref{ex:I.5.3.2}(c) 的解答), 得到同构 $V \simeq V_1 \oplus V_2 \oplus \cdots \oplus V_n$ 使得自同态 $T_p$ 对应于 $p$ 在每个 $V_i$ 上的作用. 所以我们得到
\[
\det(T_p) = \prod_{i=1}^n\prod_{r=1}^i a_r = a_1^na_2^{n-2}\cdots a_{n-1}^2a_n = \prod_{i=1}^n a_i^{n+1-i}.
\]

(c). 我们将使用习题 \ref{ex:I.5.3.6}, 取 $G = P$, $N$ 为幂幺上三角矩阵群 (即对角线为 $1$ 的上三角矩阵群), $A$ 为对角线元素为正的对角矩阵群. 设 $\alpha \colon A \times N \to P$ 是由 $\alpha(a, n) = an$ 定义的映射. 让我们证明 $\alpha$ 是同胚. 映射 $\alpha$ 显然连续, 且它单射, 因为 $N \cap A = \{1\}$. 设 $x \in P$, 设 $a \in A$ 是与 $x$ 有相同对角线元素的矩阵. 则 $n := a^{-1}x \in N$, 且 $\alpha(a, n) = x$. 所以 $\alpha$ 双射. 而且, 矩阵 $a$ 连续依赖于 $x$, 所以 $n$ 也如此, 故 $\alpha$ 的逆连续, 最后 $\alpha$ 是同胚.
我们想应用习题 \ref{ex:I.5.3.6}(c). 为此需要计算 $A$ 和 $N$ 的模函数以及函数 $\alpha \colon A \to \R_{>0}$.
首先, 由于 $A$ 交换, 有 $\Delta_A = 1$.
对 $N$, 有几种做法. 例如, 你可以注意到 $N$ (仅作为拓扑空间) 显然与 $M_n(\R)$ 中对角线为零的上三角矩阵构成的 $\R$-向量空间 $W$ 同胚. (只需忘掉矩阵的对角线项.) 而且, 对 $n \in N$, $N$ 上按 $n$ 的左平移对应于 $W$ 的线性自同态 $U_n$, 由 $U_n(X) = nX$ ($X \in W$) 给出. 注意 $W$ 是前一问中空间 $V$ 的子空间, 且 $U_n$ 是 $T_n$ 的限制. 所以可以用与前一问相同的方法计算 $\det(U_n)$, 得到 $\det(U_n) = 1$.
所以 $W$ 上的 Lebesgue 测度是 $N$ 上的左 Haar 测度. 我们可以用右平移代替左平移重做一切, 得到 $W$ 上的 Lebesgue 测度也是 $N$ 上的右 Haar 测度. 这意味着 $N$ 幺模, 所以 $\Delta_N = 1$.
最后, 我们需要计算函数 $\alpha$. 回忆它由
\[
\int_N f(ana^{-1})\,\dd n = \alpha(a)\int_N f(n)\,\dd n
\]
定义 (对每个 $f \in \mathcal{C}_c(N)$), 其中 $dn$ 是 $W$ 上的 Lebesgue 测度 (我们刚看到它是 $N$ 上的 Haar 测度). 注意 $c_a \colon X \longmapsto aXa^{-1}$ 是 $W$ 的线性自同态, 所以一旦知道 $\det(c_a)$ 就能用换元公式计算 $\int_N f(ana^{-1})\,\dd n$. 我们得到 $\det(c_a)\int_N f(ana^{-1})\,\dd n = \int_N f(n)\,\dd n$, 所以 $\alpha(a) = \det(c_a)^{-1}$. 但容易看出, 若 $a$ 的对角线元素为 $(a_1, \ldots, a_n)$, 则
\[
\det(c_a) = a_1^{n-1}a_2^{n-3}\cdots a_n^{1-n} = \prod_{i=1}^n a_i^{n-2i+1}.
\]
所以最终对 $p \in P$,
\[
\Delta_P(p) = a_1^{1-n}a_2^{3-n}\cdots a_n^{n-1} = \prod_{i=1}^n a_i^{2i-n-1},
\]
其中 $a_1, \ldots, a_n$ 是 $p$ 的对角线元素.

(d). 若你知道 (或知道如何证明) $\SL_n(\R)$ 等于它的换位子群, 这就很容易. 这里给出另一种方法: 设 $\GL_n(\R)^+$ 是行列式为正的 $n \times n$ 矩阵群. 这是 $\GL_n(\R)$ 的开子群 (它是连续群同态 $\det \colon \GL_n(\R) \to \R^{\times}$ 下 $\R_{>0}$ 的原像), 所以若 $\mu$ 是 $\GL_n(\R)$ 上的 Haar 测度 (回忆由习题 \ref{ex:I.5.3.2}(c), $\GL_n(\R)$ 幺模), 则它到 $\GL_n(\R)^+$ 的限制是非零正则 Borel 测度, 且显然是 $\GL_n(\R)^+$ 上的左、右 Haar 测度. 现在注意有拓扑群的同构 $\R_{>0} \times \SL_n(\R) \to \GL_n(\R)^+$, $(\lambda, x) \longmapsto \lambda x$ (其逆由 $x \longmapsto (\det(x)^{1/n}, \det(x)^{-1/n}x)$ 给出), 所以可以应用习题 \ref{ex:I.5.3.6}, 取 $G = \GL_n(\R)^+$, $A = \R_{>0}I_n$, $N = \SL_n(\R)$. 由于 $A$ 和 $N$ 交换, 有 $\alpha = 1$. 我们知道 $A$ 幺模 (因为它交换), 且刚才已见 $\GL_n(\R)^+$ 幺模, 所以由习题 \ref{ex:I.5.3.6}(c), $\SL_n(\R)$ 也幺模.

(e). 众所周知 $\R^2$ 中的旋转群 (即 $\SO(2)$) 与圆群 $S^1$ 同构. 同构把 $e^{2i\pi\theta} \in S^1$ 送到矩阵 $\begin{pmatrix} \cos\theta & \sin\theta \\ -\sin\theta & \cos\theta \end{pmatrix}$. 又, 我们在课堂上见过可以通过把 $f \in \mathcal{C}_c(S^1)$ 送到 $\int_0^1 f(e^{2i\pi\theta})\,\dd\theta$ 的线性泛函来定义 $S^1$ 上的 Haar 测度, 其中 $d\theta$ 是 $\R$ 上的 Lebesgue 测度.
这当然的重点是, 习题 \ref{ex:I.5.3.5} 现在允许你定义 $\SL_2(\R)$ 上的 Haar 测度. 为处理 $\SL_n(\R)$ 的情形, 我们需要 $\SO(n)$ 上的 Haar 测度. 这种测度的一个例子在习题 \ref{ex:I.5.3.9} 中给出.
\end{solution}

\begin{exercise}[$\GL_n(\Q_p)$ 上的 Haar 测度]
\label{ex:I.5.3.8}
(回忆习题 \ref{ex:I.5.1.4}, \ref{ex:I.5.3.1}, \ref{ex:I.5.3.2} 和 \ref{ex:I.5.3.4}.) 记 $dx$ 为加法群 $\Q_p$ 上的 Haar 测度. 又记 $dx$ (相应地 $dA$) 为 $\Q_p^n$ (相应地 $M_n(\Q_p) \simeq \Q_p^{n^2}$) 上的乘积测度; 注意它是相应加法群的 Haar 测度.
\begin{enumerate}
\item[(a)] 证明对每个 $f \in L^1(\Q_p)$ 和每个 $a \in \Q_p^{\times}$, $b \in \Q_p$,
\[
\int_{\Q_p}f(x)\,\dd x = |a|_p\int_{\Q_p}f(ax + b)\,\dd x.
\]
\item[(b)] 设 $n \geq 1$. 证明若 $f \in L^1(\Q_p^n)$, $A \in \GL_n(\Q_p)$, $b \in \Q_p^n$, 则
\[
\int_{\Q_p^n}f(x)\,\dd x = |\det(A)|_p\int_{\Q_p^n}f(Ax + b)\,\dd x.
\]
\item[(c)] 证明 $|\det(A)|_p^{-n}dA$ 是 $\GL_n(\Q_p)$ 上的左、右 Haar 测度.
\item[(d)] 设 $B$ 是 $\GL_n(\Q_p)$ 中上三角矩阵群. 找出 $B$ 上的左 Haar 测度并计算 $B$ 的模函数.
\end{enumerate}
\end{exercise}

\begin{solution}
(a). 首先, 利用 $dx$ 的平移不变性, 看到对每个 $f \in L^1(\Q_p)$ 和 $a, b \in \Q_p$,
\[
\int_{\Q_p}f(ax + b)\,\dd x = \int_{\Q_p}f(ax)\,\dd x.
\]
设 $a \in \Q_p^{\times}$. 我们使用习题 \ref{ex:I.5.3.4} 的记号. 若 $x \in \Q_p$, $m \in \Z$, 则
\[
aB(x, p^m) = \{ay \mid |x - y|_p \leq p^m\} = \{y \in \Q_p \mid |ax - y|_p \leq |a|_pp^m\} = B(ax, |a|_pp^m),
\]
所以由 \ref{ex:I.5.3.4}(a), $\operatorname{vol}(aB(x, p^m)) = |a|_p\operatorname{vol}(B(x, p^m))$. 使用同一题 (b), 对 $\Q_p$ 的每个 Borel 子集 $E$ 得到 $\operatorname{vol}(aE) = |a|_p\operatorname{vol}(E)$. 假设 $f = \mathbb{1}_E$, 其中 $E$ 是 $\Q_p$ 的 Borel 子集. 则
\[
\int_{\Q_p}f(ax)\,\dd x = \operatorname{vol}(a^{-1}E) = |a|_p^{-1}\int_{\Q_p}f(x)\,\dd x,
\]
所以对这个函数 $f$ 得到所要的结果. 结论现在对每个 $f \in L^1(\Q_p)$ 由积分的线性和连续性得出.

(b). 像 (a) 那样使用 $dx$ 的平移不变性, 看到只需在 $b = 0$ 的情形证明结论. 设 $A \in \GL_n(\Q_p)$. 首先注意 $A = A_1A_2$, 且若我们对 $A_1$ 和 $A_2$ 知道结论, 就对 $A$ 知道结论; 事实上, 对每个 $f \in L^1(\Q_p^n)$,
\begin{multline*}
\int_{\Q_p^n}f(x)\,\dd x = |\det(A_1)|_p\int_{\Q_p^n}f(A_1x)\,\dd x \\
= |\det(A_1)|_p|\det(A_2)|_p\int_{\Q_p^n}f(A_1(A_2x))\,\dd x = |\det(A)|_p\int_{\Q_p^n}f(Ax)\,\dd x.
\end{multline*}
Gauss 算法 (解线性方程组) 说可以通过初等行变换 (对应左乘一个下三角矩阵) 和行置换 (对应左乘一个置换矩阵) 把 $A$ 变成上三角. 所以由上面的观察, 只需对上三角矩阵、下三角矩阵和置换矩阵证明结论.
先假设 $A$ 是置换矩阵. 则存在置换 $\sigma \in S_n$ 使得对每个 $x = (x_1 \ldots, x_n) \in \Q_p^n$, $Ax = (x_{\sigma(1)}, \ldots, x_{\sigma(n)})$. 由于 $dx$ 是 $\Q_p^n$ 的 $n$ 个因子 $\Q_p$ 上的相同测度的乘积, 对每个 $f \in L^1(\Q_p)$ 有 $\int_{\Q_p^n}f(Ax)\,\dd x = \int_{\Q_p^n}f(x)\,\dd x$. 结论由 $\det(A) = \pm 1$ 得出.
假设 $A$ 是上三角的, 记 $A = (a_{ij})_{1 \leq i,j \leq n}$. 设 $f \in L^1(\Q_p)$. 则
\begin{multline*}
\int_{\Q_p^n}f(A(x_1, \ldots, x_n)) = \int_{\Q_p}\cdots\int_{\Q_p}f(a_{11}x_1 + \cdots + a_{1n}x_n, \ldots, \\
a_{n-1,n-1}x_{n-1} + a_{n-1,n}x_n, a_{nn}x_n)\,\dd x_n\,\dd x_{n-1}\cdots\dd x_1.
\end{multline*}
使用 (a), 看到最后的积分等于
\[
|a_{11}|_p^{-1}\cdots|a_{n-1,n-1}|_p^{-1}|a_{nn}|_p^{-1}\int_{\Q_p^n}f(x)\,\dd x = |\det(A)|_p^{-1}\int_{\Q_p^n}f(x)\,\dd x.
\]
下三角矩阵的情形类似 (只需把 $dx_i$ 反序).

(c). 一旦有了 (b) 的换元公式, 就可以在习题 \ref{ex:I.5.3.1} 和 \ref{ex:I.5.3.2} 中用 $\Q_p$ 代替 $\R$, 所有结果用完全相同的证明仍然成立. (习题 \ref{ex:I.5.3.2}(b) 除外, 它对 $\Q_p$ 没有意义.) 特别地, 我们得到 $|\det(A)|_p^{-n}dA$ 是 $\GL_n(\Q_p)$ 上的左、右 Haar 测度.

(d). 同样地, 我们只需应用习题 \ref{ex:I.5.3.7} 的 (b) 和 (c) 的证明 (以及习题 \ref{ex:I.5.3.1} 的 $\Q_p$ 版本) 就能得到结果. 假设习题 \ref{ex:I.5.3.7} 中没有符号错误, 则 $B$ 上的左 Haar 测度是 $\prod_{i=1}^n|a_{ii}|_p^{i-n-1}dA$, 其中 $dA$ 是上三角矩阵构成的 $\Q_p$-向量空间 $V \simeq \Q_p^{n(n+1)/2}$ 上的乘积测度, $a_{ij}$ 是矩阵的元素. 而 $B$ 的模函数由
\[
\Delta(A) = \prod_{i=1}^n|a_{ii}|_p^{2i-n-1}
\]
给出.
\end{solution}

\begin{exercise}[$\SO(n)$ 上的 Haar 测度]
\label{ex:I.5.3.9}
本题的目标是给出 $\SO(n)$ 上 Haar 测度的公式. (对酉群 $\UU(n)$ 可以做类似的事.)
\begin{enumerate}
\item[(a)] 对 $X \in M_n(\R)$, 令 $\Phi(X) = (I_n - X)(I_n + X)^{-1}$. 证明当 $-1$ 不是 $X$ 的特征值时这有意义, 且只要有意义就有 $\Phi(\Phi(X)) = X$.
\item[(b)] 记 $A_n$ 为 $n \times n$ 反对称矩阵 (即满足 $X^T = -X$ 的 $X \in M_n(\R)$) 构成的 $\R$-向量空间, $U$ 为 $\SO(n)$ 中不以 $-1$ 为特征值的元素的集合. 证明 $U$ 是 $\SO(n)$ 的开稠密子集, 且 $\Phi$ 诱导同胚 $A_n \xrightarrow{\sim} U$.
\item[(c)] 设 $X \in A_n$. 证明存在 $A_n$ 的开稠密子集 $V$ 和 $W$ 使得公式 $\Phi(L_XY) = \Phi(X)\Phi(Y)$ 定义微分同胚 $L_X \colon V \xrightarrow{\sim} W$, 且 $0 \in V$.
\item[(d)] 设 $dX$ 是 $A_n$ 上的 Lebesgue 测度. 对每个 $X \in A_n$ 和每个 $L_X$ 在其上有定义的 $Y \in A_n$, 记 $L'_X(Y)$ 为 $L_X$ 在 $Y$ 处的微分. 它是从 $A_n$ 到 $A_n$ 的线性变换, 使得对每个 $H \in A_n$,
\[
L_X(Y + tH) = L_X(Y) + tL'_X(Y)(H) + o(t).
\]
固定 $X \in A_n$. 我们要计算 $\det(L'_X(0))$. 回忆 $L'_X(0)$ 是 $A_n$ 的线性自同态, 且注意 $A_n \otimes_{\R}\CC$ 是 $M_n(\CC)$ 中反对称矩阵的空间.
\begin{enumerate}
\item[(i)] 证明 $\det(L'_X(0))$ 良定义且非零.
\item[(ii)] 证明对每个 $H \in A_n$, $L'_X(0)(H) = (I_n - X)H(I_n + X)$.
\item[(iii)] 证明 $X$ 有 (复) 特征向量基 $(v_1, \ldots, v_n)$, 使得对应特征值具有 $i\lambda_1, \ldots, i\lambda_n$ 的形式, 其中 $\lambda_1, \ldots, \lambda_n \in \R$.
\item[(iv)] 对 $j, k \in \{1, \ldots, n\}$, 令 $Y_{jk} = v_jv_k^T - v_kv_j^T$. 证明 $Y_{jk} \in A_n \otimes_{\R}\CC$, 且它是 $L'_X(0)$ 的特征向量, 对应特征值为 $(1 - i\lambda_j)(1 - i\lambda_k)$.
\item[(v)] 证明 $(Y_{jk})_{1 \leq j < k \leq n}$ 是 $A_n \otimes_{\R}\CC$ 的基.
\item[(vi)] 证明 $\det(L'_X(0)) = \det(I_n - iX)^{n-1}$.
\end{enumerate}
\item[(e)] 证明把 $f \in \mathcal{C}_c(\SO(n))$ 送到
\[
\int_{A_n}f(\Phi(X))\frac{1}{|\det L'_X(0)|}\,\dd X
\]
的线性泛函定义 $\SO(n)$ 上的左 Haar 测度. (提示: 注意 $(L_X \circ L_Y)(0) = L_X(Y)$, 并使用链式法则.)
\end{enumerate}
\end{exercise}

\begin{solution}
(a). 若 $X \in M_n(\R)$, 则 $-1$ 不是 $X$ 的特征值当且仅当 $I_n + X$ 可逆, 即当且仅当定义 $\Phi(X)$ 的公式有意义. 所以 $\Phi$ 的定义域是由方程 $\det(I_n + X) \neq 0$ 定义的开集. 又注意 $I_n - X$ 与 $I_n + X$ 交换, 所以 $I_n - X$ 与 $(I_n + X)^{-1}$ 交换 (若后者有定义), 因此也有 $\Phi(X) = (I_n + X)^{-1}(I_n - X)$.
设 $X \in M_n(\R)$ 使得 $\Phi(X)$ 有定义. 则
\[
I_n + \Phi(X) = ((I_n + X) + (I_n - X))(I_n + X)^{-1} = 2(I_n + X)^{-1}
\]
且
\[
I_n - \Phi(X) = ((I_n + X) - (I_n - X))(I_n + X)^{-1} = 2X(I_n + X)^{-1}.
\]
特别地, $I_n + \Phi(X)$ 可逆, 所以 $\Phi(\Phi(X))$ 有意义, 且
\[
\Phi(\Phi(X)) = (I_n - \Phi(X))(I_n + \Phi(X))^{-1} = 2X(I_n + X)^{-1}(2(I_n + X)^{-1})^{-1} = X.
\]

(b). 设 $g \in \SO(n)$. 则可以找到 $P \in \GL_n(\R)$ 使得
\[
PgP^{-1} = \begin{pmatrix} r_1 & 0 & \cdots & 0 \\ 0 & r_2 & 0 & 0 \\ \vdots & 0 & \ddots & 0 \\ 0 & \cdots & 0 & r_m \end{pmatrix},
\]
其中:
\begin{itemize}
\item 若 $n$ 偶, 则 $m = n/2$ 且 $r_1, \ldots, r_m$ 是形如 $\begin{pmatrix} \cos\theta_i & \sin\theta_i \\ -\sin\theta_i & \cos\theta_i \end{pmatrix}$ ($\theta_i \in [0, 2\pi)$) 的 $2 \times 2$ 矩阵;
\item 若 $n$ 奇, 则 $m = (n+1)/2$, 矩阵 $r_m$ 是 $1 \times 1$ 矩阵 $1$, 且 $r_1, \ldots, r_{m-1}$ 是形如 $\begin{pmatrix} \cos\theta_i & \sin\theta_i \\ -\sin\theta_i & \cos\theta_i \end{pmatrix}$ ($\theta_i \in [0, 2\pi)$) 的 $2 \times 2$ 矩阵.
\end{itemize}
两种情形下, $-1$ 是 $g$ 的特征值当且仅当至少有一个 $\theta_i$ 等于 $\pi$. 所以, 通过变化 $\theta_i$, 可以找到不以 $-1$ 为特征值的 $\SO(n)$ 的元素序列收敛到 $g$. 这证明 $U$ 在 $\SO(n)$ 中稠密.
其次, 由于反对称矩阵只有纯虚特征值, 函数 $\Phi$ 在 $A_n$ 上有定义. 又注意在 $\Phi$ 的定义上显然 $\Phi$ 在其定义开集上连续. 由 (a) 的第二部分, $\Phi$ 单射, 且为证明 $\Phi$ 是从 $A_n$ 到 $U$ 的同胚, 只需证明它是从 $A_n$ 到 $U$ 的双射 (因为此时它的逆将是 $\Phi$). 所以我们只需证明 $\Phi(A_n) = U$. 再次利用只要有意义就有 $\Phi(\Phi(X)) = X$, 看到只需证明 $\Phi(A_n) \subset U$ 和 $\Phi(U) \subset A_n$.
设 $X \in A_n$. 则 $X^T = -X$, 所以 $\Phi(X)^T = (I_n + X^T)^{-1}(I_n - X^T) = (I_n - X)^{-1}(I_n + X)$, 从而 $\Phi(X)^T\Phi(X) = I_n$, 这意味着 $\Phi(X) \in \Or(n)$. 由于 $\Phi$ 连续且 $A_n$ 连通, $\Phi(A_n)$ 连通. 但 $I_n = \Phi(0) \in \Phi(A_n)$, 所以 $\Phi(A_n)$ 包含于 $\SO(n)$.
设 $X \in \SO(n)$ 不以 $-1$ 为特征值. 则 $X^T = X^{-1}$, 所以
\[
\Phi(X)^T = (I_n - X^T)(I_n + X^T)^{-1} = (I_n - X^{-1})(I_n + X^{-1})^{-1} = (X - I_n)(X + I_n)^{-1} = -\Phi(X).
\]
所以我们有 $\Phi(U) \subset A_n$.

(c). 固定 $X \in A_n$. 注意公式 $\Phi(L_XY) = \Phi(X)\Phi(Y)$ 也可以写成 $L_XY = \Phi(\Phi(X)\Phi(Y))$ (由 (a)).
对 $Y \in A_n$, $\Phi(X)\Phi(Y)$ 在 $\Phi$ 下有像 (由 (b) 将自动在 $A_n$ 中) 当且仅当 $\Phi(Y) \in \Phi(X)^{-1}U$. 所以可以取 $V = \Phi(U \cap (\Phi(X)^{-1}U))$; 由于由 (b) $U \cap (\Phi(X)^{-1}U)$ 在 $\SO(n)$ 中稠密, 它在 $A_n$ 中稠密. 则 $V$ 在映射 $L_X \colon Y \longmapsto \Phi(\Phi(X)\Phi(Y))$ 下的像是 $W := \Phi((\Phi(X)U) \cap U)$.
映射 $L_X \colon V \to W$ 连续且满射. 事实上, 由于 $\Phi$ 无穷次可微 (由公式``对每个 $A \in \GL_n(\R)$, $A^{-1}$ 等于 $\det(A)^{-1}$ 乘以其余因子矩阵的转置'', 它由参数的元素的合理函数给出), 映射 $L_X$ 也无穷次可微.
设 $X' = \Phi(\Phi(X)^{-1}) \in A_n$. 则如上得到连续满射 $L_{X'} \colon W \to V$, 由公式 $L_{X'}(Y) = \Phi(\Phi(X)^{-1}\Phi(Y))$ 定义. 映射 $L_X$ 与 $L_{X'}$ 互为逆, 特别地它们都是微分同胚.
最后, 若 $Y = 0$, 则 $\Phi(Y) = I_n$. 所以 $\Phi(Y) \in U$, 且也有 $\Phi(Y) \in \Phi(X)^{-1}U$, 因为 $\Phi(X)\Phi(Y) = \Phi(X) \in U$. 这证明 $0 \in V$.

(d). (i) 设 $X \in A_n$. 由于 $L_X$ 在点 $0$ 处有定义, 微分 $L'_X(0)$ 有意义; 又, 由于 $L_X$ 是微分同胚, $\det(L'_X(0)) \neq 0$.

(ii) 注意对 $Y \in A_n$,
\[
(I_n + X + Y)^{-1} = (I_n + X)^{-1}(I_n + (I_n + X)^{-1}Y) = (I_n + X)^{-1}(I_n - Y(I_n + X)^{-1} + o(Y)),
\]
所以
\begin{align*}
\Phi(X + Y) &= (I_n - X - Y)(I_n + X + Y)^{-1} \\
&= ((I_n - X) - Y)(I_n + X)^{-1}(I_n - Y(I_n + X)^{-1} + o(Y)) \\
&= \Phi(X) - \Phi(X)Y(I_n + X)^{-1} - Y(I_n + X)^{-1} + o(Y).
\end{align*}
特别地 (取 $X = 0$),
\[
\Phi(Y) = I_n - 2Y + o(Y).
\]
所以
\[
\Phi(X)\Phi(Y) = \Phi(X) - 2\Phi(X)Y + o(Y),
\]
且
\begin{align*}
L_X(Y) &= \Phi(\Phi(X)\Phi(Y)) = \Phi(\Phi(X) - 2\Phi(X)Y + o(Y)) \\
&= \Phi(\Phi(X)) - \Phi(\Phi(X))(-2\Phi(X)Y)(I_n + \Phi(X))^{-1} - (-2\Phi(X)Y)(I_n + \Phi(X))^{-1} + o(Y).
\end{align*}
利用 $\Phi(\Phi(X)) = X$ 和 $I_n + \Phi(X) = 2(I_n + X)^{-1}$ (见 (a)), 可以把最后这个表达式化简为
\begin{multline*}
X + X\Phi(X)Y(I_n + X)^{-1} + \Phi(X)Y(I_n + X)^{-1} + o(Y) \\
= X + (I_n + X)\Phi(X)Y(I_n + X)^{-1} + o(Y) = X + (I_n - X)Y(I_n + X) + o(Y).
\end{multline*}
但然后 $L'_X(0)(Y) = (I_n - X)Y(I_n + X)$ 的结论由微分的定义立即得出.

(iii) 由于 $X$ 反对称且元素为实数, 它正规, 所以谱定理说 $X$ 在 $\CC^n$ 的标准正交基下可对角化; 换句话说, 存在酉矩阵 $P$ 使得 $PXP^{-1}$ 对角. 我们已经用过 $X$ 的特征值是纯虚的的事实, 但容易快速重验: $X^* = -X$, $P^* = P^{-1}$, 且 $(PXP^{-1})^* = (P^*)^{-1}X^*P^* = -PXP^{-1}$. 由于 $PXP^{-1}$ 对角, 这意味着它的对角线元素 (即 $X$ 的特征值) 全是纯虚的.

(iv) 由 $Y_{jk}$ 的定义直接得出 $Y_{jk}^T = -Y_{jk}$, 所以 $Y_{jk} \in A_n \otimes_{\R}\CC$. 而且, 由 (ii),
\begin{align*}
L'_X(0)(Y_{ij}) &= (I_n - X)Y_{ij}(I_n + X) \\
&= (I_n - X)(v_jv_k^T)(I_n - X^T) - (I_n - X)(v_kv_j^T)(I_n - X^T) \\
&= (1 - i\lambda_j)(v_jv_k^T)(1 - i\lambda_k) - (1 - i\lambda_k)(v_kv_j^T)(1 - i\lambda_j) \\
&= (1 - i\lambda_j)(1 - i\lambda_k)Y_{ij}.
\end{align*}

(v) 由于 $(v_1, \ldots, v_n)$ 是 $\CC^n$ 的基, 矩阵 $v_jv_k^T$ ($1 \leq j, k \leq n$) 构成 $M_n(\CC)$ 的基. 所以矩阵 $Y_{jk} = (v_jv_k^T) - (v_jv_k^T)^T$ ($1 \leq j, k \leq n$) 生成 $A_n \otimes_{\R}\CC$. 注意 $Y_{jj} = 0$ 且 $Y_{kj} = -Y_{jk}$, 所以 $A_n \otimes_{\R}\CC$ 实际上由矩阵 $Y_{jk}$ ($1 \leq j < k \leq n$) 张成. 由于这样的矩阵有 $n(n-1)/2$ 个且 $\dim_{\CC}(A_n \otimes_{\R}\CC) = \dim_{\R}(A_n) = n(n-1)/2$, 它们构成 $A_n \otimes_{\R}\CC$ 的基.

(vi) 由 (iv) 和 (v),
\[
\det(L'_X(0)) = \prod_{1 \leq j < k \leq n}(1 - i\lambda_j)(1 - i\lambda_k) = \prod_{r=1}^n(1 - i\lambda_r)^{n-1}
\]
(因为每个 $1 - i\lambda_r$ 在第一个大乘积中出现 $n-1$ 次: 作为第一个因子 $(1 - i\lambda_j)$ 出现 $n - r$ 次, 作为第二个因子 $(1 - i\lambda_k)$ 出现 $r - 1$ 次). 为得到结果, 只需注意 $I_n - iX$ 的特征值是 $1 - i\lambda_1, \ldots, 1 - i\lambda_n$, 所以
\[
\det(I_n - X) = \prod_{r=1}^n(1 - i\lambda_r).
\]

(e). 记此泛函为 $\Lambda$. 首先, 由 (d), 函数 $X \longmapsto |\det(L'_X(0))|^{-1}$ 在 $A_n$ 上处处有定义且连续, 所以定义 $\Lambda$ 的积分有意义.
我们需要验证 $\Lambda$ 正且被左平移不变. 首先验证正性. 设 $f \in \mathcal{C}_c^+(\SO(n))$. 则可以找到 $\varepsilon > 0$ 和 $\SO(n)$ 的非空开子集 $\Omega$ 使得 $f|_{\Omega} \geq \varepsilon$. 由于 $U$ 在 $\SO(n)$ 中开稠, 它与 $\Omega$ 的交开且非空, 所以 $\Phi(U \cap \Omega)$ 在 $A_n$ 中开且非空, 且
\[
\Lambda(f) \geq \varepsilon\int_{\Phi(U \cap \Omega)}\frac{1}{|\det(L'_X(0))|}\,\dd X > 0
\]
(因为函数 $X \longmapsto |\det(L'_X(0))|^{-1}$ 在 $\Phi(U \cap \Omega)$ 上连续且正).
现在验证左不变性. 固定 $f \in \mathcal{C}_c(\SO(n))$. 设 $g \in U$. 则 $\Lambda(L_gf) = \int_{A_n}f(g^{-1}\Phi(Y))\frac{1}{|\det(L'_Y(0))|}\,\dd Y$. 选取 $X, X' \in A_n$ 使得 $\Phi(X) = g^{-1}$, $\Phi(X') = g$. 则
\begin{align*}
\Lambda(L_gf) &= \int_{A_n}f(\Phi(X)\Phi(Y))\frac{1}{|\det(L'_Y(0))|}\,\dd Y \\
&= \int_V f(\Phi(X)\Phi(Y))\frac{1}{|\det(L'_Y(0))|}\,\dd Y \quad \text{(因为 } \operatorname{vol}(A_n - V) = 0\text{)} \\
&= \int_V f(\Phi(L_XY))\frac{1}{|\det(L'_Y(0))|}\,\dd Y.
\end{align*}
现在注意, 若 $Y \in V$, 则 $L_Y(0) = Y$ 也在 $V$ 中, 所以 $L_X(Y) = L_X \circ L_Y(0) = L_{L_XY}(0)$ 有意义, 且由链式法则
\[
L'_{L_XY}(0) = L'_X(Y) \circ L'_Y(0),
\]
从而特别地
\[
\frac{1}{|\det(L'_Y(0))|} = \frac{|\det(L'_X(Y))|}{|\det(L'_{L_XY}(0))|}.
\]
这蕴含
\[
\Lambda(L_gf) = \int_V f(\Phi(L_XY))\frac{|\det(L'_X(Y))|}{|\det(L'_{L_XY}(0))|}\,\dd Y.
\]
使用代换 $Z = L_XY$, 看到这等于
\[
\int_W f(\Phi(Z))\frac{1}{|\det(L'_Z(0))|}\,\dd Z.
\]
由于 $\operatorname{vol}(A_n - W) = 0$, 最后的积分等于 $\int_{A_n}f(\Phi(Z))\frac{1}{|\det(L'_Z(0))|}\,\dd Z$, 即 $\Lambda(f)$.
所以我们已证明函数 $\SO(n) \to \CC$, $g \longmapsto \Lambda(L_gf)$ 在开稠子集 $U$ 上为常数. 由于此函数连续 (它是连续函数 $\SO(n) \to \mathcal{C}_c(\SO(n))$, $g \longmapsto L_gf$ 与连续线性泛函 $\Lambda \colon \mathcal{C}_c(\SO(n)) \to \CC$ 的复合), 它在整个 $\SO(n)$ 上为常数, 这意味着对每个 $g \in \SO(n)$ 有 $\Lambda(L_gf) = \Lambda(f)$.
\end{solution}

\begin{exercise}[$\SU(2)$ 上的 Haar 测度]
\label{ex:I.5.3.10}
设 $G = \SU(2)$.
\begin{enumerate}
\item[(a)] 证明 $G$ 的每个元素都具有 $\begin{pmatrix} a & -\bar{b} \\ b & \bar{a} \end{pmatrix}$ 的形式, 其中 $a, b \in \CC$ 且 $|a|^2 + |b|^2 = 1$.
若按通常方式把 $\CC$ 等同于 $\R^2$, 前一问给出 $\SU(2)$ 与 $S^3$ ($\R^4$ 中的单位球面) 之间的同胚 $\alpha$.
\item[(a)] 若 $g \in \SU(2)$, 证明 $\SU(2)$ 上按 $g$ 的左平移通过 $\alpha$ 对应于 $\SO(4)$ 的元素在 $\R^4$ 上的作用到 $S^3$ 的限制 (即存在 $A \in \SO(4)$ 使得对每个 $h \in \SU(2)$, $gh = A\alpha(h)$).
\item[(b)] 设 $\mu$ 是 $S^3$ 上的通常球面测度; 即若 $\lambda$ 是 $\R^4$ 上的 Lebesgue 测度, 则由定义对 $S^3$ 的每个 Borel 子集 $E$,
\[
\mu(E) = \frac{2}{\pi^2}\lambda(\{tx, t \in [0,1], x \in E\})
\]
(注意 $\R^4$ 中单位球的体积是 $\frac{\pi^2}{2}$). 证明 $\mu$ 在 $\alpha$ 下的原像是 $\SU(2)$ 上的左、右 Haar 测度.
\item[(c)] 我们在 $S^3$ 上使用下列 (超球面) 坐标: 若 $(x_1, x_2, x_3, x_4) \in S^3$, 写
\[
\begin{cases} x_1 = \cos\theta, \\ x_2 = \sin\theta\cos\psi, \\ x_3 = \sin\theta\sin\psi\cos\phi, \\ x_4 = \sin\theta\sin\psi\sin\phi, \end{cases}
\]
其中 $0 \leq \theta \leq \pi$, $0 \leq \psi \leq \pi$, $0 \leq \phi \leq 2\pi$. 证明对每个 $f \in \mathcal{C}_c(S^3)$,
\[
\int_{S^3}f\,\dd\mu = \frac{1}{2\pi^2}\int_0^{\pi}\int_0^{\pi}\int_0^{2\pi}f(\cos\theta, \sin\theta\cos\psi, \sin\theta\sin\psi\cos\phi, \sin\theta\sin\psi\sin\phi)\sin^2\theta\sin\psi\,\dd\theta\,\dd\psi\,\dd\phi.
\]
(你可以自由使用计算机计算任何大行列式.)
\end{enumerate}
\end{exercise}

\begin{solution}
(a). 显然题设形式的每个矩阵都在 $\SU(2)$ 中. 让我们证明其逆. 设 $A = \begin{pmatrix} a & c \\ b & d \end{pmatrix} \in M_2(\CC)$. 则 $A \in \UU(2)$ 当且仅当 $A^*A = I_2$, 这意味着 $A$ 的两个列向量对 $\CC^2$ 上的通常 Hermitian 内积正交且范数为 $1$. 由于 $\CC^2$ 中一条线的正交补是 $1$ 维的, 这蕴含存在 $\lambda \in \CC^{\times}$ 使得 $\begin{pmatrix} c \\ d \end{pmatrix} = \lambda\begin{pmatrix} -\bar{b} \\ \bar{a} \end{pmatrix}$. 列的范数条件给出 $\bar{a}a + \bar{b}b = \lambda\bar{\lambda}(\bar{a}a + \bar{b}b) = 1$, 而 $\det(A) = 1$ 的条件给出 $\lambda(\bar{a}a + \bar{b}b) = 1$. 所以我们得到 $\lambda = 1$, 如所愿.

(b). 设 $V$ 是形如 $\begin{pmatrix} a & -\bar{b} \\ b & \bar{a} \end{pmatrix}$ ($a, b \in \CC$) 的矩阵的空间. 则 $\alpha$ 延拓成从 $V$ 到 $\CC^2$ 的 $\CC$-线性同构, 从而到从 $V$ 到 $\R^4$ 的 $\R$-线性同构, 把 $\begin{pmatrix} a & -\bar{b} \\ b & \bar{a} \end{pmatrix}$ 送到 $(\Re(a), \Im(a), \Re(b), \Im(b))$. 若 $A \in \SU(2)$, 则 $A$ 在 $V$ 上的左乘作用就是 $A$ 在 $\CC^2$ 上的通常作用, 所以它对应于 $\R^4$ 的保持通常 Euclid 范数的线性自同构, 即属于 $\Or(4)$. 又, 此作用的行列式就是 $\det(A) = 1$, 所以 $\R^4$ 的对应自同构在 $\SO(4)$ 中.

(c). 首先注意 $\mu$ 是 $S^3$ 上的正则 Borel 测度 ($S^3$ 的子集 $E$ 是 Borel 子集当且仅当 $\{tx, t \in [0,1], x \in E\}$ 是 $\R^4$ 的 Borel 子集, 它紧当且仅当 $\{tx, t \in [0,1], x \in E\}$ 紧, 开当且仅当 $\{tx, t \in (0,1], x \in E\}$ (与 $\{tx, t \in [0,1], x \in E\}$ 有相同测度) 开).
由 $\R^4$ 中的换元公式, 测度 $\mu$ 被 $\SO(4)$ 在 $S^3$ 上的作用不变. 由 (b), 它在 $\alpha$ 下的原像被 $\SU(2)$ 上的左平移不变, 所以是左 Haar 测度. 但群 $\SU(2)$ 紧, 所以每个左 Haar 测度也是右 Haar 测度.

(d). 设 $B^4$ 是 $\R^4$ 中的闭单位球. 设 $f \in \mathcal{C}_c(S^3)$. 我们定义函数 $g \in L^1(B^4)$:
\begin{multline*}
g(r\cos\theta, r\sin\theta\cos\psi, r\sin\theta\sin\psi\cos\phi, r\sin\theta\sin\psi\sin\phi) \\
= f(\cos\theta, \sin\theta\cos\psi, \sin\theta\sin\psi\cos\phi, \sin\theta\sin\psi\sin\phi)
\end{multline*}
对 $0 \leq r \leq 1$. (注意: $g$ 在 $0$ 处可能没有良好定义, 但这无所谓, 因为 $\{0\}$ 体积为 $0$.) 则由 $\mu$ 的定义, $\int_{S^3}f\,\dd\mu = \frac{2}{\pi^2}\int_{B^4}g\,\dd\lambda$. 我们可以用换元公式计算这最后一个积分 (并避开此换元不双射的集合, 反正它的体积为 $0$). 若 $\beta$ 是把 $(r, \theta, \phi, \psi)$ 送到 $(r\cos\theta, r\sin\theta\cos\psi, r\sin\theta\sin\psi\cos\phi, r\sin\theta\sin\psi\sin\phi)$ 的映射, 则
\[
|D\beta(r, \theta, \phi, \psi)| = r^3(\sin\theta)^2\sin\psi,
\]
所以 $\int_{B^4}g\,\dd\lambda$ 等于
\begin{multline*}
\int_0^1\int_0^{\pi}\int_0^{\pi}\int_0^{2\pi}f(\cos\theta, \sin\theta\cos\psi, \sin\theta\sin\psi\cos\phi, \sin\theta\sin\psi\sin\phi)r^3\sin^2\theta\sin\psi\,\dd r\,\dd\theta\,\dd\psi\,\dd\phi \\
= \frac{1}{4}\int_0^{\pi}\int_0^{\pi}\int_0^{2\pi}f(\cos\theta, \sin\theta\cos\psi, \sin\theta\sin\psi\cos\phi, \sin\theta\sin\psi\sin\phi)\sin^2\theta\sin\psi\,\dd\theta\,\dd\psi\,\dd\phi.
\end{multline*}
乘以 $\frac{2}{\pi^2}$ 就得到结果.
\end{solution}

\subsection{局部紧交换群的对偶}

\begin{exercise}[对偶群 (dual group)]
\label{ex:I.5.4.1}
设 $G$ 是交换拓扑群. 记 $\widehat{G}$ 为从 $G$ 到 $S^1$ 的连续群同态的集合.
由于从 $G$ 到 $S^1$ 的两个连续同态的乘积也是从 $G$ 到 $S^1$ 的连续同态 (因为 $S^1$ 交换), 集合 $\widehat{G}$ 有自然的群结构. 我们在 $\widehat{G}$ 上赋予紧收敛拓扑 (topology of compact convergence); 即若 $\chi \in \widehat{G}$, 则 $\chi$ 的邻域基由集合 $\{\psi \in \widehat{G} \mid \sup_{x \in K}|\chi(x) - \psi(x)| < c\}$ 给出, 其中 $K$ 遍历 $G$ 的紧子集, $c$ 遍历正实数.
\begin{enumerate}
\item[(a)] 证明 $\widehat{G}$ 是拓扑群. 它称为 $G$ 的对偶群 (dual group).
\item[(b)] 假设 $G = \R$.
  \begin{enumerate}
  \item[(i)] 设 $\rho \colon G \to \GL_n(\CC)$ 是连续群同态. 证明存在唯一的 $A \in M_n(\CC)$ 使得对每个 $t \in \R$, $\rho(t) = \exp(tA)$. (有几种方法. 一种方法是注意到若结论成立, 则 $c_0(0)$ 必须存在且等于 $A$, 然后由此反推.)
  \item[(ii)] 证明 $\rho$ 的像包含于 $\UU(n)$ 当且仅当 $A^* = -A$.
  \item[(iii)] 证明映射 $\R \to \widehat{G}$, $x \longmapsto (t \longmapsto e^{ixt})$ 是拓扑群的同构 (即是群同构也是同胚).
  \end{enumerate}
\item[(c)] 证明存在拓扑群的同构 $\widehat{S^1} \simeq \Z$, 它把 $\mathrm{id}_{S^1}$ 送到 $1$.
\item[(d)] 拓扑群 $\widehat{\Z}$ 是什么?
\item[(e)] 假设 $G = \Q_p$ (参看习题 \ref{ex:I.5.1.4}). 我们如下定义映射 $\chi_1 \colon \Q_p \to S^1$: 若 $x \in \Q_p$, 可以把 $x$ 写成 $x = \sum_{n=-\infty}^{+\infty}c_np^n$, 其中 $0 \leq c_n \leq p - 1$ 且对足够小的 $n$ 有 $c_n = 0$, 且这唯一确定 $c_n$ (见习题 \ref{ex:I.5.1.4}(h)). 我们令
  \[
  \chi_1(x) = \exp\Bigl(2\pi i\sum_{n=-\infty}^{-1}c_np^n\Bigr).
  \]
  \begin{enumerate}
  \item[(i)] 证明 $\chi_1 \colon \Q_p \to S^1$ 是连续群同态且 $\Ker(\chi_1) = \Z_p$.
  \item[(ii)] 对每个 $y \in \Q_p$, 定义 $\chi_y \colon \Q_p \to S^1$ 为 $\chi_y(x) = \chi_1(xy)$. 证明这也是连续群同态, 并求其核.
  \item[(iii)] 设 $\chi \in \widehat{\Q_p}$. 证明存在 $k \in \Z$ 使得 $\chi$ 在 $\{x \in \Q_p \mid |x|_p \leq p^{-k}\}$ 上等于 $1$.
  \item[(iv)] 设 $\chi \in \widehat{\Q_p}$ 满足 $\chi(1) = 1$ 且 $\chi(p^{-1}) \neq 1$. 证明存在整数序列 $(c_r)_{r \geq 0}$ 使得 $1 \leq c_0 \leq p - 1$, 对 $r \geq 1$ 有 $0 \leq c_r \leq p - 1$, 且对每个 $k \in \Z_{\geq 1}$,
  \[
  \chi(p^{-k}) = \exp\Bigl(2\pi i\sum_{r=1}^k c_{k-r}p^{-r}\Bigr).
  \]
  \item[(v)] 设 $\chi \in \widehat{\Q_p}$ 满足 $\chi(1) = 1$ 且 $\chi(p^{-1}) \neq 1$. 证明存在 $y \in \Q_p$ 使得 $|y|_p = 1$ 且 $\chi = \chi_y$.
  \item[(vi)] 证明映射 $\Q_p \to \widehat{\Q_p}$, $y \longmapsto \chi_y$ 是拓扑群的同构.
  \item[(vii)] 证明 $\chi_y|_{\Z_p} = \chi_{y'}|_{\Z_p}$ 当且仅当 $y - y' \in \Z_p$, 且映射 $y \longmapsto \chi_y$ 诱导拓扑群的同构 $\Q_p/\Z_p \xrightarrow{\sim} \widehat{\Z_p}$, 其中 $\Q_p/\Z_p$ 赋予离散拓扑.
  \end{enumerate}
\end{enumerate}
\end{exercise}

\begin{solution}
(a). 我们需要验证 $\widehat{G}$ 的群运算连续. 从乘法开始. 设 $\chi_1, \chi_2 \in \widehat{G}$, 选取 $\chi_1\chi_2$ 的形如 $\{\psi \in \widehat{G} \mid \sup_{x \in K}|\chi(x) - \psi(x)| < c\}$ ($K \subset G$ 紧, $c > 0$) 的邻域 $U$. 我们需要找到 $\chi_1$ 的邻域 $U_1$ 和 $\chi_2$ 的邻域 $U_2$ 使得 $U_1U_2 \subset U$. 取
\[
U_i = \{\psi \in \widehat{G} \mid \sup_{x \in K}|\chi_i(x) - \psi(x)| < c/2\}
\]
($i = 1, 2$). 设 $\psi_1 \in U_1$, $\psi_2 \in U_2$. 则若 $x \in K$,
\begin{multline*}
|(\psi_1\psi_2)(x) - (\chi_1\chi_2)(x)| = |\psi_1(x)(\psi_2(x) - \chi_2(x)) + \chi_2(x)(\psi_1(x) - \chi_1(x))| \\
\leq |\psi_1(x)||\psi_2(x) - \chi_2(x)| + |\chi_2(x)||\psi_1(x) - \chi_1(x)| = |\psi_2(x) - \chi_2(x)| + |\psi_1(x) - \chi_1(x)| < c,
\end{multline*}
因为 $\psi_1$ 和 $\chi_2$ 酉. 所以 $\psi_1\psi_2 \in U$.
取逆的证明类似. 设 $\chi \in \widehat{G}$, 选取 $\chi^{-1}$ 的形如 $\{\psi \in \widehat{G} \mid \sup_{x \in K}|\chi^{-1}(x) - \psi(x)| < c\}$ ($K \subset G$ 紧, $c > 0$) 的邻域 $U$. 我们需要找到 $\chi$ 的邻域 $V$ 使得 $V^{-1} \subset U$. 取 $V = \{\psi \in \widehat{G} \mid \sup_{x \in K}|\chi(x) - \psi(x)| < c\}$. 设 $\psi \in V$. 则对每个 $x \in K$,
\[
|\psi^{-1}(x) - \chi^{-1}(x)| = |\psi^{-1}(x)||\chi^{-1}(x)||\chi(x) - \psi(x)| = |\chi(x) - \psi(x)| < c.
\]
所以 $\psi^{-1} \in U$.

(b). (i) 在 $M_n(\CC)$ 上选取范数 $\|\cdot\|$. 由于 $\GL_n(\CC)$ 在 $M_n(\CC)$ 中开, 可以选取以 $I_n$ 为心、包含于 $\GL_n(\CC)$ 的非空开球 $B$. 我们只关心 $B$ 是 $M_n(\CC)$ 的凸子集这一事实. 由于 $\rho$ 连续且 $\rho(0) = I_n$, 可以找到 $c > 0$ 使得 $\rho([0, c]) \subset B$. 则
\[
\int_0^1\rho(cx)\,\dd x = c\int_0^c\rho(x)\,\dd x \in B,
\]
所以 $X := \int_0^c\rho(x)\,\dd x \in \GL_n(\CC)$. 对每个 $t \in \R$,
\[
X\rho(t) = \int_0^c\rho(x + t)\,\dd x = \int_t^{t+c}\rho(x)\,\dd x.
\]
特别地, $\rho$ 连续可微, 且
\[
\rho'(t) = X^{-1}(\rho(t + c) - \rho(t)) = X^{-1}(\rho(c) - I_n)\rho(t).
\]
满足初值条件 $\rho(0) = I_n$ 的这个微分方程的唯一解是 $\rho(t) = \exp(tA)$, 其中 $A = X^{-1}(\rho(c) - I_n)$. 最后, 矩阵由 $\rho$ 唯一确定, 因为必须有 $A = \rho'(0)$.

(ii) 若 $A^* = -A$, 则对每个 $t \in \R$,
\[
\rho(t)\rho(t)^* = \exp(tA)\exp(tA^*) = \exp(t(A + A^*)) = \exp(0) = I_n
\]
(我们用 $tA$ 与 $tA^*$ 交换这一事实得到等式 $\exp(tA)\exp(tA^*) = \exp(tA + tA^*)$), 所以 $\rho(t) \in \UU(n)$.
反之, 假设 $\rho(\R) \subset \UU(n)$. 注意 $A = \lim_{t \to 0}\frac{1}{t}(\rho(t) - I_n)$, 所以 $A^* = \lim_{t \to 0}\frac{1}{t}(\rho(t)^* - I_n)$. 由于
\[
\rho(t)^* - I_n = \rho(t)^{-1} - I_n = -\rho(t)^{-1}(\rho(t) - I_n)
\]
且当 $t \to 0$ 时 $\rho(t)^{-1} \to I_n$, 这蕴含 $A^* = -A$.

(iii) 记题设的映射 $\R \to \widehat{G}$ 为 $\alpha$.
我们在 (i) 和 (ii) 中已见每个连续群同态 $\rho \colon \R \to S^1$ 都具有 $\rho(t) = e^{zt}$ 的形式, 其中 $z \in \CC$ 唯一且满足 $\bar{z} = -z$; 最后一个条件意味着 $z = ix$ (对某个 $x \in \R$). 这意味着 $\alpha$ 双射. 又容易看出 $\alpha$ 是群同态, 所以我们只需证明 $\alpha$ 是同胚.
首先证明 $\alpha$ 连续. 设 $x \in \R$, 考虑 $\alpha(x)$ 的形如 $\{\rho \in \widehat{G} \mid \forall t \in K,\, |\alpha(x)(t) - \rho(t)| < c\}$ ($K \subset \R$ 紧, $c > 0$) 的邻域 $U$. 则对每个 $y, t \in \R$,
\[
|\alpha(x)(t) - \alpha(y)(t)|^2 = |e^{ixt} - e^{iyt}|^2 = |1 - e^{it(x-y)}|^2 = (1 - \cos(t(x-y)))^2 + (\sin(t(x-y)))^2.
\]
选取 $\varepsilon > 0$ 使得对每个 $t \in K$ 和 $z \in (-\varepsilon, \varepsilon)$, 有 $(1 - \cos(tz))^2 + (\sin(tz))^2 < c^2$. 则若 $|x - y| < \varepsilon$, 有 $\alpha(y) \in U$.
现在我们证明 $\alpha$ 开. 设 $x \in \R$, 选取 $x$ 的形如 $(x - \varepsilon, x + \varepsilon)$ 的邻域 $V$. 我们要证明 $\alpha(V)$ 包含 $\alpha(x)$ 的邻域. 选取 $\delta > 0$ 使得函数 $t \longmapsto \sin(t)$ 和 $t \longmapsto 1 - \cos(t)$ 在 $[0, 2\delta\varepsilon]$ 上都递增, 并设
\[
U = \{\rho \in \widehat{G} \mid \forall t \in K,\, |\alpha(x)(t) - \rho(t)| < c\},
\]
其中 $K = [-\delta, \delta]$, $c = \sup_{t \in [0, \varepsilon\delta/2]}((1 - \cos(t))^2 + (\sin(t))^2)^{1/2}$ (注意这也是 $[-\varepsilon\delta/2, \varepsilon\delta/2]$ 上的上确界, 因为我们估计的函数是偶的). 设 $y \in \R$ 满足 $|x - y| \geq \varepsilon$. 我们要证明 $\alpha(y) \notin U$. 可以找到 $t \in K$ 使得 $\varepsilon\delta \leq t(x - y) \leq 2\varepsilon\delta$. 则
\[
|\alpha(x)(t) - \alpha(y)(t)| = ((1 - \cos(t(x - y)))^2 + (\sin(t(x - y)))^2)^{1/2} > c,
\]
由 $\delta$ 和 $c$ 的选取. 所以 $\alpha(y) \notin U$.

(c). 注意有拓扑群的同构 $\R/2\pi\Z \xrightarrow{\sim} S^1$, 由 $t \longmapsto e^{it}$ 给出. 所以得到群同构
\[
\widehat{S^1} \simeq \{\rho \in \widehat{\R} \mid \rho(2\pi\Z) = \{1\}\} \simeq \{x \in \R \mid \forall t \in 2\pi\Z,\, e^{ixt} = 1\} = \Z
\]
(其中第二个同构来自 (b)). 还需证明这是拓扑群的同构, 即 $\widehat{S^1}$ 离散. 若你已提前读过, 你知道这是习题 \ref{ex:I.5.4.2}(e) 的特殊情形 (在 $S^1$ 的情形我不知道更简单的证明).

(d). 由于 $\Z$ 离散, 从 $\Z$ 到 $S^1$ 的连续群同态就是从 $\Z$ 到 $S^1$ 的群同态. 由于 $\Z$ 是由 $1 \in \Z$ 生成的自由交换群, 映射 $\rho \longmapsto \rho(1)$ 是从 $\Z \to S^1$ 的群同态集合到 $S^1$ 的同构. 所以作为群, $\widehat{\Z}$ 与 $S^1$ 同构. 记此同构为 $\beta \colon S^1 \to \widehat{\Z}$ (于是 $\beta$ 把 $z \in S^1$ 送到同态 $\Z \to S^1$, $n \longmapsto z^n$). 若我们证明 $\beta$ 连续, 则它将自动是同胚, 因为 $S^1$ 紧. 但 $\Z$ 的紧子集恰是它的有限子集, 所以 $\beta$ 的连续性由映射 $S^1 \to S^1$, $z \longmapsto z^n$ 的连续性立即得出.

(e). (i) 设 $x, x' \in \Q_p$, 写 $x = \sum_{n=-\infty}^{+\infty}c_np^n$, $x' = \sum_{n=-\infty}^{+\infty}c'_np^n$ (对 $c_n$ 和 $c'_n$ 的条件与题设相同). 则由习题 \ref{ex:I.5.1.4}(h), 对每个 $N \in \Z$, 若对每个 $n \leq N - 1$ 有 $c_n = c'_n$, 则 $|x - x'|_p \leq p^{-N}$. 特别地, 若 $|x - x'|_p \leq 1$ 则 $\chi_1(x) = \chi_1(x')$, 所以 $\chi_1$ 连续且把每个 $x \in \Z_p$ 送到 $1 = \chi_1(0)$.
我们还需证明 $\chi_1$ 是群同态. 设 $G_0$ 是 $\Q_p$ 的由满足 $x = \sum_{n=-\infty}^{+\infty}a_np^n$ (其中 $a_n \in \Z$ 且对足够大的 $|n|$ 有 $a_n = 0$) 的元素组成的子群. 这是稠密子群 (因为 $\sum_{n=-\infty}^{+\infty}c_np^n$ 是 $N \to +\infty$ 时 $\sum_{n=-N}^{N}c_np^n$ 的极限), 且包含于 $\Q$ 中. 由于我们知道 $\chi_1$ 连续, 只需对 $x, y \in G_0$ 证明 $\chi_1(x + y) = \chi_1(x)\chi_1(y)$. 但注意若 $x \in G_0$, 则 $\chi_1(x) = \exp(2\pi ix)$, 其中把 $x$ 看作 $\Q$ 的元素. 这蕴含结果.
最后, 我们必须证明 $\Ker(\chi_1) = \Z_p$. 我们已经看到 $\Z_p \subset \Ker(\chi_1)$. 反之, 设 $x \in \Q_p$, 如上写 $x = \sum_{n=-\infty}^{+\infty}c_np^n$. 假设 $x \notin \Z_p$; 则存在 $m < 0$ 使得 $c_m \neq 0$. 选取这样的 $m$. 我们有
\[
p^{-m} \leq c_mp^{-m} < \sum_{n=-\infty}^{-1}c_np^n \leq (p - 1)\sum_{r \geq 1}p^{-r} = 1
\]
(第二个不等式严格, 因为对足够小的 $n$ 有 $c_n = 0$). 所以 $\sum_{n=-\infty}^{-1}c_np^n \in (0, 1)$, 且 $\chi_1(x) = \exp\bigl(2\pi i\sum_{n=-\infty}^{-1}c_np^n\bigr) \neq 1$.

(ii) 映射 $\chi_y$ 是连续群同态, 因为它是连续群同态 $\chi_1$ 和 $m_y \colon \Q_p \to \Q_p$, $x \longmapsto xy$ 的复合. $\Q_p$ 的元素 $x$ 在 $\chi_y$ 的核中当且仅当 $xy \in \Ker(\chi_1) = \Z_p$. 所以若 $y = 0$, 则 $\Ker(\chi_y) = \Q_p$; 若 $y \neq 0$, 则
\[
\Ker(\chi_y) = y^{-1}\Z_p = |y|_p\Z_p = \{x \in \Q_p \mid |x|_p \leq |y|_p^{-1}\}.
\]

(iii) 选取 $\CC^{\times}$ 的邻域 $U$ 使得 $\CC^{\times}$ 包含于 $U$ 的唯一子群是平凡群. (见 (b).) 则 $\chi^{-1}(U \cap S^1)$ 是 $1$ 在 $\Q_p$ 中的邻域, 所以存在 $k \in \Z$ 使得 $\chi^{-1}(U \cap S^1) \supset \{x \in \Q_p \mid |x|_p \leq p^k\}$. 但由于 $\{x \in \Q_p \mid |x|_p \leq p^k\}$ 是 $\Q_p$ 的子群, 它在 $\chi$ 下的像是 $S^1$ 的包含于 $U$ 的子群, 所以等于 $\{1\}$.

(iv) 对每个整数 $r \geq 0$ 写 $z_r = \chi(p^{-r})$. 则 $z_r \in S^1$, 且对每个 $r \geq 0$,
\[
z_{r+1}^p = \chi(p^{-r-1})^p = \chi(p^{-r}) = z_r.
\]
我们将对 $r \geq 0$ 归纳构造整数 $c_r$. 首先注意由假设 $z_1 \neq 1 = z_0$, 所以可以找到 $c_0 \in \{1, \ldots, p-1\}$ 使得 $z_1 = \exp(2\pi ic_0p^{-1})$. 假设已经找到 $c_0, \ldots, c_{r-1}$ ($r \geq 1$) 使得对 $1 \leq s \leq r$ 有
\[
\chi(p^{-s}) = z_s = \exp\Bigl(2\pi i\sum_{k=1}^s c_{s-k}p^{-k}\Bigr).
\]
我们必须找到 $c_r \in \{0, \ldots, p-1\}$ 使得
\[
z_{r+1} = \exp\Bigl(2\pi i\sum_{k=1}^{r+1}c_{r+1-k}p^{-k}\Bigr) = \exp\Bigl(2\pi ip^{-(r+1)}\sum_{s=0}^r c_sp^s\Bigr).
\]
由于 $z_{r+1}^p = z_r$, 有
\[
\Bigl(z_{r+1}\exp\Bigl(2\pi ip^{-r-1}\sum_{s=0}^{r-1}c_sp^s\Bigr)\Bigr)^p = 1,
\]
所以存在 $c_r \in \{0, \ldots, p-1\}$ 使得
\[
z_{r+1}\exp\Bigl(2\pi ip^{-r-1}\sum_{s=0}^{r-1}c_sp^s\Bigr) = \exp(2\pi ip^{-1}c_r),
\]
即
\[
z_{r+1} = \exp\Bigl(2\pi ip^{-(r+1)}\sum_{s=0}^r c_sp^s\Bigr).
\]

(v) 设 $(c_r)_{r \geq 0}$ 如 (iv), 令 $y = \sum_{r=0}^{+\infty}c_rp^r$. 由于 $c_0 \in \{1, \ldots, p-1\}$, 有 $|y|_p = 1$. 又, 对每个 $r \geq 1$,
\[
\chi(p^{-r}) = \exp\Bigl(2\pi ip^{-(r+1)}\sum_{k=-r}^{-1}c_{r+k}p^k\Bigr) = \chi_1(p^{-r}y) = \chi_y(p^{-r}),
\]
因为
\[
p^{-r}y = \sum_{s \geq 0}c_sp^{s-r} = \sum_{n=-r}^{+\infty}c_{r+n}p^n.
\]
另一方面, 若 $r \geq 0$, 则
\[
\chi(p^r) = \chi(1)^{p^r} = 1 = \chi_y(p^r).
\]
由于 $\chi$ 和 $\chi_y$ 是连续群同态, 且族 $(p^r)_{r \in \Z}$ 生成 $\Q_p$ 的稠密子群, 这蕴含 $\chi = \chi_y$.

(vi) 记映射 $\Q_p \to \widehat{\Q_p}$, $y \longmapsto \chi_y$ 为 $\alpha$. 容易看出 $\alpha$ 是群同态 (这由 $\chi_1$ 是群同态和 $\Q_p$ 上乘法的分配律立即得出).
首先证明 $\Ker(\alpha) = \{0\}$. 设 $y \in \Q_p - \{0\}$. 则 $y = \sum_{n=m}^{+\infty}c_np^n$, 其中 $m \in \Z$, $0 \leq c_n \leq p - 1$ 且 $c_m \geq 1$. 所以
\[
p^{-m-1}y = c_mp^{-1} + \sum_{n \geq 0}c_{n+m+1}p^n,
\]
且 $\chi_y(p^{-m-1}) = \exp(2\pi ip^{-1}c_m) \neq 1$. 这证明 $y \notin \Ker(\alpha)$.
现在我们证明 $\alpha$ 满射. 设 $\chi \in \widehat{\Q_p}$. 若 $\chi = 1$, 则 $\chi = \chi_0$, 所以假设 $\chi \neq 1$. 由 (iii), 存在 $k \in \Z$ 使得 $\chi$ 在 $\{x \in \Q_p \mid |x|_p \leq p^{-k}\}$ 上等于 $1$. 选取最小的这种 $k$ (这是可能的, 否则 $\chi$ 在整个 $\Q_p$ 上等于 $1$, 与我们 $\chi \neq 1$ 的假设矛盾). 则存在 $a \in \Q_p$ 使得 $|a|_p = p^{-k+1}$ 且 $\chi(a) \neq 1$. 定义 $\psi \in \widehat{\Q_p}$ 为 $\psi(x) = \chi(pax)$. 则 $\psi(p^{-1}) = \chi(a) \neq 1$ 且 $\psi(1) = \chi(pa) = 1$ (因为 $|pa|_p = p^{-k}$). 由 (v), 存在 $y \in \Z_p$ 使得 $\psi = \chi_y$. 换句话说, 对每个 $x \in \Q_p$,
\[
\chi(x) = \psi(p^{-1}a^{-1}x) = \chi_1(p^{-1}a^{-1}yx),
\]
即 $\chi = \alpha(p^{-1}a^{-1}y)$.
我们证明 $\alpha$ 连续. 设 $y \in \Q_p$, 选取 $\alpha(y)$ 的形如
\[
U = \{\chi \in \widehat{\Q_p} \mid \forall x \in K,\, |\chi(x) - \chi_y(x)| < c\}
\]
的邻域, 其中 $K$ 是 $\Q_p$ 的紧子集, $c > 0$. 我们要找 $y$ 在 $\Q_p$ 中的邻域 $V$ 使得 $\alpha(V) \subset U$.
由于 $\Q_p = \bigcup_{k \in \Z}p^k\Z_p$, 可以假设 $K = p^N\Z_p$ (对某个 $N \in \Z$). 我们知道 $\chi_1$ 在 $\Q_p$ 中 $\Z_p$ 的陪集上为常数, 所以若 $x \in p^N\Z_p$, 则 $\chi_x$ 在 $\Q_p$ 中 $p^{-N}\Z_p$ 的陪集上为常数. 所以若 $y' \in y + p^{-N}\Z_p$, 则对每个 $x \in K = p^N\Z_p$,
\[
|\chi_{y'}(x) - \chi_y(x)| = |\chi_x(y') - \chi_x(y)| = 0 < c.
\]
换句话说, $\alpha(y + p^{-N}\Z_p) \subset U$.
最后, 我们证明 $\alpha$ 开. 设 $y \in \Q_p$, $V$ 是 $y$ 的邻域. 可以假设 $V$ 具有 $y + p^N\Z_p = \{y' \in \Q_p \mid |y' - y|_p \leq p^{-N}\}$ 的形式 (对某个 $N \in \Z$). 我们要证明 $\alpha(V)$ 包含 $\alpha(y)$ 的邻域. 由于 $\alpha$ 是群同态, 可以假设 $y = 0$. 设
\[
U = \{\chi \in \widehat{\Q_p} \mid \forall x \in p^{-N}\Z_p,\, |\chi(x) - 1| < c\},
\]
其中 $c = \min_{1 \leq r \leq p-1}|1 - e^{2\pi irp^{-1}}|$, 让我们证明 $\alpha(p^N\Z_p) \supset U$. 设 $y' \notin p^N\Z_p$, 我们要证明 $\chi_{y'} \notin U$. 把 $y'$ 写成 $y' = \sum_{n=m}^{+\infty}c_np^n$, 其中对每个 $n \geq m$ 有 $c_n \in \{0, \ldots, p-1\}$ 且 $c_m \geq 1$. 则对 $y'$ 的假设说 $m < N$. 设 $x = p^{-m-1}$. 则 $x \in p^{-N}\Z_p$, 且
\[
\chi_{y'}(x) = \chi_1(xy') = \exp(2\pi ip^{-1}c_m),
\]
所以 $|\chi_{y'}(x) - 1| \geq c$ 且 $\chi_{y'} \notin U$.

(vii) 由于映射 $y \longmapsto \chi_y$ 是群同态, 第一个结论等价于 $\chi_y|_{\Z_p} = 1$ 当且仅当 $y \in \Z_p$ 这一事实. 我们知道 $\Ker(\chi_1) = \Z_p$, 所以对每个 $y \in \Z_p$ 有 $\Ker(\chi_y) \supset \Z_p$. 反之, 设 $y \in \Q_p - \Z_p$. 则 $|y|_p > 1$, 所以 $|y|_p \geq p$, 从而 $|py|_p \geq 1$, 且 $p^{-1}y^{-1} \in \Z_p$. 由于 $\chi_y(p^{-1}y^{-1}) = \chi_1(p^{-1}) = \exp(2\pi ip^{-1}) \neq 1$, 有 $\Ker(\chi_y) \not\supset \Z_p$.
所以映射 $y \longmapsto \chi_y$ 诱导从 $\Q_p/\Z_p$ 到 $\widehat{\Z_p}$ 的单射群同态. 我们知道 (或很快会知道) 由习题 \ref{ex:I.5.4.2}(e), $\widehat{\Z_p}$ 离散, 所以只需证明 $\widehat{\Z_p}$ 的每个元素都具有 $\chi_y|_{\Z_p}$ 的形式 (对某个 $y \in \Q_p$).
设 $\chi \in \widehat{\Z_p}$. 同 (iii), 可以找到 $k \in \N$ 使得 $\Ker(\chi) \supset p^k\Z_p$. 设 $z = \chi(1)$. 则 $z^{p^k} = \chi(p^k) = 1$, 所以可以找到 $c \in \{0, \ldots, p^k - 1\}$ 使得 $z = e^{2\pi icp^{-k}}$. 把 $c$ 写成 $c = \sum_{r=0}^{k-1}c_rp^r$, 其中 $c_r \in \{0, \ldots, p-1\}$. 则
\[
\chi(1) = \exp\Bigl(2\pi i\sum_{r=0}^{k-1}c_rp^{r-k}\Bigr) = \exp\Bigl(2\pi i\sum_{n=-k}^{-1}c_{k+n}p^n\Bigr).
\]
设 $y = \sum_{n=-k}^{-1}c_{k+n}p^n$. 则 $\chi(1) = \chi_y(1)$. 由于 $\chi_y|_{\Z_p}$ 和 $\chi$ 是 $\Z_p$ 上的连续群同态, 且 $1$ 生成 $\Z_p$ 的稠密子群, 这蕴含 $\chi = \chi_y|_{\Z_p}$.
\end{solution}

\begin{exercise}[对偶群的局部紧性]
\label{ex:I.5.4.2}
我们使用前一题的记号, 并假设 $G$ 是交换局部紧群, 固定 $G$ 上的 Haar 测度 $\mu$.
回忆有同构 $L^{\infty}(G) \to L^1(G)^{\vee} := \Hom(L^1(G), \CC)$, 把 $f \in L^{\infty}(G)$ 送到 $L^1(G)$ 上的有界算子 $g \longmapsto \int_G fg\,\dd\mu$. (这不使用 $G$ 是交换群这一事实.) 所以我们可以在 $L^{\infty}(G)$ 上考虑弱*拓扑 (或逐点收敛拓扑): 对 $f \in L^{\infty}(G)$, $f$ 的邻域基由集合 $U_{g_1,\ldots,g_n,c} = \{f' \in L^{\infty}(G) \mid |\int_G (f - f')g_i\,\dd\mu| < c,\, 1 \leq i \leq n\}$ 给出, 其中 $n \in \Z_{\geq 1}$, $g_1, \ldots, g_n \in L^1(G)$, $c > 0$.
\begin{enumerate}
\item[(a)] 证明 $\widehat{G} \subset L^{\infty}(G)$, 且 $\widehat{G}$ 的拓扑由 $L^{\infty}(G)$ 的弱*拓扑诱导.\footnote{难题.}
\item[(b)] 证明 $L^{\infty}(G)$ 的子集 $\widehat{G} \cup \{0\}$ 对弱*拓扑闭. (提示: 把它等同于 Banach $*$-代数 $L^1(G)$ 在 $\CC$ 上的表示集合.)
\item[(c)] 证明 $\widehat{G}$ 是局部紧拓扑群. (提示: Alaoglu 定理.)
\item[(d)] 若 $G$ 离散, 证明 $\widehat{G}$ 紧.
\item[(e)] 若 $G$ 紧, 证明 $\widehat{G}$ 离散.
\end{enumerate}
\end{exercise}

\begin{solution}
(a). $\widehat{G}$ 的元素是从 $G$ 到 $S^1$ 的连续函数, 从而是 $G$ 上的连续有界函数, 因此是 $L^{\infty}(G)$ 的元素. 现在我们证明 $\widehat{G}$ 上的两个拓扑一致.
首先设 $\chi \in \widehat{G}$, 设 $f_1, \ldots, f_n \in L^1(G)$, $c > 0$. 这定义 $\chi$ 的弱*开邻域
\[
U = \Bigl\{\psi \in \widehat{G} \mid \forall i \in \{1, \ldots, n\},\, \Bigl|\int_G \chi f_i\,\dd\mu - \int_G \psi f_i\,\dd\mu\Bigr| < c\Bigr\}.
\]
我们要找紧收敛拓扑下 $\chi$ 的开邻域 $V$ 使得 $V \subset U$. 设 $\varepsilon > 0$. 选取 $G$ 的紧子集 $K$ 使得对每个 $i \in \{1, \ldots, n\}$ 有 $\int_{G-K}|f_i|\,\dd\mu < \varepsilon$ (由 $\mu$ 的内正则性这是可能的). 设
\[
V = \{\psi \in \widehat{G} \mid \forall x \in K,\, |\chi(x) - \psi(x)| < \varepsilon\}.
\]
则若 $\psi \in V$, $i \in \{1, \ldots, n\}$, 有
\begin{align*}
\Bigl|\int_G \chi f_i\,\dd\mu - \int_G \psi f_i\,\dd\mu\Bigr| &\leq \int_K |f_i(x)||\chi(x) - \psi(x)|\,\dd\mu(x) + \int_{G-K}|f_i(x)||\chi(x) - \psi(x)|\,\dd x \\
&\leq \varepsilon\int_K |f_i(x)|\,\dd\mu(x) + 2\int_{G-K}|f_i(x)|\,\dd\mu(x) \\
&\leq \varepsilon(\|f_i\|_1 + 2).
\end{align*}
所以若取 $\varepsilon$ 足够小, 将得到 $V \subset U$.
现在我们证明其逆. 设 $\chi \in \widehat{G}$, $K$ 是 $G$ 的紧子集, $c > 0$. 考虑 $\chi$ 在紧收敛拓扑下的邻域
\[
V = \{\psi \in \widehat{G} \mid \forall x \in K,\, |\chi(x) - \psi(x)| < c\}.
\]
必须找到包含于它的弱*邻域. 设 $\eta > 0$ (稍后调整), 选取 $1$ 的紧邻域 $A$ 使得对每个 $y \in A$ 有 $|\chi(y) - 1| < \eta$. 设 $f = \mathbb{1}_A$; 由于 $A$ 紧, $f \in L^1(G)$. 注意对每个 $x \in G$,
\[
|\mu(A)\chi(x) - f * \chi(x)| = \Bigl|\int_A (\chi(x) - \chi(y^{-1}x))\,\dd y\Bigr| \leq \int_A |1 - \chi(y)|\,\dd y \leq \eta\mu(A).
\]
现在我们试着找到 $\widehat{G}$ 中 $\chi$ 的弱*邻域, 使其元素 $\psi$ 对 $x \in K$ 满足类似的不等式. 注意若 $\psi \in \widehat{G}$, $x \in G$, 则
\[
f * \psi(x) = \int_A \psi(y^{-1}x)\,\dd y = \psi(x)\int_A \psi(y)\,\dd y = \int_G \psi(y^{-1})f(xy)\,\dd y = \int_G \psi(y)L_{x^{-1}}f(y)\,\dd y
\]
(用 $G$ 交换且 $\psi$ 是从 $G$ 到 $S^1$ 的群同态这一事实). 现在回忆映射 $G \to L^1(G)$, $x \longmapsto L_{x^{-1}}f$ 连续 (命题 \ref{prop:I.3.1.13}). 由于 $K$ 紧, 可以找到 $x_1, \ldots, x_n$ 使得对每个 $x \in K$, 存在 $i \in \{1, \ldots, n\}$ 满足 $\|L_{x^{-1}}f - L_{x_i^{-1}}f\|_1 < \eta\mu(A)$. 考虑 $\chi$ 的如下弱*邻域:
\[
U = \Bigl\{\psi \in \widehat{G} \mid \forall x \in \{1, x_1, \ldots, x_n\},\, \Bigl|\int_G \chi(y)L_{x^{-1}}f(y)\,\dd y - \int_G \psi(y)L_{x^{-1}}f(y)\,\dd y\Bigr| < \eta\mu(A)\Bigr\}.
\]
设 $\psi \in U$. 首先, 对每个 $x \in G$,
\begin{align*}
|\mu(A)\psi(x) - f * \psi(x)| &= \Bigl|\int_A (\psi(x) - \psi(y^{-1}x))\,\dd y\Bigr| \\
&= \Bigl|\int_A (1 - \psi(y))\,\dd y\Bigr| \\
&\leq \int_A |1 - \chi(y)|\,\dd y + \Bigl|\int_A (\chi(y) - \psi(y))\,\dd y\Bigr| \\
&\leq 2\eta\mu(A).
\end{align*}
其次, 要对 $x \in K$ 估计 $|f * \chi(x) - f * \psi(x)|$. 固定 $x \in K$. 设 $i \in \{1, \ldots, n\}$ 满足 $\|L_{x^{-1}}f - L_{x_i^{-1}}f\|_1 \leq \eta\mu(A)$. 则:
\begin{align*}
|f * \chi(x) - f * \psi(x)| &= \Bigl|\int_G (\chi(y) - \psi(y))L_{x^{-1}}f(y)\,\dd y\Bigr| \\
&\leq \Bigl|\int_G (\chi(y) - \psi(y))L_{x_i^{-1}}f(y)\,\dd y\Bigr| + \Bigl|\int_G (\chi(y) - \psi(y))(L_{x_i^{-1}}f(y) - L_{x^{-1}}f(y))\,\dd y\Bigr| \\
&< \eta\mu(A) + 2\int_G |L_{x_i^{-1}}f(y) - L_{x^{-1}}f(y)|\,\dd y \\
&\leq 3\eta\mu(A).
\end{align*}
把所有估计放在一起, 对 $x \in K$ 得到 $|\mu(A)\chi(x) - \mu(A)\psi(x)| < 6\eta\mu(A)$, 即 $|\chi(x) - \psi(x)| < 6\eta$. 在一开始选取 $\eta$ 使得 $6\eta \leq c$, 得到 $U \subset V$, 如所愿.
(b). 我们在课堂上见过 $\widehat{G} \subset L^{\infty}(G) \simeq L^1(G)^{\vee}$ 是 Banach $*$-代数 $L^1(G)$ 的非退化表示集合. 设 $\pi \colon L^1(G) \to \CC$ 是 $L^1(G)$ 在 $\CC$ 上的表示, 假设它不是非退化的. 则存在 $v \in \CC - \{0\}$ 使得对每个 $f \in L^1(G)$ 有 $\pi(f)v = 0$. 但这蕴含 $\pi = 0$. 所以我们看到 $\widehat{G} \cup \{0\} \subset L^{\infty}(G)$ 确实是 $L^1(G)$ 在 $\CC$ 上的表示集合. 但说一个有界线性泛函 $\Lambda \colon L^1(G) \to \CC$ 是表示的条件在弱*拓扑下都是闭条件 (因为它们都断言 $\Lambda$ 在 $L^1(G)$ 的一些点处的取值相等), 所以 $L^1(G)$ 的表示集合是 $L^{\infty}(G)$ 的弱*闭子集.
(c). Alaoglu 定理\footnote{参考文献 ?} 说 $L^{\infty}(G)$ 的闭单位球 (对来自 $\|\cdot\|_1$ 的算子范数, 即 $\|\cdot\|_{\infty}$) 对弱*拓扑紧 Hausdorff. 但 $\widehat{G} \cup \{0\}$ 显然包含于这个闭单位球, 所以它对弱*拓扑紧 Hausdorff. 因此它的开子集 $\widehat{G}$ 对弱*拓扑局部紧, 且已在 (a) 中看到 $\widehat{G}$ 上的弱*拓扑等于紧收敛拓扑, 所以我们完成了.
(d). 考虑映射 $\alpha \colon \widehat{G} \to (S^1)^G$, 把 $\chi$ 送到族 $(\chi(x))_{x \in G}$. 这显然单射. 由于 $G$ 离散, 它的紧子集恰是有限子集, 所以紧收敛拓扑恰是 $(S^1)^G$ 上乘积拓扑诱导的拓扑. 又, 由 Tychonoff 定理, $(S^1)^G$ 紧 Hausdorff. 所以要得到结果, 只需证明 $\alpha$ 的像在 $(S^1)^G$ 中闭. 但 $\alpha$ 的像是子集族
\[
\{(a_x)_{x \in G} \in (S^1)^G \mid a_{x_0}a_{y_0} = a_{x_0y_0}\}
\]
(对所有 $x_0, y_0 \in G$) 的交, 且每个这样的子集都闭, 所以 $\Im(\alpha)$ 闭.
(e). 假设 $G$ 紧. 则 $\widehat{G}$ 的拓扑是一致收敛拓扑 (由范数 $\|\cdot\|_{\infty}$ 诱导). 要证明 $\widehat{G}$ 离散, 只需证明其子集 $\{1\}$ 开 (因为 $\widehat{G}$ 是拓扑群). 设 $c > 0$ 使得 $\CC^{\times}$ 包含于 $\{z \in \CC^{\times} \mid |1 - z| < c\}$ 的唯一子群是平凡群. 设 $U = \{\chi \in \widehat{G} \mid \|\chi - 1\|_{\infty} < c\}$. 这是 $1$ 在 $\widehat{G}$ 中的开邻域. 另一方面, 若 $\chi \in U$, 则 $\chi(G) \subset \{z \in \CC^{\times} \mid |1 - z| < c\}$; 由于 $\chi(G)$ 是 $\CC^{\times}$ 的子群, 这意味着 $\chi(G) = \{1\}$, 即 $\chi = 1$, 于是 $U = \{1\}$.
\end{solution}

\subsection{表示}

若 $G$ 是一个群, 我们说 $G$ 的表示 $(\pi, V)$ 是忠实的 (faithful), 如果 $\pi \colon G \to \GL(V)$ 是单射.

\begin{exercise}[$\SU(2)$ 的不可约表示]
\label{ex:I.5.5.1}
设 $G = \SU(2)$. 群 $G$ 通过包含 $G \subset \GL_2(\CC)$ 作用于 $\CC^2$, 我们把这个作用记为 $(g, (z_1, z_2)) \longmapsto g(z_1, z_2)$. (这称为 $G$ 的标准表示 (standard representation).)
对每个整数 $n \geq 0$, 设 $V_n$ 为 $n$ 次齐次 (homogeneous) 多项式 $P \in \CC[t_1, t_2]$ 的空间 (即 $P(t_1, t_2) = \sum_{r=0}^n a_rt_1^rt_2^{n-r}$, 其中 $a_0, \ldots, a_n \in \CC$).
\begin{enumerate}
\item[(a)] 若 $P \in V_n$, $g \in G$, 证明函数 $\CC^2 \to \CC$, $(z_1, z_2) \longmapsto P(g^{-1}(z_1, z_2))$ 仍由 $V_n$ 中的多项式给出, 且这定义 $G$ 在 $V_n$ 上的连续表示.
\item[(b)] 证明对每个 $n \geq 0$, $G$ 的表示 $V_n$ 不可约.
\item[(c)] 对哪些 $n$, 表示 $V_n$ 忠实?
\end{enumerate}
\end{exercise}

\begin{remarkstar}
我们将在后面 (见习题 \ref{ex:IV.9.1}) 看到 $\SU(2)$ 的每个不可约酉表示都同构于某个 $V_n$.
\end{remarkstar}

\begin{solution}
(a). 先取 $P = t_1^rt_2^{n-r}$, 其中 $0 \leq r \leq n$. 设 $x = \begin{pmatrix} a & b \\ c & d \end{pmatrix} \in \SU(2)$. 由于 $\det(x) = 1$, 有 $x^{-1} = \begin{pmatrix} d & -b \\ -c & a \end{pmatrix}$. 所以
\[
P(x^{-1}(z_1, z_2)) = (dz_1 - bz_2)^r(-cz_1 + az_2)^{n-r}.
\]
这仍是 $z_1$ 和 $z_2$ 的 $n$ 次齐次多项式, 记之为 $P \circ x^{-1}$. 又, 由上式显然映射 $G \to V_n$, $x \longmapsto P \circ x^{-1}$ 连续 (这意味着 $P \circ x^{-1}$ 的系数是矩阵 $x$ 的元素的连续函数).
由于单项式 $t_1^rt_2^{n-r}$ ($0 \leq r \leq n$) 生成 $V_n$, 上一段蕴含对每个 $P \in V_n$ 和每个 $x \in \SU(2)$, 函数 $\CC^2 \to \CC$, $(z_1, z_2) \longmapsto P(x^{-1}(z_1, z_2))$ 仍由 $V_n$ 的元素给出 (记之为 $P \circ x^{-1}$); 它也蕴含映射 $G \to V_n$, $x \longmapsto P \circ x^{-1}$ 连续.
对每个 $x \in G$, 映射 $V_n \to V_n$, $P \longmapsto P \circ x^{-1}$ 显然是 $\CC$-线性的 (事实上我们已经用过这一事实). 又对每个 $P \in V_n$ 和所有 $x, y \in G$, 有 $P \circ (xy)^{-1} = (P \circ y^{-1}) \circ x^{-1}$. 所以由命题 \ref{prop:I.3.5.1}, 映射 $G \times V_n \to V_n$, $(x, P) \longmapsto P \circ x^{-1}$ 连续, 即定义 $G$ 在 $V_n$ 上的连续表示.

(b). 设 $W$ 是 $V$ 的 $G$-不变子空间. 设 $P = \sum_{r=0}^n c_rt_1^rt_2^{n-r} \in W$. 我们证明对每个使得 $c_r \neq 0$ 的 $r \in \{0, \ldots, n\}$, 有 $t_1^rt_2^{n-r} \in W$. 我们对 $P$ 的非零系数的个数归纳证明. 若 $P$ 有 $0$ 或 $1$ 个非零系数, 则已完成. 假设 $P$ 至少有 $2$ 个非零系数. 固定使得 $c_r \neq 0$ 的 $r \in \{0, \ldots, n\}$. 只需找到 $W$ 的另一个元素 $Q$, 使得 $t_1^rt_2^{n-r}$ 的系数非零, 且 $Q$ 的非零系数个数比 $P$ 少; 然后可以对 $Q$ 应用归纳假设. 选取 $s \in \{0, \ldots, n\} - \{r\}$ 使得 $c_s \neq 0$. 考虑 $x_a = \begin{pmatrix} a & 0 \\ 0 & a \end{pmatrix}$, 其中 $a \in S^1$. 则 $x_a \in \SU(2)$, 且
\[
P \circ x_a^{-1} = \sum_{i=0}^n a^ia^{n-i}c_it_1^it_2^{n-i} = \sum_{i=0}^n a^{n-2i}c_it_1^it_2^{n-i}.
\]
选取 $a, a' \in S^1$ 使得 $a^{n-2s}c_s - (a')^{n-2s}c_s = 0$ 且 $a^{n-2r}c_r - (a')^{n-2r}c_r \neq 0$. 则 $Q := P \circ x_a^{-1} - P \circ x_{a'}^{-1} \in W - \{0\}$ 具有所要的性质.

现在假设 $W \neq 0$. 由上一段, 可以找到 $r \in \{0, \ldots, n\}$ 使得 $P := t_1^rt_2^{n-r} \in W$. 设 $x = \begin{pmatrix} a & b \\ -\bar{b} & \bar{a} \end{pmatrix}$, 其中 $a, b \in \CC$ 且 $|a|^2 + |b|^2 = 1$. 则 $x \in \SU(2)$ 且 $P \circ x^{-1} = (at_1 - bt_2)^r(\bar{b}t_1 + \bar{a}t_2)^{n-r} \in W$. 若把 $P \circ x^{-1}$ 写成 $P \circ x^{-1} = \sum_{i=0}^n c_it_1^it_2^{n-i}$, 则
\[
c_i = \sum_{j=\max(0,i-n+r)}^{\min(i,r)}(-1)^{r-j}\binom{r}{j}\binom{n-r}{i-j}a^{r-j}\bar{a}^{j}\bar{b}^{n-r+j-i}b^{i-j}.
\]
若取 $a = \frac{1}{\sqrt{1+t^2}}$ 且 $b = \frac{t}{\sqrt{1+t^2}}$ (其中 $t \in [-1, 1]$), 则每个 $c_i$ 是非零多项式在 $(1+t^2)^{n/2}$ 下的商, 所以只有有限多个 $t$ 使得 $c_i = 0$. 所以可以选择 $x \in \SU(2)$ 使得 $P \circ x^{-1}$ 的所有系数非零. 由第一段, 这蕴含每个单项式 $t_1^it_2^{n-i}$ ($0 \leq i \leq n$) 都在 $W$ 中. 所以 $W = V_n$.

(c). 记 $\pi_n$ 为映射 $\SU(2) \to \GL(V_n)$. 若 $n = 0$, 则 $V_n$ 是 $\SU(2)$ 的平凡表示, 所以 $\Ker(\pi_n) = \SU(2)$. 假设 $n \geq 1$, 设 $x = \begin{pmatrix} a & b \\ -\bar{b} & \bar{a} \end{pmatrix} \in \Ker(\pi_n)$. 特别地, 若 $P = t_1^n$, 则必须有 $P \circ x^{-1} = P$. 由于 $P \circ x^{-1} = \sum_{i=0}^n(-1)^{n-i}\binom{n}{i}a^i\bar{b}^{n-i}t_1^it_2^{n-i}$, 这蕴含 $a^n = 1$ 且对 $0 \leq i \leq n-1$ 有 $a^i\bar{b}^{n-i} = 0$. 特别地, $a \neq 0$, 所以必有 $b = 0$. 则 $a \in S^1$, 且对每个 $r \in \{0, \ldots, n\}$, $(t_1^rt_2^{n-r}) \circ x^{-1} = a^{n-2r}t_1^rt_2^{n-r}$, 所以 $a^{n-2r} = 1$. 若 $n$ 奇, 这蕴含 $a = 1$, 所以 $x = I_2$ 是 $\Ker(\pi_n)$ 的唯一元素. 若 $n$ 偶, 这只能蕴含 $a = \pm 1$, 所以 $x = \pm I_2$. 事实上, 若 $n$ 偶且非零, 容易验证 $-I_2$ 在 $V_n$ 上平凡作用, 所以 $\Ker(\pi_n) = \{\pm I_2\}$.
所以回答这个问题: 表示 $V_n$ 忠实当且仅当 $n$ 奇.
\end{solution}

\begin{exercise}[$\SL_2(\R)$ 的有限维酉表示平凡]
\label{ex:I.5.5.2}
设 $(\pi, V)$ 是 $G := \SL_2(\R)$ 的有限维酉表示. 我们要证明 $V$ 平凡 (即对每个 $x \in G$ 有 $\pi(x) = \mathrm{id}_V$).
\begin{enumerate}
\item[(a)] 考虑群同态 $\alpha \colon \R \to G$, $t \longmapsto \begin{pmatrix} 1 & t \\ 0 & 1 \end{pmatrix}$. 证明存在 $V$ 的基 $\mathcal{B}$ 和 $y_1, \ldots, y_n \in \R$ (其中 $n = \dim V$) 使得对每个 $t \in \R$, 自同构 $\pi(\alpha(t))$ 在 $\mathcal{B}$ 中对角, 对角线元素为 $e^{ity_1}, \ldots, e^{ity_n}$.
\item[(b)] 证明对每个 $t \in \R$ 有 $\pi(\alpha(t)) = \mathrm{id}_V$. (提示: 若 $u \in \R^{\times}$, $x = \begin{pmatrix} u & 0 \\ 0 & u^{-1} \end{pmatrix}$, 考虑 $x\alpha(t)x^{-1}$ 在 $V$ 上的作用.)
\item[(c)] 证明对每个 $x \in G$ 有 $\pi(x) = \mathrm{id}_V$.
\item[(d)] 若 $n \geq 3$, 证明 $\SL_n(\R)$ 的每个有限维酉表示都平凡.
\end{enumerate}
\end{exercise}

\begin{solution}
(a). $\GL(V)$ 的子群 $\pi(\alpha(\R))$ 交换, 且它的所有元素可对角化 (因为它们都酉), 所以可以找到 $V$ 的基 $\mathcal{B} = (v_1, \ldots, v_n)$ 使得 $\pi(\alpha(\R))$ 的所有元素在其中对角 (若需要还可以是标准正交基). (若不喜欢同时对角化 $\GL(V)$ 的无穷子集, 只需选取生成 $\operatorname{Span}(\pi(\alpha(\R)))$ 的 $A_1, \ldots, A_m \in \pi(\alpha(\R))$ 并同时对角化它们.)
对每个 $j \in \{1, \ldots, n\}$, 子空间 $\CC v_j$ 被 $\alpha(\R) \subset G$ 的作用稳定 (由基的选取), 所以我们得到 $\R$ 在 $\CC v_j$ 上的 $1$ 维表示, 且由习题 \ref{ex:I.5.4.1}(b) 知道这种表示具有 $t \longmapsto e^{ity_j}v_j$ 的形式 (对某个 $y_j \in \CC$).

(b). 我们有 $x\alpha(t)x^{-1} = \alpha(u^2t)$, 所以 $\pi(x\alpha(t)x^{-1})$ 在基 $\mathcal{B}$ 中对角, 对角线元素为 $e^{iu^2ty_1}, \ldots, e^{iu^2ty_n}$. 另一方面,
\[
\operatorname{Tr}(\pi(x\alpha(t)x^{-1})) = \operatorname{Tr}(\pi(x)\pi(\alpha(t))\pi(x)^{-1}) = \operatorname{Tr}(\pi(\alpha(t))),
\]
所以对每个 $t \in \R$ 和每个 $u \in \R^{\times}$,
\[
\sum_{j=1}^n e^{ity_j} = \sum_{j=1}^n e^{iu^2ty_j}.
\]
假设我们知道 $L^{\infty}(\R)$ 的子集 $\widehat{\R}$ 线性无关. 则此等式告诉我们: 对每个 $u \in \R^{\times}$, 集合 $\{y_1, \ldots, y_n\}$ 和 $\{u^2y_1, \ldots, u^2y_n\}$ 相等. 这只有在 $y_1 = \cdots = y_n = 0$ 时才可能, 而这又蕴含 $\alpha(t)$ 对每个 $t \in \R$ 在 $V$ 上平凡作用.
现在证明关于 $\widehat{\R}$ 的结论. 设 $y_1, \ldots, y_m \in \R$ 两两不同, $c_1, \ldots, c_m \in \CC$ 满足对每个 $t \in \R$ 有 $\sum_{j=1}^m c_je^{ity_j} = 0$. 我们要证明 $c_1 = \cdots = c_m = 0$. 设 $r \in \R$. 取 $t = 0, r, \ldots, r(m-1)$, 利用 Vandermonde 行列式的计算, 看到必须有 $e^{iry_1} = \cdots = e^{iry_m}$. 由于这对每个 $r \in \R$ 成立, 它蕴含 $y_1 = \cdots = y_m$ (例如对前面的等式关于 $r$ 求导然后在 $r = 0$ 处取值). 所以 $m = 1$, 然后 $c_1e^{ity_1} = 0$ (对每个 $t \in \R$) 这一事实蕴含 $c_1 = 0$.

(c). 若 $x \in \SL_2(\R)$ 是剪切 (transvection) 矩阵 (即形如 $\begin{pmatrix} 1 & t \\ 0 & 1 \end{pmatrix}$ 的矩阵的共轭), 则 $x = y\alpha(t)y^{-1}$ (对某个 $t \in \R$ 和 $y \in \SL_2(\R)$), 所以由 (b), $\pi(x) = \pi(y)\pi(\alpha(t))\pi(y)^{-1} = \pi(y)\pi(y)^{-1} = \mathrm{id}_V$. 由于 $\SL_2(\R)$ 由剪切矩阵生成, 这蕴含对每个 $x \in \SL_2(\R)$ 有 $\pi(x) = \mathrm{id}_V$.

(d). 设 $\pi \colon \SL_n(\R) \to \GL(V)$ 是有限维酉表示. 设 $x \in \SL_n(\R)$ 是剪切矩阵. 我们可以模仿 (a) 和 (b) 来证明 $\pi(x) = \mathrm{id}_V$, 但也可以这样做: 选取 $\R^n$ 的基 $(v_1, \ldots, v_n)$ 使得对应于 $x$ 的 $\R^n$ 的线性自同态的矩阵为
\[
\begin{pmatrix} 1 & 1 & 0 & \cdots & 0 \\ 0 & 1 & 0 & \cdots & 0 \\ 0 & 0 & 1 & \cdots & 0 \\ \vdots & \vdots & \vdots & \ddots & \vdots \\ 0 & \cdots & 0 & \cdots & 1 \end{pmatrix}.
\]
考虑 $\SL_n(\R)$ 的由在基 $(v_1, \ldots, v_n)$ 中矩阵形如
\[
\begin{pmatrix} a & b & 0 & \cdots & 0 \\ c & d & 0 & \cdots & 0 \\ 0 & 0 & 1 & \cdots & 0 \\ \vdots & \vdots & \vdots & \ddots & \vdots \\ 0 & \cdots & 0 & \cdots & 1 \end{pmatrix}
\]
($a, b, c, d \in \R$, $ad - bc = 1$) 的元素组成的子集 $G'$. 则 $G'$ 是子群, 且与 $\SL_2(\R)$ 同构. 由于 $\pi|_{G'}$ 是 $G'$ 在 $V$ 上的酉表示, 由 (c) 有 $G' \subset \Ker(\pi)$. 特别地, $\pi(x) = \mathrm{id}_V$.
现在用 $\SL_n(\R)$ 由剪切矩阵生成的事实得出结论: 对每个 $x \in \SL_n(\R)$ 有 $\pi(x) = \mathrm{id}_V$.
\end{solution}

\begin{exercise}[酉化、小球与忠实表示]
\label{ex:I.5.5.3}
\begin{enumerate}
\item[(a)] 设 $G$ 是 $\GL_n(\CC)$ 的紧子群. 证明存在 $x \in \GL_n(\CC)$ 使得 $xGx^{-1} \subset \UU(n)$.
\item[(b)] 在 $M_n(\CC)$ 上放一个你喜欢的范数 (反正它们都等价). 证明存在 $c > 0$ 使得 $\GL_n(\CC)$ 包含于球 $\{x \in \GL_n(\CC) \mid \|x - I_n\| < c\}$ 的唯一子群是平凡群.
\item[(c)] 证明对 $\GL_n(\Q_p)$ 在有限维 $\CC$-向量空间上的每个连续表示, 存在整数 $m \geq 0$ 使得 $\GL_n(\Q_p)$ 的子群 $I_n + p^mM_n(\Z_p)$ 平凡作用.
\item[(d)] 证明若 $(\pi, V)$ 是 $\GL_n(\Z_p)$ 的不可约酉表示, 则存在 $m \geq 1$ 使得 $\pi(I_n + p^mM_n(\Z_p)) = \{1\}$.
\item[(e)] 更一般地, 证明若 $G$ 是射影有限 (profinite) 群 (即有限离散群的投射极限, 见习题 \ref{ex:I.5.1.3}), 则只有当 $G$ 有限时 $G$ 才有忠实的不可约酉表示.
\end{enumerate}
\end{exercise}

\begin{solution}
(a). 考虑 $G$ 在 $\CC^n$ 上由包含 $G \subset \GL_n(\CC)$ 给出的表示 $\rho$. 由定理 \ref{thm:I.3.2.8}, 存在 $\CC^n$ 上的 Hermitian 内积使得该表示酉. 设 $A$ 是此 Hermitian 内积在 $\CC^n$ 的典范基下的矩阵. 则 $A$ 是 Hermitian 正定矩阵, 所以可以把它写成 $A = B^*B$, 其中 $B \in \GL_n(\CC)$. (这是谱定理的简单推论. 由于 $A$ Hermitian, 有酉矩阵 $P$ 和对角矩阵 $D$ 使得 $A = P^*DP$. 由于 $A$ 正定, $D$ 的对角线元素是正实数, 所以可以写成 $D = C^2$, 其中 $C$ 是另一个对角线元素为正的对角矩阵. 取 $B = P^*CP$, 则 $B$ 是 Hermitian 正定且 $A = B^2 = B^*B$.)
$\rho$ 对 $A$ 酉意味着对每个 $X \in G$ 有 $X^*AX = A$. 由于 $A = B^*B$, 这等价于 $(BXB^{-1})^*(BXB^{-1}) = I_n$. 所以 $BGB^{-1} \subset \UU(n)$.

(b). 设 $G$ 是 $\GL_n(\CC)$ 的包含于形如 $\{x \in \GL_n(\CC) \mid \|x - I_n\| < c\}$ 的球的子群. 则闭子群 $\overline{G}$ 包含于闭球 $\{x \in \GL_n(\CC) \mid \|x - I_n\| \leq c\}$, 所以它紧, 从而由 (a) 包含于形如 $P\UU(n)P^{-1}$ 的子群. 特别地, $G$ 的每个元素可对角化且所有特征值的模为 $1$.
在 $\CC^n$ 上固定任何范数, 考虑 $M_n(\CC)$ 上对应的算子范数 $\|\cdot\|$. 我们将使用此范数. 注意若 $X \in M_n(\CC)$ 且 $\lambda$ 是 $X$ 的特征值, 则存在范数 $1$ 的向量 $v \in \CC^n$ 使得 $Xv = \lambda v$, 所以 $\|X\| \geq |\lambda|$. 现在让我们证明 $\GL_n(\CC)$ 的每个包含于开球 $B := \{x \in \GL_n(\CC) \mid \|x - I_n\| < \sqrt{2}\}$ 的子群都平凡. 设 $G$ 是这种子群, $X \in G$. 设 $\lambda_1, \ldots, \lambda_n$ 是 $X$ 的特征值. 刚才已见 $|\lambda_1| = \cdots = |\lambda_n| = 1$. 假设存在 $r \in \{1, \ldots, n\}$ 使得 $\lambda_r \neq 1$, 则可以把 $\lambda_r$ 写成 $\lambda_r = e^{i\theta}$, 其中 $-\pi/2 < \theta < \pi/2$, 因为 $|\lambda_r - 1| \leq \|X - I_n\| < \sqrt{2}$; 但此时, 若选取整数 $m \geq 1$ 使得 $\pi/2 \leq m|\theta| \leq \pi$, 将有 $\|X^m - I_n\| \geq |\lambda_r^m - 1| \geq \sqrt{2}$, 这与 $X^m \in G$ 矛盾. 所以必有 $\lambda_1 = \cdots = \lambda_n = 1$, 这意味着 $X = I_n$.

(c). 我们已在习题 \ref{ex:I.5.2.2}(d) 的解答中看到 $K_m := I_n + p^mM_n(\Z_p)$ 确实是 $\GL_n(\Q_p)$ 的子群. 我们还在 $M_n(\Q_p)$ 上放了范数使得 $K_m$ 是以 $I_n$ 为心、$p^{-m+1}$ 为半径的开球.
设 $\rho \colon \GL_n(\Q_p) \to \GL(V)$ 是 $\GL_n(\Q_p)$ 在有限维向量空间 $V$ 上的连续表示. 由命题 \ref{prop:I.3.5.1}, 同态 $\rho$ 连续. 设 $U$ 是 $\mathrm{id}_V$ 在 $\GL(V)$ 中的开邻域, 使得 $\GL(V)$ 包含于 $U$ 的唯一子群是 $\{1\}$ (由 (b) 这存在). 则 $\rho^{-1}(U)$ 是 $I_n$ 在 $\GL_n(\Q_p)$ 中的开邻域, 所以对 $m \gg 0$ 它包含 $K_m$. 但 $K_m$ 是 $\GL_n(\Q_p)$ 的子群, 所以 $\rho(K_m)$ 是 $\GL(V)$ 的子群, 故一旦 $\rho(K_m) \subset U$ 就有 $\rho(K_m) = \{1\}$.

(d). 由习题 \ref{ex:I.5.1.4}(l), 群 $\GL_n(\Z_p)$ 紧. 所以由习题 \ref{ex:I.5.5.9}, 空间 $V$ 有限维. 现在结论的证明与习题 \ref{ex:I.5.5.3}(c) 完全相同.

(e). 由习题 \ref{ex:I.5.1.3}, 我们知道 $G$ 紧 Hausdorff (注意有限离散群紧 Hausdorff). 所以由习题 \ref{ex:I.5.5.9}, $G$ 的每个不可约酉表示有限维.
假设我们知道 $G$ 完全不连通. 设 $(\pi, V)$ 是 $G$ 的连续有限维表示. 由习题 \ref{ex:I.5.2.2}(c), $G$ 的紧开子群构成 $1$ 的邻域基. 由习题 \ref{ex:I.5.5.3}(b), 可以找到 $\mathrm{id}_V$ 在 $\GL(V)$ 中的邻域 $U$ 使得 $\GL(V)$ 包含于 $U$ 的唯一子群是 $\{\mathrm{id}_V\}$. 所以若选取紧开子群 $K$ 使得 $\pi(K) \subset U$, 则必有 $K \subset \Ker(\pi)$. 所以 $\Ker(\pi) = \bigcup_{x \in \Ker(\pi)}xK$ 是 $G$ 的开子群, 从而群 $G/\Ker(\pi)$ 离散. 由于它也紧, 它是有限群. 这证明除非 $G$ 有限, $G$ 不能有忠实的不可约酉表示.
所以还需证明 $G$ 完全不连通, 即 $G$ 的唯一的非空连通子集是单点集. 取有限群的投射系统 $((G_i)_{i \in I}, (u_{ij} \colon G_i \to G_j)_{i \geq j})$ 使得 $G = \varprojlim_{i \in I}G_i$. 设 $C \subset G$ 是非空连通子集. 则 $C$ 在每个 $G_i$ 中的像是非空连通的, 从而是单点集 $\{g_i\}$. 这蕴含 $C$ 的唯一元素是族 $(g_i)_{i \in I} \in \prod_{i \in I}G_i$ (此族自动在投射极限中).
\end{solution}

\begin{remarkstar}
有限群并不总有忠实的不可约酉表示. 例如, 若 $G$ 是有限交换群, 则由 Schur 引理, $G$ 的每个不可约酉表示都是 $1$ 维的. 取 $G = \Z/2\Z \times \Z/2\Z$, 设 $\pi \colon G \to \CC$ 是 $1$ 维表示. 由于 $(1, 0)$ 和 $(0, 1)$ 在 $G$ 中阶为 $2$, 必须有 $\pi(1, 0), \pi(0, 1) \in \{\pm 1\}$. 若想要 $\pi$ 忠实, 需要 $\pi(1, 0) = \pi(0, 1) = -1$, 但此时 $\pi(1, 1) = (-1)^2 = 1$, 所以 $\pi$ 不能忠实. (更一般地, 有限交换群有忠实的 $1$ 维表示当且仅当它循环.)
\end{remarkstar}

\begin{exercise}[闭子群的 Lie 代数]
\label{ex:I.5.5.4}
本题的目标是定义 $\GL_n(\CC)$ 的闭子群 $G$ 的 Lie 代数, 并证明矩阵指数在它们恒等元附近诱导 Lie 代数与群之间的同胚.
考虑函数 $\exp \colon M_n(\CC) \to M_n(\CC)$, 定义为 $\exp(X) = \sum_{n \geq 0}\frac{1}{n!}X^n$; 也记 $e^X = \exp(X)$. 你可以假定此函数的如下基本性质: 定义它的级数绝对收敛, 它无穷次可微, 且可以逐项求导. 你也可以假定对任何满足 $AB = BA$ 的 $A, B \in M_n(\CC)$, 有 $\exp(A + B) = \exp(A)\exp(B)$; 特别地, 对每个 $X \in M_n(\CC)$ 有 $\exp(X) \in \GL_n(\CC)$ 且 $\exp(X)^{-1} = \exp(-X)$.
我们还固定 $\GL_n(\CC)$ 的闭子群 $G$.
\begin{enumerate}
\item[(a)] 计算 $\exp$ 在点 $0 \in M_n(\CC)$ 处的微分. (回忆这是从 $M_n(\CC)$ 到自身的线性算子.)
\item[(b)] 证明 $\exp$ 诱导从 $M_n(\CC)$ 中 $0$ 的邻域到 $\GL_n(\CC)$ 中 $1$ 的邻域的微分同胚.
\item[(c)] 设 $L = \{X \in M_n(\CC) \mid \forall t \in \R,\, \exp(tX) \in G\}$. 证明 $L$ 是 $M_n(\CC)$ 的 $\R$-线性子空间. (提示: 对所有 $X, Y \in M_n(\CC)$, 证明 $\exp(X + Y) = \lim_{k \to +\infty}\bigl(\exp(\frac{1}{k}X)\exp(\frac{1}{k}Y)\bigr)^k$.)\footnote{这个 $L$ 称为 $G$ 的 Lie 代数. 也容易证明它对换子括号 $[X, Y] = XY - YX$ 稳定.}
\item[(d)] 若 $G = \UU(n)$, 证明 $L = \{X \in M_n(\CC) \mid X^* = -X\}$.
\item[(e)] 若 $G = \SO(n)$, 证明 $L = \{X \in M_n(\CC) \mid X^T = -X\}$.
\item[(f)] 再次假设 $G$ 是 $\GL_n(\CC)$ 的任意闭子群. 本题的目标是证明如下结论 $(*)$: 存在 $L$ 中 $0$ 的邻域 $U$ 使得 $\exp(U)$ 是 $1$ 在 $G$ 中的邻域, 且 $\exp$ 诱导同胚 $U \xrightarrow{\sim} \exp(U)$.
  \begin{enumerate}
  \item[(i)] 设 $L'$ 是 $M_n(\CC)$ 的 $\R$-线性子空间使得 $M_n(\CC) = L \oplus L'$, 考虑函数 $\phi \colon M_n(\CC) \to \GL_n(\CC)$, 定义为 $\phi(A + B) = e^Ae^B$ (对每个 $A \in L$, $B \in L'$). 证明存在 $0$ 在 $L$ 中的邻域 $U_0$、$0$ 在 $L'$ 中的邻域 $V$ 和 $1$ 在 $\GL_n(\CC)$ 中的邻域 $W$ 使得 $\phi$ 诱导微分同胚 $U_0 \times V \xrightarrow{\sim} W$.
  \item[(ii)] 假设 $(*)$ 不成立. 证明存在 $L$ 中 $0$ 的递减邻域序列 $U_0 \supset U_1 \supset \cdots$, $L$ 的元素序列 $(A_k)_{k \geq 0}$ 和 $L'$ 的元素序列 $(B_k)_{k \geq 0}$ 使得:
  \begin{itemize}
  \item 对每个 $k \geq 0$, $A_k \in U_k$;
  \item 对每个 $k \geq 0$, $B_k \neq 0$;
  \item 对每个 $k \geq 0$, $\phi(A_k + B_k) \in G$;
  \item 序列 $(B_k)_{k \geq 0}$ 的极限是 $0$;
  \item 对 $L$ 中 $0$ 的每个邻域 $U$, 对足够大的 $k$ 有 $U_k \subset U$.
  \end{itemize}
  \item[(iii)] 证明序列 $(\frac{1}{\|B_k\|}B_k)_{k \geq 0}$ 有收敛子序列, 且该子序列的极限 $B$ 非零.
  \item[(iv)] 对每个 $t \in \R$, 证明 $\lfloor\frac{t}{\|B_k\|}\rfloor\|B_k\| \to t$ ($k \to +\infty$). (其中对每个 $c \in \R$, 记 $\lfloor c \rfloor$ 为不超过 $c$ 的最大整数.)
  \item[(v)] 证明 $B \in L$.
  \end{enumerate}
\item[(g)] 设 $(\rho, V)$ 是 $G$ 的连续有限维表示. 对每个 $X \in L$, 证明存在唯一的元素 $u(X) \in \End(V)$ 使得对每个 $t \in \R$ 有 $\rho(e^{tX}) = e^{tu(X)}$. 又证明函数 $u \colon L \to \End(V)$ 是 $\R$-线性的.\footnote{函数 $u \colon L \to \End(V)$ 称为 $\rho$ 在 $0$ 处的微分. 稍加努力可以证明它保持换子括号.}
\end{enumerate}
\end{exercise}

\begin{solution}
(a). 设 $d\exp_0$ 是 $\exp$ 在点 $0$ 处的微分. 由微分的定义, 对每个 $H \in M_n(\CC)$,
\[
d\exp_0(H) = \lim_{t \to 0}\frac{1}{t}(e^{tH} - e^0) = \lim_{t \to 0}\frac{1}{t}(e^{tH} - I_n) = \frac{d}{dt}e^{tH}\Big|_{t=0}.
\]
但
\[
\frac{d}{dt}e^{tH} = \sum_{n \geq 0}\frac{1}{n!}\frac{d}{dt}(tH)^n = H\exp(tH) = \exp(tH)H,
\]
所以 $d\exp_0(H) = H$. 最后得到 $d\exp_0 = \mathrm{id}_{M_n(\CC)}$.

(b). 由于 $\GL_n(\CC)$ 在 $M_n(\CC)$ 中开, $1$ 在 $\GL_n(\CC)$ 中的邻域与足够小的 $1$ 在 $M_n(\CC)$ 中的邻域相同. 所以结论由 $d\exp_0$ 可逆和反函数定理得出.

(c). 先证明提示. 设 $U$ 是 $0$ 在 $M_n(\CC)$ 中的邻域, $V$ 是 $1$ 在 $\GL_n(\CC)$ 中的邻域, 使得 $\exp$ 是从 $U$ 到 $V$ 的微分同胚. 记 $\log \colon V \to U$ 为它的逆. 由于当 $H \to 0$ 时 $\exp(H) = 1 + H + o(H)$, 当 $H \to 0$ 时有 $\log(1 + H) = H + o(H)$.
设 $X, Y \in M_n(\CC)$. 则 $\exp(\frac{1}{k}X) = 1 + \frac{1}{k}X + O(\frac{1}{k^2})$, $\exp(\frac{1}{k}Y) = 1 + \frac{1}{k}Y + O(\frac{1}{k^2})$, 所以 $\exp(\frac{1}{k}X)\exp(\frac{1}{k}Y) = 1 + \frac{1}{k}(X + Y) + O(\frac{1}{k^2})$. 若 $k$ 足够大, 则 $\exp(\frac{1}{k}X)\exp(\frac{1}{k}Y) \in V$, 且 $\log(\exp(\frac{1}{k}X)\exp(\frac{1}{k}Y)) = \frac{1}{k}(X + Y) + o(\frac{1}{k})$. 所以最终
\[
(\exp(\frac{1}{k}X)\exp(\frac{1}{k}Y))^k = \exp(k\log(\exp(\frac{1}{k}X)\exp(\frac{1}{k}Y))) = \exp(X + Y + o(1)) \to \exp(X + Y)
\]
当 $k \to +\infty$ 时.
集合 $L$ 由定义对纯量乘法稳定. 设 $X, Y \in L$. 则对每个 $t \in \R$ 有 $c(t) := e^{tX}e^{tY} \in G$. 由于 $G$ 在 $\GL_n(\CC)$ 中闭, 这蕴含对每个 $t \in \R$,
\[
\exp(t(X + Y)) = \lim_{k \to +\infty}(\exp(\frac{t}{k}X)\exp(\frac{t}{k}Y))^k \in G.
\]
所以 $X + Y \in L$.

(d). 设 $X \in M_n(\CC)$ 满足 $X = -X^*$. 则对每个 $t \in \R$ 有 $tX = -(tX)^*$ (特别地 $tX$ 与 $tX^*$ 交换), 所以
\[
e^{tX}(e^{tX})^* = e^{tX}e^{tX^*} = e^{tX + tX^*} = e^0 = I_n,
\]
即 $e^{tX} \in \UU(n)$. 所以 $X \in L$.
反之, 设 $X \in L$. 则对每个 $t \in \R$ 有 $e^{tX}e^{tX^*} = I_n$. 对此表达式求导 (并使用 (a) 的证明中的导数表达式) 给出
\[
0 = Xe^{tX}e^{tX^*} + e^{tX}e^{tX^*}X^* = X + X^*.
\]

(e). 这与 (d) 的证明完全相同, 只需把 ``$*$'' 换成 ``$T$''.

(f). (i) 由反函数定理, 只需证明 $\phi$ 在 $0 \in M_n(\CC)$ 处的微分可逆. 若 $A \in L$, $B \in L'$, 由微分的定义
\[
d\phi_0(A + B) = \lim_{t \to 0}\frac{1}{t}(\phi(tA + tB) - \phi(0)) = \lim_{t \to 0}\frac{1}{t}(e^{tA}e^{tB} - 1) = A + B
\]
(由 (a) 中的计算), 所以 $d\phi_0 = \mathrm{id}_{M_n(\CC)}$, 它当然可逆.

(ii) 选取 $0$ 在 $L$ (相应地 $L'$) 中的邻域序列 $U_0 \supset U_1 \supset \cdots$ (相应地 $V = V_0 \supset V_1 \supset \cdots$) 使得 $L$ (相应地 $L'$) 中 $0$ 的每个邻域 $U$ (相应地 $V'$) 对足够大的 $k$ 包含 $U_k$ (相应地 $V_k$). (例如可以取 $L$ 和 $L'$ 中半径趋于 $0$ 的球.) 对每个 $k \geq 0$, 函数 $\phi$ 是从 $U_k \times V_k$ 到 $\phi(U_k \times V_k)$ 的微分同胚, 特别地 $\phi(U_k \times V_k) \cap G$ 是 $1$ 在 $G$ 中的邻域, 包含 $\exp(U_k)$. 若 $(*)$ 不成立, 则对每个 $k$, $\phi(U_k \times V_k) \cap G$ 严格包含 $\exp(U_k)$, 所以可以找到 $A_k \in U_k$ 和 $B_k \in V_k$ 使得 $\phi(A_k + B_k) \in G$ 且 $\phi(A_k + B_k) \notin \exp(U_k)$, 即 $B_k \neq 0$. 又, 由对邻域 $V_k$ 的条件, 当 $k \to +\infty$ 时 $B_k \to 0$.

(iii) 序列 $(\frac{1}{\|B_k\|}B_k)_{k \geq 0}$ 是 $L'$ 的单位球的元素序列, 而单位球紧, 所以它有收敛子序列, 其极限仍在单位球中 (且特别地非零).

(iv) 对每个 $k \geq 0$,
\[
0 \leq \frac{t}{\|B_k\|} - \Bigl\lfloor\frac{t}{\|B_k\|}\Bigr\rfloor < 1,
\]
所以
\[
0 \leq t - \Bigl\lfloor\frac{t}{\|B_k\|}\Bigr\rfloor\|B_k\| < \|B_k\|.
\]
由于 $B_k \to 0$, 有 $\|B_k\| \to 0$, 这蕴含
\[
\Bigl\lfloor\frac{t}{\|B_k\|}\Bigr\rfloor\|B_k\| \to t.
\]

(v) 取子序列后, 可以假设 $B = \lim_{k \to +\infty}\frac{1}{\|B_k\|}B_k$. 我们必须证明对每个 $t \in \R$ 有 $e^{tB} \in G$. 设 $t \in \R$. 由 (iv) 和 $B$ 的定义,
\[
e^{tB} = \lim_{k \to +\infty}\exp\Bigl(\Bigl\lfloor\frac{t}{\|B_k\|}\Bigr\rfloor\|B_k\|\frac{1}{\|B_k\|}B_k\Bigr) = \exp\Bigl(\Bigl\lfloor\frac{t}{\|B_k\|}\Bigr\rfloor B_k\Bigr).
\]
但由 (ii), 对每个 $k \geq 0$ 有 $\phi(A_k + B_k) = e^{A_k}e^{B_k} \in G$ 且 $e^{A_k} \in G$ (因为 $A_k \in L$), 所以 $e^{B_k} \in G$; 由于 $N := \lfloor\frac{t}{\|B_k\|}\rfloor$ 是整数, 这蕴含 $e^{NB_k} = (e^{B_k})^N \in G$. 最后, 由于 $G$ 在 $\GL_n(\CC)$ 中闭, 我们推出 $e^{tB} \in G$.

(g). 设 $X \in L$. 考虑映射 $\R \to \GL(V)$, $t \longmapsto \rho(e^{tX})$. 这是连续群同态, 所以由习题 \ref{ex:I.5.4.1}(b)(i), 存在唯一的 $u(X) \in \End(V)$ 使得对每个 $t \in \R$ 有 $\rho(e^{tX}) = \exp(tu(X))$.
设 $X, Y \in L$, $a \in \R$. 对每个 $t \in \R$,
\[
e^{tu(aX)} = \rho(e^{taX}) = e^{tau(X)}.
\]
在 $t = 0$ 处求导, 得到 $u(aX) = au(X)$. 现在考虑 $c \colon \R \to \GL(V)$, $t \longmapsto \rho(e^{tX})\rho(e^{tY})\rho(e^{-t(X+Y)})$. 有 $c(t) = e^{tu(X)}e^{tu(Y)}e^{-tu(X+Y)}$, 所以 $c$ 是 $\mathcal{C}^{\infty}$ 的且 $c'(0) = u(X) + u(Y) - u(X+Y)$. 另一方面, 利用 $c$ 是 $\mathcal{C}^{\infty}$ 的这一事实, 可以像 (c) 那样证明对每个 $t \in \R$,
\[
\lim_{k \to +\infty}c(\frac{t}{k})^k = e^{tc'(0)}.
\]
所以我们只需证明这个极限对每个 $t \in \R$ 等于 $\mathrm{id}_V$. 用无穷小量做简单计算表明 (若 $t$ 固定)
\[
e^{\frac{t}{k}X}e^{\frac{t}{k}Y}e^{-\frac{t}{k}(X+Y)} = I_n + O(\frac{1}{k^2}),
\]
所以
\[
(e^{\frac{t}{k}X}e^{\frac{t}{k}Y}e^{-\frac{t}{k}(X+Y)})^k = I_n + O(\frac{1}{k}),
\]
且
\[
c(\frac{t}{k}) = \rho((e^{\frac{t}{k}X}e^{\frac{t}{k}Y}e^{-\frac{t}{k}(X+Y)})^k) \xrightarrow[k \to +\infty]{} \rho(I_n) = \mathrm{id}_V.
\]
\end{solution}

我们现在来看如何找完全没有忠实有限维表示 (在任何域上) 的离散群.
设 $\Gamma$ 是一个 (离散) 群. 我们说 $\Gamma$ 剩余有限 (residually finite), 如果对每个 $x \in \Gamma - \{1\}$, 存在 $\Gamma$ 的正规子群 $\Delta$ 使得 $\Gamma/\Delta$ 有限 (我们说 $\Delta$ 在 $\Gamma$ 中有限指标) 且 $x$ 在 $\Gamma/\Delta$ 中的像不平凡.\footnote{事实上, 可以用类似的想法证明: 若 $\operatorname{char}(k) = 0$, 这样的 $\Gamma$ 对几乎每个素数 $p$ 必须实质上剩余 $p$-有限 (virtually residually p-finite), 即它有一个有限指标子群 $\Gamma'$ 使得对每个 $x \in \Gamma' - \{1\}$, 存在 $\Gamma'$ 的有限指标正规子群 $\Delta \not\ni x$ 使得 $\Gamma'/\Delta$ 是 $p$-群; 但我所知道的唯一证明使用 Noether 正规化定理.}\footnote{加一个参考文献!}
下面两个习题的目标是证明: 若 $k$ 是域且 $\Gamma \subset \GL_n(k)$ 是有限生成子群, 则 $\Gamma$ 剩余有限.

\begin{exercise}[$\GL_n(R)$ 的剩余有限性]
\label{ex:I.5.5.5}
设 $R$ 是有限生成的 $\Z$-代数且是整环. 固定整数 $n \geq 1$. 对 $R$ 的每个理想 $I$, 令
\[
\Gamma(I) = \Ker(\GL_n(R) \to \GL_n(R/I)).
\]
\begin{enumerate}
\item[(a)] 证明 $R$ 是域当且仅当 $R$ 有限.
\item[(b)] 若 $\mathfrak{m}$ 是 $R$ 的极大理想, 证明 $\Gamma(\mathfrak{m})$ 是 $\GL_n(R)$ 的有限指标正规子群.
\item[(c)] 证明 $R$ 的所有极大理想的交是 $0$. (提示: 可以假设 $R$ 不是域. 若 $a \in R - \{0\}$, 证明局部化 $R[1/a]$ 不是域, 在 $R[1/a]$ 中取极大理想, 并与 $R$ 相交.)
\item[(d)] 证明 $\GL_n(R)$ 剩余有限.
\end{enumerate}
\end{exercise}

\begin{solution}
(a). 有限整环必是域, 这是经典事实. 证明如下. 假设 $R$ 有限, 设 $a \in R - \{0\}$. 则乘以 $a$ 是从 $R$ 到自身的加法映射, 其核为 $\{0\}$ (因为 $R$ 是整环), 所以它单射; 由于 $R$ 有限, 它满射, 这意味着存在 $b \in R$ 使得 $ab = 1$, 即 $a \in R^{\times}$.
其逆由交换代数的两个经典结果得出 (例如见 Eisenbud 的书 {[}9{]} 的习题 4.30 和 4.32):
\begin{itemize}
\item 若 $K \subset L$ 是域扩张且 $L$ 作为 $K$-代数有限生成, 则 $L$ 是有限维 $K$-向量空间 (Zariski 引理).
\item 若 $R$ 是 Noether 环, $S$ 是有限生成的 $R$-代数, $T \subset S$ 是 $R$-子代数且 $S$ 是有限 $T$-代数 (即作为 $T$-模有限生成), 则 $T$ 是有限生成的 $R$-代数 (Artin-Tate).
\end{itemize}
事实上, 若 $R$ 是域, 考虑它的素域 $k$. 则 $R$ 是有限生成的 $k$-代数, 所以由 Zariski 引理是有限维 $k$-向量空间, 这由第二个结果蕴含 $k$ 是有限生成的 $\Z$-代数. 注意 $k$ 要么是 $\Q$ 要么是某个有限域 $\mathbb{F}_p$. 但 $\Q$ 不是有限生成的 $\Z$-代数 (若 $x_1, \ldots, x_n \in \Q - \{0\}$, $P$ 是整除某个 $x_i$ 的分母的 (有限) 素数集合, 则由 $x_1, \ldots, x_n$ 生成的 $\Z$-子代数的非零元素的分母的素因子也必在 $P$ 中, 所以这个 $\Z$-子代数不能等于 $\Q$). 所以 $k$ 是有限域; 由于 $R$ 是有限维 $k$-向量空间, 它也是有限域.
为完整起见, 让我们给出我们实际需要的交换代数结果部分的证明. 假设我们知道:
$(*)$ 设 $L/K$ 是域扩张, 满足:
\begin{itemize}
\item 存在 $u \in L$ 使得 $L = K(u)$ (即 $L$ 作为域由 $u$ 生成);
\item $L$ 是有限生成的 $\Z$-代数,
\end{itemize}
则扩张有限且 $K$ 也是有限生成的 $\Z$-代数.
则可以用同样的方法证明: 若 $R$ 是域, 它必有限. (只需选取生成 $R$ (在其素域 $k$ 上) 的元素 $x_1, \ldots, x_n \in R$, 并对扩张 $k(x_1, \ldots, x_{i-1}) \subset k(x_1, \ldots, x_i)$ 应用 $(*)$ 证明 $k$ 是有限生成的 $\Z$-代数. 证明的结尾如前.)
我们现在证明 $(*)$. 设 $x_1, \ldots, x_n \in L^{\times}$ 生成 $L$ (作为 $\Z$-代数). 假设 $u$ 在 $K$ 上超越; 则 $L$ 同构于 $K$ 的有理分式域. 把 $x_i$ 写成 $x_i = \frac{P_i}{Q_i}$, 其中 $P_i, Q_i \in K[u]$. 由于 $(1 + u\prod_{i=1}^nQ_i)^{-1} = L = \Z[a_1, \ldots, a_n]$, 可以写成
\[
(1 + u\prod_{i=1}^nQ_i)^{-1} = \frac{R}{Q_1^{d_1}\cdots Q_n^{d_n}},
\]
其中 $R \in K[u]$ 与所有 $Q_i$ 互素, $d_1, \ldots, d_n \in \Z_{\geq 0}$. 得到 $Q_1^{d_1}\cdots Q_n^{d_n} = R(1 + u\prod_{i=1}^nQ_i)$, 这与 $R$ 与所有 $Q_i$ 互素矛盾. 所以 $u$ 在 $K$ 上代数. 设 $X^d + a_1X^{d-1} + \cdots + a_d \in K[X]$ 是 $u$ 在 $K$ 上的最小多项式. 对每个 $i \in \{1, \ldots, n\}$, 把 $x_i$ 写成 $x_i = \sum_{j=0}^{d-1}b_{ij}u^j$, 其中 $b_{ij} \in K$. 设 $A$ 是由 $a_1, \ldots, a_d$ 和所有 $b_{ij}$ 生成的 $K$ 的 $\Z$-子代数, 让我们证明 $A = K$. 设 $y \in K$. 则 $y$ 可以写成系数在 $\Z$ 中的 $x_1, \ldots, x_d$ 的多项式, 所以它也是系数在 $A$ 中的 $u$ 的多项式, 且可以取次数 $\leq d-1$ (可以用关系 $u^d = -a_1u^{d-1} - \cdots - a_d$ 把任何次数 $\geq d$ 的项换成更低次数的项). 换句话说, 可以把 $y$ 写成 $y = \sum_{i=0}^{d-1}c_iu^i$, 其中 $c_0, \ldots, c_{d-1} \in A$. 由于族 $(1, u, \ldots, u^{d-1})$ 在 $K$ 上线性无关, 必有 $c_1 = \cdots = c_{d-1} = 0$ 且 $y = c_0 \in A$.

(b). 首先, 群 $\Gamma(I)$ 是 $\GL_n(R)$ 的正规子群 (对任何理想 $I$), 因为它是群同态的核. 假设 $\mathfrak{m}$ 是极大理想. 则由 (a), $R/\mathfrak{m}$ 是有限域. 由于 $\GL_n(R)/\Gamma(\mathfrak{m})$ 单射进 $\GL_n(R/\mathfrak{m})$, 这蕴含 $\Gamma(\mathfrak{m})$ 在 $\GL_n(R)$ 中有限指标.

(c). 若 $R$ 是域, 则 $(0)$ 是 $R$ 的极大理想, 我们完成了. 假设 $R$ 不是域; 特别地, 由 (a), 它不有限. 设 $a \in R - \{0\}$. 局部化 $R[1/a] := R[X]/(aX - 1)$ 是有限生成的 $\Z$-代数 (因为 $R$ 是), 所以由 (a) 它只有有限时才是域. 但显然的映射 $R \to R[1/a]$ 单射 (因为 $a$ 不是零因子, 回忆 $R$ 是整环), 且 $R$ 无穷, 所以 $R[1/a]$ 也无穷, 故它不是域. 设 $\mathfrak{m}'$ 是 $R[1/a]$ 的极大理想, $\mathfrak{m}$ 是它在 $R$ 中的原像. 则映射 $R/\mathfrak{m} \to R[1/a]/\mathfrak{m}'$ 单射 (因为 $R \to R[1/a]$ 是), 且 $R[1/a]/\mathfrak{m}'$ 有限 (因为由 (a) 它是域), 所以 $R/\mathfrak{m}$ 有限且是整环, 因此是域 (再次由 (a)!), $\mathfrak{m}$ 是 $R$ 的极大理想. 又注意, 由于 $a$ 在 $R[1/a]$ 中可逆, 它不能在 $\mathfrak{m}'$ 中, 所以也不能在 $\mathfrak{m}$ 中. 所以我们找到了不包含 $a$ 的 $R$ 的极大理想.

(d). 设 $x = (x_{ij})_{1 \leq i,j \leq n} \in \GL_n(R)$ 满足 $x \neq I_n$. 选取 $i, j \in \{1, \ldots, n\}$ 使得当 $i \neq j$ 时 $x_{ij} \neq 0$, 当 $i = j$ 时 $x_{ij} \neq 1$. 由 (c), 可以找到 $R$ 的极大理想 $\mathfrak{m}$ 使得当 $i \neq j$ 时 $x_{ij} \notin \mathfrak{m}$, 当 $i = j$ 时 $x_{ij} - 1 \notin \mathfrak{m}$. 换句话说, $x$ 在 $\GL_n(R)/\Gamma(\mathfrak{m})$ 中的像不是单位元. 由于由 (b) $\Gamma(\mathfrak{m})$ 是 $\GL_n(R)$ 的有限指标正规子群, 我们完成了.
\end{solution}

\begin{exercise}[有限生成子群剩余有限]
\label{ex:I.5.5.6}
设 $k$ 是域, $\Gamma$ 是 $\GL_n(k)$ 的有限生成子群.
\begin{enumerate}
\item[(a)] 证明存在 $k$ 的有限生成 $\Z$-子代数 $R$ 使得 $\Gamma \subset \GL_n(R)$.
\item[(b)] 证明 $\Gamma$ 剩余有限.
\end{enumerate}
\end{exercise}

\begin{solution}
(a). 设 $\gamma_1, \ldots, \gamma_n$ 是 $\Gamma$ 的生成元, 设 $R$ 是由 $\gamma_i$ 及其逆的元素生成的 $k$ 的 $\Z$-子代数; 由定义这是有限生成的 $\Z$-代数. 由于 $\Gamma$ 的每个元素都是元素 $\gamma_i^{\pm 1}$ 的乘积, 有 $\Gamma \subset \GL_n(R)$.

(b). 这由习题 \ref{ex:I.5.5.5}(d) 立即得出: 若 $\gamma \in \Gamma - \{1\}$, 选取 $\GL_n(R)$ 的有限指标正规子群 $\Delta$ 使得 $\gamma$ 在 $\GL_n(R)/\Delta$ 中的像不平凡. 则 $\Gamma \cap \Delta$ 是 $\Gamma$ 的正规子群, 且 $\Gamma/(\Gamma \cap \Delta)$ 单射进 $\GL_n(R)/\Delta$, 所以 $\Gamma \cap \Delta$ 在 $\Gamma$ 中有限指标且 $\gamma$ 在 $\Gamma/(\Gamma \cap \Delta)$ 中的像不平凡.
\end{solution}

当然, 如果给不出任何有限生成的非剩余有限群的例子, 前一题的结果就不太有意思了. 我们在下面两个习题中给出这种例子.

\begin{exercise}[Baumslag-Solitar 群 $BS(2,3)$ 不剩余有限]
\label{ex:I.5.5.7}
设 $\Gamma$ 是自由群 (生成元 $a$ 和 $b$) 关于关系 $a^{-1}b^2a = b^3$ 的商. 在此题中, 我们假设 $b_1 := a^{-1}ba$ 与 $b$ 在 $\Gamma$ 中不交换, 并推出 $\Gamma$ 不剩余有限.
设 $u \colon \Gamma \to \Gamma'$ 是群同态, 其中 $\Gamma'$ 有限.
\begin{enumerate}
\item[(a)] 设 $n$ 是 $u(a)$ 在 $\Gamma'$ 中的阶. 证明 $u(b)$ 的阶整除 $3^n - 2^n$.
\item[(b)] 证明存在整数 $N \geq 0$ 使得 $u(b_1) = u(b_1)^{2^N}$. (注意 $u(b)$ 的阶与 $2$ 和 $3$ 都互素.)
\item[(c)] 证明 $u(b_1)$ 与 $u(b)$ 交换.
\item[(d)] 证明 $\Gamma$ 不剩余有限.
\end{enumerate}
\end{exercise}

\begin{solution}
(a). 我们首先证明对每个 $r \in \Z_{\geq 0}$ 有 $b^{2^r} = a^rb^{3^r}a^{-r}$. $r = 0$ 的情形显然, $r = 1$ 的情形就是定义 $\Gamma$ 的关系. 设 $r \geq 1$, 假设对 $r$ 已知结论, 证明对 $r+1$. 有
\[
b^{2^{r+1}} = (b^{2^r})^2 = (ab^{3^r}a^{-1})^2 = (ab^{2^r}a^{-1})^3 = (a^{r+1}b^{3^r}a^{-(r+1)})^3 = a^{r+1}b^{3^{r+1}}a^{-(r+1)}.
\]
应用于 $r = n$ 给出 $b^{2^n} = a^nb^{3^n}a^{-n}$, 所以 $u(b)^{3^n - 2^n} = 1$, 故 $u(b)$ 的阶整除 $3^n - 2^n$.

(b). 注意 $u(b)$ 的阶奇, 因为它整除奇数 $3^n - 2^n$. 所以存在 $N \geq 1$ 使得 $u(b)^{2^N} = u(b)$. 由于 $b_1 = a^{-1}ba$, 对每个 $r \geq 0$ 有 $b_1^r = a^{-1}b^ra$, 所以 $u(b_1)^{2^N} = u(b_1)$, 如所愿.

(c). 由定义 $\Gamma$ 的关系有 $b_1^2 = b^3$, 所以由 (b) 有 $u(b_1) = u(b)^{3^N}$. 这蕴含 $u(b)$ 与 $u(b_1)$ 交换.

(d). 设 $c = b_1^{-1}bb_1b^{-1} \in \Gamma$. 则我们假设 $c \neq 1$, 但 (c) 证明对 $\Gamma$ 的每个有限指标正规子群 $\Delta$, $c$ 在 $\Gamma/\Delta$ 中的像平凡. 所以 $\Gamma$ 不剩余有限.
\end{solution}

\begin{exercise}[$b_1$ 与 $b$ 不交换]
\label{ex:I.5.5.8}
设 $\Gamma$ 是自由群 (生成元 $a$ 和 $b$) 关于关系 $a^{-1}b^2a = b^3$ 的商. 本题的目标是证明 $b_1 := a^{-1}ba$ 与 $b$ 在 $\Gamma$ 中不交换, 即 $b_1bb_1^{-1}b^{-1}$ 在 $\Gamma$ 中不平凡.\footnote{证明它的最简单方法是找 $\Gamma$ 在其上 $b_1bb_1^{-1}b$ 非平凡作用的有限维表示, 但我们做不到. 不过, 这种想法的某种变体会奏效.}
设 $F$ 是生成元 $a$ 和 $b$ 上的自由群. 回忆 $F$ 的元素是字母 $a, a^{-1}, b, b^{-1}$ 上的既约字 (reduced words). (既约字是不含冗余对 $aa^{-1}$, $a^{-1}a$, $bb^{-1}$ 或 $b^{-1}b$ 的字.) 我们把 $F$ 的元素写成 $a^{n_1}b^{m_1}\cdots a^{n_r}b^{m_r}$, 其中 $n_1, m_1, \ldots, n_r, m_r \in \Z$, 且 $m_1, n_2, m_2, \ldots, n_{r-1}, m_{r-1}, n_r \neq 0$.
设 $\Omega$ 是如下形式的既约字的集合 $b^{r_1}a^{s_1}\cdots b^{r_m}a^{s_m}b^r$, 其中:
\begin{enumerate}
\item[(i)] $m \in \Z_{\geq 0}$, $r_i, s_i, r \in \Z$;
\item[(ii)] 对每个 $i \in \{1, \ldots, m\}$ 有 $s_i \neq 0$;
\item[(iii)] 对每个 $i \in \{2, \ldots, m\}$ 有 $r_i \neq 0$;
\item[(iv)] 对每个 $i \in \{1, \ldots, m\}$, 若 $s_i > 0$ 则 $0 \leq r_i \leq 1$;
\item[(v)] 对每个 $i \in \{1, \ldots, m\}$, 若 $s_i < 0$ 则 $0 \leq r_i \leq 2$.
\end{enumerate}
由 $\Gamma$ 的定义, 有满射群同态 $F \to \Gamma$, 记为 $\phi$.
\begin{enumerate}
\item[(a)] 证明 $\phi(\Omega) = \Gamma$.
\item[(b)] 对每个 $w \in \Omega$ 和每个 $s \in \{a, a^{-1}, b, b^{-1}\}$, 找到 $w' \in \Omega$ 使得 $\phi(w') = \phi(ws)$. 在下面把此 $w'$ 记为 $w \cdot s$.
\item[(c)] 对每个 $w \in \Omega$ 和每个 $s \in \{a, a^{-1}, b, b^{-1}\}$, 证明 $(w \cdot s) \cdot s^{-1} = w$.
\item[(d)] 证明 $(w, s) \longmapsto w \cdot s$ 延拓成 $\Gamma$ 在 $\Omega$ 上的右作用.
\item[(e)] 证明 $\phi$ 诱导双射 $\Omega \xrightarrow{\sim} \Gamma$.
\item[(f)] 证明 $b_1bb_1^{-1}b^{-1} \neq 1$ 在 $\Gamma$ 中成立.
\end{enumerate}
\end{exercise}

\begin{solution}
(a). 由自由群的定义, 可以把 $F$ 的每个元素 $w$ 写成满足条件 (i), (ii) 和 (iii) 的既约字 $b^{r_1}a^{s_1}\cdots b^{r_m}a^{s_m}b^r$. 定义 $N(w)$ 为所有满足 $r_i \notin \{0, 1\}$ 的 $s_i > 0$ 的最大值; 所以若 $w$ 满足条件 (iv), 则 $N(w) = 0$. 定义 $M(w)$ 为所有满足 $r_i \notin \{0, 1, 2\}$ 的 $s_i < 0$ 的 $|s_i|$ 的最大值; 所以若 $w$ 满足条件 (v), 则 $M(w) = 0$. 我们通过对 $N(w) + M(w)$ 归纳证明存在 $w' \in \Omega$ 使得 $\phi(w) = \phi(w')$. 若 $N(w) + M(w) = 0$, 则 $w$ 满足条件 (iv) 和 (v), 所以它在 $\Omega$ 中且结论显然.
假设 $N(w) + M(w) > 0$. 若 $N(w) > 0$, 选取满足 $s_i > 0$ 且 $r_i \notin \{0, 1\}$ 的 $i \in \{1, \ldots, m\}$. 注意定义 $\Gamma$ 的关系说 $\phi(b^2a) = \phi(ab^3)$, 从而也有 $\phi(b^{-2}a) = \phi(ab^{-3})$, 这蕴含对每个 $k \in \Z$ 有 $\phi(b^{2k}a) = \phi(ab^{3k})$. 把 $r_i$ 写成 $r_i = 2k + l$, 其中 $k \in \Z$, $l \in \{0, 1\}$, 令
\[
w' = b^{r_1}a^{s_1}\cdots b^{r_{i-1}}a^{s_{i-1}}b^la b^{3k}a^{s_i-1}b^{r_{i+1}}a^{s_{i+1}}\cdots b^{r_m}a^{s_m}b^r.
\]
则由上面的观察, $\phi(w) = \phi(w')$, 且 $N(w') < N(w)$, $M(w') = M(w)$.
类似地, 若 $M(w) > 0$, 选取满足 $s_i < 0$ 且 $r_i \notin \{0, 1, 2\}$ 的 $i \in \{1, \ldots, m\}$. 对 $k \in \Z$, 等式 $\phi(b^{2k}a) = \phi(ab^{3k})$ 也可以写成 $\phi(a^{-1}b^{2k}) = \phi(b^{3k}a^{-1})$. 把 $r_i$ 写成 $r_i = 3k + l$, 其中 $k \in \Z$, $l \in \{0, 1, 2\}$, 令
\[
w' = b^{r_1}a^{s_1}\cdots b^{r_{i-1}}a^{s_{i-1}}b^la^{-1}b^{2k}a^{s_i+1}b^{r_{i+1}}a^{s_{i+1}}\cdots b^{r_m}a^{s_m}b^r.
\]
则由上面的观察, $\phi(w) = \phi(w')$, 且 $N(w') = N(w)$, $M(w') < M(w)$. 由于 $N(w)$ 和 $M(w)$ 中必有一个 $> 0$, 总可以找到 $w' \in F$ 使得 $\phi(w') = \phi(w)$ 且 $N(w') + M(w') < N(w) + M(w)$. 对 $w'$ 应用归纳假设得到结果.

(b). 设 $w = b^{r_1}a^{s_1}\cdots b^{r_m}a^{s_m}b^r \in \Omega$; 假设条件 (i)-(v) 满足. 若 $s = b$ (相应地 $s = b^{-1}$), 则 $w' = b^{r_1}a^{s_1}\cdots b^{r_m}a^{s_m}b^{r+1}$ (相应地 $w' = b^{r_1}a^{s_1}\cdots b^{r_m}a^{s_m}b^{r-1}$) 可行. 若 $s = a$, 把 $r$ 写成 $r = 2k + l$, 其中 $k \in \Z$, $l \in \{0, 1\}$, 取 $w' = b^{r_1}a^{s_1}\cdots b^{r_m}a^{s_m}b^lab^{3k}$. 若 $s = a^{-1}$, 把 $r$ 写成 $r = 3k + l$, 其中 $k \in \Z$, $l \in \{0, 1, 2\}$, 取 $w' = b^{r_1}a^{s_1}\cdots b^{r_m}a^{s_m}b^la^{-1}b^{2k}$.

(c). 若 $s \in \{b, b^{-1}\}$ 结论显然. 假设 $s = a$, 把 $r$ 写成 $r = 2k + l$, 其中 $k \in \Z$, $l \in \{0, 1\}$. 则 $w \cdot a = b^{r_1}a^{s_1}\cdots b^{r_m}a^{s_m}b^lab^{3k}$, 所以
\[
(w \cdot a) \cdot a^{-1} = b^{r_1}a^{s_1}\cdots b^{r_m}a^{s_m}b^laa^{-1}b^{2k} = w.
\]
若 $s = a^{-1}$, 把 $r$ 写成 $r = 3k + l$, 其中 $k \in \Z$, $l \in \{0, 1, 2\}$. 则 $w \cdot a = b^{r_1}a^{s_1}\cdots b^{r_m}a^{s_m}b^la^{-1}b^{2k}$, 所以
\[
(w \cdot a^{-1}) \cdot a = b^{r_1}a^{s_1}\cdots b^{r_m}a^{s_m}b^la^{-1}ab^{3k} = w.
\]

(d). 由 (c), $(w, s) \longmapsto w \cdot s$ 延拓成 $F$ 在 $\Omega$ 上的右作用. 为证明它穿过 $\Gamma$ 在 $\Omega$ 上的右作用, 只需证明 $b^{-3}a^{-1}b^2a$ 平凡作用. 设 $w = b^{r_1}a^{s_1}\cdots b^{r_m}a^{s_m}b^r \in \Omega$; 假设条件 (i)-(v) 满足. 把 $r$ 写成 $r = 3k + l$, 其中 $k \in \Z$, $l \in \{0, 1, 2\}$. 则
\[
w \cdot b^{-3} = b^{r_1}a^{s_1}\cdots b^{r_m}a^{s_m}b^{r-3} = b^{r_1}a^{s_1}\cdots b^{r_m}a^{s_m}b^{l+3(k-1)},
\]
所以
\[
w \cdot (b^{-3}a^{-1}) = b^{r_1}a^{s_1}\cdots b^{r_m}a^{s_m}b^la^{-1}b^{2(k-1)},
\]
从而
\[
w \cdot (b^{-3}a^{-1}b^2) = b^{r_1}a^{s_1}\cdots b^{r_m}a^{s_m}b^la^{-1}b^{2k},
\]
最后
\[
w \cdot (b^{-3}a^{-1}b^2a) = b^{r_1}a^{s_1}\cdots b^{r_m}a^{s_m}b^la^{-1}ab^{3k} = b^{r_1}a^{s_1}\cdots b^{r_m}a^{s_m}b^{l+3k} = w.
\]

(e). 我们已由 (a) 知道 $\phi(\Omega) = \Gamma$, 所以只需证明 $\phi|_{\Omega}$ 单射. 由 (b) 的证明中给出的作用的显式公式, 若 $w \in \Omega$, 则 $1 \cdot w = w$. 由于由 (d), $1 \cdot w$ 只依赖于 $\phi(w)$, 这证明 $\phi(w)$ 确定 $w$.

(f). 由 (e), 只需证明 $\phi(b_1bb_1^{-1}b^{-1})$ 在 $\Omega$ 中的唯一原像不平凡. 我们已在 (a) 的证明中看到把既约字变换成 $\phi$ 下同像的 $\Omega$ 元素的算法. 把它应用于 $b_1bb_1^{-1}b^{-1} = a^{-1}baba^{-1}b^{-1}ab^{-1}$, 得到 $a^{-1}baba^{-1}bab^{-4} \neq 1$ (除了一些容易犯的错误), 所以我们完成了.
\end{solution}

\begin{exercise}[紧群的不可约酉表示有限维]
\label{ex:I.5.5.9}
本题的目标是证明定理 \ref{thm:I.3.2.13}, 即紧群的每个不可约酉表示都有限维这一事实.
设 $G$ 是紧群, $dx$ 是 $G$ 上的规范化 Haar 测度, $(\pi, V)$ 是 $G$ 的非零酉表示. 固定 $u \in V - \{0\}$, 定义 $T \colon V \to V$ 为
\[
T(v) = \int_G \langle v, \pi(x)(u) \rangle\pi(x)(u)\,\dd x.
\]
\begin{enumerate}
\item[(a)] 证明 $T$ 良定义且 $T \in \End(V)$.
\item[(b)] 证明 $T$ 是 $G$-等变的.
\item[(c)] 证明对每个 $v \in V$ 有 $\langle T(v), v \rangle \geq 0$.
\item[(d)] 证明 $T \neq 0$.
\item[(e)] 证明 $T$ 在 $\{T' \in \End(V) \mid \dim_{\CC}(\Im(T')) < +\infty\}$ 的 (对 $\|\cdot\|_{\mathrm{op}}$ 的) 闭包中; 换句话说, $T$ 在有限秩自同态空间的闭包中. (提示: $G \to V$, $x \longmapsto \pi(x)(u)$ 一致连续.)
\item[(f)] 设 $B$ 是 $V$ 中的闭单位球. 证明 $T(B)$ 紧. (换句话说, 算子 $T$ 是紧算子 (compact operator). 习题 \ref{ex:I.5.6.5} 可以帮你缩短证明.)
\item[(g)] 若 $V$ 是 $G$ 的不可约表示, 证明 $V$ 有限维.
\end{enumerate}
\end{exercise}

\begin{solution}
(a). 我们必须证明定义 $T(v)$ 的积分对每个 $v \in V$ 收敛. 设 $v \in V$. 则函数 $G \to V$, $\langle v, \pi(x)(u) \rangle\pi(x)(u)$ 连续 (因为 $x \longmapsto \pi(x)(u)$ 连续); 由于 $G$ 紧, 由习题 \ref{ex:I.5.6.3} 积分存在, 而且
\[
\|T(v)\| \leq \int_G |\langle v, \pi(x)(u) \rangle|\|\pi(x)(u)\|\,\dd x \leq \|v\|\|u\|^2.
\]
函数 $T \colon V \to V$ 是 $\CC$-线性的 (因为加法和乘以纯量在 $V$ 上连续, 所以它们与积分交换, 由习题 \ref{ex:I.5.6.1}(b)), 且上面的不等式证明 $T$ 有界且 $\|T\|_{\mathrm{op}} \leq \|u\|^2$.

(b). 设 $v \in V$, $x \in G$.
\begin{align*}
T(\pi(x)(v)) &= \int_G \langle \pi(x)(v), \pi(y)(u) \rangle\pi(y)(u)\,\dd y \\
&= \int_G \langle v, \pi(x)^*\pi(y)(u) \rangle\pi(y)(u)\,\dd y \\
&= \int_G \langle v, \pi(x^{-1}y)(u) \rangle\pi(y)(u)\,\dd y \\
&= \int_G \langle v, \pi(y)(u) \rangle\pi(x)\pi(y)(u)\,\dd y
\end{align*}
(最后一个等式用 Haar 测度的左不变性). 由于 $\pi(x) \colon V \to V$ 连续且线性, 习题 \ref{ex:I.5.6.1}(b) 蕴含最后一行等于
\[
\pi(x)\Bigl(\int_G \langle v, \pi(y)(u) \rangle\pi(y)(u)\,\dd y\Bigr) = \pi(x)(T(v)),
\]
这正是我们想要的.

(c). 设 $v \in V$. 由于 $\langle v, \cdot \rangle$ 在 $V$ 上连续且线性,
\[
\langle T(v), v \rangle = \int_G \overline{\langle \pi(x)(u), v \rangle}\langle v, \pi(x)(u) \rangle\,\dd x = \int_G |\langle v, \pi(x)(u) \rangle|^2\,\dd x \geq 0.
\]

(d). 取 $v = u$. 由于 $\langle u, \pi(x)(u) \rangle = \|u\|^2 > 0$ 且 $x \longmapsto \langle u, \pi(x)(u) \rangle$ 是从 $G$ 到 $\CC$ 的连续函数, 存在 $\varepsilon > 0$ 和 $1$ 在 $G$ 中的开邻域 $U$ 使得对 $x \in U$ 有 $|\langle u, \pi(x)(u) \rangle|^2 \geq \varepsilon$. 则 (c) 的证明中的计算给出
\[
\langle T(u), u \rangle = \int_G |\langle u, \pi(x)(u) \rangle|^2\,\dd x \geq \varepsilon\mu(U) > 0.
\]
所以 $T \neq 0$.

(e). 设 $\varepsilon > 0$. 由于 $G$ 紧, 连续函数 $G \to \CC$, $x \longmapsto \pi(x)(u)$ 一致连续, 所以存在 $1$ 的邻域 $U$ 使得对 $x \in G$ 和 $y \in xU$ 有 $\|\pi(x)(u) - \pi(y)(u)\| \leq \varepsilon$. 由于 $G$ 紧且族 $(xU)_{x \in G}$ 覆盖 $G$, 可以找到 $x_1, \ldots, x_n \in G$ 使得 $G = \bigcup_{i=1}^n x_iU$. 选取 $X$ 的 Borel 子集 $E_1, \ldots, E_n$ 使得对每个 $i \in \{1, \ldots, n\}$ 有 $x_i \in E_i \subset x_iU$ 且 $X = E_1 \sqcup \cdots \sqcup E_n$ (作为集合). 若 $x \in E_i$, $v \in V$, 则
\begin{align*}
&\|\langle v, \pi(x)(u) \rangle\pi(x)(u) - \langle v, \pi(x_i)(u) \rangle\pi(x_i)(u)\| \\
&\leq \|\langle v, (\pi(x) - \pi(x_i))(u) \rangle\pi(x)(u)\| + \|\langle v, \pi(x_i)(u) \rangle(\pi(x) - \pi(x_i))(u)\| \\
&\leq \|v\|\varepsilon\|u\| + \|v\|\|u\|\varepsilon = 2\varepsilon\|v\|\|u\|.
\end{align*}
定义 $U \in \End(V)$ 为
\[
T(v) = \sum_{i=1}^n\mu(E_i)\langle v, \pi(x_i)(u) \rangle\pi(x_i)(u) = \sum_{i=1}^n\int_{E_i}\langle v, \pi(x_i)(u) \rangle\pi(x_i)(u)\,\dd x.
\]
这个算子 $U$ 有有限秩, 因为它的像包含于 $\operatorname{Span}(\pi(x_1)(u), \ldots, \pi(x_n)(u))$. 又, 由上面的计算 (以及习题 \ref{ex:I.5.6.2}), 对每个 $v \in V$,
\begin{align*}
\|T(v) - U(v)\| &\leq \sum_{i=1}^n\Bigl\|\int_{E_i}\langle v, \pi(x)(u) \rangle\pi(x)(u)\,\dd x - \int_{E_i}\langle v, \pi(x_i)(u) \rangle\pi(x_i)(u)\,\dd x\Bigr\| \\
&\leq \sum_{i=1}^n\mu(E_i)2\varepsilon\|v\|\|u\| \\
&= 2\varepsilon\|v\|\|u\|.
\end{align*}
所以 $\|T - U\|_{\mathrm{op}} \leq 2\varepsilon\|u\|$. 由于 $\varepsilon > 0$ 任意, 这证明 $T$ 是有限秩算子的极限.

(f). 由习题 \ref{ex:I.5.6.5}(e), 只需证明 $T(B)$ 完全有界. 设 $U$ 是 $0$ 的邻域, 可以假设它是半径 $\varepsilon > 0$ 的开球. 我们必须找到 $x_1, \ldots, x_n \in B$ 使得 $T(B)$ 的每个点与某个 $T(x_i)$ 的距离 $< \varepsilon$.
由 (e), 我们知道 $T$ 是有限秩算子的极限, 所以可以找到有限秩的 $U \in \End(V)$ 使得 $\|T - U\|_{\mathrm{op}} \leq \varepsilon/4$. 由于 $U$ 有有限秩, $U(B)$ 是有限维空间 $\Im(U)$ 的闭有界子集, 所以它紧. 特别地, 可以找到 $x_1, \ldots, x_n \in B$ 使得对每个 $y \in B$, 存在 $i \in \{1, \ldots, n\}$ 满足 $\|U(y) - U(x_i)\| < \varepsilon/2$.
现在设 $y \in B$, 选取满足 $\|U(y) - U(x_i)\| < \varepsilon/2$ 的 $i \in \{1, \ldots, n\}$. 则
\[
\|T(y) - T(x_i)\| \leq \|T(y) - U(y)\| + \|U(y) - U(x_i)\| + \|U(x_i) - T(x_i)\| < \|y\|\varepsilon/4 + \varepsilon/2 + \|x_i\|\varepsilon/4 \leq \varepsilon.
\]
(回忆 $y, x_i$ 在 $V$ 的闭单位球中.)

(g). 现在把一切放在一起. 假设 $V$ 是 $G$ 的不可约酉表示. 则我们构造的算子 $T \in \End(V)$ 是 $G$-等变的, 所以由 Schur 引理存在 $\lambda \in \CC$ 使得 $T = \lambda\mathrm{id}_V$. 由于 $T \neq 0$, 有 $\lambda \neq 0$. 所以 $T(\lambda^{-1}B)$ 是 $V$ 中的闭单位球. 部分 (f) 说它紧, 由 Riesz 引理这蕴含 $V$ 有限维.
\end{solution}

\subsection{向量值积分与 Minkowski 不等式}

注意: 允许不使用证明就用下列结果:
\begin{itemize}
\item Hahn-Banach 定理.
\item 希尔伯特空间 $V$ 上的每个连续线性泛函都具有 $\langle \cdot, v \rangle$ 的形式 (其中 $v \in V$) 这一事实.
\item H\"older 不等式.
\item 若 $(X, \mu)$ 是测度空间, $1 \leq p < +\infty$, $1 < q \leq +\infty$ 满足 $p^{-1} + q^{-1} = 1$, 则映射 $L^p(X, \mu) \to \Hom(L^q(X, \mu), \CC)$, $f \longmapsto (g \longmapsto \int_X fg\,\dd\mu)$ 是保范同构 (左边是 $L^p$ 范数, 右边是算子范数) 这一事实.\footnote{技术注记: 对 $p = 1$, $q = +\infty$, 若 $\mu$ 不是 $\sigma$-有限的, 这在一般情形不成立; 但对局部紧 Hausdorff 空间上的正则 Borel 测度, 可以通过稍微修改 $L^{\infty}$ 的定义来补救. 你可以忽略这一点.}
\end{itemize}

\begin{exercise}[弱可积函数与积分]
\label{ex:I.5.6.1}
设 $(X, \mu)$ 是测度空间, $V$ 是巴拿赫空间. 记 $V^{\vee} = \Hom(V, \CC)$. 若对每个 $T \in V^{\vee}$, 函数 $T \circ f \colon X \to \CC$ 都在 $L^1(X, \mu)$ 中, 则称函数 $f \colon X \to V$ 弱可积 (weakly integrable). 若 $f$ 弱可积且存在 $V$ 的元素 $v$ 使得对每个 $T \in V^{\vee}$ 有 $T(v) = \int_X T \circ f(x)\,\dd\mu(x)$, 则称 $v$ 是 $f$ 在 $X$ 上的积分 (integral), 记为 $v = \int_X f(x)\,\dd\mu(x) = \int_X f\,\dd\mu$.
\begin{enumerate}
\item[(a)] 证明 $f$ 的积分若存在则唯一.
\item[(b)] 设 $W$ 是另一个巴拿赫空间, $u \in \Hom(V, W)$. 若 $f \colon X \to V$ 弱可积且有积分 $v$, 证明 $u \circ f \colon X \to W$ 弱可积且有积分, 且积分等于 $u(v)$.
\item[(c)] 给出一个弱可积但没有积分的函数的例子.
\end{enumerate}
\end{exercise}

\begin{solution}
(a). 由 Hahn-Banach 定理, 对每个 $v \in V$, 存在 $T \in \Hom(V, \CC)$ 使得 $T(v) = \|v\|$ 且 $\|T\|_{\mathrm{op}} \leq 1$. 特别地, $V$ 的元素 $v$ 为零当且仅当对每个 $T \in \Hom(V, \CC)$ 有 $T(v) = 0$; 换句话说, 两个元素 $v, w \in V$ 相等当且仅当对每个 $T \in \Hom(V, \CC)$ 有 $T(v) = T(w)$. 这蕴含 $f$ 的积分若存在则唯一.

(b). 我们首先证明 $u \circ f$ 弱可积. 设 $T \in \Hom(W, \CC)$. 则 $T \circ u \in \Hom(V, \CC)$, 所以函数 $T \circ u \circ f \colon X \to \CC$ 可积.
现在假设 $f$ 有积分 $v$. 则对每个 $T \in \Hom(W, \CC)$, 有 $T \circ u \in \Hom(V, \CC)$, 所以 $\int_X T \circ u \circ f\,\dd\mu = T \circ u(v)$. 这意味着 $u(v)$ 是 $u \circ f$ 的积分.

(c). 取 $X = \N$ 带计数测度 $\mu$,
\[
V = c_0(\N) := \{(x_n)_{n \geq 0} \in \CC^{\N} \mid \lim_{n \to +\infty}x_n = 0\}.
\]
我们将用 $\ell^1(\N)$ 是 $c_0(\N)$ 的连续对偶这一事实, 对偶配对为 $\ell^1(\N) \times c_0(\N) \to \CC$, $((x_n), (y_n)) \longmapsto \sum_{n \geq 0}x_ny_n$, 且 $\ell^1(\N)$ 的连续对偶是 $\ell^{\infty}(\N)$ (通过类似的映射). 从 $c_0(\N)$ 到其双重对偶的映射是通常的嵌入 $c_0(\N) \subset \ell^{\infty}(\N)$.
我们定义 $f \colon X \to V$ 为 $f(n) = \mathbb{1}_{\{n\}}$. 则对每个 $(x_n)_{n \geq 0} \in \ell^1(\N)$, 若 $T \colon c_0(\N) \to \CC$ 是对应的线性泛函, 则
\[
\int_X T(f(x))\,\dd\mu(x) = \sum_{n \geq 0}x_n,
\]
它收敛因为 $(x_n)_{n \geq 0} \in \ell^1(\N)$. 所以 $f$ 弱可积. 但 $f$ 没有积分 (至少在 $c_0(\N)$ 中没有), 因为它在 $\ell^1(\N)$ 上定义的连续线性泛函由 $\ell^{\infty}(\N)$ 的某个不在 $c_0(\N)$ 中的元素 (常值序列 $1$) 表示. 由于在 $\ell^1(\N)$ 的点上取值分离 $\ell^{\infty}(\N)$ 的元素, 不可能有 $c_0(\N)$ 的元素在 $\ell^1(\N)$ 上给出相同的线性泛函.
\end{solution}

\begin{exercise}[紧支撑连续函数的积分]
\label{ex:I.5.6.2}
在此题中, $X$ 是局部紧 Hausdorff 空间, $\mu$ 是 $X$ 上的正则 Borel 测度. 设 $V$ 是巴拿赫空间, $f \colon X \to V$ 是具有紧支撑的连续函数.
\begin{enumerate}
\item[(a)] 证明 $f$ 弱可积.
\item[(b)] 若 $\mu(\operatorname{supp} f) = 0$, 证明 $\int_X f\,\dd\mu$ 存在且等于 $0$.
\end{enumerate}
本题的目标是证明:
\begin{enumerate}
\item[(i)] $f$ 有积分 $v$;
\item[(ii)] $\|v\| \leq \int_X \|f(x)\|\,\dd\mu(x)$;
\item[(iii)] 若 $\mu(\operatorname{supp} f) \neq 0$, 则 $\mu(\operatorname{supp} f)^{-1}v$ 在 $f(X)$ 的凸包的闭包中.
\end{enumerate}
由 (b), 可以 (而且将) 假设 $\mu(\operatorname{supp} f) \neq 0$.
\begin{enumerate}
\item[(a)] 证明可以假设 $X = \operatorname{supp} f$ (特别地 $X$ 紧) 且 $\mu(X) = 1$.
\end{enumerate}
从现在起, 我们假设 $X$ 紧且 $\mu(X) = 1$.
\begin{enumerate}
\item[(a)] 设 $T_1, \ldots, T_n \colon V \to \R$ 是有界 $\R$-线性泛函 (以显然的方式把 $V$ 看作 $\R$-向量空间), 定义 $a_1, \ldots, a_n \in \R$ 为 $a_i = \int_X T_i \circ f\,\dd\mu$. 证明 $(a_1, \ldots, a_n)$ 在 $\R^n$ 的紧子集 $((T_1, \ldots, T_n) \circ f)(X)$ 的凸包中. (提示: 如果不是会怎样?)
\end{enumerate}
设 $K$ 是 $f(X)$ 的凸包的闭包. 由习题 \ref{ex:I.5.6.7}, 这是 $V$ 的紧子集. 对 $\Hom_{\R}(V, \R)$ (从 $V$ 到 $\R$ 的有界 $\R$-线性泛函空间) 的每个有限子集 $\Omega$, 记 $I_{\Omega}$ 为 $K$ 的满足对每个 $T \in \Omega$ 有 $T(v) = \int_X T \circ f\,\dd\mu$ 的元素 $v$ 的集合.
\begin{enumerate}
\item[(a)] 证明对每个 $\Omega \subset \Hom_{\R}(V, \R)$, $I_{\Omega}$ 紧.
\item[(b)] 证明若 $\Omega$ 有限, 则 $I_{\Omega}$ 非空.
\item[(c)] 证明 $f$ 的积分存在且在 $K$ 中.
\item[(d)] 证明:
\[
\Bigl\|\int_X f\,\dd\mu\Bigr\| \leq \int_X \|f(x)\|\,\dd\mu(x).
\]
(提示: Hahn-Banach.)
\end{enumerate}
\end{exercise}

\begin{solution}
(a). 设 $T \in \Hom(V, \CC)$. 则 $T \circ f \colon X \to \CC$ 连续, 且其支撑包含于 $\operatorname{supp}(f)$, 从而是紧的. 所以 $T \circ f$ 可积.

(b). 假设 $\mu(\operatorname{supp} f) = 0$. 若 $T \in \Hom(V, \CC)$, 则 $T \circ f \colon X \to \CC$ 连续且 $\mu(\operatorname{supp}(T \circ f)) = 0$, 所以 $\int_X T \circ f\,\dd\mu = 0$. 这证明 $0$ 是 $f$ 的积分.

(c). 假设当 $X = \operatorname{supp} f$ 且 $\mu(X) = 1$ 时知道结论. 设 $f \colon X \to V$ 是具有紧支撑的连续函数. 我们已见可以假设 $\mu(\operatorname{supp} f) \neq 0$, 就这样做. 设 $X' = \operatorname{supp} f$, 考虑 $X'$ 上的测度 $\mu'$, 它是 $\mu$ 的限制乘以 $\mu(\operatorname{supp} f)^{-1}$. 由我们的假设, $\int_{X'}f|_{X'}\,\dd\mu'$ 存在, 记之为 $v$, 有 $\|v\| \leq \int_{X'}\|f(x)\|\,\dd\mu'(x)$, 且 $v$ 在 $f(X')$ 的凸包的闭包中.
让我们证明 $w := \mu(\operatorname{supp} f)v$ 是 $f$ 的积分. 注意 $\mu(\operatorname{supp} f)^{-1}w$ 在 $f(X') = f(X)$ 的凸包中, 且
\[
\|w\| \leq \mu(\operatorname{supp} f)\int_{X'}|f(x')|\,\dd\mu'(x') = \int_X |f(x)|\,\dd\mu(x),
\]
所以这证明了对 $f$ 的结论.
设 $T \in \Hom(V, \CC)$. 则 $\operatorname{supp}(T \circ f) \subset X'$, 所以
\[
\int_X T \circ f(x)\,\dd\mu(x) = \mu(\operatorname{supp} f)\int_{X'}T \circ f(x)\,\dd\mu'(x) = T(w).
\]
所以 $w = \int_X f\,\dd\mu$.

(d). 设 $L = ((T_1, \ldots, T_n) \circ f)(X)$. 假设 $(a_1, \ldots, a_n)$ 不在 $L$ 的凸包中. 则由超平面分离定理, 存在线性泛函 $\lambda \colon \R^n \to \R$ 和 $c > 0$ 使得对每个 $v \in L$ 有 $\lambda(a_1, \ldots, a_n) \geq c + \lambda(v)$. 设 $\lambda_1, \ldots, \lambda_n \in \R$ 是 $\lambda$ 在 $\R^n$ 的典范基向量下的像. 则对每个 $x \in X$,
\[
\sum_{i=1}^n\lambda_ia_i = \lambda(a_1, \ldots, a_n) \geq c + \lambda \circ (T_1, \ldots, T_n) \circ f(x) = c + \sum_{i=1}^n\lambda_iT_i(f(x)).
\]
对 $X$ 取积分 (利用 $\mu(X) = 1$) 给出
\[
\sum_{i=1}^n\lambda_ia_i \geq c + \sum_{i=1}^n\lambda_i\int_X T_i \circ f(x)\,\dd\mu(x) = c + \sum_{i=1}^n\lambda_ia_i,
\]
矛盾.

(e). 对每个 $T \in \Hom_{\R}(V, \R)$, $K$ 的满足 $T(v) = \int_X T \circ f\,\dd\mu$ 的元素 $v$ 的集合是 $K$ 的闭子集. 由于集合 $I_{\Omega}$ 是这种形式的集合的交, 它也是 $K$ 的闭子集, 从而因 $K$ 紧而紧.

(f). 若 $\Omega$ 有限, 写 $\Omega = \{T_1, \ldots, T_n\}$, $T_{\Omega} = (T_1, \ldots, T_n) \colon V \to \R^n$. 我们已在 (d) 中看到 $a := \int_X T_{\Omega}(f(x))\,\dd\mu(x)$ 在 $T_{\Omega}(f(X))$ 的凸包中, 所以存在 $v \in V$ 使得 $v$ 在 $f(X)$ 的凸包中 (从而在 $K$ 中) 且 $T_{\Omega}(v) = a = \int_X T_{\Omega}(f(x))\,\dd\mu(x)$. 第二个条件正好说 $v \in I_{\Omega}$.

(g). $K$ 的子集族 $(I_{\{T\}})_{T \in \Hom_{\R}(V, \R)}$ 由 (f) 有有限交性质. 由于 $K$ 紧, 这蕴含 $\bigcap_{T \in \Hom_{\R}(V, \R)}I_{\{T\}}$ 非空. 在其中选取向量 $v$. 设 $T \in \Hom(V, \CC)$. 由于 $\Re(T)$ 和 $\Im(T)$ 在 $\Hom_{\R}(V, \R)$ 中,
\begin{align*}
T(v) &= \Re(T(v)) + i\Im(T(v)) \\
&= \int_X \Re(T(f(x)))\,\dd\mu(x) + i\int_X \Im(T(f(x)))\,\dd\mu(x) \\
&= \int_X T(f(x))\,\dd\mu(x),
\end{align*}
所以 $v = \int_X f\,\dd\mu$. 又, $v \in K$, 因为由定义所有 $I_{\{T\}}$ 都包含于 $K$.

(h). 由 Hahn-Banach 定理, 存在 $T \in \Hom(V, \CC)$ 使得 $T(v) = \|v\|$ 且 $\|T\|_{\mathrm{op}} \leq 1$. 则
\[
\|v\| = |T(v)| = \Bigl|\int_X T(f(x))\,\dd\mu(x)\Bigr| \leq \int_X |T(f(x))|\,\dd\mu(x) \leq \int_X \|f(x)\|\,\dd\mu(x).
\]
\end{solution}

\begin{exercise}[$L^1$ 函数乘有界连续函数的积分]
\label{ex:I.5.6.3}
在此题中, $X$ 是局部紧 Hausdorff 空间, $\mu$ 是 $X$ 上的正则 Borel 测度. 设 $V$ 是巴拿赫空间, $f \colon X \to \CC$ 是 $L^1(X, \mu)$ 中的函数, $G \colon X \to V$ 是有界连续函数.
本题的目标是证明:
\begin{enumerate}
\item[(i)] 函数 $fG \colon X \to V$ 有积分 $v$;
\item[(ii)] $\|v\| \leq \bigl(\sup_{x \in X}\|G(x)\|\bigr)\int_X |f(x)|\,\dd\mu(x)$;
\item[(iii)] $v \in \overline{\operatorname{Span}(G(X))}$.
\end{enumerate}
\begin{enumerate}
\item[(a)] 证明 $fG$ 弱可积.
\item[(b)] 设 $(f_n)_{n \geq 0}$ 是 $\mathcal{C}_c(X)$ 的函数序列, 在 $L^1(X, \mu)$ 中收敛到 $f$. 证明对每个 $n \geq 0$, $\int_X f_nG\,\dd\mu$ 存在, 且 $(\int_X f_nG\,\dd\mu)_{n \geq 0}$ 是 Cauchy 序列.
\item[(c)] 证明上面的断言 (i), (ii) 和 (iii).
\end{enumerate}
\end{exercise}

\begin{solution}
(a). 设 $T \in \Hom(V, \CC)$. 则对每个 $x \in X$,
\[
|T(f(x)G(x))| \leq |f(x)||T(G(x))| \leq |f(x)|\|T\|_{\mathrm{op}}\|G(x)\| \leq |f(x)|\|T\|_{\mathrm{op}}\sup_{y \in X}\|G(y)\|.
\]
由于 $\sup_{y \in X}\|G(y)\| < +\infty$ 且 $f \in L^1(X, \mu)$, 函数 $T \circ (fG)$ 可积. 所以 $fG$ 弱可积.

(b). 对每个 $h \in \mathcal{C}_c(X)$, 函数 $hG \colon X \to V$ 连续且支撑包含于 $\operatorname{supp}(h)$, 从而是紧的. 由习题 \ref{ex:I.5.6.2}, 此函数可积, 且
\[
\Bigl\|\int_X (hG)(x)\,\dd\mu(x)\Bigr\| \leq \int_X |h(x)|\|G(x)\|\,\dd\mu(x) \leq \sup_{y \in X}\|G(y)\|\|h\|_1.
\]
把它应用于 $f_n$ 证明 $f_nG$ 可积, 把它应用于 $f_n - f_m$ 证明
\[
\Bigl\|\int_X (f_nG)\,\dd\mu - \int_X (f_mG)\,\dd\mu\Bigr\| \leq \sup_{y \in X}\|G(y)\|\|f_n - f_m\|_1 \xrightarrow[n,m \to +\infty]{} 0
\]
因为 $(f_n)_{n \geq 0}$ 在 $L^1(X, \mu)$ 中收敛.

(c). 由于 $V$ 完备, Cauchy 序列 $(\int_X (f_nG)\,\dd\mu)$ 在 $V$ 中有极限, 记之为 $v$. 对每个 $T \in \Hom(V, \CC)$,
\[
T(v) = \lim_{n \to +\infty}T\Bigl(\int_X (f_nG)\,\dd\mu\Bigr) = \lim_{n \to +\infty}\int_X T(f_n(x)G(x))\,\dd\mu(x).
\]
像 (a) 那样,
\[
\Bigl\|\int_X T(f_n(x)G(x))\,\dd\mu(x) - \int_X T(f(x)G(x))\,\dd\mu(x)\Bigr\| \leq \sup_{y \in X}\|G(y)\|\|T\|_{\mathrm{op}}\|f_n - f\|_1,
\]
所以当 $n \to +\infty$ 时这收敛到 $0$, 于是
\[
T(v) = \int_X T(f(x)G(x))\,\dd\mu(x).
\]
这证明 $v$ 是 $fG$ 的积分.
而且, 由习题 \ref{ex:I.5.6.2}, 对每个 $n \geq 0$, $\int_X (f_nG)\,\dd\mu$ 在 $G(X)$ 的线性张成的闭包中. 由于 $v$ 是这些向量的极限, 它也在 $\operatorname{Span}(G(X))$ 的闭包中.
最后, 为证明对 $\|v\|$ 的界, 可以用 Hahn-Banach 定理和 (h) 中刻画 $v$ 的性质, 或用如下事实: 对每个 $n \geq 0$,
\[
\Bigl\|\int_X (f_nG)\,\dd\mu\Bigr\| \leq \int_X |f_n(x)|\|G(x)\|\,\dd\mu(x) \leq \sup_{x \in X}\|G(x)\|\|f_n\|_1
\]
且这个积分序列收敛到 $v$.
\end{solution}

\begin{exercise}[卷积的积分公式]
\label{ex:I.5.6.4}
设 $G$ 是局部紧群, $\mu$ 是 $G$ 上的左 Haar 测度, $L^1(G) = L^1(G, \mu)$. 设 $f, g \in L^1(G)$.
\begin{enumerate}
\item[(a)] 证明函数 $G \to L^1(G)$, $y \longmapsto f(y)L_yg$ 弱可积且有积分.
\item[(b)] 证明
\[
f * g = \int_G f(y)L_yg\,\dd\mu(y).
\]
\end{enumerate}
\end{exercise}

\begin{solution}
(a). 注意函数 $G \to L^1(G)$, $y \longmapsto L_yg$ 连续, 且 $\sup_{y \in G}\|L_yg\|_1 = \|g\|_1 < +\infty$. 所以结论由习题 \ref{ex:I.5.6.3} 得出.

(b). 设 $F = \int_G f(y)L_yg\,\dd\mu(y) \in L^1(G)$. 由积分的定义, 对每个 $h \in L^{\infty}(G)$,
\[
\int_G h(x)F(x)\,\dd\mu(x) = \int_{G \times G}h(x)f(y)g(y^{-1}x)\,\dd\mu(y)\,\dd\mu(x) = \int_G h(x)(f * g)(x)\,\dd\mu(x).
\]
由于 $L^{\infty}(G)$ 是 $L^1(G)$ 的连续对偶, 由习题 \ref{ex:I.5.6.1}(a) 有 $f * g = F$.
\end{solution}

\begin{exercise}[完全有界与凸包]
\label{ex:I.5.6.5}
设 $V$ 是赋范向量空间. 若对 $V$ 中 $0$ 的每个邻域 $U$, 存在有限集 $F$ 使得 $A \subset F + U$, 则称 $V$ 的子集 $A$ 完全有界 (totally bounded).
\begin{enumerate}
\item[(a)] 证明 $V$ 的有限子集的凸包紧.
\item[(b)] 证明 $V$ 的每个紧子集完全有界.
\item[(c)] 若 $A \subset V$ 完全有界, 证明 $\overline{A}$ 完全有界.
\item[(d)] 若 $A$ 是 $V$ 的完全有界子集, 证明其凸包完全有界. (提示: 开球是凸的.)
\item[(e)] 若 $V$ 完备且 $A$ 完全有界, 证明 $\overline{A}$ 紧.
\item[(f)] 若 $V$ 完备且 $K \subset V$ 紧, 证明 $K$ 的凸包的闭包紧.
\end{enumerate}
\end{exercise}

\begin{solution}
(a). 设 $F$ 是 $V$ 的有限子集. 则其凸包包含于 $\operatorname{Span}(F)$, 它是有限维的. 在有限维向量空间中, 任何紧集的凸包紧, 所以有限集 $F$ 的凸包紧.

(b). 设 $K$ 是 $V$ 的紧子集, $U$ 是 $V$ 中 $0$ 的邻域. 可以假设 $U$ 开. 则 $K \subset \bigcup_{x \in K}(x + U)$. 由于 $K$ 紧, 存在 $K$ 的有限子集 $F$ 使得 $K \subset \bigcup_{x \in F}(x + U) = F + U$.

(c). 设 $A \subset V$ 完全有界, $U$ 是 $V$ 中 $0$ 的邻域. 可以假设 $U$ 是以 $0$ 为心、正半径 (设为 $c$) 的开球. 设 $U'$ 是以 $0$ 为心、半径 $c/2$ 的开球. 由于 $A$ 完全有界, 存在有限集 $F$ 使得 $A \subset F + U'$. 由于 $F$ 有限, 集合 $F + \overline{U'}$ 闭, 所以它包含 $\overline{A}$. 但 $U \supset U'$, 所以 $\overline{A} \subset F + U$.

(d). 设 $U$ 是 $V$ 中 $0$ 的邻域. 选取凸开邻域 $U'$ (例如球) 使得 $U' + U' \subset U$, 设 $F$ 是满足 $A \subset F + U'$ 的有限集. 设 $K$ 是 $F$ 的凸包, 则 $A \subset K + U'$. 由于 $K$ 和 $U'$ 凸, $K + U'$ 也凸, 所以 $A$ 的凸包包含于 $K + U'$. 另一方面, 集合 $K$ 由 (a) 紧, 从而由 (b) 完全有界, 所以存在有限集 $F'$ 使得 $K \subset F' + U'$, 从而 $K + U' \subset F' + U' + U' \subset F' + U$. 所以我们找到了有限集 $F'$ 使得 $A$ 的凸包包含于 $F' + U$.

(e). 记 $K = \overline{A}$. 对每个 $x \in V$ 和 $c > 0$, 设 $\overline{B}(x, c)$ 是以 $x$ 为心、$c$ 为半径的闭球.
设 $(U_i)_{i \in I}$ 是 $K$ 的开子集族, 满足 $K \subset \bigcup_{i \in I}U_i$, 并假设 $(U_i)_{i \in I}$ 没有覆盖 $K$ 的有限子族. 由 (c) 我们知道 $K$ 完全有界. 我们将对 $n$ 归纳构造 $K$ 的非空闭子集的递减序列 $(K_n)_{n \geq 1}$, 使得 $K_n$ 包含于半径为 $1/n$ 的球中且 $K_n$ 不能被 $(U_i)_{i \in I}$ 的有限子族覆盖.
首先, 由于 $K$ 完全有界, 存在有限集 $F$ 使得 $K \subset F + \overline{B}(0, 1) = \bigcup_{x \in F}\overline{B}(x, 1)$. 选取 $x \in F$ 使得 $K \cap \overline{B}(x, 1)$ 非空且不能被有限多个 $U_i$ 覆盖, 取 $K_1 = K \cap \overline{B}(x, 1)$.
现在假设已经构造了 $K_1, \ldots, K_n$ ($n \geq 1$). 则由于 $K_n$ 完全有界 (作为 $K$ 的子集), 存在有限集 $F$ 使得 $K_n \subset F + \overline{B}(0, (n+1)^{-1})$. 同样地, 由于 $K_n$ 不能被有限多个 $U_i$ 覆盖, 必存在 $x \in F$ 使得 $K_n \cap \overline{B}(x, (n+1)^{-1})$ 非空且也不能被有限多个 $U_i$ 覆盖, 取 $K_{n+1} = K_n \cap \overline{B}(x, (n+1)^{-1})$.
对每个 $n \geq 1$ 选取 $x_n \in K_n$. 由 $K_n$ 包含于半径为 $1/n$ 的球中这一条件, 序列 $(x_n)_{n \geq 0}$ 是 Cauchy 列. 由于 $V$ 完备, $(x_n)_{n \geq 1}$ 有极限, 设为 $x$. 由于 $x \in K$, 存在 $i \in I$ 使得 $x \in U_i$. 但此时对足够小的 $c > 0$ 有 $\overline{B}(x, c) \subset U_i$, 所以对足够大的 $n$ 有 $K_n \subset U_i$, 这与 $K_n$ 的性质矛盾.

(f). 由 (d), $K$ 的凸包完全有界, 所以其闭包由 (e) 紧.
\end{solution}

\begin{exercise}[$L^{\infty}$ 上的卷积]
\label{ex:I.5.6.6}
设 $G$ 是局部紧群, $dx$ 是 $G$ 上的左 Haar 测度, $\mathrm{UCB}(G)$ 是 $L^{\infty}(G)$ 的由 $G$ 上的左一致连续有界函数组成的子空间.
设 $f \in L^1(G)$, $\phi \in L^{\infty}(G)$.
\begin{enumerate}
\item[(a)] 证明 $f * \phi$ 存在且左一致连续有界.
\item[(b)] 若 $\phi \in \mathrm{UCB}(G)$, 证明积分 $\int_G f(y)L_y\phi\,\dd y$ 存在且等于 $f * \phi$.
\end{enumerate}
\end{exercise}

\begin{solution}
(a). 设 $x \in G$. 则定义 $f * \phi(x)$ 的积分是
\[
\int_G f(y)\phi(y^{-1}x)\,\dd\mu(y),
\]
它收敛, 因为对每个 $y \in G$ 有 $|f(y)\phi(y^{-1}x)| \leq \|\phi\|_{\infty}|f(y)|$. 这也证明对每个 $x \in G$ 有 $|f * \phi(x)| \leq \|\phi\|_{\infty}\|f\|_1$, 所以 $f * \phi$ 有界且
\[
\|f * \phi\|_{\infty} \leq \|\phi\|_{\infty}\|f\|_1.
\]
现在我们证明 $f * \phi$ 左一致连续. 设 $x \in G$. 由命题 \ref{prop:I.4.1.3}, $L_x(f * \phi) = (L_xf) * \phi$, 所以
\[
\|L_x(f * \phi) - f * \phi\|_{\infty} = \|(L_xf - f) * \phi\|_{\infty} \leq \|L_xf - f\|_1\|\phi\|_{\infty}.
\]
由命题 \ref{prop:I.3.1.13}, 当 $x$ 趋于 $1$ 时这趋于 $0$, 这正好意味着 $f * \phi$ 左一致连续.

(b). 假设 $\phi \in \mathrm{UCB}(G)$. 则映射 $G \to \mathrm{UCB}(G)$, $y \longmapsto L_y\phi$ 连续 (见注记 \ref{rem:I.1.13}), 所以由习题 \ref{ex:I.5.6.3}, 积分 $\int_G f(y)L_y\phi\,\dd\mu(y)$ 在 $\mathrm{UCB}(G)$ 中存在. 设 $h = \int_G f(y)L_y\phi\,\dd\mu(y)$.
对每个 $g \in L^1(G)$, 映射 $\psi \longmapsto \int_G g\psi\,\dd\mu$ 是 $\mathrm{UCB}(G)$ 上的连续线性泛函. 所以由积分的定义,
\[
\int_G gh\,\dd\mu = \int_{G \times G}h(x)f(y)\phi(y^{-1}x)\,\dd\mu(x)\,\dd\mu(y) = \int_G g(f * \phi)\,\dd\mu.
\]
由于 $L^1(G)$ 的元素定义的线性泛函在 $L^{\infty}(G)$ 上分离点, 这蕴含 $h = f * \phi$.
\end{solution}

\begin{exercise}[Minkowski 不等式]
\label{ex:I.5.6.7}
设 $(X, \mu)$ 和 $(Y, \nu)$ 是测度空间, 为简化我们取其为 $\sigma$-有限的.\footnote{有办法把结果推广到不一定是 $\sigma$-紧的带 Haar 测度的局部紧群.} 设 $p \in (1, +\infty)$.\footnote{Minkowski 不等式对 $p = 1$ 仍成立, 但该情形由 Fubini-Tonelli 定理立即得出.} 设 $\phi \colon X \times Y \to \CC$ 是可测函数. 假设
\[
\int_Y\Bigl(\int_X|\phi(x, y)|^p\,\dd\mu(x)\Bigr)^{1/p}\,\dd\nu(y) < \infty.
\]
\begin{enumerate}
\item[(a)] 证明函数 $\phi(\cdot, y)$ 对几乎所有 $y \in Y$ 在 $L^p(X, \mu)$ 中.
\item[(b)] 设 $Y'$ 是 $Y$ 的可测子集, 满足 $\nu(Y - Y') = 0$ 且对每个 $y \in Y'$ 有 $\phi(\cdot, y) \in L^p(X, \mu)$. 证明函数 $Y' \to L^p(X, \mu)$, $y \longmapsto \phi(\cdot, y)$ 弱可积.
\item[(c)] 证明 (b) 中函数的积分 $h \in L^p(X, \mu)$ 存在, 且对几乎所有 $x \in X$ 有
\[
h(x) = \int_Y \phi(x, y)\,\dd\nu(y).
\]
\item[(d)] 证明 Minkowski 不等式:
\[
\Bigl(\int_X\Bigl|\int_Y\phi(x, y)\,\dd\nu(y)\Bigr|^p\,\dd\mu(x)\Bigr)^{1/p} \leq \int_Y\Bigl(\int_X|\phi(x, y)|^p\,\dd\mu(x)\Bigr)^{1/p}\,\dd\nu(y).
\]
\end{enumerate}
\end{exercise}

\begin{solution}
(a). 函数 $y \longmapsto \bigl(\int_X|\phi(x, y)|^p\,\dd\mu(x)\bigr)^{1/p}$ 由假设可积, 所以它对几乎所有 $y \in Y$ 取有限值, 这意味着对几乎所有 $y \in Y$ 有 $\int_X|\phi(x, y)|^p\,\dd\mu(x) < +\infty$.

(b). 设 $q \in (1, +\infty)$ 满足 $\frac{1}{p} + \frac{1}{q} = 1$. 我们要验证对每个 $f \in L^q(X, \mu)$, 积分 $\int_{Y'}\int_X f\phi\,\dd\mu\,\dd\nu$ 收敛 (其中 $\phi = \phi(x, y)$). 设 $f \in L^q(X, \mu)$. 由 H\"older 不等式, 对每个 $y \in Y'$, $\int_X f(x)\phi(x, y)\,\dd\mu(x)$ 绝对收敛, 且
\[
\int_X |f(x)\phi(x, y)|\,\dd\mu(x) \leq \|f\|_q\|\phi(\cdot, y)\|_p.
\]
由于 $\int_{Y'}\|\phi(\cdot, y)\|_p\,\dd\nu(y)$ 由假设收敛, 这给出 $\int_{Y'}\int_X f(x)\phi(x, y)\,\dd\mu(x)\,\dd\nu(y)$ 的收敛性, 甚至是绝对收敛性.

(c). 我们已在 (b) 中看到对每个 $f \in L^q(X, \mu)$ 有 $\int_X\int_Y |f(x)\phi(x, y)|\,\dd\mu(x)\,\dd\nu(y) < +\infty$. 由 Fubini 定理, 这蕴含对每个 $f \in L^q(X, \mu)$, 对几乎所有 $x \in X$ 有 $\int_Y |f(x)\phi(x, y)|\,\dd\nu(y) = |f(x)|\int_Y|\phi(x, y)|\,\dd\nu(y) < +\infty$. 由于 $f$ 任意 (且 $\mu$ $\sigma$-有限), 我们得到对几乎所有 $x \in X$ 有 $\int_Y|\phi(x, y)|\,\dd\nu(y) < +\infty$, 设对 $x \in X'$ (其中 $\mu(X - X') = 0$).
我们定义函数 $h \colon X' \to \CC$ 为 $h(x) = \int_Y\phi(x, y)\,\dd\nu(y)$. 我们要证明这是 $y \longmapsto \phi(\cdot, y)$ 的积分. 若 $f \in L^q(X, \mu)$, 则
\[
\int_X f(x)h(x)\,\dd\mu(x) = \int_X f(x)\phi(x, y)\,\dd\mu(x)\,\dd\nu(y),
\]
所以, 用 (b) 中的 H\"older 不等式,
\[
\Bigl|\int_X f(x)h(x)\,\dd\mu(x)\Bigr| \leq \int_{X \times Y}|f(x)\phi(x, y)|\,\dd\mu(x) \leq C\|f\|_q,
\]
其中
\[
C = \int_Y\Bigl(\int_X|\phi(x, y)|^p\,\dd\mu(x)\Bigr)^{1/p}\,\dd\nu(y).
\]
这证明 $f \longmapsto \int_X f(x)h(x)\,\dd\mu(x)$ 是 $L^q(X, \mu)$ 上的有界线性泛函, 且其算子范数被 $C$ 控制. 由于 $L^q(X, \mu)$ 的连续对偶是 $L^p(X, \mu)$, 必有 $h \in L^p(X, \mu)$ 且 $\|h\|_p \leq C$. 第一个性质连同 $\int_X f(x)h(x)\,\dd\mu(x)$ 的公式说 $h$ 确实是 $y \longmapsto \phi(\cdot, y)$ 的积分.

(d). 上面证明的 $h$ 的第二个性质正是 Minkowski 不等式.
\end{solution}


\chapter[Gelfand 理论]{Gelfand 理论}
% ============ 第 II 部分: 一些 Gelfand 理论 (Some Gelfand Theory) ============
% 原书 II.1 - II.5 (原书 pp.119-144)
% 编号: II.1.1.1 (子节级) / II.2.1 (节级) / 习题 II.5.1 (节级)

\section{Banach 代数}
在本节中, $A$ 是一个 Banach 代数. (见定义 \ref{def:I.4.1.4}.) 注意范数的次乘性蕴含乘法是从 $A \times A$ 到 $A$ 的连续映射.
我们现在假设 $A$ 有单位元 $e$, 记 $A^{\times}$ 为 $A$ 的可逆元群.

\subsection{元素的谱}

\begin{definition}[谱与谱半径]
\label{def:II.1.1.1}
设 $x \in A$.
\begin{enumerate}
\item[(i)] $x$ 在 $A$ 中的谱 (spectrum) 是
\[
\sigma(x) = \sigma_A(x) = \{\lambda \in \CC \mid \lambda e - x \notin A^{\times}\}.
\]
\item[(ii)] $x$ 的谱半径 (spectral radius) 是
\[
\rho(x) = \inf_{n \geq 1}\|x^n\|^{1/n}.
\]
\end{enumerate}
\end{definition}

我们将在下面看到 $\rho(x)$ 等于 $\sup\{|\lambda| \mid \lambda \in \sigma(x)\}$ (这证明``谱半径''这个名字是合理的).
我们从证明可逆元和谱半径的一些基本性质开始. (注意 (i) 不用 $A$ 的完备性, 所以它在任何赋范代数中都成立.)

\begin{proposition}[可逆元的基本性质]
\label{prop:II.1.1.2}
\begin{enumerate}
\item[(i)] 若 $x, y \in A^{\times}$ 满足 $\|x - y\| \leq \frac{1}{2}\|x^{-1}\|^{-1}$, 则
\[
\|x^{-1} - y^{-1}\| \leq 2\|x^{-1}\|^2\|y - x\|.
\]
特别地, 映射 $x \longmapsto x^{-1}$ 是从 $A^{\times}$ 到自身的同胚.
\item[(ii)] 对每个 $x \in A$,
\[
\rho(x) = \lim_{n \to +\infty}\|x^n\|^{1/n}
\]
(\index{Gelfand 公式}Gelfand 公式).
\item[(iii)] 设 $x \in A$. 若 $\rho(x) < 1$ (例如 $\|x\| < 1$), 则 $e - x \in A^{\times}$, 且 $(e - x)^{-1} = \sum_{n \geq 0}x^n$, 约定 $x^0 = e$. (特别地, 级数收敛.)
\item[(iv)] 群 $A^{\times}$ 在 $A$ 中开.
\end{enumerate}
\end{proposition}

\begin{proof}
(i). 有
\[
\|y^{-1}\| - \|x^{-1}\| \leq \|y^{-1} - x^{-1}\| = \|y^{-1}(x - y)x^{-1}\| \leq \|y^{-1}\|\|x - y\|\|x^{-1}\| \leq \frac{1}{2}\|y^{-1}\|.
\]
特别地, $\|y^{-1}\| \leq 2\|x^{-1}\|$. 把它与上面的不等式结合, 给出
\[
\|y^{-1} - x^{-1}\| \leq \|y^{-1}\|\|x - y\|\|x^{-1}\| \leq 2\|x^{-1}\|^2\|x - y\|,
\]
这就是第一个结论. 这也证明映射 $x \longmapsto x^{-1}$ 连续. 由于这个映射等于它自己的逆, 它是同胚.

(ii). 设 $\varepsilon > 0$. 我们要找到 $N \in \Z_{\geq 1}$ 使得对 $n \geq N$ 有 $\|x^n\|^{1/n} \leq \rho(x) + \varepsilon$. (我们已经知道由 $\rho(x)$ 的定义有 $\|x^n\|^{1/n} \geq \rho(x)$, 所以这足以确立结果.) 由 $\rho(x)$ 的定义, 可以找到 $m \geq 1$ 使得 $\|x^m\|^{1/m} \leq \rho(x) + \frac{1}{2}\varepsilon$. 对每个整数 $n \geq 1$, 把 $n$ 写成 $n = mq(n) + r(n)$, 其中 $q(n), r(n) \in \N$, $0 \leq r(n) \leq m - 1$. 注意
\[
\frac{q(n)}{n} = \frac{1}{m}\Bigl(1 - \frac{r(n)}{n}\Bigr) \xrightarrow[n \to +\infty]{} \frac{1}{m},
\]
所以
\[
\|x^m\|^{q(n)/n}\|x\|^{r(n)/n} \xrightarrow[n \to +\infty]{} \|x^m\|^{1/m}.
\]
选取 $N \geq 1$ 使得对 $n \geq N$ 有
\[
\|x^m\|^{q(n)/n}\|x\|^{r(n)/n} \leq \|x^m\|^{1/m} + \frac{\varepsilon}{2} \leq \rho(x) + \varepsilon.
\]
则若 $n \geq N$,
\[
\|x^n\|^{1/n} = \|x^{mq(n)}x^{r(n)}\|^{1/n} \leq \|x^m\|^{q(n)/n}\|x\|^{r(n)/n} \leq \rho(x) + \varepsilon,
\]
如所愿.\footnote{此证明中的推理有时称为 Fekete 引理. 参见 https://en.wikipedia.org/wiki/Subadditivity.}

(iii). 固定 $r \in \R$ 使得 $\rho(x) < r < 1$. 则由 (ii), 对足够大的 $n$ 有 $\|x^n\| \leq r^n$. 对每个 $n \in \N$, 记 $S_n = \sum_{k=0}^n x^k$. 则若 $m \geq n$ 足够大,
\[
\|S_m - S_n\| = \Bigl\|\sum_{k=n+1}^m x^k\Bigr\| \leq r^{n+1}\sum_{k \geq 0}r^k = r^{n+1}\frac{1}{1-r}.
\]
所以序列 $(S_n)_{n \geq 0}$ 是 Cauchy 列, 且因 $A$ 完备而收敛. 这意味着级数 $\sum_{n \geq 0}x^n$ 收敛. 而且, 对每个 $n \geq 0$,
\[
(e - x)S_n = S_n(e - x) = \sum_{k=0}^n x^k - \sum_{k=1}^{n+1}x^k = e - x^{n+1}.
\]
当 $n \to +\infty$ 时这趋于 $e$, 所以 $\sum_{n \geq 0}x^n$ 是 $e - x$ 的逆.

(iv). 设 $x \in A^{\times}$. 若 $y \in A$ 满足 $\|y - x\| < \|x^{-1}\|^{-1}$, 则
\[
\|e - x^{-1}y\| \leq \|x^{-1}\|\|x - y\| < 1.
\]
所以由 (iii), $x^{-1}y \in A^{\times}$, 从而 $y \in A^{\times}$.
\end{proof}

\begin{theorem}[谱的非空紧致性]
\label{thm:II.1.1.3}
对每个 $x \in A$, 谱 $\sigma_A(x)$ 是 $\CC$ 的非空紧子集, 且
\[
\rho(x) = \max\{|\lambda| \mid \lambda \in \sigma_A(x)\}.
\]
\end{theorem}

这解释了 $\rho(x)$ 被称为``谱半径''的原因. 特别地注意, 虽然 $x$ 的谱依赖于 $A$ (例如, 若考虑包含 $x$ 的 $A$ 的 Banach 子代数 $B$, 则 $\sigma_B(x) \supset \sigma_A(x)$, 但这可能不是等式), $x$ 的谱半径却不依赖.

\begin{proof}
考虑映射 $F \colon \CC \to A$, $\lambda \longmapsto \lambda e - x$. 则 $F$ 连续, 且 $\sigma_A(x)$ 是 $A$ 的闭子集 $A - A^{\times}$ 的原像, 所以 $\sigma_A(x)$ 在 $\CC$ 中闭.
其次, 设 $\lambda \in \CC$ 满足 $|\lambda| > \rho(x)$. 则 $\rho(\lambda^{-1}x) < 1$, 所以由命题 \ref{prop:II.1.1.2}(iii), $e - \lambda^{-1}x = \lambda^{-1}(\lambda e - x) \in A^{\times}$, 这立即蕴含 $\lambda \notin \sigma_A(x)$. 所以我们已证明
\[
\rho(x) \geq \sup\{|\lambda| \mid \lambda \in \sigma_A(x)\}.
\]
特别地, $\sigma_A(x)$ 是 $\CC$ 的闭且有界子集, 所以它紧.
让我们证明 $\sigma_A(x)$ 非空. 设 $T \colon A \to \CC$ 是有界线性泛函, 定义 $f \colon \CC - \sigma_A(x) \to \CC$ 为 $f(\lambda) = T((\lambda e - x)^{-1})$. 若 $\lambda, \mu \in \CC - \sigma_A(x)$, 则
\[
(\lambda e - x)^{-1} - (\mu e - x)^{-1} = (\lambda e - x)^{-1}((\mu e - x) - (\lambda e - x))(\mu e - x)^{-1} = -(\lambda - \mu)(\lambda e - x)^{-1}(\mu e - x)^{-1},
\]
所以若 $\lambda \neq \mu$,
\[
\frac{f(\lambda) - f(\mu)}{\lambda - \mu} = -T((\lambda e - x)^{-1}(\mu e - x)^{-1}).
\]
利用函数 $y \longmapsto y^{-1}$ 的连续性 (见命题 \ref{prop:II.1.1.2}(i)), 对每个 $\lambda \in \CC - \sigma_A(x)$ 得到
\[
\lim_{\mu \to \lambda}\frac{f(\lambda) - f(\mu)}{\lambda - \mu} = -T((\lambda e - x)^{-2}).
\]
特别地, 函数 $f$ 在 $\CC - \sigma_A(x)$ 上全纯. 让我们证明 $f$ 在无穷远处消失, 即当 $|\lambda| \to +\infty$ 时 $f(\lambda)$ 趋于 $0$. 设 $\lambda \in \CC$ 满足 $|\lambda| > \rho(x)$. 则由命题 \ref{prop:II.1.1.2}(iii),
\[
(\lambda e - x)^{-1} = \lambda^{-1}(e - \lambda^{-1}x)^{-1} = \frac{1}{\lambda}\sum_{n \geq 0}\frac{x^n}{\lambda^n},
\]
所以
\[
\|(\lambda e - x)^{-1}\| \leq \frac{1}{|\lambda|}\sum_{n \geq 0}\frac{\|x\|^n}{|\lambda|^n} = \frac{1}{|\lambda|}\cdot\frac{1}{1 - |\lambda|^{-1}\|x\|}.
\]
当 $|\lambda| \to +\infty$ 时这趋于 $0$; 由于 $T$ 连续, $f(\lambda)$ 也如此.
现在假设 $\sigma_A(x) = \emptyset$. 则 $f$ 是整函数, 且当 $|\lambda| \to +\infty$ 时 $f(\lambda) \to 0$. 由 Liouville 定理, 这蕴含 $f = 0$, 即对每个 $\lambda \in \CC$ 有 $T((\lambda e - x)^{-1}) = 0$. 但这在 $\Hom(A, \CC)$ 的每个 $T$ 上都成立, 且由 Hahn-Banach 定理 $A$ 上的有界线性泛函分离点, 所以对每个 $\lambda \in \CC$ 有 $(\lambda e - x)^{-1} = 0$. 这不可能, 因为 $(\lambda e - x)^{-1} \in A^{\times}$. 所以 $\sigma_A(x) \neq \emptyset$.
最后, 我们证明
\[
\rho(x) \leq \max\{|\lambda| \mid \lambda \in \sigma_A(x)\}.
\]
设 $r = \max\{|\lambda| \mid \lambda \in \sigma_A(x)\}$. 我们已经知道 $r \leq \rho(x)$. 假设 $r < \rho(x)$, 选取 $r'$ 使得 $r < r' < \rho(x)$. 设 $T \in \Hom(A, \CC)$, 如前定义 $f \colon \CC - \sigma_A(x) \to \CC$. 则我们已见 $f$ 在 $\CC - \sigma_A(x) \supset \{\lambda \in \CC \mid |\lambda| > r\}$ 上全纯. 我们还看到, 若 $|\lambda| > \rho(x)$, 则
\[
(\lambda e - x)^{-1} = \sum_{n \geq 0}\frac{x^n}{\lambda^{n+1}},
\]
所以
\[
f(\lambda) = \sum_{n \geq 0}\frac{T(x^n)}{\lambda^{n+1}}.
\]
由幂级数展开的唯一性, 这对 $|\lambda| > r$ 仍成立. 特别地, 级数 $\sum_{n \geq 0}\frac{T(x^n)}{(r')^{n+1}}$ 收敛, 所以序列 $\bigl(\frac{T(x^n)}{(r')^{n+1}}\bigr)_{n \geq 0}$ 收敛到 $0$, 从而有界. 考虑 $\Hom(A, \CC)$ 上的有界线性泛函序列 $(\alpha_n)_{n \geq 0}$, 定义为 $\alpha_n(T) = \frac{T(x^n)}{(r')^{n+1}}$. 我们刚看到对每个 $T \in \Hom(A, \CC)$, 序列 $(\alpha_n(T))_{n \geq 0}$ 有界. 由共鸣定理 (定理 \ref{thm:I.3.2.11}), 这蕴含序列 $(\|\alpha_n\|_{\mathrm{op}})_{n \geq 0}$ 有界. 但注意由 Hahn-Banach 定理, $\|\alpha_n\|_{\mathrm{op}} = \frac{\|x^n\|}{(r')^{n+1}}$. 所以序列 $\bigl(\frac{\|x^n\|}{(r')^{n+1}}\bigr)_{n \geq 0}$ 有界. 选取控制它的实数 $C$. 则
\[
\rho(x) = \lim_{n \to +\infty}\|x^n\|^{1/n} \leq \lim_{n \to +\infty}C^{1/n}(r')^{(n+1)/n} = r',
\]
矛盾. 所以 $r \geq \rho(x)$.
\end{proof}

\subsection{\index{Gelfand-Mazur 定理}Gelfand-Mazur 定理}

众所周知, 每个是域的有限维 $\CC$-代数都同构于 $\CC$. 这就是它的 Banach 代数版本.

\begin{corollary}[Gelfand-Mazur 定理]
\label{cor:II.1.2.1}
设 $A$ 是每个非零元素都可逆的 Banach 代数. 则 $A$ 同构于 $\CC$ (即 $A = \CC e$).
\end{corollary}

\begin{proof}
设 $x \in A$. 由定理 \ref{thm:II.1.1.3}, $\sigma_A(x) \neq \emptyset$. 设 $\lambda \in \sigma_A(x)$, 则 $\lambda e - x$ 不可逆, 所以由假设 $x = \lambda e$.
\end{proof}

\begin{definition}[理想]
\label{def:II.1.2.2}
我们说 $A$ 的子集 $I$ 是理想 (ideal), 如果它是通常代数意义上的理想, 即 $I$ 是 $A$ 的 $\CC$-子空间, 且被 $A$ 的每个元素的左右乘法稳定. 若 $I$ 是 $A$ 的理想且 $I \neq A$, \index{真理想}则称 $I$ 是 $A$ 的真理想 (proper ideal).
\end{definition}

若 $I$ 是 $A$ 的理想, 则容易看出 $\overline{I}$ 也是理想.
还要回忆商范数的定义.

\begin{definition}[商范数]
\label{def:II.1.2.3}
设 $V$ 是赋范向量空间, $W \subset V$ 是闭子空间. 则 $V/W$ 上的\index{商范数}商范数 (quotient norm) 定义为
\[
\|x + W\| = \inf_{w \in W}\|v + w\|.
\]
\end{definition}

若 $V$ 是巴拿赫空间, 则 $V/W$ 也是 (对商范数).

\begin{proposition}[商代数的 Banach 代数结构]
\label{prop:II.1.2.4}
\begin{enumerate}
\item[(i)] 若 $I$ 是 $A$ 的闭理想, 则 $A/I$ 对商范数是 Banach 代数.
\item[(ii)] 若 $I$ 是 $A$ 的真理想, 则其闭包 $\overline{I}$ 也是.
\end{enumerate}
\end{proposition}

\begin{proof}
(i). 我们已经知道 $A/I$ 是巴拿赫空间和代数, 所以只需验证其范数次乘. 设 $x, y \in A$. 则
\begin{align*}
\|x + I\|\|y + I\| &= \inf_{a,b \in I}\|x + a\|\|y + b\| \\
&\geq \inf_{a,b \in I}\|(x + a)(y + b)\| \\
&= \inf_{a,b \in I}\|xy + (ay + xb + ab)\| \\
&\geq \inf_{c \in I}\|xy + c\| \quad \text{(因为若 } a, b \in I \text{ 则 } ay + xb + ab \in I\text{)} \\
&= \|xy + I\|.
\end{align*}

(ii). 考虑开球 $B = \{x \in A \mid \|e - x\| < 1\}$. 则由命题 \ref{prop:II.1.1.2}, $B \subset A^{\times}$, 所以 $B \cap I = \emptyset$. 由于 $B$ 开, 这蕴含 $B \cap \overline{I} = \emptyset$, 所以 $\overline{I} \neq A$.
\end{proof}

\begin{corollary}[极大理想与 $\CC$]
\label{cor:II.1.2.5}
设 $A$ 是交换含幺 Banach 代数. 若 $\mathfrak{m}$ 是 $A$ 的极大理想, 则 $\mathfrak{m}$ 闭, 且 $A/\mathfrak{m} = \CC$.
\end{corollary}

这是 Nullstellensatz 的 Banach 代数版本.

\begin{proof}
由命题 \ref{prop:II.1.2.4}, 理想 $\overline{\mathfrak{m}}$ 也真; 由于 $\mathfrak{m}$ 极大, 必有 $\mathfrak{m} = \overline{\mathfrak{m}}$, 即 $\mathfrak{m}$ 闭. 由同一命题, $A/\mathfrak{m}$ 是 Banach 代数. 又, $A/\mathfrak{m}$ 的每个非零元素都可逆 (因为 $\mathfrak{m}$ 极大), 所以由 Gelfand-Mazur 定理 $A/\mathfrak{m} = \CC$.
\end{proof}

% ============ II.2 \index{Banach 代数的谱}Banach 代数的谱 (原书 pp.124-127) ============
\section{Banach 代数的谱}
在本节中, $A$ 仍是 Banach 代数, 但我们不假设它有单位元.

\begin{definition}[\index{乘法泛函}乘法泛函与谱]
\label{def:II.2.1}
$A$ 上的乘法泛函 (multiplicative functional) 是非零线性泛函 $\phi \colon A \to \CC$, 使得对所有 $x, y \in A$ 有 $\phi(xy) = \phi(x)\phi(y)$.
$A$ 上所有乘法泛函的集合称为 $A$ 的谱 (spectrum), 记为 $\sigma(A)$. 我们在 $\sigma(A)$ 上赋予弱*拓扑. 换句话说, 若 $\phi \in \sigma(A)$, 则 $\phi$ 的开邻域基由集合 $\{\psi \in \sigma(A) \mid \forall i \in \{1, \ldots, n\},\, |\phi(x_i) - \psi(x_i)| < c_i\}$ 给出, 其中 $n \in \Z_{\geq 1}$, $x_1, \ldots, x_n \in A$, $c_1, \ldots, c_n \in \R_{>0}$.
\end{definition}

注意我们不假设 $\phi$ 连续; 事实上这自动成立, 我们将在下面看到.

\begin{lemma}[乘法泛函在单位元上取值 $1$]
\label{lem:II.2.2}
若 $A$ 含幺, 则对每个 $\phi \in \sigma(A)$ 有 $\phi(e) = 1$ 且 $\phi(A^{\times}) \subset \CC^{\times}$.
\end{lemma}

\begin{proof}
设 $x \in A$ 使得 $\phi(x) \neq 0$. 则 $\phi(x) = \phi(xe) = \phi(x)\phi(e)$, 所以 $\phi(e) = 1$. 又, 若 $y \in A^{\times}$, 则 $1 = \phi(e) = \phi(y)\phi(y^{-1})$, 所以 $\phi(y) \in \CC^{\times}$.
\end{proof}

\begin{definition}[单位化 (adjoining an identity)]
\label{def:II.2.3}
设 $A$ 是 Banach 代数. 我们如下定义含幺 Banach 代数 $A^{\tilde{e}}$: 取 $\CC$-向量空间 $A \oplus \CC e$, 在 $A^{\tilde{e}}$ 上定义乘法 $(x + \lambda e)(y + \mu e) = (xy + \lambda y + \mu x) + \lambda\mu e$ (对 $x, y \in A$, $\lambda, \mu \in \CC$) 和范数 $\|x + \lambda e\| = \|x\| + |\lambda|$ (对 $x \in A$, $\lambda \in \CC$). 若 $A$ 是 Banach $*$-代数, 则令 $(x + \lambda e)^* = x^* + \bar{\lambda}e$ (对 $x \in A$, $\lambda \in \CC$), 把 $A^{\tilde{e}}$ 变成 Banach $*$-代数.
此构造称为给 $A$ \index{添加单位元}添加单位元 (adjoining an identity).
\end{definition}

\begin{remark}[$A$ 已含幺的情形]
\label{rem:II.2.4}
若 $A$ 已有单位元, 则 $A^{\tilde{e}}$ 不等于 $A$. 事实上, 若记 $e_A$ 为 $A$ 的单位元, 则映射 $A^{\tilde{e}} \to A \times \CC$, $x + \lambda e \longmapsto (x + \lambda e_A, \lambda)$ 是 $\CC$-代数的同构 (也是同胚).
\end{remark}

\begin{proposition}[谱的嵌入]
\label{prop:II.2.5}
对每个 $\phi \in \sigma(A)$, 通过令 $\tilde{\phi}(x + \lambda e) = \phi(x) + \lambda$ 得到 $\sigma(A^{\tilde{e}})$ 的元素 $\tilde{\phi}$. 这定义单射映射 $\sigma(A) \to \sigma(A^{\tilde{e}})$, 其像是 $\sigma(A^{\tilde{e}}) - \{\phi_{\infty}\}$, 其中 $\phi_{\infty}$ 由 $\phi_{\infty}(x + \lambda e) = \lambda$ 定义.
\end{proposition}

以后我们将等同 $\phi$ 与 $\tilde{\phi}$, 直接写 $\sigma(A^{\tilde{e}}) = \sigma(A) \cup \{\phi_{\infty}\}$.

\begin{proof}
$\tilde{\phi}$ 是乘法泛函这一事实直接由 $A^{\tilde{e}}$ 上乘法的定义得出, 且 $\tilde{\phi}$ 显然确定 $\phi$. 所以我们只需验证关于 $\sigma(A) \to \sigma(A^{\tilde{e}})$ 的像的结论.
设 $\psi \in \sigma(A^{\tilde{e}})$ 满足 $\psi \neq \phi_{\infty}$, 令 $\phi = \psi|_A$. 则对所有 $x \in A$ 和 $\lambda \in \CC$ 有 $\psi(x + \lambda e) = \phi(x) + \lambda$; 由于 $\psi \neq \phi_{\infty}$, 线性泛函 $\phi \colon A \to \CC$ 不能为零, 所以 $\phi$ 是 $A$ 上的乘法泛函, 且显然 $\psi = \tilde{\phi}$.
\end{proof}

\begin{corollary}[乘法泛函自动有界]
\label{cor:II.2.6}
设 $\phi \in \sigma(A)$. 则 $\phi$ 是 $A$ 上的有界线性函数, 且 $\|\phi\|_{\mathrm{op}} \leq 1$, 若 $A$ 含幺则等号成立.
\end{corollary}

\begin{proof}
由命题 \ref{prop:II.2.5}, 乘法泛函 $\phi$ 延拓成 $A^{\tilde{e}}$ 上的乘法泛函 $\tilde{\phi}$. 设 $x \in A$. 对每个满足 $|\lambda| > \|x\|$ 的 $\lambda \in \CC$, 元素 $x - \lambda e$ 在 $A^{\tilde{e}}$ 中可逆 (由命题 \ref{prop:II.1.1.2}), 所以 $\phi(x) - \lambda = \tilde{\phi}(x - \lambda e) \neq 0$. 这蕴含 $|\phi(x)| \leq \|x\|$, 即 $\phi$ 有界且 $\|\phi\|_{\mathrm{op}} \leq 1$.
若 $A$ 含幺, 则 $\|e\| = 1$ 且 $\phi(e) = 1$, 所以 $\|\phi\|_{\mathrm{op}} = 1$.
\end{proof}

\begin{theorem}[谱的紧致性]
\label{thm:II.2.7}
设 $A$ 是 Banach 代数.
\begin{enumerate}
\item[(i)] 若 $A$ 含幺, 则空间 $\sigma(A)$ 紧 Hausdorff.
\item[(ii)] 一般地, 空间 $\sigma(A)$ 局部紧 Hausdorff, 且 $\sigma(A^{\tilde{e}})$ 是它的 Alexandroff 紧化 (a.k.a. 单点紧化).
\end{enumerate}
\end{theorem}

回忆若 $X$ 是 Hausdorff 局部紧拓扑空间, 则其 Alexandroff 紧化是空间 $X \cup \{\infty\}$ (即加上一个点的 $X$), 且其开子集是 $X$ 的开子集和 $X \cup \{\infty\}$ 中 $X$ 的紧子集的余集.

\begin{proof}
由推论 \ref{cor:II.2.6}, 谱 $\sigma(A)$ 是 $\Hom(A, \CC)$ 的闭单位球的子集. 我们知道这个闭单位球对 $\Hom(A, \CC)$ 上的弱*拓扑紧 Hausdorff (这是 Alaoglu 定理), 且 $\sigma(A) \cup \{0\}$ 在此拓扑中闭, 因为它由闭条件 $\phi(xy) = \phi(x)\phi(y)$ (对所有 $x, y \in A$) 定义. 所以 $\sigma(A) \cup \{0\}$ 紧 (对弱*拓扑), 且它的开子集 $\sigma(A)$ 局部紧. 若 $A$ 含幺, 则 $\sigma(A)$ 在 $\sigma(A) \cup \{0\}$ 中闭, 因为它由条件 $\phi(e) = 1$ 切出, 所以 $\sigma(A)$ 紧.
现在我们证明 (ii) 的最后一句. 若 $\phi \in \sigma(A)$ (相应地 $\sigma(A^{\tilde{e}})$), $x \in A$, $c > 0$, 令
\[
U(\phi, x, c) = \{\psi \in \sigma(A) \mid |\phi(x) - \psi(x)| < c\}
\]
(相应地 $\tilde{U}(\phi, x, c) = \{\psi \in \sigma(A^{\tilde{e}}) \mid |\phi(x) - \psi(x)| < c\}$).
这些构成 $\sigma(A)$ (相应地 $\sigma(A^{\tilde{e}})$) 的拓扑的基.
若 $\phi \in \sigma(A)$, $x \in A$, $c > 0$, 则
\[
\tilde{U}(\tilde{\phi}, x, c) = \begin{cases} U(\phi, x, c) \cup \{\phi_{\infty}\} & \text{若 } |\phi(x)| < c, \\ U(\phi, x, c) & \text{否则}. \end{cases}
\]
对 $\phi_{\infty}$ 的邻域, 若 $x \in A$, $c > 0$, 则
\[
\tilde{U}(\phi_{\infty}, x, c) = \{\phi_{\infty}\} \cup \{\phi \in \sigma(A) \mid |\phi(x)| < c\} = \sigma(A^{\tilde{e}}) - \{\psi \in \sigma(A^{\tilde{e}}) \mid |\psi(x)| \geq c\}.
\]
所以 $\sigma(A)$ 的拓扑由 $\sigma(A^{\tilde{e}})$ 的拓扑诱导. 又, 由于对所有的 $x \in A$ 和 $c > 0$, 集合 $\{\psi \in \sigma(A^{\tilde{e}}) \mid |\psi(x)| \geq c\}$ 在 $\sigma(A^{\tilde{e}})$ 中闭, 从而紧, $\phi_{\infty}$ 在 $\sigma(A^{\tilde{e}})$ 中的开邻域正好是 $\sigma(A)$ 的紧子集的余集. 这意味着 $\sigma(A^{\tilde{e}})$ 是 $\sigma(A)$ 的 Alexandroff 紧化.
\end{proof}

\begin{definition}[\index{Gelfand 变换}Gelfand 变换]
\label{def:II.2.8}
设 $A$ 是 Banach 代数. 对每个 $x \in A$, 映射 $\hat{x} \colon \sigma(A) \to \CC$, $\hat{x}(\phi) = \phi(x)$ 称为 $x$ 的 Gelfand 变换 (Gelfand transform).
注意每个 $\hat{x}$ 由 $\sigma(A)$ 的拓扑定义在 $\sigma(A)$ 上连续. 由此产生的映射 $\Gamma \colon A \to \mathcal{C}(\sigma(A))$, $x \longmapsto \hat{x}$ 称为 $A$ 的 \index{Gelfand 表示}Gelfand 表示 (Gelfand representation) (有时也称为 Gelfand 变换).
\end{definition}

注意 $\Gamma$ 由 $\mathcal{C}(\sigma(A))$ 上的代数运算的定义是 $\CC$-代数的同态.

\begin{theorem}[Gelfand 表示的性质]
\label{thm:II.2.9}
\begin{enumerate}
\item[(i)] 映射 $\Gamma$ 把 $A$ 映到 $\mathcal{C}_0(\sigma(A))$, 且对每个 $x \in A$ 有 $\|\hat{x}\|_{\infty} \leq \|x\|$.
\item[(ii)] $\Gamma$ 的像分离 $\sigma(A)$ 的点.
\item[(iii)] 若 $A$ 含幺, 则 $\hat{e}$ 是 $\sigma(A)$ 上的常值函数 $1$.
\end{enumerate}
\end{theorem}

\begin{proof}
(i). 若 $A$ 含幺, 则 $\sigma(A)$ 紧, 所以 $\mathcal{C}_0(\sigma(A)) = \mathcal{C}(\sigma(A))$. 一般地, 由于 $\sigma(A^{\tilde{e}})$ 是 $\sigma(A)$ 的 Alexandroff 紧化, 只需验证对每个 $x \in A$ 有 $\hat{x}(\phi_{\infty}) = 0$; 但这由定义立即得出.
设 $x \in A$. 则
\[
\|\hat{x}\|_{\infty} = \sup_{\phi \in \sigma(A)}|\hat{x}(\phi)| = \sup_{\phi \in \sigma(A)}|\phi(x)| \leq \|x\|
\]
由推论 \ref{cor:II.2.6}.
(ii). 设 $\phi, \phi' \in \sigma(A)$ 满足 $\phi \neq \phi'$. 则存在 $x \in A$ 使得 $\phi(x) \neq \phi'(x)$, 即 $\hat{x}(\phi) \neq \hat{x}(\phi')$.
(iii). 这由引理 \ref{lem:II.2.2} 立即得出.
\end{proof}

对一般的 Banach 代数 (即使是含幺的), 谱可以为空 (见习题 \ref{ex:II.5.1}). 但这不能发生在交换 Banach 代数上.

\begin{theorem}[谱与极大理想的双射]
\label{thm:II.2.10}
设 $A$ 是交换含幺 Banach 代数. 则映射 $\phi \longmapsto \Ker(\phi)$ 诱导从 $\sigma(A)$ 到 $A$ 的极大理想集合的双射.
\end{theorem}

若你见过谱的另一个定义 (例如在代数几何中), 此定理说明它与我们的定义如何联系.

\begin{proof}
若 $\phi \in \sigma(A)$, 则 $A/\Ker(\phi) \simeq \CC$ (注意 $\phi$ 因非零而满射), 所以 $\Ker(\phi)$ 是 $A$ 的极大理想. 这证明映射良定义.
若 $\mathfrak{m}$ 是极大理想, 则由 Gelfand-Mazur 定理 $A/\mathfrak{m} \simeq \CC$ (见推论 \ref{cor:II.1.2.5}), 所以映射 $\phi \colon A \to A/\mathfrak{m} \simeq \CC$ 是满足 $\Ker(\phi) = \mathfrak{m}$ 的 $\sigma(A)$ 的元素. 这证明映射满射.
现在我们需要验证单射性. 设 $\phi, \psi \in \sigma(A)$ 满足 $\mathfrak{m} := \Ker(\phi) = \Ker(\psi)$. 设 $x \in A$. 由于 $A/\mathfrak{m} \simeq \CC$, 可以把 $x$ 写成 $x = \lambda e + y$, 其中 $\lambda \in \CC$, $y \in \mathfrak{m}$. 则
\[
\phi(x) = \lambda = \psi(y).
\]
所以 $\phi = \psi$.
\end{proof}

\begin{corollary}[谱半径的 Gelfand 公式]
\label{cor:II.2.11}
设 $A$ 是交换含幺 Banach 代数. 则对每个 $x \in A$:
\begin{enumerate}
\item[(i)] $x \in A^{\times}$ 当且仅当 $\hat{x}$ 不取零值;
\item[(ii)] $\hat{x}(\sigma(A)) = \sigma_A(x)$;
\item[(iii)] $\|\hat{x}\|_{\infty} = \rho(x)$.
\end{enumerate}
\end{corollary}

\begin{proof}
(i). 若 $x \in A^{\times}$, 则 $\hat{x}$ 不能消失, 因为有 $\widehat{x}\widehat{x^{-1}} = \hat{e} = 1$. 反之, 假设 $x$ 不可逆. 则存在包含 $x$ 的极大理想, 所以由定理 \ref{thm:II.2.10} 存在 $\phi \in \sigma(A)$ 使得 $0 = \phi(x) = \hat{x}(\phi)$.
(ii). 由 (i),
\[
\sigma_A(x) = \{\lambda \in \CC \mid x - \lambda e \notin A^{\times}\} = \{\lambda \in \CC \mid \widehat{x - \lambda e} \text{ 至少在一个点消失}\} = \hat{x}(\sigma(A)).
\]
(iii). 这由 (ii) 和定理 \ref{thm:II.1.1.3} 得出.
\end{proof}

% ============ II.3 C* 代数与 \index{Gelfand-Naimark 定理}Gelfand-Naimark 定理 (原书 pp.128-130) ============
\section{\texorpdfstring{\index{$\mathcal{C}^*$-代数}$\mathcal{C}^*$-代数与 Gelfand-Naimark 定理}{C*-代数与 Gelfand-Naimark 定理}}

\begin{definition}[$\mathcal{C}^*$-代数]
\label{def:II.3.1}
Banach $*$-代数 $A$ 称为 $\mathcal{C}^*$-代数 ($\mathcal{C}^*$-algebra), 如果对每个 $x \in A$ 有 $\|x^*x\| = \|x\|^2$.
\end{definition}

\begin{remark}[Bourbaki 的术语]
\label{rem:II.3.2}
除 Bourbaki 外大家都称它为 $\mathcal{C}^*$-代数, Bourbaki 说``星代数'' (algèbre stellaire).
\end{remark}

\begin{lemma}[对合保范]
\label{lem:II.3.3}
若 $A$ 是 $\mathcal{C}^*$-代数, 则对每个 $x \in A$ 有 $\|x\| = \|x^*\|$.
\end{lemma}

\begin{proof}
设 $x \in A - \{0\}$. 则
\[
\|x\|^2 = \|x^*x\| \leq \|x^*\|\|x\|,
\]
所以 $\|x\| \leq \|x^*\|$. 把它应用于 $x^*$ 并利用 $(x^*)^* = x$, 得到 $\|x^*\| \leq \|x\|$.
\end{proof}

\begin{example}[$\mathcal{C}^*$-代数的例子]
\label{exa:II.3.4}
例 \ref{exa:I.4.2.2} 的大多数例子实际上是 $\mathcal{C}^*$-代数.
\begin{enumerate}
\item[(a)] $\CC$ 是 $\mathcal{C}^*$-代数, 因为对每个 $\lambda \in \CC$ 有 $|\lambda\bar{\lambda}| = |\lambda|^2$.
\item[(b)] 设 $G$ 是局部紧群. 则 $L^1(G)$ 一般不是 $\mathcal{C}^*$-代数, 尽管它确实满足引理 \ref{lem:II.3.3} 的结论.\footnote{有一种方法修改 $L^1(G)$ 上的范数, 使新范数下的完备化成为 $\mathcal{C}^*$-代数, 但我们这里不需要.}
\item[(c)] 设 $X$ 是局部紧 Hausdorff 空间. 则 $\mathcal{C}_0(X)$ 是 $\mathcal{C}^*$-代数, 因为对每个 $f \in \mathcal{C}_0(X)$,
\[
\|f^*f\|_{\infty} = \sup_{x \in X}|f(x)\overline{f(x)}| = \sup_{x \in X}|f(x)|^2 = \|f\|_{\infty}^2.
\]
\item[(d)] 设 $V$ 是希尔伯特空间. 则 $\End(V)$ 是 $\mathcal{C}^*$-代数. 事实上, 设 $T \in \End(V)$. 我们要证明 $\|T^*T\|_{\mathrm{op}} = \|T\|_{\mathrm{op}}^2$. 首先注意
\[
\|T^*\|_{\mathrm{op}} = \sup_{\substack{v,w \in V\\ \|v\| = \|w\| = 1}}|\langle T^*(v), w \rangle| = \sup_{\substack{v,w \in V\\ \|v\| = \|w\| = 1}}|\langle v, T(w) \rangle| = \|T\|_{\mathrm{op}},
\]
所以 $\|T^*T\|_{\mathrm{op}} \leq \|T^*\|_{\mathrm{op}}\|T\|_{\mathrm{op}} = \|T\|_{\mathrm{op}}^2$. 另一方面,
\[
\|T^*T\|_{\mathrm{op}} = \sup_{\substack{v \in V\\ \|v\| = 1}}\|T^*T(v)\| \geq \sup_{\substack{v \in V\\ \|v\| = 1}}|\langle T^*T(v), v \rangle| = \sup_{\substack{v \in V\\ \|v\| = 1}}|\langle T(v), T(v) \rangle| = \|T\|_{\mathrm{op}}^2.
\]
\end{enumerate}
\end{example}

\begin{proposition}[乘法泛函是 $*$-同态]
\label{prop:II.3.5}
设 $A$ 是 $\mathcal{C}^*$-代数. 则 Gelfand 表示 $\Gamma \colon A \to \mathcal{C}_0(\sigma(A))$ 是 $*$-同态.
\end{proposition}

\begin{remark}[对称性]
\label{rem:II.3.6}
此命题说 $A$ 上的每个乘法泛函都是 $*$-同态. 满足此条件的 Banach $*$-代数称为\index{对称 $*$-代数}对称 (symmetric) 的. 每个 $\mathcal{C}^*$-代数都对称, 但其逆不真. (例如, 若 $G$ 是局部紧交换群, 则 $L^1(G)$ 对称, 见习题 \ref{ex:II.5.4}.)
\end{remark}

\begin{proof}
通过给 $A$ 添加单位元, 可以约化到 $A$ 含幺的情形. (见习题 \ref{ex:II.5.6} 中对 $A^{\tilde{e}}$ 范数的正确选取. 注意改变 $A^{\tilde{e}}$ 上的范数不影响 $\sigma(A^{\tilde{e}})$, 因为谱的定义不涉及范数.)
设 $x \in A$, $\phi \in \sigma(A)$. 我们要证明 $\phi(x^*) = \widehat{x^*}(\phi) = \overline{\hat{x}(\phi)} = \overline{\phi(x)}$. 把 $\phi(x)$ 写成 $\phi(x) = a + ib$, $\phi(x^*) = c + id$, 其中 $a, b, c, d \in \R$.
假设 $b + d \neq 0$. 设
\[
y = \frac{1}{b + d}(x + x^* - (a + c)e) \in A.
\]
注意 $y = y^*$, 且
\[
\phi(y) = \frac{1}{b + d}(a + ib + c + id - (a + c)) = i,
\]
所以对每个 $t \in \R$ 有 $\phi(y + ite) = (1 + t)i$, 从而
\[
|1 + t| = |\phi(y + ite)| \leq \|y + ite\|
\]
(由推论 \ref{cor:II.2.6}). 利用 $\mathcal{C}^*$-代数的定义性质和 $y = y^*$ 这一事实, 对每个 $t \in \R$,
\[
(1 + t)^2 \leq \|y + ite\|^2 = \|(y + ite)(y + ite)^*\| = \|(y + ite)(y - ite)\| = \|y^2 + t^2e\| \leq \|y^2\| + t^2,
\]
即 $1 + 2t \leq \|y^2\|$. 但这蕴含 $\|y^2\|$ 无穷, 不可能. 所以 $b + d = 0$, 即 $d = -b$.
把同样的推理应用于 $ix$ 和 $(ix)^* = -ix^*$ (注意 $\phi(ix) = -b + ia$ 且 $\phi((ix)^*) = d - ic$), 得到 $a - c = 0$, 即 $a = c$. 这完成了 $\phi(x^*) = \overline{\phi(x)}$ 的证明.
\end{proof}

\begin{proposition}[自伴元的范数等于谱半径]
\label{prop:II.3.7}
设 $A$ 是交换含幺 $\mathcal{C}^*$-代数. 则对每个 $x \in A$ 有 $\|x\| = \rho(x)$.
\end{proposition}

\begin{proof}
若 $x \in A$ 满足 $x = x^*$, 则 $\|x\|^2 = \|x^*x\| = \|x^2\|$, 所以对每个 $n \in \N$ 有 $\|x^{2^n}\| = \|x\|^{2^n}$.
现在设 $x$ 是 $A$ 的任意元素. 则 $(xx^*)^* = xx^*$, 所以第一部分适用于 $xx^*$. 又, 对每个 $n \in \N$, $(xx^*)^{2^n} = x^{2^n}(x^*)^{2^n}$ (因为 $A$ 交换). 所以若 $n \geq 0$,
\[
\|x\|^{2^{n+1}} = \|xx^*\|^{2^n} = \|(xx^*)^{2^n}\| = \|x^{2^n}(x^*)^{2^n}\| = \|x^{2^n}\|^2.
\]
这蕴含
\[
\rho(x) = \lim_{n \to +\infty}\|x^{2^n}\|^{2^{-n}} = \|x\|.
\]
\end{proof}

\begin{definition}[\index{正规元}正规元]
\label{def:II.3.8}
若 $A$ 是 Banach $*$-代数, $A$ 的元素 $x$ 称为正规 (normal) 的, 如果 $xx^* = x^*x$.
\end{definition}

\begin{corollary}[正规元的范数]
\label{cor:II.3.9}
设 $A$ 是含幺 $\mathcal{C}^*$-代数, $x \in A$ 是正规元. 则 $\rho(x) = \|x\|$.
\end{corollary}

特别地, 若 $V$ 是希尔伯特空间且 $T \in \End(V)$ 正规, 则 $\|T\|_{\mathrm{op}} = \rho(T)$.

\begin{proof}
事实上, 由于 $x$ 与 $x^*$ 交换, $A$ 中包含 $x$ 和 $x^*$ 的最小含幺 $\CC$-代数 $A'$ 的闭包是交换 $\mathcal{C}^*$-代数, 且把 $x$ 看作 $A'$ 的元素时 $\rho(x)$ 和 $\|x\|$ 不变.
\end{proof}

\begin{theorem}[Gelfand-Naimark 定理]
\label{thm:II.3.10}
设 $A$ 是交换含幺 $\mathcal{C}^*$-代数. 则 Gelfand 表示 $\Gamma \colon A \to \mathcal{C}(\sigma(A))$ 是等距 $*$-同构.
\end{theorem}

\begin{proof}
由命题 \ref{prop:II.3.5} 知道 $\Gamma$ 是 $*$-同态, 由推论 \ref{cor:II.2.11}(iii) 和命题 \ref{prop:II.3.7} 知道它是等距. 特别地, $\Gamma$ 单射. 所以只需证明它满射. 由于 $\Gamma$ 等距且 $A$ 完备, 像 $\Gamma(A)$ 在 $\mathcal{C}(\sigma(A))$ 中闭; 但它由定理 \ref{thm:II.2.9}(ii) 分离点, 且因 $\Gamma(e) = 1$ 包含常值函数, 所以由 Stone-Weierstrass 定理它等于 $\mathcal{C}(\sigma(A))$.
\end{proof}

容易看出 Gelfand-Naimark 定理蕴含下面的结果 (但我们不需要它).

\begin{corollary}[非幺情形的 Gelfand-Naimark 定理]
\label{cor:II.3.11}
设 $A$ 是交换 $\mathcal{C}^*$-代数. 则 Gelfand 表示 $\Gamma \colon A \to \mathcal{C}_0(\sigma(A))$ 是等距 $*$-同构.
\end{corollary}

% ============ II.4 \index{谱定理}谱定理 (原书 pp.131-133) ============
\section{谱定理}

\begin{theorem}[谱定理]
\label{thm:II.4.1}
设 $V$ 是希尔伯特空间, $T \in \End(V)$ 正规. 记 $A_T$ 为 $\End(V)$ 中由 $T$ 和 $T^*$ 生成的含幺子代数的闭包; 它是交换的, 因为 $T$ 与 $T^*$ 交换.
则存在等距 $*$-同构 $\Phi \colon \mathcal{C}(\sigma(T)) \to A_T$ 使得, 若 $\iota$ 是 $\sigma(T)$ 到 $\CC$ 的注入, 则 $\Phi(\iota) = T$.
\end{theorem}

注意我们直接写 $\sigma(T)$ 表示 $\sigma_{\End(V)}(T)$ (这是 $T$ 的通常谱).
这个定理看起来不太像有限维线性代数的谱定理. 见习题 \ref{ex:II.5.2} 了解两者之间如何过渡.

\begin{proof}
设 $A = A_T$. 首先我们用 $\sigma_A(T)$ 代替 $\sigma(T)$ 证明结果, 然后证明 $\sigma(T) = \sigma_A(T)$. 注意我们自动有 $\sigma(T) \subset \sigma_A(T)$ (因为若 $\lambda\mathrm{id}_V - T$ 在 $\End(V)$ 中不可逆, 则它当然在子代数中也不可逆).
考虑 $T$ (看作 $A$ 的元素) 的 Gelfand 变换, 这是连续映射 $\hat{T} \colon \sigma(A) \to \CC$. 让我们证明 $\hat{T}$ 单射. 考虑 $\phi_1, \phi_2 \in \sigma(A)$, 即 $A$ 上的两个乘法泛函, 满足 $\hat{T}(\phi_1) = \hat{T}(\phi_2)$, 即 $\phi_1(T) = \phi_2(T)$. 我们已见 Gelfand 表示是 $*$-同态, 所以
\[
\widehat{T^*}(\phi_1) = \overline{\hat{T}(\phi_1)} = \overline{\hat{T}(\phi_2)} = \widehat{T^*}(\phi_2),
\]
即 $\phi_1(T^*) = \phi_2(T^*)$. 乘法泛函 $\phi_1$ 和 $\phi_2$ 在 $e$, $T$, $T^*$ 上相等, 且它们连续, 所以它们在整个 $A$ 上相等, 这正是我们想要的.
现在回忆 $\sigma(A)$ 紧 Hausdorff, 因为 $A$ 含幺. 所以 $\hat{T}$ 诱导从 $\sigma(A)$ 到其在 $\CC$ 中的像的同胚, 由推论 \ref{cor:II.2.11} 其像就是 $\sigma_A(T)$. 所以与 $\hat{T}$ 复合给出等距 $*$-同构 $\Psi \colon \mathcal{C}(\sigma_A(T)) \to \mathcal{C}(\sigma(A))$.
回忆我们还有 $A$ 的 Gelfand 表示, 它是等距 $*$-同构 $\Gamma \colon A \to \mathcal{C}(\sigma(A))$. 所以令 $\Phi = \Gamma^{-1} \circ \Psi$ 得到等距 $*$-同构 $\Phi \colon \mathcal{C}(\sigma_A(T)) \to A$.
让我们证明 $\Phi(\iota) = T$. 由于 $\Gamma \colon A \to \mathcal{C}(\sigma(A))$ 是同构, 只需验证 $\hat{T} = \widehat{\Phi(\iota)}$, 即 $\hat{T} = \Psi(\iota)$. 设 $\phi \in \sigma(A)$. 则 $\Psi(\iota)(\phi) = \iota(\hat{T}(\phi)) = \hat{T}(\phi)$, 如所愿.
最后, 我们证明包含 $\sigma(T) \subset \sigma_A(T)$ 是等式. 设 $\lambda \in \sigma_A(T)$, 假设 $\lambda \notin \sigma(T)$. 设 $\varepsilon > 0$, 选取 $f \in \mathcal{C}(\sigma_A(T))$ 使得 $\|f\|_{\infty} = 1$, $f(\lambda) = 1$, 且若 $|\lambda - \mu| \geq \varepsilon > 0$ 则 $f(\mu) = 0$. 设 $U = \Phi(f) \in A$, 则 $\|U\|_{\mathrm{op}} = \|f\|_{\infty} = 1$. 注意 $\Phi(1) = \mathrm{id}_V$ (其中 $1$ 是取值 $1$ 的常值函数), 因为 $\Phi$ 是代数同构. 所以 $T - \lambda\mathrm{id}_V = \Phi(\iota - \lambda)$, 且 $(T - \lambda\mathrm{id}_V) \circ U = \Phi((\iota - \lambda)f)$. 由于 $\Phi$ 等距, 这蕴含
\[
\|(T - \lambda\mathrm{id}_V) \circ U\|_{\mathrm{op}} = \|(\iota - \lambda)f\|_{\infty} \leq \varepsilon
\]
(因为 $f$ 被 $1$ 控制, 且若 $|\lambda - \mu| \geq \varepsilon$ 则 $f(\mu) = 0$). 另一方面, 由于 $\lambda \notin \sigma(T)$, 算子 $T - \lambda\mathrm{id}_V$ 在 $\End(V)$ 中可逆, 所以
\[
1 = \|f\|_{\infty} = \|U\|_{\mathrm{op}} = \|(T - \lambda\mathrm{id}_V)^{-1}(T - \lambda\mathrm{id}_V)U\|_{\mathrm{op}} \leq \varepsilon\|(T - \lambda\mathrm{id}_V)^{-1}\|_{\mathrm{op}}.
\]
这对每个 $\varepsilon > 0$ 都成立, 所以它蕴含 $1 = 0$, 矛盾.
\end{proof}

\begin{corollary}[谱为单点集的条件]
\label{cor:II.4.2}
设 $V$ 是希尔伯特空间, $T \in \End(V)$ 正规. 则下列条件等价:
\begin{enumerate}
\item[(i)] $\sigma(T)$ 是单点集;
\item[(ii)] $T \in \CC\,\mathrm{id}_V$;
\item[(iii)] $A_T = \CC\,\mathrm{id}_V$.
\end{enumerate}
\end{corollary}

\begin{proof}
(i)$\Rightarrow$(ii). 若 $\sigma(T) = \{\lambda\}$, 则 $\iota$ 是 $\mathcal{C}(\sigma(T))$ 的单位元的 $\lambda$ 倍, 所以 $T = \Phi(\iota) = \lambda\mathrm{id}_V$.
(ii)$\Rightarrow$(iii). 若 $T \in \CC\,\mathrm{id}_V$, 则 $\CC\,\mathrm{id}_V$ 是包含 $T$ 和 $T^*$ 的 $A$ 的闭含幺子代数, 所以它等于 $A_T$.
(ii)$\Rightarrow$(iii). 假设 $A_T = \CC\,\mathrm{id}_V$. 设 $\lambda, \mu \in \sigma(T)$. 若 $\lambda \neq \mu$, 则可以找到 $f_1, f_2 \in \mathcal{C}(\sigma(T))$ 使得 $f_1(\lambda) = 1$, $f_2(\mu) = 1$ 且 $f_1f_2 = 0$. 但此时 $\Phi(f_1)\Phi(f_2) = 0$ 且 $\Phi(f_1), \Phi(f_2) \neq 0$, 这与 $\CC$ 是整环矛盾.
\end{proof}

\begin{definition}[\index{中心化子}中心化子]
\label{def:II.4.3}
若 $A$ 是 $\CC$-代数, $E \subset A$ 是子集, 令 $Z_A(E) = \{x \in A \mid \forall y \in E,\, xy = yx\}$. 这称为 $E$ 在 $A$ 中的中心化子 (centralizer).
\end{definition}

容易看出中心化子总是 $A$ 的子代数.

\begin{corollary}[不可约作用下的中心化子]
\label{cor:II.4.4}
设 $V$ 是希尔伯特空间, $E$ 是 $\End(V)$ 的满足 $E^* = E$ 的子集. 假设 $V$ 的被 $E$ 的所有元素稳定的唯一闭子空间是 $\{0\}$ 和 $V$. 则 $Z_{\End(V)}(E) = \CC\,\mathrm{id}_V$.
\end{corollary}

\begin{proof}
设 $A = Z_{\End(V)}(E)$. 它是 $\End(V)$ 的闭子代数. 我们证明 $A$ 对 $*$ 稳定: 若 $T \in A$, 则对每个 $U \in E$, 有 $U^* \in E$, 所以
\[
T^* \circ U = (U^* \circ T)^* = (T \circ U^*)^* = U \circ T^*,
\]
所以 $T^* \in A$. 特别地, 子代数 $A$ 由其正规元生成; 事实上, 对每个 $T \in A$, $T = \frac{1}{2}((T + T^*) + (T - T^*))$, 且 $T + T^*$ 和 $T - T^*$ 都正规.

回忆我们要证明 $A$ 的每个元素都在 $\CC\,\mathrm{id}_V$ 中; 由刚才证明的, 只需对 $A$ 的正规元证明. 所以设 $T \in A$ 正规. 由推论 \ref{cor:II.4.2}, 只需证明 $T$ 的谱 $\sigma(T)$ 是单点集. 由谱定理 (定理 \ref{thm:II.4.1}), 有等距 $*$-同构 $\Phi \colon \mathcal{C}(\sigma(T)) \to A_T$, 其中 $A_T$ 是 $\End(V)$ 中由 $T$ 和 $T^*$ 生成的含幺子代数的闭包, 且 $\Phi$ 把嵌入 $\iota \colon \sigma(T) \hookrightarrow \CC$ 送到 $T$. 注意 $A_T \subset A$. 现在假设 $\sigma(T)$ 不是单点集. 则可以找到两个非零函数 $f_1, f_2 \in \mathcal{C}(\sigma(T))$ 使得 $f_1f_2 = 0$, 且 $\Phi(f_1), \Phi(f_2)$ 是满足 $\Phi(f_1)\Phi(f_2) = 0$ 的 $\End_G(V)$ 的非零元素. 设 $W = \Im(\Phi(f_2))$; 则 $W \neq \{0\}$, 因为 $\Phi(f_2)$ 非零. 又, 由于 $\Phi(f_2)$ 与 $E$ 的每个元素交换, 子空间 $W$ 被 $E$ 的所有元素稳定, 所以由假设 $W = V$. 但又有 $\Phi(f_1)(W) = 0$, 因为 $\Phi(f_1)\Phi(f_2) = 0$, 所以 $\Phi(f_1) = 0$, 这与 $f_1, f_2$ 的选取矛盾. 所以 $\sigma(T)$ 是单点集, 我们完成了.
\end{proof}

% ============ II.5 习题 (原书 pp.133-144) ============
\section{习题}

\begin{exercise}[$\End(V)$ 的谱为空]
\label{ex:II.5.1}
\begin{enumerate}
\item[(a)] 设 $V$ 是有限维 $\CC$-向量空间且 $\dim_{\CC}(V) \geq 2$. 证明 $\sigma(\End(V)) = \emptyset$.
\item[(b)] 设 $V$ 是无穷维希尔伯特空间. 证明 $\sigma(\End(V)) = \emptyset$.
\end{enumerate}
(提示: 看幂零自同态.)
\end{exercise}

\begin{solution}
(a). 可以假设 $V = \CC^n$, 所以 $\End(V) = M_n(\CC)$. 设 $\phi \colon M_n(\CC) \to \CC$ 是乘法线性泛函. 我们要证明 $\phi = 0$. 设 $(E_{ij})_{1 \leq i,j \leq n}$ 是 $M_n(\CC)$ 的典范基 (即 $E_{ij}$ 是除 $(i,j)$ 处为 $1$ 外全为零的矩阵). 则 $E_{ij}E_{kl}$ 在 $j \neq k$ 时等于 $0$, 且 $E_{ij}E_{jl} = E_{il}$. 特别地, 若 $i \neq j$, 则 $E_{ij}^2 = 0$, 所以 $0 = \phi(E_{ij}^2) = \phi(E_{ij})^2$, 且 $\phi(E_{ij}) = 0$. 又, 对每个 $i \in \{1, \ldots, n\}$, 若选取满足 $j \neq i$ 的 $j$ (因为 $n \geq 2$ 这可能), 则 $E_{ii} = E_{ij}E_{ji}$, 所以 $\phi(E_{ii}) = \phi(E_{ij})\phi(E_{ji}) = 0$. 总之, 我们证明了 $\phi$ 在 $M_n(\CC)$ 的基上为零, 所以 $\phi = 0$.

(b). 设 $\phi \colon \End(V) \to \CC$ 是乘法线性泛函. 同 (a), 关键是注意若 $T \in \End(V)$ 满足 $T^2 = 0$, 则 $\phi(T)^2 = 0$, 所以 $\phi(T) = 0$.
现在选取两个闭子空间 $V_1$ 和 $V_2$ 使得 $V = V_1 \oplus V_2$ 且 $V_1$ 与 $V_2$ 同构. (因为 $V$ 无穷维这可能. 例如, 选取 $V$ 的 Hilbert 基 $(e_i)_{i \in I}$. 由于 $I$ 无穷, 可以找到 $I_1, I_2 \subset I$ 使得 $I = I_1 \sqcup I_2$ 且 $I_1$ 与 $I_2$ 之间存在双射. 对 $r = 1, 2$ 取 $V_r = \bigoplus_{i \in I_r}\CC e_i$.) 选取同构 $U_1 \colon V_1 \to V_2$ 和 $U_2 \colon V_2 \to V_1$. 设 $T_1 \in \End(V)$ 由 $T_1(v + w) = U_1(v)$ (若 $v \in V_1$, $w \in V_2$) 定义, $T_2 \in \End(V)$ 由 $T_2(v + w) = U_2(w)$ (若 $v \in V_1$, $w \in V_2$) 定义. 则 $T_1^2 = T_2^2 = 0$, 所以 $\phi(T_1) = \phi(T_2) = 0$, 且也有 $\phi(T_1 + T_2) = 0$. 但 $T := T_1 + T_2$ 是 $V$ 的自同构, 所以对每个 $T' \in \End(V)$, $T' = T(T^{-1}T')$, 从而 $\phi(T') = \phi(T)\phi(T^{-1}T') = 0$.
\end{solution}

\begin{exercise}[有限维谱定理]
\label{ex:II.5.2}
设 $V$ 是有限维希尔伯特空间. 本题的目标是把讲义中的谱定理 (定理 \ref{thm:II.4.1}) 与通常的有限维谱定理 (它说 $V$ 的正规自同态在标准正交基下可对角化) 联系起来.
回忆若 $R$ 是交换环, 若存在整数 $n \geq 1$ 使得 $x^n = 0$ 则称 $x \in R$ 幂零 (nilpotent); 若 $R$ 的唯一幂零元是 $0$, 则称 $R$ 约化 (reduced).
\begin{enumerate}
\item[(a)] 证明通常的有限维谱定理 (如上所述) 蕴含 $V$ 的定理 \ref{thm:II.4.1}.
\item[(b)] 设 $T \in \End(V)$, $A$ 是 $\End(V)$ 中由 $T$ 生成的含幺子代数 (即 $T$ 的多项式的空间). 证明 $T$ 可对角化当且仅当 $A$ 约化.
\item[(c)] 设 $A$ 是 $\End(V)$ 的交换含幺子代数. 若 $A$ 约化, 证明存在 $V$ 的子空间 $V_1, \ldots, V_r$ (除排序外唯一确定) 使得 $V = \bigoplus_{i=1}^r V_i$, 且
\[
A = \{T \in \End(V) \mid \forall i \in \{1, \ldots, r\},\, T(V_i) \subset V_i \text{ 且 } T|_{V_i} \in \CC \cdot \mathrm{id}_{V_i}\}.
\]
\item[(d)] 设 $A$ 如 (c). 证明 $A$ 对映射 $T \longmapsto T^*$ 稳定当且仅当 $V_i$ 两两正交.
\item[(e)] 证明定理 \ref{thm:II.4.1} 蕴含通常的有限维谱定理 (如上所述).
\end{enumerate}
\end{exercise}

\begin{solution}
(a). 设 $T \in \End(V)$ 是正规自同态. 由有限维谱定理, 可以找到 $V$ 的标准正交基 $(e_1, \ldots, e_n)$ 和 $\lambda_1, \ldots, \lambda_n \in \CC$ 使得对每个 $i \in \{1, \ldots, n\}$ 有 $T(e_i) = \lambda_ie_i$. 由于基标准正交, 对每个 $i \in \{1, \ldots, n\}$ 也有 $T^*(e_i) = \bar{\lambda}_ie_i$. 重新排列 $e_i$ 后, 还可以假设存在 $1 \leq n_0 \leq \cdots \leq n_r = n+1$ 使得: 若存在 $s \in \{0, \ldots, r-1\}$ 满足 $n_s \leq i, j \leq n_{s+1} - 1$ 则 $\lambda_i = \lambda_j$, 否则 $\lambda_i \neq \lambda_j$.
特别地, 可以假设 $V = \CC^n$, 且 $T$ 是对角线元素为 $\lambda_1, \ldots, \lambda_n$ 的对角矩阵, 对 $\lambda_i$ 的条件同上. 我断言 $A_T$ 是 $M_n(\CC)$ 中满足如下条件的对角矩阵 $x_1, \ldots, x_n$ 的子代数: 若存在 $s \in \{0, \ldots, r-1\}$ 满足 $n_s \leq i, j \leq n_{s+1} - 1$ 则 $x_i = x_j$. 首先, 这确实定义 $M_n(\CC)$ 的子代数. 也显然 $A_T$ 的每个矩阵都具有此形式, 因为 $A_T$ (作为代数) 由 $I_n$, $T$ 和 $T^*$ 生成, 且这三个矩阵都满足定义 $A_T$ 的条件. 最后, 设 $X \in A_T$, $x_1, \ldots, x_n$ 是它的对角线元素. 由 Lagrange 插值, 存在多项式 $P \in \CC[t]$ 使得对 $s \in \{0, \ldots, r-1\}$ 有 $P(\lambda_{n_s}) = x_{n_s}$, 然后 $P(T)$ 是对角线元素为 $x_1, \ldots, x_n$ 的对角矩阵, 即 $X$.

(b). 设 $P \in \CC[t]$ 是 $T$ 的最小多项式. 则由最小多项式的定义, 映射 $\CC[t] \to M_n(\CC)$, $f(t) \longmapsto f(T)$ 是以 $A$ 为像、以 $P$ 生成的理想为核的 $\CC$-代数同态. 所以 $A \simeq \CC[t]/(P)$. 若把 $P(t)$ 写成 $P(t) = \prod_{i=1}^r(t - a_i)^{n_i}$ (其中 $a_1, \ldots, a_r \in \CC$ 两两不同, $n_1, \ldots, n_r \geq 1$), 则由中国剩余定理 $A \simeq \prod_{i=1}^n \CC[t]/(t - a_i)^{n_i}$. 所以 $A$ 约化当且仅当所有 $n_i$ 都等于 $1$, 即当且仅当 $P$ 只有单根, 这等价于 $T$ 可对角化.

(c). 若 $T \in A$, 则 $\End(V)$ 中由 $T$ 生成的含幺子代数包含于 $A$, 特别地它约化; 由 (b), 这蕴含 $T$ 可对角化. 由于 $A$ 由有限多个元素生成 (因为它是有限维 $\CC$-向量空间), 且这些元素可对角化且两两交换, 可以找到基 $(e_1, \ldots, e_n)$ 使得 $A$ 的每个元素在其中对角. 对 $i \in \{1, \ldots, n\}$, 定义 $\phi_i \colon A \to \CC$ 为 $T(e_i) = \phi_i(T)e_i$ (对 $T \in A$). 则 $\phi_1, \ldots, \phi_n$ 是 $A$ 上的乘法泛函. 重新排列 $e_i$ 后, 可以假设存在 $1 \leq n_0 \leq \cdots \leq n_r = n+1$ 使得: 若存在 $s \in \{0, \ldots, r-1\}$ 满足 $n_s \leq i, j \leq n_{s+1} - 1$ 则 $\phi_i = \phi_j$, 否则 $\phi_i \neq \phi_j$. 注意所有 $\phi_i$ 非零 (因为它们把 $I_n$ 送到 $1$), 所以它们满射. 我断言 $\phi_{n_0}, \phi_{n_1}, \ldots, \phi_{n_{r-1}}$ 线性无关 (作为函数 $A \to \CC$). 这是经典结果, 但让我们快速证明. 假设不真, 选取 $\phi_{n_i}$ 的非平凡线性相关关系 $\sum_{i=0}^{r-1}a_i\phi_{n_i} = 0$ ($a_i \in \CC$), 使得非零 $a_i$ 的个数最小. 至少有 $2$ 个非零 $a_i$, 所以重新排序后可以假设 $a_0, a_1 \neq 0$. 选取 $x_0 \in A$ 使得 $\phi_{n_1}(x_0) \neq \phi_{n_1}(x_0)$. 则对每个 $x \in A$,
\[
0 = \phi_{n_0}(x_0)\sum_{i=0}^{r-1}a_i\phi_{n_i}(x) - \sum_{i=0}^{r-1}a_i\phi_{n_i}(x_0x) = \sum_{i=1}^{r-1}a_i(\phi_{n_0}(x_0) - \phi_{n_i}(x_0))\phi_{n_i}(x).
\]
所以 $\sum_{i=1}^{r-1}a_i(\phi_{n_0}(x_0) - \phi_{n_i}(x_0))\phi_{n_i} = 0$, 其中 $a_1(\phi_{n_1}(x_0) - \phi_{n_1}(x_0)) \neq 0$. 所以我们有 $\phi_{n_i}$ 的另一个非平凡线性相关关系, 且它的非零系数比第一个少, 矛盾.
既然知道 $\phi_{n_0}, \phi_{n_1}, \ldots, \phi_{n_{r-1}}$ 线性无关, 也知道 $(\phi_{n_0}, \phi_{n_1}, \ldots, \phi_{n_{r-1}}) \colon A \to \CC^r$ 满射. 对 $i \in \{1, \ldots, r\}$, 设 $V_i = \operatorname{Span}(e_{n_{i-1}}, \ldots, e_{-1+n_i}) \subset V$. 则若 $T \in A$, $T$ 在每个 $V_i$ 上按恒等映射的倍数作用, 且 $(\phi_{n_0}, \phi_{n_1}, \ldots, \phi_{n_{r-1}}) \colon A \to \CC^r$ 的满射性蕴含其逆也真. (若 $(a_1, \ldots, a_r) \in \CC^r$, 选取 $T \in A$ 使得对 $0 \leq i \leq r-1$ 有 $\phi_{n_i} = a_{i+1}$. 则 $T$ 在每个 $V_i$ 上按 $\phi_{n_{i+1}}(T) = a_i$ 倍作用.)
最后, 让我们证明 $V_1, \ldots, V_r$ 唯一确定. 设 $V = V'_1 \oplus \cdots \oplus V'_s$ 是满足同样性质的另一个分解. 选取两两不同的 $a_1, \ldots, a_r \in \CC$, 设 $T \in A$ 满足对每个 $i$ 有 $T|_{V_i} = a_i\mathrm{id}_{V_i}$. 则 $V_1, \ldots, V_r$ 是 $T$ 的特征空间, 且 $T$ 在每个 $V'_j$ 上按恒等映射的倍数作用, 所以必存在 $\{1, \ldots, r\}$ 的分拆 $I_1, \ldots, I_s$ 使得对每个 $j \in \{1, \ldots, s\}$ 有 $V'_j = \bigoplus_{i \in I_j}V_i$. 但 $V_i$ 与 $V'_j$ 的角色对称, 所以交换 $V_i$ 与 $V'_j$ 后有类似的性质. 这蕴含 $r = s$ 且 $V'_1, \ldots, V'_r$ 除排序外等于 $V_1, \ldots, V_r$.

(d). 对 $i \in \{1, \ldots, r\}$, 定义线性自同态 $\pi_i \colon V \to V$ 为: 若 $v_j \in V_j$ ($j \in \{1, \ldots, r\}$), 则 $\pi_i(v_1 \cdots + v_r) = v_i$. 则 $\Im(\pi_i) = V_i$, $\Ker(\pi_i) = \bigoplus_{j \neq i}V_j$, 且 $\pi_1 + \cdots + \pi_r = \mathrm{id}_V$. 注意 $A$ 正好是 $\End(V)$ 的子代数 $\{\sum_{i=1}^r\lambda_i\pi_i \mid \lambda_1, \ldots, \lambda_r \in \CC\}$.
若 $V_1, \ldots, V_r$ 两两正交, 则 $\pi_1, \ldots, \pi_r$ 是正交投影, 所以它们自伴, 从而 $A$ 对 $T \longmapsto T^*$ 稳定.
反之, 假设 $A$ 对 $T \longmapsto T^*$ 稳定. 若 $v \in V$, $w \in V_1^{\perp}$, 则
\[
0 = \langle \pi_1(v), w \rangle = \langle v, \pi_1^*(w) \rangle.
\]
这蕴含 $V_1^{\perp} \subset \Ker(\pi_1^*)$. 由于 $\operatorname{rk}(\pi_1^*) = \operatorname{rk}(\pi_1) = \dim(V_1)$, 实际上有 $\Ker(\pi_1^*) = V_1^{\perp}$. 但 $\pi_1^* \in A$, 所以 $\pi_1^*$ 的每个特征空间都是 $V_i$ 的和, 因此存在 $I \subset \{1, \ldots, r\}$ 使得 $\Ker(\pi_1^*) = \bigoplus_{i \in I}V_i$. 由于 $V_1^{\perp} \cap V_1 = \{0\}$, 集合 $I$ 不能包含 $1$. 但 $\Ker(\pi_1^*)$ 的维数只能是 $\dim(V) - n_1$, 所以唯一可能是 $I = \{2, \ldots, r\}$. 最后, 我们证明了
\[
V_1^{\perp} = \Ker(\pi_1^*) = V_2 \oplus \cdots \oplus V_r.
\]
对其他 $\pi_i$ 重复此过程, 看到对每个 $i \in \{1, \ldots, r\}$,
\[
V_i^{\perp} = \bigoplus_{j \neq i}V_j.
\]

(e). 设 $T \in \End(V)$, 设 $\Phi \colon \mathcal{C}(\sigma(T)) \to A_T$ 如定理 \ref{thm:II.4.1}. 特别地, $A_T$ 是 $\End(V)$ 的交换约化子代数 (因为 $\mathcal{C}(\sigma(T))$ 约化), 且它 (由定义) 对 $*$ 稳定, 所以由 (c) 和 (d), 有 $V$ 分解成两两正交子空间的直和 $V = V_1 \oplus \cdots \oplus V_r$, 使得 $A_T$ 的每个元素保持此分解并在每个 $V_i$ 上按纯量作用. 若为每个 $V_i$ 选取标准正交基并把它们合在一起, 将得到 $V$ 的标准正交基, 其中 $A_T$ 的每个元素对角. 现在只需回忆 $T \in A_T$.
\end{solution}

\begin{exercise}[$\mathcal{C}(X)$ 的谱]
\label{ex:II.5.3}
本题旨在不使用 II.3 和 II.4 节的任何结果来解.\footnote{与 Nullstellensatz 比较.}
设 $X$ 是局部紧 Hausdorff 拓扑空间. 设 $\tilde{X}$ 是 $X$ 的 Alexandroff 紧化. 这意味着 $\tilde{X} = X \cup \{\infty\}$, 且 $\tilde{X}$ 的开集是 $X$ 的开子集和形如 $(X - K) \cup \{\infty\}$ 的集合, 其中 $K$ 是 $X$ 的紧子集.
\begin{enumerate}
\item[(a)] 证明 $\tilde{X}$ 是紧 Hausdorff 拓扑空间, $X$ 在 $\tilde{X}$ 中开, 且 $X$ 在 $\tilde{X}$ 中稠密当且仅当 $X$ 不紧.
\item[(b)] 证明 $\mathcal{C}(\tilde{X})$ 同构于给 $\mathcal{C}_0(X)$ 添加单位元得到的 Banach $*$-代数. (别忘了比较拓扑.)
\item[(c)] 若 $X$ 紧, 证明 $\mathcal{C}(X)$ 的每个真理想都包含于某个理想 $\mathfrak{m}_x = \{f \in \mathcal{C}(X) \mid f(x) = 0\}$ ($x \in X$).
\item[(d)] 一般地, 证明映射 $X \to \sigma(\mathcal{C}_0(X))$, $x \longmapsto (\phi_x \colon f \longmapsto f(x))$ 是同胚.
\end{enumerate}
设 $A$ 是交换 Banach 代数. 若 $I$ 是 $A$ 的理想, 令
\[
V(I) = \{x \in \sigma(A) \mid \forall f \in A,\, \hat{f}(x) = 0\}.
\]
若 $N$ 是 $\sigma(A)$ 的子集, 令
\[
I(N) = \{f \in A \mid \forall x \in N,\, \hat{f}(x) = 0\}.
\]
\begin{enumerate}
\item[(a)] 假设 $X$ 紧. 证明对 $\mathcal{C}(X)$ 的每个闭理想 $I$ 和 $\sigma(\mathcal{C}(X)) \simeq X$ 的每个闭子集 $N$,
\[
I(V(I)) = I \quad \text{且} \quad V(I(N)) = N.
\]
\end{enumerate}
注意. 不假设 $X$ 紧时结果仍成立 (处处用 $\mathcal{C}_0(X)$).

\begin{solution}
(a). 首先证明定义确实给出 $\tilde{X}$ 上的拓扑. 设 $(U_i)_{i \in I}$ 是 $\tilde{X}$ 的开子集族. 则可以把 $I$ 写成 $I = I' \sqcup I''$, 其中对 $i \in I'$, $U_i \subset X$ 开; 对 $i \in I''$, $U_i = X - K_i$ 且 $K_i$ 紧. 有
\[
\bigcup_{i \in I}U_i = \Bigl(\bigcup_{i \in I'}U_i\Bigr) \cup \Bigl(X - \bigcap_{i \in I''}K_i\Bigr).
\]
若 $I''$ 空, 这是 $X$ 的开子集, 从而是 $\tilde{X}$ 的开子集. 否则, 这是 $\tilde{X}$ 的紧子集 $\bigcap_{i \in I''}K_i - \bigcup_{i \in I'}U_i$ 的余集, 所以它又是 $\tilde{X}$ 的开子集. 另一方面,
\[
\bigcap_{i \in I}U_i = \Bigl(\bigcap_{i \in I'}U_i\Bigr) \cap \Bigl(X - \bigcup_{i \in I''}K_i\Bigr).
\]
假设 $I$ 有限. 则若 $I'' = \emptyset$, 集合 $\bigcap_{i \in I}U_i$ 是 $X$ 的开子集 $\bigcap_{i \in I'}U_i$, 从而是 $\tilde{X}$ 的开子集. 否则, 它是 $\tilde{X}$ 的紧子集 $\bigcup_{i \in I''}K_i - \bigcap_{i \in I'}U_i$ 的余集, 所以它又是 $\tilde{X}$ 的开子集.
让我们证明 $\tilde{X}$ Hausdorff. 设 $x, y \in \tilde{X}$ 满足 $x \neq y$. 我们要找到 $x$ 和 $y$ 的不相交开邻域. 若 $x, y \in X$, 则存在 $X$ 的开子集 $U$ 和 $V$ 使得 $x \in U$, $y \in V$, $U \cap V = \emptyset$. 这些集合在 $\tilde{X}$ 中仍开, 所以我们完成了. 若 $x$ 或 $y$ 是 $\infty$, 可以假设 $x$ 是. 由于 $X$ 局部紧, 可以找到 $X$ 的紧子集 $K$ 和 $X$ 的开子集 $V$ 使得 $y \in V \subset K$. 则 $U := \tilde{X} - K$ 是包含 $x = \infty$ 的 $\tilde{X}$ 的开子集, 且 $U \cap V = \emptyset$.
让我们证明 $\tilde{X}$ 紧. 设 $(U_i)_{i \in I}$ 是 $\tilde{X}$ 的开子集族, 满足 $\tilde{X} = \bigcup_{i \in I}U_i$. 设 $i_0 \in I$ 使得 $\infty \in U_{i_0}$, 记 $K = X - U_{i_0}$. 这是 $X$ 的紧子集, 且被开子集 $U_i \cap X$ ($i \in I - \{i_0\}$) 覆盖. 所以存在 $I - \{i_0\}$ 的有限子集 $J$ 使得 $K \subset \bigcup_{i \in J}U_i$, 然后 $\tilde{X} = \bigcup_{i \in J \cup \{i_0\}}U_i$.
集合 $X$ 由 $\tilde{X}$ 的拓扑定义在 $\tilde{X}$ 中开.
假设 $X$ 不紧. 则若 $U$ 是 $\infty$ 在 $\tilde{X}$ 中的开邻域, $\tilde{X}$ 的紧子集 $X - U$ 不能等于 $X$, 这意味着 $U \cap X \neq \emptyset$. 所以 $\infty$ 在 $X$ 于 $\tilde{X}$ 中的闭包里. 反之, 假设 $X$ 紧. 则 $\{\infty\} = \tilde{X} - X$ 是 $\tilde{X}$ 的开子集, 所以 $\infty$ 是 $\tilde{X}$ 的孤立点.

(b). 设 $A$ 是给 $\mathcal{C}_0(X)$ 添加单位元得到的 Banach $*$-代数. 有 $A = \mathcal{C}_0(X) \oplus \CC e$, 其中 $\|f + \lambda e\| = \|f\|_{\infty} + |\lambda|$, $(f + \lambda e)^* = \bar{f} + \bar{\lambda}e$ (对 $f \in \mathcal{C}_0(X)$, $\lambda \in \CC$).
注意可以把每个 $f \in \mathcal{C}_0(X)$ 通过令 $f(\infty) = 0$ 延拓成 $\tilde{X}$ 上的连续函数. ($f$ 在无穷远处为零这一条件正好说明延拓后的函数连续, 由 $\tilde{X}$ 上拓扑的定义.) 这给出单射 $\CC$-代数映射 $\mathcal{C}_0(X) \to \mathcal{C}(\tilde{X})$. 所以得到映射 $\alpha \colon A \to \mathcal{C}(\tilde{X})$, $f + \lambda e \longmapsto \tilde{f} + \lambda$, 其中第二个``$\lambda$''是 $\tilde{X}$ 上的常值函数. 由 $A$ 上乘法的定义, 这个 $\alpha$ 是 $\CC$-代数的同态; 由 $A$ 上 $*$ 的定义, 它是 $*$-同态. 又, $\alpha$ 有界, 因为若 $f \in \mathcal{C}_0(X)$, $\lambda \in \CC$, 则
\[
\|\tilde{f} + \lambda\|_{\infty} \leq \|f\|_{\infty} + |\lambda| = \|f + \lambda e\|.
\]
最后, 注意 $\alpha$ 满射, 因为它有逆, 把 $f \in \mathcal{C}(\tilde{X})$ 送到 $(f|_X - \lambda) + \lambda e$. 由开映射定理 (也称为 Banach-Schauder 定理), $\alpha$ 的逆也有界, 所以 $\alpha$ 是同胚.

(c). 设 $I$ 是 $\mathcal{C}(X)$ 的理想, 假设 $I$ 不包含于任何 $\mathfrak{m}_x$. 则对每个 $x \in X$, 可以找到 $f_x \in I$ 使得 $f_x(x) \neq 0$; 由于 $f_x$ 连续, 还可以找到 $x$ 的开邻域 $U_x$ 使得对每个 $y \in U_x$ 有 $f_x(y) \neq 0$. 有 $X = \bigcup_{x \in X}U_x$, 且 $X$ 紧, 所以存在 $x_1, \ldots, x_n \in X$ 使得 $X = \bigcup_{i=1}^n U_{x_i}$. 设 $f = \sum_{i=1}^n|f_{x_i}|^2 = \sum_{i=1}^n f_{x_i}\overline{f_{x_i}}$. 则 $f \in I$ (因为 $I$ 是理想), 且 $f$ 在 $X$ 上不消失; 事实上, 若 $x \in X$, 可以找到满足 $x \in U_{x_i}$ 的 $i \in \{1, \ldots, n\}$, 然后 $f(x) \geq |f_{x_i}(x)|^2 > 0$. 所以函数 $g \colon x \longmapsto f(x)^{-1}$ 存在且在 $X$ 上连续, 且有 $gf = 1$, 这蕴含 $1 \in I$, 从而 $I = \mathcal{C}(X)$.

(d). 记此映射为 $\alpha$. 首先证明 $\alpha$ 单射. 若 $x, y \in X$ 满足 $x \neq y$, 则由 Urysohn 引理存在 $f \in \mathcal{C}_0(X)$ 使得 $f(x) \neq f(y)$, 所以 $\phi_x \neq \phi_y$.
让我们证明 $\alpha$ 满射. 设 $\phi \colon \mathcal{C}_0(X) \to \CC$ 是乘法泛函. 可以把它延拓成 $\mathcal{C}_0(X)^{\tilde{e}}$ 上的乘法泛函 $\tilde{\phi}$, 且我们在 (b) 中看到 $\mathcal{C}_0(X)^{\tilde{e}}$ 同构于 $\mathcal{C}(\tilde{X})$. 设 $I = \Ker(\tilde{\phi})$. 这是 $\mathcal{C}(\tilde{X})$ 的极大理想, 所以由 (c), 存在 $x \in \tilde{X}$ 使得 $I \subset \mathfrak{m}_x$, 且由于 $I$ 极大必有 $I = \mathfrak{m}_x$. 又注意 (b) 中构造的同构 $\mathcal{C}_0(X)^{\tilde{e}} \simeq \mathcal{C}(\tilde{X})$ 把 $\mathcal{C}_0(X)$ 等同于 $\mathfrak{m}_{\infty}$. 所以, 由于 $\phi$ 在 $\mathcal{C}_0(X)$ 上不为零, 不能有 $x = \infty$, 故 $x \in X$, 且有
\[
\Ker(\phi) = \{f \in \mathcal{C}_0(X) \mid f(x) = 0\} = \Ker(\phi_x).
\]
像定理 \ref{thm:II.2.10} 的证明那样, 这容易蕴含 $\phi = \phi_x$.
映射 $\alpha$ 由 $\sigma(\mathcal{C}_0(X))$ 上拓扑的定义连续. 若 $X$ 紧, 这蕴含 $\alpha$ 是同胚. 一般地, $\tilde{X}$ 的 Alexandroff 紧化的 $\alpha$ 版本因 $\tilde{X}$ 紧而是同胚, 且它到 $X$ 的限制是 $\alpha$ (若像 (b) 那样把 $\mathcal{C}_0(X)$ 等同于 $\mathcal{C}(\tilde{X})$ 的子代数), 所以 $\alpha$ 开, 我们完成了.

(e). 注意若 $N$ 是 $X$ 的闭子集, 则 $I(N) = \bigcap_{x \in N}\mathfrak{m}_x$, 所以 $I(N)$ 是 $\mathcal{C}(X)$ 的理想.
设 $I$ 是 $\mathcal{C}(X)$ 的闭理想, 设 $N = N(I)$. 对每个 $x \in N$ 和每个 $f \in I$, 由 $N(I)$ 的定义有 $f(x) = 0$. 所以 $I \subset \bigcap_{x \in N}\mathfrak{m}_x = I(N)$. 反之, 设 $f \in \bigcap_{x \in N}\mathfrak{m}_x$; 我们要证明 $f \in I$. 由假设, 对每个 $x \in N$ 有 $f(x) = 0$, 所以 $\operatorname{supp}(f) \cap N = \emptyset$. 对每个 $y \in \operatorname{supp}(f)$, 选取满足 $f_y(y) \neq 0$ 的 $f_y \in I$; 由于 $f_y$ 连续, 可以找到 $X$ 的开子集 $U_y \ni y$ 使得对每个 $z \in U_y$ 有 $f_y(z) \neq 0$. 有 $\operatorname{supp}(f) \subset \bigcup_{y \in \operatorname{supp}(f)}U_y$, 且 $\operatorname{supp}(f)$ 紧, 所以可以找到 $y_1, \ldots, y_n \in \operatorname{supp}(f)$ 使得 $\operatorname{supp}(f) \subset \bigcup_{i=1}^n U_{y_i}$. 设 $g = \sum_{i=1}^n|f_{y_i}|^2$. 则 $g \in I$, 且对每个 $y \in \operatorname{supp}(f)$ 有 $g(y) > 0$. 定义函数 $h \colon X \to \CC$:
\[
h(x) = \begin{cases} f(x)g(x)^{-1} & \text{若 } x \in \operatorname{supp}(f), \\ 0 & \text{否则}. \end{cases}
\]
设 $U = \{x \in X \mid g(x) \neq 0\}$, $V = X - \operatorname{supp}(f)$. 则 $U$ 和 $V$ 是 $X$ 的开子集且 $X = U \cup V$. 在 $U$ 上, 函数 $h$ 等于 $fg^{-1}$, 所以连续; 在 $V$ 上, 它等于 $0$, 所以也连续. 所以 $h \in \mathcal{C}(X)$, 且由 $h$ 的定义有 $f = gh$. 由于 $g \in I$, 这证明 $f \in I$, 如所愿.
现在设 $N$ 是 $X$ 的闭子集, $I = I(N)$. 对每个 $x \in N$ 和每个 $f \in I$, 由 $I(N)$ 的定义有 $f(x) = 0$, 所以 $N \subset V(I)$. 反之, 若 $x \notin N$, 则由 Urysohn 引理可以找到 $f \in \mathcal{C}(X)$ 使得 $f|_N = 0$ 且 $f(x) \neq 0$. 则由 $I(N)$ 的定义 $f \in I$, 所以 $x \notin V(I)$.
\end{solution}
\end{exercise}

\begin{exercise}[$\ell^1(\Z)$ 的谱]
\label{ex:II.5.4}
考虑 Banach $*$-代数 $\ell^1(\Z)$ (即离散群 $G = \Z$ 的 $L^1(G)$, 带卷积乘积和课堂上定义的对合). 把 $\ell^1(\Z)$ 的元素写成 $\CC^{\Z}$ 中的序列 $a = (a_n)_{n \in \Z}$.
\begin{enumerate}
\item[(a)] 证明 $\ell^1(\Z)$ 不是 $\mathcal{C}^*$-代数.
\item[(b)] 证明存在同胚 $\sigma(\ell^1(\Z)) \xrightarrow{\sim} S^1$ 使得 $a = (a_n)_{n \in \Z}$ 的 Gelfand 变换是函数 $S^1 \to \CC$, $e^{i\theta} \longmapsto \sum_{n=-\infty}^{+\infty}a_ne^{in\theta}$.\footnote{这意味着 Gelfand 变换是 $*$-同态, 即 Banach $*$-代数 $L^1(G)$ 对称, 即使它不是 $\mathcal{C}^*$-代数.}
\item[(c)] 更一般地, 若 $G$ 是交换局部紧群, 证明映射 $\widehat{G} \to \sigma(L^1(G))$, $\chi \longmapsto \bigl(f \longmapsto \int_G f(x)\chi(x)\,\dd x\bigr)$ 是同胚. (提示: $L^1(G)$ 的对偶是什么?)
\end{enumerate}
\end{exercise}

\begin{solution}
(a). 设 $a = (a_n)_{n \in \Z}$. 则 $a^* = (a_{-n})_{n \in \Z}$ (回忆 $\Z$ 幺模, 因为它交换 (或因为离散)). 设 $b = a^* * a$. 则对每个 $n \in \Z$,
\[
b_n = \sum_{m \in \Z}a^*_ma_{n-m} = \sum_{m \in \Z}\overline{a_{-m}}a_{n-m}.
\]
取由 $a_0 = i$, $a_1 = 1$, $a_2 = i$, 且对 $n \in \Z - \{0, 1, 2\}$ 有 $a_n = 0$ 定义的 $a$. 则 $a^*_0 = -i$, $a^*_{-1} = 1$, $a^*_2 = -i$, 且对 $n \in \Z - \{-2, -1, 0\}$ 有 $a^*_n = 0$. 所以对 $n \notin \{-2, -1, 0, 1, 2\}$ 有 $b_n = 0$, 且
\[
b_{-2} = a^*_{-2}a_0 = 1,
\]
\[
b_{-1} = a^*_{-1}a_0 + a^*_{-2}a_1 = i - i = 0,
\]
\[
b_0 = a^*_{-2}a_2 + a^*_{-1}a_1 + a^*_0a_0 = 3,
\]
\[
b_1 = a^*_0a_1 + a^*_{-1}a_2 = -i + i = 0,
\]
且
\[
b_2 = a^*_0a_2 = 1.
\]
所以 $|b|_1 = 5 \neq |a|_1^2 = 9$.

(b). 由习题 \ref{ex:I.5.4.1}(d) 知道 $\widehat{\Z} \simeq S^1$, 所以由 (c) 得到同胚 $S^1 = \widehat{\Z} \to \sigma(\ell^1(\Z))$. 展开公式, 看到它把 $z \in S^1$ 送到 $\ell^1(\Z)$ 上的乘法泛函 $a = (a_n)_{n \in \Z} \longmapsto \sum_{n \in \Z}a_nz^n$, 这正是我们想要的.

(c). 设 $G$ 是交换局部紧群. 设 $\phi \in \sigma(L^1(G))$. 我们要证明 $\phi$ 来自 $\widehat{G}$ 的元素 $\chi$. 由于 $\phi$ 是 $L^1(G)$ 上的连续线性泛函, 存在 $\chi \in L^{\infty}(G)$ 使得对每个 $f \in L^1(G)$ 有 $\phi(f) = \int_G f(x)\chi(x)\,\dd x$.
对每个 $f, g \in L^1(G)$,
\begin{align*}
\phi(f)\int_G g(y)\chi(y)\,\dd y &= \phi(f)\phi(g) \\
&= \phi(g * f) \\
&= \int_{G \times G}g(y)f(y^{-1}x)\chi(x)\,\dd x\,\dd y \\
&= \int_G g(y)\phi(L_yf)\,\dd y.
\end{align*}
由于这对每个 $g \in L^1(G)$ 成立, 函数 $\phi(f)\chi$ 和 $y \longmapsto \phi(L_yf)$ (都在 $L^{\infty}(G)$ 中) 几乎处处相等. 所以, 若选取满足 $\phi(f) \neq 0$ 的 $f \in L^1(G)$, 可以把 $\chi$ 换成 $y \longmapsto \phi(f)^{-1}\phi(L_yf)$. 由于函数 $\phi \colon L^1(G) \to \CC$ 和 $G \to L^1(G)$, $y \longmapsto L_yf$ 连续 (第二个由命题 \ref{prop:I.3.1.13}), 这个新的 $\chi$ 连续. 又, 对每个 $g \in L^1(G)$ 和每个 $y \in G$ 有 $\phi(g)\chi(y) = \phi(L_yg)$.
设 $x, y \in G$. 由于 $L_{xy}f = L_x(L_yf)$, 有
\[
\phi(L_{xy}f) = \chi(xy)\phi(f) = \chi(x)\phi(L_yf) = \chi(x)\chi(y)\phi(f),
\]
所以 $\chi(xy) = \chi(x)\chi(y)$. 所以 $\chi \in \widehat{G}$, 且我们已证明题设的映射 $\widehat{G} \to \sigma(L^1(G))$ 满射. 注意此映射也单射, 因为 $G$ 上的连续函数由它在 $L^1(G)$ 上定义的线性泛函确定. 又, $\sigma(L^1(G)) \subset L^{\infty}(G)$ 上的拓扑由定义是弱*拓扑, 且我们已在习题 \ref{ex:I.5.4.2}(a) 中看到这与 $\widehat{G}$ 上的紧收敛拓扑一致, 所以映射 $\widehat{G} \to \sigma(L^1(G))$ 是同胚.
\end{solution}

\begin{exercise}[$\mathcal{C}^*$-范数的唯一性]
\label{ex:II.5.5}
设 $A$ 是带对合反自同构 $*$ 的含幺 $\CC$-代数. 证明 $A$ 上至多有一个范数使 $A$ 成为 $\mathcal{C}^*$-代数.
\end{exercise}

\begin{solution}
设 $\|\cdot\|$ 是使 $A$ 成为 $\mathcal{C}^*$-代数的 $A$ 上的范数. 设 $x \in A$. 注意 $(x^*x)^* = x^*x$, 所以 $x^*x$ 正规. 由 $\mathcal{C}^*$-代数的定义和推论 \ref{cor:II.3.9},
\[
\|x\| = \|x^*x\|^{1/2} = \rho(x^*x)^{1/2}.
\]
但由定理 \ref{thm:II.1.1.3},
\[
\rho(x^*x) = \max\{|\lambda| \mid \lambda \in \CC,\, x^*x - \lambda e \notin A^{\times}\}.
\]
最后一个量只依赖于 $A$ 的代数结构和 $*$, 它确定 $\|x\|$.
\end{solution}

\begin{exercise}[$A^{\tilde{e}}$ 的 $\mathcal{C}^*$-范数]
\label{ex:II.5.6}
设 $A$ 是 $\mathcal{C}^*$-代数. 则 $A^{\tilde{e}}$ 是 Banach $*$-代数, 但对范数 $\|x + \lambda e\| = \|x\| + |\lambda|$ 它不总是 $\mathcal{C}^*$-代数. (见习题 \ref{ex:II.5.3}(b) 了解此现象的例子.)
我们在 $A^{\tilde{e}}$ 上定义新范数 $\|\cdot\|'$:
\[
\|x + \lambda e\|' = \sup\{\|xy + \lambda y\| \mid y \in A,\, \|y\| \leq 1\}.
\]
我们现在假设 $A$ 没有单位元且 $A \neq \{0\}$.
\begin{enumerate}
\item[(a)] 证明 $\|\cdot\|'$ 是 $A^{\tilde{e}}$ 上的次乘范数.
\item[(b)] 证明 $\|\cdot\|'$ 在 $A$ 上与 $\|\cdot\|$ 一致, $A$ 在 $A^{\tilde{e}}$ 中闭, 且 $A^{\tilde{e}}$ 对 $\|\cdot\|'$ 完备.
\item[(c)] 证明 $A^{\tilde{e}}$ 对范数 $\|\cdot\|'$ 是 $\mathcal{C}^*$-代数.
\end{enumerate}
\end{exercise}

\begin{solution}
(a). 设 $x_1 = y_1 + \lambda_1e$, $x_2 = y_2 + \lambda_2e$ 是 $A^{\tilde{e}}$ 的元素 ($y_1, y_2 \in A$, $\lambda_1, \lambda_2 \in \CC$), $c \in \CC$. 则
\begin{align*}
\|x_1 + x_2\|' &= \sup\{\|y_1y + \lambda_1y + y_2y + \lambda_2y\| \mid y \in A,\, \|y\| = 1\} \\
&\leq \sup\{\|y_1y + \lambda_1y\| \mid y \in A,\, \|y\| \leq 1\} + \sup\{\|y_2y + \lambda_2y\| \mid y \in A,\, \|y\| \leq 1\} \\
&= \|x_1\|' + \|x_2\|',
\end{align*}
\[
\|cx_1\|' = \sup\{\|cy_1y + c\lambda_1y\| \mid y \in A,\, \|y\| \leq 1\} = |c|\sup\{\|y_1y + \lambda_1y\| \mid y \in A,\, \|y\| \leq 1\} = |c|\|x_1\|',
\]
且
\begin{align*}
\|x_1x_2\|' &= \sup\{\|(y_1y_2 + \lambda_2y_1 + \lambda_1y_2)y + \lambda_1\lambda_2y\| \mid y \in A,\, \|y\| = 1\} \\
&= \sup\{\|y_1(y_2y + \lambda_2y) + \lambda_1(y_2y + \lambda_2y)\| \mid y \in A,\, \|y\| = 1\} \\
&\leq \sup\{\|y_1 + \lambda_1e\|'\|y_2y + \lambda_2y\| \mid y \in A,\, \|y\| \leq 1\} \\
&= \|x_1\|'\|x_2\|'.
\end{align*}
为证明 $\|\cdot\|'$ 是 $A^{\tilde{e}}$ 上的范数, 还需证明 $x + \lambda e \neq 0$ 时 $\|x + \lambda e\|' \neq 0$. 假设 $\|x + \lambda e\|' = 0$, 则对每个满足 $\|y\| = 1$ 的 $y \in A$ 有 $xy + \lambda y = 0$, 从而对每个 $y \in A$ 成立. 若 $x = 0$, 则 $\lambda = 0$. 若 $x \neq 0$, 则取 $y = x^*$ (注意 $xx^* \neq 0$, 因为 $\|xx^*\| = \|x^*\|^2 \neq 0$), 看到 $\lambda \neq 0$. 则对每个 $y \in A$ 有 $\lambda^{-1}xy = y$, 即 $\lambda^{-1}x$ 是 $A$ 的左单位元. 这蕴含 $(\lambda^{-1}x)^*$ 是 $A$ 的右单位元, 所以 $A$ 有单位元, 与我们假设矛盾. 所以 $x = 0$.

(b). 若 $x \in A$, 则
\[
\|x\|' = \sup\{\|xy\| \mid y \in A,\, \|y\| = 1\} \leq \|x\|.
\]
若 $x = 0$, 则 $\|x\|' = \|x\| = 0$. 否则, 也有 $x^* \neq 0$; 取 $y = \frac{1}{\|x^*\|}x^*$, 得到
\[
\|x\|' \geq \frac{1}{\|x^*\|}\|xx^*\| = \|x^*\| = \|x\|.
\]
所以 $A$ 对 $\|\cdot\|'$ 完备, 从而是 $A^{\tilde{e}}$ 的闭子集. 特别地, 商映射 $A^{\tilde{e}} \to A^{\tilde{e}}/A \simeq \CC$, $x + \lambda e \longmapsto \lambda$ 连续.
现在我们证明 $A^{\tilde{e}}$ 对 $\|\cdot\|'$ 完备. 设 $(x_n + \lambda_ne)_{n \geq 0}$ 是 $A^{\tilde{e}}$ 中的 Cauchy 序列, 其中 $x_n \in A$, $\lambda_n \in \CC$. 由上一段, 序列 $(\lambda_n)_{n \geq 0}$ 是 Cauchy 列, 所以 $A$ 中的序列 $(x_n)_{n \geq 0}$ 也是 Cauchy 列. 由于两个范数在 $A$ 上一致, 序列 $(x_n)_{n \geq 0}$ 收敛到某个 $x \in A$, 当然 $(\lambda_n)_{n \geq 0}$ 收敛到某个 $\lambda \in \CC$. 现在显然 (用显然的事实: 对 $z \in A$, $\mu \in \CC$ 有 $\|z + \mu e\|' \leq \|z\| + |\mu|$) 序列 $(x_n + \lambda_ne)_{n \geq 0}$ 在 $A^{\tilde{e}}$ 中收敛到 $x + \lambda e$.

(c). 最后, 我们证明 $A^{\tilde{e}}$ 是 $\mathcal{C}^*$-代数. 设 $x \in A$, $\lambda \in \CC$. 我们要证明 $\|(x + \lambda e)^*(x + \lambda e)\|' = (\|x + \lambda e\|')^2$. 可以假设 $x + \lambda e \neq 0$. 设 $\varepsilon > 0$. 则可以找到 $y \in A$ 使得 $\|y\| = 1$ 且
\[
\|xy + \lambda y\| \geq \|x + \lambda e\|'(1 - \varepsilon).
\]
注意 $xy + \lambda y = (x + \lambda e)y$ (在 $A^{\tilde{e}}$ 中). 所以
\begin{align*}
(1 - \varepsilon)^2(\|x + \lambda e\|')^2 &\leq \|xy + \lambda y\|^2 \\
&= \|(xy + \lambda y)^*(xy + \lambda y)\| \\
&= \|y^*(x + \lambda e)^*(x + \lambda e)y\|' \\
&\leq \|y\|^2\|(x + \lambda e)^*(x + \lambda e)\|' \\
&= \|(x + \lambda e)^*(x + \lambda e)\|'.
\end{align*}
由于这对每个 $\varepsilon$ 成立, 得到
\[
\|(x + \lambda e)^*(x + \lambda e)\|' \geq (\|x + \lambda e\|')^2.
\]
利用范数的次乘性, 推出
\[
\|x + \lambda e\|' \leq \|(x + \lambda e)^*\|'.
\]
由于 $*$ 在 $A^{\tilde{e}}$ 上双射, 最后一个不等式实际上是等式, 所以还得到
\[
(\|x + \lambda e\|')^2 \leq \|(x + \lambda e)^*(x + \lambda e)\|' \leq (\|x + \lambda e\|')^2,
\]
这就完成了证明.
\end{solution}


\chapter[Gelfand-Raikov 定理]{Gelfand-Raikov 定理}
% ============ 第 III 部分: \index{Gelfand-Raikov 定理}Gelfand-Raikov 定理 (The Gelfand-Raikov Theorem) ============
% 原书 III.1 - III.6 (原书 pp.145-194)
% 编号: III.2.1 (节级) / III.6.1.1 (子节级, 习题)

本章的目标是证明 Gelfand-Raikov 定理, 它说局部紧群的不可约酉表示分离点 (即若 $G$ 是局部紧群且 $x \in G - \{1\}$, 则存在 $G$ 的不可约酉表示使得 $\pi(x) \neq 1$).
在本章中, $G$ 是局部紧群, $\mu$ (或直接写``$dx$'') 是 $G$ 上的左 Haar 测度.

\section{\texorpdfstring{$L^{\infty}(G)$}{L-infinity(G)}}
你可以安全地忽略本节, 假设所有群都是 $\sigma$-紧的.
我们将在本章更认真地使用 $L^{\infty}(G)$, 并希望它是 $L^1(G)$ 的连续对偶, 但若 $G$ 不是 $\sigma$-紧的这不是真的. 所以我们修改 $L^{\infty}(G)$ 的定义使之成立. 见 {[}11{]} 的第 2.3 节.
更一般地, 设 $X$ 是局部紧 Hausdorff 拓扑空间, $\mu$ 是正则 Borel 测度. 若对 $X$ 的每个满足 $\mu(F) < +\infty$ 的 Borel 子集 $F$, $E \cap F$ 是 $X$ 的 Borel 子集, 则称 $E \subset X$ \index{局部 Borel 集}局部 Borel (locally Borel). 若 $E$ 局部 Borel, 且对 $X$ 的每个满足 $\mu(F) < +\infty$ 的 Borel 子集 $F$ 有 $\mu(E \cap F) = 0$, 则称 $E$ 局部零 (locally null). 若关于 $X$ 的点的断言在某个\index{局部零子集}局部零子集之外成立, 则称它\index{局部几乎处处}局部几乎处处 (locally almost everywhere) 成立. 我们已经见过: 若对 $\CC$ 的每个 Borel 子集 $A$, 集合 $f^{-1}(A)$ 局部 Borel, 则函数 $f \colon X \to \CC$ 局部可测 (locally measurable). 现在令 $L^{\infty}(X)$ 为局部几乎处处有界的\index{局部可测函数}局部可测函数 $X \to \CC$ 的空间, 商去等价关系: $f \sim g$ 若 $f - g$ 局部几乎处处为零. $L^{\infty}(X)$ 上的范数由
\[
\|f\|_{\infty} = \inf\{c \in \R_{\geq 0} \mid |f(x)| \leq c \text{ 局部几乎处处}\}
\]
给出.

\section{\index{正定型函数}正定型函数}

\begin{definition}[正定型函数 (function of positive type)]
\label{def:III.2.1}
$G$ 上的正定型函数是满足对每个 $f \in L^1(G)$ 有
\[
\int_G (f^* * f)(x)\phi(x)\,\dd x \geq 0
\]
的函数 $\phi \in L^{\infty}(G)$.
\end{definition}

注意若 $f \in L^1(G)$ 则 $f^* * f \in L^1(G)$, 所以积分收敛.

\begin{remark}[卷积与正定型]
\label{rem:III.2.2}
对每个 $\phi \in L^{\infty}(G)$ 和每个 $f, g \in L^1(G)$,
\begin{align*}
\int_G (f^* * g)(x)\phi(x)\,\dd x &= \int_{G \times G}f^*(y)g(y^{-1}x)\phi(x)\,\dd x\,\dd y \\
&= \int_{G \times G}\Delta(y)^{-1}\overline{f(y^{-1})}g(y^{-1}x)\phi(x)\,\dd x\,\dd y \\
&= \int_{G \times G}\overline{f(y)}g(yx)\phi(x)\,\dd x\,\dd y \\
&= \int_{G \times G}\overline{f(y)}g(x)\phi(y^{-1}x)\,\dd x\,\dd y.
\end{align*}
\end{remark}

\begin{example}[正定型函数的例子]
\label{exa:III.2.3}
\begin{enumerate}
\item[(1)] $0$ 是正定型函数.
\item[(2)] 若 $\phi \colon G \to S^1 \subset \CC$ 是 $1$ 维表示 (即对所有 $x, y \in G$ 有 $\phi(xy) = \phi(x)\phi(y)$), 则它是正定型函数. 事实上, 对每个 $f \in L^1(G)$, 由注记 \ref{rem:III.2.2},
\begin{align*}
\int_G (f^* * f)(x)\phi(x)\,\dd x &= \int_{G \times G}\overline{f(y)}f(yx)\phi(y^{-1}x)\,\dd x\,\dd y \\
&= \int_{G \times G}\overline{f(y)}f(x)\phi(y)\phi(x)\,\dd x\,\dd y \\
&= \Bigl|\int_G \phi(x)f(x)\,\dd x\Bigr|^2 \geq 0.
\end{align*}
\end{enumerate}
我们将在下面的命题 (ii) 中推广第二个例子.
\end{example}

\begin{proposition}[正定型函数的来源]
\label{prop:III.2.4}
\begin{enumerate}
\item[(i)] 若 $\phi \colon G \to \CC$ 是正定型函数, 则 $\bar{\phi}$ 也是.
\item[(ii)] 若 $(\pi, V)$ 是 $G$ 的酉表示且 $v \in V$, 则 $\phi \colon G \to \CC$, $x \longmapsto \langle \pi(x)(v), v \rangle$ 是连续正定型函数.
\item[(iii)] 设 $f \in L^2(G)$, 定义 $\tilde{f} \colon G \to \CC$ 为 $\tilde{f}(x) = \overline{f(x^{-1})}$. 则 $f * \tilde{f}$ 有意义, 在 $L^{\infty}(G)$ 中, 且是正定型函数.
\end{enumerate}
\end{proposition}

\begin{proof}
(i). 设 $f \in L^1(G)$. 则由注记 \ref{rem:III.2.2},
\begin{align*}
\int_G (f^* * f)\bar{\phi}\,\dd\mu &= \overline{\int_{G \times G}\overline{f(y)}f(yx)\phi(x)\,\dd x\,\dd y} \\
&= \overline{\int_{G \times G}\overline{f(y)}f(yx)\phi(x)\,\dd x\,\dd y} \\
&= \overline{\int_G (f^* * f)\phi\,\dd\mu} \geq 0.
\end{align*}

(ii). 函数 $\phi$ 连续, 因为 $G \to V$, $x \longmapsto \pi(x)(v)$ 连续. 让我们证明它是正定型的. 注意对所有 $x, y \in G$,
\[
\phi(y^{-1}x) = \langle \pi(y^{-1}x)(v), v \rangle = \langle \pi(x)(v), \pi(y)(v) \rangle.
\]
设 $f \in L^1(G)$. 则由注记 \ref{rem:III.2.2},
\begin{align*}
\int_G (f^* * f)\phi\,\dd\mu &= \int_{G \times G}\overline{f(x)}f(y)\phi(y^{-1}x)\,\dd x\,\dd y \\
&= \int_{G \times G}\langle f(x)\pi(x)(v), f(y)\pi(y)(v) \rangle\,\dd x\,\dd y \\
&= \langle \pi(f)(v), \pi(f)(v) \rangle \geq 0.
\end{align*}

(iii). 设 $x \in G$. 则定义 $f * \tilde{f}(x)$ 的积分是
\[
\int_G f(y)\overline{f(x^{-1}y)}\,\dd y.
\]
这个积分收敛, 因为 $f$ 和 $L_xf \colon y \longmapsto f(x^{-1}y)$ 都在 $L^2(G)$ 中 (由 $\mu$ 的左不变性). 又, 由 Cauchy-Schwarz 不等式,
\[
|f * \tilde{f}(x)| \leq \|f\|_2\|L_xf\|_2 = \|f\|_2^2.
\]
所以 $f * \tilde{f} \in L^{\infty}(G)$.
让我们证明 $f * \tilde{f}$ 是正定型的. 设 $\pi_L$ 是 $G$ 的左正则表示, 即 $G$ 在 $L^2(G)$ 上由 $\pi_L(x) = L_x$ 给出的酉表示. 则对每个 $x \in G$,
\[
\langle \pi_L(x)(f), f \rangle = \int_G f(x^{-1}y)\overline{f(y)}\,\dd y = \int_G \overline{\tilde{f}(y^{-1}x)}f(y)\,\dd y = f * \tilde{f}(x).
\]
所以结果由 (ii) 得出.
\end{proof}

本节的主要结果是上面的例子 (ii) 是仅有的情形.

\begin{theorem}[Gelfand-Raikov 构造]
\label{thm:III.2.5}
设 $\phi \colon G \to \CC$ 是正定型函数. 则存在 $G$ 的循环酉表示 $(\pi, V)$ 和 $V$ 的循环向量 $v$ 使得 $\phi(x) = \langle \pi(x)(v), v \rangle$ 局部几乎处处成立.
而且, 表示 $\pi$ 和向量 $v$ 由 $\phi$ 唯一确定, 意义如下: 若 $(\pi', V')$ 是另一个循环酉表示且 $v' \in V'$ 是循环向量, 满足 $\phi(x) = \langle \pi'(x)(v'), v' \rangle$ 局部几乎处处成立, 则存在 $G$-等变等距 $T \colon V \to V'$ 使得 $T(v) = v'$.
\end{theorem}

事实上, 我们将在证明中给出 $(\pi, V)$ 的一个相当显式的构造.
在证明定理之前, 让我们看一些简单推论.

\begin{corollary}[正定型函数的性质]
\label{cor:III.2.6}
设 $\phi \colon G \to \CC$ 是正定型函数. 则 $\phi$ 局部几乎处处与某个连续函数一致, $\|\phi\|_{\infty} = \phi(1)$, 且对每个 $x \in G$ 有 $\phi(x^{-1}) = \overline{\phi(x)}$.
\end{corollary}

\begin{proof}
第一个结论由命题 \ref{prop:III.2.4}(ii) 得出. 为证明其余结论, 选取 $G$ 的循环酉表示 $(\pi, V)$ 和 $v \in V$ 使得 $\phi(x) = \langle \pi(x)(v), v \rangle$. 则对每个 $x \in G$,
\[
|\phi(x)| \leq \|\pi(x)(v)\|\|v\| = \|v\|^2 = \phi(1)
\]
且
\[
\phi(x^{-1}) = \langle \pi(x^{-1})(v), v \rangle = \langle \pi(x)^*(v), v \rangle = \langle v, \pi(x)(v) \rangle = \overline{\phi(x)}.
\]
\end{proof}

现在我们回到定理的证明. 设 $\phi \colon G \to \CC$ 是正定型函数. 在 $L^1(G)$ 上定义 Hermitian 形式 $\langle \cdot, \cdot \rangle_{\phi}$:
\[
\langle f, g \rangle_{\phi} = \int_G (g^* * f)\phi = \int_{G \times G}\overline{f(x)}g(y)\phi(y^{-1}x)\,\dd x\,\dd y
\]
(见注记 \ref{rem:III.2.2}). 特别地, 显然对所有 $f, g \in L^1(G)$,
\[
|\langle f, g \rangle_{\phi}| \leq \|f\|_1\|g\|_1\|\phi\|_{\infty}.
\]
由于 $\phi$ 正定型, 有 $\langle f, f \rangle_{\phi} \geq 0$, 即我们刚定义的 Hermitian 形式半正定; 特别地, 对它适用 Cauchy-Schwarz 不等式, 它对所有 $f, g \in L^1(G)$ 给出
\[
|\langle f, g \rangle_{\phi}|^2 \leq \langle f, f \rangle_{\phi}\langle g, g \rangle_{\phi}.
\]
设 $N$ 是形式 $\langle \cdot, \cdot \rangle_{\phi}$ 的核 (或根), 即 $L^1(G)$ 的正交补, 即满足对每个 $g \in L^1(G)$ 有 $\langle f, g \rangle_{\phi} = 0$ 的 $f \in L^1(G)$ 的空间. 由 Cauchy-Schwarz 不等式, $f \in N$ 当且仅当 $\langle f, f \rangle_{\phi} = 0$. 所以形式 $\langle \cdot, \cdot \rangle_{\phi}$ 在 $L^1(G)/N$ 上定义正定 Hermitian 形式, 仍记为 $\langle \cdot, \cdot \rangle_{\phi}$; 记对应范数为 $\|\cdot\|_{\phi}$. 对每个 $f \in L^1(G)$,
\[
\|f + N\|_{\phi}^2 \leq \|\phi\|_{\infty}\|f\|_1^2.
\]
设 $V_{\phi}$ 是 $L^1(G)/N$ 对范数 $\|\cdot\|_{\phi}$ 的完备化; 这是希尔伯特空间.
我们要构造 $G$ 在 $V_{\phi}$ 上的酉作用. 我们已经用算子 $L_x$ 有 $G$ 在 $L^1(G)$ 上的连续表示. 这将魔术般地给出我们的酉表示. 首先注意, 对每个 $f \in L^1(G)$, 映射 $G \to L^1(G)$, $x \longmapsto L_xf$ 对半范数 $\|\cdot\|_{\phi}$ 连续, 因为由刚证明的不等式 $\|\cdot\|_{\phi} \leq \|\phi\|_{\infty}^{1/2}\|\cdot\|_1$.
让我们证明 $\langle \cdot, \cdot \rangle_{\phi}$ 被 $G$ 的作用不变. 设 $x \in G$, $f, g \in L^1(G)$. 则
\begin{align*}
\langle L_xf, L_xg \rangle_{\phi} &= \int_{G \times G}\overline{f(x^{-1}y)}g(x^{-1}z)\phi(z^{-1}y)\,\dd y\,\dd z \\
&= \int_{G \times G}\overline{f(y)}g(z)\phi((xz)^{-1}(xy))\,\dd y\,\dd z \\
&= \int_{G \times G}\overline{f(y)}g(z)\phi(z^{-1}y)\,\dd y\,\dd z = \langle f, g \rangle_{\phi}.
\end{align*}
特别地, 形式 $\langle \cdot, \cdot \rangle_{\phi}$ 的根 $N$ 是 $L^1(G)$ 的 $G$-不变子空间, 所以我们得到 $G$ 在 $L^1(G)/N$ 上的作用, 它保持 Hermitian 内积, 且由命题 \ref{prop:I.3.1.10} 是连续表示. 我们把这个作用按连续性延拓到 $V_{\phi}$. 这给出 $G$ 在 $V_{\phi}$ 上的酉表示, 记为 $\pi_{\phi}$.
设 $f, g \in L^1(G)$. 则由例 \ref{exa:I.4.2.7},
\[
\pi_{\phi}(f)(g + N) = f * g + N.
\]
下面的引理将蕴含定理 \ref{thm:III.2.5} 的第一个结论.

\begin{lemma}[循环向量]
\label{lem:III.2.7}
存在 $V_{\phi}$ 的循环向量 $v = v_{\phi}$ 使得:
\begin{enumerate}
\item[(i)] 对 $f \in L^1(G)$, 有 $\pi_{\phi}(f)(v) = f + N$;
\item[(ii)] 在 $G$ 上局部几乎处处有 $\phi(x) = \langle \pi_{\phi}(x)(v), v \rangle_{\phi}$.
\end{enumerate}
\end{lemma}

\begin{proof}
由 $\pi_{\phi}(f)(g + N)$ ($f, g \in L^1(G)$) 的计算 (见上面), 我们看到 $v$ 应是 $L^1(G)/N$ 中 $*$ 的单位元 (即 $1 \in G$ 处的 Dirac 测度) 的像, 如果这样的单位元存在的话. 一般地它不存在, 但我们可以逼近它, 并希望得到 $L^1(G)/N$ 中的 Cauchy 序列.
所以设 $(\psi_U)_{U \in \mathcal{U}}$ 是近似单位元 (见定义 \ref{def:I.4.1.7}). 注意 $(\psi_U^*)_{U \in \mathcal{U}}$ 也是近似单位元, 所以由命题 \ref{prop:I.4.1.9}, 对每个 $f \in L^1(G)$ 有 $\psi_U^* * f \xrightarrow[U \to \{1\}]{} f$ ($L^1(G)$ 中). 所以对每个 $f \in L^1(G)$,
\[
\langle f, \psi_U \rangle_{\phi} = \int_G (\psi_U^* * f)\phi\,\dd\mu \xrightarrow[U \to \{1\}]{} \int_G f\phi\,\dd\mu.
\]
所以 $f \longmapsto \int_G f\phi\,\dd\mu$ 是 $L^1(G)$ 上对 $\|\cdot\|_1$ 和 $\|\cdot\|_{\phi}$ 有界的线性泛函, 其核包含 $N$. 我们可以把这个有界线性泛函下降到 $L^1(G)/N$ 并按连续性延拓到 $V_{\phi}$, 得到 $V_{\phi}$ 上的有界线性泛函, 它必具有 $\langle \cdot, v \rangle_{\phi}$ 的形式 (对某个唯一确定的 $v \in V_{\phi}$), 因为 $V_{\phi}$ 是希尔伯特空间. 由 $v$ 的定义, 对每个 $f \in L^1(G)$,
\[
\langle f + N, v \rangle_{\phi} = \int_G f\phi\,\dd\mu,
\]
且这确定 $v$, 因为 $L^1(G)$ 的像在 $V_{\phi}$ 中稠密.
现在我们证明性质 (i) 和 (ii). 设 $f, g \in L^1(G)$. 则
\begin{align*}
\langle g, f \rangle_{\phi} &= \int_G (f^* * g)\phi\,\dd\mu \\
&= \int_{G \times G}\overline{g(x)}f(y)\phi(y^{-1}x)\,\dd x\,\dd y \\
&= \int_{G \times G}\overline{g(yx)}f(y)\phi(x)\,\dd x\,\dd y \\
&= \int_{G \times G}f(y)\overline{L_{y^{-1}}g(x)}\phi(x)\,\dd x\,\dd y \\
&= \int_G f(y)\langle \pi_{\phi}(y^{-1})(g + N), v \rangle_{\phi}\,\dd x \\
&= \int_G \langle g + N, f(y)\pi_{\phi}(y)(v) \rangle_{\phi}\,\dd x \\
&= \langle g + N, \pi_{\phi}(f)(v) \rangle_{\phi}.
\end{align*}
由于这对每个 $g \in L^1(G)$ 成立且 $L^1(G)$ 的像在 $V_{\phi}$ 中稠密, 得到 $\pi_{\phi}(f)(v) = f + N$. 特别地, $\{\pi_{\phi}(f)(v) \mid f \in L^1(G)\}$ 的张成在 $V_{\phi}$ 中稠密, 所以 $v$ 是循环向量 (由定理 \ref{thm:I.4.2.6}(iii)).
又, 对 $f \in L^1(G)$, 由刚才证明的:
\begin{align*}
\int_G f(x)\langle \pi_{\phi}(x)(v), v \rangle_{\phi}\,\dd x &= \Bigl\langle \int_G f(x)\pi_{\phi}(x)(v)\,\dd x, v \Bigr\rangle_{\phi} \\
&= \langle \pi_{\phi}(f)(v), v \rangle_{\phi} \\
&= \langle f + N, v \rangle_{\phi} \\
&= \int_G f(x)\phi(x)\,\dd x.
\end{align*}
由于这对每个 $f \in L^1(G)$ 成立, 它蕴含 $\phi(x) = \langle \pi_{\phi}(x)(v), v \rangle_{\phi}$ 局部几乎处处成立.
\end{proof}

为完成定理 \ref{thm:III.2.5} 的证明, 我们只需建立下面的引理.

\begin{lemma}[唯一性]
\label{lem:III.2.8}
设 $(\pi, V)$ 和 $(\pi', V')$ 是 $G$ 的两个循环酉表示, $v \in V$, $v' \in V'$ 是循环向量, 满足对每个 $x \in G$,
\[
\langle \pi(x)(v), v \rangle = \langle \pi'(x)(v'), v' \rangle.
\]
则存在 $G$-等变等距 $T \colon V \to V'$ 使得 $T(v) = v'$.
\end{lemma}

\begin{proof}
当然, 我们要通过公式 $T(\pi(x)(v)) = \pi'(x)(v')$ (对每个 $x \in G$) 定义 $T \colon V \to V'$. 我们需要使它有意义. 设 $W = \operatorname{Span}\{\pi(x)(v) \mid x \in G\}$. 由 $v$ 循环的假设, 子空间 $W$ 在 $V$ 中稠密. 让我们验证上面的公式定义等距 $T \colon W \to V'$. 设 $x_1, \ldots, x_n \in G$, $\lambda_1, \ldots, \lambda_n \in \CC$. 则
\begin{align*}
\Bigl\|\sum_{i=1}^n\lambda_i\pi(x_i)(v)\Bigr\|^2 &= \sum_{i=1}^n\sum_{j=1}^n\lambda_i\bar{\lambda}_j\langle \pi(x_j^{-1}x_i)(v), v \rangle \\
&= \sum_{i=1}^n\sum_{j=1}^n\lambda_i\bar{\lambda}_j\langle \pi'(x_j^{-1}x_i)(v'), v' \rangle \\
&= \Bigl\|\sum_{i=1}^n\lambda_i\pi'(x_i)(v')\Bigr\|^2.
\end{align*}
特别地, 若 $\sum_{i=1}^n\lambda_i\pi(x_i)(v) = 0$, 则也有 $\sum_{i=1}^n\lambda_i\pi'(x_i)(v') = 0$. 所以可以定义 $T \colon W \to V'$, $T\bigl(\sum_{i=1}^n\lambda_i\pi(x_i)(v)\bigr) = \sum_{i=1}^n\lambda_i\pi'(x_i)(v')$, 且上面的计算表明 $T$ 是等距. 所以 $T$ 连续, 从而可以延拓成连续线性算子 $T \colon V \to V'$, 它仍是等距, 所以单射且像闭. 又, 若 $x \in G$, $w \in W$, 则由 $T$ 的定义 $T(\pi(x)(w)) = \pi'(x)(T(w))$. 由于 $T$ 连续且 $W$ 在 $V$ 中稠密, 这对每个 $w \in V$ 仍成立; 换句话说, $T$ 是 $G$-等变的. 最后, 由 $T$ 的定义 $T(v) = v'$, 所以 $T$ 的像在 $V'$ 中稠密, 从而等于 $V'$.
\end{proof}

\section{正定型函数与不可约表示}

我们已经看到 $G$ 的循环酉表示 (连同固定的循环向量) 由正定型函数参数化. 下一个自然的问题是``哪些正定型函数对应不可约表示?''
\begin{definition}[$P$, $P_1$, $P_0$]
\label{def:III.3.1}
记 $P(G)$ 或 $P$ 为 $G$ 上连续正定型函数的集合. 这是 $\mathcal{C}_b(G)$ 中的凸锥.\footnote{``锥''意味着它被 $\R_{\geq 0}$ 的元素乘法稳定.}
设
\[
P_1 = \{\phi \in P \mid \|\phi\|_{\infty} = 1\} = \{\phi \in P \mid \phi(1) = 1\}
\]
且
\[
P_0 = \{\phi \in P \mid \|\phi\|_{\infty} \leq 1\} = \{\phi \in P \mid \phi(1) \leq 1\}.
\]
(回忆由推论 \ref{cor:III.2.6}, 对每个 $\phi \in P$ 有 $\|\phi\|_{\infty} = \phi(1)$.)
则 $P_1$ 和 $P_0$ 是 $\mathcal{C}_b(G)$ 的凸子集. 记 $E(P_1)$ 和 $E(P_0)$ 为它们的端点集.
\end{definition}

\begin{theorem}[端点与不可约性]
\label{thm:III.3.2}
设 $\phi \in P_1$. 则上一节构造的酉表示 $(V_{\phi}, \pi_{\phi})$ 不可约当且仅当 $\phi \in E(P_1)$.
\end{theorem}

\begin{remark}[纯量倍数]
\label{rem:III.3.3}
若 $\phi \in P$, $c \in \R_{>0}$, 则 $\langle \cdot, \cdot \rangle_{c\phi} = c\langle \cdot, \cdot \rangle_{\phi}$, 所以 $V_{c\phi} = V_{\phi}$, $\pi_{c\phi} = \pi_{\phi}$, $v_{c\phi} = v_{\phi}$. (但 $V_{\phi}$ 的恒等不是等距, 因为我们用了两个不同的内积, 即 $\langle \cdot, \cdot \rangle_{\phi}$ 和 $\langle \cdot, \cdot \rangle_{c\phi}$.) 由于 $P$ 的每个非零元 $\phi$ 都具有 $c\phi'$ 的形式 (其中 $c = \phi(1)$, $\phi' \in P_1$ 唯一), 定理回答了本节开头的问题.
\end{remark}

\begin{remark}[交换群的情形]
\label{rem:III.3.4}
若 $G$ 交换, 定理说 $\widehat{G} = E(P_1)$.
\end{remark}

\begin{proof}
在此证明中, 记 $V = V_{\phi}$ 的内积和范数为 $\langle \cdot, \cdot \rangle$ 和 $\|\cdot\|$, 记 $\pi = \pi_{\phi}$. (除非这引起混淆.)
我们首先假设 $\pi$ 不可约. 设 $0 \neq W \subsetneq V$ 是闭 $G$-不变子空间. 由于 $\pi$ 酉, $W^{\perp}$ 也 $G$-不变, 且 $V = W \oplus W^{\perp}$. 设 $v \in V$ 是引理 \ref{lem:III.2.7} 的循环向量. 由于 $v$ 循环, 它不能包含于 $W$ 或 $W^{\perp}$ (否则会有 $W = V$ 或 $W^{\perp} = V$). 所以可以把 $v$ 写成 $v = v_1 + v_2$, 其中 $v_1 \in W$, $v_2 \in W^{\perp}$, $v_1, v_2 \neq 0$. 定义 $\phi_1, \phi_2 \colon G \to \CC$ 为 $\phi_i(x) = \langle \pi(x)(v_i), v_i \rangle$. 则由命题 \ref{prop:III.2.4}(ii), $\phi_1, \phi_2 \in P$, 且有 $\phi = \phi_1 + \phi_2$. 设 $c_1 = \|v_1\|^2$, $c_2 = \|v_2\|^2$; 由勾股定理 $c_1 + c_2 = \|v\|^2 = \phi(1) = 1$, 所以 $c_1, c_2 \in (0, 1)$. 设 $\psi_i = \frac{1}{c_i}\phi_i$ ($i = 1, 2$). 则 $\phi = c_1\psi_1 + c_2\psi_2$, 且 $\psi_1, \psi_2 \in P_1$ (因为 $\psi_1(1) = \psi_2(1) = 1$). 为得出结论 $\phi$ 不是 $P_1$ 的端点, 还需证明 $\psi_1 \neq \psi_2$, 即 $\phi_2$ 不具有 $c\phi_1$ 的形式 (对 $c \in \R_{>0}$).
设 $c \in \R_{>0}$. 选取 $\varepsilon > 0$ 使得 $\varepsilon < \frac{c\|v_1\|^2}{\|v_2\| + c\|v_1\|}$, 即 $\varepsilon\|v_2\| < c\|v_1\|^2 - \varepsilon c\|v_1\|$. 由于 $v$ 是 $V$ 的循环向量, 可以找到 $x_1, \ldots, x_n \in G$ 和 $a_1, \ldots, a_n \in \CC$ 使得
\[
\Bigl\|\sum_{i=1}^n a_i\pi(x_i)(v) - v_1\Bigr\| < \varepsilon.
\]
由于 $v = v_1 + v_2$ (其中 $v_1 \in W$, $v_2 \in W^{\perp}$), 且 $W$ 和 $W^{\perp}$ 都被 $G$ 的作用稳定, 对 $x \in G$,
\[
\langle \pi(x)(v), v_1 \rangle = \langle \pi(x)(v_1) + \pi(x)(v_2), v_1 \rangle = \langle \pi(x)(v_1), v_1 \rangle.
\]
所以
\begin{align*}
\Bigl|\sum_{i=1}^n a_i\langle \pi(x_i)(v_1), v_1 \rangle - \langle v_1, v_1 \rangle\Bigr| &= \Bigl|\sum_{i=1}^n a_i\langle \pi(x_i)(v), v_1 \rangle - \langle v_1, v_1 \rangle\Bigr| \\
&= \Bigl|\Bigl\langle \sum_{i=1}^n a_i\pi(x_i)(v) - v_1, v_1 \Bigr\rangle\Bigr| \\
&< \varepsilon\|v_1\|,
\end{align*}
这蕴含
\[
\|v_1\|^2 - \varepsilon\|v_1\| < \Bigl|\sum_{i=1}^n a_i\langle \pi(x_i)(v_1), v_1 \rangle\Bigr| = \Bigl|\sum_{i=1}^n a_i\phi_1(x_i)\Bigr|.
\]
另一方面 (利用对每个 $x \in G$ 有 $\langle \pi(x)(v), v_2 \rangle = \langle \pi(x)(v_2), v_2 \rangle$ 这一事实),
\begin{align*}
\Bigl|\sum_{i=1}^n a_i\langle \pi(x_i)(v_2), v_2 \rangle\Bigr| &= \Bigl|\Bigl\langle \sum_{i=1}^n a_i\pi(x_i)(v) - v_1, v_2 \Bigr\rangle\Bigr| \\
&\leq \Bigl\|\sum_{i=1}^n a_i\pi(x_i)(v) - v_1\Bigr\|\|v_2\| \\
&< \varepsilon\|v_2\| \\
&< c\|v_1\|^2 - \varepsilon c\|v_1\| \\
&< c\Bigl|\sum_{i=1}^n a_i\phi_1(x_i)\Bigr|,
\end{align*}
即
\[
\Bigl|\sum_{i=1}^n a_i\phi_2(x_i)\Bigr| < c\Bigl|\sum_{i=1}^n a_i\phi_1(x_i)\Bigr|.
\]
所以我们不能有 $\phi_2 = c\phi_1$. 由于 $c$ 任意, 这完成 $\psi_1 \neq \psi_2$ 的证明, 从而 $\phi$ 不是 $P_1$ 的端点.
反之, 我们要证明若 $\pi_{\phi}$ 不可约则 $\phi$ 在 $P_1$ 中端点. 假设 $\phi = \phi_1 + \phi_2$, 其中 $\phi_1, \phi_2 \in P$. 对每个 $f \in L^1(G)$,
\[
\langle f, f \rangle_{\phi_1} = \langle f, f \rangle_{\phi} - \langle f, f \rangle_{\phi_2} \leq \langle f, f \rangle_{\phi}.
\]
特别地, $\langle \cdot, \cdot \rangle_{\phi}$ 的核包含于 $\langle \cdot, \cdot \rangle_{\phi_1}$ 的核, 所以 $L^1(G)$ 的恒等延拓成连续满射 $T \colon V_{\phi} \to V_{\phi_1}$, 且此映射是 $G$-等变的, 因为 $V_{\phi}$ 和 $V_{\phi_1}$ 上的 $G$ 作用都来自 $G$ 在 $L^1(G)$ 上通过左平移的作用. 又, 由于 $v_{\phi}$ (相应地 $v_{\phi_1}$) 只是近似单位元的像在 $V_{\phi}$ (相应地 $V_{\phi_1}$) 中的极限, 算子 $T$ 把 $v_{\phi}$ 送到 $v_{\phi_1}$. 由于 $\Ker T$ 是 $V_{\phi}$ 的 $G$-不变子空间, $(\Ker T)^{\perp}$ 也是, 所以 $T$ 定义 $G$-等变同构 $(\Ker T)^{\perp} \xrightarrow{\sim} V_{\phi_1}$, 从而 $V_{\phi_1}$ 同构于 $V_{\phi}$ 的子表示.
现在假设 $\pi_{\phi}$ 不可约. 则 $T^*T \in \End(V_{\phi})$ 是 $G$-等变的, 所以由 Schur 引理 (定理 \ref{thm:I.3.4.1}) 对某个 $c \in \CC$ 有 $T^*T = c\,\mathrm{id}_{V_{\phi}}$. 由于 $T(v_{\phi}) = v_{\phi_1}$, 对每个 $x \in G$,
\begin{align*}
\phi_1(x) &= \langle \pi_{\phi_1}(x)(v_{\phi_1}), v_{\phi_1} \rangle_{\phi_1} \\
&= \langle \pi_{\phi_1}(x)(T(v_{\phi})), T(v_{\phi}) \rangle_{\phi_1} \\
&= \langle T(\pi_{\phi}(x)(v_{\phi})), T(v_{\phi}) \rangle_{\phi_1} \\
&= \langle T^*T(\pi_{\phi}(x)(v_{\phi})), v_{\phi} \rangle_{\phi} \\
&= c\phi(x).
\end{align*}
由于 $\phi_1$ 和 $\phi$ 正定型, 必有 $c \in \R_{\geq 0}$. 类似地看到 $\phi_2$ 必在 $\R_{\geq 0}\phi$ 中. 所以 $\phi$ 端点.
\end{proof}

\section{\texorpdfstring{凸集 $P_1$}{凸集 P1}}
我们在前两节看到 $G$ 的不可约酉表示由 $P_1$ 的端点参数化. 回忆我们试图证明存在足够多的不可约酉表示来分离 $G$ 上的点. 所以我们要证明 $P_1$ 有大量端点. 自然的想法是用 Krein-\index{Milman 定理}Milman 定理 (定理 \ref{thm:B.5.2}, 它说紧凸集是其端点的闭凸包), 但 $P_1$ 一般不紧. 然而, 集合 $P_0$ 凸且弱*紧, 且与 $P_1$ 密切相关; 这足以把 Krein-Milman 定理的结论推广到 $P_1$.
回忆 $P$ 是 $L^{\infty}(G)$ 的子集. 我们把 $L^{\infty}(G)$ 等同于 $L^1(G)$ 的连续对偶, 并在它及其子集 $P$, $P_0$, $P_1$ 上考虑弱*拓扑. 对 $f \in L^{\infty}(G)$, $f$ 的邻域基由集合 $U_{g_1,\ldots,g_n,c} = \{f' \in L^{\infty}(G) \mid |\int_G (f - f')g_i\,\dd\mu| < c,\, 1 \leq i \leq n\}$ 给出, 其中 $n \in \Z_{\geq 1}$, $g_1, \ldots, g_n \in L^1(G)$, $c > 0$. 本节第二个主要结果是弱*拓扑与 $P_1$ 上的紧收敛拓扑一致.

\begin{theorem}[$E(P_1)$ 的凸包稠密]
\label{thm:III.4.1}
$E(P_1)$ 的凸包在 $P_1$ 中对弱*拓扑稠密.
\end{theorem}

\begin{lemma}[$E(P_0) = E(P_1) \cup \{0\}$]
\label{lem:III.4.2}
有 $E(P_0) = E(P_1) \cup \{0\}$.
\end{lemma}

\begin{proof}
首先证明 $E(P_1) \cup \{0\}$ 的每个点在 $P_0$ 中端点. 设 $\phi_1, \phi_2 \in P_0$, $c_1, c_2 \in (0, 1)$ 满足 $c_1 + c_2 = 1$. 若 $c_1\phi_1 + c_2\phi_2 = 0$, 则 $0 = c_1\phi_1(1) + c_2\phi_2(1)$, 所以 $\phi_1(1) = \phi_2(1) = 0$, 从而 $\|\phi_1\|_{\infty} = \|\phi_2\|_{\infty} = 0$, 即 $\phi_1 = \phi_2 = 0$. 这证明 $0$ 端点. 假设 $\phi := c_1\phi_1 + c_2\phi_2 \in E(P_1)$. 则 $1 = \phi(1) = c_1\phi_1(1) + \phi_2(1)$, 所以 $\phi_1(1) = \phi_2(1) = 1$, 因此 $\phi_1, \phi_2 \in P_1$; 由于 $\phi$ 在 $P_1$ 中端点, 这蕴含 $\phi_1 = \phi_2$. 所以 $\phi$ 在 $P_0$ 中也端点.
现在我们证明 $P_0$ 的每个端点都在 $E(P_1) \cup \{0\}$ 中. 设 $\phi \in P_0 - (E(P_1) \cup \{0\})$. 若 $\phi \in P_1$, 则它不端点. 若 $\phi \notin P_1$, 则 $0 < \phi(1) < 1$, 所以 $\phi = (1 - c)0 + c\frac{1}{\phi(1)}\phi$, 其中 $c = \phi(1) \in (0, 1)$ 且 $\frac{1}{\phi(1)}\phi \in P_0$; 这证明 $\phi$ 不端点.
\end{proof}

\begin{proof}[定理的证明]
注意定义 $L^{\infty}(G)$ 中 $P$ 的条件是弱*闭条件, 所以 $P$ 是 $L^{\infty}(G)$ 的弱*闭子集. 由于 $P_0$ 是 $P$ 与 $L^{\infty}(G)$ 的闭单位球的交, 它在这个闭单位球中弱*闭, 从而由 Banach-Alaoglu 定理 (定理 \ref{thm:B.4.1}) 弱*紧. 由于 $P_0$ 也凸, Krein-Milman 定理 (定理 \ref{thm:B.5.2}) 说 $E(P_0)$ 的凸包在 $P_0$ 中弱*稠密. 又, 上面的引理说 $E(P_0) = E(P_1) \cup \{0\}$.
设 $\phi \in P_1$, $U$ 是 $\phi$ 的形如 $\{\psi \in P_1 \mid |\int_G (\phi - \psi)g_i\,\dd\mu| < c,\, 1 \leq i \leq n\}$ ($n \in \Z_{\geq 1}$, $g_1, \ldots, g_n \in L^1(G)$, $c > 0$) 的弱*邻域. 我们要找到 $U$ 中属于 $E(P_1)$ 的凸包的点. 选取 $\varepsilon > 0$ (稍后看它需要多小). 由第一段和 $L^{\infty}(G)$ 中闭球弱*闭 (Hahn-Banach 定理的推论) 的事实, 可以找到 $E(P_1) \cup \{0\}$ 的凸包中的 $\psi$ 使得对每个 $i \in \{1, \ldots, n\}$ 有 $|\int_G (\phi - \psi)g_i\,\dd\mu| < c/2$ 且 $\|\psi\|_{\infty} \geq 1 - \varepsilon$. 把 $\psi$ 写成 $\psi = c_1\psi_1 + \cdots + c_r\psi_r$, 其中 $c_1, \ldots, c_r \in [0, 1]$, $\psi_1, \ldots, \psi_r \in E(P_1)$, $c_1 + \cdots + c_r \leq 1$. 设 $a = \frac{1}{\|\psi\|_{\infty}}$. 则 $a\psi = (ac_1)\psi_1 + \cdots + (ac_r)\psi_r$ 且 $ac_1 + \cdots + ac_r = 1$, 所以 $a\psi$ 在 $E(P_1)$ 的凸包中. 让我们证明 $a\psi \in U$. 若 $i \in \{1, \ldots, n\}$, 则
\begin{align*}
\Bigl|\int_G (\phi - a\psi)g_i\,\dd\mu\Bigr| &\leq \Bigl|\int_G (\phi - \psi)g_i\,\dd\mu\Bigr| + \Bigl|\int_G (\psi - a\psi)g_i\,\dd\mu\Bigr| \\
&< c/2 + |1 - a|\Bigl|\int_G \psi g_i\,\dd\mu\Bigr| \\
&< c/2 + \varepsilon\Bigl(c/2 + \Bigl|\int_G \phi g_i\,\dd\mu\Bigr|\Bigr).
\end{align*}
所以, 若选取 $\varepsilon$ 足够小使得对每个 $i \in \{1, \ldots, n\}$ 有 $\varepsilon\bigl(c/2 + |\int_G \phi g_i\,\dd\mu|\bigr) < c/2$, 函数 $a\psi$ 将在 $U$ 中.
\end{proof}

由于 $P$ 是空间 $\mathcal{C}(G)$ 的子空间, 我们也可以在 $P$ 上考虑紧收敛拓扑, 即在 $G$ 的紧子集上的收敛拓扑. 若 $\phi \in P$, 则 $\phi$ 在此拓扑下的邻域基由 $\{\psi \in P \mid \sup_{x \in K}|\phi(x) - \psi(x)| < c\}$ 给出, 其中 $K$ 遍历 $G$ 的紧子集, $c$ 遍历正实数.

\begin{theorem}[Raikov 定理]
\label{thm:III.4.3}
在 $P$ 的子集 $P_1$ 上, 紧收敛拓扑与弱*拓扑一致.
\end{theorem}

\begin{remark}[定理不适用于 $P_0$]
\label{rem:III.4.4}
此定理推广了习题 \ref{ex:I.5.4.2}(a). (见注记 \ref{rem:III.3.4}.)
注意定理对 $P_0$ 不成立. 例如, 若 $G = \R$, 则紧收敛拓扑与弱*拓扑在 $\widehat{G} \cup \{0\}$ 上不一致. 例如, 考虑 $\widehat{G}$ 的元素 $\chi_y \colon x \longmapsto e^{ixy}$. 我断言当 $|y| \to +\infty$ 时族 $(\chi_y)_{y \in \R}$ 弱收敛到 $0$. (显然, 它对紧收敛拓扑不收敛到 $0$; 事实上, 它在此拓扑下没有极限.) 回忆这个断言意味着对每个 $f \in L^1(\R)$,
\[
\lim_{|y| \to +\infty}\int_{\R}f(x)e^{ixy}\,\dd x = 0.
\]
首先假设 $f$ 是紧区间 $[a, b]$ 的特征函数. 则
\[
\int_{\R}f(x)e^{ixy}\,\dd y = \frac{1}{y}(e^{iby} - e^{iay}) \xrightarrow[|y| \to +\infty]{} 0.
\]
所以, 若 $f$ 是紧区间特征函数的 (有限) 线性组合, 结论仍成立. 现在设 $f$ 是 $L^1(\R)$ 的任意元素, $\varepsilon > 0$. 可以找到紧区间特征函数的线性组合 $g$ 使得 $\|f - g\|_1 \leq \varepsilon$. 由刚才看到的, 还可以找到 $A \in \R$ 使得对 $|y| \geq A$ 有 $|\int_{\R}g(x)e^{ixy}\,\dd y| \leq \varepsilon$. 则若 $|y| \geq A$,
\begin{align*}
\Bigl|\int_{\R}f(x)e^{ixy}\,\dd x\Bigr| &\leq \Bigl|\int_{\R}g(x)e^{ixy}\,\dd x\Bigr| + \Bigl|\int_{\R}(f(x) - g(x))e^{ixy}\,\dd x\Bigr| \\
&\leq \varepsilon + \int_{\R}|f(x) - g(x)|\,\dd x \\
&\leq 2\varepsilon.
\end{align*}
所以 $\int_{\R}f(x)e^{ixy}\,\dd x$ 在 $|y| \to +\infty$ 时收敛到 $0$.
\end{remark}

\begin{corollary}[紧收敛意义下的稠密性]
\label{cor:III.4.5}
$E(P_1)$ 的凸包在 $P_1$ 中对紧收敛拓扑稠密.
\end{corollary}

\begin{proof}[定理 \ref{thm:III.4.3} 的证明]
我们首先证明 $P_1$ 上的紧收敛拓扑细于弱*拓扑 (这是较容易的部分). 设 $\phi \in P_1$. 设 $f \in L^1(G)$, $c > 0$, 设 $U = \{\psi \in P_1 \mid \int_G f(\phi - \psi)\,\dd\mu < c\}$. 我们要找到 $\phi$ 在紧收敛拓扑下的邻域包含于 $U$. 设 $K \subset G$ 是紧子集使得 $\int_{G \setminus K}|f|\,\dd\mu < c/3$, 设 $V = \{\psi \in P_1 \mid \sup_{x \in K}|\phi(x) - \psi(x)| \leq \frac{c}{3\|f\|_1 + 1}\}$. 则若 $\psi \in V$,
\begin{align*}
\Bigl|\int_G f(\phi - \psi)\,\dd\mu\Bigr| &\leq \Bigl|\int_K f(\phi - \psi)\,\dd\mu\Bigr| + \Bigl|\int_{G \setminus K}f(\phi - \psi)\,\dd\mu\Bigr| \\
&\leq \|f\|_1\sup_{x \in K}|\phi(x) - \psi(x)| + 2\int_{G \setminus K}|f|\,\dd\mu \\
&< c
\end{align*}
所以 $\psi \in U$ (在第二行我们用 $\|\phi\|_{\infty} = \|\psi\|_{\infty} = 1$ 这一事实).
现在证明困难的方向, 即 $P_1$ 上的弱*拓扑细于紧收敛拓扑. 设 $\phi \in P_1$, $V = \{\psi \in P_1 \mid \sup_{x \in K}|\phi(x) - \psi(x)| < c\}$, 其中 $K \subset G$ 紧, $c > 0$. 设 $\delta > 0$ 使得 $\delta + 4\sqrt{\delta} < c$. 设 $Q$ 是 $1$ 在 $G$ 中的紧邻域使得
\[
\sup_{x \in Q}|\phi(x) - 1| \leq \delta.
\]
(这样的 $Q$ 存在, 因为 $\phi$ 连续且 $\phi(1) = 1$.) 由于 $Q$ 包含开集, $\mu(Q) \neq 0$. 设 $f = \frac{1}{\mu(Q)}\mathbb{1}_Q$. 由下面的第一个引理 (应用于 $V = L^1(G)$ 和 $B = P_1$) 和 $G \to L^1(G)$, $x \longmapsto L_{x^{-1}}f$ 连续 (从而 $\{L_{x^{-1}}f \mid x \in K\} \subset L^1(G)$ 紧) 的事实, 可以找到 $P_1$ 中 $\phi$ 的弱*邻域 $U_1$ 使得对每个 $x \in K$ 和每个 $\psi \in U_1$,
\[
\Bigl|\int_G (\phi - \psi)L_{x^{-1}}f\Bigr| \leq \delta.
\]
则对每个 $x \in K$ 和每个 $\psi \in U_1$,
\begin{align*}
|f * \phi(x) - f * \psi(x)| &= \Bigl|\int_G f(xy)(\phi(y^{-1}) - \psi(y^{-1}))\,\dd y\Bigr| \\
&= \Bigl|\int_G L_{x^{-1}}f(y)(\phi(y) - \psi(y))\,\dd y\Bigr| \quad \text{(见推论 \ref{cor:III.2.6})} \\
&\leq \delta.
\end{align*}
设 $U_2 = \{\psi \in P_1 \mid \int_G (\phi - \psi)f\,\dd\mu < \delta\}$. (这是 $\phi$ 的弱*邻域.) 设 $\psi \in U_1 \cap U_2$. 则
\begin{align*}
\Bigl|\int_G (1 - \psi)f\,\dd\mu\Bigr| &\leq \Bigl|\int_G (1 - \phi)f\,\dd\mu\Bigr| + \Bigl|\int_G (\phi - \psi)f\,\dd\mu\Bigr| \\
&\leq \frac{1}{\mu(Q)}\int_Q (1 - \phi(x))\,\dd x + \delta \\
&\leq \sup_{x \in Q}|1 - \phi(x)| + \delta \\
&\leq 2\delta.
\end{align*}
另一方面, 对每个 $x \in G$,
\begin{align*}
|f * \psi(x) - \psi(x)| &= \Bigl|\frac{1}{\mu(Q)}\int_G \mathbb{1}_Q(y)\psi(y^{-1}x)\,\dd y - \frac{1}{\mu(Q)}\int_Q\psi(x)\,\dd y\Bigr| \\
&= \Bigl|\frac{1}{\mu(Q)}\int_Q(\psi(y^{-1}x) - \psi(x))\,\dd y\Bigr| \\
&\leq \frac{1}{\mu(Q)}\int_Q|\psi(y^{-1}x) - \psi(x)|\,\dd y \\
&\leq \frac{1}{\mu(Q)}\int_Q\sqrt{2(1 - \Re(\psi(y)))}\,\dd y \quad \text{(见下面的第二个引理)} \\
&\leq \frac{\sqrt{2}}{\mu(Q)}\Bigl(\int_Q(1 - \Re(\psi(y)))\,\dd y\Bigr)^{1/2}\Bigl(\int_Q \,\dd y\Bigr)^{1/2} \quad \text{(Cauchy-Schwarz)} \\
&\leq \frac{\sqrt{2}}{\sqrt{\mu(Q)}}\Bigl(\int_Q(1 - \psi(y))\,\dd y\Bigr)^{1/2} \\
&= \sqrt{2}\Bigl(\int_G(1 - \psi)f\,\dd\mu\Bigr)^{1/2}
\end{align*}
由于 $\psi \in U_2$, 前面的计算表明这 $\leq \sqrt{2}\sqrt{2\delta} = 2\sqrt{\delta}$. 注意这也适用于 $\psi = \phi$, 因为显然 $\phi \in U_1 \cap U_2$. 把所有这些估计放在一起, 若 $\psi \in U_1 \cap U_2$, $x \in K$, 得到
\begin{align*}
|\psi(x) - \phi(x)| &\leq |\psi(x) - f * \psi(x)| + |f * \psi(x) - f * \phi(x)| + |f * \phi(x) - \phi(x)| \\
&\leq \delta + 4\sqrt{\delta} \\
&< c.
\end{align*}
所以 $U_1 \cap U_2 \subset V$, 我们完成了.
\end{proof}

\begin{lemma}[有界线性泛函族的拓扑]
\label{lem:III.4.6}
设 $V$ 是巴拿赫空间, $B$ 是 $\Hom(V, \CC)$ 的范数有界子集. 则 $B$ 上的紧收敛拓扑 (即在 $V$ 的紧子集上的一致收敛拓扑) 与弱*拓扑一致.
\end{lemma}

\begin{proof}
我们要比较 $V$ 上的逐点收敛拓扑 (即弱*拓扑) 与 $V$ 上的紧收敛拓扑. 第二个在 $\Hom(V, \CC)$ 的全体上细于第一个, 所以我们只需证明第一个在 $B$ 上细于第二个.
设 $T_0 \in B$, $K \subset V$ 紧, $c > 0$. 我们要找到 $B$ 中 $T_0$ 的弱*邻域包含于 $\{T \in B \mid \sup_{x \in K}|T(x) - T_0(x)| < c\}$. 设 $M = \sup_{T \in B}\|T\|_{\mathrm{op}}$ (它有限, 因为 $B$ 有界). 设 $x_1, \ldots, x_n \in K$ 使得 $K$ 包含于以 $x_i$ 为心、半径 $\frac{c}{3M}$ 的开球的并. 设 $T \in B$ 满足对 $i = 1, \ldots, n$ 有 $|(T - T_0)(x_i)| < c/3$ (这定义 $T$ 的弱*邻域). 对每个 $x \in K$, 存在 $i \in \{1, \ldots, n\}$ 使得 $\|x - x_i\| < \frac{c}{3M}$, 然后
\begin{align*}
|T(x) - T_0(x)| &\leq |T(x - x_i)| + |(T - T_0)(x_i)| + |T_0(x - x_i)| \\
&\leq \|T\|_{\mathrm{op}}\|x - x_i\| + c/3 + \|T_0\|_{\mathrm{op}}\|x - x_i\| \\
&< c/3 + 2M\frac{c}{3M} = c.
\end{align*}
所以 $T \in U$.
\end{proof}

\begin{lemma}[正定型函数的一致连续估计]
\label{lem:III.4.7}
设 $\phi \in P_1$. 则对所有 $x, y \in G$,
\[
|\phi(x) - \phi(y)|^2 \leq 2 - 2\Re(\phi(yx^{-1})).
\]
\end{lemma}

\begin{proof}
由定理 \ref{thm:III.2.5}, 可以找到 $G$ 的酉表示 $(\pi, V)$ 和 $v \in V$ 使得对每个 $x \in G$ 有 $\phi(x) = \langle \pi(x)(v), v \rangle$. 又, 由于 $\phi(1) = 1$, 有 $\|v\| = 1$. 所以对所有 $x, y \in G$,
\begin{align*}
|\phi(x) - \phi(y)|^2 &= |\langle (\pi(x) - \pi(y))(v), v \rangle|^2 \\
&= |\langle v, (\pi(x^{-1}) - \pi(y^{-1}))(v) \rangle|^2 \\
&\leq \|(\pi(x^{-1}) - \pi(y^{-1}))(v)\|^2 \\
&= \|\pi(x^{-1})(v)\|^2 + \|\pi(x^{-1})(v)\|^2 - 2\Re(\langle \pi(x^{-1})(v), \pi(y^{-1})(v) \rangle) \\
&= 2 - 2\Re(\langle \pi(x^{-1})(v), \pi(y^{-1})(v) \rangle) \\
&= 2 - 2\Re(\langle \pi(yx^{-1})(v), v \rangle) \\
&= 2 - 2\Re(\phi(yx^{-1})).
\end{align*}
\end{proof}

\section{Gelfand-Raikov 定理}

\begin{theorem}[Gelfand-Raikov 定理]
\label{thm:III.5.1}
设 $G$ 是局部紧群. 则对所有满足 $x \neq y$ 的 $x, y \in G$, 存在 $G$ 的不可约酉表示 $\pi$ 使得 $\pi(x) \neq \pi(y)$.
更精确地, 存在 $G$ 的不可约酉表示 $(\pi, V)$ 和向量 $v \in V$ 使得 $\langle \pi(x)(v), v \rangle \neq \langle \pi(y)(v), v \rangle$.
\end{theorem}

\begin{proof}
设 $x, y \in G$. 假设对 $G$ 的每个不可约酉表示 $(\pi, V)$ 和每个 $v \in V$ 都有 $\langle \pi(x)(v), v \rangle = \langle \pi(y)(v), v \rangle$. 由定理 \ref{thm:III.3.2}, 这蕴含对每个 $\phi \in E(P_1)$ 有 $\phi(x) = \phi(y)$, 从而由推论 \ref{cor:III.4.5} (以及 $\{x, y\}$ 是 $G$ 的紧子集这一事实) 对每个 $\phi \in P_1$ 有 $\phi(x) = \phi(y)$, 从而因 $P = \R_{\geq 0} \cdot P_1$ 对每个 $\phi \in P$ 成立.
设 $\pi_L$ 是 $G$ 的左正则表示, 即 $G$ 在 $L^2(G)$ 上由 $\pi_L(z)(f) = L_zf$ ($z \in G$, $f \in L^2(G)$) 定义的表示. 这是 $G$ 的酉表示, 所以由第一段和命题 \ref{prop:III.2.4}, 对每个 $f \in L^2(G)$ 有 $\langle \pi_L(x)(f), f \rangle = \langle \pi_L(y)(f), f \rangle$. 设 $f_1, f_2 \in L^2(G)$. 则
\[
\langle \pi_L(x)(f_1 + f_2), f_1 + f_2 \rangle = \langle \pi_L(x)(f_1), f_1 \rangle + \langle \pi_L(f_2), f_2 \rangle + \langle \pi_L(x)(f_1), f_2 \rangle + \langle \pi_L(x)(f_2), f_1 \rangle
\]
且
\[
\langle \pi_L(x)(f_1 + if_2), f_1 + if_2 \rangle = \langle \pi_L(x)(f_1), f_1 \rangle + \langle \pi_L(f_2), f_2 \rangle - i\langle \pi_L(x)(f_1), f_2 \rangle + i\langle \pi_L(x)(f_2), f_1 \rangle,
\]
所以
\begin{multline*}
2\langle \pi_L(x)(f_1), f_2 \rangle = \langle \pi_L(x)(f_1 + f_2), f_1 + f_2 \rangle + i\langle \pi_L(x)(f_1 + if_2), f_1 + if_2 \rangle \\
- (1 + i)(\langle \pi_L(x)(f_1), f_1 \rangle + \langle \pi_L(x)(f_2), f_2 \rangle).
\end{multline*}
对 $\pi_L(y)$ 有类似的恒等式, 这证明
\[
\langle \pi_L(x)(f_1), f_2 \rangle = \langle \pi_L(y)(f_1), f_2 \rangle.
\]
现在注意
\begin{align*}
\langle \pi_L(x)(f_1), f_2 \rangle &= \int_G L_xf_1(z)\overline{f_2(z)}\,\dd z \\
&= \int_G f_1(x^{-1}z)\overline{f_2(z)}\,\dd z \\
&= \int_G \overline{f_2(z)}\tilde{f}_1(z^{-1}x)\,\dd z \\
&= f_2 * \tilde{f}_1(x)
\end{align*}
(回忆 $\tilde{f}_1 \in L^2(G)$ 由 $\tilde{f}_1(z) = \overline{f_1(z^{-1})}$ 定义), 所以 $f_2 * \tilde{f}_1(x) = f_2 * \tilde{f}_1(y)$. 此计算也表明 $f_2 * \tilde{f}_1$ 有意义且连续.
由于 $f \longmapsto \tilde{f}$ 是 $L^2(G)$ 上的对合, 我们推出对所有 $f_1, f_2 \in L^2(G)$ 有 $f_1 * f_2(x) = f_1 * f_2(y)$, 特别地对所有 $f_1, f_2 \in \mathcal{C}_c(G)$ 成立. 设 $f \in \mathcal{C}_c(G)$, $(\psi_U)_{U \in \mathcal{U}}$ 是近似单位元. 则对每个 $U \in \mathcal{U}$ 有 $\psi_U \in \mathcal{C}_c(G)$, 且由命题 \ref{prop:I.4.1.9} (以及 $f$ 一致连续这一事实, 即命题 \ref{prop:I.1.12}), $\psi_U * f \xrightarrow[U \to \{1\}]{} f$ (对 $\|\cdot\|_{\infty}$). 由于对每个 $U \in \mathcal{U}$ 有 $\psi_U * f(x) = \psi_U * f(y)$, 这蕴含 $f(x) = f(y)$. 但此时必有 $x = y$ (由 Urysohn 引理).
\end{proof}

\section{习题}
设 $G$ 是拓扑群, $(\pi, V)$ 是 $G$ 的酉表示. $\pi$ 的矩阵系数 (matrix coefficient) 是形如 $x \longmapsto \langle \pi(x)(v), w \rangle$ ($v, w \in V$) 的函数 $G \to \CC$. 注意这些函数都连续.
还要回忆若 $G$ 局部紧 (且 $\mu$ 是 $G$ 上的左 Haar 测度), 则左正则表示 $\pi_L$ 是 $G$ 在 $L^2(G) := L^2(G, \mu)$ 上由 $\pi_L(x)(f) = L_xf$ ($x \in G$, $f \in L^2(G)$) 给出的表示. 在本节中, 我们直接把 $\pi_L$ 称为 $G$ 的正则表示.

\subsection{正则表示}

\begin{exercise}[离散群的情形]
\label{ex:III.6.1.1}
设 $G$ 是离散群, $\phi = \mathbb{1}_{\{1\}}$.
\begin{enumerate}
\item[(a)] (1) 证明 $\phi$ 是 $G$ 上的正定型函数.
\item[(b)] (2) 证明 $V_{\phi}$ 等价于 $G$ 的正则表示.
\end{enumerate}
\end{exercise}

\begin{solution}
(a). 计数测度是 $G$ 上的左 Haar 测度, 所以我们用此测度. 对每个 $f \in L^1(G)$, $\int_G f\phi\,\dd\mu = f(1)$. 所以
\[
\int_G (f^* * f)\phi\,\dd\mu = (f^* * f)(1) = \sum_{y \in G}\overline{f(y^{-1})}f(y^{-1}) \in \R_{\geq 0}.
\]
(b). 对所有 $f, g \in L^1(G)$,
\begin{align*}
\langle f, g \rangle_{\phi} &= \int_G (g^* * f)\phi\,\dd\mu \\
&= (g^* * f)(1) \\
&= \sum_{y \in G}\overline{g(y^{-1})}f(y^{-1}) \\
&= \langle f, g \rangle_{L^2(G)}.
\end{align*}
所以 $\langle \cdot, \cdot \rangle_{\phi}$ 的核等于 $\{0\}$, 且希尔伯特空间 $V_{\phi}$ 是 $L^1(G)$ 对范数 $\|\cdot\|_2$ 的完备化, 即 $L^2(G)$. $G$ 在 $V_{\phi}$ 上的作用是按连续性延拓的它在 $L^1(G)$ 上通过左平移的作用, 所以我们得到 $G$ 在 $L^2(G)$ 上通过左平移的作用.
\end{solution}

\begin{exercise}[正则表示的矩阵系数]
\label{ex:III.6.1.2}
设 $G$ 是局部紧群.
\begin{enumerate}
\item[(a)] 若 $f, g \in \mathcal{C}_c(G)$, 证明 $f * g \in \mathcal{C}_c(G)$.
\item[(b)] 证明 $G$ 的正则表示的每个矩阵系数都在无穷远处消失.
\item[(c)] 假设 $G$ 不紧. 若 $(\pi, V)$ 是 $G$ 的有限维酉表示, 证明它有不在无穷远处消失的矩阵系数.
\item[(d)] 若 $G$ 不紧, 证明其正则表示没有有限维子表示.\footnote{与 Peter-Weyl 比较.}
\end{enumerate}
\end{exercise}

\begin{solution}
(a). 首先, 我们知道 $f * g$ 存在, 因为 $f$ 和 $g$ 在 $L^1(G)$ 中. 若 $x, x' \in G$, 则
\[
|f * g(x) - f * g(x')| = \Bigl|\int_G f(y)(g(y^{-1}x) - g(y^{-1}x'))\,\dd y\Bigr| \leq \|f\|_1\sup_{y \in \operatorname{supp}(f)}|g(y^{-1}x) - g(y^{-1}x')|.
\]
由于 $g$ 右一致连续 (见命题 \ref{prop:I.1.12}), 当 $x'$ 趋于 $x$ 时这趋于 $0$, 所以 $f * g$ 连续.
设 $x \in G$ 使得 $f * g(x) \neq 0$. 有
\[
f * g(x) = \int_G f(y)g(y^{-1}x)\,\dd y,
\]
所以存在 $y \in \operatorname{supp}(f)$ 使得 $y^{-1}x \in \operatorname{supp}(g)$. 换句话说, $x \in \operatorname{supp}(f)\operatorname{supp}(g)$. 由于 $\operatorname{supp}(f)$ 和 $\operatorname{supp}(g)$ 都紧, 其乘积 $\operatorname{supp}(f)\operatorname{supp}(g)$ 也紧, 所以 $f * g$ 有紧支撑.
(b). 回忆若 $f \in L^2(G)$, 则定义 $\tilde{f} \colon G \to \CC$ 为 $\tilde{f}(x) = \overline{f(x^{-1})}$.
$G$ 的左正则表示的每个矩阵系数都具有
\[
x \longmapsto \langle L_{x^{-1}}f, g \rangle_{L^2(G)}
\]
的形式, 其中 $f, g \in L^2(G)$. 有 (见命题 \ref{prop:III.2.4}(iii))
\begin{align*}
\langle L_{x^{-1}}f, g \rangle_{L^2(G)} &= \int_G f(x^{-1}y)\overline{g(y)}\,\dd y \\
&= \int_G \overline{g(y)}\tilde{f}(y^{-1}x)\,\dd y \\
&= g * \tilde{f}(x).
\end{align*}
而且, 若 $f', g' \in L^2(G)$, 则 (利用 Cauchy-Schwarz 不等式)
\begin{align*}
|\langle L_{x^{-1}}f, g \rangle_{L^2(G)} - \langle L_{x^{-1}}f', g \rangle_{L^2(G)}| &\leq \int_G |g(y)||f(x^{-1}y) - f'(x^{-1}y)|\,\dd y \\
&\leq \|g\|_2\|L_{x^{-1}}(f - f')\|_2 \\
&= \|g\|_2\|f - f'\|_2
\end{align*}
且
\begin{align*}
|\langle L_{x^{-1}}f, g \rangle_{L^2(G)} - \langle L_{x^{-1}}f, g' \rangle_{L^2(G)}| &\leq \int_G |g(y) - g'(y)||f(x^{-1}y)|\,\dd y \\
&\leq \|g - g'\|_2\|L_{x^{-1}}f\|_2 \\
&= \|g - g'\|_2\|f\|_2.
\end{align*}
假设 $f, g \in \mathcal{C}_c(G)$. 则 $g, \tilde{f} \in \mathcal{C}_c(G)$, 所以由 (a), $g * \tilde{f} \in \mathcal{C}_c(G)$, 特别地此函数在无穷远处消失.
在一般情形, 设 $\varepsilon > 0$, 设 $f', g' \in \mathcal{C}_c(G)$ 满足 $\|f - f'\|_2 \leq \varepsilon$, $\|g - g'\|_2 \leq \varepsilon$. 则由上面的两个不等式, 对每个 $x \in G$,
\[
|g * \tilde{f}(x) - g' * \tilde{f}'(x)| \leq \varepsilon(\|f\|_2 + \|g\|_2).
\]
但我们刚看到 $g' * \tilde{f}'$ 有紧支撑, 所以存在 $G$ 的紧子集 $K$ 使得对每个 $x \notin K$,
\[
|g * \tilde{f}(x)| \leq \varepsilon(\|f\|_2 + \|g\|_2).
\]
这证明矩阵系数 $g * \tilde{f}$ 在无穷远处消失.
(c). 设 $(e_1, \ldots, e_n)$ 是 $V$ 的标准正交基. 对每个 $i \in \{1, \ldots, n\}$, 设 $f_i$ 是矩阵系数 $x \longmapsto \langle \pi(x)(e_1), e_i \rangle$. 则对每个 $x \in G$,
\[
\sum_{i=1}^n|f_i(x)|^2 = \sum_{i=1}^n|\langle \pi(x)(e_1), e_i \rangle|^2 = \|\pi(x)(e_1)\|^2 = 1.
\]
这证明 $f_i$ 中至少有一个不在无穷远处消失.
(d). 这直接由 (c) 和 (b) 得出.
\end{solution}

\subsection{\index{弱包含}弱包含}

设 $G$ 是拓扑群, $(\pi, V)$ 是 $G$ 的酉表示. 若 $v = w$, 则称函数 $G \to \CC$, $x \longmapsto \langle \pi(x)(v), w \rangle$ ($v, w \in V$) 为\index{对角矩阵系数}对角矩阵系数 (diagonal matrix coefficient); 由命题 \ref{prop:III.2.4}, 对角矩阵系数是正定型函数, 我们称它为与 $\pi$ 关联的正定型函数. 若正定型函数具有 $x \longmapsto \langle \pi(x)(v), v \rangle$ ($\|v\| = 1$) 的形式, 则称它\index{规范化矩阵系数}规范化 (normalized). 记 $P(\pi)$ 为与 $\pi$ 关联的正定型函数的集合.

我们很快将看到,\footnote{加参考文献.} 若 $G$ 紧, 则 $G$ 的正则表示包含 $G$ 的所有不可约表示 (它们都有限维); 事实上, 它是其所有不可约子表示的直和的闭包. 另一方面, 若 $G$ 交换, 则其正则表示是所有不可约表示 (它们都 $1$ 维) 的直积分, 尽管若 $G$ 不紧它不包含其中任何一个. 我们不会在这里严格定义直积分, 但将引入一个更弱的包含定义, 对这个定义, 交换局部紧群的不可约表示包含于正则表示, 并开始研究它.
设 $(\pi, V)$ 和 $(\pi', V')$ 是 $G$ 的酉表示. 若 $P(\pi)$ 包含于 $P(\pi')$ 的元素的有限和的集合对 $G$ 的紧子集上的收敛拓扑的闭包, 则称 $\pi$ 弱包含于 (weakly contained in) $\pi'$, 记 $\pi \prec \pi'$. 换句话说, $\pi \prec \pi'$ 若: 对每个 $v \in V$, 每个紧 $K \subset G$ 和每个 $c > 0$, 存在 $v'_1, \ldots, v'_n \in V'$ 使得
\[
\sup_{x \in K}\Bigl|\langle \pi(x)(v), v \rangle - \sum_{i=1}^n\langle \pi'(x)(v'_i), v'_i \rangle\Bigr| < c.
\]

\begin{exercise}[张量积表示]
\label{ex:III.6.2.1}
设 $(\pi_1, V_1)$, $(\pi_2, V_2)$ 是 $G$ 的酉表示.
\begin{enumerate}
\item[(a)] 证明代数张量积 $V_1 \otimes_{\CC}V_2$ 有 Hermitian 内积, 由 $\langle v_1 \otimes v_2, w_1 \otimes w_2 \rangle = \langle v_1, w_1 \rangle\langle v_2, w_2 \rangle$ 唯一确定.
\item[(b)] 记 $V_1 \hat{\otimes}_{\CC}V_2$ 为 $V_1 \otimes_{\CC}V_2$ 对此内积的完备化. 证明公式 $(x, v_1 \otimes v_2) \longmapsto \pi_1(x)(v_1) \otimes \pi_2(x)(v_2)$ 定义 $G$ 在 $V_1 \hat{\otimes}_{\CC}V_2$ 上的酉表示. (这称为张量积表示 (tensor product representation), 通常记为 $\pi_1 \otimes \pi_2$.)
\item[(c)] 若 $V_1$ 和 $V_2$ 有限维, 证明对每个 $x \in G$,
\[
\operatorname{Tr}(\pi_1 \otimes \pi_2(x)) = \operatorname{Tr}(\pi_1(x))\operatorname{Tr}(\pi_2(x)).
\]
\end{enumerate}
\end{exercise}

\begin{solution}
(a). 由于纯张量张成 $V_1 \otimes_{\CC}V_2$, $V_1 \otimes_{\CC}V_2$ 上至多有一个满足 $B(v_1 \otimes v_2, w_1 \otimes w_2) = \langle v_1, w_1 \rangle\langle v_2, w_2 \rangle$ 的半双线性形式 $B$. 让我们证明这样的形式存在. 设 $w_1 \in V_1$, $w_2 \in V_2$. 则 $V_1 \times V_2 \to \CC$, $(v_1, v_2) \longmapsto \langle v_1, w_1 \rangle\langle v_2, w_2 \rangle$ 是双线性形式, 所以它对应于 $V_1 \otimes_{\CC}V_2$ 上的唯一线性形式, 记为 $B_{w_1,w_2}$. 其次, $V_1 \times V_2$ 上把 $(w_1, w_2)$ 送到反线性形式 $v \longmapsto B_{w_1,w_2}(v)$ 的映射是双线性的, 所以它对应于 $V_1 \otimes_{\CC}V_2$ 上的唯一线性泛函 $T$. 最后, 映射 $B \colon (V_1 \otimes_{\CC}V_2) \times (V_1 \otimes_{\CC}V_2) \to \CC$, $(v, w) \longmapsto T(w)(v)$ 对 $v$ 线性、对 $w$ 反线性, 所以它是半双线性形式, 且由定义它在纯张量上取我们想要的值.
现在我们证明 $B$ Hermitian, 即对所有 $v, w \in V_1 \otimes_{\CC}V_2$ 有 $B(w, v) = \overline{B(v, w)}$. 由于 $B$ 半双线性, 只需对纯张量 $v$ 和 $w$ 验证此性质, 但此时它由 $V_1$ 和 $V_2$ 的内积的相应性质立即得出.
最后, 我们证明 $B$ 正定. 设 $v \in V_1 \otimes_{\CC}V_2$, 把 $v$ 写成 $v = \sum_{i=1}^n v_{1,i} \otimes v_{2,i}$, $v_{1,i} \in V_1$, $v_{2,i} \in V_2$. 则 $v$ 在 $V'_1 \otimes_{\CC}V'_2$ 中, 其中 $V'_1 = \operatorname{Span}(v_{1,1}, \ldots, v_{1,n})$, $V'_2 = \operatorname{Span}(v_{2,1}, \ldots, v_{2,n})$. 所以可以用 $V'_1$ 和 $V'_2$ 代替 $V_1$ 和 $V_2$, 从而可以假设 $V_1$ 和 $V_2$ 有限维. 若如此, 设 $(e_1, \ldots, e_r)$ (相应地 $(e'_1, \ldots, e'_s)$) 是 $V_1$ (相应地 $V_2$) 的标准正交基. 则 $(e_i \otimes e'_j)_{1 \leq i \leq r, 1 \leq j \leq s}$ 是 $V_1 \otimes_{\CC}V_2$ 的基, 且由 $V_1 \otimes_{\CC}V_2$ 上 Hermitian 形式 $B$ 的定义显然它是 $B$ 的标准正交基. 但标准正交基的存在迫使形式正定 (若 $v = \sum_{i=1}^r\sum_{j=1}^s a_{ij}e_i \otimes e'_j$, 则 $B(v, v) = \sum_{i,j}|a_{ij}|^2$).

(b). 首先注意若 $v_1 \in V_1$, $v_2 \in V_2$, 则 $\|v_1 \otimes v_2\| = \|v_1\|\|v_2\|$.
设 $x \in G$. 则映射 $V_1 \times V_2 \to V_1 \otimes_{\CC}V_2$, $(v_1, v_2) \longmapsto \pi_1(x)(v_1) \otimes \pi_2(x)(v_2)$ 双线性, 所以它诱导 $V_1 \otimes_{\CC}V_2$ 到自身的 $\CC$-线性映射 $\pi_1 \otimes \pi_2(x)$. 我们证明此映射等距 (从而连续). 设 $(e_i)_{i \in I}$ (相应地 $(f_j)_{j \in J}$) 是 $V_1$ (相应地 $V_2$) 的 Hilbert 基. 若 $v_1 \in V_1$, $v_2 \in V_2$, 可以把 $v_1$ 写成 $v_1 = \sum_{i \in I}a_ie_i$, $v_2 = \sum_{j \in J}b_jf_j$, 然后由上面的注记, 级数 $\sum_{i,j}a_ib_j e_i \otimes f_j$ 在 $V_1 \otimes_{\CC}V_2$ 中收敛到 $v_1 \otimes v_2$. 由于 $V_1 \otimes_{\CC}V_2$ 的每个元素都是 $v_1 \otimes v_2$ 形式元素的有限和, 这证明 $V_1 \otimes_{\CC}V_2$ 的每个元素 $v$ 都可以写成收敛级数 $\sum_{i \in I, j \in J}a_ib_je_i \otimes f_j$ ($a_i, b_j \in \CC$) 的极限. 则 $\pi_1 \otimes \pi_2(x)(v) = \sum_{i,j}a_ib_j\pi_1(x)(e_i) \otimes \pi_2(x)(f_j)$. 由于族 $(e_i \otimes f_j)$ 和 $(\pi_1(x)(e_i) \otimes \pi_2(x)(f_j))$ 在 $V_1 \otimes_{\CC}V_2$ 中都正交, 得到 $\|v\|^2 = \sum_{i,j}|a_i|^2|b_j|^2 = \|\pi_1 \otimes \pi_2(x)(v)\|^2$.
由于映射 $\pi_1 \otimes \pi_2(x)$ 连续, 它延拓成 $V_1 \hat{\otimes}_{\CC}V_2$ 的连续自同态, 它也是等距, 仍记为 $\pi_1 \otimes \pi_2(x)$.
若 $y$ 是 $G$ 的另一个元素, $V_1 \hat{\otimes}_{\CC}V_2$ 的自同态 $\pi_1 \otimes \pi_2(xy)$ 和 $(\pi_1 \otimes \pi_2(x)) \circ (\pi_1 \otimes \pi_2(y))$ 在纯张量上相等, 所以它们相等, 因为纯张量生成 $V_1 \hat{\otimes}_{\CC}V_2$ 的稠密子空间.
为验证这定义 $G$ 在 $V_1 \hat{\otimes}_{\CC}V_2$ 上的酉表示, 还需验证对每个 $v \in V_1 \hat{\otimes}_{\CC}V_2$, 映射 $G \to V_1 \hat{\otimes}_{\CC}V_2$, $x \longmapsto \pi_1 \otimes \pi_2(x)(v)$ 连续. 这对纯张量 $v$ 成立: 若 $v = v_1 \otimes v_2$, 则对 $x, y \in G$,
\begin{align*}
\|(\pi_1 \otimes \pi_2(x) - \pi_1 \otimes \pi_2(y))(v)\| &\leq \|\pi_1(x)(v_1) \otimes (\pi_2(x) - \pi_2(y))(v_2)\| + \|(\pi_1(x) - \pi_1(y))(v_1) \otimes \pi_2(y)(v_2)\| \\
&= \|v_1\|\|(\pi_2(x) - \pi_2(y))(v_2)\| + \|(\pi_1(x) - \pi_1(y))(v_1)\|\|v_2\|,
\end{align*}
这蕴含结果. 所以它对纯张量的有限和仍成立, 然后一个标准的论证表明它对 $V_1 \hat{\otimes}_{\CC}V_2$ 的每个元素成立.
(c). 设 $(e_1, \ldots, e_n)$ (相应地 $(f_1, \ldots, f_m)$) 是 $V_1$ (相应地 $V_2$) 的标准正交基. 则 $(e_i \otimes f_j)_{1 \leq i \leq n, 1 \leq j \leq m}$ 是 $V_1 \otimes_{\CC}V_2$ 的标准正交基. 设 $x \in G$. 则
\[
\operatorname{Tr}(\pi_1(x)) = \sum_{i=1}^n\langle \pi_1(x)(e_i), e_i \rangle
\]
且
\[
\operatorname{Tr}(\pi_2(x)) = \sum_{j=1}^m\langle \pi_2(x)(f_j), f_j \rangle,
\]
所以
\begin{align*}
\operatorname{Tr}(\pi_1 \otimes \pi_2(x)) &= \sum_{i=1}^n\sum_{j=1}^m\langle \pi_1 \otimes \pi_2(x)(e_i \otimes f_j), e_i \otimes f_j \rangle \\
&= \Bigl(\sum_{i=1}^n\langle \pi_1(x)(e_i), e_i \rangle\Bigr)\Bigl(\sum_{j=1}^m\langle \pi_2(x)(f_j), f_j \rangle\Bigr) \\
&= \operatorname{Tr}(\pi_1(x))\operatorname{Tr}(\pi_2(x)).
\end{align*}
\end{solution}

\begin{exercise}[Milman 定理]
\label{ex:III.6.2.2}
设 $V$ 是局部凸拓扑 $\CC$-向量空间, $K$ 是 $V$ 的紧凸子集, $F \subset K$ 满足 $K$ 是 $F$ 的凸包的闭包. 证明 $K$ 的每个端点都在 $F$ 的闭包中. (这称为 Milman 定理.)
\end{exercise}

\begin{solution}
若 $X \subset V$ 是开凸子集, 则 $X \subset 2X$. 事实上, 若 $p \colon V \to \R_{\geq 0}$ 是 $X$ 的规范函数 (见引理 \ref{lem:B.3.8}), 则 $X = \{v \in V \mid p(v) < 1\}$, 所以
\[
X \subset \{v \in V \mid p(v) \leq 1\} \subset \{v \in V \mid p(v) < 2\} = 2X.
\]
设 $v$ 是 $K$ 的端点, 假设 $v \notin \overline{F}$. 则可以找到 $V$ 中 $0$ 的凸邻域 $X$ 使得 $X = -X$ 且 $(v + X) \cap \overline{F} = \emptyset$. 用 $\frac{1}{2}X$ 代替 $X$ 后, 可以假设 $(v + X) \cap \overline{F} = \emptyset$.
由于 $F$ 紧 (作为 $K$ 的闭子集), 可以找到 $x_1, \ldots, x_n \in F$ 使得 $F \subset \bigcup_{i=1}^n(x_i + X)$. 对每个 $i \in \{1, \ldots, n\}$, 设 $K_i$ 是 $F \cap (x_i + X)$ 的凸包的闭包; 这是 $V$ 的紧凸子集 (它紧, 因为它在 $K$ 中闭). 由于 $K$ 是 $F$ 的凸包的闭包, 有 $K \supset K_1 \cup \cdots \cup K_n$, 所以 $K$ 包含 $K_1 \cup \cdots \cup K_n$ 的凸包 $L$. 让我们证明 $K = L$. 由于 $L \supset F$ 且 $L$ 凸, 只需证明 $L$ 紧. 设
\[
S = \{(x_1, \ldots, x_n) \in [0,1]^n \mid x_1 + \cdots + x_n = 1\},
\]
考虑函数
\[
f \colon S \times K_1 \times \cdots \times K_n \to L, \quad ((x_1, \ldots, x_n), v_1, \ldots, v_n) \longmapsto \sum_{i=1}^n x_iv_i.
\]
此映射连续, 所以其像紧. 若我们证明此像是凸的, 则它由 $L$ 的定义等于 $L$, 我们就完成了. 所以设 $a = ((x_1, \ldots, x_n), v_1, \ldots, v_n)$ 和 $a' = ((x'_1, \ldots, x'_n), v'_1, \ldots, v'_n)$ 是 $S \times K_1 \times \cdots \times K_n$ 的元素, $t \in [0,1]$. 则
\[
tf(a) + (1 - t)f(a') = \sum_{i=1}^n(tx_iv_i + (1 - t)x'_iv'_i)
\]
设 $i \in \{1, \ldots, n\}$. 若 $tx_i + (1 - t)x'_i \neq 0$, 令 $y_i = tx_i + (1 - t)x'_i$, $w_i = \frac{1}{y_i}(tx_iv_i + (1 - t)x'_iv'_i)$. 否则令 $w_i = v_i$, $y_i = 0$. 则对每个 $i$ 有 $w_i \in K_i$ (因为 $K_i$ 凸), 对每个 $i$ 有 $y_i \geq 0$, 且
\[
\sum_{i=1}^n y_i = t\sum_{i=1}^n x_i + (1 - t)\sum_{i=1}^n x'_i = 1.
\]
所以
\[
tf(a) + (1 - t)f(a') = f((y_1 \ldots, y_n), w_1, \ldots, w_n)
\]
在 $f$ 的像中, 我们完成了.
现在我们导出矛盾. 由于 $K = L$, 可以把 $v$ 写成 $v = \sum_{i=1}^n t_iv_i$, 其中 $(t_1, \ldots, t_n) \in S$, 对每个 $i$ 有 $v_i \in K_i$. 由于 $v$ 在 $K$ 中端点, 存在 $i \in \{1, \ldots, n\}$ 使得 $v = v_i$. 但此时 $v \in K_i \subset \overline{x_i + X}$ (因为 $K_i$ 包含于 $x_i + X$ 的凸包的闭包, 而这是 $x_i + \overline{X}$, 因为 $X$ 凸). 由于 $x_i \in F$ 且 $X = -X$, 这蕴含 $x_i \in (v + X) \cap F$, 与 $X$ 的选取矛盾.
\end{solution}

\begin{exercise}[弱包含的判定]
\label{ex:III.6.2.3}
设 $(\pi, V)$ 和 $(\pi', V')$ 是 $G$ 的酉表示. 设 $C \subset V$ 使得 $\operatorname{Span}(\pi(x)(v) \mid x \in G, v \in C)$ 在 $V$ 中稠密.\footnote{例如, 若 $V$ 循环, $C$ 可以只包含 $V$ 的一个循环向量.} 假设对每个 $v \in C$, 函数 $x \longmapsto \langle \pi(x)(v), v \rangle$ 在 $P(\pi')$ 的元素的有限和的集合的闭包中 (仍对 $G$ 的紧子集上的收敛拓扑). 本题的目标是证明这蕴含 $\pi \prec \pi'$.
设 $X$ 是满足如下条件的 $v \in V$ 的集合: $x \longmapsto \langle \pi(x)(v), v \rangle$ 在 $P(\pi')$ 的元素的有限和的集合的闭包中 (对与上面相同的拓扑).
\begin{enumerate}
\item[(a)] 证明 $X$ 被所有 $\pi(x)$ ($x \in G$) 稳定, 且被纯量乘法稳定.
\item[(b)] 若 $v \in X$, $x_1, x_2 \in G$, 证明 $\pi(x_1)(v) + \pi(x_2)(v) \in X$.
\item[(c)] 证明 $X$ 在 $V$ 中闭.
\item[(d)] 若 $v \in X$, 证明 $V$ 中包含 $v$ 的最小闭 $G$-不变子空间包含于 $X$.
\item[(e)] 设 $v_1, v_2 \in X$, 设 $W_1$ (相应地 $W_2$) 是 $V$ 中包含 $v_1$ (相应地 $v_2$) 的最小闭 $G$-不变子空间. 设 $W = W_1 + W_2$, 记 $T \colon W \to W_1^{\perp}$ 为正交投影, 其中 $W_1^{\perp}$ 在 $W$ 中取.
  \begin{enumerate}
  \item[(i)] 证明 $T$ 是 $G$-等变的且 $T(W_2)$ 在 $W_1^{\perp}$ 中稠密.
  \item[(ii)] 证明 $W_1^{\perp} \subset X$.
  \item[(iii)] 证明 $v_1 + v_2 \in X$. (提示: 用 $T(v_1 + v_2)$ 和 $(v_1 + v_2) - T(v_1 + v_2)$.)
  \end{enumerate}
\item[(f)] 证明 $\pi \prec \pi'$.
\end{enumerate}
\end{exercise}

\begin{solution}
(a). 对每个 $v \in V$ (相应地 $v' \in V'$), 记 $\phi_v$ 为矩阵系数 $x \longmapsto \langle \pi(x)(v), v \rangle$ (相应地 $x \longmapsto \langle \pi'(x)(v), v \rangle$). 又记 $\sum P(\pi')$ 为 $P(\pi')$ 的元素的有限和的集合.
设 $v \in X$, $y \in G$, $\lambda \in \CC$. 我们要证明 $\pi(y)(v)$ 和 $\lambda v$ 在 $X$ 中, 即 $\phi_{\pi(y)(v)}$ 和 $\phi_{\lambda v}$ 在 $\overline{\sum P(\pi')}$ 中. 若 $\lambda = 0$, 对 $\lambda v$ 结论显然 (注意 $0$ 是 $G$ 的每个表示的矩阵系数), 所以可以假设 $\lambda \neq 0$. 设 $K$ 是 $G$ 的紧子集, $\varepsilon > 0$. 选取 $v'_1, \ldots, v'_n$ 使得 $\sup_{x \in K \cup y^{-1}Ky}|\phi_v(x) - \sum_{i=1}^n\phi_{v'_i}(x)| \leq \min(\varepsilon, |\lambda|^{-2}\varepsilon)$. 则对每个 $x \in K$,
\begin{align*}
\Bigl|\phi_{\pi(y)v}(x) - \sum_{i=1}^n\phi_{\pi'(y)(v'_i)}(x)\Bigr| &= \Bigl|\langle \pi(xy)(v), \pi(y)(v) \rangle - \sum_{i=1}^n\langle \pi'(xy)(v'_i), \pi'(y)(v'_i) \rangle\Bigr| \\
&= \Bigl|\langle \pi(y^{-1}xy)(v), v \rangle - \sum_{i=1}^n\langle \pi'(y^{-1}xy)(v'_i), v'_i \rangle\Bigr| \\
&= \Bigl|\phi_v(y^{-1}xy) - \sum_{i=1}^n\phi_{v'_i}(y^{-1}xy)\Bigr| \\
&\leq \varepsilon
\end{align*}
且
\[
\Bigl|\phi_{\lambda v}(x) - \sum_{i=1}^n\phi_{\lambda v'_i}(x)\Bigr| = |\lambda|^2\Bigl|\phi_v(x) - \sum_{i=1}^n\phi_{v'_i}(x)\Bigr| \leq \varepsilon.
\]
所以 $\phi_{\pi(y)(v)}$ 和 $\phi_{\lambda v}$ 在 $\overline{\sum P(\pi')}$ 中.
(b). 设 $v \in X$, $x_1, x_2 \in G$. 对每个 $y \in G$,
\begin{align*}
\phi_{\pi(x_1)(v) + \pi(x_2)(v)}(y) &= \langle \pi(y)(\pi(x_1)(v) + \pi(x_2)(v)), \pi(x_1)(v) + \pi(x_2)(v) \rangle \\
&= \langle \pi(x_1^{-1}yx_1)(v), v \rangle + \langle \pi(x_2^{-1}yx_1)(v), v \rangle + \langle \pi(x_1^{-1}yx_2)(v), v \rangle + \langle \pi(x_2^{-1}yx_2)(v), v \rangle.
\end{align*}
换句话说,
\[
\phi_{\pi(x_1)(v) + \pi(x_2)(v)} = \sum_{i,j=1}^2 L_{x_i}R_{x_j}\phi_v.
\]
设 $K$ 是 $G$ 的紧子集, $\varepsilon > 0$. 选取 $v'_1, \ldots, v'_n \in V'$ 使得
\[
\sup_{x \in \bigcup_{i,j=1}^2 x_iKx_j^{-1}}\Bigl|\phi_v(x) - \sum_{i=1}^n\phi_{v'_i}(x)\Bigr| \leq \varepsilon.
\]
则由上面的计算 (以及它对函数 $\phi_{v'_i}$ 的类似), 对每个 $x \in K$,
\[
\Bigl|\phi_{\pi(x_1)(v) + \pi(x_2)(v)}(x) - \sum_{i=1}^n\phi_{\pi(x_1)(v_i) + \pi(x_2)(v_i)}(x)\Bigr| \leq \varepsilon.
\]
这证明 $\phi_{\pi(x_1)(v) + \pi(x_2)(v)}$ 在 $\overline{\sum P(\pi')}$ 中, 即 $\pi(x_1)(v) + \pi(x_2)(v) \in X$.
(c). 只需证明若对 $P(\pi)$ 用紧收敛拓扑, 则映射 $V \to P(\pi)$, $v \longmapsto \phi_v$ 连续. 设 $v, v' \in V$. 则对每个 $x \in G$,
\begin{align*}
|\phi_v(x) - \phi_{v'}(x)| &= |\langle \pi(x)(v), v \rangle - \langle \pi(x)(v'), v' \rangle| \\
&\leq |\langle \pi(x)(v - v'), v \rangle| + |\langle \pi(x)(v'), v - v' \rangle| \\
&\leq \|v - v'\|\|v\| + \|v'\|\|v - v'\|.
\end{align*}
所以映射 $v \longmapsto \phi_v$ 甚至对 $\|\cdot\|_{\infty}$ 给出的 $P(\pi)$ 上的拓扑连续.
(d). 设 $v \in X$. 由 (a) 和 (b), 对每个 $n \geq 1$ 和所有 $\lambda_1, \ldots, \lambda_n \in \CC$, $x_1, \ldots, x_n \in G$, 有 $\sum_{i=1}^n\lambda_i\pi(x_i)(v) \in X$. 所以 $V$ 中包含 $v$ 的最小 $G$-不变子空间 (即 $\sum_{x \in G}\pi(x)(\CC v)$) 包含于 $X$. 结论现在由 (c) 得出.
(e). (i) 由于 $W_1$ 是 $G$-不变的, 算子 $T$ 由引理 \ref{lem:I.3.4.3} 是 $G$-等变的. 由于 $W = W_1 + W_2$, $W_1 + W_2$ 在 $T$ 下的像在 $\Im(T) = W_1^{\perp}$ 中稠密. 由于 $\Ker(T) = W_1$, 有 $T(W_1 + W_2) = T(W_2)$, 所以 $T(W_2)$ 在 $W_1^{\perp}$ 中稠密.
(ii) 由于 $W = W_1 \oplus W_1^{\perp}$, 推出 $T(W_1^{\perp} \cap W_2) = W_1^{\perp} \cap W_2$ 在 $W_1^{\perp}$ 中稠密. 由于 $W_2 \subset X$, (c) 蕴含 $W_1^{\perp} \subset X$.
(iii) 令 $v = T(v_1 + v_2)$, $w = v_1 + v_2 - v$. 则 $v \in W_1^{\perp} \subset X$, $w \in \Ker(T) = W_1 \subset X$, 所以 $v, w \in X$. 另一方面, 对每个 $x \in G$,
\begin{align*}
\phi_{v_1 + v_2}(x) &= \langle \pi(x)(v_1 + v_2), v_1 + v_2 \rangle \\
&= \langle \pi(x)(v + w), v + w \rangle \\
&= \phi_v(x) + \phi_w(x).
\end{align*}
由于 $\overline{\sum P(\pi')}$ 被和稳定, 这蕴含 $v_1 + v_2 \in X$.
(f). 由 (a), (c) 和 (e), 集合 $X$ 是 $V$ 的闭 $G$-不变子空间, 所以由对 $C$ 的假设它等于 $V$. 这意味着 $\pi \prec \pi'$.
\end{solution}

\begin{exercise}[弱包含与弱*闭包]
\label{ex:III.6.2.4}
设 $(\pi, V)$ 和 $(\pi', V')$ 是满足 $\pi \prec \pi'$ 的两个酉表示. 设 $\mathcal{C}$ 是 $L^{\infty}(G)$ 上弱*拓扑中与 $\pi'$ 关联的规范化正定型函数的集合的凸包的闭包.
\begin{enumerate}
\item[(a)] 证明与 $\pi$ 关联的每个规范化正定型函数都在 $\mathcal{C}$ 中.
\item[(b)] 若 $\pi$ 不可约, 证明与 $\pi$ 关联的每个规范化正定型函数都是与 $\pi'$ 关联的规范化正定型函数在 $G$ 的紧子集上的收敛拓扑下的极限. (提示: 习题 \ref{ex:III.6.2.2}.)
\item[(c)] 若 $\pi$ 是 $G$ 的平凡表示, 证明对每个紧子集 $K \subset G$ 和每个 $c > 0$, 存在 $v' \in V'$ 使得 $\|v'\| = 1$ 且
\[
\sup_{x \in K}\|\pi'(x)(v') - v'\| < c.
\]
\item[(d)] 反之, 假设对每个紧子集 $K \subset G$ 和每个 $c > 0$, 存在 $v' \in V'$ 使得 $\|v'\| = 1$ 且
\[
\sup_{x \in K}\|\pi'(x)(v') - v'\| < c.
\]
证明平凡表示弱包含于 $\pi'$.
\end{enumerate}
\end{exercise}

\begin{solution}
(a). 设 $\phi$ 是与 $\pi$ 关联的规范化正定型函数. 设 $f \in L^1(G)$, $\varepsilon > 0$. 我们要找到与 $\pi'$ 关联的规范化正定型函数的凸组合 $\psi$ 使得 $|\int_G f(\phi - \psi)\,\dd\mu| \leq \varepsilon$. 选取 $\delta > 0$; 稍后看它需要多小. 设 $K \ni 1$ 是 $G$ 的紧子集使得 $\int_{G-K}|f|\,\dd\mu \leq \delta$. 由于 $\pi \prec \pi'$, 可以找到 $v_1, \ldots, v_n \in V'$ 使得 $\sup_{x \in K}|\phi(x) - \sum_{i=1}^n\phi_{v_i}(x)| \leq \delta$. 特别地, 在 $1$ 处取值, 得到 $|1 - \sum_{i=1}^n\|v_i\|^2| \leq \delta$. 设 $c_i = \|v_i\|^2$, $c = c_1 + \cdots + c_n$, $\phi_i = \frac{1}{c_i}\phi_{v_i} = \phi_{\frac{1}{\|v_i\|}v_i}$, $\psi = \frac{1}{c}\sum_{i=1}^n\phi_{v_i} = \frac{1}{c}\sum_{i=1}^n c_i\phi_i$. 则 $\phi_1, \ldots, \phi_n$ 是与 $\pi'$ 关联的规范化正定型函数, 且 $\psi$ 是 $\phi_1, \ldots, \phi_n$ 的凸组合. 特别地, $\|\psi\|_{\infty} \leq 1 = \|\phi\|_{\infty}$.
对每个 $x \in K$,
\[
|\phi(x) - \psi(x)| \leq \Bigl|\phi(x) - \sum_{i=1}^n\phi_{v_i}(x)\Bigr| + |1 - c||\psi(x)| \leq 2\delta.
\]
所以
\begin{align*}
\Bigl|\int_G f(\phi - \psi)\,\dd\mu\Bigr| &\leq \sup_{x \in K}|\phi(x) - \psi(x)|\int_K |f|\,\dd\mu + \sup_{x \in G-K}|\phi(x) - \psi(x)|\int_{G-K}|f|\,\dd\mu \\
&\leq 2\delta\|f\|_1 + 2\delta.
\end{align*}
取 $\delta$ 足够小可使这 $\leq \varepsilon$.
(b). 设 $F$ 是与 $\pi'$ 关联的规范化正定型函数的集合, $K$ 是其凸包的弱*闭包. 则 $F$ 包含于所有规范化正定型函数的凸集 $P_1$ 中, 所以 $K \subset P_1$. 设 $\phi$ 是与 $\pi$ 关联的规范化正定型函数. 由 (a), $\phi \in K$. 由定理 \ref{thm:III.3.2}, 函数 $\phi$ 在 $P_1$ 中端点, 从而在 $K$ 中也端点. 由习题 \ref{ex:III.6.2.2}, 这蕴含 $\phi$ 在 $F$ 的弱*拓扑闭包中. 但 $F$ 和 $\phi$ 在 $P_1$ 中, 且 $P_1$ 上的弱*拓扑与 $G$ 的紧子集上的收敛拓扑一致 (由 Raikov 定理, 即定理 \ref{thm:III.4.3}), 所以 $\phi$ 也在 $F$ 对 $G$ 的紧子集上的收敛拓扑的闭包中.
(c). 由于 $\pi$ 是平凡表示, 与 $\pi$ 关联的唯一规范化正定型函数是常值函数 $1$. 由 (c), 存在 $v' \in V'$ 使得 $\|v'\| = 1$ 且
\[
\sup_{x \in K}|1 - \langle \pi'(x)(v'), v' \rangle| \leq c^2/3.
\]
设 $x \in G$. 则
\[
\|\pi'(x)(v') - v'\|^2 = \|\pi'(x)(v')\|^2 + \|v'\|^2 - 2\Re(\langle \pi'(x)(v'), v' \rangle) \leq 2|1 - \langle \pi'(x)(v'), v' \rangle| \leq 2c^2/3,
\]
所以
\[
\sup_{x \in K}\|\pi'(x)(v') - v'\| < c.
\]
(d). 设 $\pi$ 是 $G$ 的平凡表示. 则 $P(\pi)$ 是非负常值函数的集合, 所以为证明 $\pi \prec \pi'$, 只需证明常值函数 $1$ 是 $P(\pi')$ 的元素的有限和的极限 (在 $G$ 的紧子集上的收敛拓扑中). 设 $K$ 是 $G$ 的紧子集, $c > 0$. 选取 $v' \in V'$ 使得 $\|v'\| = 1$ 且 $\sup_{x \in K}\|\pi'(x)(v') - v'\| < c$, 定义 $\phi'$ 为 $\phi'(x) = \langle \pi'(x)(v'), v' \rangle$. 则对每个 $x \in K$,
\[
|1 - \phi'(x)| = |\langle v', v' \rangle - \langle \pi'(x)(v'), v' \rangle| = |\langle v' - \pi'(x)(v'), v' \rangle| \leq \|v' - \pi'(x)(v')\| < c.
\]
\end{solution}

\begin{exercise}[有限生成离散群]
\label{ex:III.6.2.5}
设 $G$ 是有限生成离散群, $S$ 是 $G$ 的有限生成元集. 证明 $G$ 的平凡表示弱包含于 $G$ 的正则表示当且仅当: 对每个 $\varepsilon > 0$, 存在 $f \in L^2(G)$ 使得
\[
\sup_{x \in S}\|L_xf - f\|_2 < \varepsilon\|f\|_2.
\]
\end{exercise}

\begin{solution}
我们用习题 \ref{ex:III.6.2.4}(c) 和 (d) 的判定, 它说 $G$ 的平凡表示弱包含于正则表示当且仅当: 对 $G$ 的每个紧 (即有限) 子集 $K$ 和每个 $\varepsilon > 0$, 存在 $f \in L^2(G)$ 使得 $\|f\|_2 = 1$ 且
\[
\sup_{x \in K}\|L_xf - f\|_2 < \varepsilon.
\]
首先, 由于 $S$ 有限, 立即看到若平凡表示包含于正则表示, 则题设的条件满足.
反之, 假设题设的条件满足. 设 $K$ 是 $G$ 的有限子集, $\varepsilon > 0$. 设 $T = S \cup S^{-1} \cup \{1\}$. 有 $G = \bigcup_{n \geq 1}T^n$, 因为 $S$ 生成 $G$, 且这是递增并. 由于 $K$ 有限, 存在 $n \geq 1$ 使得 $K \subset T^n$. 由假设, 可以找到 $f \in L^2(G)$ 使得 $\|f\|_2 = 1$ 且
\[
\sup_{x \in S}\|L_xf - f\|_2 \leq \frac{1}{n}\varepsilon.
\]
我们要证明
\[
\sup_{x \in K}\|L_xf - f\|_2 \leq \varepsilon.
\]
只需对 $\sup_{x \in T^n}$ 证明. 设 $x \in T^n$, 把 $x$ 写成 $x = x_1\cdots x_n$, 其中 $x_1, \ldots, x_n \in T$. 我们对 $i \in \{1, \ldots, n\}$ 归纳证明 $\|L_{x_1\cdots x_i}f - f\|_2 \leq \frac{i}{n}\varepsilon$. 若 $i = 1$, 要证明 $\|L_{x_1}f - f\|_2 \leq \frac{1}{n}\varepsilon$. 若 $x_1 \in S$, 这由 $f$ 的选取成立; 若 $x_1 = 1$ 显然; 若 $x_1 \in S^{-1}$, 它由
\[
\|L_{x_1}f - f\|_2 = \|f - L_{x_1^{-1}}f\|_2
\]
得出. 现在假设对 $i \in \{1, \ldots, n-1\}$ 已知结果, 证明对 $i+1$. 有
\begin{align*}
\|L_{x_1\cdots x_{i+1}}f - f\|_2 &\leq \|L_{x_1\cdots x_i}(L_{x_{i+1}}f - f)\|_2 + \|L_{x_1\cdots x_i}f - f\|_2 \\
&= \|L_{x_{i+1}}f - f\|_2 + \|L_{x_1\cdots x_i}f - f\|_2 \\
&\leq \frac{i}{n}\varepsilon + \frac{1}{n}\varepsilon = \frac{i+1}{n}\varepsilon.
\end{align*}
\end{solution}

\begin{exercise}[$G = \Z$ 的情形]
\label{ex:III.6.2.6}
设 $G = \Z$. 证明 $G$ 的平凡表示弱包含于 $G$ 的正则表示.
\end{exercise}

\begin{solution}
我们应用习题 \ref{ex:III.6.2.5} 的结果, 取 $S = \{1\}$. 所以对每个 $\varepsilon > 0$, 必须找到 $f \in L^2(\Z)$ 使得 $\|f\|_2 = 1$ 且 $\|L_1f - f\|_2 \leq \varepsilon$. 第一个条件说 $\sum_{n \in \Z}|f(n)|^2 = 1$, 第二个条件说 $\sum_{n \in \Z}|f(n-1) - f(n)|^2 \leq \varepsilon^2$. 设 $N \in \Z_{\geq 0}$, 考虑函数 $g_N = \mathbb{1}_{[0,N]} \in L^2(\Z)$. 则 $\|g_N\|_2^2 = N + 1$, 且 $\sum_{n \in \Z}|g(n-1) - g(n)|^2 = 2$. 所以若 $f_N = \frac{1}{\sqrt{N+1}}g_N$, 有 $\|f\|_2 = 1$ 且 $\|L_1f - f\|_2 = \sqrt{\frac{2}{N+1}}$. 取 $N$ 足够大, 看到 $f_N$ 具有所要的性质.
\end{solution}

\begin{exercise}[$G = \R$ 的情形]
\label{ex:III.6.2.7}
设 $G = \R$.
\begin{enumerate}
\item[(a)] 证明 $G$ 的平凡表示弱包含于 $G$ 的正则表示.
\item[(b)] 证明 $G$ 的每个不可约酉表示弱包含于 $G$ 的正则表示.\footnote{我们将在后面看到这对每个交换局部紧群都成立.}\footnote{在哪里?}
\end{enumerate}
\end{exercise}

\begin{solution}
(a). 若 $a, b \in \R$ 满足 $a < b$, 设 $f = (b-a)^{-1/2}\mathbb{1}_{[a,b]}$. 则 $f \in L^2(\R)$ 且 $\|f\|_2 = 1$. 而且, 对每个 $t \in \R$, $L_tf = (b-a)^{-1/2}\mathbb{1}_{[a+t,b+t]}$, 所以
\[
\|L_tf - f\|_2^2 \leq \frac{2|t|}{b-a}.
\]
设 $K$ 是 $\R$ 的紧子集, $\varepsilon > 0$. 若选取 $a, b \in \R$ 使得 $b - a \geq 2\varepsilon^{-2}\sup_{t \in K}|t|$, 则上面的构造给出 $f \in L^2(\R)$ 使得 $\|f\|_2 = 1$ 且 $\sup_{x \in K}\|L_tf - f\|_2 \leq \varepsilon$. 由习题 \ref{ex:III.6.2.4}(d), $\R$ 的平凡表示包含于其正则表示.
(b). 由于 $\R$ 交换, 由 Schur 引理 (定理 \ref{thm:I.3.4.1}), 每个不可约酉表示都是 $1$ 维的. 设 $\chi \colon \R \to S^1$ 是这种表示. 设 $K$ 是 $\R$ 的紧子集, $\varepsilon > 0$. 由 (a), 存在 $f \in L^2(\R)$ 使得 $\|f\|_2 = 1$ 且 $\sup_{t \in K}\|L_tf - f\|_2 \leq \varepsilon$. 设 $g = \chi f$. 则对每个 $t \in \R$,
\[
\langle L_tg, g \rangle_{L^2(\R)} = \int_{\R}g(x-t)\overline{g(x)}\,\dd x = \chi(t)\langle L_tf, f \rangle_{L^2(\R)},
\]
所以
\[
|\chi(t) - \langle L_tg, g \rangle_{L^2(\R)}| = |1 - \langle L_tf, f \rangle_{L^2(\R)}| = |\langle f - L_tf, f \rangle_{L^2(\R)}| \leq \|L_tf - f\|_2.
\]
所以
\[
\sup_{t \in K}|\chi(t) - \langle L_tg, g \rangle_{L^2(\R)}| \leq \varepsilon.
\]
这由习题 \ref{ex:III.6.2.4}(d) 蕴含所要的结果.
\end{solution}

\begin{exercise}[自由群的情形]\index{乒乓引理}
\label{ex:III.6.2.8}
设 $G$ 是两个生成元上的 (非交换) 自由群, 赋予离散拓扑. 证明 $G$ 的平凡表示不弱包含于 $G$ 的正则表示.\footnote{应在下一节.}
\end{exercise}

\begin{solution}
设 $a, b \in G$ 是 $G$ 的两个生成元, 设 $S = \{1, a, b, a^{-1}, b^{-1}\}$. 有 $G = \bigcup_{n \geq 1}S^n$, 且这是递增并. 假设 $G$ 的平凡表示弱包含于正则表示. 则由习题 \ref{ex:III.6.2.4}(c), 对每个 $n \geq 1$, 存在 $f_n \in L^2(G)$ 使得 $\|f_n\|_2 = 1$ 且
\[
\sup_{x \in S^n}\|L_xf_n - f_n\|_2 \leq \frac{1}{n}.
\]
设 $g_n = |f_n|^2$. 则 $g_n \in L^1(G)$, $\|g_n\|_1 = 1$, 且对每个 $x \in S^n$, Cauchy-Schwarz 不等式给出
\[
\|L_xg_n - g_n\|_1 \leq \|L_xf_n - f_n\|_2\|L_xf_n + f_n\|_2 \leq \frac{2}{n}.
\]
对每个 $n \geq 1$, 定义 $L^{\infty}(G)$ 上的连续线性泛函 $\Lambda_n$ 为 $\Lambda_n(\phi) = \sum_{x \in G}g_n(x)\phi(x)$. 则 $\|\Lambda_n\|_{\mathrm{op}} = \|g_n\|_1 = 1$, 所以由 Banach-Alaoglu 定理, $(\Lambda_n)_{n \geq 1}$ 有对 $\Hom(L^{\infty}(G), \CC)$ 上弱*拓扑收敛的子序列 $(\Lambda_{n_k})_{k \geq 0}$. 设 $\Lambda$ 是它的极限. 设 $\phi \in L^{\infty}(G)$. 有
\[
\Lambda(\phi) = \lim_{k \to +\infty}\Lambda_{n_k}(\phi).
\]
设 $y \in G$. 存在 $n \geq 1$ 使得 $y^{-1} \in S^n$. 则若 $k$ 满足 $n_k \geq n$,
\begin{align*}
|\Lambda_{n_k}(L_y\phi) - \Lambda_{n_k}(\phi)| &= \Bigl|\sum_{x \in G}L_{y^{-1}}g_{n_k}(x)\phi(x) - \sum_{x \in G}g_{n_k}(x)\phi(x)\Bigr| \\
&\leq \|L_{y^{-1}}g_{n_k} - g_{n_k}\|_1\|\phi\|_{\infty} \\
&\leq \frac{2}{n_k}\|\phi\|_{\infty}.
\end{align*}
令 $k \to +\infty$ 取极限, 看到 $\Lambda(L_y\phi) = \Lambda(\phi)$. 又注意 $\Lambda(1) = 1$, 且若 $\phi$ 取非负值则 $\Lambda(\phi) \geq 0$.
回忆 $G$ 的每个元素都可以唯一地写成 $a$, $b$, $a^{-1}$, $b^{-1}$ 上的既约字. 设 $A$ 是既约表达式以 $a$ 的非零幂开始的 $G$ 的元素集合. 则对每个 $x \in G$, 若 $x \notin A$, 有 $a^{-1}x \in A$, 从而 $x \in aA$. 换句话说, $G = A \cup aA$, 所以 $\mathbb{1}_A + \mathbb{1}_{aA} - \mathbb{1}_G$ 取非负值, 因此
\[
\Lambda(\mathbb{1}_A) = \frac{1}{2}(\Lambda(\mathbb{1}_A) + \Lambda(\mathbb{1}_{aA})) \geq \frac{1}{2}\Lambda(\mathbb{1}_G) = \frac{1}{2}.
\]
另一方面, 群 $G$ 是子集 $b^nA$, $n \in \Z$ 的不相交并, 所以特别地
\[
1 = \Lambda(\mathbb{1}_G) \geq \Lambda(\mathbb{1}_A) + \Lambda(\mathbb{1}_{bA}) + \Lambda(\mathbb{1}_{b^2A}) = 3\Lambda(\mathbb{1}_A),
\]
即 $\Lambda(\mathbb{1}_A) \leq \frac{1}{3}$. 所以我们得到矛盾.
\end{solution}

\begin{exercise}[张量积与弱包含]
\label{ex:III.6.2.9}
若 $\pi_1, \pi_2, \pi'_1, \pi'_2$ 是 $G$ 的酉表示且 $\pi_1 \prec \pi'_1$, $\pi_2 \prec \pi'_2$, 证明 $\pi_1 \otimes \pi_2 \prec \pi'_1 \otimes \pi'_2$.
\end{exercise}

\begin{solution}
我们使用习题 \ref{ex:III.6.2.3} 解答中对正定型函数 $\phi_v$ 的相同记号. 对 $i = 1, 2$, 记 $V_i$ (相应地 $V'_i$) 为 $\pi_i$ (相应地 $\pi'_i$) 的空间.
若 $v_1 \in V_1$, $v_2 \in V_2$, 则由 $V_1 \otimes_{\CC}V_2$ 上内积的定义, $\phi_{v_1 \otimes v_2} = \phi_{v_1}\phi_{v_2}$. 对 $V'_1 \otimes_{\CC}V'_2$ 中的纯张量有类似结果. 所以 $\phi_{v_1 \otimes v_2}$ 在 $\overline{\sum P(\pi'_1 \otimes \pi'_2)}$ 中. 由于 $\overline{\sum P(\pi'_1 \otimes \pi'_2)}$ 被有限和稳定, 且 $V_1 \otimes_{\CC}V_2$ 的每个元素都可以写成两两正交纯张量的有限和 (见 (a) 的证明), 这蕴含对 $V_1 \otimes_{\CC}V_2$ 的每个元素 $v$, $\phi_v \in \overline{\sum P(\pi'_1 \otimes \pi'_2)}$. 最后, 我们在 (c) 中证明了映射 $v \longmapsto \phi_v$ 连续, 且 $V_1 \otimes_{\CC}V_2$ 在 $V_1 \hat{\otimes}_{\CC}V_2$ 中稠密, 所以对 $V_1 \hat{\otimes}_{\CC}V_2$ 的每个元素 $v$, $\phi_v \in \overline{\sum P(\pi'_1 \otimes \pi'_2)}$.
\end{solution}

\begin{exercise}[$G$ 离散时 $\pi \otimes \pi_L$ 的等价]
\label{ex:III.6.2.10}
假设 $G$ 离散. 对每个 $x \in G$, 记 $\delta_x \in L^2(G)$ 为 $\{x\}$ 的特征函数.
设 $(\pi, V)$ 是 $G$ 的酉表示, $(\pi_0, V)$ 是 $G$ 在 $V$ 上的平凡表示 (即对每个 $x \in G$ 有 $\pi_0(x) = \mathrm{id}_V$).
\begin{enumerate}
\item[(a)] 证明公式 $v \otimes f \longmapsto \sum_{x \in G}f(x)(\pi(x)^{-1}(v)) \otimes \delta_x$ 给出从 $V \hat{\otimes}_{\CC}L^2(G)$ 到自身的良定义且连续的 $\CC$-线性变换.
\item[(b)] 证明表示 $\pi \otimes \pi_L$ 和 $\pi_0 \otimes \pi_L$ 等价 (回忆 $\pi_L$ 是 $G$ 的左正则表示).
\end{enumerate}
\end{exercise}

\begin{solution}
(a). 首先, 映射 $V \times L^2(G) \to V \otimes_{\CC}L^2(G)$, $(v, f) \longmapsto \sum_{x \in G}f(x)(\pi(x)^{-1}(v)) \otimes \delta_x$ 双线性, 所以它定义线性映射 $\alpha \colon V \otimes_{\CC}L^2(G) \to V \otimes_{\CC}L^2(G)$. 对每个 $v, v' \in V$ 和 $f, f' \in L^2(G)$, (注意对 $V$ 的元素的每个族 $(v_x)_{x \in G}$, 族 $(v_x \otimes \delta_x)_{x \in G}$ 正交)
\begin{align*}
\langle \alpha(v \otimes f), \alpha(v' \otimes f') \rangle &= \sum_{x \in G}f(x)\overline{f'(x)}\langle \pi(x)^{-1}(v), \pi(x)^{-1}(v') \rangle \\
&= \sum_{x \in G}f(x)\overline{f'(x)}\langle v, v' \rangle \\
&= \langle v \otimes f, v' \otimes f' \rangle.
\end{align*}
利用 $V \otimes_{\CC}L^2(G)$ 的每个元素都可以写成两两正交纯张量的有限和 (见 (a) 的证明) 这一事实, 这蕴含对 $V \otimes_{\CC}L^2(G)$ 的每个元素 $v$ 有 $\|\alpha(v)\| = \|v\|$. 特别地, $\alpha$ 连续, 所以它延拓成 $V \hat{\otimes}_{\CC}L^2(G)$ 的连续自同态, 它仍是等距.
(b). 我们仍记 (a) 中构造的 $V \hat{\otimes}_{\CC}L^2(G)$ 的自同态为 $\alpha$. 我们证明它是从 $\pi \otimes \pi_L$ 到 $\pi_0 \otimes \pi_L$ 的 $G$-等变映射. 由于纯张量生成 $V \hat{\otimes}_{\CC}L^2(G)$ 的稠密子空间, 只需在纯张量上验证 $G$-等变性. 所以设 $v \in V$, $f \in L^2(G)$, $x \in G$. 有
\[
\alpha(\pi \otimes \pi_L(x)(v \otimes f)) = \alpha(\pi(x)(v) \otimes L_xf) = \sum_{y \in G}f(x^{-1}y)\pi(y^{-1}x)(v) \otimes \delta_y.
\]
另一方面,
\begin{align*}
\pi_0 \otimes \pi_L(x)(\alpha(v \otimes f)) &= \pi_0 \otimes \pi_L(x)\Bigl(\sum_{y \in G}f(y)\pi(y)^{-1}(v) \otimes \delta_y\Bigr) \\
&= \sum_{y \in G}f(y)\pi(y)^{-1}(v) \otimes L_x\delta_y \\
&= \sum_{y \in G}f(y)\pi(y)^{-1}(v) \otimes \delta_{xy} \\
&= \sum_{y \in G}f(x^{-1}z)\pi(z^{-1}x)(v) \otimes \delta_z \\
&= \alpha(\pi \otimes \pi_L(x)(v \otimes f)).
\end{align*}
我们还需验证 $\alpha$ 是向量空间的同构. 这由它有逆 $\beta$ 得出, $\beta$ 由公式 $\beta(v \otimes f) = \sum_{x \in G}f(x)\pi(x)(v) \otimes \delta_x$ 给出. (可以像 (a) 那样验证 $\beta$ 良定义且连续, 然后在纯张量上验证它是 $\alpha$ 的逆, 这是容易的验证.)
注意 $\pi \otimes \pi_L$ 和 $\pi_0 \otimes \pi_L$ 之间的同构是等距, 所以这些表示有相同的正定型函数.
\end{solution}

\begin{exercise}[非离散群的情形]
\label{ex:III.6.2.11}
把习题 \ref{ex:III.6.2.10}(b) 的结果推广到非离散局部紧群.
\end{exercise}

\begin{solution}
设 $(\pi, V)$ 是 $G$ 的酉表示. 记 $V_0$ 为带 $G$ 的平凡作用的 $V$.
首先我们定义带 $G$ 的酉作用的希尔伯特空间 $L^2(G, V_0)$. (这也常记为 $\mathrm{Ind}^G_{\{1\}}V_0$.) 考虑从 $G$ 到 $V_0$ 的具紧支撑的连续函数空间 $\mathcal{C}_c(G, V_0)$, 带范数 $\|\cdot\|_{\infty}$ (由 $\|f\|_{\infty} = \sup_{x \in G}\|f(x)\|$ 定义). 令 $G$ 通过 $(x, f) \longmapsto L_xf$ ($x \in G$, $f \in \mathcal{C}_c(G, V_0)$) 作用于这个空间. 看命题 \ref{prop:I.1.12}, 看到它的证明可以推广到从 $G$ 到 $V_0$ 的函数, 并表明 $\mathcal{C}_c(G, V_0)$ 的每个元素都左、右一致连续. 特别地, 对每个 $f \in \mathcal{C}_c(G, V_0)$, 映射 $G \to \mathcal{C}_c(G, V_0)$, $x \longmapsto L_xf$ 连续.
现在在 $\mathcal{C}_c(G, V_0)$ 上定义 Hermitian 半双线性形式:
\[
\langle f, g \rangle = \int_G \langle f(x), g(x) \rangle_{V_0}\,\dd x.
\]
容易看出这是内积, 且 $G$ 在 $\mathcal{C}_c(G, V_0)$ 上的作用保持此内积, 并对 $\mathcal{C}_c(G, V_0)$ 上由相应范数定义的拓扑在第一个变量 $x \in G$ 上连续. 记 $L^2(G, V_0)$ 为 $\mathcal{C}_c(G, V_0)$ 对 $\langle \cdot, \cdot \rangle$ 的完备化. 这是希尔伯特空间, 我们像 $V_0 = \CC$ 的情形那样证明 $G$ 在 $\mathcal{C}_c(G, V_0)$ 上的作用延拓成 $G$ 在 $L^2(G, V_0)$ 上的酉作用.
现在我们构造 $G$-等变等距 $V \hat{\otimes}_{\CC}L^2(G) \to L^2(G, V_0)$. 考虑映射 $V \times \mathcal{C}_c(G) \to L^2(G, V_0)$, 把 $(v, f)$ 送到函数 $x \longmapsto f(x)\pi(x^{-1})(v)$. 这是双线性映射, 所以它诱导 $\CC$-线性算子 $\alpha \colon V \otimes_{\CC}\mathcal{C}_c(G) \to L^2(G, V_0)$. 我们验证 $\alpha$ 是 $G$-等变的. 只需在纯张量上验证, 因为它们生成 $V \otimes_{\CC}\mathcal{C}_c(G)$. 若 $y \in G$, $v \in V$, $f \in \mathcal{C}_c(G)$, 则对每个 $x \in G$,
\[
\alpha(\pi(y)(v) \otimes L_yf)(y) = f(y^{-1}x)\pi(xy^{-1}v) = L_y(\alpha(v \otimes f))(x).
\]
我们还验证 $\alpha$ 保持内积. 如前, 由双线性只需在纯张量上验证. 设 $v, w \in V$, $f, g \in \mathcal{C}_c(G)$. 则
\begin{align*}
\langle \alpha(v \otimes f), \alpha(w \otimes g) \rangle &= \int_G \langle f(x)\pi(x)^{-1}(v), g(x)\pi(x)^{-1}(w) \rangle_{V_0}\,\dd x \\
&= \int_G f(x)\overline{g(x)}\langle v, w \rangle_{V_0}\,\dd x \\
&= \langle f, g \rangle_{L^2(G)}\langle v, w \rangle_{V_0}.
\end{align*}
这蕴含 $\alpha$ 等距, 所以它按连续性延拓成等距 $V \hat{\otimes}_{\CC}L^2(G) \to L^2(G, V_0)$ (我们用 $\mathcal{C}_c(G)$ 在 $L^2(G)$ 中稠密这一事实), 它仍是 $G$-等变的.
我们以类似 $\alpha$ 的方式定义 $G$-等变等距 $\alpha' \colon V_0 \hat{\otimes}_{\CC}L^2(G) \to L^2(G, V_0)$, 但对 $v \in V_0$ 和 $f \in \mathcal{C}_c(G)$, 取 $\alpha'(v \otimes f)$ 为函数 $x \longmapsto f(x)v$. 证明这确实定义所要的 $G$-等变等距与 $\alpha$ 的情形相同.
最后, 我们证明 $\alpha$ 和 $\alpha'$ 是同构. 我们已经知道它们单射且像闭 (因为它们是等距), 所以只需证明它们有稠密像.
设 $(e_i)_{i \in I}$ 是 $V_0$ 的 Hilbert 基. 考虑 $L^2(G, V_0)$ 的如下子空间 $W$: 其元素是具紧支撑的连续函数 $f \colon G \to V_0$, 使得存在 $J \subset I$ 有限满足 $f(G) \subset \operatorname{Span}(e_j, j \in J)$. 让我们证明 $W$ 在 $L^2(G, V_0)$ 中稠密. 只需证明 $W$ 在 $\mathcal{C}_c(G, V_0)$ 中稠密. 设 $f \in \mathcal{C}_c(G, V_0)$. 由于 $f$ 有紧支撑, $V_0$ 的子集 $f(G)$ 紧. 设 $\varepsilon > 0$. 对每个 $x \in K$, 存在 $I$ 的有限子集 $J$ 使得以 $x$ 为心、$\varepsilon$ 为半径的闭球与 $\operatorname{Span}(e_j, j \in J)$ 相交. 由于 $K$ 紧, 它可以被有限多个这样的球覆盖, 所以可以找到 $I$ 的有限子集 $J$ 使得对每个 $x \in K$, $x$ 与 $\operatorname{Span}(e_j, j \in J)$ 的距离 $\leq \varepsilon$. 换句话说, 若 $\pi_J$ 是 $\operatorname{Span}(e_j, j \in J)$ 上的正交投影, 则对每个 $x \in K$ 有 $\|\pi_J(x) - x\| \leq \varepsilon$. 则 $\pi_J \circ f \in W$, 且 $\|f - \pi_J \circ f\|_{\infty} \leq \varepsilon$, 所以 $\|f - \pi_J \circ f\|_2 \leq \operatorname{vol}(\operatorname{supp} f)\varepsilon$. 这证明 $W$ 对 $\mathcal{C}_c(G, V_0)$ 上的两个拓扑 (由 $\|\cdot\|_{\infty}$ 诱导的和由 $\|\cdot\|_2$ 诱导的; 这里只有第二个相关) 都在 $\mathcal{C}_c(G, V_0)$ 中稠密. 为完成, 只需证明 $W$ 包含于 $\alpha$ 和 $\alpha'$ 的像. 设 $f \in W$. 可以找到 $I$ 的有限子集 $J$ 使得 $f(G) \subset \operatorname{Span}(e_j, j \in J)$, 然后 $f(x) = \sum_{j \in J}f_j(x)e_j$, 其中 $f_j \in \mathcal{C}_c(G)$. (只需在 $\operatorname{Span}(e_j, j \in J)$ 的标准正交基 $(e_j)_{j \in J}$ 中取坐标.) 特别地, $f = \alpha'(\sum_{j \in J}e_j \otimes f_j)$, 所以 $f \in \Im(\alpha')$. 这证明 $\alpha'$ 是同构.
对 $\alpha$, 我们改为考虑 $\mathcal{C}_c(G, V_0)$ 的如下子空间 $W'$: 其元素 $f$ 满足存在 $J \subset I$ 有限使得对每个 $x \in G$, 向量 $\pi(x)(f(x))$ 在 $\operatorname{Span}(e_j, j \in J)$ 中. 我们如前证明 $W'$ 在 $\mathcal{C}_c(G, V_0)$ 中稠密 (对 $\|\cdot\|_{\infty}$ 和 $\|\cdot\|_2$ 两个范数): 设 $f \in \mathcal{C}_c(G, V_0)$, $\varepsilon > 0$. 由于 $f$ 有紧支撑, $V_0$ 的子集 $\{\pi(x)(f(x)) \mid x \in G\}$ 紧, 所以可以找到 $I$ 的有限子集 $J$ 使得对每个 $x \in G$, $\pi(x)(f(x))$ 与 $\operatorname{Span}(e_j, j \in J)$ 的距离至多为 $\varepsilon$. 设 $\pi_J$ 是 $\operatorname{Span}(e_j, j \in J)$ 上的正交投影, 定义 $g \in W'$ 为 $g(x) = \pi(x)^{-1}\circ\pi_J\circ\pi(x)(f(x))$. 对每个 $x \in G$,
\[
\|g(x) - f(x)\| = \|\pi(x)(g(x) - f(x))\| = \|\pi_J(\pi(x)(f(x))) - \pi(x)(f(x))\| \leq \varepsilon,
\]
所以 $\|g - f\|_{\infty} \leq \varepsilon$, $\|g - f\|_2 \leq \operatorname{vol}(\operatorname{supp} f)\varepsilon$. 最后, 我们证明 $W'$ 包含于 $\alpha$ 的像. 设 $f \in W'$, 定义 $g \in \mathcal{C}_c(G, V_0)$ 为 $g(x) = \pi(x)(f(x))$. 选取 $I$ 的有限子集 $J$ 使得 $g(G) \subset \operatorname{Span}(e_j, j \in J)$, 把 $g$ 写成 $g = \sum_{j \in J}g_je_j$, 其中 $g_j \in \mathcal{C}_c(G)$. 则对每个 $x \in G$,
\[
f(x) = \sum_{j \in J}g_j(x)\pi(x)^{-1}(e_j).
\]
换句话说, $f = \alpha(\sum_{j \in J}e_j \otimes g_j)$.
\end{solution}

\begin{exercise}[弱包含与平凡表示]
\label{ex:III.6.2.12}
证明下列条件等价:
\begin{enumerate}
\item[(i)] $G$ 的平凡表示弱包含于 $\pi_L$.
\item[(ii)] $G$ 的每个酉表示弱包含于 $\pi_L$.
\end{enumerate}
\end{exercise}

\begin{solution}
(ii)$\Rightarrow$(i) 显然. 所以让我们证明 (i)$\Rightarrow$(ii). 设 $(\pi, V)$ 是 $G$ 的酉表示, $\pi_0$ 是 $G$ 在 $V$ 上的平凡表示, $\mathbf{1}$ 是 $G$ 在 $\CC$ 上的平凡表示. 我们知道 $\mathbf{1} \prec \pi_L$, 所以由习题 \ref{ex:III.6.2.9} 和 \ref{ex:III.6.2.10},
\[
\pi \simeq \pi \otimes \mathbf{1} \prec \pi \otimes \pi_L \simeq \pi_0 \otimes \pi_L.
\]
像习题 \ref{ex:III.6.2.3} 的解答那样, 对 $G$ 的每个酉表示 $\pi'$, 记 $\overline{\sum P(\pi')}$ 为与 $\pi'$ 关联的正定型函数的有限和的集合的闭包. 让我们证明 $\overline{\sum P(\pi_L)} = \overline{\sum P(\pi_0 \otimes \pi_L)}$, 这将完成证明, 因为我们已经知道 $P(\pi) \subset \overline{\sum P(\pi_0 \otimes \pi_L)}$.
由于 $\pi_L$ 是 $\pi_0 \otimes \pi_L$ 的子表示 (对每个 $v \in V - \{0\}$, $V \hat{\otimes}_{\CC}L^2(G)$ 的子空间 $\CC v \otimes L^2(G)$ 是 $G$-不变的, 且通过映射 $v \otimes f \longmapsto f$ 等价于表示 $\pi_L$), 有 $P(\pi_L) \subset P(\pi_0 \otimes \pi_L)$, 所以 $\overline{\sum P(\pi_L)} \subset \overline{\sum P(\pi_0 \otimes \pi_L)}$. 反之, 设 $(e_i)_{i \in I}$ 是 $V$ 的标准正交基, $v \in V \hat{\otimes}_{\CC}L^2(G)$. 则可以把 $v$ 写成 $v = \sum_{i \in I}e_i \otimes f_i$, 其中和在 $V \hat{\otimes}_{\CC}L^2(G)$ 中收敛 (即 $\sum_{i \in I}\|f_i\|_2^2$ 收敛). 则对每个 $x \in G$,
\[
\langle \pi_0 \otimes \pi_L(x)(v), v \rangle = \sum_{i \in I}\langle L_xf_i, f_i \rangle_{L^2(G)},
\]
所以函数 $x \longmapsto \langle \pi_0 \otimes \pi_L(x)(v), v \rangle$ 在 $\overline{\sum P(\pi_L)}$ 中.
\end{solution}

\subsection{\index{可解群}可解群}

设 $(X, \mu)$ 是测度空间, $E$ 是 $L^{\infty}(X)$ 的包含常值函数且对映射 $\phi \longmapsto \bar{\phi}$ 封闭的闭线性子空间. $E$ 上的\index{平均}平均 (mean) 是满足下列条件的线性泛函 $M \colon E \to \CC$:
\begin{enumerate}
\item[(i)] $M(\mathbb{1}_X) = 1$;
\item[(ii)] 若 $f \geq 0$ (局部) 几乎处处, 则 $M(f) \geq 0$.
\end{enumerate}
若 $X = G$ 是局部紧群, 若对每个 $f \in E$ 和每个 $x \in G$ 有 $L_xf \in E$ 且 $M(L_xf) = M(f)$, 则称 $E$ 上的平均 $M$ 是 $G$-不变的 (G-invariant).
若 $L^{\infty}(G)$ 上存在 \index{$G$-不变平均}$G$-不变平均, 则称群 $G$ 可解 (amenable).
设 $V$ 是局部凸拓扑向量空间 (见定义 \ref{def:B.3.5}), $K$ 是 $V$ 的凸子集. 若对所有 $v, w \in K$ 和每个 $t \in [0,1]$ 有 $f(tv + (1-t)w) = tf(v) + (1-t)f(w)$, 则称映射 $f \colon K \to K$ 仿射 (affine). 设 $G \times K \to K$, $(x, v) \longmapsto x \cdot v$ 是 $G$ 在 $K$ 上的连续左作用. 若对每个 $x \in G$, 映射 $K \to K$, $v \longmapsto x \cdot v$ 仿射, 则称此作用是\index{仿射作用}仿射作用 (affine action).
若 $G$ 在局部凸拓扑向量空间的非空紧凸子集上的每个仿射作用都有不动点, 则称群 $G$ 具有\index{不动点性质}不动点性质 (fixed point property).
注意: 在这一系列习题中 Hahn-Banach 定理是你的朋友. 还有如下事实: 若 $V$ 是拓扑向量空间, 则 $\Hom(V, \CC)$ 上的每个弱*连续线性泛函都具有 $\Lambda \longmapsto \Lambda(v)$ 的形式 (对某个 $v \in V$). (见 Rudin 的书 {[}20{]} 的定理 3.10.)

\begin{exercise}[平均的连续性]
\label{ex:III.6.3.1}
\begin{enumerate}
\item[(a)] 若 $(X, \mu)$ 是测度空间, $E$ 是 $L^{\infty}(X)$ 的包含常值函数的子空间, 证明 $E$ 上的任何平均 $M$ 自动连续 (对由范数 $\|\cdot\|_{\infty}$ 给出的拓扑), 且 $\|M\|_{\mathrm{op}} = 1$.
\end{enumerate}
我们现在假设 $G$ 是局部紧群.
\begin{enumerate}
\item[(a)] 若 $G$ 紧, 证明 $\mathcal{C}(G)$ 上的左不变平均与 $G$ 上的规范化 Haar 测度自然双射对应.
\item[(b)] 设 $L^1(G)^{1,+}$ 是满足 $f \geq 0$ 几乎处处且 $\|f\|_1 = 1$ 的 $f \in L^1(G)$ 的凸子集. 我们按通常方式把 $L^1(G)$ 等同于 $L^{\infty}(G)$ 的连续对偶的子空间 (即函数 $f \in L^1(G)$ 对应于 $L^{\infty}(G)$ 上的连续线性泛函 $\phi \longmapsto \int_G f\phi\,\dd\mu$). 证明 $L^{1,+}(G)$ 在 $L^{\infty}(G)$ 上的平均集合中弱*稠密.
\item[(c)] 设 $\mathrm{UCB}(G)$ 是 $L^{\infty}(G)$ 的由 $G$ 上的左一致连续有界函数组成的子空间. 对每个 $x \in G$, 记 $\delta_x$ 为线性泛函 $\mathcal{C}(G) \to \CC$, $f \longmapsto f(x)$. 证明泛函 $\delta_x$ 的凸组合 (即和 $\sum_{i=1}^n a_i\delta_{x_i}$, 其中 $x_1, \ldots, x_n \in G$, $a_1, \ldots, a_n \in [0,1]$ 且 $a_1 + \cdots + a_n = 1$) 的集合在 $\mathrm{UCB}(G)$ 上的平均集合中弱*稠密.
\end{enumerate}
\end{exercise}

\begin{solution}
(a). 设 $M$ 是 $E$ 上的平均. 设 $\phi \in E$, 假设 $\phi(x)$ 对几乎每个 $x$ 是实数. 有 $\|\phi\|_{\infty}\mathbb{1}_X \in E$, 因为它是常值函数的倍数, 且函数 $\|\phi\|_{\infty}\mathbb{1}_X - \phi$ 和 $\|\phi\|_{\infty}\mathbb{1}_X + \phi$ 几乎处处 $\geq 0$, 所以它们在 $M$ 下的像 $\geq 0$, 即 $M(\phi) \in \R$ 且
\[
-\|\phi\|_{\infty} \leq M(\phi) \leq \|\phi\|_{\infty},
\]
即 $|M(\phi)| \leq \|\phi\|_{\infty}$.
现在设 $\phi$ 是 $E$ 的任意元素. 选取 $c \in \CC$ 使得 $|c| = 1$ 且 $M(c\phi) \in \R$. 设 $\phi_1 = \frac{1}{2}(c\phi + \overline{c\phi})$, $\phi_2 = \frac{1}{2i}(c\phi - \overline{c\phi})$. 则 $\phi_1, \phi_2$ 取实值, 且 $c\phi = \phi_1 + i\phi_2$. 有
\[
|\phi(x)| = \sqrt{\phi_1(x)^2 + \phi_2(x)^2} \geq \max(|\phi_1(x)|, |\phi_2(x)|)
\]
对每个 $x \in X$, 所以 $\|\phi\|_{\infty} \geq \max(\|\phi_1\|_{\infty}, \|\phi_2\|_{\infty})$. 另一方面, $M(c\phi) = M(\phi_1) + iM(\phi_2)$ 且 $M(\phi_1), M(\phi_2) \in \R$, 所以 $M(\phi_2) = 0$, 且
\[
|M(\phi)| = |M(c\phi)| = |M(\phi_1)| \leq \|\phi_1\|_{\infty} \leq \|\phi\|_{\infty}.
\]
这证明 $M$ 连续且 $\|M\|_{\mathrm{op}} \leq 1$. 由于 $M(\mathbb{1}_X) = 1 = \|\mathbb{1}_X\|_{\infty}$, 有 $\|M\|_{\mathrm{op}} = 1$.
(b). 这只是 Riesz 表示定理 (定理 \ref{thm:I.2.3}) 和命题 \ref{prop:I.2.6}.
(c). 设 $M$ 是 $L^{\infty}(G)$ 上的平均的集合. 它显然是 $\Hom(L^{\infty}(G), \CC)$ 的凸子集. 由 (a), 集合 $M$ 包含于 $\Hom(L^{\infty}(G), \CC)$ 的闭单位球. 又, 由于刻画平均的条件对弱*拓扑都闭, 集合 $M$ 在 $\Hom(L^{\infty}(G), \CC)$ 中弱*闭. 所以 $M$ 弱*紧.
由 $L^1(G)^{1,+}$ 的定义, 对每个 $f \in L^1(G)^{1,+}$, 对应的 $L^{\infty}(G)$ 上的线性形式是 $M$ 的元素. 又注意 $L^1(G)^{1,+}$ 是 $L^1(G)$ 的凸子集, 所以它在 $\Hom(L^{\infty}(G), \CC)$ 中的像是凸的. 设 $M'$ 是这个像的弱*闭包. 有 $M' \subset M$, 所以 $M'$ 凸且弱*紧. 假设 $M' \neq M$. 则由 Hahn-Banach 定理 (第二几何形式), 存在 $M \in M$ 和弱*连续的 $\R$-线性算子 $\Lambda \colon \Hom(L^{\infty}(G), \CC) \to \R$ 使得 $\Lambda(M) > \sup_{M' \in M'}\Lambda(M')$. 注意线性算子 $\Lambda' \colon M' \longmapsto \Lambda(M') + \frac{1}{i}\Lambda(iM')$ 弱*连续且 $\CC$-线性, 所以存在 $\phi \in L^{\infty}(G)$ 使得对每个 $M' \in \Hom(L^{\infty}(G), \CC)$ 有 $\Lambda'(M') = M'(\phi)$, 这给出 $\Lambda(M') = \Re(M'(\phi))$. 则
\[
\Re(M(\phi)) > \sup_{f \in L^1(G)^{1,+}}\Re\Bigl(\int_G f\phi\,\dd\mu\Bigr).
\]
把 $\phi$ 写成 $\phi = \phi_1 + i\phi_2$, 其中 $\phi_1 = \Re\phi$, $\phi_2 = \Im\phi$. 则
\[
M(\phi_1) > \sup_{f \in L^1(G)^{1,+}}\int_G f\phi_1\,\dd\mu
\]
(因为由 (a) 的解答, $M(\phi_1), M(\phi_2) \in \R$). 设
\[
c = \inf\{d \in \R \mid \phi_1 \leq d\mathbb{1}_G \text{ 局部几乎处处}\}.
\]
若 $\phi_1 \leq d\mathbb{1}_G$ 局部几乎处处, 则 $M(\phi_1) \leq M(d\mathbb{1}_G) = d$. 所以 $M(\phi_1) \leq c$. 设 $\delta > 0$ 使得 $M(\phi_1) - \delta > \sup_{f \in L^1(G)^{1,+}}\int_G f\phi_1\,\dd\mu$. 由 $c$ 的定义, 存在 $G$ 的可测子集 $A$ 使得 $\mu(A) > 0$ 且 $\phi_1|_A \geq (c + \delta)\mathbb{1}_A$. 设 $f = \mu(A)^{-1}\mathbb{1}_A$. 则 $f \in L^1(G)^{1,+}$ 且 $\int_G \phi_1f\,\dd\mu \geq c + \delta \geq M(\phi_1) + \delta$, 矛盾.
(d). 设 $M$ 是 $\mathrm{UCB}(G)$ 上的平均的集合. 我们像 (c) 的解答那样看到 $M$ 是 $\Hom(\mathrm{UCB}(G), \CC)$ 的凸弱*紧子集. 设 $M'$ 是 $\delta_x$ ($x \in G$) 的凸包的弱*闭包; 则 $M' \subset M$, 因为每个 $\delta_x$ 在 $M$ 中. 若 $M' \neq M$, 则由 Hahn-Banach 定理 (第二几何形式), 存在 $M$ 的元素 $M$ 和连续 $\R$-线性泛函 $\Lambda \colon \Hom(\mathrm{UCB}(G), \CC) \to \R$ 使得
\[
\Lambda(M) > \sup_{M' \in M'}\Lambda(M').
\]
像 (c) 的解答那样, 看到可以找到取实值的函数 $\phi \in \mathrm{UCB}(G)$ 使得对每个 $M' \in M$ 有 $\Lambda(M') = M'(\phi)$. 所以有
\[
M(\phi) > \sup_{M' \in M'}M'(\phi) \geq \sup_{x \in G}\delta_x(\phi) = \sup_{x \in G}\phi(x).
\]
设 $\delta > 0$ 使得 $M(\phi) - \delta \geq \sup_{x \in G}\phi(x)$. 则 $\phi \leq (M(\phi) - \delta)\mathbb{1}_G$, 从而 $M(\phi) \leq M(\phi) - \delta$, 矛盾.
\end{solution}

\begin{exercise}[可解群具有不动点性质]
\label{ex:III.6.3.2}
设 $G$ 是可解局部紧群. 本题的目标是证明 $G$ 具有不动点性质.
所以设 $V$ 是局部凸拓扑向量空间, $K$ 是 $V$ 的非空紧凸子集, $G \times K \to K$, $(x, v) \longmapsto x \cdot v$ 是连续仿射作用.
\begin{enumerate}
\item[(a)] 证明 $\mathrm{UCB}(G)$ 上存在左不变平均.
\item[(b)] 固定点 $v_0 \in K$, 定义 $t \colon G \to K$ 为 $t(x) = x \cdot v_0$. 若 $M$ 是 $\mathrm{UCB}(G)$ 上的平均, 证明存在唯一的正则 Borel 测度 $\mu_M$ 使得对每个 $f \in \mathcal{C}(K)$,
\[
\int_K f\,\dd\mu_M = M(f \circ t).
\]
\item[(c)] 证明积分 $b_M = \int_K v\,\dd\mu_M(v)$ 存在且 $b_M \in K$.
\item[(d)] 设 $M$ 是所有 $\mathrm{UCB}(G)$ 上的平均的集合, 赋予弱*拓扑 (其中 $\mathrm{UCB}(G)$ 上的拓扑由 $\|\cdot\|_{\infty}$ 给出). 证明对每个连续线性泛函 $\Lambda \colon V \to \CC$, 映射 $M \to K$, $M \longmapsto \Lambda(b_M)$ 连续.
\item[(e)] 若 $M = \delta_x$ (对某个 $x \in G$), 计算 $b_M$.
\item[(f)] 证明对每个 $M \in M$ 和每个 $x \in G$, 有 $b_{M \circ L_{x^{-1}}} = x \cdot b_M$. (提示: 习题 \ref{ex:III.6.3.1}(d). 另外, 可以假设向量值积分的形成与连续\index{仿射映射}仿射映射交换这一事实.)\footnote{删掉这个提示?}
\item[(g)] 证明 $G$ 在 $K$ 上的作用有不动点.
\end{enumerate}
\end{exercise}

\begin{solution}
(a). 只需取 $L^{\infty}(G)$ 上的左不变平均在 $\mathrm{UCB}(G)$ 上的限制.
(b). 我们首先证明对每个 $f \in \mathcal{C}(K)$, $f \circ t \in \mathrm{UCB}(G)$, 且线性算子 $\mathcal{C}(K) \to \mathrm{UCB}(G)$, $f \longmapsto f \circ t$ 连续. 所以设 $f \in \mathcal{C}(K)$. 注意函数 $t \colon G \to K$ 由假设连续, 所以 $f \circ t$ 连续. 又, 显然有 $\|f \circ t\|_{\infty} \leq \|f\|_{\infty}$. 还需证明 $f \circ t$ 左一致连续. 记 $a \colon G \times K \to K$ 为作用映射. 设 $\varepsilon > 0$. 对每个 $v \in K$, 存在 $a^{-1}(v)$ 的开邻域 $\Omega$ 使得对每个 $(x, w) \in \Omega$ 有 $|f(x \cdot w) - f(v)| < \varepsilon$; 由于 $(1, v) \in a^{-1}(v)$, 可以假设 $\Omega = U_v \times V_v$, 其中 $U_v$ 是 $1$ 在 $G$ 中的开邻域, $V_v$ 是 $v$ 在 $K$ 中的开邻域. 由于 $K$ 紧, 可以找到 $v_1, \ldots, v_n \in K$ 使得 $K = \bigcup_{i=1}^n V_{v_i}$. 设 $U = \bigcap_{i=1}^n U_{v_i}$. 设 $x \in U$, $v \in K$. 则存在 $i \in \{1, \ldots, n\}$ 使得 $v \in V_{v_i}$, 且有
\[
|f(x \cdot v) - f(v)| \leq |f(x \cdot v) - f(v_i)| + |f(v_i) - f(1 \cdot v)| < 2\varepsilon.
\]
所以对每个 $x \in U$ 和每个 $y \in G$,
\[
|f \circ t(xy) - f \circ t(y)| = |f(x \cdot t(y)) - f(t(y))| < 2\varepsilon.
\]
这证明 $f \circ t$ 一致连续.
设 $M$ 是 $\mathrm{UCB}(G)$ 上的平均. 把 $M$ 与连续线性算子 $\mathcal{C}(K) \to \mathrm{UCB}(G)$, $f \longmapsto f \circ t$ 复合, 得到 $\mathcal{C}(K)$ 上的平均. 由 Riesz 表示定理, 存在 $K$ 上唯一的正则 Borel 测度 $\mu_M$ 使得对每个 $f \in \mathcal{C}(K)$ 有 $M(f \circ t) = \int_K f\,\dd\mu_M$.
(c). 注意 $\mu(K) = \int_K 1\,\dd\mu_M = M(\mathbb{1}_G) = 1$. 函数 $\mathrm{id}_K$ 是 $K$ 上具紧支撑的连续函数, 所以由习题 \ref{ex:I.5.6.2}, 它关于 $\mu_M$ 的积分 $b_M = \int_K v\,\dd\mu_M$ 存在, 且 $\mu(K)^{-1}b_M = b_M$ 在 $K$ 的凸包的闭包中, 即在 $K$ 中.
(d). 由积分的定义, 对每个 $\Lambda \in \Hom(V, \CC)$ 和 $\mathrm{UCB}(G)$ 上的每个平均 $M$,
\[
\Lambda(b_M) = \int_G \Lambda(v)\,\dd\mu_M = M(\Lambda \circ t).
\]
这由弱*拓扑的定义对 $M$ 连续.
(e). 设 $x \in G$, $M = \delta_x$. 则对每个 $f \in \mathcal{C}_f(K)$,
\[
\int_K f\,\dd\mu_M = M(f \circ t) = f(x \cdot v_0).
\]
取 $f = \mathrm{id}_K$, 得到
\[
b_M = \int_K v\,\dd\mu_M = x \cdot v_0.
\]
(f). 设 $M$ 是 $\mathrm{UCB}(G)$ 上的平均的集合. 固定 $M \in M$. 设 $x \in G$, $\Lambda \in \Hom(V, \CC)$. 映射 $L_{x^{-1}}$ 把 $\mathrm{UCB}(G)$ 送到自身, 所以 $M \circ L_{x^{-1}}$ 有意义. 对每个 $M \in M$, 利用映射 $K \to K$, $v \longmapsto x \cdot v$ 连续且仿射这一事实, 得到
\begin{align*}
\Lambda(x \cdot b_M) &= \Lambda\Bigl(\int_K x \cdot v\,\dd\mu_M\Bigr) = \int_K \Lambda(x \cdot v)\,\dd\mu_M = M(\Lambda(x \cdot t)) \\
&= M(L_{x^{-1}}(\Lambda \circ t)) \\
&= \Lambda(b_{M \circ L_{x^{-1}}}).
\end{align*}
由于连续线性泛函分离点 (由 Hahn-Banach 定理), 这蕴含 $x \cdot b_M = b_{M \circ L_{x^{-1}}}$.
(g). 设 $M$ 是不变平均 (由 (a) 存在). 则由 (f), 点 $b_M \in K$ 是 $G$ 作用的不动点.
\end{solution}

\begin{exercise}[不动点性质蕴含可解性]
\label{ex:III.6.3.3}
设 $G$ 是局部紧群, 假设 $G$ 具有不动点性质. 本题的目标是证明 $G$ 可解. (你可能会发现习题 \ref{ex:I.5.6.6} 有用.)
\begin{enumerate}
\item[(a)] 设 $M$ 是所有 $\mathrm{UCB}(G)$ 上的平均的集合. 证明这是 $\mathrm{UCB}(G)$ 的连续对偶的非空弱*紧凸子集, 且 $G$ 在 $M$ 上由 $x \cdot M(f) = M(L_{x^{-1}}f)$ ($x \in G$, $M \in M$, $f \in \mathrm{UCB}(G)$) 给出的作用是连续且仿射的. (对 $M$ 上的弱*拓扑.)
\item[(b)] 证明 $\mathrm{UCB}(G)$ 上存在左不变平均 $m$.\footnote{若 $G$ 是一般拓扑群, 若这样的平均存在则称 $G$ 可解. 我们在本题中证明的其中一件事是: 对局部紧 $G$, 这与另一个定义等价.}\footnote{参考文献}
\item[(c)] 证明若 $f \in L^1(G)^{1,+}$, $\phi \in \mathrm{UCB}(G)$, 则 $m(f * \phi) = m(\phi)$.
\item[(d)] 证明若 $f, f' \in L^1(G)^{1,+}$, $\phi \in L^{\infty}(G)$, 则 $m(f * \phi) = m(f' * \phi)$.
\item[(e)] 设 $f_0 \in L^1(G)^{1,+}$. 证明公式 $\phi \longmapsto m(f_0 * \phi)$ 定义 $L^{\infty}(G)$ 上的平均 $\tilde{m}$, 且对每个 $f \in L^1(G)^{1,+}$ 和每个 $\phi \in L^{\infty}(G)$ 有 $\tilde{m}(f * \phi) = \tilde{m}(\phi)$.
\end{enumerate}
设 $E = \prod_{f \in L^1(G)^{1,+}}L^1(G)$. 我们在 $E$ 上考虑两个拓扑:
\begin{itemize}
\item $L^1(G)$ 上的弱*拓扑 (通过把 $L^1(G)$ 看作 $L^{\infty}(G)$ 的连续对偶的子空间得到) 的乘积. 我们称之为 $E$ 上的弱拓扑.
\item $L^1(G)$ 上由范数 $\|\cdot\|_1$ 定义的拓扑的乘积. 我们称之为 $E$ 上的强拓扑.
\end{itemize}
\begin{enumerate}
\item[(a)] 设
\[
\Sigma = \{(f * g - g)_{f \in L^1(G)^{1,+}, g \in L^1(G)^{1,+}}\} \subset E.
\]
证明 $\Sigma$ 在弱拓扑中的闭包包含 $0$.
\item[(b)] 证明 $\Sigma$ 在强拓扑中的闭包包含 $0$. (提示: $E$ 上的任何强连续线性泛函 $\Lambda$ 都可以写成 $\Lambda((g_f)_{f \in L^1(G)^{1,+}}) = \sum_{f \in L^1(G)^{1,+}}\int_G g_f\phi_f\,\dd\mu$, 其中 $\phi_f \in L^{\infty}(G)$, 且除有限多个 $f$ 外 $\phi_f = 0$.)
\item[(c)] 设 $Q \ni 1$ 是 $G$ 的紧子集, $\varepsilon > 0$, $f \in L^1(G)^{1,+}$. 证明存在 $g \in L^1(G)^{1,+}$ 使得
\[
\sup_{x \in Q}\|(L_xf) * g - g\|_1 \leq \varepsilon.
\]
\item[(d)] 找到函数 $h \in L^1(G)^{1,+}$ 使得
\[
\sup_{x \in Q}\|L_xh - h\|_1 \leq 2\varepsilon.
\]
\item[(e)] 证明 $L^{\infty}(G)$ 上存在左不变平均. (若你不习惯用网, 可以假设 $G$ $\sigma$-紧, 即紧子集的可数并.)
\end{enumerate}
\end{exercise}

\begin{solution}
(a). 我们已经在习题 \ref{ex:III.6.3.1}(d) 的解答中看到 $M$ 是 $\Hom(\mathrm{UCB}(G), \CC)$ 的弱*紧凸子集, 且 $M$ 非空, 因为它包含所有线性泛函 $\delta_x$ ($x \in G$).
若 $x \in G$, 同态 $\Lambda \longmapsto \Lambda \circ L_{x^{-1}}$ 从 $\Hom(\mathrm{UCB}(G), \CC)$ 到自身线性, 且显然保持 $M$, 所以它在 $M$ 上的限制仿射.
还需证明映射 $G \times M \to M$, $(x, M) \longmapsto M \circ L_{x^{-1}}$ 连续. 由于我们在 $M$ 上用弱*拓扑, 这意味着对每个 $\phi \in \mathrm{UCB}(G)$, 映射 $G \times M \to \CC$, $(x, M) \longmapsto M(L_{x^{-1}}\phi)$ 连续. 固定 $\phi \in \mathrm{UCB}(G)$, 设 $\varepsilon > 0$. 由于 $\phi$ 左一致连续, 存在 $1$ 在 $G$ 中的开邻域 $U$ 使得对每个 $y \in U$ 有 $\|L_{y^{-1}}\phi - \phi\|_{\infty} < \varepsilon$. 注意对每个 $y \in U$ 和每个 $x \in G$,
\[
\|L_{x^{-1}y^{-1}}\phi - L_{x^{-1}}\phi\|_{\infty} = \|L_{y^{-1}}\phi - \phi\|_{\infty} < \varepsilon.
\]
设 $(x, M) \in G \times M$. 设 $V = \{M' \in M \mid |M(L_{x^{-1}}\phi) - M'(L_{x^{-1}}\phi)| < \varepsilon\}$. 这是 $M$ 的弱*邻域, 所以 $Ux \times V$ 是 $(x, M)$ 在 $G \times M$ 中的邻域. 若 $y \in U$, $M' \in V$, 有
\begin{align*}
|M(L_{x^{-1}}\phi) - M'(L_{(yx)^{-1}}\phi)| &\leq |M(L_{x^{-1}}\phi) - M'(L_{x^{-1}}\phi)| + |M'(L_{x^{-1}}\phi) - M'(L_{x^{-1}y^{-1}}\phi)| \\
&< \varepsilon + \|M'\|_{\mathrm{op}}\|L_{x^{-1}}\phi - L_{x^{-1}y^{-1}}\phi\|_{\infty} \\
&< 2\varepsilon
\end{align*}
(用习题 \ref{ex:III.6.3.1}(a) 看到 $\|M'\|_{\mathrm{op}} = 1$). 这证明所要的结果.
(b). $\mathrm{UCB}(G)$ 上的左不变平均正是 $G$ 在 $M$ 上由 $x \cdot M = M \circ L_{x^{-1}}$ 定义的作用的不动点. 所以这种平均的存在性由 (a) 和 $G$ 具有不动点性质这一事实得出.
(c). 由习题 \ref{ex:I.5.6.6}, 有 $f * \phi = \int_G f(y)L_y\phi\,\dd y$. 由习题 \ref{ex:I.5.6.6}, $\mathrm{UCB}(G)$ 上的线性泛函 $m$ 连续. 应用积分的定义和 $m$ 的左不变性, 得到
\[
m(f * \phi) = \int_G f(y)m(L_y\phi)\,\dd y = \int_G f(y)m(\phi)\,\dd y = m(\phi)\int_G f\,\dd\mu = m(\phi).
\]
(d). 设 $(\psi_U)_{U \in \mathcal{U}}$ 是 $G$ 上的近似单位元. 注意对每个 $U \in \mathcal{U}$, $\psi_U \in L^1(G)^{1,+}$. 设 $\phi \in L^{\infty}(G)$, $f, f' \in L^1(G)^{1,+}$. 由习题 \ref{ex:I.5.6.6}(a), 对每个 $U \in \mathcal{U}$ 有 $\psi_U * \phi \in \mathrm{UCB}(G)$, 所以由 (c),
\[
m(f * \psi_U * \phi) = m(\psi_U * \phi) = m(f' * \psi_U * \phi).
\]
又, 由命题 \ref{prop:I.4.1.9}, $\lim_{U \to \{1\}}f * \psi_U = f$, $\lim_{U \to \{1\}}f' * \psi_U = f'$. 在上面的等式中取 $U \to \{1\}$ 的极限 (忽略中间项), 并利用卷积乘积 $L^1(G) \times L^{\infty}(G) \to \mathrm{UCB}(G)$ 在两个变量中都连续 (由习题 \ref{ex:I.5.6.6}(a) 的解答) 且 $m$ 连续, 得到 $m(f * \phi) = m(f' * \phi)$.
(e). 映射 $\tilde{m}$ 由习题 \ref{ex:I.5.6.6}(a) 良定义, 且显然 $\CC$-线性. 若 $\phi = \mathbb{1}_G$, 则对每个 $x \in G$,
\[
f_0 * \phi(x) = \int_G f_0(y)\,\dd\mu(y) = 1,
\]
所以 $\tilde{m}(\phi) = m(\mathbb{1}_G) = 1$. 若 $\phi \geq 0$ 局部几乎处处, 则 $f_0 * \phi \geq 0$ 几乎处处, 所以 $\tilde{m}(\phi) \geq 0$. 这证明 $\tilde{m}$ 是 $L^{\infty}(G)$ 上的平均.
设 $f \in L^1(G)^{1,+}$, $\phi \in L^{\infty}(G)$. 则
\[
\tilde{m}(f * \phi) = m(f * f_0 * \phi) \quad \text{且} \quad \tilde{m}(\phi) = m(f_0 * \phi).
\]
由 (d), 为证明它们相等, 只需证明 $f * f_0 \in L^1(G)^{1,+}$. 我们已经知道由命题 \ref{prop:I.4.1.2}, $f * f_0 \in L^1(G)$, 且 $f * f_0 \geq 0$ 几乎处处由定义 $f * f_0$ 的公式显然. 最后,
\begin{align*}
\int_G f * f_0(x)\,\dd x &= \int_{G \times G}f(y)f_0(y^{-1}x)\,\dd x\,\dd y = \int_G f(y)\Bigl(\int_G f_0(y^{-1}x)\,\dd x\Bigr)\,\dd y \\
&= \int_G f(y)\,\dd y = 1.
\end{align*}

(f). 一个有用的记号: 对每个 $f \in L^1(G)$, 记 $M_f$ 为 $L^{\infty}(G)$ 上的线性泛函 $\phi \longmapsto \int_G f\phi\,\dd\mu$.
我们要证明如下结论: 对每个 $n \geq 1$, 对所有 $f_1, \ldots, f_n \in L^1(G)^{1,+}$, 若 $U_1, \ldots, U_n$ 是 $\Hom(L^{\infty}(G), \CC)$ 中 $0$ 的弱*邻域, 则存在 $g \in L^1(G)^{1,+}$ 使得对 $i \in \{1, \ldots, n\}$ 有 $M_{f_i * g} - M_g \in U_i$.
若 $f \in L^1(G)^{1,+}$, 映射 $c_f \colon \Lambda \longmapsto \Lambda(f * (\cdot))$ 从 $\Hom(L^{\infty}(G), \CC)$ 到自身弱*连续 (因为 $\phi \longmapsto f * \phi$ 在 $L^{\infty}(G)$ 上连续, 由习题 \ref{ex:I.5.6.6}(a)). 而且, 若 $\Lambda = M_g$ (其中 $g \in L^1(G)$), 则对每个 $\phi \in L^{\infty}(G)$,
\begin{align*}
(c_{f'}\Lambda)(\phi) &= \int_G g(y)(f' * \phi)(y)\,\dd y \\
&= \int_{G \times G}g(y)\Delta(x)^{-1}\overline{f(x^{-1})}\phi(x^{-1}y)\,\dd x\,\dd y \\
&= \int_{G \times G}\Delta(x)^{-1}\overline{f(x^{-1})}g(xz)\phi(z)\,\dd x\,\dd z \\
&= \int_G (f * g)\phi\,\dd\mu \quad \text{(见命题 \ref{prop:I.4.1.3})},
\end{align*}
其中 $f' \in L^1(G)$ 由 $f'(x) = \Delta(x)^{-1}\overline{f(x^{-1})}$ 定义. 换句话说, $c_{f'}M_g = M_{f*g}$.
固定 $n$, $f_1, \ldots, f_n$, $U_1, \ldots, U_n$ 如上. 则 $U := U_1 \cap \cdots \cap U_n$ 是 $\Hom(L^{\infty}(G), \CC)$ 中 $0$ 的弱*邻域. 选取另一个弱*邻域 $V$ 使得 $V = -V$ 且 $V + V \subset U$.
则 $\tilde{m} + V$ 是 $\tilde{m}$ 的弱*邻域, 所以由习题 \ref{ex:III.6.3.1}(c) 和上一段, 存在 $g \in L^1(G)^{1,+}$ 使得 $M_g - \tilde{m}$ 和所有 $M_{f_i*g} - c_{f'_i}\tilde{m}$ ($1 \leq i \leq n$) 都在 $V$ 中. 注意 $c_{f'_i}\tilde{m} = \tilde{m}$ (由 (e)), 所以对 $1 \leq i \leq n$,
\[
M_{f_i*g} - M_g = (M_{f_i*g} - \tilde{m}) + (\tilde{m} - M_g) \in V - V \subset U_i.
\]
(g). 注意 $\Sigma$ 是 $E$ 的凸子集. 设 $\bar{\Sigma}$ 是 $\Sigma$ 对强拓扑的闭包. 若 $0 \notin \bar{\Sigma}$, 则由 Hahn-Banach 定理 (第二几何形式), 存在强连续 $\R$-线性泛函 $\Lambda' \colon E \to \R$ 使得 $0 = \Lambda'(0) > \sup_{x \in \Sigma}\Lambda'(x)$. 像习题 \ref{ex:III.6.3.1}(c) 和 (d) 的解答那样, 可以把 $\Lambda'$ 写成 $\Lambda' = \Re\Lambda$, 其中 $\Lambda \colon E \to \CC$ 是强连续 $\CC$-线性泛函 (由 $\Lambda(x) = \Lambda'(x) + \frac{1}{i}\Lambda'(ix)$ 定义).
现在一个重要的注记是: 由于我们在 $E$ 上用乘积拓扑, 直和 $\bigoplus_{f \in L^1(G)^{1,+}}L^1(G)$ 在 $E$ 中稠密.
对每个 $f_0 \in L^1(G)^{1,+}$, 考虑线性泛函 $\Lambda_{f_0} \colon L^1(G) \to \CC$, 它是 $\Lambda$ 与包含因子 $f_0$ 进 $\prod_{f \in L^1(G)^{1,+}}L^1(G) = E$ 的复合. 这是 $L^1(G)$ 上的连续线性泛函, 所以存在唯一的 $\phi_{f_0} \in L^{\infty}(G)$ 使得 $\Lambda_{f_0}$ 是对 $\phi_{f_0}$ 的积分.
现在考虑 $L^1(G)^{1,+}$ 的递增族 $(X_n)_{n \geq 0}$ 使得 $L^1(G)^{1,+} = \bigcup_{n \geq 0}X_n$. 对每个 $x = (g_f)_{f \in L^1(G)^{1,+}} \in E$, 序列 $((g_f)_{f \in X_n})_{n \geq 0}$ 在强拓扑中收敛到 $x$, 所以
\[
\Lambda(x) = \lim_{n \to +\infty}\Lambda((g_f)_{f \in X_n}) = \lim_{n \to +\infty}\sum_{f \in X_n}\int_G g_f\phi_f\,\dd\mu = \sum_{f \in L^1(G)^{1,+}}\int_G g_f\phi_f\,\dd\mu.
\]
由于和对任何 $(g_f) \in E$ 收敛, 除有限多个 $f \in L^1(G)^{1,+}$ 外必有 $\phi_f = 0$.
但此时, 若考虑任何满足 $0 > c > \sup_{x \in \Sigma}\Re(\Lambda(x))$ 的实数 $c$, 集合 $\{x \in E \mid \Re(\Lambda(x)) \leq c\}$ 在 $E$ 中弱闭, 所以包含 $\Sigma$ 的弱闭包, 从而由 (f) 包含 $0$, 矛盾.
(h). 对每个 $x \in Q$, 设 $U_x$ 是 $x$ 在 $G$ 中的邻域使得对 $y \in U_x$ 有 $\|L_yf - L_xf\|_1 \leq \varepsilon/2$. (见命题 \ref{prop:I.3.1.13}.) 由于 $Q$ 紧, 可以找到 $x_1, \ldots, x_n \in Q$ 使得 $Q \subset \bigcup_{i=1}^n U_{x_i}$. 由 (f), 存在 $g \in L^1(G)^{1,+}$ 使得对每个 $i \in \{1, \ldots, n\}$ 有 $\|(L_{x_i}f) * g - g\|_1 \leq \varepsilon/2$.
让我们证明这个 $g$ 可行. 设 $x \in Q$. 则存在 $i \in \{1, \ldots, n\}$ 使得 $x \in U_{x_i}$, 且有
\begin{align*}
\|(L_xf) * g - g\|_1 &\leq \|(L_xf) * g - (L_{x_i}f) * g\|_1 + \|(L_{x_i}f) * g - g\|_1 \\
&\leq \|L_xf - L_{x_i}f\|_1\|g\|_1 + \|(L_{x_i}f) * g - g\|_1 \\
&\leq \varepsilon.
\end{align*}
(第二个等式用命题 \ref{prop:I.4.1.2}).
(i). 设 $h = f * g$. 有 $h \in L^1(G)^{1,+}$ (见 (e) 的解答), 且由于 $1 \in Q$, 对每个 $x \in Q$,
\[
\|L_xh - h\|_1 \leq \|L_xh - g\|_1 + \|g - h\|_1 \leq \|(L_xf) * g - g\|_1 + \|(L_1f) * g - g\|_1 \leq 2\varepsilon.
\]
(j). 假设 $G$ $\sigma$-紧, 写 $G = \bigcup_{n \geq 0}Q_n$, 其中每个 $Q_n$ 是 $G$ 的包含 $1$ 的紧子集. 对每个 $n \in \Z_{\geq 0}$, 由 (h) 可以找到函数 $h_n \in L^1(G)^{1,+}$ 使得 $\sup_{x \in Q_n}\|L_xh_n - h_n\|_1 \leq 2^{-n}$. 序列 $(M_{h_n})_{n \geq 0}$ 是 $L^{\infty}(G)$ 上的平均的弱*紧子集 (我们在习题 \ref{ex:III.6.3.1}(c) 中看到这个集合弱*紧) 的元素, 所以它有收敛子序列, 可以假设它收敛. 设 $M = \lim_{n \geq 0}M_{h_n}$. 我们证明 $M$ 左不变. 设 $x \in G$. 则对 $n \gg 0$ 有 $x^{-1} \in Q_n$, 所以对每个 $\phi \in L^{\infty}(G)$,
\begin{align*}
M(L_x\phi) &= \lim_{n \to +\infty}M_{h_n}(L_x\phi) \\
&= \lim_{n \to +\infty}\int_G h_n(y)\phi(x^{-1}y)\,\dd y \\
&= \lim_{n \to +\infty}\int_G h_n(xy)\phi(y)\,\dd y \\
&= \lim_{n \to +\infty}\int_G L_{x^{-1}}h_n\phi\,\dd\mu,
\end{align*}
且
\begin{align*}
|M(L_x\phi) - M(\phi)| &= \Bigl|\lim_{n \to +\infty}\int_G (L_{x^{-1}}h_n - h_n)\phi\,\dd\mu\Bigr| \\
&\leq \lim_{n \to +\infty}\|L_{x^{-1}}h_n - h_n\|_1\|\phi\|_{\infty} \\
&= 0,
\end{align*}
即 $M(L_x\phi) = M(\phi)$.
假设 $G$ 不 $\sigma$-紧. 则写 $G = \bigcup_{Q \in \mathcal{Q}}Q$, 其中 $\mathcal{Q}$ 是 $G$ 的紧子集族, 使得若 $Q_1, Q_2 \in \mathcal{Q}$ 则 $Q_1 \cup Q_2 \in \mathcal{Q}$. 即 $\mathcal{Q}$ 对由包含给出的序关系是定向集. 对每个 $Q \in \mathcal{Q}$, 由 (i) 可以找到函数 $h_Q \in L^1(G)^{1,+}$ 使得 $\sup_{x \in Q}\|L_xh_Q - h_Q\|_1 \leq (1 + \mu(Q))^{-1}$. 若 $G$ 不紧, 则 $\mu(G) = +\infty$, 所以 $\lim_{Q \in \mathcal{Q}}(1 + \mu(Q))^{-1} = 0$. 设 $M$ 是 $(M_{h_Q})_{Q \in \mathcal{Q}}$ 的弱*极限点, 它存在, 因为 $L^{\infty}(G)$ 上的平均集合弱*紧. 则我们完全像上面那样看到 $M$ 左不变.
\end{solution}

\begin{exercise}[可解性与平凡表示的弱包含]
\label{ex:III.6.3.4}
设 $G$ 是局部紧群. 回忆习题 \ref{ex:III.6.2.4}.
\begin{enumerate}
\item[(a)] 若 $G$ 可解, 证明平凡表示弱包含于 $G$ 的左正则表示. (提示: 对所有 $a, b \in \R_{\geq 0}$, 有 $|a - b|^2 \leq |a^2 - b^2|$.)
\item[(b)] 若平凡表示弱包含于 $G$ 的左正则表示, 证明 $G$ 可解. (提示: 对所有 $a, b \in \CC$, 证明 $||a|^2 - |b|^2| \leq |a + b||a - b|$.)
\end{enumerate}
\end{exercise}

\begin{solution}
(a). 让我们先证明提示中的不等式. 设 $a, b \in \R_{\geq 0}$. 可以假设 $a \geq b$. 则
\[
|a - b|^2 = (a - b)^2 = a^2 + b^2 - 2ab \leq a^2 + b^2 - 2b^2 = a^2 - b^2 = |a^2 - b^2|.
\]
假设 $G$ 可解. 设 $K$ 是 $G$ 的紧子集, $c > 0$. 由习题 \ref{ex:III.6.3.3}(i), 存在 $h \in L^1(G)^{1,+}$ 使得 $\sup_{x \in K}\|L_xh - h\|_1 < c^2$. 设 $f = \sqrt{h}$. 则由上面的不等式, 对每个 $x \in K$,
\begin{align*}
\|L_xf - f\|_2^2 &= \int_G |f(x^{-1}y) - f(y)|^2\,\dd y \\
&\leq \int_G |h(x^{-1}y) - h(y)|\,\dd y \\
&= \|L_xh - h\|_1 \\
&< c^2,
\end{align*}
所以 $\|L_xf - f\|_2 < c$.
由习题 \ref{ex:III.6.2.4}(d), 这蕴含 $G$ 的平凡表示弱包含于正则表示.
(b). 我们验证习题 \ref{ex:III.6.3.3}(i) 的结果成立, 即对 $G$ 的每个紧子集 $Q$ 和每个 $\varepsilon > 0$, 存在 $h \in L^1(G)^{1,+}$ 使得 $\sup_{x \in Q}\|L_xh - h\|_1 \leq \varepsilon$. 事实上, 我们已在习题 \ref{ex:III.6.3.3}(j) 中看到这蕴含 $L^{\infty}(G)$ 上存在左不变平均.
设 $Q$ 是 $G$ 的紧子集, $\varepsilon > 0$. 由习题 \ref{ex:III.6.2.4}(c), 存在 $f \in L^2(G)$ 使得 $\|f\|_2 = 1$ 且 $\sup_{x \in Q}\|L_xf - f\|_2 \leq \varepsilon/2$. 设 $h = |f|^2$. 则 $\|h\|_1 = \|f\|_2^2 = 1$, 所以 $h \in L^1(G)^{1,+}$. 注意对所有 $a, b \in \CC$, 由三角不等式有 $|a^2 - b^2| \geq ||a|^2 - |b|^2|$, 所以
\[
|a + b|^2|a - b|^2 = (a^2 - b^2)(\overline{a^2 - b^2}) = |a^2 - b^2|^2 \geq ||a|^2 - |b|^2|^2.
\]
现在, 若 $x \in Q$, 得到
\begin{align*}
\|L_xh - h\|_1 &= \int_G ||L_xf(y)|^2 - |f(y)|^2|\,\dd y \\
&\leq \int_G (|L_xf(y) + f(y)|)|L_xf(y) - f(y)|\,\dd y \\
&\leq \|L_xf - f\|_2\|L_xf + f\|_2 \quad \text{(Cauchy-Schwarz)} \\
&\leq \varepsilon
\end{align*}
(因为 $\|L_xf + f\|_2 \leq \|L_xf\|_2 + \|f\|_2 = 2\|f\|_2 = 2$).
\end{solution}

\begin{exercise}[交换局部紧群的可解性]
\label{ex:III.6.3.5}
设 $G$ 是交换局部紧群. 本题的目标是证明 $G$ 具有不动点性质 (从而是可解的).
设 $V$ 是局部凸拓扑向量空间, $K$ 是 $V$ 的非空紧凸子集, $G \times K \to K$, $(x, v) \longmapsto x \cdot v$ 是 $G$ 在 $K$ 上的仿射作用.
对每个 $n \in \Z_{\geq 0}$ 和每个 $x \in G$, 定义连续仿射映射 $A_n(x) \colon K \to K$:
\[
A_n(x)(v) = \frac{1}{n+1}\sum_{i=0}^n x^i \cdot v.
\]
设 $\mathcal{G}$ 是所有 $A_n(x)$ ($n \geq 0$, $x \in G$) 生成的 $K \to K$ 的连续仿射映射的半群. (即其元素是态射 $A_n(x)$ 的有限复合的半群, 其中半群运算就是映射 $K \to K$ 的复合.)
\begin{enumerate}
\item[(a)] 设 $v \in \bigcap_{\gamma \in \mathcal{G}}\gamma(K)$. 证明 $v$ 是 $G$ 作用的不动点. (提示: 对 $V$ 上的每个连续线性泛函 $\Lambda$ 和每个 $x \in G$, 证明 $\Lambda(v) = \Lambda(x \cdot v)$.)
\item[(b)] 对所有 $\gamma_1, \ldots, \gamma_n \in \mathcal{G}$, 证明 $\bigcap_{i=1}^n\gamma_i(K) \neq \emptyset$.
\item[(c)] 证明 $G$ 在 $K$ 中有不动点.
\end{enumerate}
\end{exercise}

\begin{solution}
(a). 设 $x \in G$. 设 $\Lambda$ 是 $V$ 上的连续线性泛函. 由于 $K$ 紧, $C := \sup_{w \in K}|\Lambda(w)| < +\infty$. 若 $n \geq 0$, 有 $v \in A_n(x)(K)$, 所以存在 $w \in K$ 使得 $v = A_n(x)(w)$. 由于 $G$ 的作用仿射, 这蕴含
\[
x \cdot v = \frac{1}{n+1}\sum_{i=0}^n x^{i+1} \cdot w,
\]
所以 $v - x \cdot v = \frac{1}{n+1}(w - x^{n+1} \cdot w)$, 从而 $|\Lambda(v - x \cdot v)| \leq \frac{2C}{n+1}$. 由于这对每个 $n \geq 0$ 成立, 有 $|\Lambda(v) - \Lambda(x \cdot v)| = 0$, 即 $\Lambda(v) = \Lambda(x \cdot v)$. 由于 $V$ 上的连续线性泛函分离点, 最终得到 $x \cdot v = v$.
(b). 注意若 $x, y \in G$, $n, m \in \Z_{\geq 0}$, 则对每个 $v \in K$,
\begin{align*}
A_n(x) \circ A_m(y)(v) &= \frac{1}{n+1}\sum_{i=1}^n x^i \cdot A_m(y)(v) \\
&= \frac{1}{(n+1)(m+1)}\sum_{i=0}^n\sum_{j=0}^m x^i \cdot (y^j \cdot v) \\
&= \frac{1}{(n+1)(m+1)}\sum_{i=0}^n\sum_{j=0}^m y^j \cdot (x^i \cdot v) \quad \text{(因为 } G \text{ 交换)} \\
&= A_m(y) \circ A_n(x)(v).
\end{align*}
这蕴含半群 $\mathcal{G}$ 交换.
现在设 $\gamma_1, \ldots, \gamma_n \in \mathcal{G}$. 则对每个 $i \in \{1, \ldots, n\}$,
\[
\gamma_i(K) \supset \gamma_i(\gamma_1 \circ \cdots \circ \gamma_{i-1} \circ \cdots \circ \gamma_n(K)) = \gamma_1 \circ \cdots \circ \gamma_n(K).
\]
所以
\[
\bigcap_{i=1}^n\gamma_i(K) \supset \gamma_1 \circ \cdots \circ \gamma_n(K) \neq \emptyset.
\]
(c). 由于 $K$ 紧且每个 $\gamma \in \mathcal{G}$ 连续, 对每个 $\gamma \in \mathcal{G}$, $K$ 的子集 $\gamma(K)$ 紧, 从而在 $K$ 中闭. 由 (b), 族 $(\gamma(K))_{\gamma \in \mathcal{G}}$ 有有限交性质. 由 $K$ 的紧性, $\bigcap_{\gamma \in \mathcal{G}}\gamma(K) \neq \emptyset$. 由 (a), 此交的每个点都是 $G$ 在 $K$ 上的不动点.
\end{solution}

\begin{exercise}[$\Z$ 上的不变概率测度]
\label{ex:III.6.3.6}
\begin{enumerate}
\item[(a)] 设 $P(\Z)$ 是 $\Z$ 的子集的集合. 证明 $\Z$ 上存在有限可加左不变概率测度, 即满足下列条件的函数 $\mu \colon P(\Z) \to \R_{\geq 0}$:
\begin{enumerate}
\item[(i)] 若 $A_1, \ldots, A_n \in P(\Z)$ 满足对 $i \neq j$ 有 $A_i \cap A_j = \emptyset$, 则 $\mu(A_1 \cup \cdots \cup A_n) = \mu(A_1) + \cdots + \mu(A_n)$.
\item[(ii)] $\mu(\Z) = 1$.
\item[(iii)] 对每个 $A \in P(\Z)$ 和 $n \in \Z$, 有 $\mu(n + A) = \mu(A)$.
\end{enumerate}
\item[(b)] (a) 的测度唯一吗? (提示: 你需要某种相当显式的方法来构造 $\Z$ 上的不变平均. 例如你可以试着利用 $L^{\infty}(\Z) \to \CC$, $(x_k)_{k \in \Z} \longmapsto \frac{1}{2k+1}\sum_{k=-n}^n x_k$ 的 (非不变) 平均序列 $M_n$.)
\end{enumerate}
\end{exercise}

\begin{solution}
(a). 由于 $\Z$ 是交换局部紧群, 由习题 \ref{ex:III.6.3.3} 和 \ref{ex:III.6.3.5} 它可解. 这意味着 $L^{\infty}(\Z)$ 上存在左不变平均 $M$. 定义 $\mu$ 为 $\mu(A) = M(\mathbb{1}_A)$; 由平均的定义, 此函数取 $\R_{\geq 0}$ 中的值. 则 $\mu$ 满足 (i), 因为 $M$ 线性; 满足 (ii), 因为 $M(1) = 1$; 满足 (iii), 因为 $M$ 左不变.
反之, 注意题设中 $\mu$ 的存在蕴含不变平均的存在.
(b). 不唯一.
设 $V = L^{\infty}(\Z)$, 考虑线性泛函族 $M_n \colon L^{\infty}(\Z) \to \CC$:
\[
M_n((x_k)_{k \in \Z}) = \frac{1}{2n+1}\sum_{k=-n}^n x_k,
\]
其中 $n \in \N$. 对 $V$ 的每个元素 $x$ 有 $|M_n(x)| \leq \|x\|_{\infty}$, 所以 $M_n$ 连续. 又, 由定义显然 $M_n$ 是平均. 若 $a \in \Z$, 则对 $V$ 的每个元素 $x$ 和每个 $n \in \N$,
\[
|M_n(L_ax) - M_n(x)| \leq \frac{2|a|}{2n+1}\|x\|_{\infty}.
\]
所以, 若能使序列 $(M_n)_{n \geq 0}$ 在 $\Hom(V, \CC)$ 的弱*拓扑中收敛, 则其极限将是不变平均, 且它将定义如 (a) 的不变有限可加概率测度. 我们总可以用 Banach-Alaoglu 定理找到 $(M_n)_{n \geq 0}$ 的收敛子序列, 但我们还想证明可以得到两个不同的极限.
考虑 $V$ 的元素 $x = (x_n)_{n \in \Z}$: 对 $n \leq 0$ 取 $x_n = 0$; 若 $2k \leq n \leq 2^{k+1} - 1$ ($k \in \Z_{\geq 0}$), 取 $x_n = (-1)^k$. 则若 $n = 2^k - 1$ ($k \geq 0$),
\[
\sum_{r=-n}^n x_r = \sum_{s=0}^{k-1}(-1)^s2^s = \frac{1 - (-2)^k}{3},
\]
所以
\[
M_n(x) = \frac{1 - (-2)^k}{3(2^{k+1} - 1)}.
\]
特别地, 序列 $(M_{2^{2l}-1}(x))_{l \geq 0}$ 收敛到 $-\frac{1}{6}$, 序列 $(M_{2^{2l+1}-1}(x))_{l \geq 0}$ 收敛到 $\frac{1}{6}$.
由 Banach-Alaoglu 定理,\footnote{参考文献 ?} 序列 $(M_{2^{2l}-1})_{l \geq 0}$ 和 $(M_{2^{2l+1}-1})_{l \geq 0}$ 都有弱*极限点, 记为 $M$ 和 $M'$. $M$ 和 $M'$ 都是 $V$ 上的左不变平均, 但由上面的计算 $M(x) = -\frac{1}{6}$, $M'(x) = \frac{1}{6}$, 所以 $M \neq M'$.
\end{solution}

\begin{exercise}[$\SL_2(\R)$ 不可解]
\label{ex:III.6.3.7}
\begin{enumerate}
\item[(a)] 设 $G$ 是作用于集合 $X$ 的群. 假设有 $G$ 的子群 $G_1$, $G_2$ 和 $X$ 的子集 $X_1$, $X_2$ 满足:
\begin{itemize}
\item 集合 $X_1$ 和 $X_2$ 非空, 且 $X_1 \neq X_2$;
\item 对每个 $x \in G_1 - \{1\}$, 有 $x \cdot X_1 \subset X_2$;
\item 对每个 $x \in G_2 - \{1\}$, 有 $x \cdot X_2 \subset X_1$;
\item $G_2$ 的基数至少为 $3$.
\end{itemize}
证明不可能有等式 $1 = h_1\cdots h_n$, 其中对奇数 $i$, $h_i \in G_1 - \{1\}$; 对偶数 $i$, $h_i \in G_2 - \{1\}$; $n \geq 1$.
\item[(b)] 设 $a_1, a_2 \in \CC$ 满足 $|a_1| \geq 2$, $|a_2| \geq 2$. 定义 $x, y \in \SL_2(\CC)$ 为
\[
x = \begin{pmatrix} 1 & a_1 \\ 0 & 1 \end{pmatrix}, \qquad y = \begin{pmatrix} 1 & 0 \\ a_2 & 1 \end{pmatrix}.
\]
证明由 $x$ 和 $y$ 生成的 $\SL_2(\CC)$ 的子群同构于两个生成元上的自由群. (提示: 令 $\SL_2(\CC)$ 按通常方式作用于 $\CC^2$. 看子集 $\{(z_1, z_2) \in \CC^2 \mid |z_1| > |z_2|\}$ 和 $\{(z_1, z_2) \in \CC^2 \mid |z_1| < |z_2|\}$.)
\item[(c)] 设 $G = \SL_2(\R)$, 赋予离散拓扑. 证明 $G$ 不可解.
\end{enumerate}
\end{exercise}

\begin{solution}
(a). 若 $G_1 = \{1\}$, 结果显然. 所以可以假设 $G_1 \neq \{1\}$.
假设有 $1 = h_1\cdots h_n$, 其中对奇数 $i$, $h_i \in G_1 - \{1\}$; 对偶数 $i$, $h_i \in G_2 - \{1\}$; $n \geq 1$.
我们首先假设 $n$ 偶. 由于 $|G_2| \geq 3$, 可以找到 $h \in G_2 - \{1\}$ 使得 $h \neq h_n$. 注意 $1 = hh^{-1} = hh_1\cdots(h_nh^{-1})$, 其中 $h_nh^{-1} \in G_2 - \{1\}$. 设 $g \in G_1 - \{1\}$. 也有 $1 = gg^{-1} = ghh_1\cdots(h_nh^{-1})g^{-1}$. 所以对每个 $x \in X_2$,
\[
x = hh_1\cdots h_{n-1}(h_nh^{-1})(x) \in X_1,
\]
所以 $X_2 \subset X_1$. 另一方面, 对每个 $y \in X_1$,
\[
y = ghh_1\cdots(h_nh^{-1})g^{-1}(y) \in X_2,
\]
所以 $X_1 \subset X_2$. 这与 $X_1 \neq X_2$ 矛盾.
现在假设 $n$ 奇. 设 $h \in G_2 - \{1\}$. 则 $1 = hh^{-1} = hh_1\cdots h_nh^{-1}$. 所以对每个 $x \in X_2$,
\[
x = hh^{-1} = hh_1\cdots h_nh^{-1}(x) \in X_1,
\]
所以 $X_2 \subset X_1$. 另一方面, 对每个 $y \in X_1$,
\[
y = h_1\cdots h_n(y) \in X_2,
\]
所以 $X_1 \subset X_2$. 同样, 这与 $X_1 \neq X_2$ 矛盾.
(b). 我们想把 (a) 应用于 $X = \CC^2$, $X_1 = \{(z_1, z_2) \in \CC^2 \mid |z_1| < |z_2|\}$, $X_2 = \{(z_1, z_2) \in \CC^2 \mid |z_1| > |z_2|\}$, $G_1 = \langle x \rangle$, $G_2 = \langle y \rangle$. 我们必须验证这些子集和子群满足 (a) 的条件.
设 $g \in G_1 - \{1\}$, $(z_1, z_2) \in X_1$. 有 $g = x^n$ (其中 $n \in \Z - \{0\}$), 所以 $g = \begin{pmatrix} 1 & na_1 \\ 0 & 1 \end{pmatrix}$, 且 $g \cdot (z_1, z_2) = (z_1 + na_1z_2, z_2)$. 所以
\[
|z_1 + na_1z_2| \geq |n||a_1||z_2| - |z_1| \geq 2|z_2| - |z_1| > |z_2|,
\]
即 $g \cdot (z_1, z_2) \in X_2$. (我们用 $|n| \geq 1$ 这一事实.) 对 $g \in G_2 - \{1\}$ 和 $(z_1, z_2) \in X_2$ 证明 $g \cdot (z_1, z_2) \in X_1$ 类似.
设 $G$ 是由 $x$ 和 $y$ 生成的 $\SL_2(\CC)$ 的子群, $F$ 是生成元 $a$ 和 $b$ 上的自由群. 有满射群同态 $\phi \colon F \to G$, 把 $F$ 的元素 $a^{n_1}b^{m_1}\cdots a^{n_r}b^{m_r}$ ($r \geq 0$, $n_1, m_1, \ldots, n_r, m_r \in \Z$) 送到 $x^{n_1}y^{m_1}\cdots x^{n_r}y^{m_r} \in G$. 我们要验证 $\phi$ 单射. 这意味着它的核平凡, 即它把 $F$ 中的既约字送到 $G$ 的非平凡元素. 但这个性质正是 (a) 的结论.
(c). 假设 $G$ 可解. 则由习题 \ref{ex:III.6.3.4}, $G$ 在 $\CC$ 上的平凡表示 $\mathbf{1}$ 包含于其正则表示 $\pi_L$. 设 $H$ 是 $G$ 的子群. 由弱包含的定义立即得到 $H$ 的表示 $\mathbf{1}|_H$ (它就是 $H$ 的平凡表示) 弱包含于 $\pi_L|_H$. 设 $\pi$ 是 $H$ 的正则表示, 让我们证明 $\pi_L|_H$ 弱包含于 $\pi$. 这将蕴含 $H$ 的平凡表示包含于其正则表示.
设 $(x_i)_{i \in I}$ 是商 $H\backslash G$ 的代表系; 有 $G = \bigsqcup_{i \in I}Hx_i$. 设 $\phi$ 是与 $\pi_L|_H$ 关联的正定型函数. 这意味着存在 $f \in L^2(G)$ 使得对每个 $x \in H$,
\[
\phi(x) = \langle L_xf, f \rangle_{L^2(G)}.
\]
对每个 $i \in I$, 设 $f_i = f|_{Hx_i} \in L^2(G)$. 则级数 $\sum_{i \in I}f_i$ 在 $L^2(G)$ 中收敛到 $f$, 且若 $i \neq j$, 对每个 $x \in H$ 有 $\langle L_xf_i, f_j \rangle_{L^2(G)} = 0$ (因为 $L_xf_i$ 和 $f_j$ 有不相交的支撑). 特别地, $\|f\|_2^2 = \sum_{i \in I}\|f_i\|_2^2$. 所以对每个 $x \in H$,
\[
\phi(x) = \sum_{i \in I}\langle L_xf_i, f_i \rangle_{L^2(G)},
\]
且这个和在 $x \in H$ 上一致收敛 (因为 $|\langle L_xf_i, f_i \rangle_{L^2(G)}| \leq \|f_i\|_2^2$).
对每个 $i \in I$, 定义 $g_i \in L^2(H)$ 为 $g_i(y) = f_i(yx_i)$. 则对每个 $x \in H$ 有 $\langle L_xg_i, g_i \rangle_{L^2(H)} = \langle L_xf_i, f_i \rangle_{L^2(G)}$. 所以我们已经把 $\phi$ 写成 $H$ 的正则表示关联的正定型函数的有限和的极限, 这正是我们想要的.
总之, 我们证明了: 若 $G$ 可解, 则对 $G$ 的每个子群 $H$, $H$ 的平凡表示包含于其正则表示 (即 $H$ 也可解). 注意到目前为止我们只用 $G$ 离散这一事实.
现在若 $G = \SL_2(\R)$, (b) 说 $G$ 有同构于两个生成元上的自由群的子群 $H$ (只需在 (b) 中取 $a_1, a_2 \in \R$). 则上面的结果与习题 \ref{ex:III.6.2.8} 矛盾.
\end{solution}


\chapter[Peter-Weyl 定理]{Peter-Weyl 定理}
% ============ 第 IV 部分: \index{Peter-Weyl 定理}Peter-Weyl 定理 (The Peter-Weyl Theorem) ============
% 原书 IV.1 - IV.9 (原书 pp.195-229)
% 编号: IV.1.1 (节级) / 习题 IV.9.1 (节级)

\section{\index{紧算子}紧算子}

\begin{definition}[紧算子 (compact operator)]
\label{def:IV.1.1}
设 $V$ 和 $W$ 是巴拿赫空间, $B$ 是 $V$ 中的闭单位球. 若 $T(B)$ 紧, 则称连续线性算子 $T \colon V \to W$ 紧 (compact).
\end{definition}

\begin{example}[紧算子的例子]
\label{exa:IV.1.2}
\begin{enumerate}
\item[(1)] 若 $\Im(T)$ 有限维 (即 $T$ 有有限秩 (finite rank)), 则 $T$ 紧.
\item[(2)] 若 $T$ 是有限秩算子的极限, 则 $T$ 紧; 更一般地, 紧算子的任何极限都紧 (见习题 \ref{ex:I.5.5.9}).
  反之, 若 $W$ 是希尔伯特空间, 则每个紧算子 $T \colon V \to W$ 都是有限秩算子的极限.\footnote{一般情形这不成立, 见 Enflo 的文章 {[}10{]} 中的反例.}
\item[(3)] $V$ 的恒等算子紧当且仅当 $V$ 有限维. (这是 Riesz 引理的推论, 见定理 \ref{thm:B.4.2}.)
\end{enumerate}
\end{example}

在本课程中, 我们只需使用希尔伯特空间的自伴紧自同态. 对它们有谱定理的一个更简单的版本.

\begin{theorem}[\index{自伴紧算子的谱定理}自伴紧算子的谱定理]
\label{thm:IV.1.3}
设 $V$ 是 $\CC$ 上的希尔伯特空间, $T \colon V \to V$ 是 $V$ 的连续自同态. 假设 $T$ 紧且自伴, 对每个 $\lambda \in \CC$ 记 $V_{\lambda} = \Ker(T - \lambda\mathrm{id}_V)$.
则:
\begin{enumerate}
\item[(i)] 若 $V_{\lambda} \neq 0$, 则 $\lambda \in \R$.
\item[(ii)] 若 $\lambda, \mu \in \CC$ 且 $\lambda \neq \mu$, 则 $V_{\mu} \subset V_{\lambda}^{\perp}$.
\item[(iii)] 若 $\lambda \in \CC - \{0\}$, 则 $\dim_{\CC}V_{\lambda} < +\infty$.
\item[(iv)] $\{\lambda \in \CC \mid V_{\lambda} \neq 0\}$ 有限或可数, 且其唯一可能的极限点是 $0$.
\item[(v)] $\bigoplus_{\lambda \in \CC}V_{\lambda}$ 在 $V$ 中稠密.
\end{enumerate}
\end{theorem}

\begin{proof}
我们证明 (i). 设 $\lambda \in \CC$ 使得 $V_{\lambda} \neq 0$, 选取非零的 $v \in V_{\lambda}$. 则
\[
\lambda\|v\|^2 = \langle \lambda v, v \rangle = \langle T(v), v \rangle = \langle v, T^*(v) \rangle = \langle v, T(v) \rangle = \bar{\lambda}\|v\|^2.
\]
由于 $\|v\| \neq 0$, 这蕴含 $\lambda \in \R$.
我们证明 (ii). 设 $\lambda, \mu \in \CC$ 满足 $\lambda \neq \mu$, $v \in V_{\lambda}$, $w \in V_{\mu}$. 要证明 $\langle v, w \rangle = 0$. 由 (i), 只需处理 $\lambda, \mu \in \R$ 的情形 (否则 $v = w = 0$). 此时
\[
\lambda\langle v, w \rangle = \langle T(v), w \rangle = \langle v, T(w) \rangle = \mu\langle v, w \rangle,
\]
所以 $\langle v, w \rangle = 0$.
设 $r > 0$. 设 $W = \bigoplus_{|\lambda| \geq r}V_{\lambda}$. 要证明 $\dim W < +\infty$, 这将蕴含 (iii) 和 (iv). 选取 $W$ 的由 $T$ 的特征向量组成的 Hilbert 基 $(e_i)_{i \in I}$, 即对每个 $i \in I$, $T(e_i) = \lambda_ie_i$, 其中 $|\lambda_i| \geq r$. 若 $I$ 无穷, 则族 $(T(e_i))_{i \in I}$ 不可能有收敛的 (非稳定的) 子序列. 事实上, 若存在单射映射 $\N \to I$, $n \longmapsto i_n$, 使得 $(T(e_{i_n}))_{n \geq 0}$ 收敛到 $V$ 的某个向量 $v$, 则 $\lambda_{i_n}e_{i_n} \to v$, 所以 $v$ 在 $\operatorname{Span}(e_{i_n} \mid n \geq 0)$ 的闭包中. 但另一方面, 对每个 $n \geq 0$,
\[
\langle v, e_{i_n} \rangle = \lim_{m \to +\infty}\langle \lambda_{i_m}e_{i_m}, e_{i_n} \rangle = 0,
\]
所以 $v \in \operatorname{Span}(e_{i_n} \mid n \geq 0)^{\perp}$. 这迫使 $v = 0$. 但 $\|v\| = \lim_{n \to +\infty}\|\lambda_{i_n}e_{i_n}\| \geq r > 0$, 矛盾. 由于 $T$ 紧, 这证明 $I$ 不能无穷, 即 $\dim(W) < +\infty$.
让我们证明 (v). 设 $W' = \bigoplus_{\lambda \in \CC}V_{\lambda}$, $W = W'^{\perp}$. 要证明 $W = 0$. 所以假设 $W \neq 0$. 由于 $T$ 自伴且 $W'$ 显然被 $T$ 稳定, 有 $T(W) \subset W$. (若 $v \in W$, 则对每个 $w \in W'$, $\langle T(v), w \rangle = \langle v, T(w) \rangle = 0$.) 由 $W$ 的定义, $\Ker(T|_W) = \{0\}$, 所以 $\|T|_W\|_{\mathrm{op}} > 0$. 设 $B = \{x \in W \mid \|x\| = 1\}$. 由于由下面的引理, $\|T|_W\|_{\mathrm{op}} = \sup_{x \in B}|\langle T(x), x \rangle|$, 存在 $B$ 的元素序列 $(x_n)_{n \geq 0}$ 使得当 $n \to +\infty$ 时 $\langle T(x_n), x_n \rangle \to \lambda$, 其中 $\lambda = \pm\|T|_W\|_{\mathrm{op}}$. 则
\[
0 \leq \|T(x_n) - \lambda x_n\|^2 = \|T(x_n)\|^2 + \lambda^2\|x_n\|^2 - 2\lambda\langle T(x_n), x_n \rangle \leq 2\lambda^2 - 2\lambda\langle T(x_n), x_n \rangle
\]
在 $n \to +\infty$ 时收敛到 $0$, 所以 $T(x_n) - \lambda x_n$ 本身收敛到 $0$. 由于 $T$ 紧, 可以假设序列 $(T(x_n))_{n \geq 0}$ 在 $W$ 中有极限, 设为 $w$. 则 $T(w) - \lambda w = 0$. 由 $W$ 的定义, 必有 $w = 0$. 但此时 $T(x_n) \to 0$, 所以 $\langle T(x_n), x_n \rangle \to 0$, 从而 $\lambda = 0 = \|T|_W\|_{\mathrm{op}}$, 矛盾.
\end{proof}

\begin{lemma}[自伴算子的范数]
\label{lem:IV.1.4}
设 $V$ 是希尔伯特空间, $T \in \End(V)$ 自伴. 则
\[
\|T\|_{\mathrm{op}} = \sup_{x \in V, \|x\| = 1}|\langle T(x), x \rangle|.
\]
\end{lemma}

\begin{proof}
设 $c = \sup_{x \in V, \|x\| = 1}|\langle T(x), x \rangle|$. 由 $\|T\|_{\mathrm{op}}$ 的定义有 $c \leq \|T\|_{\mathrm{op}}$. 由于 $\|T\|_{\mathrm{op}}$ 等于在范数 $1$ 向量上的 $\sup_{x,y \in V}|\langle T(x), y \rangle|$, 为证另一个不等式只需证明对所有满足 $\|x\| = \|y\| = 1$ 的 $x, y \in V$ 有 $|\langle T(x), y \rangle| \leq c$. 设 $x, y \in V$. 用 $\CC$ 的范数 $1$ 的元素乘 $y$ (这不改变 $\|y\|$) 后, 可以假设 $\langle T(x), y \rangle \in \R$. 则
\[
\langle T(x), y \rangle = \frac{1}{4}(\langle T(x + y), x + y \rangle - \langle T(x - y), x - y \rangle).
\]
所以, 利用平行四边形恒等式,
\begin{multline*}
|\langle T(x), y \rangle| \leq \frac{c}{4}(\|x + y\|^2 + \|x - y\|^2) \\
= \frac{c}{2}(\|x\|^2 + \|y\|^2).
\end{multline*}
这证明所要的结果.
\end{proof}

这里有一些对更一般情形的紧算子成立的结果 (见 {[}20{]} 4.16-4.25).

\begin{theorem}[紧算子谱的性质]
\label{thm:IV.1.5}
设 $V$ 是巴拿赫空间, $T \in \End(V)$ 是紧自同态. 记 $\sigma(T)$ 为 $T$ 在 $\End(V)$ 中的谱, 即
\[
\sigma(T) = \{\lambda \in \CC \mid \lambda\mathrm{id}_V - T \notin \End(V)^{\times}\}.
\]
则:
\begin{enumerate}
\item[(i)] 对每个 $\lambda \neq 0$, $T - \lambda\mathrm{id}_V$ 的像闭.
\item[(ii)] 对每个 $\lambda \in \sigma(T) - \{0\}$, 有 $\Im(T - \lambda\mathrm{id}_V) \neq V$ 且 $\dim(\Ker(T - \lambda\mathrm{id}_V)) = \dim(V/\Im(T - \lambda\mathrm{id}_V))$.\footnote{注意这推广了秩-零度定理.} 特别地, $\Ker(T - \lambda\mathrm{id}_V) \neq \{0\}$.
\item[(iii)] 对每个 $\lambda \neq 0$, 递增序列 $(\Ker((T - \lambda\mathrm{id}_V)^n))_{n \geq 1}$ 稳定, 且其极限有限维.
\item[(iv)] 若 $\dim_{\CC}V = +\infty$, 则 $0 \in \sigma(T)$.
\item[(v)] $\CC - \{0\}$ 的子集 $\sigma(T) - \{0\}$ 离散. 特别地, 对每个 $r > 0$, 只有有限多个 $\lambda \in \sigma(T)$ 满足 $|\lambda| \geq r$.
\end{enumerate}
特别地, 若 $V$ 是希尔伯特空间且 $T$ 自伴, 则定理 \ref{thm:IV.1.3} 的 (v) 变成
\[
V = \bigoplus_{\lambda \in \sigma(T)}\Ker(T - \lambda\mathrm{id}_V).
\]
\end{theorem}

\section{紧群的酉表示的半单性}
本节的目标是证明下面的定理. (与命题 \ref{prop:I.3.3.3} 比较.)

\begin{theorem}[紧群酉表示的半单性]
\label{thm:IV.2.1}
设 $G$ 是紧群, $V$ 是 $G$ 的酉表示. 则存在 $V$ 的两两正交的子表示族 $(W_i)_{i \in I}$ 使得每个 $W_i$ 不可约, 且
\[
V = \bigoplus_{i \in I}W_i.
\]
\end{theorem}

我们已经在习题 \ref{ex:I.5.5.9} 中见过关键的构造. 让我们把它总结成一个命题.

\begin{proposition}[等变自伴紧算子]
\label{prop:IV.2.2}
(见习题 \ref{ex:I.5.5.9}.) 设 $G$ 是紧群, $dx$ 是 $G$ 上的规范化 Haar 测度, $(\pi, V)$ 是 $G$ 的酉表示. 若 $u \in V$, 则公式
\[
T(v) = \int_G \langle v, \pi(x)(u) \rangle\pi(x)(u)\,\dd x
\]
定义 $V$ 的连续 $G$-等变自伴紧自同态, 且 $T = 0$ 当且仅当 $u = 0$.
\end{proposition}

事实上, 我们甚至知道 $T$ 正定, 即对每个 $v \in V$ 有 $\langle T(v), v \rangle \geq 0$.

\begin{corollary}[不可约子表示的存在性]
\label{cor:IV.2.3}
设 $V$ 是紧群 $G$ 的非零酉表示. 则 $V$ 包含 $G$ 的不可约表示.
\end{corollary}

\begin{proof}
若 $V$ 有限维, 则 $V$ 的任何极小维数的非零 $G$-不变子空间必不可约.
一般情形下, 选取 $u \in V - \{0\}$, 设 $T \in \End(V)$ 是命题中构造的 $V$ 的自同态. 由自伴紧算子的谱定理 (定理 \ref{thm:IV.1.3}),
\[
V = \bigoplus_{\lambda \in \CC}\Ker(T - \lambda\mathrm{id}_V).
\]
由于 $T \neq 0$, $V$ 的闭子空间 $\Ker(T)$ 不等于 $V$. 由上面的等式, 存在 $\lambda \in \CC - \{0\}$ 使得 $W := \Ker(T - \lambda\mathrm{id}_V) \neq 0$. 则 $W$ 是 $V$ 的非零闭子空间, 且它 $G$-不变 (因为 $T$ 是 $G$-等变的) 且由定义被 $T$ 稳定. 又, 由定理 \ref{thm:IV.1.3}(iii), 空间 $W$ 有限维. 所以由证明的开头, $W$ 有不可约子表示, 我们完成了.
\end{proof}

\begin{proof}[定理的证明]
由 Zorn 引理, 可以找到 $V$ 的两两正交的不可约子表示的极大族 $(W_i)_{i \in I}$. 假设 $W_i$ 的直和不在 $V$ 中稠密, 则 $W := \bigl(\bigoplus_{i \in I}W_i\bigr)^{\perp}$ 是 $V$ 的非零闭不变子空间 (见引理 \ref{lem:I.3.2.6}). 由上面的推论, 表示 $W$ 有不可约子表示, 这与族 $(W_i)_{i \in I}$ 的极大性矛盾. 所以 $V = \bigoplus_{i \in I}W_i$.
\end{proof}

我们用关于酉表示的两种等价概念的一个注记结束本节. 回忆拓扑群 $G$ 的两个连续表示 $V_1$ 和 $V_2$ 称为等价 (equivalent) (或同构 (isomorphic)), 如果存在有连续逆的连续 $G$-等变同构 $V_1 \to V_2$.

\begin{definition}[\index{酉等价}酉等价 (unitarily equivalent)]
\label{def:IV.2.4}
拓扑群 $G$ 的两个酉表示 $V_1$ 和 $V_2$ 称为酉等价, 如果存在是等距的 $G$-等变同构 $V_1 \to V_2$.
\end{definition}

两个酉等价的表示显然等价.

\begin{example}[左右正则表示的酉等价]
\label{exa:IV.2.5}
设 $G$ 是局部紧群, $\mu$ 是 $G$ 上的左 Haar 测度, $\nu$ 是由 $\nu(E) = \mu(E^{-1})$ 定义的右 Haar 测度.
则 $G$ 的左、右正则表示酉等价, 同构把 $f \in L^2(G, \mu)$ 送到 $L^2(G, \nu)$ 的元素 $x \longmapsto \Delta(x)^{-1/2}\overline{f(x^{-1})}$. (见命题 \ref{prop:I.2.12}.)
\end{example}

\begin{proposition}[不可约表示的等价与酉等价]
\label{prop:IV.2.6}
假设 $V_1$ 和 $V_2$ 是 $G$ 的不可约酉表示. 则它们等价当且仅当它们酉等价.
\end{proposition}

\begin{proof}
假设 $V_1$ 和 $V_2$ 等价, 设 $U \colon V_1 \to V_2$ 是 $G$-等变同构. 记 $\langle \cdot, \cdot \rangle_1$ 和 $\langle \cdot, \cdot \rangle_2$ 为 $V_1$ 和 $V_2$ 的内积. 设 $B \colon V_1 \times V_1 \to \CC$, $(v, w) \longmapsto \langle U(v), U(w) \rangle_2$. 这是 $V_1$ 上的 Hermitian 半双线性形式, 且因 $U$ 有界而有界. 由下面的引理, 存在自伴自同态 $T \in \End(V_1)$ 使得对所有 $v, w \in V_1$ 有 $B(v, w) = \langle T(v), w \rangle_1$. 让我们证明 $T$ 是 $G$-等变的. 设 $v \in V_1$, $x \in G$. 对每个 $w \in V_1$,
\begin{align*}
\langle T(\pi_1(x)(v)), w \rangle_1 &= B(\pi_1(x)(v), w) \\
&= \langle U(\pi_1(x)(v)), U(w) \rangle_2 \\
&= \langle \pi_2(x)(U(v)), U(w) \rangle_2 \\
&= \langle U(v), \pi_2(x^{-1})U(w) \rangle_2 \\
&= \langle U(v), U(\pi_1(x)^{-1}(w)) \rangle_2 \\
&= B(v, \pi_1(x^{-1})(w)) \\
&= \langle T(v), \pi_1(x)^{-1}(w) \rangle_1 \\
&= \langle \pi_1(x)(T(v)), w \rangle_1,
\end{align*}
所以 $T(\pi_1(x)(v)) = \pi_1(x)(T(v))$. 由于 $V_1$ 不可约, Schur 引理 (定理 \ref{thm:I.3.4.1}) 蕴含对某个 $\lambda \in \CC$ 有 $T = \lambda\mathrm{id}_{V_1}$. 由于对每个非零 $v \in V_1$ 有 $\langle T(v), v \rangle_1 = \langle U(v), U(v) \rangle_2 > 0$, 必有 $\lambda \in \R_{>0}$. 则 $\lambda^{-1/2}U$ 是等距, 所以 $V_1$ 和 $V_2$ 酉等价.
\end{proof}

\begin{lemma}[有界半双线性形式]
\label{lem:IV.2.7}
设 $V$ 是希尔伯特空间, $B \colon V \times V \to \CC$ 是有界半双线性形式 (即 $B$ 对第一个变量 $\CC$-线性、对第二个变量 $\CC$-反线性; 有界性条件指 $\sup_{v,w \in V, \|v\| = \|w\| = 1}|B(v, w)| < +\infty$).
则存在唯一的 $T \in \End(V)$ 使得对所有 $v, w \in V$,
\[
B(v, w) = \langle T(v), w \rangle.
\]
而且, $T$ 自伴当且仅当 $B$ Hermitian (即对所有 $v, w \in V$ 有 $B(w, v) = \overline{B(v, w)}$).
\end{lemma}

\begin{proof}
$T$ 的唯一性显然 (它由 $V^{\perp} = \{0\}$ 这一事实得出.)
若 $v \in V$, 则映射 $V \to \CC$, $w \longmapsto B(v, w)$ 是 $V$ 上的连续线性泛函, 所以存在唯一的 $T(v) \in V$ 使得对每个 $w \in V$ 有 $B(v, w) = \langle T(v), w \rangle$. $T$ 的线性由 $B$ 对第一个变量线性得出. 而且, 对每个 $v \in V$,
\[
\|T(v)\| = \sup_{w \in V, \|w\| = 1}|\langle T(v), w \rangle| = \sup_{w \in V, \|w\| = 1}|B(v, w)| \leq C\|v\|,
\]
其中
\[
C = \sup_{x,y \in V, \|x\| = \|y\| = 1}|B(x, y)|.
\]
所以 $T$ 有界.
最后, $T$ 自伴当且仅当对所有 $v, w \in V$,
\[
B(v, w) = \langle T(v), w \rangle = \langle v, T(w) \rangle = \overline{B(w, v)}.
\]
这证明最后一个结论.
\end{proof}

\begin{definition}[\index{酉对偶}酉对偶 (unitary dual)]
\label{def:IV.2.8}
记 $\widehat{G}$ 为 $G$ 的不可约酉表示的等价 (或酉等价) 类的集合, 称之为 $G$ 的酉对偶.
若 $(\pi, V) \in \widehat{G}$, 记 $\dim(\pi)$ 和 $\End(\pi)$ 为 $\dim(V)$ 和 $\End(\pi)$.
\end{definition}

注意此记号与习题 \ref{ex:I.5.4.1} 中对交换群使用的记号一致.

\section{\index{矩阵系数}矩阵系数}

\begin{definition}[矩阵系数 (matrix coefficient)]
\label{def:IV.3.1}
设 $(\pi, V)$ 是拓扑群 $G$ 的酉表示. $(\pi, V)$ 的矩阵系数是形如 $x \longmapsto \langle \pi(x)(u), v \rangle$ ($u, v \in V$) 的函数 $G \to \CC$.
\end{definition}

注意矩阵系数是连续函数. 记 $E_{\pi}$ 或 $E_V$ 为 $\mathcal{C}(G)$ 中由 $\pi$ 的矩阵系数张成的子空间.
我们先证明对任意群 $G$ 成立的一些一般结果.

\begin{proposition}[矩阵系数的基本性质]
\label{prop:IV.3.2}
设 $(\pi, V)$ 是 $G$ 的酉表示.
\begin{enumerate}
\item[(i)] $\mathcal{C}(G)$ 的子空间 $E_{\pi}$ 只依赖 $\pi$ 的酉等价类, 且对每个 $x \in G$ 被算子 $L_x$ 和 $R_x$ 不变.
\item[(ii)] 若 $V$ 有限维, 则 $E_{\pi}$ 有限维且 $\dim(E_{\pi}) \leq (\dim V)^2$.
\item[(iii)] 若 $V = V_1 \oplus \cdots \oplus V_n$ 且 $V_i$ 两两正交且 $G$-不变, 则 $E_{\pi} = \sum_{i=1}^n E_{V_i}$.
\item[(iv)] 有 $E_{\pi^*} = \overline{E_{\pi}}$.
\end{enumerate}
特别地, 我们得到 $G \times G$ 在 $E_{\pi}$ 上的作用, 由 $(x, y) \in G \times G$ 按 $L_x \circ R_y = R_y \circ L_x$ 作用.
\end{proposition}

\begin{proof}
(i). 第一个结论显然. 为证明第二个结论, 设 $v, w \in V$, $x \in G$. 则对每个 $y \in G$,
\[
\langle \pi(x^{-1}y)(v), w \rangle = \langle \pi(y)(v), \pi(x)(w) \rangle
\]
且
\[
\langle \pi(yx)(v), w \rangle = \langle \pi(y)(\pi(x)(v)), w \rangle,
\]
所以函数 $y \longmapsto \langle \pi(x^{-1}y)(v), w \rangle$ 和 $y \longmapsto \langle \pi(yx)(v), w \rangle$ 也是 $\pi$ 的矩阵系数.
(ii). 设 $(e_1, \ldots, e_n)$ 是 $V$ 的基. 对 $i, j \in \{1, \ldots, n\}$, 记 $\phi_{ij}$ 为函数 $G \to \CC$, $x \longmapsto \langle \pi(x)(e_i), e_j \rangle$. 若 $v, w \in V$, 可以把 $v$ 写成 $v = \sum_{i=1}^n a_ie_i$, $w = \sum_{j=1}^n b_je_j$, 然后对每个 $x \in G$,
\[
\langle \pi(x)(v), w \rangle = \sum_{i,j=1}^n a_i\bar{b}_j\phi_{ij}(x).
\]
所以族 $(\phi_{ij})_{1 \leq i,j \leq n}$ 张成 $E_{\pi}$.
(iii). 对每个 $i \in \{1, \ldots, n\}$, 显然有 $E_{V_i} \subset E_{\pi}$. 所以 $\sum_{i=1}^n E_{V_i} \subset E_{\pi}$. 反之, 设 $v, w \in V$, 把 $v$ 写成 $v = \sum_{i=1}^n v_i$, $w = \sum_{i=1}^n w_i$, 其中 $v_i, w_i \in V_i$. 则对每个 $x \in G$,
\[
\langle \pi(x)(v), w \rangle = \sum_{i,j=1}^n\langle \pi(x)(v_i), w_j \rangle = \sum_{i=1}^n\langle \pi(x)(v_i), w_i \rangle.
\]
所以函数 $x \longmapsto \langle \pi(x)(v), w \rangle$ 在 $\sum_{i=1}^n E_{V_i}$ 中.
\end{proof}

\begin{definition}[$\Hom(V, V')$ 上的 $G \times G$ 作用]
\label{def:IV.3.3}
设 $(\pi, V)$ 和 $(\pi', V')$ 是 $V$ 的连续表示. 定义 $G \times G$ 在 $\Hom(V, V')$ 上的作用 $\rho$ 为
\[
\rho(x, y)(T) = \pi'(y) \circ T \circ \pi(x)^{-1},
\]
其中 $T \in \Hom(V, V')$, $x, y \in G$.
\end{definition}

\begin{proposition}[交错算子与不动点]
\label{prop:IV.3.4}
有
\[
\Hom_G(V, V') = \{T \in \Hom(V, V') \mid \forall x \in G,\, \rho(x, x)(T) = T\}.
\]
而且, 若映射 $G \to \End(V)$, $x \longmapsto \pi(x)$ 和 $G \to \End(V')$, $x \longmapsto \pi'(x)$ 连续 (例如若 $V$ 和 $V'$ 有限维, 见命题 \ref{prop:I.3.5.1}), 则上面定义的作用是 $G \times G$ 在 $\Hom(V, V')$ 上的连续表示.
\end{proposition}

\begin{proof}
第一个结论显然. 第二个结论由 $\Hom$ 空间上复合的连续性以及 $G$ 上取逆的连续性得出.
\end{proof}

特别地, 我们得到 $G \times G$ 在 $\End(V)$ 和 $V^* := \Hom(V, \CC)$ 上的作用 (用 $G$ 在 $\CC$ 上的平凡作用); 第二个限制到第一个因子上给出 $G$ 在 $V^*$ 上的作用 (若 $x \in G$, $\Lambda \in V^*$, 则 $(x, \Lambda)$ 被送到 $\Lambda \circ \pi(x)^{-1}$). 这将是这些空间上的默认作用.

\begin{definition}[张量积上的 $G \times G$ 作用]
\label{def:IV.3.5}
设 $(\pi, V)$ 和 $(\pi', V')$ 是 $V$ 的连续表示. 定义 $G \times G$ 在代数张量积 $V \otimes_{\CC}V'$ 上的作用 $\rho$ 为
\[
\rho(x, y)(v \otimes w) = \pi(x)(v) \otimes \pi'(y)(w),
\]
其中 $x, y \in G$, $v \in V$, $w \in V'$.
\end{definition}

这个作用良定义, 因为对所有 $x, y \in G$, 映射 $V \times V' \to V \otimes_{\CC}V'$, $(v, w) \longmapsto \pi(x)(v) \otimes \pi'(y)(w)$ 双线性, 所以它诱导映射 $\rho(x, y) \colon V \otimes_{\CC}V' \to V \otimes_{\CC}V'$. 若 $V$ 和 $V'$ 有限维, 则由命题 \ref{prop:I.3.5.1}, $G \times G$ 在 $V \otimes_{\CC}V'$ 上的作用连续.
注意, 若把 $G \times G$ 在 $V \otimes_{\CC}W$ 上的作用限制到第一个 (相应地第二个) 因子, 得到等价于 $V^{\oplus \dim(W)}$ (相应地 $W^{\oplus \dim(V)}$) 的表示.

\begin{proposition}[$V^* \otimes W$ 与 $\Hom(V, W)$]
\label{prop:IV.3.6}
设 $V, W$ 是 $G$ 的连续表示.
\begin{enumerate}
\item[(i)] 映射 $V^* \otimes_{\CC}W \to \Hom(V, W)$, 把 $\Lambda \otimes w$ ($\Lambda \in V^*$, $w \in W$) 送到线性算子 $V \to W$, $v \longmapsto \Lambda(v)w$ 良定义且 $G \times G$-等变. 若 $V$ 和 $W$ 有限维, 它是连续表示的等价.
\item[(ii)] 映射 $V^* \otimes_{\CC}V \to \mathcal{C}(G)$, 把 $\Lambda \otimes v$ ($\Lambda \in V^*$, $v \in V$) 送到函数 $G \to \CC$, $x \longmapsto \Lambda(\pi(x)(v))$ 良定义且 $G \times G$-等变, 且若 $V$ 酉, 其像是 $E_V$.
\end{enumerate}
特别地, 若 $V$ 有限维且酉, 得到满的 $G \times G$-等变映射 $\End(V) \to E_V$.
\end{proposition}

\begin{remark}[非酉情形的矩阵系数]
\label{rem:IV.3.7}
(ii) 提示了把矩阵系数的定义推广到非酉情形的方法: 只把矩阵系数定义为映射 $V^* \otimes_{\CC}V \to \mathcal{C}(G)$ 下纯张量的像.
\end{remark}

\begin{proof}
在此证明中, 我们用 $\cdot$ 记 $G$ 和 $G \times G$ 的所有作用 (这不应引起混淆, 因为每个空间至多有一个作用).
(i). 映射良定义, 因为映射 $V^* \otimes_{\CC}W \to \Hom(V, W)$, 把 $(\Lambda, w)$ 送到 $(v \longmapsto \Lambda(v)w)$ 双线性. 记之为 $\phi$. 为验证它是 $G \times G$-等变的, 只需在纯张量上验证 (因为它们生成 $V^* \otimes_{\CC}W$). 所以设 $\Lambda \in V^*$, $w \in W$, $x, y \in G$. 对每个 $v \in V$,
\[
\phi((x, y) \cdot (\Lambda \otimes w))(v) = \phi((y \cdot \Lambda) \otimes (x \cdot w))(v) = \Lambda(y^{-1} \cdot v)(x \cdot w)
\]
且
\[
((x, y) \cdot \phi(\Lambda \otimes w))(v) = x \cdot (\phi(\Lambda \otimes w)(y^{-1} \cdot v)) = x \cdot (\Lambda(y^{-1} \cdot v)w) = \Lambda(y^{-1} \cdot v)(x \cdot w).
\]
所以
\[
\phi(x \cdot (\Lambda \otimes w)) = x \cdot \phi(\Lambda \otimes w).
\]
假设 $V$ 有限维, 设 $(e_1, \ldots, e_n)$ 是 $V$ 的基, $(e_1^*, \ldots, e_n^*)$ 是对偶基. 定义 $\psi \colon \Hom(V, W) \to V^* \otimes_{\CC}W$, 把 $T$ 送到 $\sum_{i=1}^n e_i^* \otimes T(e_i)$. 让我们证明 $\psi$ 是 $\phi$ 的逆.
若 $j \in \{1, \ldots, m\}$, $w \in W$, 则
\[
\psi(\phi(e_j^* \otimes w)) = \sum_{i=1}^n e_i^* \otimes (\phi(e_j^* \otimes w)(e_i)) = e_j^* \otimes w.
\]
由于元素 $e_j^* \otimes w$ ($j \in \{1, \ldots, n\}$, $w \in W$) 生成 $V^* \otimes_{\CC}W$, 这证明 $\psi \circ \phi = \mathrm{id}$.
反之, 若 $T \in \Hom(V, W)$, 则对每个 $v \in V$,
\[
\phi(\psi(T)) = \sum_{i=1}^n\phi(e_i^* \otimes T(e_i))(v) = \sum_{i=1}^n e_i^*(v)T(e_i) = T(v),
\]
因为 $v = \sum_{i=1}^n e_i^*(v)e_i$. 所以 $\phi(\psi(T)) = T$.
这证明若 $V$ 有限维, 映射 $V^* \otimes_{\CC}W \to \Hom(V, W)$ 是同构. 最后一个结论由此立即得出.
(ii). 映射良定义, 因为映射 $V^* \times V \to \mathcal{C}(G)$, 把 $(\Lambda, v)$ 送到函数 $x \longmapsto \Lambda(\pi(x)(v))$ 双线性. 记之为 $\alpha$. 我们证明 $\alpha$ 是 $G \times G$-等变的. 如前, 只需在纯张量上验证. 所以设 $\Lambda \in V^*$, $v \in V$, $x, y \in G$. 对每个 $z \in G$,
\[
\alpha((x, y) \cdot (\Lambda \otimes v))(z) = \Lambda(x^{-1} \cdot (z \cdot (y \cdot v))) = \Lambda((x^{-1}zy) \cdot v) = ((L_x \circ R_y)(\alpha(\Lambda \otimes v)))(z),
\]
所以 $\alpha((x, y) \cdot (\Lambda \otimes v)) = (L_x \circ R_y)(\alpha(\Lambda \otimes v))$.
最后, 我们证明若 $V$ 酉, $\alpha$ 的像是 $E_V$. 设 $\Lambda \in V^*$. 由于 $V$ 是希尔伯特空间, 存在唯一的 $v \in V$ 使得 $\Lambda = \langle \cdot, v \rangle$. 所以对每个 $w \in V$ 和每个 $x \in G$,
\[
\alpha(\Lambda \otimes w)(x) = \langle \pi(x)(w), v \rangle.
\]
这证明 $\alpha(\Lambda \otimes w)$ 是 $\pi$ 的矩阵系数, 也证明我们以这种方式得到 $\pi$ 的所有矩阵系数.
\end{proof}

现在我们证明只对紧群成立的更强结果. 若 $G$ 是紧群, 固定 $G$ 上的规范化 Haar 测度, 记 $L^p(G)$ 为此测度的 $L^p$ 空间. 注意对每个 $p$ 有 $\mathcal{C}(G) \subset L^p(G)$.

\begin{theorem}[Schur 正交性]
\label{thm:IV.3.8}
设 $G$ 是紧群, $(\pi, V)$ 是 $G$ 的不可约酉表示. 回忆 $V$ 有限维 (由习题 \ref{ex:I.5.5.9}).
\begin{enumerate}
\item[(i)] (Schur 正交性 (Schur orthogonality)) 若 $(\pi', V')$ 是 $G$ 的另一个不等价于 $(\pi, V)$ 的不可约酉表示, 则 $E_{\pi}$ 和 $E_{\pi'}$ 作为 $L^2(G)$ 的子空间正交.
\item[(ii)] 有 $\dim(E_{\pi}) = (\dim V)^2$. 更精确地, 若 $(e_1, \ldots, e_d)$ 是 $V$ 的标准正交基, 记 $\phi_{ij}$ 为函数 $G \to \CC$, $x \longmapsto \langle \pi(x)(e_j), e_i \rangle$, 则 $\{\sqrt{d}\phi_{ij}, 1 \leq i, j \leq d\}$ 是 $E_{\pi}$ 对 $L^2$ 内积的标准正交基.
\item[(iii)] 上面定义的 $G \times G$-等变映射 $\End(V) \to E_{\pi}$ 是同构.
\end{enumerate}
\end{theorem}

\begin{proof}
注意 (iii) 由 (ii) 立即得出, 因为 $\End(V) \to E_{\pi}$ 满射且 (ii) 说 $\dim(E_{\pi}) = (\dim V)^2 = \dim(\End(V))$.
我们证明 (i) 和 (ii). 设 $(\pi', V')$ 是 $G$ 的不可约酉表示, 可以等于 $(\pi, V)$. 若 $A \in \Hom(V, V')$, 定义 $\tilde{A} \in \Hom(V, V')$ 为
\[
\tilde{A} = \int_G \pi'(x)^{-1} \circ A \circ \pi(x)\,\dd x
\]
(注意积分没有问题, 因为表示有限维). 则对每个 $y \in G$,
\begin{align*}
\tilde{A} \circ \pi(y) &= \int_G \pi'(x)^{-1} \circ A \circ \pi(xy)\,\dd x \\
&= \int_G \pi'(xy^{-1})^{-1} \circ A \circ \pi(x)\,\dd x \quad \text{($dx$ 的右不变性)} \\
&= \pi'(y) \circ \tilde{A}.
\end{align*}
换句话说, $\tilde{A}$ 是 $G$-等变的.
设 $v \in V$, $v' \in V'$, 定义 $A \in \Hom(V, V')$ 为 $A(u) = \langle u, v \rangle v'$. 则对所有 $u \in V$ 和 $u' \in V'$,
\begin{align*}
\langle \tilde{A}(u), u' \rangle &= \int_G \langle (\pi'(x)^{-1} \circ A \circ \pi(x))(u), u' \rangle\,\dd x \\
&= \int_G \langle \langle \pi(x)(u), v \rangle\pi'(x)^{-1}(v'), u' \rangle\,\dd x \\
&= \int_G \overline{\langle \pi(x)(u), v \rangle}\langle \pi'(x)(u'), v' \rangle\,\dd x.
\end{align*}
假设 $\pi$ 和 $\pi'$ 不等价. 则由 Schur 引理, 对每个 $A \in \Hom(V, V')$ 有 $\tilde{A} = 0$, 所以由上面的计算, 对所有 $u, v \in V$ 和 $u', v' \in V'$,
\[
\int_G \overline{\langle \pi(x)(u), v \rangle}\langle \pi'(x)(u'), v' \rangle\,\dd x = 0.
\]
这证明 (i).
假设 $\pi = \pi'$, 用 (ii) 的记号. 取 $v = e_i$, $v' = e_{i'}$ ($i, i' \in \{1, \ldots, d\}$), 如前定义 $A$. 再次由 Schur 引理, 存在 $c \in \CC$ 使得 $\tilde{A} = c\,\mathrm{id}_V$. 所以取 $u = e_j$, $u' = e_{j'}$, 由上面的计算得到
\[
\langle \phi_{i,j}, \phi_{i',j'} \rangle_{L^2(G)} = \langle ce_j, e_{j'} \rangle = \begin{cases} c & \text{若 } j = j', \\ 0 & \text{否则}. \end{cases}
\]
另一方面,
\[
cd = \operatorname{Tr}(\tilde{A}) = \int_G \operatorname{Tr}(\pi(x)^{-1} \circ A \circ \pi(x))\,\dd x = \int_G \operatorname{Tr}(A)\,\dd x = \operatorname{Tr}(A).
\]
由于 $A$ 由 $A(w) = \langle w, e_i \rangle e_{i'}$ 定义, 若 $i \neq i'$ 则 $\operatorname{Tr}(A) = 0$, 若 $i = i'$ 则 $\operatorname{Tr}(A) = 1$. 这完成 $\{\sqrt{d}\phi_{ij}, 1 \leq i, j \leq d\}$ 是 $E_{\pi}$ 的标准正交基的证明.
\end{proof}

\section{Peter-Weyl 定理}
设 $G$ 是紧群. 我们把 $L^2(G)$ 看作 $G \times G$ 的表示, 让 $(x, y) \in G \times G$ 按 $L_x \circ R_y = R_y \circ L_x$ 作用. 这限制到第一个 (相应地第二个) 因子是 $G$ 的左 (相应地右) 正则表示.

\begin{definition}[矩阵系数的空间 $E$]
\label{def:IV.4.1}
设 $G$ 是紧群. 记 $E$ 为 $\mathcal{C}(G)$ 中矩阵系数的子空间.
由定理 \ref{thm:IV.3.8}, $E = \bigoplus_{\pi \in \widehat{G}}E_{\pi}$.
\end{definition}

\begin{theorem}[Peter-Weyl 定理: $E$ 稠密]
\label{thm:IV.4.2}
若 $G$ 紧, 则 $E$ 是 $\mathcal{C}(G)$ 的稠密子代数. (对通常的逐点乘法和范数 $\|\cdot\|_{\infty}$.)
\end{theorem}

\begin{proof}
让我们证明 $E$ 被乘法稳定. 注意由命题 \ref{prop:IV.3.2}(iii) 和定理 \ref{thm:IV.2.1}, 对 $G$ 的每个有限维酉表示 $\pi$, 有 $E_{\pi} \subset E$. 设 $(\pi, V)$ 和 $(\pi', V')$ 是 $G$ 的不可约酉表示, $v, w \in V$, $v', w' \in V'$. 回忆我们在习题 \ref{ex:III.6.2.1} 中定义了 $G$ 在 $V \otimes_{\CC}V'$ 上的作用 $\pi \otimes \pi'$\footnote{这只是上面定义的 $G \times G$ 的作用限制到 $G \times G$ 的对角线上.} 以及 $V \otimes_{\CC}V'$ 上的内积.\footnote{这里我们不需要完备化张量积, 因为 $V$ 和 $V'$ 有限维.} 由这些定义, 对每个 $x \in G$,
\[
\langle (\pi \otimes \pi')(x)(v \otimes w), v' \otimes w' \rangle = \langle \pi(x)(v), w \rangle\langle \pi'(x)(v'), w' \rangle.
\]
这证明 $\pi$ 的矩阵系数与 $\pi'$ 的矩阵系数的乘积是 $\pi \otimes \pi'$ 的矩阵系数. 由上面的观察, $\pi \otimes \pi'$ 的每个矩阵系数都在 $E$ 中, 我们完成了.
现在我们证明 $E$ 在 $\mathcal{C}(G)$ 中稠密. 我们已经证明 $E$ 是子代数, 它由命题 \ref{prop:IV.3.2}(iv) 被复共轭稳定, 它包含常数 (它们是 $G$ 在 $\CC$ 上的平凡表示的矩阵系数), 且它由 Gelfand-Raikov 定理 (定理 \ref{thm:III.5.1}) 分离 $G$ 上的点. 所以由 Stone-Weierstrass 定理它在 $\mathcal{C}(G)$ 中稠密.
\end{proof}

\begin{corollary}[$L^p$ 中的稠密性]
\label{cor:IV.4.3}
对每个 $p \in [1, +\infty)$, 子空间 $E$ 在 $L^p(G)$ 中对 $L^p$ 范数稠密.
特别地, 有典则 $G \times G$-等变同构
\[
L^2(G) = \bigoplus_{(\pi, V) \in \widehat{G}}\End(V).
\]
\end{corollary}

最后一个结论通常称为 Peter-Weyl 定理. 它蕴含 $G$ 的左、右正则表示都同构于 $\bigoplus_{\pi \in \widehat{G}}\pi^{\oplus \dim(\pi)}$ 的完备化.

\begin{remark}[经典证明]
\label{rem:IV.4.4}
Peter-Weyl 定理实际上早于 Gelfand-Raikov 定理, 且原始证明用算子 $f * \cdot$ ($f \in L^2(G)$) 在 $L^2(G)$ 上紧这一事实.
\end{remark}

\section{\index{特征标}特征标}

\begin{definition}[特征标 (character)]
\label{def:IV.5.1}
设 $(\pi, V)$ 是拓扑群 $G$ 的连续有限维表示. $\pi$ 的特征标 (character) 是连续映射 $\chi_V = \chi_{\pi} \colon G \to \CC$, $x \longmapsto \operatorname{Tr}(\pi(x))$.
\end{definition}

\begin{remark}[特征标与矩阵系数]
\label{rem:IV.5.2}
若 $(\pi, V)$ 是 $G$ 的有限维表示, $(e_1, \ldots, e_n)$ 是 $V$ 的标准正交基, 则对每个 $x \in G$,
\[
\chi_{\pi}(x) = \sum_{i=1}^n\langle \pi(x)(e_i), e_i \rangle.
\]
所以 $\chi_{\pi} \in E_{\pi}$.
\end{remark}

\begin{definition}[\index{中心函数}中心函数 (central function)]
\label{def:IV.5.3}
若对所有 $x, y \in G$ 有 $f(xyx^{-1}) = f(y)$, 则称函数 $f \colon G \to \CC$ 中心 (central) 或\index{类函数}类函数 (class function).
\end{definition}

这些函数被称为中心的, 因为它们是卷积乘积的中心元, 我们将在第 IV.7 节看到.

\begin{proposition}[特征标的性质]
\label{prop:IV.5.4}
设 $G$ 是拓扑群, $(\pi, V)$ 和 $(\pi', V')$ 是 $G$ 的连续有限维表示. 则:
\begin{enumerate}
\item[(i)] $\chi_{\pi}$ 是中心函数, 且只依赖 $\pi$ 的等价类.
\item[(ii)] $\chi_{V \oplus V'} = \chi_V + \chi_{V'}$.
\item[(iii)] 对每个 $x \in G$, $\chi_{V^*}(x) = \chi_V(x^{-1})$.
\item[(iv)] 对所有 $x, y \in G$,
\[
\chi_{V \otimes_{\CC}V'}(x, y) = \chi_V(x)\chi_{V'}(y) \quad \text{且} \quad \chi_{\Hom(V,V')}(x, y) = \chi_V(x^{-1})\chi_{V'}(y).
\]
\item[(v)] 若 $(\pi, V)$ 可酉化 (例如若 $G$ 紧), 则对每个 $x \in G$ 有 $\chi_V(x^{-1}) = \overline{\chi_V(x)}$.
\end{enumerate}
\end{proposition}

\begin{proof}
(i) 由迹的性质, 即对所有 $A, B \in M_n(\CC)$ 有 $\operatorname{Tr}(AB) = \operatorname{Tr}(BA)$ 得出.
在 $V$ 和 $V'$ 上放任意 Hermitian 内积. 设 $(e_1, \ldots, e_n)$ (相应地 $(e'_1, \ldots, e'_m)$) 是 $V$ (相应地 $V'$) 的标准正交基. 则 $(e_1, \ldots, e_n, e'_1, \ldots, e'_m)$ 是 $V \oplus V'$ 的标准正交基, 所以对每个 $x \in G$,
\[
\chi_{V \oplus V'}(x) = \sum_{i=1}^n\langle \pi(x)(e_i), e_i \rangle + \sum_{j=1}^m\langle \pi'(x)(e'_j), e'_j \rangle = \chi_V(x) + \chi_{V'}(x).
\]
这证明 (ii).
设 $(e_1^*, \ldots, e_n^*)$ 是 $(e_1, \ldots, e_n)$ 的对偶基. 设 $x \in G$. 则若 $A$ 是 $\pi(x^{-1})$ 在基 $(e_1, \ldots, e_n)$ 中的矩阵, $V^*$ 的自同态 $\Lambda \longmapsto \Lambda \circ \pi(x^{-1})$ 在基 $(e_1^*, \ldots, e_n^*)$ 中的矩阵是 $A^T$, 且有
\[
\chi_{V^*}(x) = \operatorname{Tr}(A^T) = \operatorname{Tr}(A) = \chi_V(x^{-1}).
\]
这证明 (iii).
我们证明 $\chi_{V \otimes_{\CC}V'}$ 的公式. 我们在习题 \ref{ex:III.6.2.1} 中见过如何在 $V \otimes_{\CC}V'$ 上放使 $(e_i \otimes e'_j)_{1 \leq i \leq n, 1 \leq j \leq m}$ 成为标准正交基的内积. 所以对所有 $x, y \in G$,
\begin{align*}
\chi_{V \otimes_{\CC}V'}(x, y) &= \sum_{i=1}^n\sum_{j=1}^m\langle \pi(x)(e_i) \otimes \pi'(y)(e'_j), e_i \otimes e'_j \rangle \\
&= \sum_{i=1}^n\sum_{j=1}^m\langle \pi(x)(e_i), e_i \rangle\langle \pi'(y)(e'_j), e'_j \rangle \\
&= \chi_V(x)\chi_{V'}(y).
\end{align*}
现在 $\chi_{\Hom(V,V')}$ 的公式由它、由 (iii) 和命题 \ref{prop:IV.3.6}(i) 得出.
最后, 我们证明 (v). 若 $V$ 可酉化, 可以选择 $V$ 上被 $G$ 不变的 Hermitian 内积. 则对每个 $x \in G$,
\[
\chi_V(x^{-1}) = \sum_{i=1}^n\langle \pi(x^{-1})(e_i), e_i \rangle = \sum_{i=1}^n\langle e_i, \pi(x)(e_i) \rangle = \sum_{i=1}^n\overline{\langle \pi(x)(e_i), e_i \rangle} = \overline{\chi_V(x)}.
\]
\end{proof}

\begin{notation}[不动点子空间]
\label{not:IV.5.5}
若 $(\pi, V)$ 是拓扑群 $G$ 的表示 (连续或不连续), 记
\[
V^G = \{v \in V \mid \forall x \in G,\, \pi(x)(v) = v\}.
\]
这是 $V$ 的闭 $G$-不变子空间.
\end{notation}

\begin{example}[交错算子作为不动点]
\label{exa:IV.5.6}
若 $V$ 和 $W$ 是 $G$ 的两个表示, 则
\[
\Hom(V, W)^G = \Hom_G(V, W).
\]
\end{example}

\begin{theorem}[特征标的积分]
\label{thm:IV.5.7}
设 $G$ 是紧群, $(\pi, V)$ 是 $G$ 的有限维连续表示. 则
\[
\int_G \chi_V(x)\,\dd x = \dim(V^G).
\]
\end{theorem}

\begin{proof}
由于 $V$ 有限维, 可以找到 $V$ 的不可约子表示的有限族 $(V_i)_{i \in I}$ 使得 $V = \bigoplus_{i \in I}V_i$. (参见推论 \ref{cor:I.3.2.9}.) 由命题 \ref{prop:IV.5.4} 有 $\chi_V = \sum_{i \in I}\chi_{V_i}$, 且 $V^G = \bigoplus_{i \in I}V_i^G$. 所以只需对 $V$ 不可约证明定理.
假设 $V$ 是 $G$ 的不可约表示. 由于 $V^G$ 是 $V$ 的 $G$-不变子空间, 有 $V^G = V$ 或 $V^G = \{0\}$. 若 $V^G = V$, 则 $G$ 在 $V$ 上平凡作用, 所以 $V$ 的每个线性子空间都被 $G$ 不变, 故必有 $\dim V = 1$. 另一方面, 对每个 $x \in G$ 有 $\chi_V(x) = \operatorname{Tr}(1) = 1$, 所以 $\int_G\chi_V(x)\,\dd x = 1$. 假设 $V$ 不可约且 $V^G = \{0\}$. 设 $\pi_0$ 是 $G$ 在 $\CC$ 上的平凡表示. 则由定理 \ref{thm:IV.3.8}(i), $L^2(G)$ 的子空间 $E_{\pi}$ 和 $E_{\pi_0}$ 正交. 但 $E_{\pi_0}$ 是常值函数的子空间, 且我们上面 (注记 \ref{rem:IV.5.2}) 看到 $\chi_V \in E_{\pi}$. 所以 $\chi_V$ 与常值函数 $1$ 正交, 这正好意味着 $\int_G\chi_V(x)\,\dd x = 0$.
\end{proof}

\begin{corollary}[特征标的内积]
\label{cor:IV.5.8}
设 $G$ 是紧群, $(\pi, V)$ 和 $(\sigma, W)$ 是 $G$ 的两个连续有限维表示.
\begin{enumerate}
\item[(i)] 有 $\langle \chi_W, \chi_V \rangle_{L^2(G)} = \dim_{\CC}(\Hom_G(V, W))$.
\item[(ii)] 若 $V$ 和 $W$ 不可约且不等价, 则 $\langle \chi_V, \chi_W \rangle_{L^2(G)} = 0$.
\item[(iii)] 表示 $V$ 不可约当且仅当 $\|\chi\|_{L^2(G)} = 1$.
\end{enumerate}
\end{corollary}

\begin{proof}
(i). 让 $G$ 作用于 $\Hom(V, W)$: $x \cdot T = \sigma(x) \circ T \circ \pi(x)^{-1}$. 我们知道 (参见命题 \ref{prop:IV.3.4}) $\Hom_G(V, W) = \Hom(V, W)^G$. 把定理应用于表示 $\Hom(V, W)$, 并用命题 \ref{prop:IV.5.4} 的 (iv) 和 (v) 计算此表示的特征标, 得到:
\begin{align*}
\dim_{\CC}(\Hom_G(V, W)) &= \dim_{\CC}(\Hom(V, W)^G) \\
&= \int_G \chi_{\Hom(V,W)}(x)\,\dd x \\
&= \int_G \chi_V(x)\overline{\chi_W(x)}\,\dd x \\
&= \langle \chi_W, \chi_V \rangle_{L^2(G)}.
\end{align*}
(ii). 这由 (i) 和 Schur 引理 (定理 \ref{thm:I.3.4.1}) 得出, 或由 $\chi_V \in E_V$, $\chi_W \in E_W$ 且 $E_V$ 和 $E_W$ 在 $L^2(G)$ 中正交 (见定理 \ref{thm:IV.3.8}) 得出.
(iii). 若 $V$ 不可约, 则 Schur 引理蕴含 $\End_G(V)$ 是 $1$ 维的, 所以由 (i) 有 $\|\chi_V\|_{L^2(G)} = 1$. 反之, 假设 $\|\chi_V\|_{L^2(G)} = 1$. 把 $V$ 写成 $V = \bigoplus_{i \in I}V_i$, 其中 $I$ 有限, $V_i$ 是 $V$ 的不可约子表示. 由 (ii), $G$ 的不同构不可约表示的特征标在 $L^2(G)$ 中正交, 所以有
\[
\|\chi_V\|_{L^2(G)}^2 = \sum_{W \in \widehat{G}}n_W\|\chi_W\|_{L^2(G)}^2 = \sum_{W \in \widehat{G}}n_W,
\]
其中对每个 $W \in \widehat{G}$,
\[
n_W = |\{i \in I \mid V_i \simeq W\}|.
\]
由于 $\|\chi_V\|_{L^2(G)} = 1$, 存在唯一的 $W \in \widehat{G}$ 使得 $n_W \neq 0$, 且必有 $n_W = 1$. 由整数 $n_W$ 的定义, 这意味着 $V \simeq W$, 所以 $V$ 不可约.
\end{proof}

\begin{corollary}[特征标的线性无关性]
\label{cor:IV.5.9}
设 $G$ 是紧群. 则族 $(\chi_V)_{V \in \widehat{G}}$ (作为 $L^2(G)$ 或 $\mathcal{C}(G)$ 的元素) 线性无关.
\end{corollary}

\begin{proof}
这由前一推论的 (ii) 得出.
\end{proof}

\begin{corollary}[特征标确定表示]
\label{cor:IV.5.10}
设 $\pi$ 和 $\pi'$ 是紧群 $G$ 的两个连续有限维表示. 则 $\pi$ 和 $\pi'$ 等价当且仅当 $\chi_{\pi} = \chi_{\pi'}$.
\end{corollary}

\begin{proof}
若 $\pi$ 和 $\pi'$ 等价, 已知 $\chi_{\pi} = \chi_{\pi'}$. 反之, 假设 $\chi_{\pi} = \chi_{\pi'}$. 把 $\pi$ 和 $\pi'$ 分解成不可约表示的直和:
\[
\pi \simeq \bigoplus_{\rho \in \widehat{G}}\rho^{n_{\rho}}
\]
且
\[
\pi' \simeq \bigoplus_{\rho \in \widehat{G}}\rho^{m_{\rho}},
\]
其中 $n_{\rho}, m_{\rho} \in \Z_{\geq 0}$, 且除有限多个 $\rho \in \widehat{G}$ 外 $n_{\rho} = m_{\rho} = 0$. 由推论 \ref{cor:IV.5.8}, 有 $\chi_{\pi} = \sum_{\rho \in \widehat{G}}n_{\rho}\chi_{\rho}$ 和 $\chi_{\pi'} = \sum_{\rho \in \widehat{G}}m_{\rho}\chi_{\rho}$ (且这是有限和). 由 $\chi_{\rho}$ 的线性无关性, 等式 $\chi_{\pi} = \chi_{\pi'}$ 蕴含对每个 $\rho \in \widehat{G}$ 有 $n_{\rho} = m_{\rho}$, 这进而蕴含 $\pi$ 和 $\pi'$ 等价.
\end{proof}

\section{\index{傅里叶变换}傅里叶变换}

我们仍假设 $G$ 是紧群.
由命题 \ref{prop:I.4.3.4} 和 \ref{prop:I.4.1.3}, 空间 $L^2(G)$ 对卷积乘积实际上是 Banach 代数. 本节回答``如何在 Peter-Weyl 定理给出的分解中看到代数结构?''这个问题.

\begin{definition}[傅里叶变换 (Fourier transform)]
\label{def:IV.6.1}
设 $f \in L^2(G)$. 对每个 $(\pi, V) \in \widehat{G}$, $f$ 在 $\pi$ 处的傅里叶变换是 $V$ 的自同态
\[
\hat{f}(\pi) = \int_G f(x)\pi(x^{-1})\,\dd x = \int_G f(x)\pi(x)^*\,\dd x.
\]
\end{definition}

这显然是 $V$ 的 $\CC$-线性自同态.

\begin{example}[$G = S^1$ 的情形]
\label{exa:IV.6.2}
假设 $G = S^1$. 则我们在习题 \ref{ex:I.5.4.1} 中见过 $\widehat{G} \simeq \Z$, 其中 $n \in \Z$ 对应于表示 $G \to \CC^{\times}$, $e^{2\pi it} \longmapsto e^{2\pi int}$ ($t \in \R$). 所以若 $f \in L^1(G)$, 其傅里叶变换是函数 $\hat{f} \colon \Z \to \CC$, 把 $n$ 送到
\[
\hat{f}(n) = \int_0^1 f(e^{2\pi it})e^{-2\pi int}\,\dd t.
\]
\end{example}

\begin{theorem}[傅里叶变换的性质]
\label{thm:IV.6.3}
\begin{enumerate}
\item[(i)] 对每个 $\pi \in \widehat{G}$, 映射 $L^2(G) \to \End(\pi)$, $f \longmapsto \hat{f}(\pi)$ 是从 $L^2(G)$ 到 $\End(\pi)$ 的相反代数的 $G \times G$-等变 $*$-同态. (注意 $L^2(G) \subset L^1(G)$, 因为 $G$ 紧. 例 \ref{exa:I.4.2.2}(b) 中定义的 $L^1(G)$ 的对合限制为 $L^2(G)$ 上的对合.)
换句话说, 对 $f, g \in L^2(G)$ 和 $x \in G$:
\begin{align*}
\widehat{f * g}(\pi) &= \hat{g}(\pi) \circ \hat{f}(\pi), \\
\widehat{f^*}(\pi) &= (\hat{f}(\pi))^*, \\
\widehat{L_xf}(\pi) &= \hat{f}(\pi) \circ \pi(x)^{-1}, \quad \widehat{R_xf}(\pi) = \pi(x) \circ \hat{f}(\pi).
\end{align*}
(与定理 \ref{thm:I.4.2.6}(i) 比较.)
\item[(ii)] 设 $f \in L^2(G)$. 则对每个 $\pi \in \widehat{G}$, 函数 $\dim(\pi)\operatorname{Tr}(\hat{f}(\pi) \circ \pi(\cdot)) \in L^2(G)$ 是 $f$ 在 $E_{\pi}$ 上的正交投影, 且级数
\[
\sum_{\pi \in \widehat{G}}\dim(\pi)\operatorname{Tr}(\hat{f}(\pi) \circ \pi(\cdot))
\]
在 $L^2(G)$ 中收敛到 $f$ (Fourier 反演公式 (Fourier inversion formula)).
\item[(iii)] 对每个 $f \in L^2(G)$,
\[
\|f\|_2^2 = \sum_{\pi \in \widehat{G}}\dim(\pi)\operatorname{Tr}(\hat{f}(\pi)^* \circ \hat{f}(\pi))
\]
(Parseval 公式 (Parseval formula)).
\end{enumerate}
\end{theorem}

\begin{example}[$S^1$ 的 Fourier 级数]
\label{exa:IV.6.4}
取 $G = S^1$. 则 (ii) 和 (iii) 说: 对每个 $f \in L^2(S^1)$, 级数 $\sum_{n \in \Z}\hat{f}(n)e^{2\pi int}$ 在 $L^1(S^1)$ 中收敛到 $f$, 且
\[
\int_0^1|f(e^{2\pi it})|^2\,\dd t = \sum_{n \in \Z}|\hat{f}(n)|^2.
\]
\end{example}

\begin{proof}
(i). 有
\begin{align*}
\widehat{f * g}(\pi) &= \int_G (f * g)(x)\pi(x^{-1})\,\dd x \\
&= \int_{G \times G}f(y)g(y^{-1}x)\pi(x^{-1})\,\dd x\,\dd y \\
&= \int_{G \times G}f(y)g(x)\pi(x^{-1}y^{-1})\,\dd x\,\dd y \quad \text{(代换 } x' = y^{-1}x\text{)} \\
&= \hat{g}(\pi) \circ \hat{f}(\pi).
\end{align*}
回忆 $f^*(x) = \overline{f(x^{-1})}$, 因为 $\Delta = 1$. 所以
\begin{align*}
\widehat{f^*}(\pi) &= \int_G \overline{f(x^{-1})}\pi(x)^*\,\dd x \\
&= \int_G \overline{f(x)}\pi(x^{-1})^*\,\dd x \\
&= (\hat{f}(\pi))^*.
\end{align*}
最后,
\begin{align*}
\widehat{L_xf}(\pi) &= \int_G f(x^{-1}y)\pi(y^{-1})\,\dd y \\
&= \int_G f(y)\pi(y^{-1}x^{-1})\,\dd y \\
&= \hat{f}(\pi) \circ \pi(x^{-1})
\end{align*}
且
\begin{align*}
\widehat{R_xf}(\pi) &= \int_G f(yx)\pi(y^{-1})\,\dd y \\
&= \int_G f(y)\pi(xy^{-1})\,\dd y \\
&= \pi(x) \circ \hat{f}(\pi).
\end{align*}
(ii). 只需证明第一个结论 (第二个将由 Peter-Weyl 定理得出). 设 $(\pi, V) \in \widehat{G}$. 像定理 \ref{thm:IV.3.8} 中那样, 固定 $V$ 的标准正交基 $(e_1, \ldots, e_d)$, 记 $\phi_{ij}$ 为函数 $G \to \CC$, $x \longmapsto \langle \pi(x)(e_j), e_i \rangle$. 则我们 (在定理 \ref{thm:IV.3.8} 的 (ii) 中) 见过 $\{\sqrt{d}\phi_{ij}, 1 \leq i, j \leq d\}$ 是 $E_{\pi}$ 对 $L^2$ 内积的标准正交基. 所以 $f$ 在 $E_{\pi}$ 上的正交投影是
\[
d\sum_{i,j=1}^d\langle f, \phi_{ij} \rangle_{L^2(G)}\phi_{ij}.
\]
对所有 $i, j \in \{1, \ldots, d\}$,
\begin{align*}
\langle f, \phi_{i,j} \rangle_{L^2(G)} &= \int_G f(x)\overline{\langle e_i, \pi(x)(e_j) \rangle}\,\dd x \\
&= \int_G f(x)\langle \pi(x)^*(e_i), e_j \rangle\,\dd x \\
&= \langle \hat{f}(\pi)(e_i), e_j \rangle.
\end{align*}
设 $y \in G$, 记 $(\hat{f}(\pi)_{i,j})$ 和 $(\pi(y)_{i,j})$ 为 $\hat{f}(\pi)$ 和 $\pi(y)$ 在基 $(e_1, \ldots, e_d)$ 中的矩阵. 则
\[
\hat{f}(\pi)_{i,j} = \langle \hat{f}(\pi)(e_j), e_i \rangle = \langle f, \phi_{j,i} \rangle_{L^2(G)}
\]
且
\[
\pi(y)_{i,j} = \langle \pi(y)(e_j), e_i \rangle = \phi_{ij}(y),
\]
所以
\[
\operatorname{Tr}(\hat{f}(\pi) \circ \pi(y)) = \sum_{i,j=1}^d\hat{f}(\pi)_{j,i}\pi(y)_{i,j} = \sum_{i,j=1}^d\langle \hat{f}(\pi), \phi_{i,j} \rangle_{L^2(G)}\phi_{i,j}(y).
\]
这给出 $f$ 在 $E_{\pi}$ 上的正交投影的所要公式.
(iii). 设 $\pi \in \widehat{G}$, 用 (ii) 证明中的记号. 设 $g = d\operatorname{Tr}(\hat{f}(\pi)^* \circ \hat{f}(\pi))$. 只需证明 $\|g\|_2^2 = d\operatorname{Tr}(\hat{f}(\pi)^* \circ \hat{f}(\pi))$ (因为不同构 $\pi$ 的 $E_{\pi}$ 正交, 由定理 \ref{thm:IV.3.8}). 有
\[
\operatorname{Tr}(\hat{f}(\pi)^* \circ \hat{f}(\pi)) = \sum_{i,j=1}^d|\hat{f}(\pi)_{i,j}|^2 = \sum_{i,j=1}^d|\langle f, \phi_{i,j} \rangle_{L^2(G)}|^2.
\]
另一方面, 由于 $g = d\sum_{i,j=1}^d\langle f, \phi_{ij} \rangle_{L^2(G)}\phi_{ij}$, 得到
\[
\|g\|_{L^2(G)}^2 = d^2\sum_{i,j=1}^d|\langle f, \phi_{i,j} \rangle_{L^2(G)}|^2 = d \cdot d\operatorname{Tr}(\hat{f}(\pi)^* \circ \hat{f}(\pi)).
\]
\end{proof}

\section{特征标与傅里叶变换}
为完成本章, 我们把特征标与傅里叶变换联系起来, 并解释``中心函数''这个名字.
\begin{proposition}[特征标与卷积]
\label{prop:IV.7.1}
设 $f \in L^2(G)$. 则对每个 $x \in G$,
\[
\operatorname{Tr}(\hat{f}(\pi) \circ \pi(x)) = f * \chi_{\pi}(x).
\]
\end{proposition}

由定理 \ref{thm:IV.6.3}, 这说 $f$ 在 $E_{\pi}$ 上的正交投影是 $\dim(\pi)f * \chi_{\pi}$, 所以在 $L^2(G)$ 中有
\[
f = \sum_{\pi \in \widehat{G}}\dim(\pi)f * \chi_{\pi}.
\]

\begin{proof}
有
\begin{align*}
\operatorname{Tr}(\hat{f}(\pi) \circ \pi(x)) &= \int_G f(y)\operatorname{Tr}(\pi(y)^{-1}\pi(x))\,\dd y \\
&= \int_G f(y)\chi_{\pi}(y^{-1}x)\,\dd y \\
&= f * \chi_{\pi}(x).
\end{align*}
\end{proof}

\begin{corollary}[特征标的卷积]
\label{cor:IV.7.2}
对所有 $\pi, \pi' \in \widehat{G}$,
\[
\chi_{\pi} * \chi_{\pi'} = \begin{cases} \dim(\pi)^{-1}\chi_{\pi} & \text{若 } \pi \simeq \pi', \\ 0 & \text{否则}. \end{cases}
\]
\end{corollary}

\begin{proof}
我们知道对每个 $\pi \in \widehat{G}$ 有 $\chi_{\pi} \in E_{\pi}$, $\pi \not\simeq \pi'$ 时 $E_{\pi}$ 和 $E_{\pi'}$ 正交, 且命题说 $\dim(\pi)\chi_{\pi} * \chi_{\pi'}$ 是 $\chi_{\pi}$ 在 $E_{\pi'}$ 上的正交投影. 这立即蕴含推论的公式.
\end{proof}

\begin{definition}[中心函数空间]
\label{def:IV.7.3}
对 $1 \leq p < +\infty$, 记 $ZL^p(G)$ 为 $L^p(G)$ 中中心函数的子空间. 又记 $Z\mathcal{C}(G)$ 为 $\mathcal{C}(G)$ 中中心函数的子空间.
\end{definition}

\begin{proposition}[$L^p$ 的中心]
\label{prop:IV.7.4}
空间 $L^p(G)$ ($1 \leq p < +\infty$) (相应地 $\mathcal{C}(G)$) 对卷积乘积是 Banach 代数, 且 $ZL^p(G)$ (相应地 $Z\mathcal{C}(G)$) 是它的中心.
\end{proposition}

\begin{proof}
设 $p \in [1, +\infty)$, $q \in [1, +\infty)$ 满足 $p^{-1} + q^{-1} = 1$. 由于 $G$ 紧, 常值函数 $1$ 在 $L^q(G)$ 中且 $L^q$ 范数为 $1$, 所以由 H\"older 不等式, $f = f \cdot 1$ 在 $L^1(G)$ 中, 且 $\|f\|_1 \leq \|f\|_p$. 现在推论 \ref{cor:I.4.3.2} 说: 对每个 $g \in L^p(G)$, 函数 $f * g$ 存在且在 $L^p(G)$ 中, 且有 $\|f * g\|_p \leq \|f\|_1\|g\|_p \leq \|f\|_p\|g\|_p$. 这证明 $L^p(G)$ 对 $*$ 是 Banach 代数.
我们证明 $\mathcal{C}(G)$ 对 $*$ 也是 Banach 代数. 若 $f, g \in \mathcal{C}(G)$, 则 $f * g$ 显然存在, 且对每个 $x \in G$,
\[
|f * g(x)| \leq \int_G |f(y)||g(y^{-1}x)|\,\dd y \leq \|f\|_{\infty}\|g\|_{\infty}\int_G 1\,\dd y = \|f\|_{\infty}\|g\|_{\infty}.
\]
所以 $\|f * g\|_{\infty} \leq \|f\|_{\infty}\|g\|_{\infty}$.
最后, 我们证明关于中心的结论. 设 $f \in L^p(G)$, 假设对每个 $g \in L^p(G)$ 有 $f * g = g * f$. 则对每个 $x \in G$ 和每个 $g \in L^p(G)$,
\[
\int_G f(xy)g(y^{-1})\,\dd y = \int_G g(y)f(y^{-1}x)\,\dd y = \int_G f(yx)g(y^{-1})\,\dd y.
\]
这当且仅当 $f(xy) = f(yx)$ 在 $G \times G$ 上几乎处处成立. 对 $f \in \mathcal{C}(G)$ 的证明相同.
\end{proof}

\begin{corollary}[$ZL^2(G)$ 的基]
\label{cor:IV.7.5}
族 $(\chi_{\pi})_{\pi \in \widehat{G}}$ 是 $ZL^2(G)$ 的标准正交基.
\end{corollary}

\begin{proof}
我们已经知道 $\chi_{\pi}$ 在 $ZL^2(G)$ 中且两两正交, 所以只需证明正交于所有 $\chi_{\pi}$ 的中心函数必为 $0$. 设 $f \in ZL^2(G)$. 由下面的引理, 对每个 $\pi \in \widehat{G}$ 有 $(\dim\pi)f * \chi_{\pi} = \langle f, \chi_{\pi} \rangle_{L^2(G)}\chi_{\pi}$, 所以若 $f$ 与每个 $\chi_{\pi}$ 正交, 则由命题 \ref{prop:IV.7.1}, 它在所有空间 $E_{\pi}$ 上的投影为 $0$, 从而由定理 \ref{thm:IV.4.2} 有 $f = 0$.
\end{proof}

\begin{lemma}[中心函数的傅里叶变换]
\label{lem:IV.7.6}
若 $f \in ZL^1(G)$, $\pi \in \widehat{G}$, 则 $(\dim\pi)f * \chi_{\pi} = \langle f, \chi_{\pi} \rangle_{L^2(G)}\chi_{\pi}$.
\end{lemma}

\begin{proof}
我们知道由命题 \ref{prop:IV.7.1} 有 $f * \chi_{\pi} = \operatorname{Tr}(\hat{f}(\pi) \circ \pi(\cdot))$. 对每个 $x \in G$,
\begin{align*}
\hat{f}(\pi) \circ \pi(x) &= \int_G f(y)\pi(y^{-1}x)\,\dd y \\
&= \int_G f(xy^{-1})\pi(y)\,\dd y \\
&= \int_G f(y^{-1}x)\pi(y)\,\dd y \quad \text{(因为 } f \text{ 中心)} \\
&= \int_G f(y)\pi(xy^{-1})\,\dd y \\
&= \pi(x) \circ \hat{f}(\pi).
\end{align*}
所以 $\hat{f}(\pi) \in \End(\pi)$ 是 $G$-等变的. 由 Schur 引理, 这蕴含 $\hat{f}(\pi) = c\,\mathrm{id}$, 其中 $c \in \CC$. 取迹给出
\[
c(\dim\pi) = \operatorname{Tr}(\hat{f}(\pi)) = \int_G f(y)\operatorname{Tr}(\pi(y^{-1}))\,\dd y = \langle f, \chi_{\pi} \rangle_{L^2(G)}.
\]
所以
\[
\langle f, \chi_{\pi} \rangle_{L^2(G)}\chi_{\pi} = (\dim\pi)\operatorname{Tr}(\hat{f}(\pi) \circ \pi(\cdot)) = (\dim\pi)f * \chi_{\pi}.
\]
\end{proof}

\begin{remark}[$ZL^p(G)$ 中的稠密性]
\label{rem:IV.7.7}
事实上, 我们甚至可以证明族 $(\chi_{\pi})_{\pi \in \widehat{G}}$ 在 $ZL^p(G)$ (对每个 $p \in [1, +\infty)$) 和 $ZC(G)$ 中张成稠密子空间. (见 {[}11{]} 的命题 5.25.)
\end{remark}

\begin{remark}[有限群的情形]
\label{rem:IV.7.8}
若 $G$ 有限, 则 $L^2(G)$ 是所有从 $G$ 到 $\CC$ 的函数的空间, $ZL^2(G)$ 是在 $G$ 的共轭类上取常值的函数的空间. 所以上面的命题说 $|\widehat{G}|$ 等于 $G$ 的共轭类数, 且 Peter-Weyl 定理说 $|G| = \sum_{\pi \in \widehat{G}}(\dim\pi)^2$.
\end{remark}

\begin{remark}[$ZL^p(G)$ 的谱]
\label{rem:IV.7.9}
我们特别证明了 Banach 代数 $(ZL^p(G), *)$ ($1 \leq p < +\infty$) 和 $(Z\mathcal{C}(G), *)$ 交换. 我们可以问它们的谱是什么. 事实上, 答案非常简单 (见 {[}11{]} 的定理 5.26): 对每个 $\pi \in \widehat{G}$, 公式 $f \longmapsto (\dim\pi)\int_G f\chi_{\pi}\,\dd\mu$ 定义 $ZL^p(G)$ (相应地 $Z\mathcal{C}(G)$) 上的乘法泛函, 且这诱导从离散集 $\widehat{G}$ 到 $ZL^p(G)$ (相应地 $Z\mathcal{C}(G)$) 的谱的同胚.
\end{remark}

\section{Peter-Weyl 定理的经典证明}
在第 IV.4 节中, 我们给出了使用 Gelfand-Raikov 和 Stone-Weierstrass 定理的 Peter-Weyl 定理的证明. 我们现在解释原始证明. 我们固定紧群 $G$. 回忆这蕴含 $L^2(G) \subset L^1(G)$.
由推论 \ref{cor:I.4.3.2}, 若 $f \in L^1(G)$, $g \in L^2(G)$, 则定义 $f * g$ 和 $g * f$ 的积分收敛, 并定义 $L^2(G)$ 的函数, 使得 $\|f * g\|_2 \leq \|f\|_1\|g\|_2$, $\|g * f\|_2 \leq \|f\|_1\|g\|_2$.

\begin{definition}[$L^2(G)$ 上的卷积算子]
\label{def:IV.8.1}
若 $f \in L^1(G)$, 定义 $L^2(G)$ 的连续线性自同态 $L_f$ 和 $R_f$ 为 $L_f(g) = f * g$, $R_f(g) = g * f$.
\end{definition}

事实上, 由习题 \ref{ex:I.5.6.4} (及其对右乘法的显然类似), $L^1(G)$ 在 $L^2(G)$ 上的这些作用正是定理 \ref{thm:I.4.2.6} 中定义的 $G$ 在 $L^2(G)$ 上的左、右正则表示向群代数 $L^1(G)$ 的延拓. 特别地, $\|L_f\|_{\mathrm{op}} \leq \|f\|_1$, $\|R_f\|_{\mathrm{op}} \leq \|f\|_1$, 这也是推论 \ref{cor:I.4.3.2} 已知的.
在例 \ref{exa:I.4.2.2}(b) 中, 我们在 $L^1(G)$ 上定义了对合 $*$; 由于 $G$ 的模函数为 $1$, 此对合把 $f \in L^1(G)$ 送到由 $f^*(x) = \overline{f(x^{-1})}$ 定义的函数 $f^*$. 由定理 \ref{thm:I.4.2.6}, 对每个 $f \in L^1(G)$ 有 $(L_f)^* = L_{f^*}$, $(R_f)^* = R_{f^*}$.

\begin{theorem}[卷积算子紧]
\label{thm:IV.8.2}
对每个 $f \in L^1(G)$, $L^2(G)$ 的自同态 $R_f$ 和 $L_f$ 紧.
\end{theorem}

\begin{remark}[迹类与 Hilbert-Schmidt]
\label{rem:IV.8.3}
我们还可以证明 $L_f$ 和 $R_f$ 是迹类算子 (trace class operator) (对 $f \in L^1(G)$ 迹为 $\int_G f(x)\,\dd x$), 且对 $f \in L^2(G)$ 是 Hilbert-Schmidt 算子 (Hilbert-Schmidt norm 为 $\|f\|_2$), 但我们不需要.
\end{remark}

\begin{lemma}[核算子]
\label{lem:IV.8.4}
对每个 $F \in L^2(G \times G)$, 公式
\[
T_F(h)(x) = \int_G F(x, y)h(y)\,\dd y
\]
定义连续线性算子 $T_F \colon L^2(G) \to L^2(G)$, 且有 $\|T_F\| \leq \|F\|_2$.\footnote{事实上, 可以证明 $T_F$ 是 Hilbert-Schmidt 算子且其 Hilbert-Schmidt 范数等于 $\|F\|_2$.}
\end{lemma}

\begin{proof}
设 $h \in L^2(G)$. 由 Minkowski 不等式和 Cauchy-Schwarz 不等式,
\begin{align*}
\Bigl(\int_G\Bigl|\int_G F(x, y)h(y)\,\dd y\Bigr|^2\,\dd x\Bigr)^{1/2} &\leq \int_G\Bigl(\int_G|F(x, y)|^2|h(y)|^2\,\dd x\Bigr)^{1/2}\,\dd y \\
&\leq \int_G |h(y)|\Bigl(\int_G|F(x, y)|^2\,\dd x\Bigr)^{1/2}\,\dd y \\
&\leq \|h\|_2\Bigl(\int_G\Bigl(\int_G|F(x, y)|^2\,\dd x\Bigr)^{1/2 \times 2}\,\dd y\Bigr)^{1/2} \\
&= \|h\|_2\|F\|_2.
\end{align*}
这证明 $T_F(h)$ 良定义, 在 $L^2(G)$ 中, 且 $\|T_F(h)\|_2 \leq \|h\|_2\|F\|_2$.
\end{proof}

\begin{lemma}[$L^2(G) \otimes L^2(G) \simeq L^2(G \times G)$]
\label{lem:IV.8.5}
若 $f_1, f_2 \colon G \to \CC$, 定义函数 $u(f_1 \otimes f_2) \colon G \times G \to \CC$ 为 $(f_1 \otimes f_2)(x_1, x_2) = f_1(x_1)f_2(x_2)$. 这诱导 $\CC$-线性等距
\[
u \colon L^2(G) \otimes_{\CC}L^2(G) \to L^2(G \times G),
\]
它单射且像稠密, 从而给出等距同构 $L^2(G)\hat{\otimes}_{\CC}L^2(G) \xrightarrow{\sim} L^2(G \times G)$.
\end{lemma}

证明此引理后, 我们直接把 $L^2(G)\hat{\otimes}_{\CC}L^2(G)$ 与 $L^2(G \times G)$ 等同, 并写 $f_1 \otimes f_2$ 代替 $u(f_1 \otimes f_2)$.

\begin{proof}
$\CC$-线性映射 $u$ 的存在性由张量积的性质得出. 若 $f_1, f_2, g_1, g_2 \in L^2(G)$, 则
\begin{align*}
\langle u(f_1 \otimes f_2), u(g_1 \otimes g_2) \rangle_{L^2(G \times G)} &= \int_{G \times G}u(f_1 \otimes f_2)(x_1, x_2)\overline{u(g_1 \otimes g_2)(x_1, x_2)}\,\dd x_1\,\dd x_2 \\
&= \int_{G \times G}f_1(x_1)f_2(x_2)\overline{g_1(x_1)}\overline{g_2(x_2)}\,\dd x_1\,\dd x_2 \\
&= \langle f_1, g_1 \rangle_{L^2(G)}\langle f_2, g_2 \rangle_{L^2(G)}.
\end{align*}
这蕴含 $u$ 等距.
取 $L^2(G)$ 的 Hilbert 基 $(e_i)_{i \in I}$. 则由上面的计算, $L^2(G \times G)$ 的族 $(u(e_i \otimes e_j))_{i,j \in I}$ 标准正交, 且其张成是 $u$ 的像. 所以我们只需证明正交于每个 $u(e_i \otimes e_j)$ 的 $L^2(G \times G)$ 中的函数为 $0$. 设 $F \in L^2(G \times G)$ 是这种函数. 设 $i \in I$, 考虑函数 $f_i \colon G \to \CC$, $f_i(x) = \int_G F(x, y)\overline{e_i(y)}\,\dd y$. 则 $f_i \in L^2(G)$ (由引理 \ref{lem:IV.8.4}), 且 $f_i$ 与所有 $e_j$ 正交 (由 $F$ 的选取). 由于 $(e_j)_{j \in I}$ 是 $L^2(G)$ 的 Hilbert 基, 必有 $f_i = 0$. 这对每个 $i \in I$ 成立, 所以再次用 $(e_i)_{i \in I}$ 是 $L^2(G)$ 的 Hilbert 基这一事实, 看到函数 $F(x, \cdot)$ 对几乎每个 $x$ 几乎处处为零, 即 $F = 0$ ($L^2(G \times G)$ 中).
\end{proof}

\begin{lemma}[核算子紧]
\label{lem:IV.8.6}
对每个 $F \in L^2(G \times G)$, 引理 \ref{lem:IV.8.4} 的算子 $T_F \colon L^2(G) \to L^2(G)$ 紧.
\end{lemma}

\begin{proof}
由引理 \ref{lem:IV.8.5}, 存在 $L^2(G)$ 的函数族 $(g_i)_{i \in I}$ 和 $(h_i)_{i \in I}$ 使得和 $\sum_{i \in I}g_i \otimes h_i$ 在 $L^2(G \times G)$ 中收敛到 $F$.
对 $I$ 的每个有限子集 $J$, 设 $S_J = \sum_{i,j \in J}g_i \otimes h_j \in L^2(G \times G)$, $T_J = T_{S_J}$. 则对每个 $h \in L^2(G)$, 每个 $x \in G$,
\[
(T_Jh)(x) = \sum_{(i,j) \in J^2}\int_G h(y)g_i(x)\overline{h_j(y)}\,\dd y = \sum_{i \in J}\Bigl(\sum_{j \in J}\int_G h(y)\overline{h_j(y)}\,\dd y\Bigr)g_i(x).
\]
换句话说, 对每个 $h \in L^2(G)$, $T_Jh$ 在由 $g_i$ ($i \in J$) 张成的 $L^2(G)$ 的有限维子空间中. 所以算子 $T_J$ 有有限秩.
为证明 $T_F$ 紧, 由习题 \ref{ex:I.5.5.9}(f) 只需证明它是 $J$ 变大时算子 $T_J$ 的极限. 但这由引理 \ref{lem:IV.8.4} 和 $S_J$ 在 $L^2(G \times G)$ 中收敛到 $F$ 的事实得出.
\end{proof}

\begin{proof}[定理 \ref{thm:IV.8.2} 的证明]
我们证明 $L_f$ 的结果; $R_f$ 的证明类似.
首先假设 $f \in L^2(G)$. 考虑函数 $K \colon G \times G \to \CC$, $(x, y) \longmapsto f(xy^{-1})$. 则 $K \in L^2(G \times G)$, 且 $L_f = T_K$. 所以结果由引理 \ref{lem:IV.8.6} 得出.
一般情形, 结果由紧算子的极限紧 (见习题 \ref{ex:I.5.5.9}(f))、$L^2(G)$ 在 $L^1(G)$ 中稠密以及 $f \longmapsto L_f$ (相应地 $f \longmapsto R_f$) 是从 $L^1(G)$ 到 $\End(L^2(G))$ 的连续映射这一事实得出.
\end{proof}

回忆 $\mathcal{C}(G)$ 的子空间 $E$ 在第 IV.4.1 节中定义; 这是 $G$ 的有限维表示的矩阵系数的空间.

\begin{proposition}[有限维子表示]
\label{prop:IV.8.7}
设 $f \in L^2(G)$. 则下列条件等价:
\begin{enumerate}
\item[(i)] $\operatorname{Span}\{L_xf, x \in G\}$ 有限维;
\item[(ii)] $\operatorname{Span}\{R_xf, x \in G\}$ 有限维;
\item[(iii)] $f \in E$.
\end{enumerate}
\end{proposition}

\begin{proof}
我们首先证明 (i) 蕴含 (ii) 和 (iii). 若 $f \in E$, 则 $f$ 是 $\sum_{\pi \in \widehat{G}}f_{\pi}$ 的有限和, 其中 $f_{\pi} \in E_{\pi}$. 所以只需证明对 $\pi \in \widehat{G}$, $E_{\pi}$ 的每个元素在左 (相应地右) 正则表示中生成有限维子空间; 但这由 $E_{\pi}$ 有限维且被算子 $L_x$ (相应地 $R_x$) ($x \in G$) 稳定这一事实立即得出.
现在我们证明 (ii) 蕴含 (i). (证明 (iii) 蕴含 (i) 类似.) 所以设 $f \in L^2(G)$, 假设存在包含 $f$ 且被所有 $L_x$ ($x \in G$) 稳定的 $L^2(G)$ 的有限维子空间 $V$. 我们将证明 $f$ 是 $V^*$ 的矩阵系数. 设 $\langle \cdot, \cdot \rangle_V$ 是 $L^2(G)$ 的内积到 $V$ 的限制, $\rho \colon G \to \GL(V)$ 是 $G$ 在 $V$ 上的作用, $\Pi \colon L^2(G) \to V$ 是正交投影. 假设 $\psi \in \mathcal{C}(G)$ 取实值且对每个 $x \in G$ 有 $\psi(x^{-1}) = \psi(x)$. 特别地, $\psi \in L^2(G)$, 所以由命题 \ref{prop:I.4.3.4}, $f * \psi$ 连续. 对每个 $x \in G$,
\begin{align*}
f * \psi(x) &= L_{x^{-1}}(f * \psi)(1) \\
&= (L_{x^{-1}}f) * \psi(1) \quad \text{(由命题 \ref{prop:I.4.1.3})} \\
&= \int_G (L_{x^{-1}}f)(y)\overline{\psi(y)}\,\dd y \quad \text{(因为 } \psi \text{ 取实值且 } \psi(y^{-1}) = \psi(y)\text{)} \\
&= \langle L_{x^{-1}}f, \psi \rangle_{L^2(G)} \\
&= \langle \rho(x^{-1})\Pi(f), \Pi(\psi) \rangle_V \quad \text{(因为 } L_{x^{-1}}f \in V\text{)} \\
&= \langle \rho(x)\Pi(\psi), \Pi(f) \rangle_V.
\end{align*}
由命题 \ref{prop:IV.3.2}(iv), $f * \psi$ 是 $V^*$ 的矩阵系数, 即 $f * \psi \in E_{V^*}$. 现在我们选取近似单位元 $(\psi_U)_{U \in \mathcal{U}}$ (定义 \ref{def:I.4.1.7}, 由命题 \ref{prop:I.4.1.8} 存在). 由我们刚做的计算, 对每个 $U \in \mathcal{U}$ 有 $f * \psi_U \in E_{V^*}$. 但由推论 \ref{cor:I.4.3.3}, $f$ 是族 $(f * \psi_U)_{U \in \mathcal{U}}$ 在 $L^2(G)$ 中的极限. 由于 $E_{V^*}$ 有限维 (由命题 \ref{prop:IV.3.2}(ii)), 它是 $L^2(G)$ 的闭子空间, 所以 $f$ 也在 $E_{V^*}$ 中.
\end{proof}

\begin{proof}
现在我们解释如何不用 Gelfand-Raikov 定理证明定理 \ref{thm:IV.4.2} (即 $E$ 在 $\mathcal{C}(G)$ 中稠密这一事实). 只需证明 $E$ 在 $L^2(G)$ 中稠密. 考虑 $G$ 在 $L^2(G)$ 上的左正则表示. 由命题 \ref{prop:IV.8.7}, $E$ 包含 $L^2(G)$ 的每个有限维子表示, 事实上它是 $L^2(G)$ 的所有有限维子表示的和. 由于 $L^2(G)$ 是希尔伯特空间, 为证明 $E$ 在 $L^2(G)$ 中稠密, 只需证明 $E^{\perp} = \{0\}$. 所以设 $f \in E^{\perp}$. 选取近似单位元 $(\psi_U)_{U \in \mathcal{U}}$ (定义 \ref{def:I.4.1.7}, 由命题 \ref{prop:I.4.1.8} 存在). 若 $U \in \mathcal{U}$, 则 $\psi_U^* = \psi_U$, 所以 $R_{\psi_U}$ 自伴, 且由定理 \ref{thm:IV.8.2} 它紧. 所以由自伴紧算子的谱定理 (定理 \ref{thm:IV.1.3}), $\Ker(R_{\psi_U})$ 是 $\bigoplus_{\lambda \in \CC^{\times}}\Ker(R_{\psi_U} - \lambda\mathrm{id})$ 的和的闭包的正交补, 且对每个 $\lambda \in \CC^{\times}$, $\Ker(R_{\psi_U} - \lambda\mathrm{id})$ 有限维. 又, $R_{\psi_U}$ 与 $G$ 的左作用交换 (由命题 \ref{prop:I.4.1.3}), 所以每个空间 $\Ker(R_{\psi_U} - \lambda\mathrm{id})$ 都是 $L^2(G)$ 的子表示. 由于由假设 $f$ 与 $L^2(G)$ 的每个有限维子表示正交, 这蕴含 $f \in \Ker(R_{\psi_U})$, 即 $f * \psi_U = 0$. 但 $f$ 是 $(f * \psi_U)_{U \in \mathcal{U}}$ 在 $L^2(G)$ 中的极限 (由推论 \ref{cor:I.4.3.3}), 所以我们得出结论 $f = 0$.
\end{proof}

\section{习题}

\begin{exercise}[$\SU(d)$ 的表示]
\label{ex:IV.9.1}
设 $d$ 是正整数.
\begin{enumerate}
\item[(a)] 设 $T$ 是 $\UU(d)$ 与对角矩阵集的交. 证明
\[
T = \left\{\begin{pmatrix} z_1 & \cdots & 0 \\ \vdots & \ddots & \vdots \\ 0 & \cdots & z_d \end{pmatrix},\, z_1, \ldots, z_d \in S^1\right\}.
\]
\item[(b)] 证明 $\UU(d)$ 的每个元素都在 $\UU(d)$ 中与 $T$ 的元素共轭.
\item[(c)] 证明 $\SU(d)$ 的每个元素都在 $\SU(d)$ 中与 $T_0 := T \cap \SU(d)$ 的元素共轭.
\item[(d)] 证明 $\SU(d)$ 的有限维表示 $V$ 由 $\chi_V|_{T_0}$ 在等价意义下唯一确定.
\end{enumerate}
我们现在取 $d = 2$. 回忆习题 \ref{ex:I.5.5.1} 中定义的 $\SU(2)$ 的不可约表示 $V_n$ ($n \geq 0$).
\begin{enumerate}
\item[(a)] 计算 $\chi_{V_n}$ 到 $T_0$ 的限制.
\item[(b)] 设 $(\rho, V)$ 是 $\SU(2)$ 的有限维表示. 证明存在 $m \geq 1$ 和非负整数 $a_0, \ldots, a_m$ 使得对每个 $z \in S^1$,
\[
\chi_V\left(\begin{pmatrix} z & 0 \\ 0 & z^{-1} \end{pmatrix}\right) = a_0 + \sum_{i=1}^m a_i(z^i + z^{-i}).
\]
\item[(c)] 证明存在整数 $c_n \in \Z$ ($n \geq 0$) 使得对足够大的 $n$ 有 $c_n = 0$, 且
\[
\chi_V = \sum_{n \geq 0}c_n\chi_{V_n}.
\]
\item[(d)] 证明 (c) 的整数 $c_n$ 都非负.
\item[(e)] 若 $V$ 不可约, 证明存在 $n \geq 0$ 使得 $V \simeq V_n$.
\end{enumerate}
\end{exercise}

\begin{solution}
(a). 设 $x = \begin{pmatrix} z_1 & \cdots & 0 \\ \vdots & \ddots & \vdots \\ 0 & \cdots & z_d \end{pmatrix} \in M_d(\CC)$, 其中 $z_1, \ldots, z_d \in \CC$. 则 $x \in T$ 当且仅当 $xx^* = I_d$. 由于 $x^* = \begin{pmatrix} \bar{z}_1 & \cdots & 0 \\ \vdots & \ddots & \vdots \\ 0 & \cdots & \bar{z}_d \end{pmatrix}$, 此条件等价于 $|z_1| = \cdots = |z_d| = 1$.
(b). 设 $x \in \UU(d)$. 则 $x$ 正规, 所以由谱定理, 它可以在 $\CC^n$ 的标准正交基中对角化. 这意味着存在 $y \in \UU(d)$ 使得 $yxy^{-1}$ 对角, 即 $yxy^{-1} \in T$.
(c). 设 $x \in \SU(d)$. 由 (b), 存在 $y \in \UU(d)$ 使得 $yxy^{-1} \in T$. 有 $\det(yxy^{-1}) = \det(x) = 1$, 所以 $yxy^{-1}$ 实际上在 $T_0$ 中. 设 $c = \det(y) \in \CC^{\times}$. 选取 $c' \in \CC$ 使得 $(c')^d = c$; 由于 $|c| = 1$, 也有 $|c'| = 1$. 则 $y' := (c')^{-1}y$ 的行列式为 $1$, 所以在 $\SU(d)$ 中, 且 $y'x(y')^{-1} = yxy^{-1}$.
(d). 设 $V, W$ 是 $\SU(d)$ 的两个有限维表示, 假设 $\chi_V = \chi_W$ 在 $T_0$ 上. 由 (c) 和 $\chi_V$, $\chi_W$ 是中心函数这一事实, 这蕴含 $\chi_V = \chi_W$ 在整个 $\SU(d)$ 上. 但此时 $V$ 和 $W$ 由推论 \ref{cor:IV.5.10} 等价.
(e). 设 $x = \begin{pmatrix} x_1 & 0 \\ 0 & x_2 \end{pmatrix} \in T_0$. 注意 $x_1x_2 = 1$, 所以 $x_2 = x_1^{-1} = \bar{x}_1$. 我们计算 $x$ 在 $V_n$ 的基 $(t_1^kt_2^{n-k})_{0 \leq k \leq n}$ 上的作用. 对 $0 \leq k \leq n$,
\[
x \cdot t_1^kt_2^{n-k} = (x_1^{-1}t_1)^k(x_2^{-1}t_2)^{n-k} = x_1^{n-2k}t_1^kt_2^{n-k}.
\]
所以
\[
\chi_{V_n}(x) = \sum_{k=0}^n x_1^{n-2k}.
\]
(f). 我们把 $S^1$ 通过连续群态射 $z \longmapsto \begin{pmatrix} z & 0 \\ 0 & z^{-1} \end{pmatrix}$ 嵌入 $\SU(2)$. 注意这诱导拓扑群的同构 $S^1 \xrightarrow{\sim} T_0$. 则 $\rho|_{S^1}$ 是 $S^1$ 的有限维表示, 所以它是不可约表示的有限直和. 我们知道 (从习题 \ref{ex:I.5.4.1}) $S^1$ 的每个不可约表示都具有 $\rho_m \colon z \longmapsto z^m$ ($m \in \Z$) 的形式, 所以存在非负整数 $a_m$ ($m \in \Z$) 使得除有限多个 $m$ 外 $a_m = 0$, 且 $\rho|_{S^1} \simeq \bigoplus_{m \in \Z}\rho_m^{a_m}$. 特别地, 对每个 $z \in S^1$,
\[
\chi_V\left(\begin{pmatrix} z & 0 \\ 0 & z^{-1} \end{pmatrix}\right) = \sum_{m \in \Z}a_mz^m.
\]
设 $y = \begin{pmatrix} 0 & i \\ i & 0 \end{pmatrix}$. 则 $y \in \SU(2)$, 且对每个 $z \in S^1$,
\[
y\begin{pmatrix} z & 0 \\ 0 & z^{-1} \end{pmatrix}y^{-1} = \begin{pmatrix} z^{-1} & 0 \\ 0 & z \end{pmatrix}.
\]
由于 $V$ 是 $\SU(2)$ 的表示, 函数 $\chi_V$ 在 $\SU(2)$ 上中心, 所以
\[
\chi_V\begin{pmatrix} z & 0 \\ 0 & z^{-1} \end{pmatrix} = \chi_V\begin{pmatrix} z^{-1} & 0 \\ 0 & z \end{pmatrix}.
\]
这蕴含对每个 $m \in \Z$ 有 $a_{-m} = a_m$, 所以我们得到所要的结论.
(g). 设 $M$ 是可以写成 $\chi(z) = a_0 + \sum_{m \geq 1}a_m(z^m + z^{-m})$ ($a_0, a_1, \ldots \in \Z$ 且对足够大的 $m$ 有 $a_m = 0$) 的函数 $\chi \colon S^1 \to \Z$ 的 $\Z$-模. 由 (f), $\chi_V$ 到 $S^1$ 的限制在 $M$ 中.
$M$ 在 $\Z$ 上的基由函数 $(\chi_0 = 1, \chi_1 = z + z^{-1}, \chi_2 = z^2 + z^{-2}, \ldots)$ 组成. 另一方面, 我们在 (e) 中见过 $\chi_{V_n} = \sum_{0 \leq k \leq n}\chi_k$. 所以表示 $(\chi_{V_n})_{n \geq 0}$ 在基 $(\chi_n)_{n \geq 0}$ 中的 (无穷) 矩阵是上三角的, 对角线上为 $1$, 这意味着它可逆, 即 $(\chi_{V_n})_{n \geq 0}$ 也是 $M$ 在 $\Z$ 上的基. (如果你不喜欢这样, 也很容易从把 $\chi_{V_n}$ 在基 $(\chi_m)_{m \geq 0}$ 中表示的公式对 $n$ 归纳证明 $(\chi_{V_0}, \ldots, \chi_{V_n})$ 线性无关且张成与 $(\chi_0, \ldots, \chi_n)$ 相同的 $\Z$-子模.)
题目的结论由此立即得出.
(h). 我们知道函数 $\chi_{V_n}$ 在 $L^2(\SU(2))$ 中两两正交 (由推论 \ref{cor:IV.5.8}). 所以对每个 $n \geq 0$,
\[
c_n = \langle \chi_V, \chi_{V_n} \rangle_{L^2(\SU(2))}.
\]
由同一个推论, 右端也等于 $\dim_{\CC}(\Hom_{\SU(2)}(V, V_n))$, 它是非负整数.
(i). 若 $V$ 不可约, 则由 (h) 的证明中的最后一个公式 (和 Schur 引理), 除非 $V \simeq V_n$ 否则 $c_n = 0$. 所以, 若不存在 $n \geq 0$ 使得 $V \simeq V_n$, 则有 $\chi_V = 0$, 从而 $V = 0$, 这不可能.
\end{solution}

\begin{exercise}[有限群: 忠实的有限维表示]
\label{ex:IV.9.2}
设 $G$ 是紧群, $(\pi, V)$ 是 $G$ 的忠实有限维连续表示. (回忆这意味着 $\pi \colon G \to \GL(V)$ 单射.) 本题的目标是证明: 若 $G$ 有限, 则 $G$ 的每个不可约表示都是形如 $V^{\otimes n} \otimes (V^*)^{\otimes m}$ (对某个 $n, m \geq 1$) 的表示的直和项, 其中记号 $V^{\otimes n}$ 表示 $\underbrace{V \otimes \cdots \otimes V}_{n \text{ 次}}$, 对 $(V^*)^{\otimes m}$ 类似.
\begin{enumerate}
\item[(a)] 设 $\mathbf{1}$ 是 $G$ 在 $\CC$ 上的平凡表示. 证明只需证明 $G$ 的每个不可约表示都是形如 $(V \oplus V^* \oplus \mathbf{1})^{\otimes N}$ (对某个 $N \geq 1$) 的表示的直和项.
\item[(b)] 设 $W$ 是 $G$ 的不可约表示. 证明 $W$ 是 $(V \oplus V^* \oplus \mathbf{1})^{\otimes N}$ 的直和项当且仅当 $\int_G (1 + 2\Re\chi_V(x))^N\chi_W(x)\,\dd x \neq 0$.
\end{enumerate}
从现在起, 我们假设 $G$ 有限, 固定 $G$ 的有限维表示 $W$, 对每个 $N \in \Z_{\geq 0}$ 写
\[
S_N = \sum_{x \in G}(1 + 2\Re\chi_V(x))^N\chi_W(x).
\]
设 $d = \dim V$.
\begin{enumerate}
\item[(c)] 若 $x \neq 1$, 证明当 $N \to +\infty$ 时 $(1 + 2\Re\chi_V(x))^N\chi_W(x) = o((1 + 2d)^N)$.
\item[(d)] 若 $W \neq 0$, 证明对足够大的 $N$ 有 $S_N \neq 0$.
\end{enumerate}
\end{exercise}

\begin{solution}
(a). 让我们对 $N$ 归纳证明: 对每个 $N \in \Z_{\geq 1}$, 有 $G$-等变同构
\[
(\mathbf{1} \oplus V \oplus V^*)^{\otimes N} \simeq \bigoplus_{\substack{k,l,m \geq 0\\ k+l+m=N}}V^{\otimes k} \otimes (V^*)^{\otimes l}\,^{\oplus \frac{N!}{k!l!m!}}.
\]
这显然蕴含 (a) 的结果.
对 $N = 1$, 它由对 $G$ 的每个表示 $W$ 有 $\mathbf{1} \otimes W \simeq W$ 这一事实得出. 假设对 $N$ 已知结果, 证明对 $N+1$. 有
\begin{align*}
(\mathbf{1} \oplus V \oplus V^*)^{\otimes N+1} &\simeq (\mathbf{1} \oplus V \oplus V^*)^{\otimes N} \otimes (\mathbf{1} \oplus V \oplus V^*) \\
&\simeq (\mathbf{1} \oplus V \oplus V^*) \otimes \bigoplus_{\substack{k,l,m \geq 0\\ k+l+m=N}}V^{\otimes k} \otimes (V^*)^{\otimes l}\,^{\oplus \frac{N!}{k!l!m!}} \\
&\simeq \bigoplus_{\substack{k,l,m \geq 0\\ k+l+m=N}}V^{\otimes k} \otimes (V^*)^{\otimes l}\,^{\oplus \frac{N!}{k!l!m!}} \oplus \bigoplus_{\substack{k,l,m \geq 0\\ k+l+m=N}}V^{\otimes k+1} \otimes (V^*)^{\otimes l}\,^{\oplus \frac{N!}{k!l!m!}} \\
&\qquad \oplus \bigoplus_{\substack{k,l,m \geq 0\\ k+l+m=N}}V^{\otimes k} \otimes (V^*)^{\otimes l+1}\,^{\oplus \frac{N!}{k!l!m!}} \\
&\simeq \bigoplus_{\substack{k,l,m \geq 0\\ k+l+m=N+1}}V^{\otimes k} \otimes (V^*)^{\otimes l}\,^{\oplus \bigl(\frac{N!}{k!l!(m-1)!} + \frac{N!}{(k-1)!l!m!} + \frac{N!}{k!(l-1)!m!}\bigr)} \\
&\simeq \bigoplus_{\substack{k,l,m \geq 0\\ k+l+m=N+1}}V^{\otimes k} \otimes (V^*)^{\otimes l}\,^{\oplus \frac{(N+1)!}{k!l!m!}}.
\end{align*}
(b). 由 $G$ 的有限维表示的半单性 (推论 \ref{cor:I.3.2.9}) 和 Schur 引理 (定理 \ref{thm:I.3.4.1}), 表示 $W$ 是 $(\mathbf{1} \oplus V \oplus V^*)^{\otimes N}$ 的直和项当且仅当 $\Hom_G((\mathbf{1} \oplus V \oplus V^*)^{\otimes N}, W) \neq 0$. 由推论 \ref{cor:IV.5.8}, 这当且仅当 $\langle \chi_{(\mathbf{1} \oplus V \oplus V^*)^{\otimes N}}, \chi_W \rangle_{L^2(G)} \neq 0$. 所以结论由 $\chi_{(\mathbf{1} \oplus V \oplus V^*)^{\otimes N}} = (1 + \chi_V + \bar{\chi}_V)^N = (1 + 2\Re\chi_V)^N$ 得出, 这是命题 \ref{prop:IV.5.4} 的直接推论.
(c). 由于 $G$ 紧, 表示 $(\pi, V)$ 可酉化, 所以可以选择同构 $V \simeq \CC^d$ 使得 $\pi(G) \subset \UU(d)$. 设 $z_1, \ldots, z_d$ 是 $\pi(x)$ 的特征值. 由于 $\pi$ 忠实, $\pi(x) \neq 1$, 所以 $z_i$ 中至少有一个不等于 $1$ (我们用 $\pi(x)$ 可对角化这一事实); 所以可以假设 $z_1 \neq 1$. 由于 $|z_1| = 1$, 这蕴含 $-1 \leq \Re z_1 < 1$, 所以 $1 - 2d \leq 1 + 2\sum_{i=1}^d\Re(z_i) < 1 + 2d$, 且 $1 + 2\sum_{i=1}^d\Re z_i < 1 + 2d$.
最后得到
\[
|(1 + 2\Re\chi_V(x))^N\chi_W(x)| \leq (\dim W)\Bigl|1 + 2\sum_{i=1}^d\Re z_i\Bigr|^N = o((1 + 2d)^N).
\]
(d). 由于 $G$ 有限, (c) 蕴含
\[
\sum_{x \in G-\{1\}}(1 + 2\Re\chi_V(x))^N\chi_W(x) = o((1 + 2d)^N).
\]
另一方面, $(1 + 2\Re\chi_V(1))^N\chi_W(1) = (\dim W)(1 + 2d)^N$. 所以
\[
S_N = (\dim W)(1 + 2d)^N + o((1 + 2d)^N),
\]
这蕴含对足够大的 $N$ 有 $S_N \neq 0$.
\end{solution}

\begin{exercise}[任意紧群: 忠实表示与 $V^{\otimes n} \otimes (V^*)^{\otimes m}$]
\label{ex:IV.9.3}
本题的目标是把习题 \ref{ex:IV.9.2} 推广到任意紧群 $G$, 假设关于 Haar 测度的某些条件. 在下一个习题中, 我们给出用矩阵系数的同一结果的另一种方法.
设 $(\rho, V)$ 是 $G$ 的忠实有限维连续表示. 我们要证明 $G$ 的任何不可约表示都是某个 $V^{\otimes N} \otimes (V^*)^{\otimes M}$ 的直和项. 固定 $G$ 上的规范化 Haar 测度 $\mu$.
\begin{enumerate}
\item[(a)] 证明存在同构 $V \simeq \CC^n$ 使得 $\rho$ 诱导 $G$ 与 $\UU(n)$ 的闭子群之间的 (拓扑群的) 同构.
\end{enumerate}
从现在起, 我们假设 $G$ 是 $\UU(n)$ 的闭子群, $V = \CC^n$, $\rho \colon G \to \GL_n(\CC)$ 是包含. 设 $(\pi, W)$ 是 $G$ 的连续非零有限维表示. 定义 $f \colon \UU(n) \to \CC$ 和 $g \colon G \to \CC$ 为 $f(x) = 1 + \operatorname{Tr}(x) + \overline{\operatorname{Tr}(x)}$, $g(x) = \chi_W(x)$.
像习题 \ref{ex:I.5.5.4} 那样定义
\[
L_0 = \{X \in M_n(\CC) \mid \forall t \in \R,\, e^{tX} \in \UU(n)\}
\]
和
\[
L = \{X \in M_n(\CC) \mid \forall t \in \R,\, e^{tX} \in G\}.
\]
回忆我们在习题 \ref{ex:I.5.5.4} 中证明了: 若 $\Omega$ 是 $L$ 中 $0$ 的足够小的邻域, 则 $\exp$ 诱导 $\Omega$ 与 $\exp(\Omega)$ 之间的同胚, 且 $\exp(\Omega)$ 是 $1$ 在 $G$ 中的邻域. 选取同构 $L_0 \simeq \R^m$, 设 $dX$ 是此同构给出的 $L_0$ 上的 Lebesgue 测度. 我们假设: ($*$) 对足够小的 $\Omega$, 存在 $c \in \R_{>0}$ 使得同胚 $\exp \colon \Omega \xrightarrow{\sim} \exp(\Omega)$ 下 Haar 测度 $\mu$ 的原像具有 $h(X)dX$ 的形式, 其中当 $X \to 0$ 时 $h(X) = c + O(\|X\|)$.
注记. 这总是真的, 但我们没有工具证明它. 事实上, $\GL_n(\CC)$ 的闭子群 $G$ 是 Lie 群, 所以 Haar 测度由 $G$ 上的左不变微分形式 $\omega$ 给出, 即 $\omega_g$ 是 $\omega_1$ 经 $g$ 的左平移的拉回. 这和一番努力给出所要的渐近公式.
用 Weyl 积分公式 (见例如 {[}22{]} 的定理 7.16) 来证明本习题中的所有估计实际上会更简单, 因为我们积分的函数是中心函数. 然而, 这需要更多理论 (特别是极大环面).
\begin{enumerate}
\item[(a)] 证明对每个 $x \in \UU(n)$, 有 $f(x) = 1 + 2\sum_{i=1}^n\cos\theta_i$, 其中 $e^{i\theta_1}, \ldots, e^{i\theta_n}$ 是 $x$ 的特征值.
\item[(b)] 若 $\Omega$ 是 $L_0$ 中 $0$ 的邻域, 证明存在 $\delta > 0$ 使得对每个 $x \notin \exp(\Omega)$ 和每个 $N \geq 1$,
\[
|f(x)^N| \leq (1 + 2n - \delta)^N.
\]
\item[(c)] 若 $\Omega$ 是 $L$ 中 $0$ 的邻域, $U = \exp(\Omega)$, 证明存在 $\delta > 0$ 和 $C \in \R_{>0}$ 使得对每个 $N \geq 1$,
\[
\Bigl|\int_{G-U}f(x)^Ng(x)\,\dd\mu(x)\Bigr| \leq C(1 + 2n - \delta)^N.
\]
\item[(d)] 证明当 $X \to 0$ ($L_0$ 中) 时
\[
f(e^X) = (2n + 1)e^{-K(X) + O(\|X\|^4)},
\]
其中 $K(X) = \frac{1}{1+2n}\|X\|^2 = \frac{1}{1+2n}\operatorname{Tr}(X^*X)$.
\item[(e)] 证明若 $\Omega$ 是 $L$ 中以 $0$ 为中心 (有限半径) 的球, 则存在 $D \in \R_{>0}$ 使得当 $N \to +\infty$ 时
\[
\int_{\Omega}e^{-NK(X)}g(e^X)\,\dd X \sim D \cdot N^{-\frac{1}{2}\dim L}.
\]
(提示: 证明当 $X \to 0$ ($L$ 中) 时 $g(e^X) = \dim W + O(\|X\|)$.)
\item[(f)] 证明存在 $1$ 在 $G$ 中的邻域 $U$ 和 $E \in \R_{>0}$ 使得当 $N \to +\infty$ 时
\[
\int_U f(x)^Ng(x)\,\dd\mu(x) \sim E\frac{(2n + 1)^N}{N^{\frac{1}{2}\dim L}}.
\]
\item[(g)] 证明若 $N$ 足够大则 $\int_G f(x)^Ng(x)\,\dd\mu(x) \neq 0$.
\end{enumerate}
\end{exercise}

\begin{solution}
(a). 由于 $G$ 紧, 表示 $(\rho, V)$ 可酉化. 这意味着存在同构 $V \simeq \CC^n$ 使得 $\rho(G) \subset \UU(n)$. 由于表示 $(\rho, V)$ 忠实, 态射 $\rho$ 单射, 所以 $\rho \colon G \to \UU(n)$ 是单射连续映射. 由于 $G$ 紧, 此映射是到其像的同胚.
(b). 设 $D$ 是对角线元素为 $e^{i\theta_1}, \ldots, e^{i\theta_n}$ 的对角矩阵. 由于 $x$ 与 $x^* = x^{-1}$ 交换, 谱定理蕴含存在 $A \in \UU(n)$ 使得 $D = AxA^{-1}$. 由于 $f$ 显然是 $\UU(n)$ 上的中心函数 (甚至在 $\GL_n(\CC)$ 上), 有 $f(x) = f(D)$. 但
\[
f(e^D) = 1 + \sum_{j=1}^n\Re(e^{i\theta_j}) = 1 + 2\sum_{j=1}^n\cos(\theta_j).
\]
(c). 由 (b), 对每个 $x \in \UU(n)$, $f(x) \in \R$ 且 $1 - 2n \leq f(x) \leq 1 + 2n$. 而且, 等式 $f(x) = 1 + 2n$ 只在 $x$ 的所有特征值都等于 $1$ 时可能, 这进而蕴含 $x = 1$, 因为 $x$ 可对角化.
由习题 \ref{ex:I.5.5.4}(f), 我们知道 $\exp(\Omega)$ 包含 $1$ 在 $\UU(n)$ 中的开邻域 $V$. 由于 $\UU(n)$ 紧, 连续函数 $f$ 在 $\UU(n) - V$ 上达到上确界, 且由上段它 $< 1 + 2n$. 所以 $\sup_{x \in \UU(n)-\exp(\Omega)}|f(x)| < 1 + 2n$, 这蕴含所要的结果.
(d). 对每个 $x \in G$, 有 $|g(x)| \leq \dim W$. 所以可以取 $C = \operatorname{vol}(G - U)(\dim W)$ 并应用 (c).
(e). 设 $i\theta_1, \ldots, i\theta_n$ 是 $X$ 的特征值. 由于 $X$ 与 $X^* = -X$ 交换, 存在 $A \in \UU(n)$ 使得 $AXA^{-1} = D$, 其中 $D$ 是对角线元素为 $i\theta_1, \ldots, i\theta_n$ 的对角矩阵. 则 $Ae^XA^{-1} = e^D$ 是对角线元素为 $e^{i\theta_1}, \ldots, e^{i\theta_n}$ 的对角矩阵, 所以 $e^X$ 的特征值是 $e^{i\theta_1}, \ldots, e^{i\theta_n}$, 且
\[
f(e^X) = 1 + 2\sum_{j=1}^n\cos(\theta_j) = 1 + 2n - \sum_{j=1}^n\theta_j^2 + O\Bigl(\sum_{j=1}^n\theta_j^4\Bigr).
\]
有 $X = A^{-1}DA$, 所以 $X^* = -X = -A^{-1}DA$, 从而 $X^*X = -A^{-1}D^2A$, 最后 $\operatorname{Tr}(X^*X) = -\operatorname{Tr}(D^2) = \sum_{j=1}^n\theta_j^2$. 所以 $f(e^X) = 1 + 2n - \sum_{j=1}^n\theta_j^2 + O(\|X\|^4)$. 另一方面,
\[
(2n+1)e^{-K(X)+O(\|X\|^4)} = (2n+1)\Bigl(1 - \frac{1}{2n+1}\|X\|^2 + O(\|X\|^4)\Bigr) = 2n+1 - \sum_{j=1}^n\theta_j^2 + O(\|X\|^4).
\]
(f). 我们先证明提示. 由习题 \ref{ex:I.5.5.4}(g), 存在 $\R$-线性映射 $u \colon L \to \End(W)$ 使得对每个 $X \in L$, $\pi(e^X) = e^{tu(X)}$. 由于 $u$ 是 $\R$-线性的, 它是 $\mathcal{C}^{\infty}$ 的, 所以映射 $U \colon L \to \CC$, $X \longmapsto g(e^X) = \operatorname{Tr}(e^{u(X)})$ 也是 $\mathcal{C}^{\infty}$ 的. 还有 $U(0) = \operatorname{Tr}(\mathrm{id}_W) = \dim W$. 所以得到 $U(X) = \dim W + O(\|X\|)$.
现在我们计算积分. 做代换 $Y = N^{1/2}X$ (并注意 $NK(X) = K(Y)$), 得到
\[
\int_{\Omega}e^{-NK(X)}U(X)\,\dd X = N^{-\frac{1}{2}\dim L}\int_{N^{1/2}\Omega}e^{-K(Y)}U(N^{-1/2}Y)\,\dd Y.
\]
还需证明 $N^{1/2}\Omega$ 上的积分 $\int e^{-K(Y)}U(N^{-1/2}Y)\,\dd Y$ 在 $N \to +\infty$ 时收敛到一个正实数. 首先注意, 由于 $\int_Le^{-K(Y)}\,\dd Y$ 收敛且函数 $g$ 被 $\dim W$ 控制, 有
\[
\Bigl|\int_{N^{1/2}\Omega - N^{1/4}\Omega}e^{-K(Y)}U(N^{-1/2}Y)\,\dd Y\Bigr| \leq \int_{L - N^{1/4}\Omega}e^{-K(Y)}U(N^{-1/2}Y)\,\dd Y \xrightarrow[N \to +\infty]{} 0.
\]
另一方面, 用 $Y \in N^{1/4}\Omega$ 时 $U(N^{-1/2}Y) = \dim W + O(N^{-1/4})$ 这一事实, 得到
\begin{align*}
\lim_{N \to +\infty}\int_{N^{1/4}\Omega}e^{-K(Y)}U(N^{-1/2}Y)\,\dd Y &= \lim_{N \to +\infty}(\dim W)\int_{N^{1/4}\Omega}e^{-K(Y)}\,\dd Y \\
&= (\dim W)\int_Le^{-K(Y)}\,\dd Y.
\end{align*}
由于函数 $Y \longmapsto e^{-K(Y)}$ 在 $L$ 上取正值, 最后一个积分正且实.
(g). 若存在 $L$ 中 $0$ 的邻域 $\Omega$ 使得 $\exp$ 是从 $\Omega$ 到 $U$ 的微分同胚 (通过把 $U$ 取得足够小总可以假设), 则
\[
\int_U f(x)^Ng(x)\,\dd x = (2n + 1)^N\int_{\Omega}e^{-NK(X) + NO(\|X\|^4)}g(e^X)h(X)\,\dd X,
\]
其中 $h(X) = c + O(\|X\|)$, $c \in \R_{>0}$. 这等于
\[
\frac{(2n + 1)^N}{N^{\dim L/2}}\int_{N^{1/2}\Omega}e^{-K(Y) + O(N^{-1}\|Y\|^4)}U(N^{-1/2}Y)h(N^{-1/2}Y)\,\dd Y.
\]
我们可以像 (f) 中那样证明: 若选择 $\Omega$ 是以 $0$ 为中心的球 (我们可以), 则当 $N \to +\infty$ 时积分收敛到 $c(\dim W)\int_Le^{-K(Y)}\,\dd Y$, 这给出结论.
(h). 由 (d) 和 (f), 可以把 $\int_G f(x)^Ng(x)\,\dd x$ 分解成两项的和, 其中一项等价于 $\frac{(2n+1)^N}{(\dim L)N/2}$ 的正倍数, 另一项被 $(1 + 2n - \delta)^N$ 控制 ($\delta > 0$). 当 $N$ 趋于 $+\infty$ 时, 第二项相对于第一项变得可忽略, 所以对足够大的 $N$, 两项的和不能为 $0$.
\end{solution}

\begin{exercise}[矩阵系数方法的推广]
\label{ex:IV.9.4}
本题的目标也是把习题 \ref{ex:IV.9.2} 推广到具有忠实有限维表示 $(\rho, V)$ 的紧群 $G$. 像第 IV.9.3 节开头那样, 可以假设 $G$ 是 $\UU(N)$ 的闭子群, $V = \CC^N$, $\rho \colon G \to \UU(N) \subset \GL_N(\CC)$ 是包含.
回忆 $G$ 的所有矩阵系数的代数 $E \subset \mathcal{C}(G)$ 在第 IV.4.1 节中定义. 像第 IV.4 节那样, 把 $\mathcal{C}(G)$ 和 $E$ 看作 $G \times G$ 的表示, 让 $(x, y) \in G \times G$ 按 $L_x \circ R_y$ 作用.
对 $i, j \in \{1, \ldots, N\}$, 记 $e_{ij} \colon G \to \CC$ 为把 $x \in G \subset \GL_N(\CC)$ 送到其 $(i, j)$ 项的函数.
\begin{enumerate}
\item[(a)] 设 $E'$ 是由 $1$ 和函数 $e_{ij}$ 与 $\overline{e_{ij}}$ ($1 \leq i, j \leq N$) 生成的 $E$ 的子代数. 证明 $E'$ 被 $G \times G$ 的作用稳定.
\item[(b)] 证明存在 $\widehat{G}$ 的子集 $A$ 使得 $E' = \bigoplus_{\pi \in A}E_{\pi}$.
\item[(c)] 证明 $E' = E$.
\item[(d)] 证明 $G$ 的每个不可约表示都是形如 $V^{\otimes n} \otimes (V^*)^{\otimes m}$ (对某个 $n, m \geq 1$) 的表示的直和项.
\end{enumerate}
\end{exercise}

\begin{solution}
(a). 设 $x \in G$, 写 $x^{-1} = (a_{ij})_{1 \leq i,j \leq N}$. 则对所有 $i, j \in \{1, \ldots, N\}$, $L_xe_{ij} = \sum_{k=1}^N a_{ik}e_{kj}$, 所以 $L_xe_{ij} \in E'$. $E'$ 被算子 $R_y$ 稳定的证明类似.
(b). 回忆由定理 \ref{thm:IV.3.8}, $E = \bigoplus_{\pi \in \widehat{G}}E_{\pi}$, 且若写 $\pi \colon G \to \GL(V_{\pi})$, 则 $E_{\pi}$ 作为 $G \times G$ 的表示同构于 $\End(V_{\pi})$. 我们证明 $E_{\pi}$ 作为 $G \times G$ 的表示不可约且两两不同构. 设 $\pi, \pi' \in \widehat{G}$. 则对 $(x, y) \in G \times G$, 由命题 \ref{prop:IV.5.4} 有 $\chi_{E_{\pi}}(x, y) = \chi_{\pi}(y)\chi_{\pi}(x)$, 对 $\pi'$ 类似, 所以
\[
\langle \chi_{E_{\pi}}, \chi_{E_{\pi'}} \rangle_{L^2(G \times G)} = \int_{G^4}\chi_{\pi}(x_1)\chi_{\pi}(x_2)\overline{\chi_{\pi'}(x_3)}\overline{\chi_{\pi'}(x_4)}\,\dd x_1\,\dd x_2\,\dd x_3\,\dd x_4.
\]
由特征标的 Schur 正交性 (推论 \ref{cor:IV.5.8}(ii)), 若 $\pi \not\simeq \pi'$ 这是 $0$; 由推论 \ref{cor:IV.5.8}(iii), 若 $\pi \simeq \pi'$ 这是 $1$. $E_{\pi}$ 的不可约性由此和推论 \ref{cor:IV.5.8}(iii) 得出, $\pi \not\simeq \pi'$ 时 $E_{\pi} \not\simeq E_{\pi'}$ 也由此得出.
对每个 $\pi \in \widehat{G}$, 设 $W_{\pi}$ 是 $E'$ 在正交投影 $L^2(G) \to E_{\pi}$ 下的像 $W$. 由于 $E'$ 被 $G \times G$ 的作用稳定, $W_{\pi}$ 也如此, 所以 $W$ 是 $\{0\}$ 或 $E_{\pi}$, 因为 $E_{\pi}$ 是 $G \times G$ 的不可约表示. 所以, 若 $A = \{\pi \in \widehat{G} \mid W_{\pi} \neq \{0\}\}$, 则 $E' = \bigoplus_{\pi \in A}E_{\pi}$.
(c). $E$ 的子代数 $E'$ 被复共轭稳定, 显然分离 $G$ 上的点, 且包含常值函数 $1$, 所以由 Stone-Weierstrass 定理它在 $\mathcal{C}(G)$ 中稠密. 特别地, 它在 $E$ 中稠密. 但若 $A \subset \widehat{G}$ 是 (b) 的子集, 有 $E' = \bigoplus_{\pi \in A}E_{\pi}$, 所以 $E$ 在 $E$ 中稠密这一事实蕴含 $A = \widehat{G}$, 即 $E' = E$.
(d). 设 $W$ 是 $G$ 的不可约表示, 假设它不是任何形如 $V^{\otimes n} \otimes (V^*)^{\otimes m}$ ($n, m \geq 1$) 的表示的直和项. 由习题 \ref{ex:IV.9.2}(a), 这蕴含 $W$ 不是任何 $(\mathbf{1} \oplus V \oplus V^*)^{\otimes n}$ ($n \geq 1$) 的直和项, 即由 Schur 正交性 (定理 \ref{thm:IV.3.8}(i)), $\chi_W$ 与 $(\mathbf{1} \oplus V \oplus V^*)^{\otimes n}$ ($n \geq 1$) 的任何矩阵系数正交. 但 $(\mathbf{1} \oplus V \oplus V^*)^{\otimes n}$ 的矩阵系数的空间正好是 $E$ 中由至多 $n$ 个函数 $e_{ij}$ 和 $\overline{e_{ij}}$ ($1 \leq i, j \leq N$) 的乘积生成的子空间, 所以所有这些表示的矩阵系数生成的空间是 $E'$, 我们刚刚看到它等于 $E$. $\chi_W$ 不可能与 $E'$ 的每个元素正交, 所以我们得到结果.
\end{solution}


\chapter[Gelfand 对]{Gelfand 对}
% ============ 第 V 部分: \index{Gelfand 对}Gelfand 对 (Gelfand Pairs) ============
% 原书 V.1 - V.8 (原书 pp.231-310)
% 编号: V.1.1 (节级) / V.3.1.1 (子节级) / V.8.1.1 (习题)

在本章中, $G$ 总是一个局部紧群, $K$ 是 $G$ 的紧子群. 我们固定 $G$ 上的左 Haar 测度 $\mu = \mu_G$ 和 $K$ 上的规范化 Haar 测度 $\mu_K$.

\section{不变与双不变函数}

\begin{definition}[左/右/双不变函数]
\label{def:V.1.1}
若对每个 $x \in K$ 有 $L_xf = f$ (相应地 $R_xf = f$, 相应地 $L_xf = R_xf = f$), 则称 $G$ 上的函数 $f$ 被 $K$ 左不变 (left invariant) (相应地右不变 (right invariant), 相应地双不变 (bi-invariant)).
若 $F(G)$ 是 $G$ 上的函数空间 (例如 $\mathcal{C}_c(G)$), 记 $F(K\backslash G)$ (相应地 $F(G/K)$, 相应地 $F(K\backslash G/K)$) 为其左不变 (相应地右不变, 相应地双不变) 函数的子空间.
\end{definition}

设 $\Delta_G$ 是 $G$ 的模函数. 由于 $K$ 紧, 有 $\Delta_G|_K = 1$, 所以可以使用习题 \ref{ex:I.5.3.5} 的结果. 特别地:

\begin{proposition}[平均化算子]
\label{prop:V.1.2}
设 $f \in \mathcal{C}(G)$, 定义两个函数 $f^K \colon G \to \CC$ 和 ${}^Kf \colon G \to \CC$ 为
\[
f^K(x) = \int_K f(xk)\,\dd k
\]
且
\[
{}^Kf(x) = \int_K f(kx)\,\dd k.
\]
则 $f^K$ 右不变, ${}^Kf$ 左不变.
\end{proposition}

\begin{proposition}[商测度的存在性]
\label{prop:V.1.3}
存在 $G/K$ (相应地 $K\backslash G$) 上唯一的正则 Borel 测度 $\mu_{G/K}$ (相应地 $\mu_{K\backslash G}$) 使得对每个 $f \in \mathcal{C}_c(G)$,
\[
\int_G f(x)\,\dd x = \int_{G/K}f^K(x)\,\dd\mu_{G/K}(x)
\]
(相应地
\[
\int_G f(x)\,\dd x = \int_{K\backslash G}{}^Kf(x)\,\dd\mu_{K\backslash G}(x)).
\]
\end{proposition}

\begin{definition}[双平均化]
\label{def:V.1.4}
若 $f$ 是 $G$ 上的连续函数, 记
\[
{}^Kf^K = {}^K(f^K) = ({}^Kf)^K.
\]
换句话说, 这是由
\[
{}^Kf^K(x) = \int_{K \times K}f(kxk')\,\dd k\,\dd k'
\]
定义的 $G$ 上的连续函数.
\end{definition}

注意 ${}^Kf^K$ 显然是双不变函数.

\begin{proposition}[不变性的刻画]
\label{prop:V.1.5}
设 $f \in \mathcal{C}(G)$. 则 $f$ 左不变 (相应地右不变, 相应地双不变) 当且仅当 $f = {}^Kf$ (相应地 $f = f^K$, 相应地 $f = {}^Kf^K$).
\end{proposition}

\begin{proof}
这由命题 \ref{prop:V.1.2} 和 $K$ 上的测度规范化这一事实立即得出.
\end{proof}

\begin{lemma}[双平均化的积分]
\label{lem:V.1.6}
对每个 $f \in \mathcal{C}_c(G)$,
\[
\int_G {}^Kf^K(x)\,\dd x = \int_G f(x)\,\dd x.
\]
\end{lemma}

\begin{proof}
有
\[
\int_G {}^Kf^K(x)\,\dd x = \int_{G \times K^2}f(kxk')\,\dd x\,\dd k\,\dd k' = \int_G f(x)\,\dd x,
\]
因为对所有 $k, k' \in K$,
\[
\int_G f(kxk')\,\dd x = \Delta(k')\int_G f(x)\,\dd x = \int_G f(x)\,\dd x
\]
(由命题 \ref{prop:I.2.8}).
\end{proof}

\begin{proposition}[投影算子]
\label{prop:V.1.7}
设 $(\pi, V)$ 是 $G$ 的酉表示, $P_K \colon V \to V$ 是 $V^K$ 上的正交投影. 则对每个 $v \in V$,
\[
P_K(v) = \int_K \pi(k)(v)\,\dd k.
\]
而且, 若 $f \in \mathcal{C}_c(G)$, $v \in V$, 则 $\pi(f)(P_K(v)) = \pi(f^K)(v)$ 且 $P_K(\pi(f)(v)) = \pi({}^Kf)(v)$. 特别地:
\begin{enumerate}
\item[(i)] 若 $f \in \mathcal{C}_c(G)$, $v \in V^K$, 则 $\pi(f)(v) = \pi(f^K)(v)$.
\item[(ii)] 若 $f \in \mathcal{C}_c(K\backslash G)$, $v \in V$, 则 $\pi(f)(v) \in V^K$.
\end{enumerate}
(回忆 $\pi \colon L^1(G) \to \End(V)$ 在定理 \ref{thm:I.4.2.6} 中定义.)
\end{proposition}

\begin{proof}
设 $v \in V$. 积分 $w := \int_K\pi(k)(v)\,\dd k$ 的存在性由习题 \ref{ex:I.5.6.2} 得出. 若 $x \in K$, 则
\[
\pi(x)(w) = \int_K\pi(xk)(v)\,\dd k = \int_K\pi(k)(v)\,\dd k = w,
\]
所以 $w \in V^K$. 又, 若 $w' \in V^K$, 则
\[
\langle w, w' \rangle = \int_K\langle \pi(k)(v), w' \rangle\,\dd k = \int_K\langle v, \pi(k)^{-1}(w') \rangle\,\dd k = \int_K\langle v, w' \rangle\,\dd k = \langle v, w' \rangle.
\]
所以 $w$ 是 $v$ 在 $V^K$ 上的正交投影.
现在我们证明最后一个结论. 设 $f \in \mathcal{C}_c(G)$, $v \in V$. 则:
\begin{align*}
\pi(f^K)(v) &= \int_G f^K(x)\pi(x)(v)\,\dd x = \int_{G \times K}f(xk)\pi(x)(v)\,\dd x\,\dd k \\
&= \int_{G \times K}f(x)\pi(x)\pi(k)^{-1}(v)\,\dd x\,\dd k \\
&= \int_{G \times K}f(x)\pi(x)\pi(k)(v)\,\dd x\,\dd k \quad \text{($K$ 幺模)} \\
&= \pi(f)(P_K(v)).
\end{align*}
另一方面:
\begin{align*}
P_K(\pi(f)(v)) &= \int_{K \times G}f(x)\pi(kx)(v)\,\dd k\,\dd x \\
&= \int_{K \times G}f(k^{-1}x)\pi(x)(v)\,\dd k\,\dd x \\
&= \int_{K \times G}f(kx)\pi(x)(v)\,\dd k\,\dd x \\
&= \pi({}^Kf)(v).
\end{align*}
\end{proof}

同样的证明给出:

\begin{proposition}[卷积与平均化]
\label{prop:V.1.8}
设 $f, g \in \mathcal{C}_c(G)$. 则
\[
{}^K(f * g) = ({}^Kf) * g \quad \text{且} \quad (f * g)^K = f * (g^K).
\]
特别地, 若 $f$ 和 $g$ 双不变, 则 $f * g$ 也双不变, 所以 $\mathcal{C}_c(K\backslash G/K)$ 对卷积乘积是 $\mathcal{C}_c(G)$ 的子代数.
\end{proposition}

\begin{remark}[$L^p$ 空间的情形]
\label{rem:V.1.9}
设 $L^p(K\backslash G/K)$ 是 $L^p(G)$ 中双不变函数的子空间. 则若 $1 \leq p < +\infty$, $f \in L^1(K\backslash G/K)$, $g \in L^p(K\backslash G/K)$, 它们的卷积乘积 $f * g$ 在 $L^p(K\backslash G/K)$ 中. 这在定义 $f * g$ 的公式 (见命题 \ref{prop:I.4.1.3}) 上显然; 事实上,
\[
f * g(x) = \int_G f(y)g(y^{-1}x)\,\dd x = \int_G f(xy^{-1})g(y)\,\dd y
\]
(第一个公式表明 $f * g$ 右不变, 第二个表明 $f * g$ 左不变).
特别地, $L^1(G)$ 的子空间 $L^1(K\backslash G/K)$ 是子代数, 且若 $G$ 紧, 对 $L^2$ 空间有类似结果.
\end{remark}

\begin{remark}[用测度代数理解]
\label{rem:V.1.10}
如果我们把卷积乘积和表示 $\pi$ 延拓到 $G$ 上 Radon 测度的空间 $\mathcal{M}(G)$ 上, 这一切更容易记住. (见注记 \ref{rem:I.4.1.6}.) 可以通过把 $\mu_K$ 等同于 Radon 测度 $\mathcal{C}_c(G) \to \CC$, $f \longmapsto \int_K f(x)\,\dd\mu_K(x)$ 把它看成 $\mathcal{M}(G)$ 的元素. 则 $\mu_K * \mu_K = \mu_K$, $f^K = f * \mu_K$, ${}^Kf = \mu_K * f$, $P_K = \pi(\mu_K)$, 所以例如命题 \ref{prop:V.1.7} 的最后部分只是由 $\pi$ 是 $*$-同态得出.
\end{remark}

\section{Gelfand 对的定义}

\begin{definition}[Gelfand 对 (Gelfand pair)]
\label{def:V.2.1}
若代数 $\mathcal{C}_c(K\backslash G/K)$ 对卷积乘积交换, 则称 $(G, K)$ 为 Gelfand 对.
\end{definition}

\begin{remark}[定义的等价形式]
\label{rem:V.2.2}
若 $p \in [1, +\infty)$, $f \in L^p(K\backslash G/K)$, $g \in \mathcal{C}_c(G)$, 则
\begin{align*}
\|{}^Kg^K - {}^Kf^K\|_p^p &= \int_G\Bigl|f(x) - \int_{K \times K}g(kxk')\,\dd k\,\dd k'\Bigr|^p\,\dd x \\
&= \int_G\Bigl|\int_{K \times K}(f(kxk') - g(kxk'))\,\dd k\,\dd k'\Bigr|^p\,\dd x.
\end{align*}
所以, 由 Minkowski 不等式 (见习题 \ref{ex:I.5.6.7}),
\[
\|{}^Kf^K - {}^Kg^K\|_p \leq \int_{K \times K}\|L_kR_{k'}f - L_kR_{k'}g\|_p\,\dd k\,\dd k' = \|f - g\|_p.
\]
由于 $\mathcal{C}_c(G)$ 在 $L^p(G)$ 中稠密, $L^p(K\backslash G/K)$ 的每个函数都可以由 $\mathcal{C}_c(G)$ 的元素逼近, 从而由上面的计算, 可以由 $\mathcal{C}_c(K\backslash G/K)$ 的元素逼近. 换句话说, 空间 $\mathcal{C}_c(K\backslash G/K)$ 在 $L^p(K\backslash G/K)$ 中稠密. 所以在 Gelfand 对的定义中, 可以把条件``$\mathcal{C}_c(K\backslash G/K)$ 对卷积乘积交换''换成条件``$L^1(K\backslash G/K)$ 对卷积乘积交换'' (或对紧 $G$, 换成``$L^2(K\backslash G/K)$ 对卷积乘积交换'').
\end{remark}

\begin{example}[Gelfand 对的例子]
\label{exa:V.2.3}
若 $G$ 交换, 则 $(G, \{1\})$ 是 Gelfand 对.
以下是其他例子 (但我们还不证明它们是 Gelfand 对):
\begin{itemize}
\item $(\SO(n+1), \SO(n))$, 其中 $\SO(n)$ 通过把 $x \in \SO(n)$ 送到 $(n+1) \times (n+1)$ 矩阵 $\begin{pmatrix} x & 0 \\ 0 & 1 \end{pmatrix}$ 等同于 $\SO(n+1)$ 的子群;
\item $(S_{n+m}, S_n \times S_m)$;
\item $(\GL_n(\Q_p), \GL_n(\Z_p))$.
\end{itemize}
\end{example}

\begin{proposition}[Gelfand 对中的群幺模]
\label{prop:V.2.4}
设 $(G, K)$ 是 Gelfand 对. 则 $G$ 幺模.
\end{proposition}

\begin{proof}
由命题 \ref{prop:I.2.12}, 对每个 $f \in \mathcal{C}_c(G)$,
\[
\int_G f(x)\,\dd x = \int_G\Delta(x)^{-1}f(x^{-1})\,\dd x.
\]
所以只需证明对每个 $f \in \mathcal{C}_c(G)$ 有 $\int_G f(x)\,\dd x = \int_G f(x^{-1})\,\dd x$. 首先注意
\[
\int_G {}^Kf^K(x)\,\dd x = \int_G f(x)\,\dd x
\]
且
\[
\int_G {}^Kf^K(x^{-1})\,\dd x = \int_G f(x^{-1})\,\dd x,
\]
由引理 \ref{lem:V.1.6}. 所以只需证明对每个 $f \in \mathcal{C}_c(K\backslash G/K)$ 有 $\int_G f(x)\,\dd x = \int_G f(x^{-1})\,\dd x$. 固定 $f \in \mathcal{C}_c(K\backslash G/K)$. 可以找到 $g \in \mathcal{C}_c(K\backslash G/K)$ 使得 $g$ 在 $(\operatorname{supp}f) \cup (\operatorname{supp}f)^{-1}$ 上等于 $1$ (因为 $\operatorname{supp}f = K(\operatorname{supp}f)K$). 则
\[
f * g(1) = \int_G f(y)g(y^{-1})\,\dd y = \int_{\operatorname{supp}f}f(y)\,\dd y = \int_G f(y)\,\dd y
\]
且
\[
g * f(1) = \int_G g(y)f(y^{-1})\,\dd y = \int_{(\operatorname{supp}f)^{-1}}f(y^{-1})\,\dd y = \int_G f(y^{-1})\,\dd y.
\]
但 $f * g = g * f$, 因为 $(G, K)$ 是 Gelfand 对, 所以这蕴含所要的结果.
\end{proof}

下面的判据将帮助我们找到更多 Gelfand 对.

\begin{proposition}[对合判据]
\label{prop:V.2.5}
假设存在连续自同构 $\theta \colon G \to G$ 使得:
\begin{enumerate}
\item[(a)] $\theta^2 = \mathrm{id}_G$ (即 $\theta$ 是对合);
\item[(b)] 对每个 $x \in G$, $\theta(x) \in Kx^{-1}K$.
\end{enumerate}
则 $(G, K)$ 是 Gelfand 对.
\end{proposition}

\begin{proof}
考虑线性泛函 $\mathcal{C}_c(G) \to \CC$, $f \longmapsto \int_G f(\theta(x))\,\dd x$. 这是 $\mathcal{C}_c(G)$ 上的左不变正线性泛函, 所以由定理 \ref{thm:I.2.7} 中的唯一性, 存在 $c \in \R_{>0}$ 使得对每个 $f \in \mathcal{C}_c(G)$,
\[
\int_G f(\theta(x))\,\dd x = c\int_G f(x)\,\dd x.
\]
由于 $\theta^2 = \mathrm{id}_G$, 必有 $c^2 = 1$, 所以 $c = 1$.
设 $f, g \in \mathcal{C}_c(G)$. 一方面, 对每个 $x \in G$,
\begin{align*}
(f \circ \theta) * (g \circ \theta)(x) &= \int_G f(\theta(y))g(\theta(y)^{-1}\theta(x))\,\dd y \\
&= \int_G f(y)g(y^{-1}\theta(x))\,\dd y \\
&= (f * g) \circ \theta(x)
\end{align*}
(第二个等式由本证明的第一段得出). 另一方面, 对每个 $x \in G$,
\begin{align*}
(g * f)(x^{-1}) &= \int_G g(x^{-1}y)f(y)\,\dd y = \int_G f'(y^{-1})g'(yx)\,\dd y = (f' * g')(x),
\end{align*}
其中 $f'(z) = f(z^{-1})$, $g'(z) = g(z^{-1})$. (我们用 $G$ 幺模这一事实做代换 $y \longmapsto y^{-1}$.)
假设 $f$ 和 $g$ 双不变. 则由条件 (b), $f(\theta(x)) = f(x^{-1})$, $g(\theta(x)) = g(x^{-1})$, 且对 $g * f$ 有类似的等式, 因为 $g * f$ 也双不变, 所以对每个 $x \in G$,
\[
(f * g)(\theta(x)) = ((f \circ \theta) * (g \circ \theta))(x) = (f' * g')(x) = (g * f)(x^{-1}) = (g * f)(\theta(x)).
\]
由于 $\theta$ 是自同构, 这蕴含 $f * g = g * f$.
\end{proof}

\begin{example}[对合判据的应用]
\label{exa:V.2.6}
\begin{enumerate}
\item[(1)] 若 $G$ 交换, 则取 $\theta \colon x \longmapsto x^{-1}$, 所以对任何紧子群 $K$, 特别地对 $K = \{1\}$, $(G, K)$ 是 Gelfand 对.
\item[(2)] 若 $G$ 紧, 则 $(G \times G, \{(x, x) \mid x \in G\})$ 是 Gelfand 对. 事实上, 只需把上面的命题应用于 $\theta(x, y) = (y, x)$. 对每个 $(x, y) \in G \times G$, 有 $\theta(x, y) = (x, x)(x^{-1}, y^{-1})(y, y)$.
\end{enumerate}
\end{example}

\section{Gelfand 对与表示}
在本节中, 我们给出 $(G, K)$ 成为 Gelfand 对的两个表示论判据, 一个在一般情况下有效, 一个对紧 $G$ 有效.

\subsection{\texorpdfstring{Gelfand 对与 $K$-固定向量}{Gelfand 对与 K-固定向量}}

\begin{theorem}[$K$-固定向量判据]
\label{thm:V.3.1.1}
$(G, K)$ 是 Gelfand 对当且仅当: 对 $G$ 的每个不可约酉表示 $(\pi, V)$, 有 $\dim(V^K) \leq 1$.
\end{theorem}

我们需要下面这个 Gelfand-Raikov 定理的变体.

\begin{lemma}[积分分离]
\label{lem:V.3.1.2}
设 $f \in \mathcal{C}_c(G)$. 若 $f \neq 0$, 则存在 $\phi \in E(P_1)$ (见第 III.3 节) 使得 $\int_G f(x)\phi(x)\,\dd x \neq 0$.
\end{lemma}

\begin{proof}
假设对每个 $\phi \in E(P_1)$ 有 $\int_G f\phi\,\dd\mu = 0$. 由定理 \ref{thm:III.4.1}, 对每个正定型函数 $\phi$ 有 $\int_G f\phi\,\dd\mu = 0$. 由定理 \ref{thm:III.2.5}, 对 $G$ 的每个酉表示 $(\pi, V)$ 和任何 $v \in V$, 有 $\langle \pi(f)(v), v \rangle = 0$. 把它应用于 $G$ 的左正则表示, 得到对每个 $g \in L^2(G)$ 有 $\langle f * g, g \rangle_{L^2(G)} = 0$. 像定理 \ref{thm:III.5.1} 的证明那样, 看到这蕴含对所有 $g_1, g_2 \in L^2(G)$ 有 $\langle f * g_1, g_2 \rangle_{L^2(G)} = 0$. 又像该定理的证明那样, 看到对所有 $g_1, g_2 \in L^2(G)$ 有 $\langle f * g_1, g_2 \rangle_{L^2(G)} = \langle f, g_2 * \tilde{g}_1 \rangle_{L^2(G)}$, 其中 $\tilde{g}_1(x) = \overline{g_1(x^{-1})}$. 所以对所有 $g_1, g_2 \in L^2(G)$ 有 $\langle f, g_1 * g_2 \rangle_{L^2(G)} = 0$. 把它应用于 $g_1 = f$ 和 $g_2 = \psi_U$ (其中 $(\psi_U)_{U \in \mathcal{U}}$ 是近似单位元), 最终得到 $\langle f, f \rangle_{L^2(G)} = 0$, 所以 $f = 0$.
\end{proof}

我们还需要 Schur 引理的下述变体.

\begin{lemma}[交换代数的表示]
\label{lem:V.3.1.3}
设 $A$ 是交换 Banach $*$-代数, $u \colon A \to \End(V)$ 是 $A$ 在非零希尔伯特空间 $V$ 上的表示. 假设被所有 $u(x)$ ($x \in A$) 固定的 $V$ 的闭子空间只有 $\{0\}$ 和 $V$. 则 $\dim V = 1$.
\end{lemma}

\begin{proof}
由假设, 子集 $u(A)$ 满足推论 \ref{cor:II.4.4} 的假设, 所以它在 $\End(V)$ 中的中心化子等于 $\CC\mathrm{id}_V$. 但由于 $A$ 交换, $u(A)$ 的每个元素都在 $u(A)$ 的中心化子中, 所以这蕴含 $\Im(u) \subset \CC\mathrm{id}_V$. 特别地, $V$ 的每个子空间都被 $u(A)$ 的所有元素不变, 所以 $V$ 没有非平凡闭子空间, 这只有 $\dim V \leq 1$ 才可能.
\end{proof}

\begin{lemma}[$V^K$ 上的作用]
\label{lem:V.3.1.4}
设 $(\pi, V)$ 是 $G$ 的酉表示. 则对每个 $f \in L^1(K\backslash G/K)$, $\pi(f)$ 把 $V^K$ 送到自身. 若此外 $\pi$ 不可约, 则被所有 $\pi(f)$ ($f \in L^1(K\backslash G/K)$) 稳定的 $V^K$ 的闭子空间只有 $\{0\}$ 和 $V^K$.
\end{lemma}

\begin{proof}
由命题 \ref{prop:V.1.7}, 对每个 $f \in \mathcal{C}_c(K\backslash G/K)$ 和每个 $v \in V^K$, 有 $\pi(f)(v) \in V^K$. 第一个结论由 $\mathcal{C}_c(K\backslash G/K)$ 在 $L^1(K\backslash G/K)$ 中稠密得出.
为证明第二个结论, 只需证明对每个 $v \in V^K - \{0\}$, 空间 $\{\pi(f)(v) \mid f \in \mathcal{C}_c(K\backslash G/K)\}$ 在 $V^K$ 中稠密. 设 $w \in V^K$, $\varepsilon > 0$. 由于 $V$ 不可约, 空间 $\{\pi(f)(v) \mid f \in L^1(G)\}$ 在 $V$ 中稠密. 由于 $\mathcal{C}_c(G)$ 在 $L^1(G)$ 中稠密, 存在 $f \in \mathcal{C}_c(G)$ 使得 $\|\pi(f)(v) - w\| \leq \varepsilon$. 再次由命题 \ref{prop:V.1.7}, 有 $\pi(f)(v) = \pi(f^K)(v)$, 从而 $\pi({}^Kf^K)(v) = P_K(\pi(f)(v))$, 其中 $P_K$ 是 $V$ 在 $V^K$ 上的正交投影. 由于 $w \in V^K$, 得到 $\|\pi({}^Kf^K)(v) - w\| = \|P_K(\pi(f)(v)) - w\| \leq \|\pi(f)(v) - w\| \leq \varepsilon$.
\end{proof}

\begin{proof}[定理 \ref{thm:V.3.1.1} 的证明]
假设 $(G, K)$ 是 Gelfand 对. 设 $(\pi, V)$ 是 $G$ 的不可约酉表示. 由引理 \ref{lem:V.3.1.4}, $\pi$ 定义从 $L^1(K\backslash G/K)$ 到 $\End(V^K)$ 的 $*$-同态, 且被 $L^1(K\backslash G/K)$ 的所有元素稳定的 $V^K$ 的闭子空间只有 $\{0\}$ 和 $V^K$. 由于 $L^1(K\backslash G/K)$ 交换, 引理 \ref{lem:V.3.1.3} 蕴含 $\dim(V^K) \leq 1$.
我们证明其逆. 假设对 $G$ 的每个不可约酉表示 $(\pi, V)$ 有 $\dim(V^K) \leq 1$. 设非零的 $f \in \mathcal{C}_c(K\backslash G/K)$. 由引理 \ref{lem:V.3.1.2}, 存在 $\phi \in E(P_1)$ 使得 $\int_G f\phi\,\dd\mu \neq 0$. 设 $(\pi, V)$ 是 $G$ 的循环酉表示, $v \in V$ 是循环向量, 满足对每个 $x \in G$ 有 $\phi(x) = \langle \pi(x)(v), v \rangle$ (见定理 \ref{thm:III.2.5}). 则
\[
\int_G f(x)\phi(x)\,\dd x = \int_G f(x)\langle \pi(x)(v), v \rangle\,\dd x = \langle \pi(f)(v), v \rangle,
\]
所以 $\pi(f)(v) \neq 0$. 由定理 \ref{thm:III.3.2}, 表示 $(\pi, V)$ 不可约. 由引理 \ref{lem:V.3.1.4}, $V$ 的自同态 $\pi(f)$ 保持 $V^K$, 且由命题 \ref{prop:V.1.7}, 若 $w$ 是 $v$ 在 $V^K$ 上的正交投影, 则 $\pi(f)(w) = \pi(f)(v) \neq 0$. 特别地, $V$ 的子空间 $V^K$ 非零, 所以由假设 $\dim(V^K) = 1$. 所以 $\End(V^K) = \CC$, 这意味着我们找到了 $*$-同态 $u \colon \mathcal{C}_c(K\backslash G/K) \to \CC$ (把 $g$ 送到 $\pi(g)|_{V^K}$) 使得 $u(f) \neq 0$.
现在设 $f_1, f_2 \in \mathcal{C}_c(K\backslash G/K)$. 由于 $\CC$ 交换, 对每个代数态射 $u \colon \mathcal{C}_c(K\backslash G/K) \to \CC$ 有 $u(f_1 * f_2 - f_2 * f_1) = u(f_1)u(f_2) - u(f_2)u(f_1) = 0$. 由上一段, 这蕴含 $f_1 * f_2 - f_2 * f_1 = 0$, 我们完成了.
\end{proof}

\subsection{Gelfand 对与\index{无重数表示}无重数表示}

\begin{definition}[无重数表示 (multiplicity-free representation)]
\label{def:V.3.2.1}
设 $(\pi, V)$ 是 $G$ 的酉表示, 假设可以把 $V$ 写成 $V = \bigoplus_{i \in I}V_i$, 其中 $V_i$ 是 $V$ 的闭 $G$-不变子空间, 作为 $V$ 的表示不可约.\footnote{若 $G$ 紧这总成立, 见定理 \ref{thm:IV.2.1}.} 则称 $(\pi, V)$ 无重数 (multiplicity-free), 如果对 $G$ 的每个不可约酉表示 $W$, 使得 $V_i$ 与 $W$ 等价的 $i \in I$ 的集合的基数 $\leq 1$.
\end{definition}

注意群 $G$ 通过左平移作用于齐性空间 $G/K$ 上, 所以若 $x \in G$, $f$ 是 $G/K$ 上的函数, 可以定义 $L_xf$ 为 $L_xf(y) = f(x^{-1}y)$.

\begin{definition}[\index{拟正则表示}拟正则表示 (quasiregular representation)]
\label{def:V.3.2.2}
$G$ 在 $L^2(G/K)$ 上的拟正则表示是由 $x \cdot f = L_xf$ (对每个 $x \in G$ 和每个 $f \in L^2(G/K)$) 定义的表示.
\end{definition}

\begin{proposition}[拟正则表示是酉表示]
\label{prop:V.3.2.3}
上面的定义有意义, 且给出 $G$ 的酉表示.
\end{proposition}

\begin{proof}
由 $G/K$ 上测度的定义, 对每个 $f \in \mathcal{C}_c(G/K)$ 和每个 $x \in G$ 有 $\int_{G/K}f\,\dd\mu_{G/K} = \int_{G/K}L_xf\,\dd\mu_{G/K}$. 由于 $\mathcal{C}_c(G/K)$ 在 $L^2(G/K)$ 中稠密, 这蕴含算子 $L_x$ 保持 $L^2(G/K)$ 且是等距. 由命题 \ref{prop:I.3.1.10}, 只需证明对每个 $f \in L^2(G/K)$, 映射 $G \to L^2(G/K)$, $x \longmapsto L_xf$ 连续. 像命题 \ref{prop:I.3.1.13} 的证明那样, 只需对 $f \in \mathcal{C}_c(G/K)$ 证明, 此时它由命题 \ref{prop:I.1.12} 得出.
\end{proof}

\begin{remark}[另一种拟正则表示]
\label{rem:V.3.1}
若让 $G$ 通过右正则表示作用于 $L^2(G)$, 则 $L^2(G/K)$ 是 $L^2(G)$ 中 $K$-不变向量的空间. 拟正则表示正是左正则表示在 $L^2(G/K)$ 上的限制.
我们也可以定义 $L^2(K\backslash G)$ 上的拟正则表示 (这是 $K$ 通过左正则表示作用时 $L^2(G)$ 中 $K$-不变向量的空间, 它通过右正则表示得到 $G$ 的作用). 我们得到的表示酉等价于 $L^2(G/K)$ 上的拟正则表示.
\end{remark}

\begin{theorem}[紧群情形]
\label{thm:V.3.2.4}
假设 $G$ 紧. 则 $(G, K)$ 是 Gelfand 对当且仅当 $G$ 在 $L^2(G/K)$ 上的拟正则表示无重数.
又, 若 $(G, K)$ 是 Gelfand 对, 则有 $G$-等变同构
\[
L^2(G/K) \simeq \bigoplus_{\substack{(\pi, V) \in \widehat{G}\\ V^K \neq 0}}V.
\]
\end{theorem}

\begin{proof}
首先注意, 若 $K$ 通过右正则表示作用, $L^2(G/K)$ 是 $L^2(G)$ 中 $K$-不变向量的空间. Peter-Weyl 定理 (推论 \ref{cor:IV.4.3}) 说, 作为 $G \times G$ 的表示, $L^2(G)$ 同构于 $\bigoplus_{(\pi, V) \in \widehat{G}}\End(V) = \bigoplus_{(\pi, V) \in \widehat{G}}V \otimes_{\CC}V^*$ 的完备化. 所以 $L^2(G/K)$ 作为 $G$ 的表示同构于
\[
\bigoplus_{\substack{(\pi, V) \in \widehat{G}\\ V^K \neq 0}}(V^*)^{\dim(V)}
\]
的完备化.
注意对每个 $(\pi, V) \in \widehat{G}$, 表示 $V^*$ 也不可约; 这例如由推论 \ref{cor:IV.5.8}(iii) 得出, 因为 $\chi_{V^*} = \overline{\chi_V}$, 所以 $\|\chi_{V^*}\|_2 = \|\chi_V\|_2$. 所以表示 $L^2(G/K)$ 无重数当且仅当: 对 $G$ 的每个不可约酉表示 $(\pi, V)$, 有 $V^K = 0$ 或 $\dim(V^K) = 1$. 所以定理的第一个结论由定理 \ref{thm:V.3.1.1} 得出.
现在我们证明第二个结论. 我们已经看到
\[
L^2(G/K) \simeq \bigoplus_{\substack{(\pi, V) \in \widehat{G}\\ V^K \neq 0}}V^*,
\]
所以只需证明: 若 $(\pi, V)$ 是 $G$ 的有限维表示, 则 $V^K \neq 0$ 当且仅当 $(V^*)^K \neq 0$. 把定理 \ref{thm:IV.5.7} 应用于表示 $V$ 和 $V^*$ 到 $K$ 的限制, 得到
\[
\dim(V^K) = \int_K\chi_V(k)\,\dd k
\]
且
\[
\dim((V^*)^K) = \int_K\chi_{V^*}(k)\,\dd k = \int_K\chi_V(k)\,\dd k = \dim(V^K).
\]
\end{proof}

\section{\index{球函数}球函数}
在本节中, 我们假设 $(G, K)$ 是 Gelfand 对.

\begin{definition}[球函数 (spherical function)]
\label{def:V.4.1}
设 $\phi \in \mathcal{C}(K\backslash G/K)$. 若线性泛函 $\chi_{\phi} \colon \mathcal{C}_c(K\backslash G/K) \to \CC$, $f \longmapsto f * \phi(1) = \int_G f(x)\phi(x^{-1})\,\dd x$ 是乘法泛函 (其中 $\mathcal{C}_c(K\backslash G/K)$ 上的乘法是卷积乘积), 则称 $\phi$ 为球函数.
换句话说, 若 $\phi \neq 0$ 且对所有 $f, g \in \mathcal{C}_c(K\backslash G/K)$ 有 $\chi_{\phi}(f * g) = \chi_{\phi}(f)\chi_{\phi}(g)$, 则称函数 $\phi$ 球.
\end{definition}

\begin{example}[交换群的情形]
\label{exa:V.4.2}
若 $G$ 交换且 $K = \{1\}$, 则每个连续群态射 $\phi \colon G \to \CC^{\times}$ 都是球函数. 事实上, 设 $f, g \in \mathcal{C}_c(G)$. 则:
\begin{align*}
\int_G (f * g)(x)\phi(x^{-1})\,\dd x &= \int_{G \times G}f(y)g(y^{-1}x)\phi(x^{-1})\,\dd x\,\dd y \\
&= \int_{G \times G}f(y)g(z)\phi(z^{-1}y^{-1})\,\dd y\,\dd z \\
&= \Bigl(\int_G f(y)\phi(y^{-1})\,\dd y\Bigr)\Bigl(\int_G g(z)\phi(z^{-1})\,\dd z\Bigr).
\end{align*}
这些实际上是该情形下的唯一球函数. (这由下一个命题立即得出.)
\end{example}

\begin{proposition}[球函数的刻画]
\label{prop:V.4.3}
设 $\phi \in \mathcal{C}(K\backslash G/K)$. 下列条件等价:
\begin{enumerate}
\item[(i)] 函数 $\phi$ 球.
\item[(ii)] 函数 $\phi$ 非零, 且对所有 $x, y \in G$,
\[
\int_K\phi(xky)\,\dd k = \phi(x)\phi(y).
\]
\item[(iii)] 有:
\begin{enumerate}
\item[(a)] $\phi(1) = 1$;
\item[(b)] 对每个 $f \in \mathcal{C}_c(K\backslash G/K)$, 存在 $\chi(f) \in \CC$ 使得 $f * \phi = \chi(f)\phi$.
\end{enumerate}
\end{enumerate}
\end{proposition}

\begin{proof}
可以把 $\chi_{\phi}$ 用同样的公式延拓到 $\mathcal{C}_c(G)$ 上, 即 $\chi_{\phi}(f) = \int_G f(x)\phi(x^{-1})\,\dd x$. 注意对所有 $f, g \in \mathcal{C}_c(G)$,
\[
\chi_{\phi}(f * g) = \int_{G \times G}f(y)g(y^{-1}x)\phi(x^{-1})\,\dd x\,\dd y = \int_{G \times G}f(y)g(z)\phi(z^{-1}y^{-1})\,\dd y\,\dd z,
\]
所以
\[
\chi_{\phi}(f * g) - \chi_{\phi}(f)\chi_{\phi}(g) = \int_{G \times G}f(y)g(z)(\phi(z^{-1}y^{-1}) - \phi(z^{-1})\phi(y^{-1}))\,\dd y\,\dd z.
\]
设 $f, g \in \mathcal{C}_c(G)$. 把上面的计算应用于 $f' := {}^Kf^K$ 和 $g' := {}^Kg^K$, 并用 $\phi$ 的双不变性 (以及 $K$ 上测度规范化这一事实), 得到 $\chi_{\phi}(f' * g') - \chi_{\phi}(f')\chi_{\phi}(g')$ 等于
\begin{align*}
&\int_{G^2 \times K^4}f(k_1xk_2)g(k_3yk_4)(\phi(y^{-1}x^{-1}) - \phi(y^{-1})\phi(x^{-1}))\,\dd x\,\dd y\,\dd k_1\,\dd k_2\,\dd k_3\,\dd k_4 \\
&= \int_{G^2 \times K^2}f(x)g(y)(\phi(y^{-1}k_3k_2x^{-1}) - \phi(x^{-1})\phi(y^{-1}))\,\dd x\,\dd y\,\dd k_2\,\dd k_3 \\
&= \int_{G^2}f(x)g(y)\Bigl(\int_K\phi(y^{-1}kx^{-1})\,\dd k - \phi(x^{-1})\phi(y^{-1})\Bigr)\,\dd x\,\dd y.
\end{align*}
这证明 $\chi_{\phi}$ 乘法当且仅当对所有 $x, y \in G$ 有 $\int_G\phi(y^{-1}kx^{-1})\,\dd k = \phi(y^{-1})\phi(x^{-1})$. 由于 $\chi_{\phi} \neq 0$ 当且仅当 $\phi \neq 0$, 这证明 (i) 和 (ii) 等价.
假设 $\phi$ 满足 (iii) 的条件 (a) 和 (b). 则对每个 $f \in \mathcal{C}_c(K\backslash G/K)$,
\[
\chi_{\phi}(f) = f * \phi(1) = \chi(f).
\]
由于 $f \longmapsto \chi(f)$ 乘法 (由卷积乘积的结合性), 这蕴含 $\phi$ 球.
最后, 假设 $\phi$ 球. 要证明满足 (iii) 的条件 (a) 和 (b). 设 $f \in \mathcal{C}_c(K\backslash G/K)$. 则对每个 $x \in G$,
\begin{align*}
f * \phi(x) &= \int_G f(y)\phi(y^{-1}x)\,\dd y \\
&= \int_{G \times K}f(y)\phi(y^{-1}kx)\,\dd y\,\dd k \quad \text{($f$ 的左不变性)} \\
&= \int_G f(y)\phi(y)^{-1}\phi(x)\,\dd y \quad \text{(由 (ii))} \\
&= \chi_{\phi}(f)\phi(x).
\end{align*}
这证明条件 (b). 选取满足 $\chi_{\phi}(f) \neq 0$ 的 $f \in \mathcal{C}_c(K\backslash G/K)$, 把上面的等式应用于 $x = 1$, 得到 $\chi_{\phi}(f) = \chi_{\phi}(f)\phi(1)$, 所以 $\phi(1) = 1$.
\end{proof}

回忆 $L^1(G)$ 对卷积乘积和对合 $f^*(x) = \overline{f(x^{-1})}$\footnote{由于 $(G, K)$ 是 Gelfand 对, 由命题 \ref{prop:V.2.4} 群 $G$ 自动幺模, 所以不需要因子 $\Delta(x)^{-1}$.} 是 Banach $*$-代数. 我们已经看到 $L^1(K\backslash G/K)$ 是 $L^1(G)$ 的交换 Banach 子代数, 且显然它也被对合保持. 所以自然要问 $L^1(K\backslash G/K)$ 的谱是什么.
若 $\phi \in \mathcal{C}_b(K\backslash G/K)$ (注意有界性条件), 则积分 $\int_G f(x)\phi(x^{-1})\,\dd x$ 对每个 $f \in L^1(G)$ 收敛, 所以可以把 $\mathcal{C}_c(K\backslash G/K)$ 上的线性泛函 $\chi_{\phi}$ 延拓成 $L^1(K\backslash G/K)$ 上的有界线性泛函, 仍记为 $\chi_{\phi}$.

\begin{theorem}[球函数与谱]
\label{thm:V.4.4}
映射 $\phi \longmapsto \chi_{\phi}$ 把有界球函数的集合等同于 $\sigma(L^1(K\backslash G/K))$.
\end{theorem}

\begin{example}[交换群的情形]
\label{exa:V.4.5}
若 $G$ 交换且 $K = \{1\}$, 有界球函数就是有界连续群态射 $G \to \CC^{\times}$, 即连续群态射 $G \to S^1$, 也就是 $G$ 的不可约酉表示. 所以我们得到典则双射 $\widehat{G} \xrightarrow{\sim} \sigma(L^1(G))$.
特别地, 此时 $L^1(G)$ 上的每个乘法泛函都是 $*$-同态, 即 Banach $*$-代数 $L^1(G)$ 对称. 这重现了习题 \ref{ex:II.5.4}(c) 的结果.
若 $G$ 紧, 我们将在 (定理 \ref{thm:V.7.1}) 中看到仍然真的每个球函数都定义 $L^1(K\backslash G/K)$ 的 $*$-同态, 即 $L^1(K\backslash G/K)$ 对称. 但一般情形这不是真的.
\end{example}

\begin{proof}[定理 \ref{thm:V.4.4} 的证明]
若 $\phi$ 是有界球函数, 则 $\chi_{\phi}$ 在 $\mathcal{C}_c(K\backslash G/K)$ 上乘法, 从而因 $\mathcal{C}_c(K\backslash G/K)$ 在 $L^1(K\backslash G/K)$ 中稠密也在 $L^1(K\backslash G/K)$ 上乘法.
反之, 设 $\chi \colon L^1(K\backslash G/K) \to \CC$ 是乘法泛函. 由推论 \ref{cor:II.2.6}, 线性泛函 $\chi$ 连续且范数 $\leq 1$.
由注记 \ref{rem:V.2.2}, 线性算子 $\mathcal{C}_c(G) \to \mathcal{C}_c(K\backslash G/K)$, $f \longmapsto {}^Kf^K$ 减小 $L^1$ 范数, 所以它延拓成连续线性算子 $L^1(G) \to L^1(K\backslash G/K)$, 仍记为 $f \longmapsto {}^Kf^K$. 则 $f \longmapsto \chi({}^Kf^K)$ 是 $L^1(G)$ 上的连续线性泛函, 其范数等于 $\chi$ 的范数, 所以存在唯一的 $\phi \in L^{\infty}(G)$ 使得 $\|\phi\|_{\infty} = \|\chi\|_{\mathrm{op}}$ 且对每个 $f \in L^1(G)$,
\[
\int_G f(x)\phi(x^{-1})\,\dd x = \chi({}^Kf^K).
\]
特别地, 对所有 $k, k' \in K$ 和每个 $f \in L^1(G)$,
\begin{align*}
\int_G f(x)\phi(kx^{-1}k')\,\dd x &= \int_G f((k')^{-1}xk^{-1})\phi(x^{-1})\,\dd x \\
&= \chi({}^K(L_{k'}R_{k^{-1}}f)^K) \\
&= \chi({}^Kf^K) \\
&= \int_G f(x)\phi(x^{-1})\,\dd x.
\end{align*}
所以 $\phi$ 双不变.
设 $f, g \in L^1(K\backslash G/K)$. 有
\begin{align*}
\chi(f * g) &= \int_G (f * g)(x)\phi(x^{-1})\,\dd x \\
&= \int_{G \times G}f(y)g(y^{-1}x)\phi(x^{-1})\,\dd x\,\dd y \\
&= \int_{G \times G}f(y)\phi(y^{-1}z)g(z^{-1})\,\dd y\,\dd z \\
&= \int_G (f * \phi)(z)g(z^{-1})\,\dd z.
\end{align*}
由于 $\chi(f * g) = \chi(f)\chi(g) = \chi(f)\int_G \phi(z)g(z^{-1})\,\dd z$, 这蕴含
\[
\int_G ((f * \phi) - \chi(f)\phi)(z)g(z^{-1})\,\dd z = 0.
\]
所以对每个 $f \in L^1(K\backslash G/K)$, 有 $f * \phi = \chi(f)\phi$. 选取满足 $\chi(f) \neq 0$ 的 $f \in \mathcal{C}_c(K\backslash G/K)$. 则 $\chi(f) = f * \phi(1) = \chi(f)\phi(1)$, 所以 $\phi(1) = 1$. 又, 函数 $f * \phi$ 连续, 因为它左一致连续 (注意对每个 $x \in G$, $\|L_x(f * \phi) - f * \phi\|_{\infty} = \|(L_xf - f) * \phi\|_{\infty} \leq \|L_xf - f\|_1\|\phi\|_{\infty}$, 并应用命题 \ref{prop:I.3.1.13}). 所以 $\phi$ 局部几乎处处等于某个双不变连续有界函数, 且这个连续有界函数由命题 \ref{prop:V.4.3} 是球函数.
最后, 设 $\phi'$ 是另一个有界球函数, 使得对每个 $f \in L^1(K\backslash G/K)$,
\[
\int_G f(x)\phi'(x^{-1})\,\dd x = \int_G f(x)\phi(x^{-1})\,\dd x.
\]
我们上面已经看到, 对每个 $f \in L^1(G)$,
\[
\int_G f(x)\phi(x^{-1})\,\dd x = \chi({}^Kf^K) = \int_G {}^Kf^K(x)\phi(x^{-1})\,\dd x,
\]
对 $\phi'$ 有类似的等式. 所以
\[
\int_G f(x)\phi'(x^{-1})\,\dd x = \int_G f(x)\phi(x^{-1})\,\dd x
\]
对每个 $f \in L^1(G)$ 成立, 这蕴含 $\phi' = \phi$.
\end{proof}

\section{正定型球函数}
对第一个结果, 我们不需要假设 $(G, K)$ 是 Gelfand 对.

\begin{proposition}[循环向量与双不变性]
\label{prop:V.5.1}
设 $\phi$ 是 $G$ 上的正定型函数, $(\pi_{\phi}, V_{\phi})$ 和 $v_{\phi} \in V_{\phi}$ 是与 $\phi$ 关联的酉表示和循环向量. (见第 III.2 节.)
则 $v_{\phi} \in V_{\phi}^K$ 当且仅当 $\phi$ 双不变.
\end{proposition}

\begin{proof}
对所有 $k, k' \in K$ 和 $x \in G$,
\[
\phi(kxk') = \langle \pi_{\phi}(kxk')(v_{\phi}), v_{\phi} \rangle = \langle \pi_{\phi}(x)(\pi_{\phi}(k')(v_{\phi})), \pi_{\phi}(k^{-1})(v_{\phi}) \rangle.
\]
所以, 若 $v_{\phi} \in V_{\phi}^K$, 得到 $\phi(kxk') = \phi(x)$. 反之, 假设 $\phi$ 双不变. 在上面的等式中取 $k' = 1$, 看到对每个 $k \in K$ 和每个 $x \in G$,
\[
\phi(x) = \langle \pi_{\phi}(x)(v_{\phi}), v_{\phi} \rangle = \phi(k^{-1}x) = \langle \pi_{\phi}(x)(v_{\phi}), \pi_{\phi}(k)(v_{\phi}) \rangle.
\]
由于 $v_{\phi}$ 是循环向量, $\{\pi_{\phi}(x)(v_{\phi}) \mid x \in G\}$ 的张成在 $V_{\phi}$ 中稠密, 所以这蕴含对每个 $k \in K$ 有 $\pi_{\phi}(k)(v_{\phi}) = v_{\phi}$.
\end{proof}

\begin{theorem}[球函数与\index{端点}端点]
\label{thm:V.5.2}
再次假设 $(G, K)$ 是 Gelfand 对. 设 $\phi$ 是 $G$ 上的连续双不变函数.
若 $\phi$ 是规范化正定型函数 (即 $\phi \in P_1$), 则 $\phi$ 球当且仅当 $\phi \in E(P_1)$, 即当且仅当表示 $(\pi_{\phi}, V_{\phi})$ 不可约.
\end{theorem}

\begin{proof}
记 $(\pi, V)$ 和 $v$ 为 $(\pi_{\phi}, V_{\phi})$ 和 $v_{\phi}$. 由于 $\phi$ 双不变, 由命题 \ref{prop:V.5.1} 知道 $v \in V^K$. 首先假设 $\phi \in E(P_1)$, 即 $\pi$ 不可约. 由定理 \ref{thm:V.3.1.1}, 有 $\dim(V^K) = 1$. 设 $f \in \mathcal{C}_c(K\backslash G/K)$. 则对每个 $x \in G$,
\begin{align*}
f * \phi(x) &= \int_G f(y)\langle \pi(y^{-1}x)(v), v \rangle\,\dd y \\
&= \int_G f(y)\langle \pi(x)(v), \pi(y)v \rangle\,\dd y \\
&= \langle \pi(x)(v), \pi(f)(v) \rangle.
\end{align*}
由于 $\pi(f)(v) \in V^K$, 可以把 $\pi(f)(v)$ 写成 $\pi(f)(v) = \chi(f)v$, 其中 $\chi(f) \in \CC$, 得到对每个 $x \in G$,
\[
f * \phi(x) = \chi(f)\langle \pi(x)(v), v \rangle = \chi(f)\phi(x).
\]
由命题 \ref{prop:V.4.3}, 这蕴含 $\phi$ 球.
反之, 假设 $\phi$ 球. 则再次由命题 \ref{prop:V.4.3}, 存在映射 $\chi \colon \mathcal{C}_c(K\backslash G/K) \to \CC$ 使得对每个 $f \in \mathcal{C}_c(K\backslash G/K)$ 有 $f * \phi = \chi(f)\phi$. 换句话说, 对每个 $f \in \mathcal{C}_c(K\backslash G/K)$ 和每个 $x \in G$,
\begin{align*}
f * \phi(x) &= \int_G f(y)\langle \pi(y^{-1}x)(v), v \rangle\,\dd y = \langle \pi(x)(v), \pi(f)(v) \rangle \\
&= \chi(f)\phi(x) \\
&= \chi(f)\langle \pi(x)(v), v \rangle.
\end{align*}
由于 $v$ 是循环向量, 这蕴含 $\pi(f)(v) = \chi(f)v \in V^K$. 但我们在引理 \ref{lem:V.3.1.4} 的证明中看到空间 $\{\pi(f)(v) \mid f \in \mathcal{C}_c(K\backslash G/K)\}$ 在 $V^K$ 中稠密 (若 $v$ 循环), 所以 $\dim(V^K) = 1$. 由引理 \ref{lem:V.5.3}, 这蕴含 $(\pi, V)$ 不可约.
\end{proof}

\begin{lemma}[$K$-固定循环向量与不可约性]
\label{lem:V.5.3}
我们不假设 $(G, K)$ 是 Gelfand 对. 设 $(\pi, V)$ 是 $G$ 的酉表示, 假设 $V^K$ 中存在循环向量. 若 $\dim(V^K) \leq 1$, 则 $(\pi, V)$ 不可约.
\end{lemma}

\begin{proof}
只需证明 $\End_G(V) = \CC\mathrm{id}_V$. 事实上, 若 $V$ 有闭 $G$-不变子空间 $W$ 满足 $W \neq \{0\}, V$, 则到 $W$ 上的正交投影是 $V$ 的 $G$-等变自同态 (由引理 \ref{lem:I.3.4.3}), 它不是 $\mathrm{id}_V$ 的倍数.
所以设 $T \in \End_G(V)$. 则由命题 \ref{prop:V.1.7}, 算子 $T$ 与到 $V^K$ 上的正交投影交换, 所以它保持 $V^K$. 选取循环向量 $v \in V^K$. 由于 $\dim(V^K) = 1$, 有 $T(v) = \lambda v$, 其中 $\lambda \in \CC$. 由于 $T$ 是 $G$-等变的, 得到对每个 $x \in G$ 有 $T(\pi(x)(v)) = \lambda\pi(x)(v)$. 由于 $v$ 循环, 这蕴含 $T = \lambda\mathrm{id}_V$.
\end{proof}

\begin{corollary}[球函数与不可约表示的对应]
\label{cor:V.5.4}
假设 $(G, K)$ 是 Gelfand 对. 则 $\phi \longmapsto (\pi_{\phi}, V_{\phi})$ 诱导从 $E(P_1)$ 中球函数的集合到满足 $V^K \neq \{0\}$ 的 $G$ 的不可约酉表示 $(\pi, V)$ 的酉等价类的集合的双射.
\end{corollary}

\begin{proof}
唯一不立即由命题 \ref{prop:V.5.1} 和定理 \ref{thm:V.5.2} 得出的结论是: 若 $E(P_1)$ 中两个球函数给出酉等价的表示, 则它们必相等. 设 $\phi_1, \phi_2 \in E(P_1)$ 球, 假设存在等距 $G$-等变同构 $T \colon V_{\phi_1} \to V_{\phi_2}$. 由命题 \ref{prop:V.5.1}, 向量 $v_{\phi_1}$ 和 $v_{\phi_2}$ 是 $K$-不变的. 又, 由于 $(G, K)$ 是 Gelfand 对, 空间 $V_{\phi_1}^K$ 和 $V_{\phi_2}^K$ 的维数都 $\leq 1$, 从而因它们包含非零向量而为 $1$. 但 $T$ 是 $G$-等变的, 所以 $T(V_{\phi_1}^K) \subset V_{\phi_2}^K$, 这蕴含对某个 $\lambda \in \CC$ 有 $T(v_{\phi_1}) = \lambda v_{\phi_2}$. 由于 $\|v_{\phi_1}\| = \|v_{\phi_2}\| = 1$, 必有 $|\lambda| = 1$. 所以对每个 $x \in G$,
\begin{align*}
\phi_2(x) &= \langle \pi_{\phi_2}(x)(v_{\phi_2}), v_{\phi_2} \rangle \\
&= \langle \pi_{\phi_2}(x)(\lambda^{-1}T(v_{\phi_1})), \lambda^{-1}T(v_{\phi_1}) \rangle \\
&= \langle T(\pi_{\phi_1}(x)(v_{\phi_1})), T(v_{\phi_1}) \rangle \\
&= \langle \pi_{\phi_1}(x)(v_{\phi_1}), v_{\phi_1} \rangle.
\end{align*}
\end{proof}

\section{\index{对偶空间}对偶空间与\index{球傅里叶变换}球傅里叶变换}
在本节中, 我们假设 $(G, K)$ 是 Gelfand 对. 我们将不加证明地陈述关于 (球)傅里叶变换的一些结果. 在下一节中, 若 $G$ 紧, 我们将给出这些结果的某种版本的证明.

\begin{definition}[对偶空间 (dual space)]
\label{def:V.6.1}
$(G, K)$ 的对偶空间是 $E(P_1)$ 中球函数的集合 $\mathcal{Z}$, 赋予由嵌入 $E(P_1) \subset L^{\infty}(G) \simeq \Hom(L^1(G), \CC)$ 给出的弱*拓扑.
\end{definition}

\begin{example}[交换群的情形]
\label{exa:V.6.2}
若 $G$ 交换且 $K = \{1\}$, 则 $\mathcal{Z} = \widehat{G}$, 即 $G$ 的对偶群. (见习题 \ref{ex:I.5.4.1}.)
\end{example}

\begin{proposition}[$\mathcal{Z}$ 局部紧]
\label{prop:V.6.3}
空间 $\mathcal{Z}$ 局部紧, 且其拓扑与 $G$ 的紧子集上的收敛拓扑一致.
\end{proposition}

\begin{proof}
对第一个结论, 首先注意 $P_0 = \{\psi \text{ 正定型} \mid \psi(1) \leq 1\}$ 弱*紧, 因为它在 $L^{\infty}(G)$ 的闭单位球中弱*闭. 由定理 \ref{thm:V.4.4} 的证明, 子集 $P_0 \cap \mathcal{C}(K\backslash G/K)$ 是满足如下条件的 $\phi \in P_0$ 的集合: 对每个 $f \in L^1(G)$, $\int_G f(x)\phi(x^{-1})\,\dd x = \int_G {}^Kf^K(x)\phi(x^{-1})\,\dd x$. 这些都是弱*闭条件, 所以 $P_0 \cap \mathcal{C}(K\backslash G/K)$ 在 $P_0$ 中弱*闭, 从而弱*紧. 最后, 由定理 \ref{thm:V.5.2}, 集合 $\mathcal{Z} \cup \{0\}$ 是 $P_0 \cap \mathcal{C}(K\backslash G/K)$ 中满足如下条件的 $\phi$ 的集合: 对所有 $f, g \in L^1(K\backslash G/K)$ 有 $\int_G (f * g)(x)\phi(x^{-1})\,\dd x = \bigl(\int_G f(x)\phi(x^{-1})\,\dd x\bigr)\bigl(\int_G g(x)\phi(x^{-1})\,\dd x\bigr)$. 这是弱*闭条件, 所以 $\mathcal{Z} \cup \{0\}$ 弱*紧, 且 $\mathcal{Z}$ 局部紧. 注意这也证明 $\mathcal{Z} \cup \{0\}$ 是 $\mathcal{Z}$ 的 Alexandroff 紧化.
第二个结论由 Raikov 定理 (定理 \ref{thm:III.4.3}) 立即得出.
\end{proof}

\begin{definition}[球傅里叶变换 (spherical Fourier transform)]
\label{def:V.6.4}
设 $f \in L^1(K\backslash G/K)$. $f$ 的 (球)傅里叶变换是函数 $\hat{f} \colon \mathcal{Z} \to \CC$:
\[
\hat{f}(\phi) = \int_G f(x)\phi(x^{-1})\,\dd x = \chi_{\phi}(f).
\]
\end{definition}

\begin{proposition}[傅里叶变换的性质]
\label{prop:V.6.5}
傅里叶变换有下列性质:
\begin{enumerate}
\item[(i)] 对每个 $f \in L^1(K\backslash G/K)$, 函数 $\hat{f} \in \mathcal{C}_0(\mathcal{Z})$, 且 $\|\hat{f}\|_{\infty} \leq \|f\|_1$.
\item[(ii)] 映射 $L^1(K\backslash G/K) \to \mathcal{C}_0(\mathcal{Z})$, $f \longmapsto \hat{f}$ 是 $\CC$-线性的且有稠密像.
\item[(iii)] 对所有 $f, g \in L^1(K\backslash G/K)$, $\widehat{f * g} = \hat{f}\hat{g}$.
\item[(iv)] 对每个 $f \in L^1(K\backslash G/K)$, $\widehat{f^*} = \overline{\hat{f}}$.
\end{enumerate}
\end{proposition}

\begin{proof}
(i). $\hat{f}$ 的连续性由弱*拓扑的定义立即得出. 事实上, 可以用同样的公式把 $\hat{f}$ 延拓成整个 $L^{\infty}(G)$ 上的连续线性泛函. 但我们在命题 \ref{prop:V.6.3} 的证明中看到 $\mathcal{Z} \cup \{0\}$ 是 $\mathcal{Z}$ 的 Alexandroff 紧化, 所以这蕴含 $\hat{f} \in \mathcal{C}(\mathcal{Z})$ 在无穷远处消失. 不等式 $\|\hat{f}\|_{\infty} \leq \|f\|_1$ 由对每个 $\phi \in \mathcal{Z}$ 有 $\|\phi\|_{\infty} = 1$ 这一事实得出.
(iii) 和 (iv). 这只是在表达对每个 $\phi \in \mathcal{Z}$, $\chi_{\phi}$ 是从 $L^1(K\backslash G/K)$ 到 $\CC$ 的 $*$-同态这一事实.
(ii). 线性显然. 第二个结论由 Stone-Weierstrass 定理得出: 事实上, 球傅里叶变换的像是由 (iii) 给出的 $\mathcal{C}_0(\mathcal{Z})$ 的 $\CC$-子代数, 它由 (iv) 被复共轭稳定, 它分离点 (因为由定理 \ref{thm:V.4.4}, 映射 $\mathcal{Z} \to \sigma(L^1(K\backslash G/K))$, $\phi \longmapsto (f \longmapsto \hat{f}(\phi))$ 单射), 且它不在任何地方消失 (对每个 $\phi \in \mathcal{Z}$, 映射 $f \longmapsto \hat{f}(\phi)$ 是 $L^1(K\backslash G/K)$ 上的乘法泛函, 所以它非零).
\end{proof}

\begin{theorem}[Fourier 反演]\footnote{见 {[}25{]} 定理 6.4.5.}
\label{thm:V.6.6}
设 $V^1(K\backslash G/K)$ 是 $G$ 上正定型双不变函数的复线性组合的 $L^1$ 函数的空间.
则存在 $\mathcal{Z}$ 上唯一的测度 $\nu$, 称为 \index{Plancherel 测度}Plancherel 测度, 使得对每个 $f \in V^1(K\backslash G/K)$, $\hat{f} \in L^1(\mathcal{Z}, \nu)$, 且对每个 $x \in G$,
\[
f(x) = \int_{\mathcal{Z}}\phi(x)\hat{f}(\phi)\,\dd\nu.
\]
\end{theorem}

\begin{theorem}[Plancherel 公式]\footnote{见 {[}25{]} 定理 6.4.6.}
\label{thm:V.6.7}
对每个 $f \in \mathcal{C}_c(K\backslash G/K)$, $\hat{f} \in L^2(\mathcal{Z}, \nu)$, 且
\[
\int_G |f(x)|^2\,\dd x = \int_{\mathcal{Z}}|\hat{f}(\phi)|^2\,\dd\nu(\phi).
\]
特别地, 映射 $f \longmapsto \hat{f}$ 延拓成等距 $L^2(K\backslash G/K) \to L^2(\mathcal{Z}, \nu)$, 且它是同构.
\end{theorem}

\begin{remark}[Pontrjagin 对偶]
\label{rem:V.6.8}
若 $G$ 交换且 $K = \{1\}$, 则 $\mathcal{Z} = \widehat{G}$ 是局部紧群, 测度 $\nu$ 是 $\widehat{G}$ 上的 Haar 测度, 且 Pontrjagin 对偶定理说典则映射 $G \to \widehat{\widehat{G}}$, $x \longmapsto (\phi \longmapsto \phi(x))$ 是拓扑群的同构. (见例如 {[}11{]} 定理 4.22 和 4.32, 或 {[}25{]} 定理 5.5.1 和 5.7.1.)
但一般地, Plancherel 测度 $\nu$ 可以支撑在 $\mathcal{Z}$ 的真子集上.
\end{remark}

\section{紧群的情形}
在本节中, 我们假设 $(G, K)$ 是 Gelfand 对, 且 $G$ 紧. 还假设 $G$ 上的 Haar 测度规范化.

\begin{theorem}[紧 Gelfand 对]
\label{thm:V.7.1}
\begin{enumerate}
\item[(i)] $(G, K)$ 的对偶空间 $\mathcal{Z}$ 离散, 且是 $L^2(G)$ 的正交子集.
\item[(ii)] $G$ 上的每个球函数正定型 (从而由定理 \ref{thm:V.5.2} 在 $E(P_1)$ 中). 换句话说, 集合 $\mathcal{Z}$ 通过 $\phi \longmapsto (\pi_{\phi}, V_{\phi})$ 与满足 $\dim(V_{\phi}^K) = 1$ 的 $G$ 的不可约酉表示的等价类的集合典则双射对应.
\item[(iii)] 对每个 $\phi \in \mathcal{Z}$,
\[
\phi(x) = \int_K\chi_{\pi_{\phi}}(xk)\,\dd k
\]
对 $x \in G$, 且
\[
\int_G|\phi(x)|^2 = \frac{1}{\dim V_{\phi}}.
\]
\item[(iv)] 若 $f \in L^2(K\backslash G/K)$, $(\pi, V) \in \widehat{G}$, 则 $V^K = \{0\}$ 时 $f * \chi_{\pi} = 0$, 否则 $f * \chi_{\pi}$ 是由推论 \ref{cor:V.5.4} 对应 $\pi$ 的 $\mathcal{Z}$ 的元素 $\phi_{\pi}$ 的倍数.
\end{enumerate}
\end{theorem}

\begin{proof}
设 $\phi, \phi' \in \mathcal{Z}$ 满足 $\phi \neq \phi'$. 由推论 \ref{cor:V.5.4} (和命题 \ref{prop:IV.2.6}) 知道表示 $V_{\phi}$ 和 $V_{\phi'}$ 酉且不等价. 又知道 (由表示的构造) $\phi$ 和 $\phi'$ 分别是 $V_{\phi}$ 和 $V_{\phi'}$ 的矩阵系数. 由 Schur 正交性 (定理 \ref{thm:IV.3.8}), 这蕴含 $\langle \phi, \phi' \rangle_{L^2(G)} = 0$.
我们证明 $\mathcal{Z}$ 离散. 设 $\phi \in \mathcal{Z}$, 考虑 $U = \{\phi' \in \mathcal{Z} \mid \|\phi - \phi'\|_{\infty} < \|\phi\|_2\}$. 这在 $G$ 的紧子集上的收敛拓扑中开 (因为 $G$ 紧), 从而由 Raikov 定理 (定理 \ref{thm:III.4.3}) 是 $\mathcal{Z}$ 的开子集. 又, 若 $\phi' \in U$, 则
\begin{align*}
|\langle \phi, \phi' \rangle_{L^2(G)}| &= |\langle \phi, \phi \rangle_{L^2(G)} - \langle \phi, \phi' - \phi \rangle_{L^2(G)}| \\
&\geq \|\phi\|_2^2 - \|\phi\|_2\|\phi - \phi'\|_2 \\
&\geq \|\phi\|_2^2 - \|\phi\|_2\|\phi - \phi'\|_{\infty} \\
&> 0,
\end{align*}
所以由第一段, $\phi' = \phi$. 这意味着 $U = \{\phi\}$, 即 $\phi$ 是 $\mathcal{Z}$ 的孤立点.
设 $(\pi, V)$ 是 $G$ 的不可约酉表示, $f \in L^2(K\backslash G/K)$. 要计算 $f * \chi_{\pi}$. 设 $(v_1, \ldots, v_d)$ 是标准正交基; 则对每个 $x \in G$ 有 $\chi_{\pi}(x) = \sum_{i=1}^d\langle \pi(x)(e_i), e_i \rangle$. 所以对每个 $x \in G$,
\[
f * \chi_{\pi}(x) = \int_G f(y)\sum_{i=1}^d\langle \pi(y^{-1}x)(e_i), e_i \rangle = \sum_{i=1}^d\langle \pi(x)(e_i), \pi(f)(e_i) \rangle.
\]
设 $P_K \in \End(V)$ 是 $V^K$ 上的正交投影. 由于 $f$ 双不变, 由命题 \ref{prop:V.1.7} 有 $\pi(f) = P_K \circ \pi(f) \circ P_K$. 首先假设 $V^K = \{0\}$. 则上面的公式给出 $f * \chi_{\pi} = 0$. 现在假设 $V^K \neq \{0\}$. 则由推论 \ref{cor:V.5.4}, 存在唯一正定型球函数 $\phi_{\pi}$, 其关联表示是 $(\pi, V)$, 且有酉循环向量 $v \in V^K$ 使得 $\phi_{\pi}(x) = \langle \pi(x)(v), v \rangle$. 可以选择标准正交基使得 $v_1 = v$. 则对 $i \geq 2$ 有 $P_K(v_i) = 0$, $P_K(v_1) = v_1$, 且对每个 $x \in G$,
\[
f * \chi_{\pi}(x) = \sum_{i=1}^d\langle \pi(x)(v_i), P_K(\pi(f)(P_K(v_i))) \rangle = \langle \pi(x)(v_1), P_K(\pi(f)(v_1)) \rangle.
\]
由于 $V^K$ 是 $1$ 维的, 向量 $P_K(\pi(f)(v_1))$ 是 $v_1$ 的倍数, 所以 $f * \chi_{\pi}$ 是 $\phi_{\pi}$ 的倍数. 这证明 (iv). 又注意对每个 $x \in G$,
\begin{align*}
\int_K\chi_{\pi}(kx)\,\dd k &= \int_K\sum_{i=1}^d\langle \pi(kx)(v_i), v_i \rangle\,\dd k \\
&= \sum_{i=1}^d\Bigl\langle \pi(x)(v_i), \int_K\pi(k^{-1})(v_i)\,\dd k\Bigr\rangle \\
&= \sum_{i=1}^d\langle \pi(x)(v_i), P_K(v_i) \rangle \quad \text{(由命题 \ref{prop:V.1.7})} \\
&= \langle \pi(x)(v_1), v_1 \rangle \\
&= \phi_{\pi}(x),
\end{align*}
这给出 (iii) 的第一部分. (iii) 的第二部分包含在命题 \ref{thm:IV.3.8}(ii) 中.
现在考虑 $G$ 上的球函数 $\phi$. 由命题 \ref{prop:IV.7.1} (即 Fourier 反演公式), 在 $L^2(G)$ 中有等式
\[
\phi = \sum_{\pi \in \widehat{G}}\dim(\pi)\phi * \chi_{\pi}.
\]
由上面的计算, 只有具有非零 $K$-不变向量的 $\pi \in \widehat{G}$ 出现在上面的和中, 且此时 $\phi * \chi_{\pi}$ 是上一段记为 $\phi_{\pi}$ 的函数的倍数. 换句话说, 再次用推论 \ref{cor:V.5.4}, 得到
\[
\phi = \sum_{\psi \in \mathcal{Z}}c_{\psi}\psi,
\]
其中 $c_{\psi} \in \CC$. 若记 $\chi_{\phi}$ (相应地 $\chi_{\psi}$) 为 $L^1(K\backslash G/K)$ 上的线性泛函 $f \longmapsto f * \phi(1)$ (相应地 $f \longmapsto f * \psi(1)$), 知道它是乘法的 (对 $\chi_{\psi}$, 这用定理 \ref{thm:V.5.2}). 又, 由于 $\phi = \sum_{\psi \in \mathcal{Z}}c_{\psi}\psi$, 有 $\chi_{\phi} = \sum_{\psi \in \mathcal{Z}}c_{\psi}\chi_{\psi}$. 设 $\psi, \psi' \in \mathcal{Z}$ 满足 $\psi \neq \psi'$. 则 $\psi * \psi' = \psi' * \psi$ 既是 $\psi$ 又是 $\psi'$ 的倍数 (由命题 \ref{prop:V.4.3}), 所以 $\psi * \psi' = 0$. 特别地, $\chi_{\psi}(\psi') = \chi_{\psi'}(\psi) = 0$. 这蕴含对每个 $\psi \in \mathcal{Z}$ 有 $\chi_{\phi}(\psi) = c_{\psi}\chi_{\psi}(\psi)$; 又注意
\[
\chi_{\psi}(\psi) = \int_G \psi(x)\psi(x^{-1})\,\dd x = \int_G \psi(x)\overline{\psi(x)}\,\dd x > 0.
\]
所以, 若 $\psi, \psi' \in \mathcal{Z}$ 且 $\psi \neq \psi'$, 则
\[
0 = \chi_{\phi}(\psi * \psi') = \chi_{\phi}(\psi)\chi_{\phi}(\psi') = c_{\psi}c_{\psi'}\chi_{\psi}(\psi)\chi_{\psi'}(\psi'),
\]
所以 $c_{\psi}c_{\psi'} = 0$. 所以 $c_{\psi}$ 中至多一个非零, 即存在 $\psi \in \mathcal{Z}$ 使得 $\phi = c_{\psi}\psi$. 由于 $\phi(1) = 1 = \psi(1)$, 必有 $c_{\psi} = 1$, 所以最终看到 $\phi = \psi$ 正定型. 这完成 (ii) 的证明.
\end{proof}

\begin{corollary}[紧 Gelfand 对的调和分析]
\label{cor:V.7.2}
\begin{enumerate}
\item[(i)] 有 $G$-等变同构
\[
L^2(G/K) \simeq \bigoplus_{\phi \in \mathcal{Z}}V_{\phi}.
\]
\item[(ii)] 族 $((\dim V_{\phi})^{1/2}\phi)_{\phi \in \mathcal{Z}}$ 是 $L^2(K\backslash G/K)$ 的 Hilbert 基.
\item[(iii)] 对每个 $f \in L^2(K\backslash G/K)$,
\[
f = \sum_{\phi \in \mathcal{Z}}\dim(V_{\phi})\hat{f}(\phi)\phi
\]
(Fourier 反演公式) 且
\[
\|f\|_{L^2(G)}^2 = \sum_{\phi \in \mathcal{Z}}\dim(V_{\phi})|\hat{f}(\phi)|^2
\]
(Parseval 公式).
\end{enumerate}
\end{corollary}

\begin{proof}
(i) 只是定理 \ref{thm:V.3.2.4} 最后一个结论的重述.
对 (ii), 我们已经知道族 $(\sqrt{\dim(V_{\phi})}\phi)_{\phi \in \mathcal{Z}}$ 在 $L^2(G)$ 中标准正交. 又, 若 $f \in L^2(K\backslash G/K)$, 由命题 \ref{prop:IV.7.1} 有
\[
f = \sum_{(\pi, V) \in \widehat{G}}\dim(V)f * \chi_{\pi},
\]
所以 $f$ 在 $\operatorname{Span}(\mathcal{Z})$ 的闭包中 (由定理的 (iv)), 这意味着 $\operatorname{Span}(\mathcal{Z})$ 在 $L^2(K\backslash G/K)$ 中稠密.
(iii) 的第二个公式由第一个公式和 (ii) 得出. 为证明第一个公式, 只需证明对每个 $f \in L^2(K\backslash G/K)$ 和每个 $\phi \in \mathcal{Z}$ 有 $f * \chi_{\pi_{\phi}} = \hat{f}(\phi)\phi$. 由于我们已经知道 $f * \chi_{\pi_{\phi}}$ 是 $\phi$ 的倍数, 只需验证 $f * \chi_{\pi_{\phi}}(1) = \hat{f}(\phi)$. 由定理的 (iii), 对每个 $x \in G$ 有 $\phi(x) = \int_K\chi_{\pi_{\phi}}(kx)\,\dd k$. 所以:
\begin{align*}
f * \chi_{\pi_{\phi}}(1) &= \int_G f(x)\chi_{\pi_{\phi}}(x^{-1})\,\dd x \\
&= \int_{G \times K}f(k^{-1}x)\chi_{\pi_{\phi}}(x)\,\dd x\,\dd k \quad \text{($f$ 左不变且 } \operatorname{vol}(K) = 1\text{)} \\
&= \int_{G \times K}f(x)\chi_{\pi_{\phi}}(kx)\,\dd x\,\dd k \\
&= \int_G f(x)\phi(x)\,\dd x \\
&= \int_G f(x)\overline{\phi(x^{-1})}\,\dd x \\
&= \hat{f}(\phi).
\end{align*}
\end{proof}

\begin{remark}[Plancherel 测度]
\label{rem:V.7.3}
推论特别地说 $\mathcal{Z}$ 上的 Plancherel 测度 $\nu$ 由 $|\nu(\{\phi\})| = \dim(V_{\phi})$ 给出.
\end{remark}

\section{习题}

\subsection{\texorpdfstring{Gelfand 对 $(\SO(n), \SO(n-1))$}{Gelfand 对}}

这一系列习题的材料是经典的, 但这里的叙述深受 {[}22{]} 的第 2.3.2 节和 {[}25{]} 的第 7.3 节的影响.

\begin{exercise}[调和多项式]
\label{ex:V.8.1.1}
固定正整数 $n$. 对每个 $m \in \Z_{\geq 0}$, 记 $V_m(\R^n)$ 为 $\R^n$ 上 $m$ 次齐次复值多项式函数的向量空间. 定义 $\Or(n)$ 在 $V_m(\R^n)$ 上的作用为 $(x \cdot f)(v) = f(x^{-1}v)$ ($x \in \Or(n)$, $f \in V_m(\R^n)$, $v \in \R^n$; 换句话说 $x \cdot f = L_xf$).
对 $i \in \{1, \ldots, n\}$, 记 $\partial_{x_i}$ 为 $\mathcal{C}^{\infty}(\R^n)$ (从 $\R^n$ 到 $\CC$ 的光滑函数空间) 的自同态 $f \longmapsto \frac{\partial f}{\partial x_i}$, 令 $\Delta = \sum_{i=1}^n(\partial_{x_i})^2$ (这称为 Laplace 算子 (Laplacian operator)).
$\R^n$ 上 $m$ 次调和多项式 (harmonic polynomial) 的空间是
\[
\mathcal{H}_m(\R^n) = \{f \in V_m(\R^n) \mid \Delta(f) = 0\}.
\]
\begin{enumerate}
\item[(a)] 计算 $\dim(V_m(\R^n))$.
\item[(b)] 证明 $\Or(n)$ 在 $V_m(\R^n)$ 上的作用是连续表示.
\item[(c)] 证明对每个 $x \in \Or(n)$ 和每个 $f \in \mathcal{C}^{\infty}(\R^n)$, $\Delta(L_xf) = L_x(\Delta(f))$. (用习题 \ref{ex:V.8.1.2}(c) 可以帮你解决此题.)
\item[(d)] 证明 $V_m(\R^n)$ 的子空间 $\mathcal{H}_m(\R^n)$ 是 $\Or(n)$-不变的.
\end{enumerate}
\end{exercise}

\begin{solution}
(a). 对每个 $i \in \{1, \ldots, n\}$, 记 $x_i \in V_1(\R^n)$ 为函数 $(z_1, \ldots, z_n) \longmapsto z_i$. 则 $\{x_1^{i_1}\cdots x_n^{i_n} \mid i_1, \ldots, i_n \in \Z_{\geq 0}, i_1 + \cdots + i_n = m\}$ 是 $V_m(\R^n)$ 的基. 所以
\[
\dim(V_m(\R^n)) = |\{(i_1, \ldots, i_n) \in (\Z_{\geq 0})^n \mid i_1 + \cdots + i_n = m\}|.
\]
这也等于
\[
|\{(j_1, \ldots, j_n) \in (\Z_{\geq 1})^n \mid j_1 + \cdots + j_n = m + n\}|
\]
(取 $j_r = i_r + 1$). 在上面的集合中选择 $(j_1, \ldots, j_n)$ 等价于选择数 $j_1, j_1 + j_2, \ldots, j_1 + \cdots + j_{n-1}$, 它们构成 $\{1, \ldots, n + m - 1\}$ 的基数为 $n - 1$ 的子集. 所以得到
\[
\dim(V_m(\R^n)) = \binom{n + m - 1}{n - 1} = \binom{n + m - 1}{m}.
\]
(b). 若用 (a) 的 $V_m(\R^n)$ 的基, $x \in \Or(n)$ 的作用由系数是 $x$ 的元素的多元多项式的矩阵给出. 所以对每个 $f \in V_m(\R^n)$, 映射 $\Or(n) \to V_m(\R^n)$, $x \longmapsto x \cdot f$ 连续. 由于 $V_m(\R^n)$ 有限维, 这蕴含作用连续.
(c). 为避免计算, 用习题 \ref{ex:V.8.1.2}(c). 注意 $\Delta = \partial_{x_1^2 + \cdots + x_n^2}$. 所以由习题 \ref{ex:V.8.1.2}(c), 对每个 $f \in \mathcal{C}^{\infty}(\R^n)$ 和每个 $x \in G$,
\[
\Delta(x \cdot f) = x \cdot (\partial_g f),
\]
其中 $g = L_{x^T}(x_1^2 + \cdots + x_n^2)$. 所以我们只需证明 $x_1^2 + \cdots + x_n^2 \in V_2(\R^n)$ 被 $\Or(n)$ 的所有元素不变, 这直接由 $\Or(n)$ 的定义得出.
(d). 问题 (c) 蕴含 $\Delta \colon V_m(\R^n) \to V_{m-2}(\R^n)$ 是 $\Or(n)$-等变的, 且 $\mathcal{H}_m(\R^n)$ 是它的核.
\end{solution}

\begin{exercise}[调和多项式的内积]
\label{ex:V.8.1.2}
我们保留习题 \ref{ex:V.8.1.1} 的记号. 对 $i \in \{1, \ldots, n\}$, 记 $x_i$ 为 $\R^n$ 上的第 $i$ 个坐标函数.
\begin{enumerate}
\item[(a)] 证明映射 $x_i \longmapsto \partial_{x_i}$ 延拓成从 $\bigoplus_{m \geq 0}V_m(\R^n)$ ($\R^n$ 上复值多项式函数的代数) 到 $\End(\mathcal{C}^{\infty}(\R^n))$ 的唯一 $\CC$-代数态射. 记此态射为 $f \longmapsto \partial_f$.
\end{enumerate}
若 $f, g \in V_m(\R^n)$, 令 $\langle f, g \rangle = \partial_g(f)$. (注意 $g$ 仍是 $\R^n$ 上的多项式函数.)
\begin{enumerate}
\item[(a)] 证明 $\langle \cdot, \cdot \rangle$ 是 $V_m(\R^n)$ 上的内积. (提示: 能找到正交基吗?)
\item[(b)] 证明对每个 $f \in V_m(\R^n)$ 和每个 $y \in \Or(n)$, 在 $\End(\mathcal{C}^{\infty}(\R^n))$ 中有 $\partial_f \circ L_y = L_y \circ \partial_{L_{y^T}f}$.
\item[(c)] 证明习题 \ref{ex:V.8.1.1} 中定义的 $\Or(n)$ 在 $V_m(\R^n)$ 上的连续表示对内积 $\langle \cdot, \cdot \rangle$ 酉.
\item[(d)] 若 $m \leq 1$, 证明 $V_m(\R^n) = \mathcal{H}_m(\R^n)$.
\item[(e)] 若 $m \geq 2$, 证明 $\mathcal{H}_m(\R^n)^{\perp} = |x|^2V_{m-2}(\R^n)$, 其中 $|x|^2$ 是函数 $\sum_{i=1}^n x_i^2 \in V_2(\R^n)$.
\item[(f)] 证明
\[
V_m(\R^n) = \bigoplus_{k=0}^{\lfloor m/2 \rfloor}|x|^{2k}\mathcal{H}_{m-2k}(\R^n),
\]
且这诱导 $\Or(n)$-等变同构
\[
V_m(\R^n) = \bigoplus_{k=0}^{\lfloor m/2 \rfloor}\mathcal{H}_{m-2k}(\R^n).
\]
\item[(g)] 若 $S \subset \R^n$ 是单位球面, 证明映射 $\bigoplus_{m \geq 0}\mathcal{H}_m(\R^n) \to \mathcal{C}(S)$, $f \longmapsto f|_S$ 单射.
\item[(h)] 证明对每个 $f \in V_m(\R^n)$, 存在唯一的 $g \in \bigoplus_{k=0}^{\lfloor m/2 \rfloor}\mathcal{H}_{m-2k}(\R^n)$ 使得 $f|_S = g|_S$.
\end{enumerate}
\end{exercise}

\begin{solution}
(a). 注意 $\bigoplus_{m \geq 0}V_m(\R^n)$ 同构于多项式代数 $\CC[x_1, \ldots, x_n]$. 所以只需验证对所有 $i, j \in \{1, \ldots, n\}$, $\partial_{x_i}$ 和 $\partial_{x_j}$ 交换. 但这是 $\mathcal{C}^2$ 函数的偏导数的众所周知的性质.
(b). 首先, 由定义显然 $\langle \cdot, \cdot \rangle$ 对第一个变量线性、对第二个变量反线性. 我们在习题 \ref{ex:V.8.1.1}(a) 的基中计算此形式的矩阵.
设 $f = x_1^{i_1}\cdots x_n^{i_n}$, 其中 $i_1, \ldots, i_n \in \Z_{\geq 0}$, $i_1 + \cdots + i_n = m$. 若 $r \in \{1, \ldots, n\}$, $a \in \Z_{\geq 0}$, 有
\[
\partial_{x_r}^a f = \begin{cases} 0 & \text{若 } a > i_r, \\ i_r(i_r - 1)\cdots(i_r - a + 1)x_r^{i_r - a}\prod_{s \neq r}x_s^{i_s} & \text{若 } a \leq i_r. \end{cases}
\]
设 $g = x_1^{j_1}\cdots x_n^{j_n}$, 其中 $j_1, \ldots, j_n \in \Z_{\geq 0}$, $j_1 + \cdots + j_n = m$. 由于 $i_1 + \cdots + i_n = j_1 + \cdots + j_n$, 要么存在 $r \in \{1, \ldots, n\}$ 使得 $j_r > i_r$, 要么对所有 $r \in \{1, \ldots, n\}$ 有 $i_r = j_r$. 第一种情形 $\langle f, g \rangle = \partial_g f = 0$. 第二种情形,
\[
\langle f, g \rangle = \partial_g f = i_1!i_2!\cdots i_n!.
\]
所以 $\langle \cdot, \cdot \rangle$ 在 $V_m(\R^n)$ 的基 $\{x_1^{i_1}\cdots x_n^{i_n} \mid i_1, \ldots, i_n \in \Z_{\geq 0}, i_1 + \cdots + i_n = m\}$ 中的矩阵是对角矩阵, 对角线元素为正实数, 特别地 Hermitian 正定. 这蕴含 $\langle \cdot, \cdot \rangle$ 是内积.
(c). 此结论实际上对每个 $y \in \GL_n(\R)$ 成立, 我们将证明这一点.
首先注意命题中的恒等式对代数 $V(\R^n) := \bigoplus_{m \geq 0}V_m(\R^n)$ 中的 $f$ 有意义, 且对 $f$ 线性. 又, 若它对 $f, g \in V(\R^n)$ 成立, 则对 $y \in \GL_n(\R)$,
\begin{align*}
\partial_{fg} \circ L_y &= \partial_f \circ \partial_g \circ L_y \\
&= \partial_f \circ L_y \circ \partial_{L_{y^T}g} \\
&= L_y \circ \partial_{L_{y^T}f} \circ \partial_{L_{y^T}g} \\
&= L_y \circ \partial_{L_{y^T}fL_{y^T}g} \\
&= L_y \circ \partial_{L_{y^T}(fg)},
\end{align*}
即恒等式对 $fg$ 也成立. 总之, 只需对函数 $x_1, \ldots, x_n$ 证明.
设 $i \in \{1, \ldots, n\}$, $y \in \GL_n(\R)$, 写 $y^{-1} = a = (a_{ij}) \in \GL_n(\R)$. 则对每个 $(z_1, \ldots, z_n) \in \R^n$,
\[
(L_{y^T}x_i)(z_1, \ldots, z_n) = x_i(a^T(z_1, \ldots, z_n)) = \sum_{j=1}^n a_{ji}z_j.
\]
换句话说,
\[
L_{y^T}x_i = \sum_{j=1}^n a_{ji}x_j,
\]
所以
\[
\partial_{L_{y^T}x_i} = \sum_{j=1}^n a_{ji}\partial_{x_j}.
\]
设 $\phi \in \mathcal{C}^{\infty}(\R^n)$, $z = (z_1, \ldots, z_n) \in \R^n$. 有
\[
\partial_{L_{y^T}x_i}\phi = \sum_{j=1}^n a_{ji}\frac{\partial\phi}{\partial x_j},
\]
所以
\[
L_y\partial_{L_{y^T}x_i}\phi(z) = \sum_{j=1}^n a_{ji}\frac{\partial\phi}{\partial x_j}(az).
\]
另一方面, $L_y\phi(z) = \phi(az)$, 其中
\[
az = \Bigl(\sum_{j=1}^n a_{rj}z_r\Bigr)_{1 \leq r \leq n},
\]
所以
\[
(\partial_{x_i}L_y\phi)(z) = \sum_{r=1}^n a_{ri}\partial_{x_r}\phi(az).
\]
我们看到 $L_y\partial_{L_{y^T}x_i}\phi(z)$ 和 $(\partial_{x_i}L_y\phi)(z)$ 确实得到相同的结果.
(d). 设 $f, g \in V_m(\R^n)$, $y \in \Or(n)$. 由 (c),
\[
\langle L_yf, L_yg \rangle = \partial_{L_yg}L_yf = L_y\partial_{L_{y^T}L_yg}f = L_y\partial_g f.
\]
由于 $\partial_g f$ 是常值函数, 有
\[
L_y\partial_g f = \partial_g f = \langle f, g \rangle,
\]
这正是我们想要的.
(e). 若 $m \leq 1$, 则 $\Delta = 0$ 在 $V_m(\R^n)$ 上, 所以 $\mathcal{H}_m(\R^n) = \Ker(\Delta) = V_m(\R^n)$.
(f). 注意 $\Delta = \partial_{|x|^2}$. 所以若 $f \in V_m(\R^n)$, $g \in V_{m-2}(\R^n)$, 有
\[
\langle f, |x|^2g \rangle = \partial_{|x|^2g}f = \partial_g(\partial_{|x|^2}f) = \langle \Delta f, g \rangle.
\]
换句话说, 映射 $V_{m-2}(\R^n) \to V_m(\R^n)$, $g \longmapsto |x|^2g$ 是 $\Delta \colon V_m(\R^n) \to V_{m-2}(\R^n)$ 的伴随, 所以它的像是 $\Ker\Delta = \mathcal{H}_m(\R^n)$ 的正交补.
(g). 第一个公式只是由 (f) 做简单归纳得出. 对第二个公式, 注意对每个 $k \in \{0, \ldots, \lfloor\frac{m}{2}\rfloor\}$, 单射线性变换 $V_{m-2}(\R^n) \to V_m(\R^n)$, $g \longmapsto |x|^{2k}g$ 是 $\Or(n)$-等变的, 因为函数 $|x|^{2k}$ 被 $\Or(n)$ 不变.
(h). 我们在 $\R^n$ 上用极坐标: $\R^n$ 的点 $z$ 可以写成 $z = rs$, 其中 $r \in \R_{\geq 0}$, $s \in S$, 且 $z \neq 0$ 时 $r$ 和 $s$ 唯一确定.
设 $f \in \mathcal{H}_m(\R^n)$, $g \in \mathcal{H}_p(\R^n)$. 则若 $r \in \R_{\geq 0}$, $s \in S$, 有 $f(rs) = r^mf(s)$, $g(rs) = r^pg(s)$. 由 Green 第二公式,
\[
\int_B (f\Delta g - g\Delta f)\,\dd\lambda = c\int_S\Bigl(f\frac{\partial g}{\partial r} - g\frac{\partial f}{\partial r}\Bigr)\,\dd\mu,
\]
其中 $B$ 是闭单位球, $\lambda$ 是 $\R^n$ 上的 Lebesgue 测度, $c$ 是正常数, $\mu$ 是习题 \ref{ex:I.5.3.5}(f) 中定义的 $S$ 上的测度. 由于 $f$ 和 $g$ 在 $\Delta$ 的核中, 这给出
\[
0 = c(p - m)\int_S f(s)\overline{g(s)}\,\dd\mu(s).
\]
若 $m \neq p$, 得到 $\int_S f(s)\overline{g(s)}\,\dd\mu(s) = 0$. 所以对 $m \neq p$, 子空间 $\mathcal{H}_m(\R^n)|_S$ 和 $\mathcal{H}_p(\R^n)|_S$ 对 $L^2(S, \mu)$ 的内积正交. 特别地, 若 $f \in \bigoplus_{m \geq 0}\mathcal{H}_m(\R^n)$ 满足 $f|_S = 0$, 则把 $f$ 写成 $f = \sum_{m \geq 0}f_m$ (其中 $f_m \in \mathcal{H}_m(\R^n)$), 必有对每个 $m \geq 0$ 有 $f_m|_S = 0$. 但 $f_m$ 是 $m$ 次齐次的, 所以对每个 $r \in \R_{\geq 0}$ 和 $s \in S$ 有 $f_m(rs) = r^mf_m(s)$, 因此 $f_m|_S = 0$ 蕴含 $f_m = 0$.
(i). 存在性由 (g) 的第一个恒等式得出 (因为在 $S$ 上 $|x|^2 = 1$), 唯一性由 (h) 得出.
\end{solution}

\begin{exercise}[$(\SO(n), \SO(n-1))$ 是 Gelfand 对]
\label{ex:V.8.1.3}
设 $n \geq 2$, 通过映射 $x \longmapsto \begin{pmatrix} 1 & 0 \\ 0 & x \end{pmatrix}$ 把 $\Or(n-1)$ 嵌入 $\Or(n)$. 设 $G = \SO(n)$, $K$ 是 $\SO(n-1)$ 在 $G$ 中通过刚定义的嵌入的像.
\begin{enumerate}
\item[(a)] 设 $A$ 是 $\SO(n)$ 中形如
\[
\begin{pmatrix} \cos\theta & 0 & \sin\theta \\ 0 & I_{n-2} & 0 \\ -\sin\theta & 0 & \cos\theta \end{pmatrix}
\]
($\theta \in \R$) 的矩阵的子集.
证明 $A$ 是 $G$ 的子群且 $G = KAK$.
\item[(b)] 证明 $(G, K)$ 是 Gelfand 对. (你可能想用 $G$ 的对合 $\theta(x) = JxJ$, 其中 $J$ 是对角线系数为 $-1, 1, \ldots, 1$ 的对角矩阵.)
\end{enumerate}
\end{exercise}

\begin{solution}
(a). 对每个 $\theta \in \R$, 记 $A_{\theta} = \begin{pmatrix} \cos\theta & 0 & \sin\theta \\ 0 & I_{n-2} & 0 \\ -\sin\theta & 0 & \cos\theta \end{pmatrix}$. 有 $A = \{A_{\theta} \mid \theta \in \R\}$. 容易验证 $A_{\theta}A_{\theta'} = A_{\theta + \theta'}$, 所以 $A$ 是 $G$ 的子群.
设 $v_0 = (1, 0, \ldots, 0) \in S$. 则 $\Or(n)$ 在 $\R^n$ 上的作用保持 $S$ (且 $\Or(n)$ 在 $S$ 上传递作用), $K$ 是 $v_0$ 在 $\Or(n)$ 中的稳定化子. 设 $x \in \Or(n)$, 写 $z = x \cdot v_0 = (z_1, \ldots, z_n)$. 可以找到 $y \in K$ 使得 $y \cdot z = (z_1, 0, \ldots, 0, c)$, 其中 $c^2 = z_2^2 + \cdots + z_n^2$, 然后可以找到 $a \in A$ 使得 $a \cdot (y \cdot z) = (1, 0, \ldots, 0) = v_0$. 则 $(ayx) \cdot v_0 = v_0$, 所以 $ayx \in K$, 且 $x \in Ka^{-1}y^{-1} \subset KAK$.
(b). 想对 $\theta$ 应用命题 \ref{prop:V.2.5}, 其中 $\theta$ 把 $x \in G$ 送到 $JxJ$, $J = \begin{pmatrix} -1 & 0 \\ 0 & I_{n-1} \end{pmatrix}$. 注意 $J^2 = I_n$, 所以 $J = J^{-1}$, 因此 $\theta$ 是群态射且是对合. 显然也有 $\theta(K) = K$. 需要验证对每个 $x \in G$ 有 $\theta(x) \in Kx^{-1}K$. 设 $x \in G$, 写 $x = kak'$, 其中 $k, k' \in K$, $a \in A$. 则 $\theta(x) = \theta(k)\theta(a)\theta(k')$, $\theta(k), \theta(k') \in K$, 且若 $a = A_{\theta}$, 则 $\theta(a) = A_{-\theta} = a^{-1}$. 所以 $\theta(x) = \theta(k)k'x^{-1}k\theta(k') \in Kx^{-1}K$.
\end{solution}

\begin{exercise}[调和表示]
\label{ex:V.8.1.4}
我们用习题 \ref{ex:V.8.1.1} 和 \ref{ex:V.8.1.2} 的记号, 以及习题 \ref{ex:V.8.1.3} 中定义的嵌入 $\Or(n-1) \subset \Or(n)$.
\begin{enumerate}
\item[(a)] 证明有 $\Or(n-1)$-等变同构
\[
V_m(\R^n) \simeq \bigoplus_{k=0}^m V_{m-k}(\R^{n-1}).
\]
\item[(b)] 证明有 $\Or(n-1)$-等变同构
\[
\mathcal{H}_m(\R^n) \simeq \bigoplus_{k=0}^m\mathcal{H}_{m-k}(\R^{n-1}).
\]
\item[(c)] 若 $m \geq 2$, 证明 $\mathcal{H}_m(\R^2)$ 是 $\Or(2)$ 的不可约表示, 但作为 $\SO(2)$ 的表示不可约.
\end{enumerate}
从现在起, 我们假设 $n \geq 3$.
\begin{enumerate}
\item[(a)] 若 $m \geq 1$, 证明 $\mathcal{H}_m(\R^n)^{\SO(n)} = \{0\}$.
\item[(b)] 证明对每个 $m \geq 0$, 空间 $\mathcal{H}_m(\R^n)^{\SO(n-1)}$ 是 $1$ 维的.
\item[(c)] 设 $S \subset \R^n$ 是单位球面, $v_0 = (1, 0, \ldots, 0) \in S$. 证明映射 $\SO(n) \to S$, $x \longmapsto x \cdot v_0$ 诱导同胚 $\SO(n)/\SO(n-1) \xrightarrow{\sim} S$.
\item[(d)] 证明习题 \ref{ex:I.5.3.5}(f) 中定义的 $S$ 上的测度 $\mu$ (用 $\SO(n)$ 和 $\SO(n-1)$ 上的规范化 Haar 测度) 由 $\mu(E) = c\lambda(\{tx \mid t \in [0,1], x \in E\})$ 给出, 其中 $\lambda$ 是 $\R^n$ 上的 Lebesgue 测度, $c^{-1}$ 是单位球 (对 $\lambda$) 的体积.
\item[(e)] 由上一问, 有 $\SO(n)$ 在 $L^2(S)$ 上的拟正则表示, 且它保持连续函数的子空间. 若 $V \subset \mathcal{C}(S)$ 是非零有限维 $\SO(n)$-稳定子空间, 证明 $V^{\SO(n-1)} \neq \{0\}$. (提示: 从满足 $f(v_0) \neq 0$ 的函数 $f \in V$ 开始.)
\item[(f)] 证明 $\SO(n)$ 的表示 $\mathcal{H}_m(\R^n)$ 不可约.
\item[(g)] 证明 $\SO(n)$ 的表示 $\mathcal{H}_m(\R^n)$ 和 $\mathcal{H}_{m'}(\R^n)$ 在 $m \neq m'$ 时不等价. (提示: 比较维数.)
\item[(h)] 若 $m \geq 2$, 证明 $\mathcal{H}_m(\R^n)$ 由函数 $(z_1, \ldots, z_n) \longmapsto (a_1z_1 + \cdots + a_nz_n)^m$ ($a_1, \ldots, a_n \in \CC$ 满足 $a_1^2 + \cdots + a_n^2 = 0$) 张成.
\end{enumerate}
\end{exercise}

\begin{solution}
(a). 若 $f \in V_m(\R^n)$, 则可以把 $f$ 写成 $f = \sum_{k=0}^m x_1^kf_k$, 其中 $f_k \in V_{m-k}(\R^{n-1})$ 唯一确定. 由于 $\Or(n-1) \subset \Or(n)$ 平凡作用于 $x_1$, 这给出 $\Or(n-1)$-等变同构 $V_m(\R^n) \simeq \bigoplus_{k=0}^m V_{m-k}(\R^{n-1})$.
(b). 在此证明中, 约定 $m < 0$ 时 $V_m(\R^n) = 0$. 固定 $n$ 和 $m$. 由习题 \ref{ex:V.8.1.2}(f) 和 (g), 有 $\Or(n)$-等变同构 $V_m(\R^n) \simeq \mathcal{H}_m(\R^n) \oplus V_{m-2}(\R^n)$. 用 (a), 由此推出 $\Or(n-1)$-等变同构
\[
V_m(\R^n) \simeq \mathcal{H}_m(\R^n) \oplus \bigoplus_{k=0}^{m-2}V_{m-2-k}(\R^{n-1}).
\]
另一方面, 把 (a) 应用于 $V_m(\R^n)$ 给出 $\Or(n-1)$-等变同构 $V_m(\R^n) \simeq \bigoplus_{k=0}^m V_{m-k}(\R^{n-1})$. 对每个和项用习题 \ref{ex:V.8.1.2}(f) 或 (g), 得到 $\Or(n-1)$-等变同构
\begin{align*}
V_m(\R^n) &\simeq \bigoplus_{k=0}^m(\mathcal{H}_{m-k}(\R^{n-1}) \oplus V_{m-2-k}(\R^{n-1})) \\
&= \Bigl(\bigoplus_{k=0}^m\mathcal{H}_{m-k}(\R^{n-1})\Bigr) \oplus \Bigl(\bigoplus_{k=0}^{m-2}V_{m-2-k}(\R^{n-1})\Bigr).
\end{align*}
定义 $\Or(n-1)$ 的表示 $V_1, V_2, V_3$ 为 $V_1 = \mathcal{H}_m(\R^n)$, $V_2 = \bigoplus_{k=0}^m\mathcal{H}_{m-k}(\R^{n-1})$, $V_3 = \bigoplus_{k=0}^{m-2}V_{m-2-k}(\R^{n-1})$. 我们刚看到 $V_1 \oplus V_3 \simeq V_2 \oplus V_3$, 所以 $\chi_{V_1} + \chi_{V_3} = \chi_{V_2} + \chi_{V_3}$, 从而 $\chi_{V_1} = \chi_{V_2}$. 由推论 \ref{cor:IV.5.10}, 这蕴含 $V_1$ 和 $V_2$ 等价.
(c). 注意 $\SO(2)$ 是交换群 (它同构于 $S^1$), 所以它的不可约表示都是 $1$ 维的. 另一方面, 由习题 \ref{ex:V.8.1.2}(f) 和习题 \ref{ex:V.8.1.1}(a),
\[
\dim\mathcal{H}_m(\R^2) = \dim V_m(\R^2) - \dim V_{m-2}(\R^2) = \binom{m+1}{m} - \binom{m-1}{m-2} = 2 > 1
\]
若 $m \geq 2$, 所以 $\mathcal{H}_m(\R^n)$ 不能是 $\SO(2)$ 的不可约表示.
若 $m = 2$, $\mathcal{H}_m(\R^2)$ 的基由函数 $x_1^2 - x_2^2$ 和 $x_1x_2$ 给出, 它们都张成被 $\Or(2)$ 的作用稳定的直线, 所以 $\mathcal{H}_2(\R^2)$ 不是 $\Or(2)$ 的不可约表示.
假设 $m \geq 3$. 由于 $\dim\mathcal{H}_m(\R^2) = 2$, 若 $\mathcal{H}_m(\R^2)$ 不是 $\Or(2)$ 的不可约表示, 必有非零的 $f \in \mathcal{H}_m(\R^2)$ 使得对每个 $x \in \Or(2)$ 有 $L_xf \in \CC f$. 按通常方式把 $\R^2$ 等同于复平面 $\CC$. 则对 $z \in \CC$, $f(z)$ 可以写成 $f(z) = \sum_{r=0}^m a_rz^r\bar{z}^{m-r}$, 其中 $a_0, \ldots, a_m \in \CC$. $\SO(2)$ 的作用变成 $S^1$ 在 $\CC$ 上通过乘法的作用, $\begin{pmatrix} 1 & 0 \\ 0 & -1 \end{pmatrix}$ 的作用对应于复共轭. 由对 $f$ 的假设, 对每个 $u \in S^1$, 函数 $f(uz) = \sum_{r=0}^m a_ru^{2r-m}z^r\bar{z}^{m-r}$ 是 $f$ 的倍数. 这只有在存在 $r \in \{0, \ldots, m\}$ 使得对 $s \neq r$ 有 $a_s = 0$ 才可能. 所以可以假设 $f(z) = z^r\bar{z}^{m-r}$. 函数 $\bar{f}(z) = z^{m-r}\bar{z}^r$ 也是 $f$ 的倍数, 所以必有 $m = 2r$, $f = |x|^m$. 则 $\Delta f = m(m-1)|x|^{m-2}$, 与 $\Delta f = 0$ 矛盾.
(d). 设 $f \in V_m(\R^n)^{\SO(n)}$. 由于 $f$ 被 $\SO(n)$ 不变, 它在以 $0$ 为中心的球面上是常数, 所以对每个 $z \in \R^n$ 有 $f(z) = f(\|z\|v_0)$. 由于 $f$ 是 $m$ 次齐次的, 得到对每个 $z \in \R^n$ 有 $f(z) = \|z\|^mf(v_0)$. 所以 $f$ 是多项式当且仅当 $m$ 偶. 又, 容易验证 $\Delta f = m(m-1)|x|^{-2}f$, 所以 $f \in \mathcal{H}_m(\R^n)$ 当且仅当 $f = 0$ 或 $m = 0$.
(e). 由 (b), 有 $\SO(n-1)$-等变同构 $\mathcal{H}_m(\R^n) \simeq \bigoplus_{k=0}^m\mathcal{H}_{m-k}(\R^{n-1})$. 所以由 (d),
\[
\mathcal{H}_m(\R^n)^{\SO(n-1)} \simeq \mathcal{H}_0(\R^{n-1})^{\SO(n-1)} \simeq \CC.
\]
(f). 记映射 $\SO(n) \to S$, $x \longmapsto x \cdot v_0$ 为 $\phi$. 首先, 此映射显然连续, 且它满射, 因为 $\SO(n)$ 在 $S$ 上传递作用. (若 $v_1, v'_1 \in S$, 要找 $x \in \SO(n)$ 使得 $x \cdot v_1 = v'_1$. 可以把 $v_1$ 和 $v'_1$ 完备成 $\R^n$ 的两个标准正交基 $(v_1, \ldots, v_n)$ 和 $(v'_1, \ldots, v'_n)$. 这两个基之间的换基矩阵在 $\Or(n)$ 中. 若它在 $\SO(n)$ 中, 我们完成了. 否则, $(v_1, \ldots, v_n)$ 和 $(v'_1, \ldots, v'_{n-1}, -v'_n)$ 之间的换基矩阵将在 $\SO(n)$ 中.)
又, $v_0$ 在 $\SO(n)$ 中的稳定化子是 $\SO(n)$ 中第一列为 $\begin{pmatrix} 1 \\ 0 \\ \vdots \\ 0 \end{pmatrix}$ 的元素构成的子群, 容易看到这就是 $\SO(n-1)$. 所以 $\phi$ 诱导连续双射 $\SO(n)/\SO(n-1) \to S$. 由于 $\SO(n)/\SO(n-1)$ 紧, 此映射是同胚.
(g). 定义 $S$ 上的正则 Borel 测度 $\nu$:
\[
\nu(E) = c\lambda(\{tx \mid t \in [0,1], x \in E\}).
\]
定义线性泛函 $I \colon \mathcal{C}(\SO(n)) \to \CC$:
\[
I(f) = \int_S f^{\SO(n-1)}(s)\,\dd\nu(s),
\]
其中
\[
f^{\SO(n-1)}(x) = \int_{\SO(n-1)}f(xy)\,\dd y
\]
(用 $\SO(n-1)$ 上的规范化 Haar 测度), 对每个 $s \in \SO(n)$; 函数 $f^{\SO(n-1)}$ 被 $\SO(n-1)$ 右不变, 所以可以由 (f) 等同于 $S$ 上的函数.
这是 $\mathcal{C}(\SO(n))$ 上的正泛函, 所以它来自 $\SO(n)$ 上的正则 Borel 测度, 设为 $\rho$. 要证明 $\rho$ 是 $\SO(n)$ 上的规范化 Haar 测度.
注意若 $f$ 是常值函数 $1$, 则 $I(f) = 1$. 所以为证明 $\rho$ 是规范化 Haar 测度, 只需证明它左不变. 设 $f \in \mathcal{C}(\SO(n))$, $y \in \SO(n)$. 则由定义立即得到 $(L_yf)^{\SO(n-1)} = L_y(f^{\SO(n-1)})$, 所以只需证明 $S$ 上的测度 $\nu$ 被 $\SO(n)$ 的作用左不变. 但这由 Lebesgue 测度 $\lambda$ 被 $\SO(n)$ 的作用左不变 (这可以用换元公式看到) 这一事实立即得出.
(h). 设 $V$ 如题设. 由于 $\SO(n)$ 在 $S$ 上传递作用且 $V$ 被 $\SO(n)$ 稳定, 可以找到 $f \in V$ 使得 $f(v_0) \neq 0$. 设 $(f_1, \ldots, f_r)$ 是 $V$ 的基. 由于 $V$ 被 $\SO(n)$ 稳定且 $\SO(n)$ 在 $V$ 上的作用连续, 可以找到连续函数 $c_1, \ldots, c_r \colon \SO(n) \to \CC$ 使得对每个 $x \in \SO(n)$ 和每个 $s \in S$, $f(x \cdot s) = \sum_{i=1}^r c_i(x)f_i(s)$. 定义 $\tilde{f} \colon S \to \CC$:
\[
\tilde{f}(s) = \int_{\SO(n-1)}f(x \cdot s)\,\dd x.
\]
则 $\tilde{f} = \sum_{i=1}^r\bigl(\int_{\SO(n-1)}c_i(x)\,\dd x\bigr)f_i$, 所以 $\tilde{f} \in V$. 又, $\tilde{f}$ 由构造是 $\SO(n-1)$-不变的. 最后, 由于对每个 $x \in \SO(n-1)$ 有 $x \cdot v_0 = v_0$, 有 $\tilde{f}(v_0) = f(v_0) \neq 0$, 所以 $\tilde{f} \neq 0$.
(i). 由习题 \ref{ex:V.8.1.2}(h), 从 $\R^n$ 到 $S$ 的限制在 $\mathcal{H}_m(\R^n)$ 上单射, 所以 $\mathcal{H}_m(\R^n)$ 作为 $\SO(n)$ 的表示不可约当且仅当 $\mathcal{H}_m = \mathcal{H}_m(\R^n)|_S \subset \mathcal{C}(S)$ 作为 $\SO(n)$ 的表示不可约. 由于 $\SO(n)$ 紧, 若 $\mathcal{H}_m$ 不可约, 则可以把 $\mathcal{H}_m$ 写成 $\mathcal{H}_m = V \oplus V'$, 其中 $V$ 和 $V'$ 是 $\mathcal{H}_m$ 的非零 $\SO(n)$-不变子空间. 由 (h), 这蕴含 $\dim(\mathcal{H}_m^{\SO(n)}) \geq 2$, 与 (d) 矛盾. 所以 $\mathcal{H}_m$ 不可约.
(j). 设 $d_m = \dim\mathcal{H}_m(\R^n)$. 将证明对每个 $m \in \Z_{\geq 0}$ 有 $d_{m+1} > d_m$, 这蕴含 $m \neq m'$ 时 $d_m \neq d_{m'}$, 所以 $\mathcal{H}_m(\R^n)$ 和 $\mathcal{H}_{m'}(\R^n)$ 不等价.
若 $m \leq 1$, 由习题 \ref{ex:V.8.1.1}(a) 和习题 \ref{ex:V.8.1.2}(e), 有 $d_m = \dim V_m(\R^n) = \binom{m+n-1}{m}$. 若 $m \geq 2$, 由习题 \ref{ex:V.8.1.1}(a) 和习题 \ref{ex:V.8.1.2}(f),
\begin{align*}
d_m &= \dim V_m(\R^n) - \dim V_{m-2}(\R^n) \\
&= \binom{m+n-1}{m} - \binom{m+n-3}{m-2} \\
&= \frac{(m+n-3)!}{(m-2)!(n-1)!}\Bigl(\frac{(m+n-1)(m+n-2)}{m(m-1)} - 1\Bigr) \\
&= \frac{(2m+n-2)(m+n-3)!}{m!(n-2)!}.
\end{align*}
特别地, $d_0 = 1$, $d_1 = n$, $d_2 = 2n-1$, 所以 $d_2 > d_1 > d_0$. 设 $m \geq 2$. 则
\begin{align*}
d_{m+1} - d_m &= \frac{(2(m+1)+n-2)(m+1+n-3)!}{(m+1)!(n-2)!} - \frac{(2m+n-2)(m+n-3)!}{m!(n-2)!} \\
&= \frac{(m+n-3)!}{m!(n-2)!}\Bigl(\frac{(2m+n)(m+n-2)}{m+1} - (2m+n-2)\Bigr) \\
&= \frac{(m+n-3)!}{m!(n-2)!}\cdot\frac{(2m+n)(m+n-2) - (2m+n-2)(m+1)}{m+1} \\
&> 0
\end{align*}
(因为 $m + n - 2 \geq m + 1$ 且 $2m + n > 2m + n - 2$).
(k). 设 $a_1, \ldots, a_n \in \CC$, 考虑 $f = (a_1x_1 + \cdots + a_nx_n)^m \in V_m(\R^n)$. 则 $\Delta f = m(m-1)(a_1^2 + \cdots + a_n^2)(a_1x_1 + \cdots + a_nx_n)^{m-2}$, 所以 $f \in \mathcal{H}_m(\R^n)$ 当且仅当 $a_1^2 + \cdots + a_n^2 = 0$. 设
\[
W = \operatorname{Span}\{(a_1x_1 + \cdots + a_nx_n)^m \mid a_1, \ldots, a_n \in \CC, a_1^2 + \cdots + a_n^2 = 0\} \subset \mathcal{H}_m(\R^n).
\]
由于 $W \neq 0$ 且 $\mathcal{H}_m(\R^n)$ 作为 $\SO(n)$ 的表示不可约, 为证明 $W = \mathcal{H}_m(\R^n)$ 只需证明 $W$ 被 $\SO(n)$ 不变. 设 $x \in \SO(n)$, $a_1, \ldots, a_n \in \CC$, 设 $f = (a_1x_1 + \cdots + a_nx_n)^m$. 写 $(b_1, \ldots, b_n) = (a_1, \ldots, a_n)x^T$ (把 $(a_1, \ldots, a_n)$ 看作行向量). 则 $b_1^2 + \cdots + b_n^2 = 0$, 因为 $x \in \Or(n)$, 且 $L_xf = (b_1x_1 + \cdots + b_nx_n)^m$, 所以 $L_xf \in W$.
\end{solution}

\begin{exercise}[球函数]
\label{ex:V.8.1.5}
我们保留习题 \ref{ex:V.8.1.1}-\ref{ex:V.8.1.4} 的记号, 假设 $n \geq 3$.
\begin{enumerate}
\item[(a)] 证明空间 $\sum_{m \geq 0}\mathcal{H}_m(\R^n)|_S$ 在 $L^2(S)$ 中稠密且和是直和.
\item[(b)] 证明 $L^2(S)$ 的子空间 $\mathcal{H}_m(\R^n)|_S$ 和 $\mathcal{H}_{m'}(\R^n)|_S$ 在 $m \neq m'$ 时正交 (对 $L^2$ 内积).
\item[(c)] 证明 $\SO(n)$ 的每个具有非零 $\SO(n-1)$-不变向量的不可约酉表示都同构于某个 $\mathcal{H}_m(\R^n)$.
\end{enumerate}
\end{exercise}

\begin{solution}
(a). 由习题 \ref{ex:V.8.1.2}(i),
\[
\sum_{m \geq 0}\mathcal{H}_m(\R^n)|_S = \sum_{m \geq 0}V_m(\R^n)|_S,
\]
且右端由 Stone-Weierstrass 定理在 $\mathcal{C}(S)$ 中稠密 (从而在 $L^2(S)$ 中稠密). 又, 我们在习题 \ref{ex:V.8.1.2}(h) 的证明中看到空间 $\mathcal{H}_m(\R^n)|_S$ 在 $L^2(S)$ 中两两正交, 所以它们在直和中.
(b). 见 (a).
(c). 设 $V$ 是 $\SO(n)$ 的不可约酉表示, 满足 $V^{\SO(n-1)} \neq 0$. 由定理 \ref{thm:V.3.2.4}, $V$ 是 $L^2(S)$ 的子表示. 但我们已经看到
\[
L^2(S) = \bigoplus_{m \geq 0}\mathcal{H}_m(\R^n)|_S
\]
且所有这些和项都不可约, 所以 $V$ 同构于其中一个.
\end{solution}

\begin{exercise}[\index{Gegenbauer 多项式}Gegenbauer 多项式]
\label{ex:V.8.1.6}
我们保留习题 \ref{ex:V.8.1.1}-\ref{ex:V.8.1.5} 的记号. 若函数 $\phi \in \mathcal{C}(S)$ 被 $\SO(n-1)$ 左不变, 则称它带状 (zonal). (由于 $S = \SO(n)/\SO(n-1)$, 也可以把 $\phi$ 看成 $\SO(n)$ 上的双不变函数.) 假设 $n \geq 3$.
\begin{enumerate}
\item[(a)] 证明 $\phi \in \mathcal{C}(S)$ 带状当且仅当存在连续函数 $f \colon [-1,1] \to \CC$ 使得对每个 $z = (z_1, \ldots, z_n) \in S$, $\phi(z) = f(z_1)$.
\item[(b)] 证明存在 $c \in \R_{>0}$ 使得对每个带状 $\phi \in \mathcal{C}(S)$, 若定义 $f \colon [-1,1] \to \CC$ 如 (a), 则
\[
\int_S\phi(z)\,\dd\mu(z) = c\int_{-1}^1 f(t)(1-t^2)^{(n-3)/2}\,\dd t.
\]
(提示: 可以试着用球坐标, 如 https://en.wikipedia.org/wiki/N-sphere\#Spherical\_coordinates.)
\item[(c)] 设 $m \geq 0$. 若 $t \in S$, 设 $f_t$ 是 $\mathcal{H}_m(\R^n)$ 中满足对每个 $g \in \mathcal{H}_m(\R^n)$ 有 $\langle g, f_t \rangle = g(t)$ 的唯一元素. (注意我们用习题 \ref{ex:V.8.1.2} 的内积.)
证明函数 $Z_m = f_{v_0}|_S$ (其中 $v_0 = (1, 0, \ldots, 0)$) 是\index{带状函数}带状函数.
\item[(d)] 设 $f_m \colon [-1,1] \to \CC$ 是 (a) 中对应于 $Z_m$ 的连续函数. 证明 $f_m$ 是次数 $\leq m$ 的多项式函数.
\item[(e)] 若 $m \neq m'$, 证明 $\int_{-1}^1 f_m(t)f_{m'}(t)(1-t^2)^{(n-3)/2}\,\dd t = 0$.
\item[(f)] 证明 $f_m$ 的次数是 $m$.
\item[(g)] 证明 $x \longmapsto \frac{1}{Z_m(v_0)}Z_m(x \cdot v_0)$ 是 $\SO(n)$ 上的球函数, 且每个球函数都具有此形式.
\end{enumerate}
多项式 $\frac{1}{f_m(1)}f_m$ 称为 Gegenbauer 多项式 (若 $n = 3$ 也称为 \index{Legendre 多项式}Legendre 多项式).
我们现在给出球函数的另一个公式.
\begin{enumerate}
\item[(h)] 考虑函数 $h_m \in V_m(\R^n)$, $h_m(z_1, \ldots, z_n) = (z_1 + iz_2)^m$. 证明 $h_m \in \mathcal{H}_m(\R^n)$.
\item[(i)] 定义函数 $\psi_m \colon S \to \CC$: $\psi_m(z) = \int_{\SO(n-1)}h_m(k \cdot z)\,\dd k$. 证明 $\psi_m$ 被 $\SO(n-1)$ 左不变, $\psi_m \in \mathcal{H}_m(\R^n)|_S$ 且 $\psi_m(v_0) = 1$.
\item[(j)] 证明 $\SO(n)$ 上的每个球函数都具有 $x \longmapsto \psi_m(x \cdot v_0)$ 的形式, 其中 $m \geq 0$ 唯一.
\end{enumerate}
我们可以计算定义 $\psi_m$ 的积分, 得到
\[
\psi_m(\cos\phi, z_2, \ldots, z_n) = \frac{\Gamma(\frac{n-1}{2})}{\sqrt{\pi}\Gamma(\frac{n-2}{2})}\int_0^{\pi}(\cos\phi + i\sin\phi\cos\theta)^m\sin^{n-3}\theta\,\dd\theta.
\]
\end{exercise}

\begin{solution}
(a). 若存在连续函数 $f \colon [-1,1] \to \CC$ 使得 $\phi(z_1, \ldots, z_n) = f(z_1)$, 则 $\phi$ 显然带状.
反之, 假设 $\phi$ 带状, 定义 $f \colon [-1,1] \to \CC$:
\[
f(z_1) = \phi(z_1, 0, \ldots, 0, \sqrt{1-z_1^2}).
\]
则 $f$ 显然连续. 设 $s = (z_1, \ldots, z_n) \in S$. 则存在 $x \in \SO(n-1)$ 使得 $x \cdot s = (z_1, 0, \ldots, 0, \sqrt{1-z_1^2})$ (用 $\SO(n-1)$ 在 $\R^{n-1}$ 的任何球面上传递作用这一事实). 由于 $\phi$ 带状, 有 $\phi(s) = f(z_1)$.
(b). 我们在 $\R^n$ 上用球坐标. 即给定 $(z_1, \ldots, z_n) \in \R^n$, 写
\begin{align*}
z_1 &= r\cos\varphi_1 \\
z_2 &= r\sin\varphi_1\cos\varphi_2 \\
&\vdots \\
z_{n-1} &= r\sin\varphi_1\sin\varphi_2\cdots\sin\varphi_{n-2}\cos\varphi_{n-1} \\
z_n &= r\sin\varphi_1\sin\varphi_2\cdots\sin\varphi_{n-2}\sin\varphi_{n-1},
\end{align*}
其中 $r = \sqrt{z_1^2 + \cdots + z_n^2} \in \R_{\geq 0}$, $\varphi_1, \ldots, \varphi_{n-2} \in [0, \pi]$, $\varphi_{n-1} \in [0, 2\pi)$ ($\varphi_1, \ldots, \varphi_{n-1}$ 一般不唯一确定, 但例如 $z_1, \ldots, z_n$ 都非零时唯一). 若 $dz$ 是 $\R^n$ 上的 Lebesgue 测度, 则
\[
dz = r^{n-1}\sin^{n-2}\varphi_1\sin^{n-3}\varphi_2\cdots\sin\varphi_{n-2}\,\dd r\,\dd\varphi_1\cdots\dd\varphi_{n-1}.
\]
设 $\phi \in \mathcal{C}(S)$ 带状. 差一个正实常数, $\int_S\phi(s)\,\dd\mu(s)$ 等于 $\int_{B-\{0\}}\psi(z)\,\dd z$, 其中 $B$ 是闭单位球, $\psi(z) = \phi(\frac{z}{\|z\|})$ ($z \neq 0$). 这等于 $\int_0^1 r^{n-1}\,\dd r$ (另一个正实常数) 与
\begin{multline*}
\int_0^{\pi}\cdots\int_0^{\pi}\int_0^{2\pi}\phi(\cos\varphi_1, \sin\varphi_1\cos\varphi_2, \ldots, \sin\varphi_1\cdots\sin\varphi_{n-2}\sin\varphi_{n-1}) \\
\times \sin^{n-2}\varphi_1\sin^{n-3}\varphi_2\cdots\sin\varphi_{n-2}\,\dd\varphi_1\cdots\dd\varphi_{n-1}
\end{multline*}
的乘积. 由于 $\phi$ 带状, 上面的大积分等于
\[
\int_0^{\pi}\cdots\int_0^{\pi}\int_0^{2\pi}f(\cos\varphi_1)\sin^{n-2}\varphi_1\sin^{n-3}\varphi_2\cdots\sin\varphi_{n-2}\,\dd\varphi_1\cdots\dd\varphi_{n-1}.
\]
差常数
\[
\int_0^{\pi}\cdots\int_0^{\pi}\int_0^{2\pi}\sin^{n-3}\varphi_2\cdots\sin\varphi_{n-2}\,\dd\varphi_2\cdots\dd\varphi_{n-1}
\]
(它必须为正, 因为它计算的是常值函数 $1$ 的积分, 差一个正常数), 这等于
\[
\int_0^{\pi}f(\cos\varphi_1)\sin^{n-2}\varphi_1\,\dd\varphi_1.
\]
最后用代换 $t = \cos\varphi_1$. 有 $dt = -\sin\varphi_1\,\dd\varphi_1$, 所以
\[
\int_0^{\pi}f(\cos\varphi_1)\sin^{n-2}\varphi_1\,\dd\varphi_1 = \int_1^{-1}f(t)(-\sqrt{1-t^2})^{n-3}\,\dd t = \int_{-1}^1 f(t)(1-t^2)^{(n-3)/2}\,\dd t.
\]
(c). 设 $y \in \SO(n-1)$. 由 $f_{v_0}$ 的定义, 对每个 $g \in \mathcal{H}_m(\R^n)$,
\[
\langle g, L_yf_{v_0} \rangle = \langle L_{y^{-1}}g, f_{v_0} \rangle = L_yg(v_0) = g(y^{-1}v_0) = g(v_0) = \langle g, f_{v_0} \rangle
\]
(因为 $\SO(n-1)$ 是 $v_0$ 在 $\SO(n)$ 中的稳定化子). 所以 $f_{v_0} - L_yf_{v_0}$ 与 $\mathcal{H}_m(\R^n)$ 的每个元素正交, 这蕴含 $f_{v_0} - L_yf_{v_0} = 0$, 即 $f_{v_0} = L_yf_{v_0}$. 所以 $Z_m$ 带状.
(d). 函数 $Z_m$ 是 $f_{v_0}$ 的限制, 而 $f_{v_0} \in \mathcal{H}_m(\R^n)$, 特别地是 $m$ 次多项式函数. 由于 $f_m$ 由 $f_m(t) = Z_m(t, 0, \ldots, 0, \sqrt{1-t^2})$ 定义, 可以把 $f_m$ 写成
\[
f_m(t) = \sum_{k=0}^m c_kt^k(1-t^2)^{(m-k)/2},
\]
其中 $c_0, \ldots, c_m \in \CC$. 但又有 $f_m(t) = Z_m(t, 0, \ldots, 0, -\sqrt{1-t^2})$, 所以
\[
\sum_{k=0}^m c_kt^k(1-t^2)^{(m-k)/2} = \sum_{k=0}^m(-1)^{m-k}c_kt^k(1-t^2)^{(m-k)/2}.
\]
这迫使除非 $m - k$ 偶, $c_k = 0$, 所以 $f(t)$ 确实是 $t$ 的次数 $\leq m$ 的多项式.
(e). 函数 $s \longmapsto Z_m(s)Z_{m'}(s)$ 带状, 所以
\[
\int_S Z_m(s)Z_{m'}(s)\,\dd\mu(s) = c\int_{-1}^1 f_m(t)f_{m'}(t)(1-t^2)^{(n-3)/2}\,\dd t
\]
由 (b). 由习题 \ref{ex:V.8.1.5}(b), 左端为 $0$.
(f). 对每个 $m \geq 0$, 函数 $Z_m$ 由定义非零 (且因为存在满足 $g(v_0) \neq 0$ 的 $f \in \mathcal{H}_m(\R^n)$, 见习题 \ref{ex:V.8.1.4}(h) 的证明), 所以由习题 \ref{ex:V.8.1.2}(h), $f_m \neq 0$. 由 (e), 函数 $(f_m)_{m \geq 0}$ 在 $L^2([-1,1], (1-t^2)^{(n-3)/2}\,\dd t)$ 中正交, 所以它们也线性无关. 固定 $m \geq 0$. 次数 $\leq m$ 的多项式空间 $P_m$ 维数为 $m + 1$ 且包含线性无关族 $(f_0, \ldots, f_m)$, 所以此族是 $P_m$ 的基. 但 $f_0, \ldots, f_{m-1}$ 的次数 $\leq m - 1$, 所以 $f_m$ 的次数必为 $m$.
(g). 由推论 \ref{cor:V.7.2}, 若 $\mathcal{Z}$ 是 $\SO(n)$ 上的球函数的集合, 则
\[
L^2(S) = \bigoplus_{\phi \in \mathcal{Z}}V_{\phi}
\]
且 $\phi$ 生成 $V_{\phi}^{\SO(n-1)}$. 所以用习题 \ref{ex:V.8.1.5}, 看到球函数正好是空间 $(\mathcal{H}_m(\R^n)|_S)^{\SO(n-1)}$ 的生成元中把 $v_0$ 送到 $1$ 的那些.
对每个 $m \geq 0$, 函数 $Z_m \in \mathcal{H}_m(\R^n)|_S$ 被 $\SO(n-3)$ 不变, 所以它生成 $\mathcal{H}_m(\R^n)|_S$ 中 $\SO(n-1)$-不变向量的空间, 且有倍数 $x \longmapsto \frac{1}{Z_m(v_0)}Z_m(x \cdot v_0)$ 是球函数. 由于球函数必把 $v_0$ 送到 $1$, 这个倍数就是球函数.
(h). 这由习题 \ref{ex:V.8.1.4}(k) 得出. 直接证明也很容易.
(i). 这与习题 \ref{ex:V.8.1.4}(h) 的证明中的构造完全相同 (取 $V = \mathcal{H}_m(\R^n)|_S$). 同样的证明表明 $\psi_m \in \mathcal{H}_m(\R^n)|_S$, $\psi_m$ 被 $\SO(n-1)$ 左不变且 $\psi_m(v_0) = f_m(v_0) = 1$.
(j). 与 (g) 的证明思路相同: 在每个 $\mathcal{H}_m(\R^n)|_S$ 中有一个球函数, 它是此空间中把 $v_0$ 送到 $1$ 的唯一 $\SO(n-1)$-不变元素. 由 (i), 函数 $\psi_m$ 满足所有要求的性质.
\end{solution}

\subsection{\texorpdfstring{Gelfand 对 $(S_n, S_r \times S_{n-r})$}{Gelfand 对 (S_n, S_r x S_{n-r})}}

这一系列习题的目标是研究 Gelfand 对 $(S_n, S_r \times S_{n-r})$, 摘录自 {[}7{]} 的第 5.1, 6.1 和 6.2 节. 我们按如下方式把 $S_r \times S_{n-r}$ 嵌入 $S_n$: 若 $\sigma \in S_r$, $\tau \in S_{n-r}$, 则 $\sigma \times \tau \in S_n$ 由
\[
(\sigma \times \tau)(i) = \begin{cases} \sigma(i) & \text{若 } 1 \leq i \leq r, \\ \tau(i - r) + r & \text{若 } r + 1 \leq i \leq n \end{cases}
\]
给出.
若 $E$ 是有限集, 记 $L(E)$ 为函数 $f \colon E \to \CC$ 的空间, 带 $L^2$ 内积 $\langle f, f' \rangle = \sum_{x \in E}f(x)f'(x)$.

\begin{exercise}[距离传递作用]
\label{ex:V.8.2.1}
在此题中, 固定传递作用 (左) 在集合 $E$ 上的有限群 $G$. 设 $x_0 \in E$, $K \subset G$ 是 $x_0$ 的稳定化子.
\begin{enumerate}
\item[(a)] 证明下列条件等价:
  \begin{enumerate}
  \item[(i)] 对所有 $x, y \in E$, 存在 $g \in G$ 使得 $g \cdot (x, y) = (y, x)$.
  \item[(ii)] 对每个 $g \in G$, 有 $g^{-1} \in KgK$.
  \end{enumerate}
\item[(b)] 若 (a) 的条件满足, 证明 $(G, K)$ 是 Gelfand 对.
\end{enumerate}
现在我们假设有度量 $d \colon E \times E \to \R_{\geq 0}$, 且群 $G$ 保距作用. 假设 $G$ 在 $E$ 上的作用是距离传递 (distance-transitive) 的, 即对所有满足 $d(x, y) = d(x', y')$ 的 $(x, y), (x', y') \in E \times E$, 存在 $g \in G$ 使得 $g \cdot (x, y) = (x', y')$.
\begin{enumerate}
\item[(c)] 证明 $(G, K)$ 是 Gelfand 对.
\item[(d)] 证明 $K$ 在 $E$ 上的轨道是球面 $\{x \in E \mid d(x, x_0) = j\}$, 其中 $j \in \R_{\geq 0}$.
\item[(e)] 设 $\Omega_r$ 是 $\{1, \ldots, n\}$ 的基数为 $r$ 的子集的集合. 证明公式 $d(A, B) = r - |A \cap B|$ 定义 $\Omega_r$ 上的度量.
\item[(f)] 证明 $(S_n, S_r \times S_{n-r})$ 是 Gelfand 对.
\end{enumerate}
\end{exercise}

\begin{solution}
(a). 假设 (i) 成立. 设 $g \in G$. 由 (i), 存在 $h \in G$ 使得 $h \cdot (K, gK) = (gK, K)$. 则第一个分量处的等式给出 $g^{-1}h \in K$, 第二个分量处的等式给出 $g^{-1}hg \in g^{-1}K$, 所以 $g \in (g^{-1}h)^{-1}g^{-1}K \subset Kg^{-1}K$.
假设 (ii) 成立. 设 $g, h \in G$. 由 (ii), 可以找到 $k_1, k_2 \in K$ 使得 $g^{-1}h = k_1h^{-1}gk_2$, 然后
\begin{align*}
(gK, hK) &= g \cdot (K, g^{-1}hK) \\
&= g \cdot (K, k_1h^{-1}gK) \\
&= gk_1^{-1} \cdot (K, h^{-1}gK) \\
&= gk_1^{-1}h^{-1} \cdot (hK, gK).
\end{align*}
(b). 这由命题 \ref{prop:V.2.5} 得出, 取 $\theta = \mathrm{id}_G$.
(c). 设 $x, y \in X$. 则 $d(y, x) = d(x, y)$, 所以存在 $g \in G$ 使得 $g \cdot (x, y) = (y, x)$. 换句话说, (a) 的条件 (i) 满足, 所以条件 (ii) 也满足. 由 (b), 这蕴含 $(G, K)$ 是 Gelfand 对.
(d). 对 $j \in \Z_{\geq 0}$ 写 $S_j = \{x \in X \mid d(x, x_0) = j\}$. 由于 $G$ 在 $X$ 上保距作用且 $K$ 固定 $x_0$, 集合 $S_j$ 被 $K$ 稳定. 为证明它们是 $K$ 在 $X$ 上的轨道, 需要证明 $K$ 在每个非空 $S_j$ 上传递作用. 所以设 $j \geq 0$, 假设有 $x, y \in S_j$. 则 $d(x_0, x) = d(x_0, y) = j$, 所以由假设, 存在 $g \in G$ 使得 $g \cdot (x_0, x) = (x_0, y)$. $g \cdot x_0 = x_0$ 这一事实蕴含 $g \in K$, 所以 $x$ 和 $y$ 在同一 $K$-轨道中.
(e). 显然对所有 $A, B \in \Omega_r$ 有 $d(A, B) = d(B, A)$. 设 $A, B, C \in \Omega_r$. 首先, 若 $d(A, B) = 0$, 则 $|A \cap B| = r = |A| = |B|$, 所以 $A \cap B = A$ 且 $A \cap B = B$, 从而 $A = B$. 证明三角不等式. 有
\[
|A \cap B| + |B \cap C| = |(A \cap B) \cup (B \cap C)| + |A \cap B \cap C| \leq |B| + |A \cap C| = r + |A \cap C|,
\]
所以
\[
d(A, C) = r - |A \cap C| \leq 2r - |A \cap B| - |B \cap C| = d(A, B) + d(B, C).
\]
(f). 让 $S_n$ 按 $\sigma \cdot A = \sigma(A)$ 作用于 $\Omega_r$. 这个作用传递, 且 $S_r \times S_{n-r}$ 是 $\{1, \ldots, r\}$ 的稳定化子. 又, 显然 $S_n$ 在 $\Omega_r$ 上保距作用. 所以由 (c), 只需验证作用是距离传递的. 设 $A, B, A', B' \in \Omega_r$ 满足 $d(A, B) = d(A', B')$, 即 $|A \cap B| = |A' \cap B'|$. 选取双射 $\phi \colon A \to A'$ 把 $A$ 的子集 $A \cap B$ 送到 $A' \cap B'$; 这由 $|A \cap B| = |A' \cap B'|$ 可能. 选取双射 $\psi \colon B - (A \cap B) \to B' - (A' \cap B')$; 这也可能, 因为 $|B - (A \cap B)| = r - |A \cap B| = |B' - (A' \cap B')|$. 把 $\phi$ 和 $\psi$ 放在一起给出双射 $A \cup B \xrightarrow{\sim} A' \cup B'$, 它把 $A$ 送到 $A'$, 把 $B$ 送到 $B'$, 且此双射的任何延拓 $\sigma \in S_n$ 将满足 $\sigma \cdot (A, B) = (A', B')$.
\end{solution}

\begin{exercise}[\index{距离正则图}距离正则图]
\label{ex:V.8.2.2}
设 $(\Omega_r, d)$ 是习题 \ref{ex:V.8.2.1}(e) 的有限度量空间. 设 $N$ 是 $\Omega_r$ 的\index{直径}直径, 即
\[
N = \max\{d(A, B) \mid A, B \in \Omega_r\}.
\]
对 $i \in \{0, \ldots, N\}$, 定义线性算子 $\Delta_i \colon L(\Omega_r) \to L(\Omega_r)$:
\[
\Delta_if(A) = \sum_{B \in \Omega_r, d(A,B)=i}f(B)
\]
对每个 $f \in L(\Omega_r)$. 又记 $A_f$ 为 $\End(L(\Omega_r))$ 中由 $\Delta_0, \ldots, \Delta_N$ 生成的子代数.
\begin{enumerate}
\item[(a)] 证明 $N = \min(r, n - r)$.
\item[(b)] 证明存在整数 $b_0, \ldots, b_N, c_0, c_1, \ldots, c_N$ 使得对每个 $i$ 和所有满足 $d(A, B) = i$ 的 $A, B \in \Omega_r$,
\[
|\{C \in \Omega_r \mid d(A, C) = 1 \text{ 且 } d(B, C) = i + 1\}| = b_i
\]
且
\[
|\{C \in \Omega_r \mid d(A, C) = 1 \text{ 且 } d(B, C) = i - 1\}| = c_i.
\]
(当然, $c_0 = 0$ 且 $c_1 = 1$.)\footnote{换句话说, 以 $\Omega_r$ 为顶点集、在满足 $d(A, B) = 1$ 的 $A, B \in \Omega_r$ 之间连边的图是距离正则的, 见 {[}7{]} 的定义 5.1.1.}
\item[(c)] 证明 $c_2, \ldots, c_N > 0$.
\item[(d)] 若 $i \in \{1, \ldots, N\}$, 证明
\[
\Delta_i\Delta_1 = b_{i-1}\Delta_{i-1} + (b_0 - b_i - c_i)\Delta_i + c_{i+1}\Delta_{i+1},
\]
约定 $\Delta_{N+1} = 0$.
\item[(e)] 证明存在多项式 $p_0, \ldots, p_N \in \R[t]$ 使得 $\deg(p_i) = i$ 且 $\Delta_i = p_i(\Delta_1)$.
\item[(f)] 证明 $A_f$ 是 $\End(L(\Omega_r))$ 中由 $\Delta_1$ 生成的子代数.
\item[(g)] 证明 $A_f$ 作为 $\CC$-向量空间由 $\Delta_0, \ldots, \Delta_N$ 张成.
\item[(h)] 证明 $\dim_{\CC}A_f = N + 1$.
\item[(i)] 证明 $L(\Omega_r)$ 的自同态 $\Delta_1$ 自伴.
\item[(j)] 证明有分解成两两正交子空间的直和 $L(\Omega_r) = \bigoplus_{i=0}^N V_i$, 其中 $V_0, \ldots, V_N$ 是 $\Delta_1$ 的特征空间. (提示: 习题 \ref{ex:II.5.2}.)
\end{enumerate}
\end{exercise}

\begin{solution}
(a). 注意对所有 $A, B \in \Omega_r$ 有 $d(A, B) \leq r$, 所以 $N \leq r$. 又, 对所有 $A, B \in \Omega_r$,
\[
d(A, B) = r - |A \cap B| = r - (|A| + |B| - |A \cup B|) = |A \cup B| - r \leq n - r,
\]
所以 $N \leq n - r$, 且 $N \leq \min(r, n - r)$.
现在取 $A = \{1, \ldots, r\}$, $B = \{n - r + 1, \ldots, n\}$. 则 $A, B \in \Omega_r$, 且 $|A \cap B| = \max(0, 2r - n)$. 所以 $N \geq d(A, B) = r - \max(0, 2r - n) = \min(r, n - r)$.
(b). 固定 $i \in \{0, \ldots, N\}$. 对所有满足 $d(A, B) = i$ 的 $A, B \in \Omega_r$, 设
\[
X_i(A, B) = \{C \in \Omega_r \mid d(A, C) = 1 \text{ 且 } d(B, C) = i + 1\}
\]
且
\[
Y_i(A, B) = \{C \in \Omega_r \mid d(A, C) = 1 \text{ 且 } d(B, C) = i - 1\}.
\]
若 $\sigma \in S_n$ 满足 $\sigma(A, B) = (A', B')$, 则 $\sigma$ 诱导双射 $X_i(A, B) \xrightarrow{\sim} X_i(A', B')$ 和 $Y_i(A, B) \xrightarrow{\sim} Y_i(A', B')$, 因为 $S_n$ 在 $\Omega_r$ 上保距作用. 所以结论由我们在习题 \ref{ex:V.8.2.1}(f) 的证明中证明的 $S_n$ 在 $\Omega_r$ 上的作用是距离传递的这一事实得出.
(c). 设 $i \in \{2, \ldots, N\}$. 取 $A = \{1, \ldots, r\}$, $B = \{i + 1, i + 2, \ldots, i + r\}$. (注意 $i + r \leq N + r \leq n$ 由 (a).) 需要证明存在至少一个 $C \in \Omega_r$ 使得 $d(A, C) = 1$ (即 $|A \cap C| = r - 1$) 且 $d(B, C) = i - 1$ (即 $|B \cap C| = r - i + 1$). 这对 $C = \{2, 3, \ldots, r + 1\}$ 成立.
(d). 设 $f \in L(\Omega_r)$, $A \in \Omega_r$. 则
\[
\Delta_i\Delta_1f(A) = \sum_{B \in \Omega_r, d(A,B)=i}\ \sum_{C \in \Omega_r, d(B,C)=1}f(C).
\]
设 $C \in \Omega_r$. 若存在 $B \in \Omega_r$ 使得 $d(A, B) = i$ 且 $d(B, C) = 1$, 则由三角不等式必有 $i - 1 \leq d(A, C) \leq i + 1$.
假设 $d(A, C) = i + 1$. 则
\[
\{B \in \Omega_r \mid d(A, B) = i \text{ 且 } d(B, C) = 1\} = Y_{i+1}(C, A)
\]
(用 (b) 的证明中的记号). 假设 $d(A, C) = i - 1$. 则
\[
\{B \in \Omega_r \mid d(A, B) = i \text{ 且 } d(B, C) = 1\} = X_{i-1}(C, A).
\]
最后, 假设 $d(A, C) = i$. 考虑集合
\begin{multline*}
\{B \in \Omega_r \mid d(A, B) = i \text{ 且 } d(B, C) = 1\} \cup \{B \in \Omega_r \mid d(A, B) = i + 1 \text{ 且 } d(B, C) = 1\} \\
\cup \{B \in \Omega_r \mid d(A, B) = i - 1 \text{ 且 } d(B, C) = 1\}.
\end{multline*}
这个并显然不相交. 要计算第一个集合的基数, 第二个集合是 $X_i(C, A)$, 第三个集合是 $Y_i(C, A)$. 又, 由三角不等式, 这个并正好是
\[
\{B \in \Omega_r \mid d(B, C) = 1\} = X_0(C, C).
\]
所以得到
\[
|\{B \in \Omega_r \mid d(A, B) = i \text{ 且 } d(B, C) = 1\}| + b_i + c_i = b_0.
\]
最后看到
\begin{align*}
\Delta_i\Delta_1f(A) &= c_{i+1}\sum_{C, d(A,C)=i+1}f(C) + b_{i-1}\sum_{C, d(A,C)=i-1}f(C) + (b_0 - b_i - c_i)\sum_{C, d(A,C)=i}f(C) \\
&= c_{i+1}\Delta_{i+1}f(A) + b_{i-1}\Delta_{i-1}f(A) + (b_0 - b_i - c_i)\Delta_if(A).
\end{align*}
(e). 对 $i$ 归纳证明此结论. $i = 0$ 时显然 (注意 $\Delta_0 = \mathrm{id}$, 所以取 $p_0 = 1$), $i = 1$ 时也显然 (取 $p_1(t) = t$). 假设对某个 $i \geq 1$ 已知结果, 证明对 $i + 1$. 由 (c) 和 (d),
\[
\Delta_{i+1} = c_{i+1}^{-1}(\Delta_i\Delta_1 - b_{i-1}\Delta_{i-1} - (b_0 - b_i - c_i)\Delta_i),
\]
所以 $\Delta_{i+1} = p_{i+1}(\Delta_1)$, 其中
\[
p_{i+1}(t) = c_{i+1}^{-1}(tp_i(t) - p_{i-1}(t) - (b_0 - b_i - c_i)p_i(t)).
\]
显然也有 $\deg(p_{i+1}(t)) = i + 1$.
(f). 设 $A'_f$ 是 $\End(L(\Omega_r))$ 中由 $\Delta_1$ 生成的子代数. 则 $A'_f \subset A_f$ 由 $A_f$ 的定义. 由 (e), 有 $\Delta_0, \ldots, \Delta_N \in A'_f$, 所以 $A_f \subset A'_f$.
(g). 对 $i \geq 0$ 归纳证明 $\Delta_1^i \in \operatorname{Span}(\Delta_0, \Delta_1, \ldots, \Delta_i)$. (结论将由 (f) 得出.) $i = 0$ 和 $i = 1$ 时显然. 假设对 $i \geq 1$ 成立, 证明对 $i + 1$. 由 (e), 存在非零 $a \in \R$ 和 $a_0, \ldots, a_i \in \R$ 使得 $\Delta_{i+1} = a\Delta_1^{i+1} + \sum_{j=0}^i a_j\Delta_1^j$. 由于由归纳假设对每个 $j \in \{0, \ldots, i\}$ 有 $\Delta_1^j \in \operatorname{Span}(\Delta_0, \ldots, \Delta_j)$, 推出 $\Delta_1^{i+1} \in \operatorname{Span}(\Delta_0, \ldots, \Delta_{i+1})$.
(h). 知道 $A_f = \operatorname{Span}(\Delta_0, \ldots, \Delta_N)$, 所以只需证明族 $(\Delta_0, \ldots, \Delta_N)$ 线性无关. 设 $c_0, \ldots, c_N \in \CC$. 若 $A, B \in \Omega_r$, 记 $\delta_A$ 为 $\{A\}$ 的特征函数, 则 $d(A, B) = i$ 时 $\Delta_i\delta_A(B) \neq 0$, 且
\[
\sum_{i=0}^N c_i\Delta_i\delta_A(B) = c_{d(A,B)}.
\]
由于存在满足 $d(A, B) = i$ (对每个 $i \in \{0, \ldots, N\}$) 的偶对 $(A, B) \in \Omega_r^2$, 若 $\sum_{i=0}^N c_i\Delta_i = 0$ 则 $c_0 = \cdots = c_N = 0$.
(i). 设 $f, g \in L(\Omega_r)$. 则
\begin{align*}
\langle f, \Delta_1g \rangle &= \sum_{A \in \Omega_r}f(A)(\Delta_1g)(A) \\
&= \sum_{A \in \Omega_r}\ \sum_{B \in \Omega_r, d(A,B)=1}f(A)g(B) \\
&= \sum_{A,B \in \Omega_r, d(A,B)=1}f(A)g(B) \\
&= \sum_{B \in \Omega_r}(\Delta_1f)(B)g(B) \\
&= \langle \Delta_1f, g \rangle.
\end{align*}
(j). 由于 $\Delta_1$ 自伴, 谱定理说它可对角化且其特征空间两两正交. 所以只需证明有 $N + 1$ 个特征空间. 由习题 \ref{ex:II.5.2}(b), 知道 $\End(L(\Omega_r))$ 中由 $\Delta_1$ 生成的子代数, 即 $A_f$ (见 (f)), 约化. 由习题 \ref{ex:II.5.2}(c), 知道 $A_f$ (即 $\Delta_1$) 的特征空间数是 $\dim(A_f)$, 由 (h) 知道 $\dim(A_f) = N + 1$.
\end{solution}

\begin{exercise}[对称幂的分解]
\label{ex:V.8.2.3}
用习题 \ref{ex:V.8.2.2} 的记号. 注意有 $G := S_n$ 在 $\Omega_r$ 上的作用, 且 $\{1, \ldots, r\}$ 的稳定化子是 $K := S_r \times S_{n-r}$. 设 $M^{n-r,r} = L(\Omega_r)$, 看作 $S_n$ 的表示, 通过拟正则表示 (即若 $g \in G$, $f \in M^{n-r,r}$, $A \in \Omega_r$, 有 $(g \cdot f)(A) = f(g^{-1}A)$).
定义 $d \colon M^{n-r,r} \to M^{n-r+1,r-1}$ 和 $d^* \colon M^{n-r+1,r-1} \to M^{n-r,r}$:
\[
(df)(A) = \sum_{B \in \Omega_r, A \subset B}f(B)
\]
且
\[
(d^*f)(B) = \sum_{A \in \Omega_r, A \subset B}f(A).
\]
(若 $r = 0$, 取 $d = 0$ 和 $d^* = 0$.)
又记 $\Delta$ 为习题 \ref{ex:V.8.2.2} 的算子 $\Delta_1$; 即对每个 $f \in M^{n-r,r}$, 函数 $\Delta f \in M^{n-r,r}$ 由
\[
(\Delta f)(A) = \sum_{B \in \Omega_r, d(A,B)=1}f(B)
\]
定义.\footnote{注意函数 $d$, $d^*$ 和 $\Delta$ 对每个 $r$ 都有定义; 记号中不标出 $r$, 它应该由上下文清楚. 这与\index{离散 Laplace 算子}离散 Laplace 算子密切相关. 事实上, $M^{n-r,r}$ 上的离散 Laplace 算子最常见的定义是 $\frac{1}{r}\Delta - \mathrm{id}$.}
最后, 若 $a \in \CC$, $i \in \Z_{\geq 0}$, 记
\[
(a)_i = a(a+1)\cdots(a+i-1).
\]
例如有 $(1)_n = n!$ 和 $\binom{n}{k} = \frac{(n-k+1)_k}{k!}$.
\begin{enumerate}
\item[(a)] 证明 $A \longmapsto \{1, \ldots, n\} - A$ 诱导 $G$-等变同构 $M^{r,n-r} \xrightarrow{\sim} M^{n-r,r}$.
\item[(b)] 证明 $d$ 和 $\Delta$ 是 $G$-等变的.
\item[(c)] 证明 $d^*$ 是 $d$ 的伴随.
\item[(d)] 若 $f \in M^{n-r,r}$, 证明
\[
dd^*f = \Delta f + (n-r)f \quad \text{且} \quad d^*df = \Delta f + rf.
\]
\item[(e)] 设 $f \in M^{n-r,r}$, $1 \leq p \leq q \leq n - r$. 证明
\[
d(d^*)^qf = (d^*)^qdf + q(n - 2r - q + 1)(d^*)^{q-1}f.
\]
若此外 $df = 0$, 证明
\[
d^p(d^*)^qf = (q - p + 1)_p(n - 2r - q + 1)_p(d^*)^{q-p}f.
\]
\end{enumerate}
假设 $0 \leq r \leq n/2$. 若 $r > 0$, 设 $S^{n-r,r} = \Ker(d \colon M^{n-r,r} \to M^{n-r+1,r-1})$; 若 $r = 0$, 设 $S^{n-r,r} = M^{n-r,r}$. 这是 $M^{n-r,r}$ 的 $G$-稳定子空间.
\begin{enumerate}
\item[(f)] 若 $0 \leq m \leq n$ 且 $0 \leq r \leq \min(m, n - m)$, 证明 $(d^*)^{m-r} \colon S^{n-r,r} \to M^{n-m,m}$ 单射. (提示: 计算 $\|(d^*)^{m-r}f\|_2^2$.)
\item[(g)] 在 (f) 的假设下, 证明 $(d^*)^{m-r}(S^{n-r,r})$ 包含于 $\Delta$ 的特征值 $m(n-m) - r(n-r+1)$ 的特征空间中.
\item[(h)] 证明若 $1 \leq m \leq n/2$, 则 $M^{n-m,m}$ 中 $S^{n-m,m}$ 的正交补是 $d^*(M^{n-m+1,m-1})$.
\item[(i)] 证明对满足 $0 \leq r \leq n/2$ 的每个 $r$, $S^{n-r,r} \neq 0$.
\item[(j)] 若 $0 \leq m \leq n$, 证明
\[
M^{n-m,m} = \bigoplus_{r=0}^{\min(m,n-m)}(d^*)^{m-r}(S^{n-r,r}),
\]
其中和项两两正交且正好是 $\Delta$ 的特征空间.
\item[(k)] 证明若 $r > 0$, $\dim_{\CC}(S^{n-r,r}) = \binom{n}{r} - \binom{n}{r-1}$.
\item[(l)] 证明表示 $S^{n-r,r}$ ($0 \leq r \leq n/2$) 不可约且两两不等价. (提示: $M^{m,n-m}$ 有多少个不可约成分?)
\end{enumerate}
\end{exercise}

\begin{solution}
(a). 映射 $A \longmapsto \{1, \ldots, n\} - A$ 是从 $\Omega_r$ 到 $\Omega_{n-r}$ 的双射, 且与 $G$ 在这两个集合上的作用交换. 结论立即得出.
(b). 设 $f \in M^{n-r,r}$, $\sigma \in S_n$. 若 $A \in \Omega_{r-1}$, 则
\begin{align*}
(dL_\sigma f)(A) &= \sum_{B \in \Omega_r, B \supset A}f(\sigma^{-1}(B)) \\
&= \sum_{B' \in \Omega_r, B' \supset \sigma^{-1}(A)}f(B') \\
&= (df)(\sigma^{-1}(A)) \\
&= L_\sigma(df)(A).
\end{align*}
若 $A \in \Omega_r$, 则
\begin{align*}
(\Delta L_\sigma f)(A) &= \sum_{B \in \Omega_r, d(A,B)=1}f(\sigma^{-1}(B)) \\
&= \sum_{B' \in \Omega_r, d(\sigma^{-1}(A),B')=1}f(B') \\
&= \Delta f(\sigma^{-1}(A)) \\
&= L_\sigma(\Delta f)(A).
\end{align*}
(c). 设 $f \in M^{n-r,r}$, $g \in M^{n-r+1,r-1}$. 则
\begin{align*}
\langle df, g \rangle &= \sum_{A \in \Omega_{r-1}}df(A)g(A) \\
&= \sum_{A \in \Omega_{r-1}, B \in \Omega_r, A \subset B}f(B)g(A) \\
&= \sum_{B \in \Omega_r}f(B)d^*g(B).
\end{align*}
(d). 设 $f \in M^{n-r,r}$, $A \in \Omega_r$. 则
\begin{align*}
dd^*f(A) &= \sum_{B \in \Omega_{r+1}, B \supset A}d^*f(B) \\
&= \sum_{B \in \Omega_{r+1}, C \in \Omega_r, C \subset B \supset A}f(C).
\end{align*}
设 $C \in \Omega_r$. 若存在 $B \in \Omega_{r+1}$ 使得 $C \subset B \supset A$, 则 $d(A, C) \leq 1$. 若 $C = A$, 则
\[
|\{B \in \Omega_{r+1} \mid C \subset B \supset A\}| = |\{B \in \Omega_{r+1} \mid A \subset B\}| = n - r.
\]
若 $d(A, C) = 1$, 则同时包含 $A$ 和 $C$ 的 $\Omega_{r+1}$ 的唯一元素是 $A \cup C$. 最后得到
\begin{align*}
dd^*f(A) &= (n - r)f(A) + \sum_{C \in \Omega_r, d(A,C)=1}f(C) \\
&= (n - r)f(A) + \Delta f(A).
\end{align*}
类似地,
\begin{align*}
d^*df(A) &= \sum_{B \in \Omega_{r-1}, B \subset A}df(B) \\
&= \sum_{C \in \Omega_r, B \in \Omega_{r-1}, C \supset B \subset A}f(C).
\end{align*}
设 $C \in \Omega_r$. 若存在 $B \in \Omega_{r-1}$ 使得 $C \supset B \subset A$, 则 $d(A, C) \leq 1$. 若 $A = C$, 则
\[
|\{B \in \Omega_{r-1} \mid C \supset B \subset A\}| = r.
\]
若 $d(A, C) = 1$, 则同时包含于 $A$ 和 $C$ 的唯一 $B \in \Omega_{r-1}$ 是 $B = A \cap C$. 所以得到
\begin{align*}
d^*df(A) &= rf(A) + \sum_{C \in \Omega_r, d(A,C)=1}f(C) \\
&= rf(A) + \Delta f(A).
\end{align*}
(e). 对 $q$ 归纳证明第一个恒等式. 若 $q = 1$, 由 (d),
\[
dd^*f = \Delta f + (n-r)f = \Delta f + rf + (n-2r)f = d^*df + q(n-2r-q+1)(d^*)^{q-1}f
\]
对每个 $f \in M^{n-r,r}$. 现在假设对 $q \in \{1, \ldots, n-r-1\}$, 每个 $s$ 和 $M^{n-s,s}$ 的每个元素已知恒等式, 证明对 $q+1$. 若 $f \in M^{n-r,r}$, 有
\begin{align*}
d(d^*)^{q+1}f &= d(d^*)^q(d^*f) \\
&= (d^*)^qd(d^*f) + q(n-2(r+1)-q+1)(d^*)^{q-1}(d^*f) \quad \text{(对 } d^*f \in M^{n-r-1,r+1} \text{ 用归纳假设)} \\
&= (d^*)^q(d^*df + (n-2r)f) + q(n-2r-q-1)(d^*)^qf \quad \text{(用 } q=1 \text{ 的情形)} \\
&= (d^*)^{q+1}df + (n-2r + q(n-2r-q-1))(d^*)^qf \\
&= (d^*)^{q+1}df + (q+1)(n-2r-(q+1)+1)(d^*)^qf.
\end{align*}
现在对 $p$ 归纳证明第二个恒等式. 若 $p = 1$, 它归结为第一个恒等式 (用 $df = 0$). 假设对某个 $p \in \{1, \ldots, q-1\}$ (和所有 $s$, 所有满足 $df = 0$ 的 $f \in M^{n-s,s}$) 已证明, 证明对 $p+1$. 设 $f \in M^{n-r,r}$ 满足 $df = 0$. 则
\begin{align*}
d^{p+1}(d^*)^qf &= d(d^p(d^*)^qf) \\
&= d((q-p+1)_p(n-2r-q+1)_p(d^*)^{q-p}f) \\
&= (q-p+1)_p(n-2r-q+1)_p(q-p)(n-2r-(q-p)+1)(d^*)^{q-p-1}f \\
&\qquad \text{(用第一个恒等式和 } df = 0\text{)} \\
&= (q-p)_{p+1}(n-2r-q+1)_{p+1}(d^*)^{q-p-1}f.
\end{align*}
(f). 设 $f \in S^{n-r,r}$, $f \neq 0$. 用 (c) 然后用 (e) 的第二个恒等式, 看到
\begin{align*}
\langle (d^*)^{m-r}f, (d^*)^{m-r}f \rangle &= \langle d^{m-r}(d^*)^{m-r}f, f \rangle \\
&= (1)_{m-r}(n-2r-(m-r)+1)_{m-r}\langle f, f \rangle \\
&\neq 0,
\end{align*}
所以 $(d^*)^{m-r}f \neq 0$.
(g). 设 $f \in S^{n-r,r}$. 用 (d) 的第二个公式在 $M^{n-m,m}$ 上计算 $\Delta$, 用 (e) 的第二个公式 (取 $p = 1$), 得到
\begin{align*}
\Delta((d^*)^{m-r}f) &= d(d^*)^{m-r+1}f - (n-m)(d^*)^{m-r}f \\
&= (m-r+1)(n-2r-(m-r+1)+1)(d^*)^{m-r}f - (n-m)(d^*)^{m-r}f \\
&= (m(n-m) - r(n-r+1))(d^*)^{m-r}f.
\end{align*}
(h). 这是 $S^{n-m,m}$ 的定义和 (c) 的直接推论.
(i). 空间 $S^{n-r,r}$ 是 $d \colon M^{n-r,r} \to M^{n-r+1,r-1}$ 的核, 且 $\dim(M^{n-r+1,r-1}) = \binom{n}{r-1} < \binom{n}{r} = \dim(M^{n-r,r})$ (因为 $r \leq n/2$), 所以 $d$ 不能单射.
(j). 子空间 $(d^*)^{m-r}(S^{n-r,r})$ ($0 \leq r \leq \min(m, n-m)$) 由 (g) 包含于 $\Delta$ 的不同特征空间. 它们都非零由 (f) 和 (i). 知道 $\Delta$ 自伴 (由习题 \ref{ex:V.8.2.2}(i)), 所以这些空间两两正交. 又知道 $\Delta \in \End(M^{n-m,m})$ 恰有 $1 + \min(m, n-m)$ 个特征值 (由习题 \ref{ex:V.8.2.2}(a) 和 \ref{ex:V.8.2.2}(j)), 所以这些特征值必是 $m(n-m) - r(n-r+1)$, $0 \leq r \leq \min(m, n-m)$. 还需证明
\[
M^{n-m,m} = \bigoplus_{r=0}^{\min(m,n-m)}(d^*)^{m-r}(S^{n-r,r}).
\]
对 $m$ 归纳证明. $m = 0$ 时显然. 假设对 $m - 1$ (其中 $n/2 \geq m \geq 1$) 已知结果, 证明对 $m$. 由 (h),
\[
M^{n-m,m} = S^{n-m,m} \oplus d^*(M^{n-m+1,m-1}).
\]
由归纳假设,
\[
M^{n-m+1,m-1} = \bigoplus_{r=0}^{m-1}(d^*)^{m-1-r}(S^{n-r,r}).
\]
$m$ 的结果由这两个事实立即得出.
最后处理 $m \geq n/2$ 的情形. 设 $m' = n - m$. 已经看到
\[
M^{n-m',m'} = \bigoplus_{r=0}^{m'}(d^*)^{m'-r}(S^{n-r,r}).
\]
由 (k), 这蕴含 $\dim(M^{n-m',m'}) = \sum_{r=0}^{m'}\dim(S^{n-r,r})$. 又看到
\[
M^{n-m,m} \supset \bigoplus_{r=0}^{m'}(d^*)^{m-r}(S^{n-r,r}),
\]
且再次由 (k), 这蕴含 $\dim(M^{n-m,m}) \geq \sum_{r=0}^{m'}\dim(S^{n-r,r}) = \dim(M^{m',n-m'})$. 但 $\dim(M^{n-m,m}) = \dim(M^{n-m',m'})$ (由 (a)), 所以上面的不等式是等式, 且
\[
M^{n-m,m} = \bigoplus_{r=0}^{m'}(d^*)^{m-r}(S^{n-r,r}).
\]
(k). 由 (j) 和 (f), 映射 $d^* \colon M^{n-r+1,r-1} \to M^{n-r,r}$ 单射. 由 (h), 这蕴含
\[
\dim(S^{n-r,r}) = \dim(M^{n-r,r}) - \dim(M^{n-r+1,r-1}) = \binom{n}{r} - \binom{n}{r-1}.
\]
(l). 设 $m = \lfloor n/2 \rfloor$. 由于映射 $d$ 和 $d^*$ 是 $S_n$-等变的 (对 $d$ 见 (b), $d^*$ 是 $d$ 的伴随, 所以也等变), 子空间 $S^{n-r,r} \subset M^{n-r,r}$ 对每个 $r \leq n/2$ 是 $S_n$-稳定的, (j) 的分解是 $S_n$-子空间的分解. 其次, 由习题 \ref{ex:V.8.2.1}(f) 知道 $(S_n, S_m \times S_{n-m})$ 是 Gelfand 对, 所以由定理 \ref{thm:V.3.2.4}, 对应的拟正则表示, 即 $M^{n-m,r}$, 分解成互不相同的不可约表示的直和. 由推论 \ref{cor:V.7.2}, $M^{n-m,m}$ 中不可约和项的数是 Gelfand 对的球函数数, 它等于 $S_n$ 上双不变函数的空间维数 (因为球函数构成这些双不变函数的基, 由同一推论的 (iii)), 即 $(S_m \times S_{n-m}) \backslash S_n /(S_m \times S_{n-m})$ 的基数, 这也等于 $S_m \times S_{n-m}$ 在 $\Omega_m$ 上的轨道数. 但由习题 \ref{ex:V.8.2.1}(d) 看到, $S_m \times S_{n-m}$ 在 $\Omega_m$ 上的轨道是以 $A_0 := \{1, \ldots, m\}$ 为中心的球面. 这些球面的可能半径是 $0, 1, \ldots, \min(m, n-m) = m$ (由习题 \ref{ex:V.8.2.2}(a)), 且容易看到所有球面非空 (在习题 \ref{ex:V.8.2.2}(f) 的证明中已经用过). 最后得到 $M^{n-m,m}$ 的不可约成分数是 $m + 1$. 由于 (j) 的分解是 $M^{n-m,m}$ 分解成 $m + 1$ 个非零子表示, 它必是其不可约成分的分解, 所以得到结论. (注意 $(d^*)^{m-r}(S^{n-r,r})$ 作为 $S_n$ 的表示与 $S^{n-r,r}$ 等价, 由 (f).)
\end{solution}

\begin{exercise}[置换群 Gelfand 对的球函数]
\label{ex:V.8.2.4}
我们保留习题 \ref{ex:V.8.2.3} 的记号. 若 $m, h \in \{0, \ldots, n\}$, $A \in \Omega_m$, $\max(0, h-m) \leq \ell \leq \min(n-m, h)$, 记 $\sigma_{\ell,h-\ell}(A) \in L(\Omega_h)$ 为集合 $\{C \in \Omega_h \mid |A \cap C| = h - \ell\}$ 的特征函数. 又写 $\sigma_{-1,h+1}(A) = 0$. 固定 $A \in \Omega_m$.
\begin{enumerate}
\item[(a)] 若 $h = m$, 证明对每个 $\ell \in \{0, \ldots, \min(m, n-m)\}$, 函数 $\sigma_{\ell,m-\ell}(A)$ 是球面 $\{C \in \Omega_m \mid d(A, C) = \ell\}$ 的特征函数.
\item[(b)] 证明
\[
d(\sigma_{\ell,h-\ell}(A)) = (n - m - \ell + 1)\sigma_{\ell-1,h-\ell}(A) + (m - h + \ell + 1)\sigma_{\ell,h-\ell-1}(A).
\]
\item[(c)] 若 $k \leq h$ 且 $\max(0, k-m) \leq i \leq \min(k, n-m)$, 证明
\[
\frac{1}{(h-k)!}(d^*)^{h-k}\sigma_{i,k-i}(A) = \sum_{\ell = \max(i,h-m)}^{\min(h-k+i,n-m)}\binom{\ell}{i}\binom{h-\ell}{k-i}\sigma_{\ell,h-\ell}(A).
\]
\end{enumerate}
从现在起, 取 $A = \{1, \ldots, m\}$.
\begin{enumerate}
\item[(a)] 若 $0 \leq h \leq \min(m, n-m)$, 证明 $M^{n-h,h}$ 中 $S_m \times S_{n-m}$-不变向量的空间由函数 $\sigma_{\ell,h-\ell}(A)$, $0 \leq \ell \leq h$, 张成.
\item[(b)] 若 $0 \leq h \leq \min(m, n-m)$, 证明 $S^{n-h,h}$ 中 $S_m \times S_{n-m}$-不变向量的空间由函数
\[
\sum_{\ell=0}^h\frac{(n - m - h + 1)_{h-\ell}}{(-m)_{h-\ell}}\sigma_{\ell,h-\ell}(A)
\]
张成.
\item[(c)] 对 $0 \leq h \leq \min(m, n-m)$, 设 $\phi_h \in M^{n-m,m}$ 是 (d) 的和项 $(d^*)^{m-h}(S^{n-h,h})$ 中包含的唯一球函数. 证明
\[
\phi_h = \sum_{\ell=0}^{\min(m,n-m)}\phi(n, m, h; \ell)\sigma_{\ell,m-\ell}(A),
\]
其中
\[
\phi(n, m, h; \ell) = (-1)^h\binom{n-m}{h}^{-1}\sum_{i=\max(0,\ell-m+h)}^{\min(\ell,h)}\binom{m-\ell}{h-i}\binom{\ell}{i}\frac{(n-m-h+1)_{h-i}}{(-m)_{h-i}}.
\]
\item[(d)] 固定 $h$ 使得 $0 \leq h \leq \min(m, n-m)$. 证明 $\phi_h$ 中 $\sigma_{1,m-1}(A)$ 的系数是 $1 - \frac{h(n-h+1)}{m(n-m)}$. (注记: 有一种用最少计算解决此题的方法.)
\end{enumerate}
\end{exercise}

\begin{solution}
(a). 函数 $\sigma_{\ell,m-\ell}(A)$ 是集合 $\{C \in \Omega_m \mid |A \cap C| = m - \ell\}$ 的特征函数. 由于 $d(A, C) = m - |A \cap C|$, 这个集合正好是以 $A$ 为心、$\ell$ 为半径的球面.
(b). 设 $B \in \Omega_{h-1}$. 有
\[
d(\sigma_{\ell,h-\ell}(A))(B) = \sum_{C \in \Omega_h, C \supset B}\sigma_{\ell,h-\ell}(A)(C) = |\{C \in \Omega_h \mid C \supset B \text{ 且 } |A \cap C| = h - \ell\}|.
\]
若存在至少一个 $C \in \Omega_h$ 使得 $C \supset B$ 且 $|A \cap C| = h - \ell$, 则 $A \cap C \supset A \cap B$, 且这两个集合至多相差一个元素, 所以 $|A \cap B| \in \{h - \ell, h - \ell - 1\}$.
假设 $|A \cap B| = h - \ell$. 则对每个满足 $C \supset B$ 且 $|A \cap C| = h - \ell$ 的 $C \in \Omega_h$, 必有 $A \cap C = A \cap B$. 通过向 $B$ 添加 $\{1, \ldots, n\} - (A \cup B)$ 的一个元素得到每个这样的 $C$, 所以 $C$ 的可能数是
\[
n - |A \cup B| = n - (|A| + |B| - |A \cap B|) = n - (m + h - 1 - (h - \ell)) = n - m - \ell + 1.
\]
假设 $|A \cap B| = h - \ell + 1$. 则对每个满足 $C \supset B$ 且 $|A \cap C| = h - \ell$ 的 $C \in \Omega_h$, $C - B$ 的唯一元素必须是 $A \cap C - A \cap B$ 的元素. 所以 $C$ 的可能数是
\[
|A - A \cap B| = m - h + \ell - 1.
\]
最后得到
\[
d(\sigma_{\ell,h-\ell}(A))(B) = (n - m - \ell + 1)\sigma_{\ell-1,h-\ell}(B) + (m - h + \ell + 1)\sigma_{\ell,h-\ell-1}(B),
\]
如所愿.
(c). 对每个 $i$ 和每个 $D \in \Omega_i$, 记 $\delta_D \in L(\Omega_i)$ 为 $\{D\}$ 的特征函数. 设 $S$ 是集合 $\{C \in \Omega_k \mid |A \cap C| = k - i\}$. 则 $\sigma_{i,k-i}(A)$ 是 $S$ 的特征函数. 若 $C \in S$, 则对每个 $D \in \Omega_h$,
\[
(d^*)^{h-k}\delta_C(D) = \sum_{C \subset D_1 \subset \cdots \subset D_{h-k-1} \subset D_{h-k} = D, D_i \in \Omega_{k+i}}1.
\]
若 $C \not\subset D$, 集合 $\{C \subset D_1 \subset \cdots \subset D_{h-k-1} \subset D_{h-k} = D \mid D_i \in \Omega_{k+i}\}$ 空; 若 $C \subset D$, 这个集合有 $(h-k)!$ 个元素. 所以看到
\[
(d^*)^{h-k}\delta_C = (h-k)!\sum_{C \subset D \in \Omega_h}\delta_D.
\]
所以若 $D \in \Omega_h$, $\frac{1}{(h-k)!}(d^*)^{h-k}\sigma_{i,k-i}(A)$ 中 $\delta_D$ 的系数是集合 $\{C \in S \mid C \subset D\}$ 的基数. 写 $|A \cap D| = h - \ell$, 其中 $0 \leq \ell \leq h$; 注意有
\[
h - \ell = |A \cap D| = |A| + |D| - |A \cup D| \geq m + h - |A \cup D|
\]
且 $h \leq |A \cup D| \leq n$, 所以 $h - m \leq \ell \leq n - m$, 且此范围内的所有非负 $\ell$ 都能出现. 通过从 $A \cap D$ 中移除 $(h - \ell) - (k - i)$ 个元素、从 $D - (A \cap D)$ 中移除 $\ell - i$ 个元素, 得到满足 $C \subset D$ 的 $C \in S$. 这只有在 $h - \ell - (k - i) \geq 0$ 且 $0 \leq \ell - i$ 时才可能, 且有 $\binom{h-\ell}{(h-\ell)-(k-i)}\binom{\ell}{\ell-i} = \binom{h-\ell}{k-i}\binom{\ell}{i}$ 种不同的可能选择. 把所有放在一起, 看到若 $|A \cap D| = h - \ell$, 则 $\frac{1}{(h-k)!}(d^*)^{h-k}\sigma_{i,k-i}(A)$ 中 $\delta_D$ 的系数在 $\max(h-m, i) \leq \ell \leq \min(n-m, h-k+i)$ 时为 $\binom{h-\ell}{k-i}\binom{\ell}{i}$, 否则为 $0$. 这给出结果.
(d). 此结论等价于: $S_m \times S_{n-m}$ 在 $\Omega_h$ 中的轨道是 $S_{\ell} := \{C \in \Omega_h \mid |A \cap C| = h - \ell\}$, $0 \leq \ell \leq h$. 证明这个事实. 由于 $S_m \times S_{n-m}$ 固定 $A$, 集合 $S_{\ell}$ 被 $S_m \times S_{n-m}$ 不变, 所以只需证明 $S_m \times S_{n-m}$ 在这些集合上传递作用. 固定 $\ell \in \{0, \ldots, h\}$, 取 $C, C' \in S_{\ell}$. 由于 $|A \cap C| = |A \cap C'|$, 可以找到 $\sigma \in S_m$ 把 $A \cap C$ 送到 $A \cap C'$. 又, $|\{m+1, \ldots, n\} \cap C| = |\{m+1, \ldots, n\} \cap C'| = \ell$, 所以可以找到 $\tau \in S_{n-m}$ 把 $\{m+1, \ldots, n\} \cap C$ 送到 $\{m+1, \ldots, n\} \cap C'$. 则 $\sigma \times \tau \in S_m \times S_{n-m}$ 把 $C$ 送到 $C'$.
(e). 设 $f \in S^{n-h,h}$ 是 $S_m \times S_{n-m}$-不变向量. 由 (d), 不变性条件等价于可以把 $f$ 写成 $f = \sum_{\ell=0}^h a_{\ell}\sigma_{\ell,h-\ell}(A)$, 其中 $a_0, \ldots, a_h \in \CC$. $f \in S^{n-h,h}$ 这一事实意味着 $df = 0$. 用 (b), 可以把此条件改写成
\begin{align*}
0 &= \sum_{\ell=0}^h a_{\ell}((n - m - \ell + 1)\sigma_{\ell-1,h-\ell}(A) + (m - h + \ell + 1)\sigma_{\ell,h-1-\ell}(A)) \\
&= \sum_{\ell=0}^{h-1}a_{\ell+1}(n - m - \ell)\sigma_{\ell,h-\ell+1}(A) + \sum_{\ell=0}^h a_{\ell}(m - h + \ell + 1)\sigma_{\ell,h-1-\ell}(A).
\end{align*}
由于函数 $\sigma_{\ell,h-1-\ell}(A)$, $0 \leq \ell \leq h - 1$, 线性无关 (因为它们有不相交的支撑), 此等式等价于对每个 $\ell \in \{0, \ldots, h-1\}$,
\[
a_{\ell} = a_{\ell+1}\frac{n-m-\ell}{-m+h-\ell-1}.
\]
对 $\ell$ 做简单的下降归纳表明这等价于对每个 $\ell \in \{0, \ldots, h\}$,
\[
a_{\ell} = \frac{(n - m - h + 1)_{h-\ell}}{(-m)_{h-\ell}}a_h.
\]
这蕴含所要的结果.
(f). 由 (e) (和 $d^*$ 的 $S_n$-等变性), 函数 $\phi_h$ 是
\[
\psi := \frac{1}{(m-h)!}(d^*)^{m-h}\Bigl(\sum_{i=0}^h\frac{(n - m - h + 1)_{h-i}}{(-m)_{h-i}}\sigma_{i,h-i}(A)\Bigr)
\]
的倍数. 用 (c) 的公式计算 $\psi$. 得到
\begin{align*}
\psi &= \sum_{i=0}^h\frac{(n - m - h + 1)_{h-i}}{(-m)_{h-i}}\sum_{\ell=i}^{\min(m-h+i,n-m)}\binom{\ell}{i}\binom{m-\ell}{h-i}\sigma_{\ell,m-\ell}(A) \\
&= \sum_{\ell=0}^{\min(m,n-m)}\ \sum_{i=\max(0,\ell-m+h)}^{\min(\ell,h)}\binom{\ell}{i}\binom{m-\ell}{h-i}\frac{(n-m-h+1)_{h-i}}{(-m)_{h-i}}\sigma_{\ell,m-\ell}(A).
\end{align*}
这几乎就是我们要的公式, 除了开头的常数 $(-1)^h\binom{n-m}{h}^{-1}$. 球函数 $\phi_h$ 由 $\phi_h(A_0) = 1$ 规范化, 所以有 $\phi_h = \frac{1}{\psi(A_0)}\psi$. 所以为完成证明, 只需证明 $\psi(A_0) = (-1)^h\binom{n-m}{h}$.
注意 $\sigma_{\ell,m-\ell}(A)(A_0)$ 除非 $\ell = 0$ 为 $0$, 且 $\sigma_{0,m}(A)(A_0) = 1$, 所以
\begin{align*}
\psi(A_0) &= \binom{m}{h}\frac{(n - m - h + 1)_h}{(-m)_h} \\
&= \frac{m!}{h!(m-h)!}\cdot\frac{(n-m)!}{(n-m-h)!}\cdot\frac{(-1)^h(m-h)!}{m!} \\
&= (-1)^h\binom{n-m}{h}.
\end{align*}
(g). 试着忽略 (c), (d), (e) 和 (f).
函数 $\phi_h$ 球, 所以它在 $\Omega_m$ 中的 $S_m \times S_{n-m}$-轨道上为常数, 而这些轨道由习题 \ref{ex:V.8.2.1}(d) 是以 $A$ 为中心的球面. 由 (a) 和习题 \ref{ex:V.8.2.2}(a), 这意味着 $\phi_h$ 是函数 $\sigma_{\ell,m-\ell}(A)$ ($0 \leq \ell \leq \min(m, n-m)$) 的线性组合. 又, 由习题 \ref{ex:V.8.2.3}(g), 函数 $\phi_h$ 是 $\Delta$ 的特征值 $m(n-m) - h(n-h+1)$ 的特征向量. 设 $S = \{B \in \Omega_m \mid d(A, B) = 1\}$. 则 $\sigma_{1,m-1}(A)$ 是 $S$ 的特征函数, 且看到 $\phi_h$ 在 $S$ 上为常数, $\phi_h$ 中 $\sigma_{1,m-1}(A)$ 的系数是 $\frac{1}{|S|}\sum_{B \in S}\phi_h(B)$.
另一方面, $\sum_{B \in S}\phi_h(B)$ 由 $\Delta$ 的定义等于 $\Delta\phi_h(A)$, 而这等于 $(m(n-m) - h(n-h+1))\phi_h(A)$. 又, 由于 $\phi_h$ 球, 必有 $\phi_h(A) = 1$. 最后, $\phi_h$ 中 $\sigma_{1,m-1}(A)$ 的系数是 $\frac{1}{|S|}(m(n-m) - h(n-h+1))$. 为完成计算, 只需证明 $|S| = m(n-m)$. 这由从 $A$ 中移除一个元素 ($m$ 种选择) 并添加 $\{1, \ldots, n\} - A$ 的一个元素 ($n - m$ 种选择) 得到 $S$ 的每个元素 $B$ 这一事实得出.
\end{solution}

\subsection{问题: \index{Satake 同构}Satake 同构}
\label{sec:V.8.3}

此问题摘录自 Cartier 在 Corvallis 的综述 {[}6{]}, 尤其是第四章, 以及 Kottwitz 未发表的笔记.
我们使用下列约定:
\begin{itemize}
\item 若 $G$ 是局部紧群, 记 $\mu_G$ 为 $G$ 上的左 Haar 测度 (有时按特定方式规范化). 记对应的 $L^p$ 空间为 $L^p(G)$.
\item 若 $G$ 是局部紧群且 $H \subset G$ 是闭子群, 记 $\mu_{G/H}$ (相应地 $\mu_{H\backslash G}$) 为由测度 $\mu_G$ 和 $\mu_H$ 诱导的 $G/H$ (相应地 $H\backslash G$) 上左 (相应地右) $G$-不变测度. 见命题 \ref{prop:V.1.3}. 为构造测度 $\mu_{H\backslash G}$, 需要 $\mu_G$ 右不变.
\end{itemize}
为使事情更具体, 我为 $\GL_3(\Q_p)$ 写了此问题. 一切都推广到 $\GL_n(\Q_p)$ (对任意 $n \geq 1$), 事实上大多数问题都有只做最小修改就适用于 $\GL_n(\Q_p)$ 的解法.
设 $n = 3$. 写 $G = \GL_3(\Q_p)$, $K = \GL_3(\Z_p)$. 记 $B$ 为 $G$ 中上三角矩阵的子群, $T$ 为 $G$ 中对角矩阵的子群, $N$ 为 $G$ 中幂幺上三角矩阵 (即所有对角线元素等于 $1$ 的上三角矩阵) 的子群. 写 $X = \Z^3$ 和
\[
X^+ = \{(\lambda_1, \lambda_2, \lambda_3) \in \Z^3 \mid \lambda_1 \geq \lambda_2 \geq \lambda_3\}.
\]
若 $\lambda = (\lambda_1, \lambda_2, \lambda_3) \in X$, 令
\[
p^{\lambda} = \begin{pmatrix} p^{\lambda_1} & 0 & 0 \\ 0 & p^{\lambda_2} & 0 \\ 0 & 0 & p^{\lambda_3} \end{pmatrix} \in T.
\]
固定 $G$, $B$, $T$, $N$ 和 $K$ 上的左 Haar 测度 $\mu_G$, $\mu_B$, $\mu_T$, $\mu_N$ 和 $\mu_K$; 若没有混淆的危险, 在积分中只写 $dx$ 代替 $d\mu_G(x)$ 等. 我们要求 $H \in \{G, B, T, N, K\}$ 时 $\operatorname{vol}(H \cap K) = 1$ 来规范化所有 Haar 测度. 在此问题中, 记 $\mathcal{C}_c(X)$ 为具紧支撑的局部常值函数 $f \colon X \to \CC$ 的空间. (其中``局部常值''就是你想象的意思: $f \colon X \to \CC$ 局部常值, 若对每个 $x \in X$, 存在 $X$ 中 $x$ 的邻域 $U$ 使得 $f|_U$ 为常数. 注意这蕴含 $f$ 连续.)
设 $\mathcal{H} = \mathcal{C}_c(K\backslash G/K)$, $\mathcal{H}_T = \mathcal{C}_c(T/(T \cap K))$, 带由卷积乘积给出的代数结构. (很容易验证卷积乘积保持``局部常值''条件.) 此问题的目标是理解代数 $\mathcal{H}$ 的结构.
\begin{enumerate}
\item[(1)] 定义 $\delta \colon B \to \R_{>0}$:
\[
\delta\begin{pmatrix} a_1 & * & * \\ 0 & a_2 & * \\ 0 & 0 & a_3 \end{pmatrix} = \prod_{i=1}^n|a_i|_p^{n+1-2i}.
\]
证明 $\delta^{-1}$ 是 $B$ 的模函数, 且对每个 $f \in \mathcal{C}_c(B)$ 和每个 $b' \in B$,
\[
\int_B f(b'b)\,\,\dd b = \delta^{-1}(b')\int_B f(b)\,\,\dd b.
\]
\item[(2)] 证明 $G = BK$, 且对每个 $f \in \mathcal{C}_c(G)$,
\[
\int_G f(g)\,\dd g = \int_B\int_K f(bk)\,\dd b\,\dd k = \int_T\int_N\int_K f(tnk)\,\dd t\,\dd n\,\dd k.
\]
\item[(3)] 设 $(e_1, e_2, e_3)$ 是 $\Q_p^3$ 的典范基. 证明映射 $g \longmapsto \Z_pg(e_1) + \Z_pg(e_2) + \Z_pg(e_3)$ 诱导 $G/K$ 与 $\Q_p^3$ 的秩 $3$ 自由 $\Z_p$-子模的集合之间的双射.
\item[(4)] 证明
\[
G = \bigsqcup_{\lambda \in X^+}Kp^{\lambda}K.
\]
(这称为 Cartan 分解.)
\item[(5)] 证明 $(G, K)$ 是 Gelfand 对. (提示: $\theta(x) = (x^T)^{-1}$.)
\end{enumerate}
设 $z = (z_1, z_2, z_3) \in (\CC^{\times})^3$. 定义态射 $\chi_z \colon B \to \CC^{\times}$:
\[
\chi_z\begin{pmatrix} a_1 & * & * \\ 0 & a_2 & * \\ 0 & 0 & a_3 \end{pmatrix} = \prod_{i=1}^n z_i^{v_p(a_i)},
\]
其中对每个 $a \in \Q_p^{\times}$ 写 $|a|_p = p^{-v_p(a)}$ (所以 $v_p(a) \in \Z$). 记 $V_z$ 为满足下列条件的函数 $f \colon G \to \CC$ 的空间:
\begin{enumerate}
\item[(i)] 对每个 $b \in B$ 和每个 $g \in G$, $f(bg) = \chi_z(b)\delta(b)^{1/2}f(g)$;
\item[(ii)] 存在 $G$ 的开紧子群 $K_1$ 使得对每个 $g \in G$ 和每个 $k \in K_1$, $f(gk) = f(g)$.
\end{enumerate}
让 $G$ 按 $g \cdot f = R_g(f)$ 作用于 $V_z$.
\begin{enumerate}
\item[(6)] 证明 $V_z$ 是 $G$ 的表示 (注意我没有说关于连续性的任何事) 且 $\dim(V_z^K) = 1$.
\end{enumerate}
注记. (不必做此题.) $V_z$ 的元素局部常值但一般没有紧支撑, 因为条件 (i). 又, 我们没有在 $V_z$ 上定义拓扑, 所以不能说 $V_z$ 是 $G$ 的连续表示. 事实上, 它是所谓的光滑表示 (smooth representation), 这意味着 $V_z$ 的每个元素都被 $G$ 的一个开紧子群固定 (这由条件 (ii) 自动). 这意味着若在 $V_z$ 上放离散拓扑, 表示连续, 而且对完全不连通群的表示这是更自然的条件. 若想要希尔伯特空间上的连续表示, 会在 $V_z$ 上放内积 $\langle f_1, f_2 \rangle = \int_K f_1(k)\overline{f_1(k)}\,\dd k$ 并对这个内积完备化. 注意即使如此, 表示也不会酉, 除非 $|z_1| = |z_2| = |z_3| = 1$.
记 $\phi_z$ 为满足 $\phi_z(1) = 1$ 的 $V_z^K$ 的唯一元素.\footnote{它是球函数 (除非一般地无界), 因此有记号 $\phi$. 但除非你先证明它, 不能使用这一事实!} 对每个 $f \in \mathcal{H}$, 记 $f^{\vee} \colon (\CC^{\times})^3 \to \CC$ 为函数:\footnote{这道德上是球傅里叶变换, 但不幸的是 $p$-进群的通常约定不同于我们在课堂上使用的约定.}
\[
z \longmapsto \int_G f(g)\phi_z(g)\,\dd g.
\]
若 $f \in \mathcal{H}_T$, 定义 $f^{\vee} \colon (\CC^{\times})^3 \to \CC$:
\[
f^{\vee}(z) = \int_T f(t)\chi_z(t)\,\dd t.
\]
\begin{enumerate}
\item[(7)] 证明对所有 $f_1, f_2 \in \mathcal{H}$, $(f_1 * f_2)^{\vee} = f_1^{\vee}f_2^{\vee}$.
\item[(8)] 证明
\[
T = \bigsqcup_{\lambda \in X}(T \cap K)p^{\lambda}.
\]
\item[(9)] 证明 $f \longmapsto f^{\vee}$ 是从 $\mathcal{H}_T$ 到 $(\CC^{\times})^3$ 上在 $z_1^{\pm 1}, z_2^{\pm 1}, z_3^{\pm 1}$ 中多项式的函数的代数的 $\CC$-代数同构; 把这个代数等同于 $\CC[z_1^{\pm 1}, z_2^{\pm 1}, z_3^{\pm 1}]$.
\end{enumerate}
对每个 $f \in \mathcal{C}_c(G)$, 定义函数 $f^{(B)} \colon T \to \CC$:
\[
f^{(B)}(t) = \delta(t)^{1/2}\int_N f(tn)\,\dd n.
\]
\begin{enumerate}
\item[(10)] 证明对每个 $f \in \mathcal{C}_c(G)$, $f^{(B)} \in \mathcal{C}_c(T)$, 且若 $f \in \mathcal{H}$ 则 $f^{(B)} \in \mathcal{H}_T$.
\item[(11)] 证明对每个 $f \in \mathcal{H}$ 和每个 $z \in (\CC^{\times})^3$, $f^{\vee}(z) = (f^{(B)})^{\vee}(z)$.
\end{enumerate}
定义 $D \colon T \to \R_{>0}$:
\[
D\begin{pmatrix} a_1 & 0 & 0 \\ 0 & a_2 & 0 \\ 0 & 0 & a_3 \end{pmatrix} = \prod_{i,j \in \{1,\ldots,n\}, i \neq j}|1 - a_ia_j^{-1}|_p.
\]
若 $f \in \mathcal{C}_c(G)$, $t \in T$ 使得 $D(t) \neq 0$, 令
\[
\mathcal{O}_t(f) = \int_{T\backslash G}f(g^{-1}tg)\,\dd\mu_{T\backslash G}.
\]
\begin{enumerate}
\item[(12)] 证明对每个 $f \in \mathcal{C}_c(N)$ 和每个满足 $D(t) \neq 0$ 的 $t \in T$,
\[
\int_N f(n)\,\dd n = D(t)^{1/2}\delta(t)^{1/2}\int_N f(ntn^{-1}t^{-1})\,\dd n.
\]
(提示: 考虑由
\[
U = \left\{\begin{pmatrix} 1 & 0 & * \\ 0 & 1 & 0 \\ 0 & 0 & 1 \end{pmatrix}\right\}
\]
定义的子群 $U \subset N$, 并用公式 $\int_N f(n)\,\dd n = \int_{N/U}\int_U f(nu)\,\dd\mu_{N/U}(n)\,\dd\mu_U(u)$.)
\item[(13)] 证明对每个 $f \in \mathcal{H}$ 和每个满足 $D(t) \neq 0$ 的 $t \in T$,
\[
\mathcal{O}_t(f) = D(t)^{-1/2}f^{(B)}(t).
\]
\item[(14)] 证明对每个 $f \in \mathcal{H}$, 函数 $f^{\vee} \colon (\CC^{\times})^3 \to \CC$ 是 $\CC[z_1^{\pm 1}, z_2^{\pm 1}, z_3^{\pm 1}]^{S_3}$ 的元素, 即 $z_1^{\pm 1}, z_2^{\pm 1}, z_3^{\pm 1}$ 的对称多项式函数的代数:
\[
\{p(z_1, z_2, z_3) \in \CC[z_1^{\pm 1}, z_2^{\pm 1}, z_3^{\pm 1}] \mid \forall \sigma \in S_3,\, p(z_{\sigma(1)}, z_{\sigma(2)}, z_{\sigma(3)}) = p(z_1, z_2, z_3)\}.
\]
\end{enumerate}
在 $X = \Z^3$ 上定义偏序, 说 $(\mu_1, \mu_2, \mu_3) \leq (\lambda_1, \lambda_2, \lambda_3)$ 当且仅当 $\mu_1 \leq \lambda_1$, $\mu_1 + \mu_2 \leq \lambda_1 + \lambda_2$ 且 $\mu_1 + \mu_2 + \mu_3 = \lambda_1 + \lambda_2 + \lambda_3$. (注意最后一个关系是等式!) 对每个 $\lambda = (\lambda_1, \lambda_2, \lambda_3) \in X$, 定义 $z^{\lambda} \in \CC[z_1^{\pm 1}, z_2^{\pm 1}, z_3^{\pm 1}]$:
\[
z^{\lambda} = z_1^{\lambda_1}z_2^{\lambda_2}z_3^{\lambda_3}
\]
和 $f_{\lambda} \in \mathcal{H}$:
\[
f_{\lambda} = \mathbb{1}_{Kp^{\lambda}K}.
\]
\begin{enumerate}
\item[(15)] 证明 $(f_{\lambda})_{\lambda \in X^+}$ 是 $\mathcal{H}$ 的基.
\item[(16)] 对每个 $\lambda = (\lambda_1, \lambda_2, \lambda_3) \in X^+$, 设
\[
c_{\lambda} = \sum_{\sigma \in S_3}z^{(\lambda_{\sigma(1)}, \lambda_{\sigma(2)}, \lambda_{\sigma(3)})}.
\]
证明 $(c_{\lambda})_{\lambda \in X^+}$ 是 $\CC[z_1^{\pm 1}, z_2^{\pm 1}, z_3^{\pm 1}]^{S_3}$ 的基.
\item[(17)] 对每个 $\lambda \in X$, 写
\[
(f_{\lambda})^{\vee} = \sum_{\mu \in X}c_{\lambda}(\mu)z^{\mu},
\]
其中 $c_{\lambda}(\mu) \in \CC$. 此题的目标是证明下面两个事实:
\begin{enumerate}
\item[(i)] 若 $\lambda, \mu \in X^+$ 且 $\mu \nleq \lambda$, 则 $c_{\lambda}(\mu) = 0$.
\item[(ii)] 对每个 $\lambda \in X^+$,
\[
c_{\lambda}(\lambda) = p^{\frac{1}{2}\sum_{i=1}^n(n+1-2i)\lambda_i}.
\]
\end{enumerate}
若 $\mu = (\mu_1, \mu_2, \mu_3) \in X$, 写 $\check{\mu} = (\mu_3, \mu_2, \mu_1)$.
\begin{enumerate}
\item[(a)] 对每个 $g = (g_{ij}) \in M_3(\Q_p)$, 设 $\|g\| = \sup_{i,j}|g_{ij}|_p$. 证明对每个 $g \in M_3(\Q_p)$ 和 $k, k' \in K$, $\|g\| = \|kgk'\|$.
\item[(b)] 设 $r \in \{1, \ldots, n\}$, $\Omega_r$ 是 $\{1, \ldots, n\}$ 的基数为 $r$ 的子集的集合. 若 $g = (g_{ij}) \in M_3(\Q_p)$, 定义 $\Lambda_rg \colon \Omega_r \times \Omega_r \to \Q_p$:
\[
\Lambda_r(g)(A, A') = \det((g_{ij})_{i \in A, j \in A'}),
\]
且 $\|\Lambda_rg\|$:
\[
\|\Lambda_rg\| = \sup_{A,A' \in \Omega_r}|\Lambda_rg(A, A')|_p.
\]
证明对每个 $g \in M_3(\Q_p)$ 和 $r \in \{1, \ldots, n\}$:
\begin{enumerate}
\item[$(\alpha)$] 若 $k, k' \in K$, 则 $\|\Lambda_r(kgk')\| = \|\Lambda_rg\|$.
\item[$(\beta)$] 若 $t \in T$, $n \in N$, 则 $\|\Lambda_r(tn)\| \geq \|\Lambda_r(t)\|$.
\item[$(\gamma)$] 若 $g = p^{\lambda}$, $\lambda = (\lambda_1, \lambda_2, \lambda_3) \in X^+$, 则
\[
\|\Lambda_rg\| = p^{-(\lambda_{3-r+1} + \cdots + \lambda_3)}.
\]
\end{enumerate}
\item[(c)] 对所有 $\lambda, \mu \in X$, 证明
\[
c_{\lambda}(\mu) = f_{\lambda}^{(B)}(p^{\mu}) = f_{\lambda}^{(B)}(p^{\check{\mu}}).
\]
\item[(d)] 设 $\lambda, \mu \in X^+$. 若存在 $n \in N$ 使得 $p^{\mu}n \in Kp^{\lambda}K$, 证明 $\mu \leq \lambda$.
\item[(e)] 证明 (i).
\item[(f)] 设 $\lambda \in X^+$, $t = p^{\lambda}$. 对每个 $n \in N$, 证明 $tn \in KtK$ 当且仅当 $n \in N \cap K$.
\item[(g)] 证明 (ii).
\end{enumerate}
\item[(18)] 证明 $f \longmapsto f^{\vee}$ 诱导 $\CC$-代数同构 $\mathcal{H} \xrightarrow{\sim} \CC[z_1^{\pm 1}, z_2^{\pm 1}, z_3^{\pm 1}]^{S_3}$. (此同构称为 Satake 同构.)
\end{enumerate}

\begin{solution}
(1) 我们在习题 \ref{ex:I.5.3.8}(d) 中看到 $\delta^{-1}$ 是 $B$ 的模函数. 第二个结论由 $db$ 的左不变性和命题 \ref{prop:I.2.8} 得出.
(2) 注意, 由于 $B$ 和 $K$ 是 $G$ 的子群, 证明 $G = BK$ 与证明 $G = KB$ 等价. (此分解称为 Iwasawa 分解.) 我们用高斯消去算法 (a.k.a. 行化简) 的变体证明第二个结论. 在此算法中, 对 $\GL_3(\Q_p)$ 的矩阵应用三种操作:
\begin{itemize}
\item[(a)] 交换两行;
\item[(b)] 把一行乘以 $\Q_p^{\times}$ 的元素 $\alpha$;
\item[(c)] 把一行的 $\Q_p$ 的元素 $\beta$ 的倍数加到另一行.
\end{itemize}
每个这样的操作对应于左乘一个初等矩阵. 而且, 此矩阵在情形 (a) 属于 $\GL_3(\Z_p)$, 在情形 (b) 中若 $\alpha \in \Z_p^{\times}$ 属于 $\GL_3(\Z_p)$, 在情形 (c) 中若 $\beta \in \Z_p$ 属于 $\GL_3(\Z_p)$. 所以只需证明: 若对操作应用刚才对操作注意到的限制的操作序列 (a), (b), (c), 总可以得到上三角矩阵. 所以设 $g = (g_{ij})$ 是 $\GL_3(\Q_p)$ 的元素. 由于 $g$ 可逆, 它第一列有非零元素. 交换 $g$ 的行 (操作 (a)) 后, 可以假设 $g_{11} \neq 0$ 且对 $i = 2, 3$ 有 $|g_{11}|_p \geq |g_{i1}|_p$. 把 $g$ 的第一行的 $-g_{11}^{-1}g_{21}$ (相应地 $-g_{11}^{-1}g_{31}$) 倍加到 $g$ 的第二 (相应地第三) 行 (操作 (b); 注意由第一步我们是用 $\Z_p$ 的元素乘以第一行) 后, 可以假设 $g_{11}$ 是 $g$ 第一列的唯一非零元素. 则 $g_{22}$ 和 $g_{23}$ 至少有一个非零 (否则 $g$ 的秩 $\leq 2$), 所以交换 $g$ 的后两行后, 可以假设 $g_{22} \neq 0$ 且 $|g_{22}|_p \geq |g_{23}|_p$. 然后, 把第二行的 $-g_{22}^{-1}g_{23}$ 倍加到第三行后, 可以假设 $g_{23} = 0$. 此时矩阵 $g$ 上三角, 完成了.
现在证明第二个结论. 考虑群 $G_0 = B \times K$, 及其子群 $H = \{(x, x) \mid x \in B \cap K\}$. 考虑 $G_0$ 上的测度 $\mu_{G_0} = \mu_B \times \mu_K$; 这显然是左 Haar 测度. 还在紧群 $H$ 上放规范化 Haar 测度 $\mu_H$. 考虑 $G_0$ 在 $G$ 上的连续作用 $(b, k) \cdot g = bgk^{-1}$. 由此题的第一部分, 这个作用传递. 又, $1 \in G$ 的稳定化子显然是 $H$, 所以得到满连续映射 $\phi \colon G_0/H \to G$. 证明 $\phi$ 也开. 若 $x \in G_0/H$, 则 $((B \cap K) \times K)x$ 是 $x$ 的开邻域 (因为 $B \cap K$ 在 $B$ 中开), 且它在 $\phi$ 下的像是 $(B \cap K)\phi(x)K$, 这在 $G$ 中开, 因为 $K$ 是 $G$ 的开子群. 所以 $\phi$ 开, 从而是同胚. 设 $\mu$ 是 $G$ 上固定 Haar 测度在 $\phi^{-1}$ 下的原像. 由于 $G$ 幺模 (由习题 \ref{ex:I.5.3.8}(c)), 测度 $\mu$ 被 $G_0$ 左不变, 所以正线性泛函 $\mathcal{C}_c(G_0) \to \CC$, $h \longmapsto \int_{G_0/H}h^H(x)\,\dd\mu(x)$ 左不变, 因此它是 $\mu_{G_0}$ 的正倍数, 设为 $c\mu_{G_0}$ ($c > 0$). 设 $f \in \mathcal{C}_c(G)$. 定义函数 $h \in \mathcal{C}(G_0)$: $h(b, k) = f(bk^{-1})$; 有 $h(b, k) = 0$ 若 $b \notin (\operatorname{supp} f)K$, 所以 $h$ 有紧支撑. 应用习题 \ref{ex:I.5.3.6}, 得到
\[
\int_{G_0}h(b, k)\,\dd\mu_{G_0}(b, k) = \int_B\int_K f(bk^{-1})\,\dd b\,\dd k = c\int_{G_0/H}h^H(x)\,\dd\mu(x).
\]
设 $x \in G_0/H$. 选取 $x$ 的原像 $(b, k) \in G_0$, 使得 $\phi(x) = bk$. 则
\[
h^H(x) = \int_H h((b, k)h)\,\dd h = \int_{B \cap K}f(bhh^{-1}k)\,\dd h = f(bk).
\]
所以
\[
\int_B\int_K f(bk)\,\dd b\,\dd k = c\int_G f(g)\,\dd g,
\]
对每个 $f \in \mathcal{C}_c(G)$. 为计算 $c$, 把它应用于 $f = \mathbb{1}_K$ (这是具紧支撑的连续函数, 因为由习题 \ref{ex:I.5.1.4}(m), $K$ 是 $G$ 的开紧子群). 左端是 $\mu_B(B)\mu_K(K) = 1$, 右端是 $\mu_G(K) = 1$ (由 $\mu_G$ 的选取). 所以 $c = 1$.
最后, 把习题 \ref{ex:I.5.3.6} 应用于 $B$ 的子群 $T$ 和 $N$, 得到正常数 $c'$ 使得对每个 $f \in \mathcal{C}_c(B)$,
\[
\int_B f\,\dd\mu_B = c'\int_T\int_N f(tn)\,\dd\mu_T(t)\,\dd\mu_N(n).
\]
为证明 $c' = 1$, 只需对 $f = \mathbb{1}_{B \cap K}$ 计算两边, 并用 (容易的) 事实 $B \cap K = (T \cap K)(N \cap K)$.
(3) 首先注意如下容易的事实: 设 $(x_1, x_2, x_3)$ 是 $\Q_p^3$ 的元素族; 则此族在 $\Z_p$ 上线性无关当且仅当它在 $\Q_p$ 上线性无关. 事实上, 在 $\Q_p$ 上线性无关显然蕴含在 $\Z_p$ 上线性无关 (因为 $\Z_p \subset \Q_p$). 反之, 假设 $(x_1, x_2, x_3)$ 在 $\Z_p$ 上线性无关; 若有 $a_1x_1 + a_2x_2 + a_3x_3 = 0$, $a_1, a_2, a_3 \in \Q_p$, 则存在 $m \in \Z$ 使得 $p^ma_1, p^ma_2, p^ma_3 \in \Z_p$, 且关系 $p^ma_1x_1 + p^ma_2x_2 + p^ma_3x_3 = 0$ 蕴含 $p^ma_1 = p^ma_2 = p^ma_3 = 0$, 所以 $a_1 = a_2 = a_3 = 0$.
设 $\mathcal{L}$ 是 $\Q_p^3$ 的秩 $3$ 自由 $\Z_p$-子模的集合, $L_0 = \Z_pe_1 + \Z_pe_2 + \Z_pe_3 = \Z_p^3 \in \mathcal{L}$. 若 $L \in \mathcal{L}$, $g \in \GL_3(\Q_p)$, 则 $g(L)$ 也是 $\mathcal{L}$ 的元素. 事实上, 若 $(x_1, x_2, x_3)$ 是 $L$ 的 $\Z_p$-基, 则族 $(x_1, x_2, x_3)$ 也在 $\Q_p$ 上线性无关, 所以族 $(g(x_1), g(x_2), g(x_3))$ 在 $\Q_p^3$ 上线性无关, 从而在 $\Z_p$ 上线性无关, 所以它生成的 $\Q_p^3$ 的 $\Z_p$-子模是 $\mathcal{L}$ 的元素. 这证明 $(g, L) \longmapsto g(L)$ 定义 $G$ 在 $\mathcal{L}$ 上的左作用. 又, $L_0$ 的稳定化子显然是 $\GL_3(\Z_p)$.
证明这个作用传递, 这将给出所要的双射 $G/K \xrightarrow{\sim} \mathcal{L}$. 设 $L \in \mathcal{L}$, $(x_1, x_2, x_3)$ 是 $L$ 在 $\Z_p$ 上的基. 则看到族 $(x_1, x_2, x_3)$ 也在 $\Q_p$ 上线性无关. 所以 $(x_1, x_2, x_3)$ 是 $\Q_p$-向量空间 $\Q_p^3$ 的基, 因此存在 $g \in \GL_3(\Q_p)$ 使得对 $i = 1, 2, 3$ 有 $x_i = g(e_i)$, 然后 $L = g(L_0)$.
(4) 首先证明 $G = \bigcup_{\lambda \in X^+}Kp^{\lambda}K$. 设 $g \in G$. 设 $x_1, x_2, x_3 \in \Q_p^3$ 是 $g$ 的列, $L$ 是 $(x_1, x_2, x_3)$ 生成的 $\Q_p^3$ 的 $\Z_p$-子模; 在 (3) 的解答的记号中, 这是 $g(L_0)$. 由于 $\Q_p = \bigcup_{m \in \Z}p^m\Z_p$, 可以选择 $m \in \Z$ 使得矩阵 $p^mg$ 的元素都在 $\Z_p$ 中. 则 $(p^mx_1, p^mx_2, p^mx_3)$ 生成的 $\Q_p^3$ 的 $\Z_p$-子模 $L'$ 满足 $L' \supset L_0 = p^mL$. 由 Dummit 和 Foote 第 12 章定理 4, 存在 $L'$ 的基 $(y_1, y_2, y_3)$ 和元素 $a_1, a_2, a_3 \in \Z_p - \{0\}$ 使得 $a_3|a_2|a_1$ 且 $(a_1y_1, a_2y_2, a_3y_3)$ 是 $L_0$ 的基. 由于 $\Z_p - \{0\}$ 的每个元素 $a$ 都可以 (唯一地) 写成 $a = p^nu$, $n \in \Z_{\geq 0}$, $u \in \Z_p^{\times}$, 可以假设 $a_i = p^{n_i}$, 其中 $n_1, n_2, n_3 \in \Z_{\geq 0}$, $n_1 \geq n_2 \geq n_3$. 设 $g'$ (相应地 $k$) 是列 $a_1y_1, a_2y_2, a_3y_3$ (相应地 $y_1, y_2, y_3$) 的矩阵, $t$ 是对角线元素 $a_1, a_2, a_3$ 的对角矩阵; 换句话说, $t = p^{(n_1,n_2,n_3)}$. 则 $k \in K$, $g' = kt$. 又, 由于 $L_0 = g'(L_0) = p^mL = p^mg(L_0)$, 问题 (3) 蕴含存在 $k' \in K$ 使得 $p^mg = g'k'$. 所以得到
\[
g = p^{-m}g'k' = kp^{(n_1-m,n_2-m,n_3-m)}k' \in Kp^{(n_1-m,n_2-m,n_3-m)}K.
\]
注意还证明了: 若 $m = 0$ (例如若 $g$ 的所有元素都已在 $\Z_p$ 中), 使得 $L \subset L_0$, 则 $L_0/L \simeq (\Z/p^{n_1}\Z) \times (\Z/p^{n_2}\Z) \times (\Z/p^{n_3}\Z)$.
现在证明这个并是不相交并. 设 $\lambda, \lambda' \in X^+$, 假设 $Kp^{\lambda}K \cap Kp^{\lambda'}K \neq \emptyset$. 把 $p^{\lambda}$ 和 $p^{\lambda'}$ 乘以 $m$ 足够大的 $p^m$ 后, 可以假设 $\lambda$ 和 $\lambda'$ 在 $\Z_{\geq 0}^3$ 中. 设 $g \in Kp^{\lambda}K \cap Kp^{\lambda'}K$, $L = g(L_0)$. 由于 $g$ 的元素是 $\Z_p$ 的元素 (由对 $\lambda$ 的假设), 有 $L \subset L_0$. 写 $\lambda = (n_1, n_2, n_3)$, $\lambda' = (n'_1, n'_2, n'_3)$, 其中 $n_1 \geq n_2 \geq n_3$, $n'_1 \geq n'_2 \geq n'_3$. 由上一段,
\[
L_0/L \simeq (\Z/p^{n_1}\Z) \times (\Z/p^{n_2}\Z) \times (\Z/p^{n_3}\Z) \simeq (\Z/p^{n'_1}\Z) \times (\Z/p^{n'_2}\Z) \times (\Z/p^{n'_3}\Z).
\]
由 Dummit 和 Foote 第 12 章定理 9, 这蕴含对 $i = 1, 2, 3$ 有 $n_i = n'_i$, 即 $\lambda = \lambda'$.
(5) 考虑映射 $\theta \colon G \to G$, $x \longmapsto (x^T)^{-1}$. 这是群的连续同构, 因为 $\theta^{-1} = \theta$ 而是同胚. 也有 $\theta(K) = K$. 设 $g \in G$. 由 (4), 可以把 $g$ 写成 $g = kp^{\lambda}k'$, 其中 $k, k' \in K$, $\lambda \in \Z^3$. 则
\[
\theta(g) = \theta(k)p^{-\lambda}\theta(k') = \theta(k)k'g^{-1}k\theta(k') \in Kg^{-1}K.
\]
所以由命题 \ref{prop:V.2.5}, $(G, K)$ 是 Gelfand 对.
(6) 所有函数 $f \colon G \to \CC$ 的空间 (带 $G$ 的作用 $x \cdot f = R_xf$) 是 $G$ 的表示. 验证 $V_z$ 是这个空间的子表示. 首先, 它是向量子空间:
\begin{itemize}
\item 函数 $0$ 在 $V_z$ 中;
\item 条件 (i) 显然被线性组合保持;
\item 设 $f, f' \in V_z$, $K_1$ 和 $K'_1$ 是 $G$ 的开紧子群, 使得 $f$ (相应地 $f'$) 被 $K_1$ (相应地 $K'_1$) 右不变; 则 $K_1 \cap K'_1$ 是 $G$ 的紧子群, 且 $f$ 和 $f'$ 的任何线性组合被 $K_1 \cap K'_1$ 右不变, 所以满足条件 (ii).
\end{itemize}
现在验证 $V_z$ 是 $G$-不变的. 设 $f \in V_z$, $x \in G$. 对每个 $b \in B$ 和每个 $g \in G$,
\[
(R_xf)(bg) = f(bgx) = \chi_z(b)\delta(b)^{1/2}f(gx) = \chi_z(b)\delta(b)^{1/2}R_xf(g),
\]
所以 $R_xf$ 满足条件 (i). 设 $K_1$ 是 $G$ 的开紧子群, $f$ 被 $K_1$ 右不变. 则对每个 $g \in G$ 和 $k \in K_1$,
\[
R_xf(gxkx^{-1}) = f(gxk) = f(gx) = R_xf(g),
\]
所以 $R_xf$ 被 $G$ 的开紧子群 $xK_1x^{-1}$ 右不变, 因此满足条件 (ii).
最后计算 $\dim(V_z^K)$. 回忆映射 $G_0 := B \times K \to G$, $(b, k) \longmapsto bk^{-1}$ 把 $G$ 等同于商 $G_0/H$, 其中 $H = \{(x, x) \mid x \in B \cap K\}$. 考虑函数 $\varphi_z^0 \colon G_0 \to \R_{>0}$, $(b, k) \longmapsto \chi_z(b)\delta(b)^{1/2}$. 这是连续的群态射, 所以它在紧子群 $H$ 上平凡, 因此它下降到商 $G_0/H$, 并定义从 $G$ 到 $\R_{>0}$ 的连续函数, 记为 $\phi_z$. 由构造, 函数 $\phi_z$ 满足
\[
\phi_z(bk) = \chi_z(b)\delta(b)^{1/2}
\]
对每个 $b \in B$ 和每个 $k \in K$, 所以 $\phi_z(1) = 1$ 且 $\phi_z \in V_z^K$. 为完成证明, 证明 $V_z^K$ 由 $\phi_z$ 生成. 设 $f \in V_z^K$. 若 $g \in G$, 由问题 (2), 可以找到 $b \in B$, $k \in K$ 使得 $g = bk$, 然后
\[
f(g) = f(b) = \chi_z(b)\delta(b)^{1/2}f(1) = \phi_z(g)f(1).
\]
所以 $f$ 是 $\phi_z$ 的倍数.
(7) 固定 $z \in (\CC^{\times})^3$.
设 $f \in \mathcal{H}$. 定义 $\tilde{f} \in \mathcal{H}$: $\tilde{f}(x) = f(x^{-1})$. 证明 $\phi_z * \tilde{f} = f^{\vee}(z)\phi_z$. 首先注意 $\phi_z * \tilde{f}$ 被 $K$ 右不变, 因为 $\tilde{f}$ 被 $K$ 右不变 (这由命题 \ref{prop:V.1.8} 得出). 设 $x \in G$, $b \in B$. 则
\[
\phi_z * \tilde{f}(bx) = \int_G\phi_z(bxy)f(y)\,\dd y = \int_G\chi_z(b)\delta(b)^{1/2}\phi_z(xy)f(y)\,\dd y = \chi_z(b)\delta(b)^{1/2}\phi_z * \tilde{f}(x)
\]
(因为 $\phi_z \in V_z$). 所以 $\phi_z * \tilde{f} \in V_z^K$, 这由问题 (6) 蕴含 $\phi_z * \tilde{f}$ 是 $\phi_z$ 的纯量倍数. 为找到纯量, 在两个函数处取值 $1$: 有 $\phi_z(1) = 1$, 且
\[
\phi_z * \tilde{f}(1) = \int_G\phi_z(y)f(y)\,\dd y = f^{\vee}(z).
\]
现在设 $f_1, f_2 \in \mathcal{H}$. 则对每个 $x \in G$,
\begin{align*}
\widehat{f_1 * f_2}(x) &= \widetilde{f_1 * f_2}(x^{-1}) \\
&= \int_G f_1(x^{-1}y^{-1})f_2(y)\,\dd y \\
&= \int_G \tilde{f}_1(yx)\tilde{f}_2(y^{-1})\,\dd y = (\tilde{f}_2 * \tilde{f}_1)(x) \quad \text{(因为 } G \text{ 幺模)} \\
&= (\tilde{f}_1 * \tilde{f}_2)(x) \quad \text{(因为 } (G, K) \text{ 是 Gelfand 对)},
\end{align*}
所以 $\widehat{f_1 * f_2} = \tilde{f}_1 * \tilde{f}_2$. (实际上, 事实 $\widehat{f_1 * f_2} = \tilde{f}_2 * \tilde{f}_1$ 对我们的目的已足够好.) 用上一段的计算, 得到
\[
(f_1 * f_2)^{\vee}(z) = (\phi_z * \widehat{f_1 * f_2})(1) = (\phi_z * \tilde{f}_1 * \tilde{f}_2)(1) = ((f_1^{\vee}(z)\phi_z) * \tilde{f}_2)(1) = f_1^{\vee}(z)f_2^{\vee}(z).
\]
(8) 注意 $T \cap K$ 是使得 $t$ 和 $t^{-1}$ 的所有元素都在 $\Z_p$ 中的可逆对角矩阵的群, 即所有元素都在 $\Z_p^{\times}$ 中的群. 所以结论由 $\Q_p^{\times} = \bigsqcup_{m \in \Z}p^m\Z_p^{\times}$ (应用于矩阵的每个对角线元素) 这一事实得出, 而这个事实由习题 \ref{ex:I.5.1.4}(i) 得出.
(9) 映射 $f \longmapsto f^{\vee}$ 显然对 $f \in \mathcal{H}_T$ 线性. 证明它是代数态射. 设 $f_1, f_2 \in \mathcal{H}_T$, $z \in (\CC^{\times})^3$. 则
\begin{align*}
(f_1 * f_2)^{\vee}(z) &= \int_T (f_1 * f_2)(t)\chi_z(t)\,\dd t \\
&= \int_{T \times T}f_1(t')f_2((t')^{-1}t)\chi_z(t)\,\dd t\,\dd t' \\
&= \int_{T \times T}f_1(t')f_2(t)\chi_z(t't)\,\dd t\,\dd t' \\
&= \Bigl(\int_T f_1(t')\chi_z(t')\,\dd t'\Bigr)\Bigl(\int_T f_2(t)\chi_z(t)\,\dd t\Bigr) \quad \text{(因为 } \chi_z \text{ 乘法)} \\
&= f_1^{\vee}(z)f_2^{\vee}(z),
\end{align*}
如所愿.
由问题 (8), 族 $(\mathbb{1}_{(T \cap K)p^{\lambda}})_{\lambda \in X}$ 是 $\mathcal{H}_T$ 的基. 设 $\lambda = (\lambda_1, \lambda_2, \lambda_3) \in X$, $g_{\lambda} = \mathbb{1}_{(T \cap K)p^{\lambda}}$. 则对每个 $z = (z_1, z_2, z_3) \in (\CC^{\times})^3$,
\[
g_{\lambda}^{\vee}(z) = \int_{(T \cap K)p^{\lambda}}\chi_z(t)\,\dd t = \operatorname{vol}(T \cap K)\chi_z(p^{\lambda}) = z_1^{\lambda_1}z_2^{\lambda_2}z_3^{\lambda_3}.
\]
由于形如 $(z_1, z_2, z_3) \longmapsto z_1^{\lambda_1}z_2^{\lambda_2}z_3^{\lambda_3}$ ($(\lambda_1, \lambda_2, \lambda_3) \in X$) 的函数构成 $\CC[z_1^{\pm 1}, z_2^{\pm 1}, z_3^{\pm 1}]$ 的基, 这证明线性变换 $f \longmapsto f^{\vee}$ 把 $\mathcal{H}_T$ 的基送到 $\CC[z_1^{\pm 1}, z_2^{\pm 1}, z_3^{\pm 1}]$ 的基, 所以它是从 $\mathcal{H}_T$ 到 $\CC[z_1^{\pm 1}, z_2^{\pm 1}, z_3^{\pm 1}]$ 的同构.
(10) 由习题 \ref{ex:I.5.2.2}, $G$ 的开紧子群构成 $1$ 的邻域基. 设 $f \in \mathcal{C}_c(G)$. 由 $\mathcal{C}_c(G)$ 的定义, 对每个 $x \in G$, 存在 $G$ 的开紧子群 $K_x$ 使得 $f$ 在 $K_xxK_x$ 上为常数. 有 $\operatorname{supp} f \subset \bigcup_{x \in \operatorname{supp} f}K_xxK_x$, 且 $\operatorname{supp} f$ 紧, 所以存在 $x_1, \ldots, x_m \in \operatorname{supp} f$ 使得 $\operatorname{supp} f \subset \bigcup_{i=1}^m K_{x_i}x_iK_{x_i}$. 设 $L = \bigcap_{i=1}^m K_{x_i}$; 这是 $G$ 的开紧子群. 证明 $f$ 是双 $L$-不变的. 设 $x \in G$. 若 $x \in \operatorname{supp} f$, 则存在 $i \in \{1, \ldots, m\}$ 使得 $x \in K_{x_i}x_iK_{x_i}$. 则 $LxL \subset K_{x_i}^2x_iK_{x_i}^2$, 所以 $f$ 在 $LxL$ 上为常数. 若 $x \notin \operatorname{supp} f$, 要证明 $LxL \cap (\operatorname{supp} f) = \emptyset$. 但假设存在 $l, l' \in L$ 使得 $lxl' \in \operatorname{supp} f$, 则由刚才看到的, 知道 $f$ 在 $L(lxl')L = LxL$ 上为常数, 特别地 $f$ 在 $x$ 处不可能为零. 所以 $LxL \cap (\operatorname{supp} f) = \emptyset$, 特别地 $f$ 在 $LxL$ 上为常数 (等于 $0$).
现在证明 $f^{(B)}$ 被 $L \cap T$ 双不变; 这蕴含 $f^{(B)}$ 局部常值, 因为 $L \cap T$ 是 $T$ 的开紧子群. 由于 $T$ 交换, 只需证明 $f^{(B)}$ 被 $K \cap T$ 左不变, 但这立即由定义和 $\delta \colon T \to \R_{>0}$ 乘法 (因为 $\delta$ 必须在 $T$ 的紧子群 $K \cap T$ 上平凡) 这一事实得出.
最后, 需要证明 $f^{(B)}$ 有紧支撑. 设 $t \in T$ 使得 $f^{(B)}(t) \neq 0$. 则存在 $n \in N$ 使得 $tn \in B \cap (\operatorname{supp} f)$. 特别地, $t$ 在紧集 $B \cap (\operatorname{supp} f)$ 在连续投影映射 $B \to T \simeq B/N$ 下的像中. 这个像是紧的且包含 $f^{(B)}$ 的支撑, 所以 $f^{(B)}$ 有紧支撑.
若 $f \in \mathcal{H}$, 则可以在第一个断言的证明中取 $L = K$, 得到 $f^{(B)}$ 被 $K \cap T$ 双不变, 即 $f^{(B)} \in \mathcal{H}_T$.
(11) 设 $z \in (\CC^{\times})^3$. 注意乘法函数 $\delta$ 和 $\chi_z$ 都在 $N$ 上平凡. 所以若 $t \in T$, $n \in N$, 则 $\phi_z(tn) = \chi_z(t)\delta(t)^{1/2}$.
设 $f \in \mathcal{H}$. 则用问题 (2) 的积分公式, 得到
\begin{align*}
f^{\vee}(z) &= \int_T\int_N\int_K f(tnk)\phi_z(tnk)\,\dd t\,\dd n\,\dd k \\
&= \int_T\int_N f(tn)\phi_z(tn)\,\dd t\,\dd n \quad \text{(因为 } f \text{ 和 } \phi_z \text{ 都被 } K \text{ 右不变)} \\
&= \int_T\int_N f(tn)\chi_z(t)\delta(t)^{1/2}\,\dd t\,\dd n \\
&= \int_T f^{(B)}(t)\chi_z(t)\,\dd t \\
&= (f^{(B)})^{\vee}(z).
\end{align*}
(12) 实际上可以直接计算 (通过写出 $ntn^{-1}t^{-1}$ 的公式并用换元), 但采用提示, 因为它更好地推广到任意 $n$ 的 $\GL_n(\Q_p)$.
首先注意 $N$ 的子群 $U$ 正规 (甚至中心) 且同构于加法群 $\Q_p$. 特别地, 有由习题 \ref{ex:I.5.3.4} 给出的 $U$ 上的 Haar 测度 $\mu_U = \lambda$ (它由要求 $\Z_p$ 体积为 $1$ 规范化).
其次, 注意映射 $N \to \Q_p^2$, $\begin{pmatrix} 1 & a & b \\ 0 & 1 & c \\ 0 & 0 & 1 \end{pmatrix} \longmapsto (a, c)$ 诱导拓扑群同构 $N/U \xrightarrow{\sim} \Q_p^2$, 且 $N$ 在 $N/U$ 上通过左平移的作用在此同构下对应于 $N$ 在 $\Q_p^2$ 上的作用 $\begin{pmatrix} 1 & a & b \\ 0 & 1 & c \\ 0 & 0 & 1 \end{pmatrix} \cdot (x, y) = (a + x, c + y)$. 特别地, $\Q_p^2$ 上的测度 $\lambda^2$ 给出 $N/U$ 上被 $N$ 左不变的规则 Borel 测度 $d\mu_{N/U}$. 则正线性函数 $\mathcal{C}_c(N) \to \CC$, $f \longmapsto \int_{N/U}\int_U f(nu)\,\dd\mu_{N/U}(n)\,\dd\mu_U(u)$ 左不变, 所以它是 Haar 测度 $\mu_N$ 的倍数. 又, 取 $f = \mathbb{1}_{K \cap N}$ (并用 $K \cap U = \Z_p$ 以及 $(K \cap N)/(K \cap U)$ 在上述同构下对应于 $\Q_p^2$ 的子群 $\Z_p^2$ 这一事实), 看到这个正线性函数实际上等于 $\mu_N$, 这证明提示中的积分公式.
在下面的计算中, 把 $U$ 等同于 $\Q_p$, 把 $N/U$ 等同于 $\Q_p^2$. 设 $t = \begin{pmatrix} x_1 & 0 & 0 \\ 0 & x_2 & 0 \\ 0 & 0 & x_3 \end{pmatrix} \in T$. 注意映射 $N \to N$, $n \longmapsto ntn^{-1}t^{-1}$ 保持 $U$, 所以它诱导映射 $U \to U$ 和 $N/U \to N/U$. 若 $u \in U = \Q_p$, 有 $utu^{-1}t^{-1} = (1 - x_1x_3^{-1})u$. 另一方面, 若 $n \in N$, 且 $n$ 在 $N/U = \Q_p^2$ 中的像是 $(a, b)$, 则 $ntn^{-1}t^{-1}$ 在 $N/U$ 中的像是 $(a(1 - x_1x_2^{-1}), b(1 - x_2x_3^{-1}))$.
假设 $D(t) \neq 0$. 由习题 \ref{ex:I.5.3.8} 证明的换元公式, 有:
\begin{itemize}
\item 对每个 $f \in \mathcal{C}_c(U)$,
\[
\int_U f(u)\,\dd u = |1 - x_1x_3^{-1}|_p\int_U f(utu^{-1}t^{-1})\,\dd u;
\]
\item 对每个 $f \in \mathcal{C}_c(N/U)$,
\[
\int_{N/U}f(n)\,\dd n = |1 - x_1x_2^{-1}|_p|1 - x_2x_3^{-1}|_p\int_{N/U}f(ntn^{-1}t^{-1})\,\dd n.
\]
\end{itemize}
又, 用 $U$ 在 $N$ 中中心且 $t^{-1}Ut = U$ 这一事实, 看到对每个 $n \in N$ 和每个 $u \in U$,
\[
(ntn^{-1}t^{-1})(utu^{-1}t^{-1}) = ntn^{-1}(t^{-1}ut)u^{-1}t^{-1} = nt(t^{-1}ut)u^{-1}n^{-1}t^{-1} = nut(nu)^{-1}t^{-1}.
\]
把积分公式和上面的计算放在一起, 看到对每个 $f \in \mathcal{C}_c(N)$,
\begin{align*}
\int_N f(n)\,\dd n &= \int_{N/U}\int_U f(nu)\,\dd n\,\dd u \\
&= |1 - x_1x_2^{-1}|_p|1 - x_2x_3^{-1}|_p\int_{N/U}\int_U f(ntn^{-1}t^{-1}utu^{-1}t^{-1})\,\dd n\,\dd u \\
&= |1 - x_1x_2^{-1}|_p|1 - x_2x_3^{-1}|_p\int_{N/U}\int_U f(nut(nu)^{-1}t^{-1})\,\dd n\,\dd u.
\end{align*}
现在做代换 $(n, u) \longmapsto nu$ (这给出从 $N/U \times U$ 到 $N$ 的同胚), 得到
\[
\int_N f(n)\,\dd n = |1 - x_1x_2^{-1}|_p|1 - x_2x_3^{-1}|_p\int_N f(ntn^{-1}t^{-1})\,\dd n.
\]
这正是所要的, 因为
\[
D(t) = |1 - x_1x_2^{-1}|_p|1 - x_1x_3^{-1}|_p|1 - x_2x_3^{-1}|_p|1 - x_2x_1^{-1}|_p|1 - x_3x_1^{-1}|_p|1 - x_3x_2^{-1}|_p
\]
且 $|1 - x_ix_j^{-1}|_p = |1 - x_jx_i^{-1}|_p$ (因为 $|1 - x_ix_j^{-1}|_p = |x_jx_i^{-1}|_p|1 - x_jx_i^{-1}|_p = |1 - x_jx_i^{-1}|_p$), 所以 $D(t) = (|1 - x_1x_2^{-1}|_p|1 - x_2x_3^{-1}|_p)^2$, 且 $D(t)^{1/2} = |1 - x_1x_2^{-1}|_p|1 - x_2x_3^{-1}|_p$. 而 $\delta(t) = |x_1|_p^2|x_2|_p^0|x_3|_p^{-2}$, 所以 $\delta(t)^{1/2} = |x_1|_p|x_3|_p^{-1} = |1 - x_1x_3^{-1}|_p$ 在 $D(t) \neq 0$ 时由上一段最后一个等式等于 $|1 - x_1x_3^{-1}|_p$.
实际上, 上面的计算给出
\[
\int_N f(n)\,\dd n = D(t)^{1/2}\delta(t)^{1/2}\int_N f(ntn^{-1}t^{-1})\,\dd n
\]
对每个 $f \in \mathcal{C}_c(N)$, 如所愿.
(13) 设 $f \in \mathcal{H}$, $t \in T$ 满足 $D(t) \neq 0$. 回忆 $f^{(B)}(t) = \delta(t)^{1/2}\int_N f(tn)\,\dd n$. 用 (12) (把 $f$ 换成 $n \longmapsto f(tn)$ 并做代换), 得到
\begin{align*}
f^{(B)}(t) &= \delta(t)^{1/2}D(t)^{-1/2}\delta(t)^{-1/2}\int_N f(tntn^{-1}t^{-1})\,\dd n \\
&= D(t)^{-1/2}\int_N f(tntn^{-1}t^{-1})\,\dd n.
\end{align*}
另一方面, 需要把 $\mathcal{O}_t(f) = \int_{T\backslash G}f(g^{-1}tg)\,\dd\mu_{T\backslash G}$ 用积分公式 $\int_G = \int_{T\backslash G}\int_N\int_T$ 重写. 用 (2) 的分解 $G = TNK$ 和双不变性, 得到
\[
\mathcal{O}_t(f) = \int_N f(tntn^{-1}t^{-1})\,\dd n
\]
(这里我们用了 $\int_G f = \int_{T\backslash G}\int_N\int_T f(tnk)$, 而 $\mu_{T\backslash G}$ 与 $\mu_T \times \mu_N$ 的配合给出测量分解). 组合两式给出 $\mathcal{O}_t(f) = D(t)^{-1/2}f^{(B)}(t)$.
(14) 设 $f \in \mathcal{H}$. 由 (5), $(G, K)$ 是 Gelfand 对, 所以由定理 \ref{thm:V.4.4}, $\chi_{\phi_z} \colon \mathcal{H} \to \CC$ 是乘法泛函. 但由 (7), $\chi_{\phi_z}(f) = f^{\vee}(z)$, 所以对 $f_1, f_2 \in \mathcal{H}$, $(f_1 * f_2)^{\vee}(z) = f_1^{\vee}(z)f_2^{\vee}(z)$, 且 $\mathcal{H}$ 交换 (因为 $(G, K)$ 是 Gelfand 对). 现在 $\mathcal{H} = \mathcal{C}_c(K\backslash G/K)$, 且由 (10), 映射 $f \longmapsto f^{(B)}$ 是 $\mathcal{H}$ 到 $\mathcal{H}_T$ 的代数同态 (对卷积). 所以 $f \mapsto f^{(B)}$ 诱导 $\mathcal{H}$ 在 $\CC[z_1^{\pm 1}, z_2^{\pm 1}, z_3^{\pm 1}]$ 的表示, 而 (11) 说 $f^{\vee}(z) = (f^{(B)})^{\vee}(z)$. 由 (9), $(f^{(B)})^{\vee}$ 在 $\CC[z_1^{\pm 1}, z_2^{\pm 1}, z_3^{\pm 1}]$ 中. 为证明对称性: 由 (13), 对满足 $D(t) \neq 0$ 的 $t \in T$, $\mathcal{O}_t(f) = D(t)^{-1/2}f^{(B)}(t)$. 但 $\mathcal{O}_t(f) = \int_{T\backslash G}f(g^{-1}tg)$ 是对 $t$ 的共轭类积分, 且显然 $\mathcal{O}_t(f) = \mathcal{O}_{t'}(f)$ 若 $t$ 与 $t'$ 共轭. 因为 $f$ 双不变, 函数 $t \longmapsto f^{(B)}(t)$ 在 $T$ 的 $W = S_3$ (Weyl 群) 轨道上为常数. 因此 $f^{\vee}(z) = (f^{(B)})^{\vee}(z)$ 是对 $z$ 的坐标置换对称的, 即 $f^{\vee} \in \CC[z_1^{\pm 1}, z_2^{\pm 1}, z_3^{\pm 1}]^{S_3}$.
(15) 由 (4), $\mathcal{H}$ 作为向量空间由 $f_{\lambda} = \mathbb{1}_{Kp^{\lambda}K}$, $\lambda \in X^+$ 张成 (这些是双不变函数的自然基). 为证明线性无关: 若 $\sum_{\lambda}c_{\lambda}f_{\lambda} = 0$, 在 $p^{\mu}$ ($\mu \in X^+$) 处取值, 得到 $c_{\mu} = 0$ (因为 $p^{\mu} \in Kp^{\mu}K$ 且 $p^{\mu} \notin Kp^{\lambda}K$ 对 $\lambda \neq \mu$, 由 (4) 的不相交性).
(16) 由 (8) 和 (9), 更准确地说由 (9) 的证明, $\CC[z_1^{\pm 1}, z_2^{\pm 1}, z_3^{\pm 1}]$ 有基 $(z^{\lambda})_{\lambda \in X}$. $S_3$ 通过置换变量作用于其上, 且由定义 $c_{\lambda} = \sum_{\sigma \in S_3}z^{\sigma \cdot \lambda}$. 由于 $S_3$ 有限, 投影算子 $p \longmapsto \frac{1}{|S_3|}\sum_{\sigma \in S_3}\sigma \cdot p$ 把 $\CC[z_1^{\pm 1}, z_2^{\pm 1}, z_3^{\pm 1}]$ 映到其 $S_3$-不变部分, 且它是满射. 由 (17) 后面的标准论证 (用偏序和三角形论证) 或者更简单地: 由对称函数的基本定理 (对称多项式 $\Leftrightarrow$ 基本对称多项式的多项式), 或者直接论证: 把 $z^{\lambda}$ 在 $S_3$-轨道上的和 $c_{\lambda}$ 记为 $\lambda \in X^+$; 它们在偏序 $\leq$ 下形成一个三角形系 (每个 $c_{\lambda}$ 有首项 $z^{\lambda}$), 所以构成不变部分的基.
(17) 参考上面 (14) 和 (15) 以及 Cartier 的论证. 由于篇幅, 我们把细节留给读者并给出核心步骤:
(a) 证明 $\|g\| = \|kgk'\|$ 对 $k, k' \in K$ (因为 $K = \GL_3(\Z_p)$ 保持 sup 范数).
(b) $\Lambda_r$ 是 $r \times r$ 子行列式; 它们在 $K$ 的左、右乘法下不变 (因为子行列式的列/行组合在 $\Z_p$ 中, 范数 $\leq 1$), 给出 $(\alpha)$. 对 $(\beta)$: 对 $t \in T$, $n \in N$, 用行阶梯: $\Lambda_r(tn)$ 的每个子行列式是 $t$ 的某个主子行列式与 $N$-因子之积, 且 $N$-因子的 $p$-进绝对值 $\leq 1$, 所以 $\|\Lambda_r(tn)\| \geq \|\Lambda_r(t)\|$. 对 $(\gamma)$: 对 $g = p^{\lambda}$, $\lambda \in X^+$, 最大的 $r \times r$ 子行列式 (绝对值) 是主子式 $\prod_{i=3-r+1}^3 p^{\lambda_i}$, 给出 $\|\Lambda_r g\| = p^{-(\lambda_{3-r+1} + \cdots + \lambda_3)}$.
(c) 用 (13): $c_{\lambda}(\mu) = (f_{\lambda})^{\vee}$ 中 $z^{\mu}$ 的系数 $= f_{\lambda}^{(B)}(p^{\mu})$ (用 (11) 和 (9) 的基), 且 $f_{\lambda}^{(B)}(p^{\mu}) = f_{\lambda}^{(B)}(p^{\check{\mu}})$, 因为 $f_{\lambda}^{(B)}$ 在 $T$ 的 Weyl 群轨道上为常数 (由 (14) 的论证, 因为 $f_{\lambda}$ 双不变).
(d) 若 $p^{\mu}n \in Kp^{\lambda}K$, 则 $\|\Lambda_r(p^{\mu}n)\| = \|\Lambda_r(p^{\lambda})\|$ 对每个 $r$ (由 (a)-$(\gamma)$), 所以由 $(\beta)$, $\|\Lambda_r(p^{\mu})\| \geq \|\Lambda_r(p^{\lambda})\|$ 对每个 $r$. 由 $(\gamma)$, 这蕴含 $\sum_{i=3-r+1}^3\mu_i \leq \sum_{i=3-r+1}^3\lambda_i$ 对每个 $r$, 即 $\mu \leq \lambda$.
(e) (i) 由 (d) 和 (c) 得出: 若 $\mu \nleq \lambda$, 则不存在 $n \in N$ 使得 $p^{\mu}n \in Kp^{\lambda}K$; 但 $c_{\lambda}(\mu) = f_{\lambda}^{(B)}(p^{\mu}) = \delta(p^{\mu})^{1/2}\int_N f_{\lambda}(p^{\mu}n)\,\dd n = 0$, 因为 $f_{\lambda}$ 是 $Kp^{\lambda}K$ 的特征函数且 $p^{\mu}n \notin Kp^{\lambda}K$.
(f) 若 $tn \in KtK$ 且 $n \in N$, 则 $p^{\mu}n \in Kp^{\lambda}K$ (取 $\mu = \lambda$), 所以用 (e) 的论证的逆: $\|\Lambda_r(tn)\| = \|\Lambda_r(t)\|$ 对每个 $r$; 由 $(\beta)$ 的等号情形 (需要 $n$ 的因子都已在 $\Z_p^{\times}$ 中), 得到 $n \in N \cap K$. 反之显然.
(g) (ii) 由 (c), (f) 和 (2) 的积分公式得出: $c_{\lambda}(\lambda) = f_{\lambda}^{(B)}(p^{\lambda}) = \delta(p^{\lambda})^{1/2}\int_N f_{\lambda}(p^{\lambda}n)\,\dd n = \delta(p^{\lambda})^{1/2}\operatorname{vol}(N \cap K) = \delta(p^{\lambda})^{1/2}$, 因为 $f_{\lambda}$ 是 $Kp^{\lambda}K$ 的特征函数且 $p^{\lambda}n \in Kp^{\lambda}K$ 当且仅当 $n \in N \cap K$ (由 (f)), 而 $\operatorname{vol}(N \cap K) = 1$. 由 (1) 中 $\delta$ 的定义, $\delta(p^{\lambda})^{1/2} = \prod_{i=1}^3|p^{\lambda_i}|_p^{(3+1-2i)/2} = p^{\frac{1}{2}\sum_{i=1}^3(4-2i)\lambda_i}$, 即所要的公式.
(18) 把 $f \longmapsto f^{\vee}$ 记为 $\mathcal{S}$. 在 (14) 中看到它良定义, 在 (7) 中看到它是 $\CC$-代数态射. 所以只需证明它是 $\CC$-向量空间的同构. 在 (15) 中给出 $\mathcal{H}$ 的基 $(f_{\lambda})_{\lambda \in X^+}$, 在 (16) 中给出 $\CC[z_1^{\pm 1}, z_2^{\pm 1}, z_3^{\pm 1}]^{S_3}$ 的基 $(c_{\lambda})_{\lambda \in X^+}$. 又, 集合 $X^+$ 有偏序, $\mathcal{S}$ 在两个基中的矩阵由 (17)(i) 对这个偏序上三角, 由 (17)(ii) 有非零对角线元素. 所以 $\mathcal{S}$ 可逆.
若想把论证做得更显式, 可以这样说: 由 (17)(i) 和 (ii), 对每个 $\lambda \in X^+$,
\[
\mathcal{S}(f_{\lambda}) = \sum_{\mu \in X^+, \mu \leq \lambda}c_{\lambda}(\mu)c_{\mu},
\]
其中 $c_{\lambda}(\lambda) \in \CC^{\times}$. 递归构造 $\mathcal{S}$ 的逆 $T \colon \CC[z_1^{\pm 1}, z_2^{\pm 1}, z_3^{\pm 1}]^{S_3} \to \mathcal{H}$. 只需给出 $T$ 在基元素 $c_{\lambda}$, $\lambda \in X^+$ 上的值. 若不存在满足 $\mu \leq \lambda$ 且 $\mu \neq \lambda$ 的 $\mu \in X^+$, 设 $T(c_{\lambda}) = c_{\lambda}(\lambda)^{-1}f_{\lambda}$. 否则, 递归定义 $T(c_{\lambda})$:
\[
T(c_{\lambda}) = c_{\lambda}(\lambda)^{-1}f_{\lambda} - c_{\lambda}(\lambda)^{-1}\sum_{\mu \in X^+ - \{\lambda\}, \mu \leq \lambda}c_{\lambda}(\mu)T(c_{\mu}).
\]
由于 $\{\mu \in X^+ \mid \mu \leq \lambda\}$ 有限 (事实上, 若 $\mu \in X^+$ 且 $\mu \leq \lambda$, 则 $\mu_3 = (\lambda_1 + \lambda_2 + \lambda_3) - (\mu_1 + \mu_2) \geq \lambda_3$, 所以 $\lambda_3 \leq \mu_3 \leq \mu_2 \leq \mu_1 \leq \lambda_1$), 这个过程总会终止.
\end{solution}

\subsection{问题}
\label{sec:V.8.4}
\footnote{这个结果叫什么名字? 有参考文献吗?}\footnote{归功于 H. T. G. Dospinescu.}
在此问题中, 写 $G = \SL_2(\R)$, $K = \SO(2)$. 记 $B$ 为 $G$ 中上三角矩阵的子群, $N$ 为幂幺上三角矩阵 (即对角线元素全为 $1$ 的上三角矩阵) 的子群.\footnote{本问题的一切都对 $\SL_n(\R)$ 成立, 但推广比问题 V.8.3 更麻烦.}\footnote{参考文献.}
在所有群上固定左 Haar 测度. 回忆 $G$ 幺模.
\begin{enumerate}
\item[(1)] 此题的目标是证明如下事实 $(*)$: 若 $(\pi, V)$ 是 $G$ 的酉表示, $v \in V$ 被 $G$ 的形如 $\begin{pmatrix} 1 & u \\ 0 & 1 \end{pmatrix}$ ($u \neq 0$) 的元素固定, 则 $v$ 被 $G$ 的每个元素固定.
我们固定 $G$ 的酉表示 $(\pi, V)$.
\begin{enumerate}
\item[(a)] 设 $v \in V$, 假设存在连续群态射 $\chi \colon \R \to S^1$ 使得对每个 $u \in \R$ 有 $\pi\begin{pmatrix} 1 & u \\ 0 & 1 \end{pmatrix}(v) = \chi(u)v$. 定义 $\phi \colon G \to \CC$: $\phi(x) = \langle \pi(x)(v), v \rangle$.
  \begin{enumerate}
  \item[(i)] 证明若 $x = \begin{pmatrix} a & b \\ c & d \end{pmatrix}$, $x' = \begin{pmatrix} a' & b' \\ c' & d' \end{pmatrix}$ 满足 $c = c'$, 则 $|\phi(x)| = |\phi(x')|$.
  \item[(ii)] 证明若 $x \in B$, 则 $|\phi(x)| = |\phi(1)|$.
  \item[(iii)] 证明存在连续群态射 $\psi \colon B \to S^1$ 使得对每个 $x \in B$, $\pi(x)(v) = \psi(x)v$.
  \item[(iv)] 证明函数 $x \longmapsto |\phi(x)|$ 在 $G$ 上为常数.
  \item[(v)] 证明对每个 $x \in G$ 有 $\pi(x)(v) = v$.
  \end{enumerate}
\item[(b)] 证明 $(*)$.
\end{enumerate}
\item[(2)] 设
\[
A^+ = \left\{\begin{pmatrix} a & 0 \\ 0 & a^{-1} \end{pmatrix}, a \in \R_{\geq 1}\right\} \subset G.
\]
\begin{enumerate}
\item[(a)] 证明
\[
G = \bigcup_{a \in A^+}KaK.
\]
(作为集合, 不作为拓扑空间.)
\item[(b)] 证明 $(G, K)$ 是 Gelfand 对.
\item[(c)] 把问题 (a) 和 (b) 推广到 $\SL_n(\R)$, 其中 $K = \SO(n)$, $A^+$ 是满足 $a_1 \geq a_2 \geq \cdots \geq a_n > 0$ 的对角线系数 $a_1, \ldots, a_n$ 的对角矩阵的集合.
\end{enumerate}
\item[(3)] 此题的目标是证明如下事实 $(**)$: 对满足 $V^G = \{0\}$ 的每个酉表示 $(\pi, V)$, 对所有 $v, w \in V$, $\lim_{x \to \infty}\langle \pi(x)(v), w \rangle = 0$, 其中 $\infty$ 是 $G$ 的 Alexandroff 紧化的无穷远点. (所以 $\lim_{x \to \infty}\langle \pi(x)(v), w \rangle = 0$ 意味着: 对每个 $\varepsilon > 0$, 存在 $G$ 的紧子集 $X$ 使得对每个 $x \in G - X$, $|\langle \pi(x)(v), w \rangle| \leq \varepsilon$.)
我们固定 $G$ 的酉表示 $(\pi, V)$, 假设它不满足 $(**)$ 的结论.
\begin{enumerate}
\item[(a)] 证明存在 $v, v' \in V$, $\R_{\geq 1}$ 的元素序列 $(a_n)_{n \geq 0}$ 和 $\alpha \in \CC - \{0\}$ 使得:
\begin{itemize}
\item 当 $n \to +\infty$ 时 $a_n \to +\infty$;
\item 当 $n \to +\infty$ 时 $\langle \pi(t_n)(v), v' \rangle \to \alpha$, 其中 $t_n = \begin{pmatrix} a_n & 0 \\ 0 & a_n^{-1} \end{pmatrix}$.
\end{itemize}
\item[(b)] 证明取 $(a_n)_{n \geq 0}$ 的子序列后, 可以假设存在 $v_0 \in V$ 使得对每个 $w \in V$, 当 $n \to +\infty$ 时 $\langle \pi(t_n)(v), w \rangle \to \langle v_0, w \rangle$.
\item[(c)] 证明 $v_0 \neq 0$ 且 $v_0$ 被 $N$ 的每个元素固定. (提示: 若 $x \in N$, $w \in V$, 当 $n \to +\infty$ 时 $\langle \pi(xt_n)(v), w \rangle - \langle \pi(t_n)(v), w \rangle \to 0$ 的行为如何?)
\item[(d)] 得出结论.
\end{enumerate}
\item[(4)] 此题的目标是证明商 $\SL_2(\Z)\backslash G$ 有有限体积.
设 $\mathfrak{h} := \{z \in \CC \mid \Im(z) > 0\}$.
对 $g = \begin{pmatrix} a & b \\ c & d \end{pmatrix} \in G$ 和 $z \in \mathfrak{h}$, 令
\[
g \cdot z = \frac{az + b}{cz + d}.
\]
\begin{enumerate}
\item[(a)] 证明这定义 $G$ 在 $\mathfrak{h}$ 上的左作用.
\item[(b)] 证明 $i \in \mathfrak{h}$ 的稳定化子是 $K$, 且这诱导同胚 $G/K \xrightarrow{\sim} \mathfrak{h}$.
\item[(c)] 证明 (b) 的同胚把测度 $\mu_{G/K}$ 送到 $\frac{dx\,dy}{y^2}$ 的倍数, 其中 $x$ 和 $y$ 是 $z \in \mathfrak{h}$ 的实部和虚部, $dx$ 和 $dy$ 是 $\R$ 上的 Lebesgue 测度.
\item[(d)] 设 $\Omega = \{z \in \mathfrak{h} \mid |\Re(z)| \leq \frac{1}{2} \text{ 且 } |z| \geq 1\}$. 证明对每个 $z \in \mathfrak{h}$, 存在 $g \in \SL_2(\Z)$ 使得 $g \cdot z \in \Omega$.
\item[(e)] 证明 $\int_{\Omega}\frac{dx\,dy}{y^2} < +\infty$.
\item[(f)] 证明 $\operatorname{vol}(\SL_2(\Z)\backslash G) < +\infty$.
\end{enumerate}
\item[(5)] 在此题中, 固定 $G$ 的闭离散子群 $\Gamma$, 使得 $\operatorname{vol}(\Gamma\backslash G) < +\infty$. 注意由习题 \ref{ex:I.5.3.7}, $G$ 幺模, 所以 $G$ 上的 Haar 测度也右不变. 写 $X = \Gamma\backslash G$, 记 $\mu$ 为测度 $\mu_{\Gamma\backslash G}$; 规范化 $G$ 和 $\Gamma$ 上的 Haar 测度使得 $\mu(X) = 1$. 考虑空间 $L^2(X) = L^2(X, \mu)$, 带 $G$ 的右拟正则表示 $\pi$ (即对 $g \in G$, $f \in L^2(X)$, $\pi(g)(f) = R_gf$).
证明对所有 $f_1, f_2 \in L^2(X)$,
\[
\lim_{g \to \infty}\int_X f_1(xg)f_2(x)\,\dd\mu(x) = \int_X f_1(x)\,\dd\mu(x)\int_X f_2(x)\,\dd\mu(x).
\]
(提示: $V := \{f \in L^2(X) \mid \int_X f(x)\,\dd\mu(x) = 0\}$ 中 $G$-不变向量是什么?)
\item[(6)] 用问题 (5) 的记号. 设 $Y = (\Gamma \cap K)\backslash K$. 在 $Y$ 上放测度 $\mu_Y = \mu_{(\Gamma\cap K)\backslash K}$, 规范化 $\Gamma \cap K$ 和 $K$ 上的 Haar 测度使得 $Y$ 体积为 $1$.
此题的目标是证明如下事实 $(***)$: 对每个 $f \in \mathcal{C}_c(X)$,
\[
\lim_{g \to \infty}\int_Y f(yg)\,\dd\mu_Y(y) = \int_X f(x)\,\dd\mu(x).
\]
\footnote{换句话说, 当 $g \to \infty$ ($G$ 中) 时集合 $Yg$ 在 $X$ 中变得等分布. 注意 $g \to \infty$ 当且仅当 $g^{-1} \to \infty$.}
把 $\mu_K$ 看作 $\mathcal{C}_c(G)$ 上的连续线性泛函:
\[
\mu_K(f) = \int_K f(x)\,\dd\mu_K(x),
\]
对 $f \in \mathcal{C}_c(G)$. 如果 $\mu_K$ 可以由 $\mathcal{C}_c(G)$ 的元素表示 (即如果有 $h \in \mathcal{C}_c(G)$ 使得对每个 $f \in \mathcal{C}_c(G)$, $\int_K f(x)\,\dd\mu_K(x) = \int_G f(x)h(x)\,\dd\mu_G(x)$), 那么证明 $(***)$ 会相对容易 (从已经做过的). 但事实并非如此. 不过, 可以尝试用 $\mathcal{C}_c(G)$ 的元素逼近 $\mu_K$. 你可能记得我们已经多次用过这种技巧来逼近 Dirac 测度.
\begin{enumerate}
\item[(a)] 若 $\psi \in \mathcal{C}_c(G)$, 定义 $\mathcal{C}_c(G)$ 上的连续线性泛函 $\psi * \mu_K$:
\[
\psi * \mu_K(f) = \int_{G \times K}f(xy)\psi(x)\,\dd\mu_G(x)\,\dd\mu_K(y).
\]
证明存在 $h \in \mathcal{C}_c(G)$ 使得对每个 $f \in \mathcal{C}_c(G)$,
\[
\psi * \mu_K(f) = \int_G f(x)h(x)\,\dd\mu_G(x).
\]
\item[(b)] 证明存在 $\mathcal{C}_c(G)$ 的元素序列 $(\psi_n)_{n \geq 0}$ 使得 $\mu_K$ 是序列 $(\psi_n * \mu_K)_{n \geq 0}$ 在 $\Hom(\mathcal{C}_c(G), \CC)$ 的弱*拓扑中的极限.
\item[(c)] 证明 $(***)$.
\end{enumerate}
\end{enumerate}

\begin{solution}
(1)(a). 注意若 $n = \begin{pmatrix} 1 & u \\ 0 & 1 \end{pmatrix}$, $n' = \begin{pmatrix} 1 & u' \\ 0 & 1 \end{pmatrix} \in N$, $x \in G$, 则
\[
|\phi(nxn')| = |\langle \pi(x)(\pi(n')(v)), \pi(n)^{-1}(v) \rangle| = |\chi(u')\chi(u)\langle \pi(x)(v), v \rangle| = |\phi(x)|.
\]
(i) 首先假设 $c = c' \neq 0$. 设 $n = \begin{pmatrix} 1 & c^{-1}(d' - d) \\ 0 & 1 \end{pmatrix}$, $n' = \begin{pmatrix} 1 & c^{-1}(a' - a) \\ 0 & 1 \end{pmatrix}$. 则
\[
n'xn = \begin{pmatrix} a' & w \\ c' & d' \end{pmatrix},
\]
对某个 $w \in \R$. 由于 $1 = \det(n'xn) = a'd' - wc' = a'd' - b'c'$ 且 $c' \neq 0$, 必有 $w = b'$, 所以 $n'xn = x'$, 且上面的注记蕴含 $|\phi(x)| = |\phi(x')|$.
现在假设 $c = c' = 0$. 则可以找到 $G$ 中的序列 $(x_n)_{n \geq 1}$ 和 $(x'_n)_{n \geq 1}$ 使得 $x_n \to x$, $x'_n \to x'$, 且对每个 $n \geq 1$, $x_n$ 和 $x'_n$ 的 $(2,1)$ 项都等于 $\frac{1}{n}$. 例如可以取 $x_n = \begin{pmatrix} a + nd & b \\ \frac{1}{n} & d \end{pmatrix}$, 这有意义因为 $\det(x) = ad = 1$ 蕴含 $d \neq 0$, 用类似的公式定义 $x'_n$. 把第一种情形应用于 $x_n$ 和 $x'_n$ 给出对每个 $n \geq 1$ 有 $|\phi(x_n)| = |\phi(x'_n)|$; 由于 $\phi$ 连续, 可以在 $n$ 中取极限, 得到 $|\phi(x)| = |\phi(x')|$.
(ii) 这由问题 (i) 的``$c = c' = 0$''情形得出.
(iii) 设 $x \in B$. 由 (ii), $|\langle \pi(x)(v), v \rangle| = |\langle v, v \rangle|$, 所以由 Cauchy-Schwarz 不等式的等号情形, 必有 $\pi(x)(v) = \psi(x)v$, 其中 $\psi(x) \in S^1$. 由于 $x \longmapsto \pi(x)(v)$ 对 $x$ 连续且 $\psi(x) = \langle \pi(x)(v), v \rangle$, 函数 $x \longmapsto \psi(x)$ 也连续. 它群态射这一事实立即由对所有 $x, y \in B$ 有 $\pi(xy) = \pi(x) \circ \pi(y)$ 得出.
(iv) 重复 (a) 的证明开头所做的, 用 (iii), 看到对每个 $x \in G$ 和所有 $b, b' \in G$, $|\phi(bxb')| = |\phi(x)|$. 设 $x_0 = \begin{pmatrix} 0 & 1 \\ -1 & 0 \end{pmatrix}$. 设 $x = \begin{pmatrix} a & b \\ c & d \end{pmatrix} \in G$. 若 $c \neq 0$, 则
\[
x = \begin{pmatrix} 1 & c^{-1}a \\ 0 & 1 \end{pmatrix}x_0\begin{pmatrix} c & -d \\ 0 & b - c^{-1}ad \end{pmatrix},
\]
所以 $|\phi(x)| = |\phi(x_0)|$. 若 $c = 0$, 则可以用 $(2,1)$ 项非零的 $G$ 的元素逼近 $x$ (如在 (i) 的证明中), 所以由第一种情形和 $\phi$ 的连续性, 也得到 $|\phi(x)| = |\phi(x_0)|$.
(v) 重复 (iii) 的推理, 证明存在连续群态射 $\psi \colon G \to S^1$ (延拓 (iii) 的映射 $\psi$) 使得对每个 $x \in G$ 有 $\pi(x)(v) = \psi(x)v$. 但不存在从 $G$ 到交换群的非平凡群态射 (因为 $G$ 等于其换位子群, 而到交换群的群态射在每个换位子上平凡), 所以 $\psi = 1$, 因此 $v \in V^G$.
(b). 记 $\pi^0$ 为 $\R$ 在 $V$ 上的酉表示 $\pi^0(t) = \pi\begin{pmatrix} 1 & t \\ 0 & 1 \end{pmatrix}$, 对每个 $t \in \R$. 由假设 $\pi^0(u)(v) = v$. 由于 $\R$ 交换, 这蕴含 $\pi^0(t)(w) = w$ 对每个 $w \in \{\pi(t)(v) \mid t \in \R\}$ 成立, 从而 (因为两边对 $w$ 都线性且连续) 对每个 $w \in W := \overline{\operatorname{Span}\{\pi(t)(v) \mid t \in \R\}}$ 成立. 换句话说, $\R$ 在 $W$ 上的酉表示 $\pi^0|_W$ 在子群 $\Z u$ 上平凡, 所以它诱导紧群 $\R/\Z u$ 在 $W$ 上的酉表示. 由定理 \ref{thm:IV.2.1} 和推论 \ref{cor:I.3.4.4}, 空间 $W$ 是 $\R/\Z u$ 的 $1$ 维表示的直和的闭包, 即存在线性无关族 $(v_i)_{i \in I}$ 生成 $W$ 的稠密子空间, 且每个 $v_i$ 是所有 $\pi(n)$, $n \in N$ 的特征向量. 由问题 (a), 这蕴含所有 $v_i$ 都在 $V^G$ 中, 所以 $W \subset V^G$, 特别地 $v \in V^G$.
(2)(a). 直接对 $\SL_n(\R)$ 证明结果, 所以写 $G = \SL_n(\R)$, $K = \SO(n)$.
对 $\GL_n(\R)$ 的元素用极分解 (见 Rudin {[}20{]} 的定理 12.35\footnote{更好的参考文献.}, 虽然这对我们的目的来说是杀鸡用牛刀, 因为我们只需要有限维情形): 对每个 $g \in \GL_n(\R)$, 存在唯一偶对 $(u, p)$ 使得 $u$ 正交, $p$ 对称正定且 $g = up$; 又, $p^2 = g^Tg$. 若 $g \in G$, 则 $\det(p)^2 = \det(g^Tg) = 1$; 由于 $p$ 正定, 有 $\det(p) > 0$, 所以 $\det(p) = 1$, 然后也推出 $\det(u) = 1$, 即 $u \in K$.
证明 $G = \bigcup_{a \in A^+}KaK$. 设 $g \in G$. 设 $g = up$ 是 $g$ 的极分解. 由于 $p$ 对称, 由谱定理它可以在标准正交基中对角化, 所以可以把 $p$ 写成 $p = kak^{-1}$, 其中 $k \in \Or(n)$, $a$ 对角. $a$ 的对角线元素是 $p$ 的特征值, 可以假设它们降序排列, 所以 $a \in A^+$. 又, 若 $k$ 不在 $\SO(n)$ 中, 则改变特征基的一个向量的符号; 这不影响 $a$, 新的 $k$ 在 $\SO(n)$ 中. 最后得到 $g = (uk)ak^{-1} \in KaK$.
还要证明集合 $KaK$, $a \in A^+$, 两两不相交. 设 $a, a' \in A^+$, 假设 $KaK \cap Ka'K \neq \emptyset$. 则存在 $k, l \in K$ 使得 $a' = kal$. 所以
\[
(a')^2 = (a')^Ta' = k^Ta^Tk^Tkak^T = ka^2k^{-1},
\]
这蕴含 $(a')^2$ 和 $a^2$ 有相同的特征值, 即它们有相同的对角线元素 (可能重排). 由于 $(a')^2$ 和 $a^2$ 的对角线元素递减, 这两个矩阵实际上有相同的对角线元素, 即 $(a')^2 = a^2$. 最后, 由于 $a$ 和 $a'$ 是对角矩阵且元素为正, $(a')^2 = a^2$ 这一事实蕴含 $a' = a$.
(b). (仍取 $G = \SL_n(\R)$, $K = \SO(n)$.) 考虑连续群自同构 $\theta \colon G \to G$, $g \longmapsto (g^T)^{-1}$. 有 $\theta^2 = \mathrm{id}_G$, 且 $\theta|_K = \mathrm{id}_K$. 对每个 $g \in G$, 若写 $g = kal$, 其中 $k, l \in K$, $a \in A^+$, 则 $\theta(g) = ka^{-1}l = klg^{-1}kl \in Kg^{-1}K$. 由命题 \ref{prop:V.2.5}, $(G, K)$ 是 Gelfand 对.
(3) 现在回到 $\SL_2(\R)$, 所以释放字母 $n$.
(a). 由于 $V$ 不满足 $(**)$ 的结论, 存在 $v_1, v_2 \in V$ 和 $\varepsilon > 0$ 使得对 $G$ 的每个紧子集 $X$, 存在 $x \in G - X$ 使得 $|\langle \pi(x)(v_1), v_2 \rangle| > \varepsilon$.
对每个 $n \geq 1$, 设 $A_n^+$ 是满足 $a \leq n$ 的 $\begin{pmatrix} a & 0 \\ 0 & a^{-1} \end{pmatrix} \in A^+$ 的集合. 则 $\bigcup_{t \in A_n^+}KtK$ 在 $G$ 中相对紧, 因为若 $g \in KtK$, $t \in A_n^+$, 则 $g$ 和 $g^{-1}$ 的所有元素在绝对值上都以 $4n$ 为界 (注意若 $k \in K$, 则 $k$ 的所有元素在绝对值上都以 $1$ 为界). 所以可以找到 $x_n \in G - \bigcup_{t \in A_n^+}KtK$ 使得 $|\langle \pi(x_n)(v_1), v_2 \rangle| > \varepsilon$. 写 $x_n = k_nt_nl_n$, 其中 $k_n, l_n \in K$, $t_n = \begin{pmatrix} a_n & 0 \\ 0 & a_n^{-1} \end{pmatrix} \in A_n^+$. 由于 $K$ 紧, 取 $(k_n)$ 和 $(l_n)$ 的子序列后, 可以假设 $(k_n)$ (相应地 $(l_n)$) 有极限 $k \in K$ (相应地 $l \in L$). 设 $v = \pi(l)(v_1)$, $v' = \pi(k)^{-1}(v_2)$. 由于 $\CC$ 的序列 $(\langle \pi(t_n)(v), v' \rangle)$ 有界 (被 $\|v\|\|v'\|$), 取子序列后, 可以假设它收敛到某个 $\alpha \in \CC$. 为完成证明, 只需证明 $\alpha \neq 0$. 注意对每个 $n \geq 1$,
\[
\langle \pi(t_n)v, v' \rangle = \langle \pi(kk_n^{-1})\pi(k_nt_nl_n)\pi(l_n^{-1}l)v_1, v_2 \rangle = \langle \pi(x_n)\pi(l_n^{-1}l)v_1, \pi(k_nk^{-1})v_2 \rangle,
\]
所以
\[
|\langle \pi(t_n)v, v' \rangle - \langle \pi(x_n)v_1, v_2 \rangle| \leq \|v_1\|\|v_2\|(\|\pi(l_n^{-1}l) - \mathrm{id}\|_{\mathrm{op}} + \|\pi(k_nk^{-1}) - \mathrm{id}\|_{\mathrm{op}}).
\]
当 $n \to \infty$ 时这趋于 $0$. 由于对每个 $n \geq 1$ 有 $|\langle \pi(x_n)v_1, v_2 \rangle| \geq \varepsilon$, 序列 $(\langle \pi(t_n)(v), v' \rangle)$ 的极限 $\alpha$ 不能是 $0$.
(b). 结论意味着 $V$ 上的有界线性泛函序列 $(\Lambda_n \colon w \longmapsto \langle \pi(t_n)(v), w \rangle)_{n \geq 1}$ 在 $V^*$ 的弱*拓扑中收敛到有界线性泛函 $w \longmapsto \langle v_0, w \rangle$. 由于每个 $\Lambda_n$ 在 $V^*$ 的闭单位球中 (对算子范数), 且由 Banach-Alaoglu 定理这个闭单位球弱*紧, 取子序列后, 确实可以假设序列 $(\Lambda_n)_{n \geq 1}$ 在弱*拓扑中收敛到某个 $\Lambda \in V^*$. 但知道 $\Lambda$ 必具有 $w \longmapsto \langle v_0, w \rangle$ 的形式, 其中 $v_0 \in V$ 唯一确定, 所以我们完成了.
(c). 有 $\langle v_0, v' \rangle = \lim_{n \to +\infty}\langle \pi(t_n)(v), v' \rangle = \alpha \neq 0$, 所以 $v_0 \neq 0$.
设 $x \in N$, 写 $x = \begin{pmatrix} 1 & u \\ 0 & 1 \end{pmatrix}$. 设 $w \in V$. 则
\[
\langle \pi(xt_n)(v), w \rangle - \langle \pi(t_n)(v), w \rangle = \langle \pi(t_n)(\pi(t_n^{-1}xt_n)(v) - v), w \rangle,
\]
所以
\[
|\langle \pi(xt_n)(v), w \rangle - \langle \pi(t_n)(v), w \rangle| \leq \|(\pi(t_n^{-1}xt_n) - \pi(1))(v)\|\|w\|.
\]
由于 $t_n^{-1}xt_n = \begin{pmatrix} 1 & a_n^{-2}u \\ 0 & 1 \end{pmatrix}$ 且 $\lim_{n \to +\infty}a_n = +\infty$, 有 $\lim_{n \to +\infty}t_n^{-1}xt_n = 1$. 用表示 $\pi$ 的连续性, 推出 $\lim_{n \to +\infty}(\pi(t_n^{-1}xt_n)(v) - v) = 0$, 所以 $\lim_{n \to +\infty}|\langle \pi(xt_n)(v), w \rangle - \langle \pi(t_n)(v), w \rangle| = 0$. 但知道 $\lim_{n \to +\infty}\langle \pi(t_n)(v), w \rangle = \langle v_0, w \rangle$, 所以得到 $\lim_{n \to +\infty}\langle \pi(xt_n)(v), w \rangle = \langle v_0, w \rangle$. 另一方面,
\[
\lim_{n \to +\infty}\langle \pi(xt_n)(v), w \rangle = \lim_{n \to +\infty}\langle \pi(t_n)(v), \pi(x)^{-1}(w) \rangle = \langle v_0, \pi(x)^{-1}(w) \rangle = \langle \pi(x)(v_0), w \rangle.
\]
由于这对每个 $w \in V$ 成立, 推出 $\pi(x)(v_0) = v_0$.
(d). 由问题 (1), (c) 蕴含 $v_0$ 被 $G$ 的每个元素固定. 所以 $V^G \neq \{0\}$, 这完成了 $(**)$ 的逆否命题的证明.
(4)(a). 首先验证 $g \cdot z$ 良定义且在 $\mathfrak{h}$ 中 (若 $z \in \mathfrak{h}$). 写 $g = \begin{pmatrix} a & b \\ c & d \end{pmatrix}$. 若 $c \neq 0$, 则 $\Im(cz + d) = c\,\Im(z) \neq 0$, 所以商 $\frac{az+b}{cz+d}$ 有意义. 若 $c = 0$, 则 $ad = \det(g) = 1$, 所以 $d \neq 0$, 故 $cz + d = d \neq 0$, 因此商 $\frac{az+b}{cz+d}$ 有意义. 计算这个商的虚部. 有
\[
\frac{az + b}{cz + d} = \frac{(az + b)(cz + d)}{(cz + d)(cz + d)} = \frac{ac|z|^2 + bd + adz + bc\bar{z}}{(cz + d)(cz + d)},
\]
所以
\[
\Im\Bigl(\frac{az + b}{cz + d}\Bigr) = \frac{1}{(cz + d)(cz + d)}(ac\,\Im(z) - bd\,\Im(z)) = \frac{\Im(z)}{(cz + d)(cz + d)} > 0,
\]
且 $\frac{az+b}{cz+d} \in \mathfrak{h}$.
若 $g = 1$, 显然 $g \cdot z = z$. 设 $h = \begin{pmatrix} a' & b' \\ c' & d' \end{pmatrix}$ 是 $G$ 的另一个元素. 必须验证 $(gh) \cdot z = g \cdot (h \cdot z)$. 这是直接计算.
(b). 设 $g = \begin{pmatrix} a & b \\ c & d \end{pmatrix} \in G$ 使得 $g \cdot i = i$. 这等价于 $ai + b = i(ci + d) = -c + id$, 即 $a = d$ 且 $b = -c$. 若 $a = d$ 且 $b = -c$, 则 $1 = ad - bc = a^2 + b^2 = d^2 + c^2$, 且 $ac + bd = 0$, 所以 $g$ 是正交矩阵, 因为 $\det(g) = 1$ 而在 $K$ 中. 反之, 若 $g \in K$, 则 $g$ 具有 $\begin{pmatrix} \cos\theta & \sin\theta \\ -\sin\theta & \cos\theta \end{pmatrix}$ 的形式, 所以确实有 $a = d$ 和 $b = -c$, 因此 $g \cdot i = i$.
若 $z$ 是 $\mathfrak{h}$ 的任意元素, 写 $z = x + iy$, $x \in \R$, $y \in \R_{>0}$. 则 $z = g(z) \cdot i$, 其中 $g(z) = \begin{pmatrix} \sqrt{y} & \sqrt{y}^{-1}x \\ 0 & \sqrt{y}^{-1} \end{pmatrix} \in G$. 所以 $G$ 在 $\mathfrak{h}$ 上的作用传递, 且诱导连续双射 $\alpha \colon G/K \xrightarrow{\sim} \mathfrak{h}$. 为证明 $\alpha$ 是同胚, 注意刚才构造的映射 $\mathfrak{h} \to G$, $z \longmapsto g(z)$ 连续, 且 $z \longmapsto g(z)K$ 是 $\alpha$ 的逆.
(c). 设 $\mu$ 是 $\frac{dx\,dy}{y^2}$ 在 $\alpha^{-1}$ 下的原像. 像习题 \ref{sec:V.8.3} 的 (2) 的证明那样推理, 看到只需证明 $\mu$ 被 $G$ 左不变. 换句话说, 只需证明测度 $\frac{dx\,dy}{y^2}$ 被 $G$ 在 $\mathfrak{h}$ 上的作用不变.
设 $g = \begin{pmatrix} a & b \\ c & d \end{pmatrix} \in G$, $z = x + iy \in \mathfrak{h}$. 在 (a) 的证明中看到
\[
\Im(g \cdot z) = \frac{y}{(cz + d)(cz + d)} = \frac{y}{c^2(x^2 + y^2) + 2cdx + d^2}.
\]
类似地,
\[
\Re(g \cdot z) = \frac{ac(x^2 + y^2) + (ad + bc)x + bd}{c^2(x^2 + y^2) + 2cdx + d^2}.
\]
若写 $x' = \Re(g \cdot z)$, $y' = \Im(g \cdot z)$, 要证明 $\frac{dx'\,dy'}{(y')^2} = \frac{dx\,dy}{y^2}$.
设 $g_0 = \begin{pmatrix} 0 & 1 \\ -1 & 0 \end{pmatrix}$. 在 (1)(a)(iv) 的证明中看到 $G = B \cup Bg_0B$, 所以只需证明 $\frac{dx\,dy}{y^2}$ 被 $B$ 的元素和 $g_0$ 的作用不变. 两种情形都用 $\R^2$ 上通常的换元计算 $dx'\,dy'$.
若 $g \in B$ (即 $c = 0$), 则 $x' = \frac{1}{d^2}(x + bd)$, $y' = \frac{y}{d^2}$, 所以 $dx'\,dy' = \frac{1}{d^4}dx\,dy$, 且得到 $\frac{dx'\,dy'}{(y')^2} = \frac{dx\,dy}{y^2}$.
若 $g = g_0$, 则 $x' = -\frac{x}{x^2 + y^2}$, $y' = \frac{y}{x^2 + y^2}$, 所以
\[
dx'\,dy' = \det\begin{pmatrix} \frac{-(x^2+y^2)+2x^2}{(x^2+y^2)^2} & \frac{-2xy}{(x^2+y^2)^2} \\ \frac{2xy}{(x^2+y^2)^2} & \frac{(x^2+y^2)-2y^2}{(x^2+y^2)^2} \end{pmatrix}dx\,dy = \frac{dx\,dy}{(x^2 + y^2)^2},
\]
且再次这立即蕴含 $\frac{dx'\,dy'}{(y')^2} = \frac{dx\,dy}{y^2}$.
(d). 若 $g = \begin{pmatrix} a & b \\ c & d \end{pmatrix} \in \SL_2(\Z)$, 则 $ad - bc = 1$, 所以 $c$ 和 $d$ 在 $\Z$ 中互素. 反之, 若 $c$ 和 $d$ 在 $\Z$ 中互素, 则存在 $a, b \in \Z$ 使得 $ad - bc = 1$, 然后 $g := \begin{pmatrix} a & b \\ c & d \end{pmatrix} \in \SL_2(\Z)$.
设 $z \in \mathfrak{h}$. 则群 $\Z + \Z z$ 在 $\CC$ 中离散 (因为 $z \notin \R$), 所以它与每个球相交有限. 特别地, $|cz + d|$ ($c, d \in \Z$ 互素) 的下确界达到. 由上一段和 $\Im(g \cdot z) = \frac{\Im(z)}{|cz+d|^2}$ (若 $g = \begin{pmatrix} a & b \\ c & d \end{pmatrix}$) 这一事实, 看到 $\sup\{|\Im(g \cdot z)| \mid g \in \SL_2(\Z)\}$ 也达到. 要证明存在 $g \in \SL_2(\Z)$ 使得 $g \cdot z \in \Omega$. 用某个 $g \in \SL_2(\Z)$ 替换 $z$ 后, 可以假设 $|\Im(z)| = \sup\{|\Im(g \cdot z)| \mid g \in \SL_2(\Z)\}$. 特别地, 取 $g = \begin{pmatrix} 0 & -1 \\ 1 & 0 \end{pmatrix}$, 得到 $|\Im(z)| \geq \frac{|\Im(z)|}{|z|^2}$, 所以 $|z| \geq 1$. 若 $g_n = \begin{pmatrix} 1 & n \\ 0 & 1 \end{pmatrix}$ ($n \in \Z$), 则 $g_n \cdot z = z + n$, 所以 $\Im(g_n \cdot z) = \Im(z)$, 且仍有 $|\Im(g_n \cdot z)| = \sup\{|\Im((hg_n) \cdot z)| \mid h \in \SL_2(\Z)\}$, 特别地仍有 $|g_n \cdot z| \geq 1$. 可以选择 $n \in \Z$ 使得 $-\frac{1}{2} \leq \Re(z) + n \leq \frac{1}{2}$, 然后有 $g_n \cdot z \in \Omega$, 完成了.
(e). 有:
\[
\int_{\Omega}\frac{dx\,dy}{y^2} = \int_{-1/2}^{1/2}\int_{\sqrt{1-x^2}}^{+\infty}\frac{dx\,dy}{y^2} = \int_{-1/2}^{1/2}\frac{1}{\sqrt{1-x^2}}\,dx < +\infty.
\]
(f). 对每个 $\varepsilon \in (0, 1/4)$, 设
\[
\Omega_{\varepsilon} = \{z \in \mathfrak{h} \mid |z| > 1 - \varepsilon \text{ 且 } |\Im(z)| < \frac{1}{2} + \varepsilon\}.
\]
如在 (e) 中, 看到
\[
I_{\varepsilon} := \int_{\Omega_{\varepsilon}}\frac{dx\,dy}{y^2} = \int_{-1/2-\varepsilon}^{1/2+\varepsilon}\frac{1}{\sqrt{1-x^2-\varepsilon}}\,dx < +\infty.
\]
设
\[
\Omega'_{\varepsilon} = \{g \in G \mid g \cdot i \in \Omega_{\varepsilon}\}.
\]
这是 $\Omega_{\varepsilon}$ 在连续映射 $G \to G/K$ 下的原像. 特别地, 它是 $G$ 的开子集, 所以由正则 Borel 测度的定义 (和 Urysohn 引理),
\[
\mu_G(\Omega'_{\varepsilon}) = \sup\Bigl\{\int_G f\,\dd\mu_G \mid f \in \mathcal{C}_c(G),\, 0 \leq f \leq 1 \text{ 且 } \operatorname{supp}(f) \subset \Omega'_{\varepsilon}\Bigr\}.
\]
若 $f \in \mathcal{C}_c(G)$ 满足 $0 \leq f \leq 1$ 且 $\operatorname{supp}(f) \subset \Omega'_{\varepsilon}$, 则对每个 $g \in G$, $f^K(gK) = \int_K f(gk)\,\dd\mu_K(k) \in [0,1]$ (因为 Haar 测度 $\mu_K$ 的总体积为 $1$), 且 $f^K(gK) = 0$ 若 $g \notin \Omega'_{\varepsilon}$, 即 $gK \notin \Omega_{\varepsilon}$; 所以 $\int_{G/K}f^K\,\dd\mu_{G/K} \leq \mu_{G/K}(\Omega_{\varepsilon})$. 由于对每个 $f \in \mathcal{C}_c(G)$ 有 $\int_G f\,\dd\mu_G = \int_{G/K}f^K\,\dd\mu_{G/K}$, 且 $\mu_{G/K}$ 由 (c) 是 $\frac{dx\,dy}{y^2}$ 的倍数, 看到 $\mu_G(\Omega'_{\varepsilon})$ 被 $I_{\varepsilon}$ 的倍数控制, 所以有限.
另一方面, 由于 $\Omega_{\varepsilon} \supset \Omega$, 投影 $\pi \colon G \to \SL_2(\Z)\backslash G$ 到 $\Omega'_{\varepsilon}$ 的限制满射; 事实上, 若 $g \in G$, 由 (d), 存在 $h \in \SL_2(\Z)$ 使得 $hg \cdot i \in \Omega$, 然后 $hg \in \Omega'_{\varepsilon}$ 且 $\pi(hg) = \pi(g)$. 再次由正则 Borel 测度的定义和 Urysohn 引理, 存在 $f \in \mathcal{C}_c^+(G)$ 使得 $f|_{\Omega'_{\varepsilon}} = 1$ 且 $\int_G f\,\dd\mu_G < +\infty$. 设 $h = \sum_{\gamma \in \SL_2(\Z)}f \in \mathcal{C}_c(\SL_2(\Z)\backslash G)$. 有 $\int_{\SL_2(\Z)\backslash G}h\,\dd\mu_{\SL_2(\Z)\backslash G} = \int_G f\,\dd\mu_G < +\infty$. 但若 $x \in \SL_2(\Z)\backslash G$, $g$ 是满足 $\pi(g) = x$ 的 $\Omega'_{\varepsilon}$ 的元素, 则
\[
h(x) = \sum_{\gamma \in \SL_2(\Z)}f(\gamma g) \geq f(g) = 1
\]
(用 $f \geq 0$ 这一事实). 所以
\[
\operatorname{vol}(\SL_2(\Z)\backslash G) \leq \int_{\SL_2(\Z)\backslash G}h\,\dd\mu_{\SL_2(\Z)\backslash G} < +\infty.
\]
(5) 设 $V := \{f \in L^2(X) \mid \int_X f(x)\,\dd\mu(x) = 0\}$. 这是 $L^2(X)$ 的 $G$-不变子空间, 因为测度 $\mu$ 被 $G$ 左不变. 证明 $V^G = 0$. 设 $f \in L^2(X)^G$. 则对每个 $g \in G$ 有 $\|R_gf - f\|_{L^2(X)} = 0$. 把 Minkowski 不等式 (见习题 \ref{ex:I.5.6.7}) 应用于函数 $\phi \colon G \times X \to \CC$, $(g, x) \longmapsto |f(xg) - f(x)|$, 看到
\[
\int_X\Bigl(\int_G|f(xg) - f(x)|\,\dd g\Bigr)^2\,\dd x = 0,
\]
所以函数 $x \longmapsto \int_G|f(xg) - f(x)|\,\dd g$ 在 $x$ 上几乎处处为 $0$. 选择满足 $\int_G|f(x_0g) - f(x_0)|\,\dd g = 0$ 的 $x_0 \in X$. 则函数 $g \longmapsto f(x_0g) - f(x_0)$ 在 $G$ 上几乎处处为 $0$, 所以 $f$ 在 $L^2(X)$ 中等于常值函数 $f(x_0)$ 的类. 若此外 $f \in V$, 这迫使 $f = 0$.
现在把问题 (4) 应用于 $V$. 它的结论说对所有 $f_1, f_2 \in V$, $\lim_{g \to +\infty}\langle R_gf_1, f_2 \rangle = 0$.
设 $f_1, f_2 \in L^2(X)$. 写 $c_1 = \int_X f_1\,\dd\mu$, $c_2 = \int_X f_2\,\dd\mu$. 由于 $\operatorname{vol}(X) = 1$, 函数 $f_1 - c_1$ 和 $f_2 - c_2$ 在 $V$ 中, 所以由刚才看到的, $\lim_{g \to +\infty}\langle R_gf_1 - c_1, f_2 - c_2 \rangle = 0$.
对每个 $g \in G$,
\begin{align*}
\langle R_gf_1 - c_1, f_2 - c_2 \rangle &= \int_X (f_1(xg) - c_1)(f_2(x) - c_2)\,\dd x \\
&= \int_X f_1(xg)f_2(x)\,\dd x - c_1\int_X f_2(x)\,\dd x - c_2\int_X f_1(xg)\,\dd x + c_1c_2 \\
&= \int_X f_1(xg)f_2(x)\,\dd x - c_1c_2.
\end{align*}
最后得到
\[
\lim_{g \to +\infty}\int_X f_1(xg)f_2(x)\,\dd x = c_1c_2,
\]
如所愿.
(6)(a). 定义函数 $h \colon G \to \CC$:
\[
h(x) = \int_K\psi(xy^{-1})\,\dd\mu_K(y),
\]
对每个 $x \in G$. 证明 $h$ 连续. 设 $\varepsilon > 0$. 由于 $\psi \in \mathcal{C}_c(G)$, 命题 \ref{prop:I.1.12} 蕴含 $\psi$ 左一致连续, 所以存在 $G$ 中 $1$ 的邻域 $U$ 使得对每个 $x \in G$ 和每个 $x_0 \in U$, $|\psi(x_0x) - \psi(x)| \leq \varepsilon$. 则若 $x \in G$, $x_0 \in U$, 有
\[
|h(x_0x) - h(x)| \leq \int_K|\psi(x_0xy^{-1}) - \psi(xy^{-1})|\,\dd\mu_K(y) \leq \int_K\varepsilon\,\dd\mu_K = \varepsilon,
\]
所以 $h$ 也左一致连续 (特别地连续). 又, 由于 $\psi(xy^{-1}) = 0$ 除非 $x \in (\operatorname{supp}\psi)y \subset (\operatorname{supp}\psi)K$, $\psi$ 的支撑包含于 $(\operatorname{supp}\psi)K$, 所以紧.
最后, 对每个 $f \in \mathcal{C}_c(G)$:
\begin{align*}
\psi * \mu_K(f) &= \int_{G \times K}f(xy)\psi(x)\,\dd\mu_G(x)\,\dd\mu_K(y) \\
&= \int_{G \times K}f(x)\psi(xy^{-1})\,\dd\mu_G(x)\,\dd\mu_K(y) \quad \text{(因为 } \Delta_G|_K = 1\text{)} \\
&= \int_G f(x)h(x)\,\dd\mu_G(x).
\end{align*}
(b). 设 $(U_n)_{n \geq 0}$ 是 $G$ 中 $1$ 的递减邻域序列, 构成邻域基; 例如可以取 $U_n$ 为 $G$ 与 $M_2(\R)$ 中半径 $2^{-n}$ 的球的交, 对 $M_2(\R)$ 上任何范数. 由命题 \ref{prop:I.4.1.8}, 存在近似单位元 $(\psi_n)_{n \geq 0}$ 使得对每个 $n$ 有 $\operatorname{supp}(\psi_n) \subset U_n$.
像命题 \ref{prop:I.4.1.9} 中那样证明序列 $(\psi_n)_{n \geq 0}$ 可行. 设 $f \in \mathcal{C}_c(G)$. 则对每个 $n$, $\mu_K(f) = \int_K f(y)\,\dd\mu_K(y) = \int_{G \times K}f(y)\psi_n(x)\,\dd\mu_G(x)\,\dd\mu_K(y)$, 所以
\[
|\psi_n * \mu_K(f) - \mu_K(f)| \leq \int_{G \times K}|f(xy) - f(y)|\psi(x)\,\dd\mu_G(x)\,\dd\mu_K(y) \leq \sup_{x \in U_n}\|L_{x^{-1}}f - f\|_{\infty}.
\]
由于 $f$ 左一致连续, 当 $n \to +\infty$ 时这趋于 $0$.
(c). 由 (b), 只需对 $f \in \mathcal{C}_c(G)$ 和泛函 $\psi * \mu_K$ ($\psi \in \mathcal{C}_c(G)$) 证明 $(***)$ 的类似物, 即证明
\[
\lim_{g \to \infty}\int_X \psi * \mu_K(\cdot g)f\,\dd x = \int_X f\,\dd\mu
\]
(其中把 $\psi * \mu_K$ 与 $L^2(X)$ 的元素等同). 用 (a) 的显式 $h$, 这由问题 (5) 的等分布结果 ($\psi * \mu_K$ 现在是 $\mathcal{C}_c(G)$ 的元素 $h$) 得出, 因为 $Yg$ 在 $X$ 中的像正是 $y \cdot g$ ($y \in Y$), 而 $h * f$ 的积分极限由 (5) 给出. 详细说来, 由 (5) 应用于 $f_1 = h$ (通过 $\pi$ 下降) 和 $f_2 = f$, 看到 $\lim_{g \to \infty}\int_X h(xg)f(x)\,\dd x = \int_X h\,\dd\mu \int_X f\,\dd\mu$. 而 $\int_X h\,\dd\mu = \mu_K$ 的规范化总质量 $= 1$, 所以得到 $\lim_{g \to \infty}\int_X h(xg)f(x)\,\dd x = \int_X f\,\dd\mu$, 这正是 $(***)$ (把 $f$ 换成 $\tilde{f}$, 用 Fubini 和 $h$ 的定义把 $y$-积分与 $x$-积分交换).
\end{solution}


\chapter[Fourier 分析在群上随机游走中的应用]{Fourier 分析在群上随机游走中的应用}
% ============ 第 VI 部分: Fourier 分析在群上随机游走中的应用 ============
% (Application of Fourier Analysis to Random Walks on Groups)
% 原书 VI.1 - VI.7 (原书 pp.311-338)

在本章中, 我们主要关注有限群的情形, 但在最后一节中给出一些对更一般群的结果.

\section{有限 \index{Markov 链}Markov 链}

我们一次性固定概率空间 $\Omega$ (即总体积为 $1$ 的测度空间).

\begin{definition}[\index{随机变量}随机变量与分布]
\label{def:VI.1.1}
设 $X$ 是可测空间 (即带 $\sigma$-代数的空间). 取值于 $X$ 的随机变量 (random variable) 是可测函数 $X \colon \Omega \to X$.
对 $X$ 的每个可测子集 $A$, 记 $\mathbb{P}(X \in A)$ 为 $X^{-1}(A)$ 的测度. (我们把它看作 $X$ 在 $A$ 中的概率.) $X$ 的分布 (distribution) 是 $X$ 上由 $\mu(A) = \mathbb{P}(X \in A)$ 定义的概率分布 $\mu$.
\end{definition}

我们把随机变量看作某个实验或观察的结果. 概率空间 $\Omega$ 通常不指明 (可以把它想成类似``所有可能的宇宙''的东西). 例如, 可以把掷硬币的结果看作取值于有限集 $\{$正面, 反面$\}$ 的随机变量. 若硬币无偏, 该\index{随机变量的分布}随机变量的分布由 $\mu(\{$正面$\}) = \mu(\{$反面$\}) = \frac{1}{2}$ 给出.
在这些讲义中, 我们只关心 $X$ 有限且其 $\sigma$-代数为 $X$ 的所有子集的集合的情形. 可以 (而且将要) 把 $X$ 上的测度看作函数 $\mu \colon X \to \R_{\geq 0}$.
从现在起, 假设 $X$ 有限.

\begin{definition}[\index{随机矩阵}随机矩阵 (stochastic matrix)]
\label{def:VI.1.2}
矩阵 $P = (P_{ij}) \in M_n(\R)$ 称为随机的 (stochastic), 若对所有 $i, j \in \{1, \ldots, n\}$ 有 $P_{ij} \geq 0$ 且对每个 $i \in \{1, \ldots, n\}$ 有 $\sum_{j=1}^n P_{ij} = 1$.
若 $P \colon X \times X \to \R$ 是函数, 把它看作大小为 $|X| \times |X|$ 的矩阵, 若对所有 $x, y \in X$ 有 $P(x, y) \geq 0$ 且对每个 $x \in X$ 有 $\sum_{y \in X}P(x, y) = 1$, 则称它随机.
\end{definition}

\begin{definition}[Markov 链 (Markov chain)]
\label{def:VI.1.3}
设 $P \colon X^2 \to \R$ 是随机函数, $\nu$ 是 $X$ 上的概率分布. 状态空间为 $X$、初始分布为 $\nu$、转移矩阵 (transition matrix) 为 $P$ 的 (离散时间齐次) Markov 链是取值于 $X$ 的随机变量序列 $(X_n)_{n \geq 0}$, 使得:
\begin{enumerate}
\item[(a)] $X_0$ 的分布是 $\nu$.
\item[(b)] 对每个 $n \geq 0$ 和所有 $x_0, \ldots, x_{n+1} \in X$, 若 $\mathbb{P}(X_n = x_n, \ldots, X_0 = x_0) > 0$, 则
\[
\mathbb{P}(X_{n+1} = x_{n+1} \mid X_n = x_n, \ldots, X_0 = x_0) = P(x_n, x_{n+1}).
\]
\end{enumerate}
\end{definition}

设 $P, Q \colon X^2 \to \R$ 是两个函数. 记 $PQ$ 为 $X^2 \to \R$ 的函数:
\[
PQ(x, y) = \sum_{z \in X}P(x, z)Q(z, y).
\]
(若把 $X^2$ 上的函数看作矩阵, 这就是通常的矩阵乘积.)
特别地, 记 $P^n$ 为乘积 $PP\cdots P$ ($n$ 次); 按约定, $P^0$ 是对角线的特征函数.

\begin{lemma}[$X_n$ 的分布]
\label{lem:VI.1.4}
设 $(X_n)_{n \geq 0}$ 是 $X$ 上初始分布为 $\nu$、转移矩阵为 $P$ 的 Markov 链. 则对每个 $x \in X$,
\[
\mathbb{P}(X_n = x) = \sum_{y \in X}\nu(y)P^n(y, x).
\]
换句话说, 若把 $P$ 看作矩阵、$\nu$ 看作行向量, 则 $X_n$ 的分布是 $\nu P^n$.
\end{lemma}

\begin{proof}
对 $n$ 归纳证明. $n = 0$ 时显然. 假设对某个 $n$ 已知, 证明对 $n+1$. 设 $x \in X$. 则
\begin{align*}
\mathbb{P}(X_{n+1} = x) &= \sum_{y \in X, \mathbb{P}(X_n = y) \neq 0}\mathbb{P}(X_{n+1} = x \mid X_n = y) \\
&= \sum_{y \in X}\mathbb{P}(X_n = y)\mathbb{P}(X_{n+1} = x \mid X_n = y) \\
&= \sum_{y \in X}\mathbb{P}(X_n = y)P(y, x).
\end{align*}
用归纳假设, 得到
\[
\mathbb{P}(X_{n+1} = x) = \sum_{y \in X}P(y, x)\sum_{z \in X}\nu(z)P^n(z, y) = \sum_{z \in X}\nu(z)P^{n+1}(z, x).
\]
\end{proof}

\begin{example}[Markov 链的例子]
\label{exa:VI.1.5}
(1) 离散圆上的随机游走 (random walk on the discrete circle): 取 $X = \Z/r\Z$, $\nu = \delta_0$, $P$ 由
\[
P(x, y) = \begin{cases} \frac{1}{2} & \text{若 } x - y \in \{\pm 1\}, \\ 0 & \text{否则} \end{cases}
\]
定义. 此 Markov 链在``离散圆'' $\Z/n\Z$ 上建模随机游走, 以概率 $1$ 从 $0$ 出发, 然后每个时刻 $n$ 有 $50\%$ 的概率走到离散圆上的前一个点, $50\%$ 的概率走到下一个点.
(2) 用随机对换洗牌 (mixing a deck of cards using random transpositions): 我们试图理解如下情形: 有一副 $N$ 张牌. 每个时刻 $n$, 随机 (均匀且独立地) 选两张牌并交换它们在牌组中的位置. 需要多久才能洗匀?
此问题用状态空间为 $S_N$ (表示牌组所有可能的排列), 初始分布为集中在起始位置上的 \index{Dirac 测度}Dirac 测度, 转移矩阵 $P$ 由
\[
P(\tau\sigma, \sigma) = \begin{cases} \frac{1}{N} & \text{若 } \tau = 1, \\ \frac{2}{N^2} & \text{若 } \tau \text{ 是对换}, \\ 0 & \text{否则} \end{cases}
\]
给出的 Markov 链建模.
(3) Bernoulli-Laplace 扩散模型 (Bernoulli-Laplace diffusion model): 有标号 $0$ 和 $1$ 的两个瓮. 开始时, 瓮 $0$ 装有 $r$ 个红球, 瓮 $1$ 装有 $b$ 个蓝球. 每个时刻 $n$, 在每个瓮中 (均匀且独立地) 选一个球并交换. 需要多久才能混匀?
用状态空间为 $S_N/S_r \times S_b$ (其中 $N = r + b$, $S_r \times S_b$ 通过显然的双射 $\{1, \ldots, r\} \times \{1, \ldots, b\} \simeq \{1, \ldots, N\}$ 嵌入 $S_N$) 的 Markov 链建模. 事实上, 可以把 $N$ 个球看作集合 $\{1, \ldots, N\}$, 其中前 $r$ 个球红、后 $b$ 个球蓝. 上述过程的状态是满足 $|A| = r$ 的 $\{1, \ldots, N\}$ 的子集 $A$ (瓮 $0$ 的内容); 注意在两个瓮之间交换两个球不改变每个瓮中的球数. 群 $S_N$ 在 $\{1, \ldots, N\}$ 的基数为 $r$ 的子集的集合 $\Omega_r$ 上传递作用, 其子群 $S_r \times S_b$ 是 $\{1, \ldots, r\}$ 的稳定化子, 所以状态集确实与 $S_N/S_r \times S_b$ 双射对应. 初始分布是集中在 $\{1, \ldots, r\}$ 上的 Dirac 测度. 转移矩阵 $P$ 由
\[
P(A', A) = \begin{cases} \frac{(r-1)!(b-1)!}{(r+b)!} & \text{若 } r - |A \cap A'| = 1, \\ 0 & \text{否则} \end{cases}
\]
给出. 事实上, 需要计算满足 $r - |A \cap A'| = 1$ 的 $\{1, \ldots, N\}$ 的基数为 $r$ 的子集偶对 $(A, A')$ 的数目; 注意此条件意味着 $A' - A$ 和 $A - A'$ 都恰好有一个元素. 有 $\frac{(r+b)!}{r!b!}$ 种选择 $A$ 的方式, $b$ 种选择 $A' - A$ 的元素的方式, $r$ 种选择 $A - A'$ 的元素的方式.
\end{example}

我们一直在问例子中的链是否收敛, 但第一个问题应该是: 它们能收敛到什么分布?

\begin{definition}[平稳分布 (stationary distribution)]
\label{def:VI.1.6}
考虑随机函数 $P \colon X^2 \to \R$. $P$ 的平稳分布是 $X$ 上满足对每个 $y \in X$ 有
\[
\mu(y) = \sum_{x \in X}\mu(x)P(x, y)
\]
的概率分布 $\mu$.
若把 $P$ 看作 $|X| \times |X|$ 矩阵、$\mu$ 看作大小为 $|X|$ 的行向量, 则条件变成 $\mu P = \mu$.
若转移矩阵为 $P$ 的 Markov 链以任何合理意义收敛, 则其极限的分布应是 $P$ 的平稳分布.
最后, 定义将对随机变量使用的距离. 注意若 $X$ 是一般测度空间, 此定义同样有意义, 且后面的引理在本质上用同样的证明仍然成立.
\end{definition}

\begin{definition}[\index{全变差距离}全变差距离 (total variation distance)]
\label{def:VI.1.7}
设 $\mu$ 和 $\nu$ 是 $X$ 上的两个概率分布. 它们的全变差距离 (total variation distance) 是
\[
\|\mu - \nu\|_{TV} = \max_{A \subset X}|\mu(A) - \nu(A)|.
\]
这显然是概率分布集合上的度量, 事实上它与 $L^1$ 度量密切相关.
\end{definition}

\begin{lemma}[全变差距离的公式]
\label{lem:VI.1.8}
设 $\mu$ 和 $\nu$ 是 $X$ 上的两个概率分布. 则
\[
\|\mu - \nu\|_{TV} = \frac{1}{2}\sum_{x \in X}|\mu(x) - \nu(x)|.
\]
\end{lemma}

\begin{proof}
设 $B = \{x \in X \mid \mu(x) \geq \nu(x)\}$. 对每个 $A \subset X$,
\begin{multline*}
\mu(A) - \nu(A) = \mu(A \cap B) - \nu(A \cap B) + \sum_{x \in A - A \cap B}(\mu(x) - \nu(x)) \leq \mu(A \cap B) - \nu(A \cap B) \\
= \mu(B) - \nu(B) - \sum_{x \in B - A \cap B}(\mu(x) - \nu(x)) \leq \mu(B) - \nu(B).
\end{multline*}
类似地,
\[
\nu(A) - \mu(A) \leq \nu(X - B) - \mu(X - B) = \mu(B) - \nu(B).
\]
所以 $|\mu(A) - \nu(A)| \leq \mu(B) - \nu(B)$, 且 $A = B$ 或 $A = X - B$ 时取等号, 得到
\begin{align*}
\|\mu - \nu\|_{TV} &= \mu(B) - \nu(B) = \frac{1}{2}(\mu(B) - \nu(B) + \nu(X - B) - \mu(X - B)) \\
&= \frac{1}{2}\sum_{x \in B}|\mu(x) - \nu(x)| + \frac{1}{2}\sum_{x \in X - B}|\mu(x) - \nu(x)| \\
&= \frac{1}{2}\|\mu - \nu\|_1.
\end{align*}
\end{proof}

\section{Perron-Frobenius 定理与 Markov 链的收敛}

\begin{notation}[矩阵的序]
\label{not:VI.2.1}
设 $A, B \in M_{nm}(\R)$. 若对每个 $(i, j) \in \{1, \ldots, n\} \times \{1, \ldots, m\}$ 有 $A_{ij} \geq B_{ij}$ (相应地 $A_{ij} > B_{ij}$), 则说 $A \geq B$ (相应地 $A > B$). 还用 $|A|$ 记 $n \times m$ 矩阵 $(|A_{ij}|)$.
\end{notation}

\begin{definition}[\index{正矩阵}正矩阵与遍历矩阵]
\label{def:VI.2.2}
若 $P > 0$, 则说矩阵 $P = (P_{ij}) \in M_n(\R)$ 正 (positive).
\end{definition}

\begin{definition}[遍历矩阵 (ergodic matrix)]
\label{def:VI.2.3}
若存在正整数 $r$ 使得 $P^r$ 正, 则说随机矩阵 $P \in M_n(\R)$ 遍历 (ergodic).
\end{definition}

回忆线性代数中下列经典定理:

\begin{theorem}[Perron-Frobenius 定理]
\label{thm:VI.2.4}
设 $P = (P_{ij}) \in M_n(\R)$ 是遍历随机矩阵. 则:
\begin{enumerate}
\item[(i)] 矩阵 $P$ 以 $1$ 为单特征值, 且 $P$ 的每个复特征值 $\lambda$ 满足 $|\lambda| < 1$.
\item[(ii)] 满足 $wP = w$ 的行向量 $w \in M_{1n}(\R)$ 的空间是 $1$ 维的, 且它有生成元 $v = (v_1, \ldots, v_n)$ 使得对每个 $i$ 有 $v_i > 0$ 且 $v_1 + \cdots + v_n = 1$.
\item[(iii)] 设 $P_{\infty}$ 是行都等于 (ii) 的向量 $v$ 的 $n \times n$ 矩阵. 则 $r \to +\infty$ 时 $P^r \to P_{\infty}$. 更精确地, 设 $\rho = \max\{|\lambda| \mid \lambda \neq 1 \text{ 是 } P \text{ 的特征值}\}$; 由 (i) 知道 $\rho < 1$. 固定 $M_n(\R)$ 上的任何范数 $\|\cdot\|$. 则存在多项式 $f \in \Z[t]$ 使得
\[
\|P^k - P_{\infty}\| \leq f(k)\rho^k.
\]
\end{enumerate}
\end{theorem}

\begin{lemma}[Perron-Frobenius 引理]
\label{lem:VI.2.5}
设 $A = (A_{ij}) \in M_n(\R)$ 是正矩阵, 设
\[
Z = \{x = (x_1, \ldots, x_n) \in \R^n \mid x \geq 0 \text{ 且 } x_1 + \cdots + x_n = 1\},
\]
设
\[
\Lambda = \{\lambda \in \R \mid \exists x \in Z,\, Ax \geq \lambda x\}.
\]
则实数 $\lambda_0 = \sup\Lambda$ 为正, 是 $A$ 的特征多项式的单根, 且它有一个全元素为正的特征向量. 又, 对 $A$ 的每个复特征值 $\lambda \neq \lambda_0$, $|\lambda| < \lambda_0$.
\end{lemma}

\begin{proof}
注意 $\Lambda \neq \emptyset$, 因为 $0 \in \Lambda$, 且 $\Lambda$ 被 $A$ 的所有元素之和控制上界. 所以 $\lambda_0$ 良定义且非负. 设 $(\mu_n)_{n \geq 0}$ 是 $\Lambda$ 中收敛到 $\lambda_0$ 的序列; 对每个 $n \geq 0$, 选择 $x^{(n)} \in Z$ 使得 $Ax^{(n)} \geq \mu_nx^{(n)}$. 由于 $Z$ 紧, 可以假设序列 $(x^{(n)})_{n \geq 0}$ 收敛到某个 $x \in Z$, 然后有 $Ax \geq \lambda_0x$. 假设 $Ax \neq \lambda_0x$, 则因 $A > 0$ 得到 $A(Ax) > \lambda_0Ax$. 由于 $Ax \geq 0$ 且 $Ax \neq 0$, 可以把 $Ax$ 乘以正纯量得到向量 $y \in Z$ 使得 $Ay > \lambda_0y$, 这与 $\lambda_0$ 的定义矛盾. 所以 $Ax = \lambda_0x$. 又, 由于 $x$ 至少有一个正元素且 $A > 0$, 向量 $\lambda_0x = Ax$ 的所有元素为正, 这蕴含 $\lambda_0 > 0$ 且 $x > 0$.
其次证明 $A$ 的每个复特征值 $\lambda \neq \lambda_0$ 满足 $|\lambda| < \lambda_0$. 设 $\lambda$ 是 $A$ 的复特征值. 则存在非零向量 $y = (y_1, \ldots, y_n) \in \CC^n$ 使得 $Ay = \lambda y$. 对每个 $i \in \{1, \ldots, n\}$,
\[
|\lambda||y_i| = \Bigl|\sum_{j=1}^n A_{ij}y_j\Bigr| \leq \sum_{j=1}^n A_{ij}|y_j|.
\]
换句话说, 有 $A|y| \geq |\lambda||y|$. 由于可以把 $|y|$ 规范化得到 $Z$ 的元素, 这证明 $|\lambda| \leq \lambda_0$. 假设 $|\lambda| = \lambda_0$. 由于 $A > 0$, 存在正实数 $\delta$ 使得 $A_0 := A - \delta I_n > 0$. 则 $\mu \longmapsto \mu - \delta$ 诱导 $A$ 的特征值与 $A_0$ 的特征值之间的双射, 特别地 $\lambda_0 - \delta$ 是 $A_0$ 的最大实特征值 (且它为正, 因为 $A_0 > 0$). 把本段开头应用于 $A_0$, 看到 $|\lambda - \delta| \leq \lambda_0 - \delta$. 但此时
\[
\lambda_0 = |\lambda| = |\lambda - \delta + \delta| \leq |\lambda - \delta| + \delta \leq \lambda_0,
\]
所以 $|\lambda - \delta| + \delta = |\lambda|$, 因此 $\lambda \in \R_{\geq \delta}$, 且必有 $\lambda = \lambda_0$.
证明特征空间 $E_{\lambda_0} := \Ker(A - \lambda_0I_n)$ 维数为 $1$. 假设存在 $y = (y_1, \ldots, y_n) \in E_{\lambda_0}$ (元素为实数) 使得族 $\{x, y\}$ 线性无关. 可以假设 $y$ 至少有一个正元素. 写 $x = (x_1, \ldots, x_n)$, 设 $\mu = \sup\{\nu \in \R \mid \forall i \in \{1, \ldots, n\},\, x_i \geq \nu y_i\}$. 则 $x - \mu y \geq 0$ 且 $x - \mu y \ngeq 0$. 向量 $x - \mu y$ 非零, 因为 $x$ 和 $y$ 线性无关, 且有 $A(x - \mu y) = \lambda_0(x - \mu y)$. 由于 $A > 0$, $x - \mu y \geq 0$, $\lambda_0 > 0$, 这蕴含 $x - \mu y > 0$, 与 $\nu$ 的选取矛盾.
现在证明 $\lambda_0$ 是 $A$ 的特征多项式 $\chi_A(t)$ 的单根. 可以找到 $g \in \GL_n(\R)$ 使得 $g^{-1}Ag$ 具有 $\begin{pmatrix} \lambda_0 & * \\ 0 & B \end{pmatrix}$ 的形式, 其中 $B \in M_{n-1}(\R)$. 有 $\chi_A(t) = (t - \lambda_0)\chi_B(t)$. 假设 $\lambda_0$ 作为 $\chi_A(t)$ 的根的重数 $\geq 2$. 则 $\lambda_0$ 是 $\chi_B(t)$ 的根, 所以存在 $z \in \R^{n-1}$ 使得 $Bz = \lambda_0z$. 设 $y = g\begin{pmatrix} 0 \\ z \end{pmatrix} \in \R^n$, 则对某个 $\alpha \in \R$ 有 $Ay = \lambda_0y + \alpha x$. 由于 $\dim(E_{\lambda_0}) \neq 0$ 是 $1$, 向量 $y$ 不可能是 $A$ 的特征向量, 所以 $\alpha \neq 0$. 简单归纳 (用 $Ax = \lambda_0x$ 这一事实) 表明对每个正整数 $r$, $A^ry = \lambda_0^ry + r\alpha\lambda_0^{r-1}x$. 由于 $A^r > 0$, 这蕴含
\[
A^r|y| \geq |A^ry| = |\lambda_0^ry + r\alpha\lambda_0^{r-1}x| \geq |r\alpha\lambda_0^{r-1}x| - \lambda_0^r|y| = \lambda_0^{r-1}(r|\alpha x| - \lambda_0|y|).
\]
由于 $\alpha \neq 0$ 且 $x > 0$, 存在正整数 $r$ 使得 $r|\alpha x| - \lambda_0|y| > \lambda_0|y|$, 然后 $A^r|y| > \lambda_0^r|y|$. 由于 $A^r > 0$, 把本证明的开头应用于 $A^r$, 看到这蕴含 $A^r$ 有实特征值 $> \lambda_0^r$. 但这不可能, 因为 $A^r$ 的特征值是 $A$ 的特征值的 $r$ 次幂, 所以它们的绝对值都 $\leq \lambda_0^r$.
\end{proof}

\begin{proof}[定理的证明]
证明 (i). 设 $v_0 = (1, \ldots, 1) \in \R^n$. 则 $P$ 随机这一事实等价于 $P \geq 0$ 且 $Pv_0 = v_0$. 由于所有矩阵 $P^r$ ($r \geq 1$) 的元素非负且满足 $P^rv_0 = v_0$, 它们都随机. 又, 若 $x = (x_1, \ldots, x_n) \in (\R_{\geq 0})^n$, $Q = (Q_{ij}) \in M_n(\R)$ 随机, 则对每个 $i \in \{1, \ldots, n\}$,
\[
(Qx)_i = \sum_{j=1}^n Q_{ij}x_j \leq \sup_{1 \leq j \leq n}x_j.
\]
固定整数 $r \geq 1$ 使得 $P^r > 0$. 由引理, 矩阵 $P^r$ 有简单实正特征值 $\lambda_0$, 使得 $P^r$ 的每个复特征值 $\lambda \neq \lambda_0$ 满足 $|\lambda| < \lambda_0$. 由引理中 $\lambda_0$ 的定义和上面关于随机矩阵的观察, 有 $\lambda_0 \leq 1$. 另一方面, $Pv_0 = v_0$, 所以 $1$ 是 $P$ 的特征值, 从而也是 $P^r$ 的, 因此 $\lambda_0 = 1$. 设 $\lambda \in \CC$ 是 $P$ 的特征值, $y \in \CC^n$ 是此特征值的特征向量. 则 $P^ry = \lambda^ry$, 所以 $\lambda^r$ 是 $P^r$ 的特征值. 若 $\lambda^r \neq 1$, 则由引理 $|\lambda^r| < 1$, 所以 $|\lambda| < 1$. 若 $\lambda^r = 1$, 则特征向量 $y$ 必在 $\Ker(P^r - I_n)$ 中, 且知道 (再次由引理) 这个空间是 $1$ 维的. 由于 $v_0 \in \Ker(P^r - I_n)$, 向量 $y$ 必是 $v_0$ 的倍数, 然后 $\lambda = 1$.
最后, 若 $P$ 的特征多项式是 $\chi_P(t) = (t - \lambda_1)\cdots(t - \lambda_n)$, 则 $P^r$ 的特征多项式是 $\chi_{P^r}(t) = (t - \lambda_1^r)\cdots(t - \lambda_n^r)$. 所以 $1$ 在 $\chi_P(t)$ 中的重数至多是它在 $\chi_{P^r}(t)$ 中的重数, 而由引理知道后者是 $1$. 这完成 (i) 的证明.
证明 (ii). 由于 $P$ 和 $P^T$ 有相同的特征多项式, 由 (i) 知道 $1$ 是 $P^T$ 的单特征值, 所以满足 $wP = w$ 的行向量空间维数为 $1$. 设 $w = (w_1, \ldots, w_n)$ 是这个空间的非零向量. 则也有 $|w|P = |w|$. 事实上, 对每个 $j \in \{1, \ldots, n\}$,
\[
|w_j| = \Bigl|\sum_{i=1}^n w_iP_{ij}\Bigr| \leq \sum_{i=1}^n |w_i|P_{ij}
\]
(因为所有 $P_{ij}$ 非负). 假设 $|w| \neq |w|P$. 则存在 $j_0 \in \{1, \ldots, n\}$ 使得 $|w_{j_0}| < \sum_{i=1}^n |w_i|P_{ij_0}$, 这蕴含
\[
\sum_{i=1}^n |w_i| = \sum_{i=1}^n\sum_{j=1}^n |w_i|P_{ij} > \sum_{j=1}^n |w_j|,
\]
矛盾. 由于 $w \neq 0$, $|w_i|$ 中至少有一个为正. 若如前选择 $r \geq 1$ 使得 $P^r > 0$, 则 $|w| = P^r|w|$, 所以对每个 $j \in \{1, \ldots, n\}$, $|w_j| = \sum_{i=1}^n (P^r)_{ij}|w_i| > 0$. 这完成 (ii) 的证明.
最后证明 (iii). 由于 $M_n(\R)$ 上的所有范数等价, 只需对某个特定范数证明结论. 由 Jordan 标准形的存在性 (实际上是 Jordan-Chevalley 分解), 存在矩阵 $g \in \GL_n(\R)$ 使得 $g^{-1}Pg = A = \begin{pmatrix} 1 & 0 \\ 0 & B \end{pmatrix}$, 其中 $B \in M_{n-1}(\R)$ 满足 $B = D + N$, $D$ 是对角矩阵, $N$ 是幂零矩阵且 $DN = ND$. 在 $M_n(\R)$ 上选择由 $\R^n$ 上通常 Euclid 范数给出的算子范数 $\|\cdot\|$. $D$ 的元素是 $P$ 的异于 $1$ 的特征值, 所以 $\|D\| = \rho$. 由于 $D$ 和 $N$ 交换, 对每个 $k \in \Z_{\geq 0}$,
\[
B^k = (D + N)^k = \sum_{j=0}^k\binom{k}{j}D^{k-j}N^j.
\]
若 $k \geq n$ (事实上 $k \geq n - 1$ 就够), 这化简为 $\sum_{j=0}^n\binom{k}{j}D^{k-j}N^j$, 因为 $j \geq n$ 时 $N^j = 0$. 所以若 $k \geq n$,
\[
\|B^k\| \leq \sum_{j=0}^n\binom{k}{j}\|D\|^{k-j}\|N\|^j \leq \rho^{k-n}\sum_{j=0}^n\binom{k}{j}\|N\|^j.
\]
设 $A_{\infty} = \begin{pmatrix} 1 & 0 \\ 0 & B_{\infty} \end{pmatrix}$, 其中 $B_{\infty} = 0 \in M_{n-1}(\R)$. 则对每个 $k \geq 0$ 有 $\|A^k - A_{\infty}\| = \|B^k\|$, 所以 $A^k \to A_{\infty}$ (因为 $\rho < 1$). 这蕴含 $P^k \to P_0 := gA_{\infty}g^{-1}$ (当 $k \to +\infty$). 又,
\[
\|P^k - P_0\| = \|g^{-1}(A^k - A_{\infty})g\| \leq \|g\|\|g^{-1}\|\|B^k\|,
\]
它被 $\rho^k$ 和 $k$ 的多项式的乘积控制. 所以只需证明 $P_0 = P_{\infty}$. 由于 $P_0 = \lim_{k \to +\infty}P^k$, 有 $P_0P = PP_0 = P_0$. 回忆 $1$ 是 $P$ 和 $P^T$ 的单特征值. 所以 $P_0$ 的所有行都是 $P^T$ 的对应特征向量的倍数, 即 $v$ 的倍数. 又, 由于对每个 $k \geq 0$ 有 $P^k$ 随机, 其极限 $P_0$ 随机. 所以 $P_0$ 的所有行的元素非负且和为 $1$, 这意味着它们都等于 $v$, 且 $P_0 = P_{\infty}$.
\end{proof}

\begin{definition}[\index{遍历 Markov 链}遍历 Markov 链]
\label{def:VI.2.6}
若 $P$ 遍历, 则称转移矩阵为 $P$ 的 Markov 链遍历 (ergodic).
\end{definition}

在例子 \ref{exa:VI.1.5} 中, 除了 (1) 中 $r$ 为偶数时的 Markov 链, 所有链都遍历.

\begin{corollary}[遍历链的收敛]
\label{cor:VI.2.7}
设 $(X_n)_{n \geq 0}$ 是转移矩阵为 $P$ 的遍历 Markov 链. 则 $P$ 有唯一平稳分布 $\mu$, 且若 $\mu_n$ 是 $X_n$ 的分布, 有
\[
\|\mu_n - \mu\|_{TV} \leq f(n)\rho^n,
\]
其中 $f$ 是多项式, $\rho = \max\{|\lambda| \mid \lambda \neq 1 \text{ 是 } P \text{ 的特征值}\} < 1$.
\end{corollary}

\begin{proof}
设 $\nu$ 是 Markov 链的初始分布. 由引理 \ref{lem:VI.1.4}, $\mu_n = \nu P^n$. 设 $P_{\infty}$ 是序列 $(P^n)_{n \geq 0}$ 的极限. $P_{\infty}$ 的所有行都等于 $\mu$, 所以 $\nu P_{\infty} = \mu$. 若用 $X$ 到 $\R$ 的函数空间上的 $L^1$ 范数 (对 $X$ 上的计数测度) 定义矩阵空间上的算子范数 $\|\cdot\|$, 则
\[
\|\mu_n - \mu\|_{TV} = \frac{1}{2}\|\mu_n - \mu\|_1 = \frac{1}{2}\|\nu P^n - \nu P_{\infty}\|_1 \leq \frac{1}{2}\|\nu\|_1\|P^n - P_{\infty}\|,
\]
所以 $\|\mu_n - \mu\|_{TV}$ 的界立即由定理的 (iii) 得出.
\end{proof}

\begin{remark}[收敛界的缺陷]
\label{rem:VI.2.8}
虽然 $\|\mu_n - \mu\|_{TV}$ 的界看起来相当好 (它是指数的), 但若想知道 $\mu_n$ 何时恰好接近平稳分布, 它毫无用处. 需要更仔细地分析问题才能回答这类问题. 这就是我们现在将在一些特殊情形中尝试做的.
\end{remark}

\begin{example}[随机对换链遍历]
\label{exa:VI.2.9}
例子 \ref{exa:VI.1.5}(2) 的链遍历. 事实上, 设 $T \subset S_n$ 是 $\{1\}$ 和对换集合的并. 则对 $r \geq 1$ 和 $\sigma, \sigma' \in S_n$, $P^r(\sigma', \sigma) > 0$ 当且仅当 $\sigma'\sigma^{-1}$ 可以写成 $T$ 的恰好 $r$ 个元素的乘积; 由于 $1 \in T$, 这等价于条件: $\sigma'\sigma^{-1}$ 可以写成 $s$ 个对换的乘积, 其中 $s \leq r$. 所以若 $r \geq \frac{n(n-1)}{2}$ ($S_n$ 的最长元素的长度), 则对所有 $\sigma, \sigma' \in S_n$, $P^r(\sigma', \sigma) > 0$.
\end{example}

\section{遍历性判据}
本节的定义和结果在后续各节中不使用.
回忆下列定义:

\begin{definition}[图 (graph)]
\label{def:VI.3.1}
(有限无向) 图是偶对 $\mathcal{G} = (X, E)$, 其中 $X$ 是有限集, $E$ 是 $X$ 的不同元素的无序偶对 $\{x, y\}$ 的集合. 说 $X$ 是 $\mathcal{G}$ 的顶点集 (set of vertices), $E$ 是边集 (set of edges).
设 $x, y \in X$. $\mathcal{G}$ 中连接 $x$ 和 $y$ 的路径 (path) 是 $\mathcal{G}$ 的边序列 $p = (e_0, \ldots, e_n)$, 使得可以把 $e_i$ 写成 $e_i = \{x_i, y_i\}$, 其中 $x_0 = x$, $y_n = y$, 对每个 $i \in \{0, \ldots, n-1\}$ 有 $y_i = x_{i+1}$. 称整数 $n + 1$ 为路径 $p$ 的长度 (length), 记为 $|p|$. 若 $x = y$, 则说路径是闭路径 (closed path) 或基于 $x$ 的回路 (loop).
说图 $\mathcal{G}$ 连通 (connected), 若对每个 $x, y \in X$, 存在连接 $x$ 和 $y$ 的路径. 说 $\mathcal{G}$ 二部 (bipartite), 若存在满射 $\varphi \colon X \to \{-1, 1\}$ 使得对 $\mathcal{G}$ 的每条边 $e = \{x, y\}$, $\varphi(x) \neq \varphi(y)$. (换句话说, 可以把 $X$ 分成两个非空子集 $X_0$ 和 $X_1$, 使得每条边连接 $X_0$ 的元素和 $X_1$ 的元素.)
对每个 $x, y \in X$, $x$ 和 $y$ 之间的距离 (distance) $d(x, y)$ 是连接 $x$ 和 $y$ 的最短路径的长度; 若不存在这样的路径, 则设 $d(x, y) = +\infty$. 注意若 $\mathcal{G}$ 连通, 这定义 $X$ 上的度量.
\end{definition}

下列结果是经典的.

\begin{proposition}[二部图判据]
\label{prop:VI.3.2}
设 $\mathcal{G} = (X, E)$ 是连通图, 满足 $|X| \geq 2$. 则下列条件等价:
\begin{enumerate}
\item[(i)] $\mathcal{G}$ 二部;
\item[(ii)] $\mathcal{G}$ 中的每个回路长度偶;
\item[(iii)] 存在 $x_0$ 使得基于 $x_0$ 的每个回路长度偶.
\end{enumerate}
\end{proposition}

\begin{proof}
证明 (i) 蕴含 (ii). 假设 $\mathcal{G}$ 二部, $\varphi \colon X \to \{-1, 1\}$ 如定义. 设 $(e_0, \ldots, e_n)$ 是 $\mathcal{G}$ 中的回路. 写 $e_i = \{x_i, y_i\}$, 其中 $0 \leq i \leq n - 1$ 时 $x_i = y_{i+1}$, $y_n = x_0$. 则对 $i$ 简单归纳表明: $i$ 偶时 $\varphi(x_i) = \varphi(x_0)$ 且 $\varphi(y_i) \neq \varphi(x_0)$; $i$ 奇时 $\varphi(x_i) \neq \varphi(x_0)$ 且 $\varphi(y_i) = \varphi(x_0)$. 但 $y_n = x_0$, 所以 $\varphi(y_n) = \varphi(x_0)$, 因此 $n$ 奇, 回路长度偶.
(ii) 蕴含 (iii) 显然. 现在假设 (iii), 证明 (i). 选择 $x_0 \in X$ 使得基于 $x_0$ 的每个回路长度偶. 要定义函数 $\varphi \colon X \to \{0, 1\}$. 设 $y \in X$. 由于 $\mathcal{G}$ 连通, 存在连接 $x_0$ 和 $x$ 的路径 $p = (e_0, \ldots, e_n)$, 设 $\varphi(x) = (-1)^{|p|}$. 需要证明这不依赖路径. 设 $q = (f_0, \ldots, f_m)$ 是连接 $x_0$ 和 $x$ 的另一条路径. 则 $(e_0, \ldots, e_n, f_m, \ldots, f_0)$ 是基于 $x_0$ 的回路, 所以由假设它长度偶, 因此 $|p| + |q|$ 偶, 且 $(-1)^{|p|} = (-1)^{|q|}$. 注意 $\varphi(x_0) = 1$, 且若 $\{x_0, x\}$ 是边 (这样的边必存在, 因为 $\mathcal{G}$ 连通且 $|X| \geq 2$), 则 $\varphi(x) = -1$. 所以 $\varphi$ 满射. 设 $e = \{x, y\}$ 是 $\mathcal{G}$ 的边. 设 $p = (e_0, \ldots, e_n)$ 是连接 $x_0$ 和 $x$ 的路径. 则 $p' := (e_0, \ldots, e_n, e)$ 是连接 $x_0$ 和 $y$ 的路径, $|p'| = |p| + 1$, 所以 $\varphi(x) \neq \varphi(y)$. 这证明 $\mathcal{G}$ 二部.
\end{proof}

现在来看与 Markov 链的联系.

\begin{proposition}[图的连通性与遍历性]
\label{prop:VI.3.3}
设 $X$ 是有限集, $P \colon X \times X \to \R$ 是随机函数. 按如下方式定义图 $\mathcal{G} = (X, E)$: $X$ 的不同元素偶对 $\{x, y\}$ 是 $\mathcal{G}$ 的边当且仅当 $P(x, y) > 0$.
假设 $\mathcal{G}$ 连通且不二部. 则函数 $P$ 遍历.
\end{proposition}

\begin{proof}
注意对每个 $x, y \in X$ 和每个 $n \geq 1$, $P^n(x, y) > 0$ 当且仅当存在连接 $x$ 和 $y$ 的长度为 $n$ 的路径.
由命题 \ref{prop:VI.3.2}, 对每个 $x \in X$, 存在基于 $x$ 的奇数长度回路 $p_x$. 写 $2m + 1 = \max_{x \in X}|p_x|$, 其中 $m \in \Z_{\geq 0}$. 设 $x \in X$. 证明对每个 $n \geq 2m$, 存在基于 $x$ 的长度为 $n$ 的回路. 设 $\{x, z\}$ 是边. 对每个 $r \geq 0$, 记 $q_{2r}$ 为长度 $2r$ 的回路 $q_{2r} = (\{x, z\}, \{z, x\}, \ldots, \{x, z\}, \{z, x\})$. 设 $n \geq 2m$. 若 $n$ 偶, 则 $q_n$ 是基于 $x$ 的长度为 $n$ 的回路. 若 $n$ 奇, 则 $r := \frac{n - |p_x|}{2}$ 是非负整数, 且把 $p_x$ 和 $q_{2r}$ 拼接得到的回路长度 $n$ 且包含 $x$.
设 $\delta = \max_{x,y \in X}d(x, y)$. (这称为图 $\mathcal{G}$ 的直径 (diameter).) 设 $x, y \in X$, $n \geq 2m + \delta$, 证明存在连接 $x$ 和 $y$ 的长度为 $n$ 的路径 (这将完成证明). 设 $p$ 是连接 $x$ 和 $y$ 的任何路径. 则 $|p| \leq \delta$, 所以由上一段, 存在基于 $x$ 的长度为 $n - |p|$ 的回路 $q$. $p$ 和 $q$ 的拼接就是所要的路径.
\end{proof}

\begin{corollary}[例子中的遍历性]
\label{cor:VI.3.4}
\begin{enumerate}
\item[(i)] 例子 \ref{exa:VI.1.5}(1) 的链遍历当且仅当 $r$ 奇.
\item[(ii)] 例子 \ref{exa:VI.1.5}(3) 的链在 $r \leq n - 1$ 时遍历.
\end{enumerate}
我们将在第 VI.5 节中用另一种方法重新证明 (ii).
\end{corollary}

\begin{proof}
(i) 对应链的图以 $\Z/r\Z$ 为顶点集, 且 $a, b \in \Z/r\Z$ 之间连边当且仅当 $a - b \in \{\pm 1\}$. 此图显然连通, 容易看到它二部当且仅当 $r$ 偶. 特别地, 若 $r$ 奇, 则命题蕴含链遍历.
现在假设 $r$ 偶. 对 $n$ 简单归纳表明, 对每个 $n \geq 1$ 和所有 $a, b \in \Z/r\Z$, 若 $n + a + b$ 在 $\Z/2\Z$ 中的像非零, 则 $P^n(a, b) = 0$. 事实上, 这由 $n = 1$ 时 $P$ 的定义得出. 假设结果对某个 $n \geq 1$ 已知, 证明对 $n + 1$. 设 $a, b \in \Z/r\Z$ 使得 $P^{n+1}(a, b) \neq 0$. 由于 $P^{n+1}(a, b) = \sum_{c \in \Z/r\Z}P(a, c)P^n(c, b)$, 存在 $c \in \Z/r\Z$ 使得 $P(a, c) \neq 0$ 且 $P^n(c, b) \neq 0$. 由归纳假设和 $n = 1$ 的情形, 这蕴含 $a + c \neq 0 \bmod 2$ 且 $n + c + b \neq 0 \bmod 2$, 然后 $n + a + b + 2c = n + a + b = 0 \bmod 2$.
(ii) 对应 Markov 链的图以 $\{1, \ldots, n\}$ 的基数为 $r$ 的子集的集合 $\Omega_r$ 为顶点集, $A, A' \in \Omega_r$ 之间连边当且仅当 $|A \cap A'| = r - 1$. 设 $A_0 = \{1, \ldots, r\}$. 首先证明图连通. 设 $A \in \Omega_r$. 把 $A$ 写成 $A = \{n_1, \ldots, n_r\}$, 并选择元素的顺序使得 $A \cap A_0 = \{1, \ldots, n_s\}$, 其中 $s = |A \cap A_0|$. 设 $m_1, \ldots, m_{r-s}$ 是 $A_0 - A$ 的元素. 对 $0 \leq i \leq r - s$, 设 $B_i = \{n_1, \ldots, n_{s+i}, m_{i+1}, \ldots, m_{r-s}\}$. 则 $B_0 = A_0$, $B_{r-s} = A$, 且对每个 $i \in \{0, \ldots, r - s - 1\}$, $B_i$ 和 $B_{i+1}$ 之间有边. 所以图连通.
现在通过找奇数长度回路证明图不二部. 设 $A = \{1, \ldots, r - 1, r + 1\}$, $B = \{2, \ldots, r, r + 1\}$. 则 $\{A_0, A\}$, $\{A, B\}$, $\{B, A_0\}$ 都是边, 所以找到了长度 $3$ 的回路.
\end{proof}

\section{齐性空间上的随机游走}
现在假设有限群 $G$ 传递 (左) 作用于有限集 $X$. 固定 $x_0 \in X$, 设 $K$ 是 $x_0$ 在 $G$ 中的稳定化子, 于是有双射 $G/K \simeq X$, $g \longmapsto g \cdot x_0$.
警告: 在本节中, 用 $G$ 上的计数测度定义卷积乘积和 $L^p$ 范数. 小心常数! (这样选择的原因是要让两个概率分布的卷积仍是概率分布.)

\begin{definition}[转移概率 $P_{\pi}$]
\label{def:VI.4.1}
若 $\pi$ 是 $G$ 上的概率分布, 记 $P_{\pi} \colon X \times X \to \R$ 为由
\[
P_{\pi}(xK, yK) = \pi(yKx^{-1})
\]
对所有 $x, y \in G$ 定义的函数.
\end{definition}

\begin{definition}[左不变随机游走]
\label{def:VI.4.2}
由 $\pi$ 驱动 (driven by $\pi$)、初始分布为 $\nu$ 的 $X$ 上左不变随机游走 (left-invariant random walk) 是状态空间为 $X$、初始分布为 $\nu$、转移矩阵为 $P_{\pi}$ 的 Markov 链.
\end{definition}

下面是用语言描述这个 Markov 链 $(X_n)_{n \geq 0}$: 按概率分布 $\nu$ 在 $X$ 上选择起点. 在时刻 $n$, 用概率分布 $\pi$ 选择 $G$ 的元素并作用于当前位置, 得到时刻 $n+1$ 的位置.

\begin{remark}[$P_{\pi}$ 双随机]
\label{rem:VI.4.3}
矩阵 $P_{\pi}$ 实际上双随机 (bistochastic), 即 $P_{\pi}$ 和它的转置都随机. 事实上, 对每个 $y \in G$,
\begin{align*}
\sum_{x \in G/K}P_{\pi}(xK, yK) &= \sum_{x \in G/K}\pi(yKx^{-1}) \\
&= \sum_{x \in G}\pi(yx^{-1}) \\
&= 1.
\end{align*}
特别地, $X$ 上的均匀概率分布是 $P_{\pi}$ 的不变分布. 若 $P_{\pi}$ 遍历, 它是唯一的不变分布.
\end{remark}

若齐性空间是 $G$ 本身, 可以给出简单的遍历性判据. (见 {[}1{]} 的引理 16.20 和命题 16.21.)

\begin{proposition}[遍历性判据]
\label{prop:VI.4.4}
假设 $X = G$, 设 $S = \operatorname{supp}(\pi)$. 记 $G_S$ 为可以写成 $g_1\cdots g_{2r}$ ($r \geq 0$) 的元素集合, 其中 $g_i$ 中恰好 $r$ 个在 $S$ 中, $r$ 个在 $S^{-1}$ 中.
则 $G_S$ 是 $G$ 的子群, 且函数 $P_{\pi}$ 遍历当且仅当 $G = G_S$.
特别地, 若 $\pi(1) \neq 0$, 则 $P_{\pi}$ 遍历当且仅当 $S$ 生成 $G$. 更一般地, 若 $S$ 生成 $G$ 且不包含于 $G$ 的严格子群的陪集中, 则 $P_{\pi}$ 遍历. (注意对每个 $g \in S$ 有 $S \subset gG_S$.)
\end{proposition}

\begin{proposition}[转移矩阵与卷积]
\label{prop:VI.4.5}
对每个 $n \geq 1$, $P_{\pi}^n = P_{\pi^{*n}}$, 其中 $\pi^{*n}$ 是 $\pi$ 的 $n$ 重卷积乘积.
\end{proposition}

\begin{proof}
对 $n$ 归纳证明. $n = 1$ 时它只是 $P_{\pi}$ 的定义. 假设等式对某个 $n \geq 1$ 已知, 证明对 $n+1$. 设 $x, y \in X$. 则
\begin{align*}
P_{\pi}^{n+1}(xK, yK) &= \sum_{z \in G/K}P_{\pi}(xK, zK)P_{\pi}^n(zK, yK) \\
&= \sum_{z \in G/K}\pi(zKx^{-1})\pi^{*n}(yKz^{-1}) \\
&= \sum_{z \in G, h \in K}\pi(zx^{-1})\pi^{*n}(yhz^{-1}) \\
&= \sum_{h \in K}\pi^{*(n+1)}(yhx^{-1}) \\
&= \pi^{*(n+1)}(x, y).
\end{align*}
\end{proof}

\begin{corollary}[上界引理 (upper bound lemma)]
\label{cor:VI.4.6}
设 $\pi$ 是 $G$ 上的概率测度, 假设 $\pi$ 被 $K$ 右不变. 考虑由 $\pi$ 驱动、初始分布为集中在 $x_0 \in X$ 上的 Dirac 测度的左不变随机游走 $(X_n)_{n \geq 0}$. 设 $\mu_n$ 是 $X_n$ 的分布, $\mu$ 是 $X$ 上的均匀概率分布.
则对每个 $n \geq 0$,
\[
\|\mu_n - \mu\|_{TV}^2 \leq \frac{1}{4}\sum_{\substack{(\rho, V) \in \widehat{G}\\ V^K \neq 0 \text{ 且 } \rho \not\simeq \mathbf{1}}}\dim(V)\operatorname{Tr}((\hat{\pi}(\rho)^*)^n \circ \hat{\pi}(\rho)^n),
\]
其中用 $\mathbf{1}$ 记 $G$ 的平凡表示.
\end{corollary}

回忆若 $(\rho, V) \in \widehat{G}$ 是 $G$ 的不可约酉表示, $f \colon G \to \CC$ 是函数, 则 $\hat{f}(\rho)$ 是由
\[
\hat{f}(\rho) = \sum_{x \in G}f(x)\rho(x^{-1})
\]
定义的 $V$ 的自同态.

\begin{proof}
固定 $n \geq 0$. 对每个 $x \in G$, 由引理 \ref{lem:VI.1.4} 和命题 \ref{prop:VI.4.5},
\[
\mu_n(x) = P_{\pi}^n(x_0, x) = \pi^{*n}(xK).
\]
设 $\pi_0$ 是 $G$ 上的均匀概率分布. 由引理 \ref{lem:VI.1.8},
\begin{align*}
\|\mu_n - \mu\|_{TV}^2 &= \frac{1}{4}\Bigl(\sum_{x \in G/K}|\mu_n(x) - \mu(x)|\Bigr)^2 \\
&= \frac{1}{4}\Bigl(\sum_{x \in G}|\pi^{*n}(x) - \pi_0(x)|\Bigr)^2 \\
&= \frac{1}{4}\|\pi^{*n} - \pi_0\|_1^2 \\
&\leq \frac{|G|}{4}\|\pi^{*n} - \pi_0\|_2^2,
\end{align*}
其中最后一个不等式是 Cauchy-Schwarz 不等式. (注意用 $G$ 上的计数测度定义 $L^p$ 范数.) 设 $f = \pi^{*n} - \pi_0 \in L^2(G)$. 由 Parseval 公式 (定理 \ref{thm:IV.6.3}(iii); 注意来自未规范化 Haar 测度的因子 $\frac{1}{|G|}$),
\[
\|f\|_2^2 = \frac{1}{|G|}\sum_{(\rho, V) \in \widehat{G}}\dim(V)\operatorname{Tr}(\hat{f}(\rho)^* \circ \hat{f}(\rho)).
\]
所以需要计算 $\hat{f}(\rho)$. 注意对每个 $\rho \in \widehat{G}$,
\[
\hat{f}(\rho) = \hat{\pi}(\rho)^n - \hat{\pi}_0(\rho).
\]
首先假设 $\rho = \mathbf{1}$. 则 $\hat{\pi}(\rho) = \hat{\pi}_0(\rho) = 1$, 所以 $\hat{f}(\rho) = 0$.
设 $(\rho, V) \in \widehat{G}$, 假设 $\rho \not\simeq \mathbf{1}$. 则 $\hat{\pi}_0(\rho) = \sum_{x \in G}\rho(x^{-1})$ 是 $\End(V)$ 的 $G$-等变元素, 所以由 Schur 引理是 $\mathrm{id}_V$ 的倍数, 且其迹等于 $\frac{1}{|G|}\sum_{x \in G}\chi(x) = 0$ (由推论 \ref{cor:IV.5.8}). 所以 $\hat{\pi}_0(\rho) = 0$, $\hat{f}(\rho) = \hat{\pi}(\rho)^n$. 为完成证明, 只需证明 $V^K = 0$ 时 $\hat{\pi}(\rho) = 0$. 设 $T = \hat{\pi}(\rho) = \sum_{x \in G}\pi(x)\rho(x^{-1})$, $P_K = \sum_{x \in K}\rho(x)$. 由于 $\pi$ 被 $K$ 右不变, 对每个 $x \in K$ 有 $\rho(x) \circ T = T$, 所以 $P_K \circ T = |K|T$. 但由命题 \ref{prop:V.1.7}, $P_K$ 是 $V^K$ 上的正交投影, 所以 $\Im(T) \subset V^K$, 因此 $V^K = 0$ 时 $T = 0$.
\end{proof}

\begin{corollary}[Gelfand 对情形的上界引理]
\label{cor:VI.4.7}
用前一推论的记号, 假设 $(G, K)$ 是 Gelfand 对且 $\pi$ 双 $K$-不变. 如第 V.6 节, 设 $\mathcal{Z}$ 是 $(G, K)$ 的对偶空间 (即由定理 \ref{thm:V.7.1} 的球函数集合).
则对每个 $n \geq 0$,
\[
\|\mu_n - \mu\|_{TV}^2 \leq \frac{1}{4}\sum_{\phi \in \mathcal{Z}, \phi \neq 1}\dim(V_{\phi})|\hat{\pi}(\phi)|^{2n},
\]
其中现在若 $f \in \mathcal{C}(K\backslash G/K)$, $\phi \in \mathcal{Z}$, 纯量 $\hat{f}(\phi) \in \CC$ 是由
\[
\hat{f}(\phi) = \sum_{x \in G}f(x)\phi(x^{-1})
\]
定义的球傅里叶变换.
\end{corollary}

\begin{proof}
证明与前一推论的证明几乎相同, 只是用推论 \ref{cor:V.7.2} 的 Parseval 公式计算 $\|\pi^{*n} - \pi_0\|_2^2$. 由该公式,
\[
\|\pi^{*n} - \pi_0\|_2^2 = \frac{1}{|G|}\sum_{\phi \in \mathcal{Z}}\dim(V_{\phi})|\hat{f}(\phi)|^2,
\]
其中 $f = \pi^{*n} - \pi_0$. 若 $\phi = 1$ 是对应于平凡表示的球函数, 则 $\hat{\pi}(\phi) = \hat{\pi}_0(\phi) = 1$, 所以 $\hat{f}(\phi) = 0$. 若 $\phi \neq 1$, 则
\[
\hat{\pi}_0(\phi) = \sum_{x \in G}\phi(x^{-1}) = \langle 1, \phi \rangle_{L^2(G)} = 0
\]
(例如由定理 \ref{thm:V.7.1}(i)). 所以 $\hat{f}(\phi) = \hat{\pi}(\phi)^n$, 完成证明.
\end{proof}

\section{Bernoulli-Laplace 扩散模型的应用}
回忆 Bernoulli-Laplace 扩散模型在例子 \ref{exa:VI.1.5}(3) 中描述过. 有两个正整数 $r$ 和 $b$. 此模型是集合 $\Omega_r$ (即 $\{1, \ldots, r + b\}$ 的基数为 $r$ 的子集) 上、初始分布为集中在 $\{1, \ldots, r\}$ 上的 Dirac 分布的 Markov 链 $(X_n)_{n \geq 0}$. 群 $G := S_{r+b}$ 在 $\Omega_r$ 上传递作用, $A_0 := \{1, \ldots, r\}$ 的稳定化子是 $K := S_r \times S_b$. 链的转移矩阵 $P$ 由
\[
P(A', A) = \begin{cases} \frac{(r-1)!(b-1)!}{(r+b)!} & \text{若 } r - |A \cap A'| = 1, \\ 0 & \text{否则} \end{cases}
\]
给出.
回忆在习题 \ref{ex:V.8.2.1}(e) 中定义了 $\Omega_r$ 上的度量 $d$, $d(A, A') = r - |A \cap A'|$, 且在 \ref{ex:V.8.2.1}(d) (和 \ref{ex:V.8.2.1}(f)) 中证明 $K$ 在 $G/K \simeq \Omega_r$ 上的轨道是以 $A_0$ 为中心的球面. $G$ 上的双不变概率分布 $\pi$ 与 $K\backslash G/K$ (即 $G/K$ 的 $K$-轨道) 上的概率分布双射对应, 且对 $P$ 的描述容易蕴含 $P = P_{\pi}$, 其中 $\pi$ 是对应于以 $A_0$ 为中心、半径 $1$ 的球面上的均匀分布的双不变概率分布.
若 $\mu_n$ 是 $X_n$ 的分布, $\mu$ 是 $\Omega_r$ 的均匀分布, 则由推论 \ref{cor:VI.4.7},
\[
\|\mu_n - \mu\|_{TV}^2 \leq \frac{1}{4}\sum_{\phi \in \mathcal{Z} - \{1\}}\dim(V_{\phi})|\hat{\pi}(\phi)|^{2n}.
\]
在第 V.8.2 节的习题中计算了所有这些项. 例如假设 $r \leq b$ (否则只需交换 $r$ 和 $b$, 得到等价的问题). 则在习题 \ref{ex:V.8.2.3} 中看到如何把 $G$ 在 $L^2(\Omega_r)$ 上的拟正则表示分解成不可约子表示 (见 \ref{ex:V.8.2.3}(j) 和 \ref{ex:V.8.2.3}(k)), 且恰好有 $r + 1$ 个. 记对应的球函数为 $\phi_0, \ldots, \phi_r$, 如在习题 \ref{ex:V.8.2.4} 中. 特别地, 函数 $\phi_0$ 就是常值函数 $1$. 在 \ref{ex:V.8.2.3}(f) 中计算了这些函数, 但实际上只需要 \ref{ex:V.8.2.3}(g). 事实上, 只关心 $\hat{\pi}(\phi_s)$, $1 \leq s \leq r$. 由于 $\pi$ 对应于以 $A_0$ 为中心、半径 $1$ 的球面上的均匀分布, 数 $\hat{\pi}(\phi_s)$ 就是 $\phi_s$ 中 $\sigma_{1,r-1}(A_0)$ 的系数 (用习题 \ref{ex:V.8.2.3} 的记号), 即
\[
\hat{\pi}(\phi_s) = 1 - \frac{s(r + b - s + 1)}{rb}.
\]
又, \ref{ex:V.8.2.3}(f) 说
\[
\dim(V_{\phi_s}) = \binom{r+b}{s} - \binom{r+b}{s-1}
\]
若 $1 \leq s \leq r$.
所以推论 \ref{cor:VI.4.7} 给出
\[
\|\mu_n - \mu\|_{TV}^2 \leq \frac{1}{4}\sum_{s=1}^r\Bigl(\binom{r+b}{s} - \binom{r+b}{s-1}\Bigr)\Bigl(1 - \frac{s(r + b - s + 1)}{rb}\Bigr)^{2n}.
\]
再花一些功夫, 可以得到下列结果.

\begin{theorem}[Bernoulli-Laplace 混合时间上界]
\label{thm:VI.5.1}
(见 {[}8{]} 第 3F 章定理 10.) 存在普适常数 $a \in \R_{>0}$ 使得若 $n = \frac{r+b}{4}(\log(2(r+b)) + c)$ 且 $c \geq 0$, 则
\[
\|\mu_n - \mu\|_{TV} \leq ae^{-c/2}.
\]
\end{theorem}

另一个不同的计算 (仍用球函数) 给出下列定理:

\begin{theorem}[Bernoulli-Laplace 混合时间下界]
\label{thm:VI.5.2}
(见 {[}7{]} 定理 6.3.2.) 若 $r = b$ 足够大, 则对 $n = \frac{r+b}{4}(\log(2(r+b)) - c)$ 且 $0 < c < \log(2(r+b))$,
\[
\|\mu_n - \mu\|_{TV} \geq 1 - 32e^{-c}.
\]
\end{theorem}

\section{局部紧群上的随机游走}
在本节中, 将看到关于更一般群上随机游走的一些结果 (大多不加证明). 这里没有触及的许多问题的好参考是 Breuillard 的综述 {[}5{]}.
固定局部紧群 $G$ 和 $G$ 上的左 Haar 测度.

\subsection{设定}

\begin{definition}[\index{Radon 测度}Radon 测度 (Radon measure)]
\label{def:VI.6.1.1}
(见注记 \ref{rem:I.4.1.6}.) $G$ 上的 (复) Radon 测度是 $\mathcal{C}_0(G)$ (带范数 $\|\cdot\|_{\infty}$) 上的有界线性泛函. 记 $\mathcal{M}(G)$ 为 Radon 测度的空间, $\|\cdot\|$ 为其范数 (即算子范数); 这是巴拿赫空间. 若 $\mu$ 是 Radon 测度, 记 $f \longmapsto \int_G f(x)\,\dd\mu(x)$ 为 $\mathcal{C}_0(G)$ 上对应的线性泛函.
\end{definition}

\begin{example}[Radon 测度的例子]
\label{exa:VI.6.1.2}
(1) 任何正则 Borel 测度都是 $G$ 上的 Radon 测度 (想把它们与一般 Radon 测度区分时, 这样的测度称为``正''的).
(2) 若 $\phi \in L^1(G)$, 则线性泛函 $f \longmapsto \int_G f(x)\phi(x)\,\dd x$ 是 $G$ 上的 Radon 测度, 常记为 $\phi(x)dx$ 或 $\phi dx$.
(3) 对每个 $x \in G$, 线性泛函 $\mathcal{C}_0(G) \to \CC$, $f \longmapsto f(x)$ 是 $G$ 上的 Radon 测度, 称为 $x$ 处的 Dirac 测度.
\end{example}

定义两个 Radon 测度 $\mu$ 和 $\nu$ 的卷积乘积 $\mu * \nu$ 为线性泛函
\[
f \longmapsto \int_{G \times G}f(xy)\,\dd\mu(x)\,\dd\nu(y).
\]
则不难验证 $\|\mu * \nu\| \leq \|\mu\|\|\nu\|$ 且 $\mathcal{M}(G)$ 上的卷积乘积结合. 这使得 $\mathcal{M}(G)$ 成为 Banach 代数, 且 $1$ 处的 Dirac 测度是 $\mathcal{M}(G)$ 的单位元.
若 $\mu = \phi dx$, $\mu' = \phi' dx$, 则容易验证 $\mu * \mu' = (\phi * \phi')dx$, 其中 $\phi * \phi'$ 是 $L^1(G)$ 中的通常卷积.
记 $\widehat{G}$ 为 $G$ 的不可约酉表示的酉等价类的集合. 可以把傅里叶变换 (普通和球两种版本) 延拓到 $\mathcal{M}(G)$:
\begin{enumerate}
\item[(1)] 若 $\mu \in \mathcal{M}(G)$, $(\pi, V) \in \widehat{G}$, 定义 $\hat{\mu}(\pi) \in \End(V)$:
\[
\hat{\mu}(\pi)(v) = \int_G\pi(x^{-1})(v)\,\dd\mu(x).
\]
\item[(2)] 假设 $G$ 是 Gelfand 对 $(G, K)$ 的第一项, $\phi$ 是 $G$ 上正定型球函数. 则对每个 $\mu \in \mathcal{M}(G)$, 定义 $\hat{\mu}(\phi) \in \CC$:
\[
\hat{\mu}(\phi) = \int_G\phi(x^{-1})\,\dd\mu(x).
\]
\end{enumerate}
对两种版本的傅里叶变换, 等式
\[
\widehat{\mu * \mu'} = \hat{\mu}\hat{\mu'}
\]
对所有 $\mu, \mu' \in \mathcal{M}(G)$ 成立 (其中右端的乘积在第一种情形是自同态的复合, 第二种情形是乘法).
下列定理是 L\'evy 收敛判据的推广. 说 Radon 测度的序列 $(\mu_n)_{n \geq 0}$ 弱收敛 (weakly converges), 若它在 $\mathcal{M}(G)$ 的弱*拓扑中收敛.

\begin{theorem}[\index{Lévy 收敛判据}L\'evy 判据]\footnote{见 {[}12{]} 第 5.2 节, 定理 5.2.}
\label{thm:VI.6.1.3}
\begin{enumerate}
\item[(i)] 若 $\mu, \mu' \in \mathcal{M}(G)$ 满足对每个 $\pi \in \widehat{G}$ 有 $\hat{\mu}(\pi) = \hat{\mu}'(\pi)$, 则 $\mu = \mu'$.
\item[(ii)] 设 $(\mu_n)_{n \geq 0}$ 是 $G$ 上的 (正) 概率测度序列, $\mu$ 是 $G$ 上的另一个概率测度. 若 $(\mu_n)_{n \geq 0}$ 弱收敛到 $\mu$, 则对每个 $(\pi, V) \in \widehat{G}$ 和每个 $v \in V$, $\lim_{n \to +\infty}\hat{\mu}_n(\pi)(v) = \hat{\mu}(\pi)(v)$. 反之, 若对每个 $(\pi, V) \in \widehat{G}$ 和所有 $v, w \in V$, $\lim_{n \to +\infty}\langle \hat{\mu}_n(\pi)(v), w \rangle = \langle \hat{\mu}(\pi)(v), w \rangle$, 则 $(\mu_n)_{n \geq 0}$ 弱收敛到 $\mu$.
\end{enumerate}
\end{theorem}

\subsection{随机游走}

固定 $G$ 上的正则 Borel 概率测度 $\mu$, 想理解 $\mu^{*n}$ 当 $n \to +\infty$ 时的行为.
与随机游走的联系是: $\mu^{*n}$ 是状态空间为 $G$、初始分布为 $\delta_1$、``转移矩阵''为 $\mu(yx^{-1})$ 的 Markov 链第 $n$ 步的分布. (为简化记号选择 $\delta_1$ 作为初始分布, 但对大多数结果这不是真正必要的.) 换句话说, 考虑取值于 $G$、分布为 $\mu$ 的独立随机变量序列 $(g_n)_{n \geq 1}$. Markov 链 $(X_n)_{n \geq 0}$ 由 $X_n = g_1\cdots g_n$ 定义 (所以 $X_0$ 是常值函数 $1$). 也可以考虑空间 $G/K$ 上的随机游走, 其中 $K$ 是 $G$ 的子群: 如前取 $(g_n)_{n \geq 1}$, 固定取值于 $G/K$ 的某个初始随机变量 $X_0$ (例如常值函数), 设 $X_n = g_n\cdots g_1X_0$.

\subsection{紧群}
在本节中, 假设 $G$ 紧. 从 Ito 和 Kawada 的一般收敛结果开始 ({[}14{]}, 也见 {[}5{]} 定理 2.3).
回忆 $\mu$ 的支撑 (support) 由定义是满足如下条件的 $x \in G$ 的集合: 对 $x$ 的每个邻域 $U$, $\mu(U) > 0$.

\begin{theorem}[Ito-Kawada 定理]\footnote{见 {[}14{]}, 也见 {[}5{]} 定理 2.3.}
\label{thm:VI.6.3.1}
假设 $\mu$ 的支撑生成 $G$ 的稠密子群, 且不包含于 $G$ 的真闭子群的任何 (左或右) 陪集中. 则序列 $(\mu^{*n})_{n \geq 0}$ 弱收敛到 $G$ 上的规范化 Haar 测度.
\end{theorem}

证明基于定理 \ref{thm:VI.6.1.3}(ii) 的收敛判据. 必须证明对每个非平凡的 $(\pi, V) \in \widehat{G}$, 序列 $\widehat{\mu^{*n}}(\pi) = \hat{\mu}(\pi)^n$ 在 $\End(V)$ 中收敛到 $0$. 注意 $V$ 有限维 (因为 $G$ 紧), 所以 $\End(V)$ 中的所有收敛概念等价, 只需证明 $\hat{\mu}(\pi)$ 的所有特征值的绝对值 $< 1$. 假设事实并非如此, 则存在单位向量 $v \in V$ 使得
\[
\int_G\pi(x^{-1})(v)\,\dd\mu(x) = \lambda v,
\]
其中 $|\lambda| = 1$. 不难看到这迫使 $\pi(x^{-1})(v)$ $\mu$-几乎处处等于 $\lambda v$, 与定理的假设矛盾.
注意此结果比命题 \ref{prop:VI.4.4} (和 Perron-Frobenius 定理) 弱得多, 因为它只保证 $(\mu^{*n})_{n \geq 0}$ 的弱收敛, 对其它拓扑 (如全变差距离诱导的拓扑) 的收敛或收敛速度只字未提. 若 $G$ 有限, 概率测度集合上的所有收敛概念一致 (它只是 $G$ 上函数空间 (有限维) 的凸子集); 又, 由上界引理 (推论 \ref{cor:VI.4.6}) 收敛速度指数且被最大的 $\neq 1$ 的 $\hat{\mu}(\pi)$ 特征值控制. 但若 $G$ 无穷, 则 $\widehat{G}$ 也无穷, 所以即使 $\hat{\mu}(\pi)$ 的所有特征值的绝对值 $< 1$, 也可以得到任意接近 $1$ 的特征值. 事实上, 当这种情况不发生时有专门的名字:

\begin{definition}[\index{谱隙}谱隙 (spectral gap)]
\label{def:VI.6.3.2}
若存在 $\varepsilon > 0$ 使得对每个非平凡的 $\pi \in \widehat{G}$ 和 $\hat{\mu}(\pi)$ 的每个特征值 $\lambda$, $|\lambda| < 1 - \varepsilon$, 则说 $G$ 上的概率测度 $\mu$ 有谱隙 (spectral gap).
\end{definition}

先看一些例子.

\begin{example}[$L^2$ 密度有谱隙]
\label{exa:VI.6.3.3}
若 $\mu = \phi dx$, $\phi \in L^2(G)$, 则 $\mu$ 有谱隙. 事实上, 上界引理 (推论 \ref{cor:VI.4.6}) 用本质上相同的证明仍成立: 对每个 $n \geq 0$,
\[
\|\mu^{*n} - \mu_G\|_{TV}^2 \leq \frac{1}{4}\sum_{\substack{(\rho, V) \in \widehat{G}\\ \rho \not\simeq \mathbf{1}}}\dim(V)\operatorname{Tr}((\hat{\pi}(\rho)^*)^n \circ \hat{\pi}(\rho)^n),
\]
其中用 $\mathbf{1}$ 记 $G$ 的平凡表示, $\mu_G$ 记 $G$ 上的规范化 Haar 测度. (也可以证明空间 $G/K$ 上的随机游走版本.) 所以此情形下有全变差距离的收敛且速度指数.
\end{example}

在另一个极端, 有有限支撑的测度.

\begin{example}[无谱隙的测度]
\label{exa:VI.6.3.4}
取 $G = S^1$. 设 $\lambda_1, \ldots, \lambda_r \in \R$, 考虑 $G$ 上的测度
\[
\mu = \frac{1}{2r}\sum_{s=1}^r(\delta_{e^{2i\pi\lambda_s}} + \delta_{e^{-2i\pi\lambda_s}}).
\]
回忆 $\widehat{G} = \Z$ (其中 $n \in \Z$ 对应于 $G$ 的表示 $z \longmapsto z^n$). 对每个 $n \in \Z$,
\[
\hat{\mu}(n) = \frac{1}{2r}\sum_{s=1}^r(e^{2i\pi n\lambda_s} + e^{-2i\pi n\lambda_s}).
\]
假设族 $(1, \lambda_1, \ldots, \lambda_r)$ 在 $\Q$ 上线性无关. 则 Kronecker 定理 (见例如 {[}13{]} 第 XXIII 章) 说集合 $\{(e^{2i\pi n\lambda_1}, \ldots, e^{2i\pi n\lambda_r}) \mid n \in \Z\}$ 在 $(S^1)^r$ 中稠密. 所以可以找到 $n \neq 0$ 使得 $\hat{\mu}(n)$ 任意接近 $1$. 换句话说, 测度 $\mu$ 没有谱隙.
\end{example}

$SU(d)$ 等``好''群上的哪些测度有谱隙是非常困难且活跃的研究领域. 给出一些 (困难的) 近期结果, 对 $G = \SU(d)$ 归功于 Bourgain 和 Gamburd (cf. {[}4{]} 和 {[}3{]}), 对一般简单紧 Lie 群归功于 Benoist 和 de Saxc\'e (cf. {[}2{]}).

\begin{theorem}[谱隙判据]\footnote{见 {[}4{]} 和 {[}3{]} (对 $\SU(d)$), {[}2{]} (一般简单紧 Lie 群).}
\label{thm:VI.6.3.5}
设 $G$ 是简单紧 Lie 群 (例如 $d \geq 2$ 时 $G = \SU(d)$ 或 $d = 3$ 或 $d \geq 5$ 时 $G = \SO(d)$), $\mu$ 是 $G$ 上的概率测度. 说 $\mu$ 几乎 Diophantine (almost Diophantine), 若存在 $c_1, c_2 > 0$ 使得对 $G$ 的每个真闭子群 $H$ 和每个足够大的 $n \in \Z_{\geq 0}$, $\mu^{*n}(\{x \in G \mid d(x, H) \leq e^{-c_1n}\}) \leq e^{-c_2n}$ (其中 $d$ 是 $G$ 上的任何度量).
则 $\mu$ 有谱隙当且仅当它几乎 Diophantine.
\end{theorem}

虽然下一版本对任何简单紧 Lie 群都有推广, 为简单起见只对 $\SU(d)$ 陈述.

\begin{theorem}[代数支撑的谱隙]
\label{thm:VI.6.3.6}
设 $G = \SU(d)$, $\mu$ 是 $G$ 上的概率测度, 使得 $\mu$ 的支撑生成 $G$ 的稠密子群 (这样的测度有时称为``适应的'' (adapted)).
若 $\mu$ 的支撑的每个元素有代数元素, 则 $\mu$ 有谱隙.
\end{theorem}

事实上, Benoist 和 de Saxc\'e 猜想代数性条件不必要 (所以每个适应测度都应有谱隙), 见 {[}2{]} 的引言.

\begin{remark}[与 Banach-Ruziewicz 问题的联系]
\label{rem:VI.6.3.7}
谱隙问题也与 Banach-Ruziewicz 问题相关 (联系见 Sarnak 的书 {[}21{]} 第 2 章; Banach-Ruziewicz 问题的另一个好参考是 Lubotzky 的书 {[}15{]}). 此问题问 Lebesgue 测度是否是球面 $S^n \subset \R^n$ 的 Lebesgue 可测子集上唯一的 (差一个常数) 有限可加 $\SO(n+1)$-不变测度. 已知 $n = 1$ 时答案``否'', $n \geq 2$ 时``是''. 对 $n \geq 4$, 它归功于 Margulis 和 Sullivan, 用 $n \geq 4$ 时 $\SO(n+1)$ 有满足 Kazhdan 性质 (T) 的有限生成稠密子群这一事实 (事实上, 同样的方法将证明 Haar 测度是任何非 $\SO(n)$ ($n \geq 4$) 的简单紧 Lie 群上唯一的左不变平均). 对 $n = 2, 3$, 解最初归功于 Drinfeld, 用 Jacquet-Langlands 对应和 Ramanujan-Petersson 猜想. (所有这些及更多内容见 {[}15{]}.)
\end{remark}

\subsection{用 Fourier 分析的随机游走收敛}
现在给出紧群 (或齐性空间) 上可以用 Fourier 分析分析的随机游走的一些例子, 精神与第 VI.5 节相同.
如前所述 (在例子 \ref{exa:VI.6.3.3} 中), 上界引理 (推论 \ref{cor:VI.4.6}) 对一般紧群仍成立.
与有限群一样, 当测度 $\mu$ 共轭不变, 或 $\mu$ 双 $K$-不变且 $(G, K)$ 是 Gelfand 对时, Fourier 分析效果最好.

\paragraph{$\SO(n)$ 中的随机反射 (Random reflections in $\SO(n)$)}
此结果的参考文献是 Rosenthal 的论文 {[}17{]}. 固定 $n \geq 2$, $\theta \in (0, 2\pi)$. 设
\[
R_{\theta} = \begin{pmatrix} \cos\theta & \sin\theta & & 0 \\ -\sin\theta & \cos\theta & & \\ & & 1 & \\ & & & \ddots \\ 0 & & & 1 \end{pmatrix} \in \SO(n),
\]
设 $\mu_{\theta}$ 是集中在 $R_{\theta}$ 的共轭类上的唯一共轭不变概率测度 (换句话说, 测度 $\mu_{\theta}$ 是 $\SO(n)$ 的规范化 Haar 测度在映射 $\SO(n) \to \SO(n)$, $x \longmapsto xR_{\theta}x^{-1}$ 下的像).

\begin{theorem}[随机反射的混合时间]
\label{thm:VI.6.4.1}
\begin{enumerate}
\item[(i)] 存在 $\Gamma, \Delta > 0$ (其中 $\Delta$ 不依赖 $\theta$) 使得对每个 $n \geq 1$ 和每个 $c > 0$, 若 $k = \frac{1}{2(1-\cos\theta)}(n\log n - cn)$, 则
\[
\|\mu_{\theta}^{*k} - dx\|_{TV} \geq 1 - \Gamma e^{-2c} - \Delta\frac{\log n}{n}.
\]
\item[(ii)] 假设 $\theta = \pi$. 则存在 $\Lambda, \gamma > 0$ 使得对每个 $n \geq 3$ 和每个 $c > 0$, 若 $k = \frac{1}{4}n\log n + cn$, 则
\[
\|\mu_{\theta}^{*k} - dx\|_{TV} \leq \Lambda e^{-\gamma c}.
\]
\end{enumerate}
\end{theorem}

\paragraph{Gelfand 对情形 (The Gelfand pair case)}
此结果的参考文献是 Su 的论文 {[}24{]}.
固定 $\theta \in (0, \pi)$, 考虑 $S^2 \simeq \SO(3)/\SO(2)$ 上的下列随机过程:
\begin{itemize}
\item $X_0$ 恒为北极;
\item 从 $X_n$ 到 $X_{n+1}$, 选择 (独立且均匀地) 一个方向, 并沿该方向的测地线 (= 大圆) 移动距离 $\theta$.
\end{itemize}
此随机游走不由 $\SO(3)$ 上的测度驱动, 但它等价于这样的游走 (见 {[}24{]} 第 3 节). 设 $\mu_n$ 是 $X_n$ 的分布, $\lambda$ 是 $S^2$ 上唯一的 $\SO(3)$-不变概率测度. 则有下列结果:

\begin{theorem}[球面上的随机游走]
\label{thm:VI.6.4.2}
若 $n = \frac{1}{\sin^2\theta}c$ 且 $c \geq 0$, 则
\[
0.433e^{-c/2} \leq \|\mu_l - \lambda\|_{DD} \leq 4.442e^{-c/8}.
\]
\end{theorem}

在此定理中, $\|\cdot\|_{DD}$ 是\index{偏差距离}偏差距离 (discrepancy distance): 若 $X$ 是度量空间, $\mu, \mu'$ 是 $X$ 上的两个 (Borel) 概率测度, 则
\[
\|\mu - \mu'\|_{DD} = \sup_{B \subset X \text{ 球}}|\mu(B) - \mu'(B)|.
\]
它被全变差距离控制上界, 但可以看到全变差距离看不到的一些现象 (见下一小节).

\paragraph{关于不同类型收敛的注记 (Remark about the different types of convergence)}
本小节的参考文献是 Su 的论文 {[}23{]}.
考虑圆 $S^1$ 上由测度 $\mu = \frac{1}{2}(\delta_{e^{2i\pi\alpha}} + \delta_{e^{-2i\pi\alpha}})$ (对某个无理数 $\alpha \in \R$) 驱动的随机游走 $(X_n)_{n \geq 0}$, 设 $\mu_n$ 是 $X_n$ 的分布. 则:
\begin{itemize}
\item Ito-Kawada 的一般收敛结果 (定理 \ref{thm:VI.6.3.1}) 说 $(\mu_n)_{n \geq 0}$ 弱收敛到 $S^1$ 上的规范化 Haar 测度 $dx$.
\item 另一方面, 在例子 \ref{exa:VI.6.3.4} 中看到 $\mu$ 没有谱隙, 所以收敛不能太好. 事实上, $(\mu_n)_{n \geq 0}$ 不全变差距离收敛到 $dx$.
\item 第三方面, 在许多情形 $(\mu_n)_{n \geq 0}$ 确实在偏差距离中收敛到 $dx$ (但不指数快). 更精确地有:
\end{itemize}

\begin{theorem}[无理旋转的偏差]
\label{thm:VI.6.4.3}
设 $\eta$ 是 $\alpha$ 的类型 (type), 即
\[
\eta = \sup\{\gamma > 0 \mid \liminf_{m \to +\infty}m^{\gamma}\{m\alpha\} = 0\}
\]
(其中 $\{\cdot\}$ 是分数部分). 则对每个 $\varepsilon > 0$,
\[
\|\mu_n - dx\|_{DD} = O(n^{-1/2\eta + \varepsilon})
\]
且
\[
\|\mu_n - dx\|_{DD} = \Omega(n^{-1/2\eta - \varepsilon}).
\]
若 $\alpha$ 是二次无理数, 可以做得更好: 存在常数 $c_1, c_2 > 0$ 使得对每个 $n \geq 1$,
\[
\frac{c_1}{\sqrt{n}} \leq \|\mu_n - dx\|_{DD} \leq \frac{c_2}{\sqrt{n}}.
\]
已知 $\alpha$ 代数时 $\eta = 1$, 且 $[0,1]$ 中类型 $1$ 的元素子集的 Lebesgue 测度为 $1$.
\end{theorem}

\subsection{非紧群上的随机游走}
不再假设 $G$ 紧. 固定 $G$ 上的概率测度 $\mu$. 可以问的众多问题之一是: $G$ 上由 $\mu$ 驱动的随机游走是否走向无穷, 若是, 多快.
本节的一个参考是 Quint 的出色课程笔记 ({[}16{]}).
首先定义连续线性算子 $P_{\mu} \colon L^2(G) \to L^2(G)$:
\[
P_{\mu}(f)(x) = \int_G f(yx)\,\dd\mu(y)
\]
若 $f \in \mathcal{C}_c(G)$, $x \in G$; 通过连续性延拓到 $L^2(G)$. (若 $\mu = \phi dx$, $\phi \in L^1(G)$, 这正好是定理 \ref{thm:I.4.2.6}(i) 的构造应用于 $G$ 的右正则表示.)
记 $\rho(P_{\mu})$ 为 $P_{\mu}$ 作为 Banach 代数 $\End(L^2(G))$ 的元素的谱半径. 总有 $\rho(P_{\mu}) \leq 1$ (因为 $\mu$ 是概率测度).

\begin{theorem}[Kesten 判据]
\label{thm:VI.6.5.1}
(定理 5.2 of {[}16{]}.)
\begin{enumerate}
\item[(i)] 若 $G$ 可解, 则 $\rho(P_{\mu}) = 1$.
\item[(ii)] 设 $H$ 是由 $\mu$ 的支撑生成的子群在 $G$ 中的闭包. 若 $\rho(P_{\mu}) = 1$, 则 $H$ 可解.
\end{enumerate}
\end{theorem}

\begin{definition}[紧生成 (compactly generated)]
\label{def:VI.6.5.2}
若存在生成 $G$ 的紧子集 $K$, 则说 $G$ 紧生成 (compactly generated).
\end{definition}

若 $G$ 离散, 这只是说 $G$ 有限生成.

\begin{definition}[$j_K$]
\label{def:VI.6.5.3}
假设 $G$ 紧生成, 设 $K$ 是生成 $G$ 的对称紧子集. 定义 $j_K \colon G \to \Z$:
\[
j_K(x) = \min\{n \in \Z_{\geq 0} \mid x \in K^n\}
\]
(约定 $K^0 = \{1\}$).
\end{definition}

\begin{lemma}[$j_K$ 的等价性]
\label{lem:VI.6.5.4}
若 $K$ 和 $L$ 是生成 $G$ 的两个对称紧子集, 则存在 $a > 0$ 使得 $j_L \leq aj_K$.
\end{lemma}

\begin{corollary}[非可解群的逃逸速度]
\label{cor:VI.6.5.5}
(推论 7.3 of {[}16{]}.)
假设 $G$ 紧生成, 设 $K$ 是生成 $G$ 的对称紧子集. 设 $\mu$ 是 $G$ 上的概率测度, $(g_n)_{n \geq 1}$ 是取值于 $G$、分布为 $\mu$ 的独立随机变量序列.
设 $H$ 是由 $\mu$ 的支撑生成的子群在 $G$ 中的闭包. 若 $H$ 不可解, 则存在 $\alpha, \varepsilon > 0$ 使得对每个 $n \geq 1$,
\[
\mathbb{P}(j_K(g_n\cdots g_1) \leq \varepsilon n) = o(e^{-\alpha n}).
\]
特别地, 由 Borel-Cantelli 引理 (见 {[}18{]} 第 17.1 节), 若 $n$ 足够大, 几乎必然 $j_K(g_n\cdots g_1) \geq \varepsilon n$.
\end{corollary}

以一个例子结束. 说 $\SL_2(\R)$ 的子群 $H$ 非初等 (non-elementary), 若 $H$ 的共轭不包含于 $\SO(2)$、$\left\{\begin{pmatrix} a & b \\ 0 & a^{-1} \end{pmatrix}, a \in \R^{\times}, b \in \R\right\}$ 或 $\left\{\begin{pmatrix} a & 0 \\ 0 & a^{-1} \end{pmatrix}, a \in \R^{\times}\right\} \cup \left\{\begin{pmatrix} 0 & a \\ a^{-1} & 0 \end{pmatrix}, a \in \R^{\times}\right\}$. (等价条件是 $H$ 不紧且不固定 $\R^2$ 中的直线或两条直线的并. 这里 $\SL_2(\R)$ 在 $\R^2$ 上的作用由包含 $\SL_2(\R) \subset \GL_2(\R)$ 给出的标准作用.)

\begin{proposition}[非初等与不可解]
\label{prop:VI.6.5.6}
(命题 8.6 of {[}16{]}.)
$\SL_2(\R)$ 的闭子群不可解当且仅当它非初等.
\end{proposition}

\begin{example}[非初等子群的随机游走]
\label{exa:VI.6.5.7}
若 $t \in \R^{\times}$, 设 $a_t = \begin{pmatrix} t & 0 \\ 0 & t^{-1} \end{pmatrix}$. 若 $\theta \in \R$, 设 $r_{\theta} = \begin{pmatrix} \cos\theta & \sin\theta \\ -\sin\theta & \cos\theta \end{pmatrix}$.
则若 $s, t > 1$, $0 < \theta < \pi/2$, 由 $a_s$ 和 $r_{\theta}a_tr_{\theta}^{-1}$ 张成的 $\SL_2(\R)$ 的子群非初等, 所以推论 \ref{cor:VI.6.5.5} 适用于由测度 $\mu = \frac{1}{2}(\delta_{a_t} + \delta_{r_{\theta}a_tr_{\theta}^{-1}})$ 驱动的随机游走.
\end{example}

\section{问题: 非可解群上的随机游走}
设 $G$ 是局部紧群. 固定 $G$ 上的左不变 Haar 测度 $\mu_G$. 在此问题中, $G$ 上的所有测度都假设是正则 Borel 测度, 可以假设 $G$ $\sigma$-紧.\footnote{我们主要关心\index{紧生成群}紧生成群, 它们显然 $\sigma$-紧.}
\begin{enumerate}
\item[(1)] 设 $H$ 是 $G$ 的闭子群, $\mu_H$ 是 $H$ 上的左不变 Haar 测度. 此题的目标是证明: 若 $L^{\infty}(G)$ 上存在被 $H$ 左不变的平均, 则 $H$ 可解.\footnote{其逆也真且容易证明得多, 但不需要它.}
  \begin{enumerate}
  \item[(a)] 要证明: 对 $G$ 中 $1$ 的每个开邻域 $U$, 存在 $G/H$ 的子集 $T$ 使得:
    \begin{itemize}
    \item $G/H = \bigcup_{t \in T}Ut$;
    \item 对 $G/H$ 的每个紧子集 $K$, 集合 $K \cap T$ 有限.
    \end{itemize}
    为简化证明 (也因为这是后面唯一需要的情形), 假设 $G$ 是紧子集的可数族的并. 选择 $G$ 中 $1$ 的开邻域 $V$ 使得 $V^{-1}V \subset U$.
    \begin{enumerate}
    \item[(i)] 证明若 $K$ 是 $G/H$ 的紧子集, 则不存在 $K$ 的元素序列 $(t_n)_{n \geq 0}$ 使得 $n \neq m$ 时 $Vt_n \cap Vt_m = \emptyset$.
    \item[(ii)] 写 $G/H = \bigcup_{n \geq 0}K_n$, 其中 $K_n$ 是 $G/H$ 的紧子集, 且对每个 $n \geq 0$ 有 $K_n \subset K_{n+1}$. 归纳定义 $G/H$ 的有限子集族 $(T_n)_{n \geq 0}$:
      \begin{itemize}
      \item $T_0 = \emptyset$;
      \item 对 $n \geq 1$, 取 $T_n$ 为 $L_n := K_n - (UT_0 \cup \cdots \cup UT_{n-1})$ 的极大有限子集, 使得对所有满足 $t \neq t'$ 的 $t, t' \in T_n$, $Vt \cap Vt' = \emptyset$.
      \end{itemize}
      证明定义有意义且对每个 $n \geq 1$, $K_n \subset \bigcup_{i=0}^n UT_i$.
    \item[(iii)] 证明 $T := \bigcup_{n \geq 0}T_n$ 满足两个要求的性质.
    \end{enumerate}
  \item[(b)] 证明存在有界连续函数 $\theta \colon G \to \R_{\geq 0}$ 使得:
    \begin{itemize}
    \item 对 $G$ 的任何紧子集 $K$, 函数 $\theta|_{KH}$ 有紧支撑;
    \item 对每个 $x \in G$, $\int_H\theta(xy)\,\dd\mu_H(y) = 1$.
    \end{itemize}
    (提示: 取 (a) 中的 $T$, 其中 $U$ 相对紧, 取满足 $\phi|_U = 1$ 的 $\phi \in \mathcal{C}_c^+(G)$, 定义 $\psi(x) = \sum_{t \in T}\phi(xg_t^{-1})$, 其中 $g_t \in G$ 是 $t$ 的提升, 然后稍微修改 $\psi$.)
  \item[(c)] 对每个 $\phi \in L^{\infty}(H)$, 定义 $\phi_{\theta} \colon G \to \CC$:
    \[
    \phi_{\theta}(x) = \int_H\phi(y)\theta(x^{-1}y)\,\dd\mu_H(y).
    \]
    证明 $\phi_{\theta} \in L^{\infty}(G)$.
  \item[(d)] 若 $M$ 是 $L^{\infty}(G)$ 上被 $H$ 左不变的平均, 证明 $\phi \longmapsto M(\phi_{\theta})$ 是 $L^{\infty}(H)$ 上的左不变平均.
  \end{enumerate}
\item[(2)] 设 $T \colon V \to W$ 是两个赋范 $\CC$-向量空间之间的有界线性算子.
  \begin{enumerate}
  \item[(a)] 若 $V$ 和 $W$ 是希尔伯特空间且 $\Im(T)$ 不在 $W$ 中稠密, 证明 $T^*$ 不单射.
  \item[(b)] 若 $V$ 完备且 $\Im(T)$ 不在 $W$ 中闭, 证明存在 $V$ 的范数 $1$ 向量序列 $(v_n)_{n \geq 0}$ 使得 $\lim_{n \to +\infty}T(v_n) = 0$.
  \end{enumerate}
\item[(3)] 设 $\mu$ 是 $G$ 上的概率测度. 对每个 $f \in \mathcal{C}_c(G)$, 定义函数 $P_{\mu}f \colon G \to \CC$:
    \[
    (P_{\mu}f)(x) = \int_G f(yx)\,\dd\mu(y).
    \]
  \begin{enumerate}
  \item[(a)] 对每个 $f \in \mathcal{C}_c(G)$, 证明 $P_{\mu}f$ 连续且 $\|P_{\mu}f\|_2 \leq \|f\|_2$.
  \item[(b)] 证明 $P_{\mu}$ 延拓成连续线性算子 $P_{\mu} \colon L^2(G) \to L^2(G)$ 且 $\|P_{\mu}\|_{\mathrm{op}} \leq 1$.
  \item[(c)] 证明 $(P_{\mu})^* = P_{\nu}$, 其中 $\nu$ 是由 $\nu(E) = \mu(E^{-1})$ 定义的概率测度.
  \end{enumerate}
\item[(4)] 假设 $G$ 可解, 设 $\mu$ 是 $G$ 上的概率测度. 要证明 $\rho(P_{\mu}) = 1$. (其中 $\rho(P_{\mu})$ 是 $P_{\mu}$ 作为 $\End(L^2(G))$ 的元素的谱半径.)
  \begin{enumerate}
  \item[(a)] 对每个 $\varepsilon > 0$, 证明存在 $G$ 的紧子集 $K$ 和函数 $f \in L^2(G)$ 使得:
    \begin{itemize}
    \item $\mu(K) \geq 1 - \varepsilon$;
    \item $\|f\|_2 = 1$;
    \item $\sup_{x \in K}\|L_xf - f\|_2 \leq \varepsilon$.
    \end{itemize}
  \item[(b)] 对每个 $\varepsilon > 0$, 证明存在函数 $f \in L^2(G)$ 使得 $\|f\|_2 = 1$ 且 $\|P_{\mu}f - f\|_2 \leq \varepsilon$.
  \item[(c)] 证明 $\rho(P_{\mu}) = 1$.
  \end{enumerate}
\item[(5)] 设 $\mu$ 是 $G$ 上的概率测度. 假设 $\rho(P_{\mu}) = 1$, 设 $H$ 是由 $\mu$ 的支撑生成的子群的闭包. (其中 $\mu$ 的支撑是满足 $x$ 的每个开邻域对 $\mu$ 有正体积的 $x \in G$ 的集合.) 要证明 $H$ 可解.
  \begin{enumerate}
  \item[(a)] 证明下列条件至少有一个成立:
    \begin{enumerate}
    \item[($\alpha$)] 存在 $\lambda \in \CC$ 使得 $|\lambda| = 1$ 且 $P_{\mu} - \lambda\mathrm{id}_{L^2(G)}$ 不单射.
    \item[($\beta$)] 存在 $\lambda \in \CC$ 使得 $|\lambda| = 1$ 且 $\Im(P_{\mu} - \lambda\mathrm{id}_{L^2(G)})$ 不在 $L^2(G)$ 中稠密.
    \item[($\gamma$)] 存在 $\lambda \in \CC$ 使得 $|\lambda| = 1$ 且 $\Im(P_{\mu} - \lambda\mathrm{id}_{L^2(G)})$ 不在 $L^2(G)$ 中闭.
    \end{enumerate}
  \item[(b)] 假设条件 ($\alpha$) 成立.
    \begin{enumerate}
    \item[(i)] 证明 $P_{\mu} - \mathrm{id}_{L^2(G)}$ 不单射.
    \item[(ii)] 证明存在 $L^2(G)$ 的非零元素被 $H$ 左不变.
    \item[(iii)] 证明 $H$ 紧.
    \end{enumerate}
  \item[(c)] 若条件 ($\beta$) 成立, 证明 $H$ 紧.
  \item[(d)] 假设条件 ($\gamma$) 成立.
    \begin{enumerate}
    \item[(i)] 证明存在 $L^2(G)$ 的范数 $1$ 元素序列 $(f_n)_{n \geq 0}$ 使得 $\lim_{n \to +\infty}\|P_{\mu}f_n - \lambda f_n\|_2 = 0$.
    \item[(ii)] 定义 $g_n \colon G \to \R_{\geq 0}$:
      \[
      g_n(x) = \|L_xf_n - \lambda f_n\|_2^2.
      \]
      证明序列 $(g_n)_{n \geq 0}$ 在 $L^1(G, \mu)$ 中收敛到 $0$ (注意测度!).
    \item[(iii)] 证明可以假设 $\lim_{n \to +\infty}g_n(x) = 0$ 对 $\mu$-几乎每个 $x \in G$.
    \item[(iv)] 定义 $h_n \colon G \to \R_{\geq 0}$: $h_n(x) = |f_n(x)|^2$. 证明对每个 $n$ 有 $\int_G h_n(x)\,\dd\mu_G(x) = 1$, 且 $\lim_{n \to +\infty}\|L_xh_n - h_n\|_1 = 0$ 对 $\mu$-几乎每个 $x \in G$. (注意回到 $L^1(G)$.)
    \item[(v)] 证明存在 $L^{\infty}(G)$ 上被 $H$ 左不变的平均 $M$.
    \end{enumerate}
\end{enumerate}
\item[(6)] 若 $K$ 是生成 $G$ 的紧对称 (即 $K^{-1} = K$) 子集, 定义 $j_K \colon G \to \Z_{\geq 0}$:
  \[
  j_K(x) = \min\{n \in \Z_{\geq 0} \mid x \in K^n\}
  \]
  (约定 $K^0 = \{1\}$).
  证明若 $L$ 是生成 $G$ 的另一个紧对称子集, 则存在 $a > 0$ 使得 $j_L \leq aj_K$.
  设 $K$ 是 $G$ 的非空内部紧对称子集. 此题的目标是证明序列 $(\mu_G(K^n)^{1/n})_{n \geq 1}$ 收敛.
  \begin{enumerate}
  \item[(a)] 若 $A$, $B$ 和 $C$ 是 $G$ 的紧子集, 证明
    \[
    \mu_G(AB)\mu_G(C) \leq \mu_G(AC)\mu_G(C^{-1}B).
    \]
    (提示: 看 $\mathbb{1}_{AC} * \mathbb{1}_{C^{-1}B}$.)
  \item[(b)] 证明对所有 $n, m \in \Z_{\geq 1}$,
    \[
    \mu_G(K^{n+m}) \leq \mu_G(K)^{-1}\mu_G(K^{n+1})\mu_G(K^{m+1}).
    \]
  \item[(c)] 得出结论.
  \end{enumerate}
\item[(7)] 设 $\mu$ 是 $G$ 上的概率测度. 假设由 $\mu$ 的支撑生成的子群的闭包 $H$ 不可解, 且 $G$ 紧生成. 固定生成 $G$ 的对称紧子集 $K$. 要证明存在 $\varepsilon, \alpha > 0$ 使得, 若 $(g_n)_{n \geq 1}$ 是取值于 $G$、分布为 $\mu$ 的独立随机变量序列, 则对每个 $n \geq 0$,
  \[
  \mathbb{P}(j_K(g_n\cdots g_1) \leq \varepsilon n) = o(e^{-\alpha n}).
  \]
  换句话说, 要证明
  \[
  \mu^{*n}(K^{\lfloor \varepsilon n \rfloor}) = o(e^{-\alpha n}),
  \]
  其中 $\mu^{*n}$ 是 $G^n$ 上的测度 $\mu^n$ 在乘法映射 $(x_1, \ldots, x_n) \longmapsto x_1\cdots x_n$ 下的像. (这两个陈述的等价基本上是``独立''的定义, 不需要证明.)
  \begin{enumerate}
  \item[(a)] 证明可以假设 $K$ 有非空内部.
  \item[(b)] 证明对 $G$ 的每个紧子集 $L$ 和每个 $n \geq 1$,
    \[
    \mu^{*n}(L)\mu_G(K) \leq \langle P_{\mu}^n\mathbb{1}_{LK}, \mathbb{1}_K \rangle_{L^2(G)}.
    \]
  \item[(c)] 设 $\varepsilon > 0$. 证明存在 $\alpha > 0$ 使得对每个 $n \geq 1$,
    \[
    \langle P_{\mu}^n\mathbb{1}_{K^{\lfloor \varepsilon n \rfloor + 1}}, \mathbb{1}_K \rangle_{L^2(G)} = o(e^{-\alpha n}\|\mathbb{1}_{K^{\lfloor \varepsilon n \rfloor + 1}}\|_{L^2(G)}\|\mathbb{1}_K\|_{L^2(G)}).
    \]
  \item[(d)] 证明若在 (c) 中选择 $\varepsilon$ 足够小, 则
    \[
    \langle P_{\mu}^n\mathbb{1}_{K^{\lfloor \varepsilon n \rfloor + 1}}, \mathbb{1}_K \rangle_{L^2(G)} = o(e^{-\alpha n/2}).
    \]
  \item[(e)] 得出结论.
  \end{enumerate}
\end{enumerate}


% 附录 A、B (章号切换为字母; 先重置 chapter 计数器使 \Alph 从 A 开始)
\setcounter{chapter}{0}
\renewcommand{\thechapter}{\Alph{chapter}}
\chapter[Urysohn 引理及其推论]{Urysohn 引理及其推论}
% ============ 附录 A: \index{Urysohn 引理}Urysohn 引理及其推论 ============
% (Urysohn's Lemma and Some Consequences)
% 原书 A.1 - A.3 (原书 pp.339-340)

\section{Urysohn 引理}

\begin{definition}[正规空间 (normal space)]
\label{def:A.1.1}
拓扑空间 $X$ 称为正规 (normal) 的, 若每当有 $X$ 的两个不相交闭子集 $A$ 和 $B$, 存在 $X$ 的开子集 $U$ 和 $V$ 使得 $A \subset U$ 且 $B \subset V$.
\end{definition}

\begin{proposition}[紧 Hausdorff 或可度量空间正规]
\label{prop:A.1.2}
任何紧 Hausdorff 或可度量的拓扑空间都正规.
\end{proposition}

\begin{theorem}[Urysohn 引理 (Urysohn's lemma)]
\label{thm:A.1.3}
设 $X$ 是\index{正规拓扑空间}正规拓扑空间, $A$, $B$ 是 $X$ 的两个不相交闭子集. 则存在连续函数 $f \colon X \to [0, 1]$ 使得对每个 $x \in A$ 有 $f(x) = 0$, 对每个 $x \in B$ 有 $f(x) = 1$.
\end{theorem}

\section{\index{Tietze 延拓定理}Tietze 延拓定理}

\begin{corollary}[Tietze 延拓定理 (Tietze extension theorem)]
\label{cor:A.2.1}
设 $X$ 是正规拓扑空间, $A$ 是 $X$ 的闭子集, $f \colon A \to \CC$ 是连续函数. 则存在连续函数 $F \colon X \to \R$ 使得 $F|_A = f$ 且 $\sup_{x \in X}|F(x)| = \sup_{x \in A}|f(x)|$.
\end{corollary}

\section{应用}

\begin{corollary}[紧集与开集的分离]
\label{cor:A.3.1}
设 $X$ 是局部紧 Hausdorff 拓扑空间, $K \subset U$ 是 $X$ 的两个子集, 其中 $K$ 紧, $U$ 开. 则存在具紧支撑的连续函数 $f \colon X \to [0, 1]$ 使得 $f|_K = 1$ 且 $\operatorname{supp} f \subset U$.
\end{corollary}

\begin{proof}
由于 $X$ 局部紧, 对每个 $x \in K$, 可以找到 $x$ 的开邻域 $V_x$ 使得 $\overline{V_x}$ 紧且包含于 $U$. 有 $K \subset \bigcup_{x \in K}V_x$; 由于 $K$ 紧, 可以找到 $x_1, \ldots, x_n \in K$ 使得 $K \subset \bigcup_{i=1}^n V_{x_i}$. 设 $K' = \bigcup_{i=1}^n \overline{V_{x_i}}$. 则 $K'$ 是 $X$ 的紧子集, 它包含于 $U$ 且其内部包含 $K$. 对 $K' \subset U$ 应用同样的过程, 可以找到 $X$ 的紧子集 $K'' \subset U$ 使得其内部包含 $K'$.
空间 $K''$ 紧, 所以正规, 且其子集 $K$ 和 $K'' - \mathring{K'}$ 闭且不相交, 所以由 Urysohn 引理, 有连续函数 $f \colon K'' \to [0, 1]$ 使得 $f|_K = 1$ 且 $f|_{K'' - \mathring{K'}} = 0$. 通过设 $x \in X - K''$ 时 $f(x) = 0$ 把 $f$ 延拓到 $X$. 则 $f$ 在 $X - K'$ 上等于 $0$ (因此连续), 且它在 $\mathring{K''}$ 上也连续. 由于 $X - K'$ 和 $\mathring{K''}$ 是并集为 $X$ 的开子集, 函数 $f$ 在 $X$ 上连续. 由 $f$ 的构造显然满足所有要求的性质.
\end{proof}

\begin{corollary}[紧集上的函数延拓]
\label{cor:A.3.2}
设 $X$ 是局部紧 Hausdorff 拓扑空间, $K \subset U$ 是 $X$ 的两个子集, 其中 $K$ 紧, $U$ 开. 则对每个连续函数 $f \colon K \to \CC$, 存在具紧支撑的连续函数 $F \colon X \to \CC$ 使得:
\begin{enumerate}
\item[(a)] $\operatorname{supp}(F) \subset U$;
\item[(b)] $F|_K = f$;
\item[(c)] $\sup_{x \in X}|F(x)| = \sup_{x \in K}|f(x)|$.
\end{enumerate}
\end{corollary}

\begin{proof}
由推论 \ref{cor:A.3.1}, 可以找到具紧支撑的连续函数 $\psi \colon X \to [0, 1]$ 使得 $\psi|_K = 1$ 且 $\operatorname{supp}(\psi) \subset U$. 另一方面, 如在推论 \ref{cor:A.3.1} 的证明中那样, 可以找到紧集 $K' \subset U$ 使其内部包含 $\operatorname{supp}\psi$. 对正规空间 $K'$ 应用 Tietze 延拓, 得到连续函数 $f' \colon K' \to \CC$ 使得 $f'|_K = f$ 且 $\sup_{x \in K'}|f'(x)| = \sup_{x \in K}|f(x)|$. 定义函数 $F \colon X \to \CC$: $x \in K'$ 时 $F(x) = f'(x)\psi(x)$, $x \in X - K'$ 时 $F(x) = 0$. 此函数 $F$ 显然满足条件 (a)-(c), 所以只需验证它连续. 但这由 $F$ 在开集 $X - \operatorname{supp}(\psi)$ (其上为 $0$) 和 $\mathring{K'}$ 上连续, 且这些开集的并集是 $X$ 得出.
\end{proof}


\chapter[赋范向量空间的有用事实]{赋范向量空间的有用事实}
% ============ 附录 B: 赋范向量空间的有用事实 ============
% (Useful Things about Normed Vector Spaces)
% 原书 B.1 - B.6 (原书 pp.341-349)

\section{商范数}
见 {[}19{]} 18.15 或 {[}20{]} 1.40, 1.41.

\begin{definition}[商范数 (quotient norm)]
\label{def:B.1.1}
设 $V$ 是赋范向量空间, $W \subset V$ 是子空间. 则 $V/W$ 上的商半范数 (quotient seminorm) 由
\[
\|x + W\| = \inf_{w \in W}\|v + w\|
\]
定义.
若 $W$ 闭, 则称它为商范数 (quotient norm).
\end{definition}

\begin{proposition}[商空间的性质]
\label{prop:B.1.2}
\begin{enumerate}
\item[(i)] 前一命题的公式给出 $V/W$ 上的半范数, 当且仅当 $W$ 在 $V$ 中闭时它是范数.
\item[(ii)] 若 $V$ 是巴拿赫空间且 $W$ 在 $V$ 中闭, 则 $V/W$ 对商范数是巴拿赫空间.
\end{enumerate}
\end{proposition}

\begin{proof}
(i) 设 $v, v' \in V$, $\lambda \in \CC$. 则
\[
\|v + v' + W\| = \inf_{w \in W}\|v + v' + w\| \leq \inf_{w \in W}\|v + w\| + \inf_{w \in W}\|v' + w\| = \|v + W\| + \|v' + W\|.
\]
若 $\lambda = 0$, 则 $\lambda v \in W$, 所以 $\|\lambda v + W\| = 0$; 否则,
\[
\|\lambda v + W\| = \inf_{w \in W}\|\lambda v + w\| = \inf_{w \in W}\|\lambda(v + w)\| = |\lambda|\inf_{w \in W}\|v + w\| = |\lambda|\|v + W\|.
\]
这证明商半范数确实是 $V/W$ 上的半范数. 现在证明 $\|v + W\| = 0$ 当且仅当 $v \in W$, 这将蕴含最后的结论. 由 $\|v + W\|$ 的定义 (以及 $W$ 是子空间这一事实), $\|v + W\| = 0$ 当且仅当对每个 $\varepsilon > 0$, 存在 $w \in W$ 使得 $\|v - w\| < \varepsilon$. 这等价于 $v \in \overline{W}$.

(ii) 设 $(v_n)_{n \geq 0}$ 是 $V$ 中的序列使得 $(v_n + W)_{n \geq 0}$ 是 $V/W$ 中的 Cauchy 序列. 取 $(v_n)_{n \geq 0}$ 的子序列后, 可以假设对每个 $n \geq 0$ 有 $\|v_{n+1} - v_n + W\| < 2^{-n}$. 如下定义另一个序列 $(v'_n)_{n \geq 0}$ 使得 $n \geq 0$ 时 $v'_n \in v_n + W$, 且 $n \geq 1$ 时 $\|v'_n - v'_{n-1}\| < 2^{-n+1}$:
\begin{itemize}
\item 取 $v'_0 = v_0$.
\item 假设已有满足两个条件的 $v'_0, \ldots, v'_n$ ($n \geq 0$). 则 $\|v_{n+1} - v'_n + W\| = \|v_{n+1} - v_n + W\| < 2^{-n}$, 所以由商范数的定义, 可以找到 $w \in W$ 使得 $\|v_{n+1} - v'_n + w\| < 2^{-n}$. 取 $v'_{n+1} = v_{n+1} + w$.
\end{itemize}
由第二个条件, $(v'_n)_{n \geq 0}$ 是 Cauchy 序列, 所以它有极限 $v \in V$. 由第一个条件, 对每个 $n \geq 0$ 有 $v'_n + W = v_n + W$, 所以 $v + W$ 是序列 $(v_n + W)_{n \geq 0}$ 在 $V/W$ 中的极限.
\end{proof}

\section{开映射定理}
这也称为 Banach-Schauder 定理. 见例如 {[}19{]} 的定理 5.10.

\begin{theorem}[开映射定理 (open mapping theorem)]
\label{thm:B.2.1}
设 $V$ 和 $W$ 是巴拿赫空间, $T \colon V \to W$ 是双射有界线性变换. 则 $T^{-1} \colon W \to V$ 也有界.
\end{theorem}

\section{Hahn-Banach 定理}
见 {[}19{]} 定理 5.16 或 {[}20{]} 定理 3.2-3.7.

\begin{theorem}[Hahn-Banach 定理, 解析版本, 实情形]
\label{thm:B.3.1}
设 $V$ 是 $\R$ 上的向量空间, 设 $p \colon V \to \R$ 满足:
\begin{enumerate}
\item[(a)] 对所有 $v, v' \in V$ 有 $p(v + v') \leq p(v) + p(v')$ (即 $p$ 次可加);
\item[(b)] 对每个 $v \in V$ 和每个 $\lambda \in \R_{>0}$ 有 $p(\lambda v) = \lambda p(v)$.
\end{enumerate}
设 $E \subset V$ 是子空间, $f \colon E \to \R$ 是满足对每个 $x \in E$ 有 $f(x) \leq p(x)$ 的线性泛函.
则存在线性函数 $F \colon V \to \R$ 使得 $F|_E = f$ 且对每个 $x \in V$ 有 $F(x) \leq p(x)$.
\end{theorem}

注意在此版本中, $V$ 上没有任何范数或拓扑, 对线性泛函也没有连续性条件.

\begin{proof}
考虑偶对 $(W, g)$ 的集合 $\mathcal{X}$, 其中 $W \supset E$ 是 $V$ 的子空间, $g \colon W \to \R$ 是满足 $g|_E = f$ 且对每个 $x \in W$ 有 $g(x) \leq p(x)$ 的线性泛函. 通过说 $(W, g) \leq (W', g')$ 若 $W \subset W'$ 且 $g'|_W = g$ 来排序此集合. 假设 $(W_i, g_i)_{i \in I}$ 是 $\mathcal{X}$ 中的非空全序族, 证明它有上界. 设 $W = \bigcup_{i \in I}W_i$; 这是 $V$ 的子空间, 因为 $(W_i)_{i \in I}$ 全序 (所以对所有 $i, j \in I$, $W_i \subset W_j$ 或 $W_j \subset W_i$). 如下定义 $g \colon W \to \R$: 若 $v \in W$, 则存在 $i \in I$ 使得 $v \in W_i$, 设 $g(v) = g_i(v)$. 显然有 $g(v) \leq p(v)$. 又, 若 $j \in I$ 是满足 $v \in W_j$ 的另一个元素, 则有 $(W_i, g_i) \leq (W_j, g_j)$ 或 $(W_j, g_j) \leq (W_i, g_i)$, 两种情形都迫使 $g_i(v) = g_j(v)$, 所以定义有意义. 也容易看到 $g$ 是 $\R$-线性的, 所以 $(W, g) \in \mathcal{X}$. 这是族的上界.
所以可以对集合 $\mathcal{X}$ 应用 Zorn 引理. 设 $(W, g)$ 是 $\mathcal{X}$ 的极大元素, 证明 $W = V$. 假设 $W \neq V$, 选择 $v \in V - W$. 要把 $g$ 延拓成 $W \oplus \R v$ 上满足 $h \leq p$ 的线性泛函 $h$, 这将与 $(W, g)$ 的极大性矛盾. 这只是意味着必须选择 $h(v)$ 的值, 设为 $h(v) = \alpha \in \R$. 条件 $h \leq p$ 意味着想要对每个 $w \in W$ 和每个 $t \in \R$:
\[
h(w + tv) = g(w) + t\alpha \leq p(w + tv).
\]
若上面的不等式对某个 $t \in \R$ 和所有 $w \in W$ 成立, 则它对所有 $ct$ ($c \in \R_{>0}$) 和所有 $w \in W$ 也成立 (因为 $W$ 是子空间, 且当参数乘以 $c$ 时 $g$ 和 $p$ 的值都乘以 $c$). 所以只需对 $t = \pm 1$ 验证, 这意味着想要对每个 $w \in W$:
\[
g(w) + \alpha \leq p(w + v) \quad \text{且} \quad g(w) - \alpha \leq p(w - v).
\]
换句话说, 想要:
\[
\sup_{w \in W}(g(w) - p(w - v)) \leq \alpha \leq \inf_{w \in W}(p(w + v) - g(w)).
\]
能找到这样的 $\alpha$, 因为对所有 $w, w' \in W$,
\[
g(w) + g(w') = g(w + w') \leq p(w + w') \leq p(w + v) + p(w' - v),
\]
即
\[
g(w') - p(w' - v) \leq p(w + v) - g(w).
\]
所以得到矛盾, 可以得出结论 $W$ 本来就等于 $V$, 完成了.
\end{proof}

\begin{theorem}[Hahn-Banach 定理, 解析版本, 复情形]
\label{thm:B.3.2}
设 $V$ 是 $\CC$ 上的向量空间, 设 $p \colon V \to \R_{\geq 0}$ 是半范数,\footnote{这意味着对所有 $x, y \in V$ 和所有 $\lambda \in \CC$, $p(x + y) \leq p(x) + p(y)$ 且 $p(\lambda x) = |\lambda|p(x)$.} 设 $E \subset V$ 是 $\CC$-子空间, $f \colon E \to \CC$ 是满足对每个 $x \in E$ 有 $|f(x)| \leq p(x)$ 的线性泛函.
则存在线性函数 $F \colon V \to \CC$ 使得 $F|_E = f$ 且对每个 $x \in V$ 有 $|F(x)| \leq p(x)$.
\end{theorem}

\begin{proof}
把 $V$ 和 $E$ 看作 $\R$-向量空间, 定义 $\R$-线性泛函 $h \colon E \to \R$:
\[
h(v) = \frac{1}{2}(f(v) + \overline{f(v)}).
\]
则对每个 $v \in E$,
\[
h(v) \leq \frac{1}{2}(|f(v)| + |f(v)|) \leq p(v).
\]
注意 $p$ 满足定理 \ref{thm:B.3.1} 的条件 (a) 和 (b). 由该定理, 可以找到 $\R$-线性泛函 $H \colon V \to \R$ 使得 $H|_E = h$ 且对每个 $v \in V$ 有 $H(v) \leq p(v)$. 定义 $F \colon V \to \CC$:
\[
F(v) = H(v) + \frac{1}{i}H(iv),
\]
证明它具有所有要求的性质.
\begin{enumerate}
\item[(i)] $F$ 由构造是 $\R$-线性的, 且对每个 $v \in V$,
\[
F(iv) = H(iv) + \frac{1}{i}H(i(iv)) = iF(v).
\]
所以 $F$ 是 $\CC$-线性的.
\item[(ii)] 若 $v \in E$, 则
\[
F(v) = h(v) + \frac{1}{i}h(iv) = \frac{1}{2}(f(v) + \overline{f(v)} - if(iv) - i\overline{f(iv)}) = f(v)
\]
(因为 $f$ 是 $\CC$-线性的), 所以 $F|_E = f$.
\item[(iii)] 设 $v \in V$, 选择 $\theta \in \R$ 使得 $e^{i\theta}F(v) \in \R_{\geq 0}$. 则
\[
|F(v)| = e^{i\theta}F(v) = F(e^{i\theta}v) = H(e^{i\theta}v) - iH(ie^{i\theta}v) \in \R.
\]
由于 $H(e^{i\theta}v) \in \R$ 且 $iH(ie^{i\theta}v) \in i\R$, 必有 $iH(ie^{i\theta}v) = 0$. 所以
\[
|F(v)| = H(e^{i\theta}v) \leq p(e^{i\theta}v) = p(v).
\]
\end{enumerate}
\end{proof}

\begin{corollary}[保范延拓]
\label{cor:B.3.3}
设 $V$ 是 (实或复) 赋范向量空间, $W$ 是 $V$ 的子空间, $T_W$ 是 $W$ 上的有界线性泛函. 则存在 $V$ 上的有界线性泛函 $T$ 使得 $T|_W = T_W$ 且 $\|T\|_{\mathrm{op}} = \|T_W\|_{\mathrm{op}}$.
\end{corollary}

\begin{proof}
设 $C = \|T_W\|_{\mathrm{op}}$. 对 $p(v) = C\|v\|$ 应用 Hahn-Banach 定理. 得到延拓 $T_W$ 且对每个 $v \in V$ 满足 $|T(v)| \leq C\|v\|$ 的线性泛函 $T \colon V \to \CC$, 这意味着 $T$ 有界且 $\|T\|_{\mathrm{op}} \leq \|T_W\|_{\mathrm{op}}$. 由于不等式 $\|T_W\|_{\mathrm{op}} \leq \|T\|_{\mathrm{op}}$ 显然, 完成了.
\end{proof}

\begin{corollary}[双重对偶嵌入]
\label{cor:B.3.4}
(见 {[}19{]} 定理 5.20 和注记 5.21.) 设 $V$ 是 $K = \R$ 或 $\CC$ 上的赋范向量空间. 记 $V^* = \Hom(V, K)$. 则映射 $V \to V^{**}$, 把 $v \in V$ 送到线性泛函 $\hat{v} \colon V^* \to \CC$, $T \longmapsto T(v)$, 是等距.
特别地, 此映射单射, 这意味着有界线性泛函分离点.
\end{corollary}

现在可以推出 Hahn-Banach 定理的几何版本. (在有限维中, 它们有时称为``超平面分离定理''.)

\begin{definition}[\index{拓扑向量空间}拓扑向量空间与\index{局部凸}局部凸]
\label{def:B.3.5}
设 $V$ 是域 $K$ ($K = \R$ 或 $\CC$) 上的向量空间. 若它有拓扑使得:
\begin{itemize}
\item $(V, +)$ 是拓扑群;
\item 映射 $K \times V \to V$, $(a, v) \longmapsto av$ 连续,
\end{itemize}
则说 $V$ 是 $K$ 上的拓扑向量空间 (topological vector space).
说拓扑向量空间局部凸 (locally convex), 若每个点有凸邻域基.
\end{definition}

\begin{example}[局部凸的例子]
\label{exa:B.3.6}
(a) 任何赋范向量空间都是局部凸拓扑向量空间, 其对偶对弱*拓扑也是.
(b) 设 $(X, \mu)$ 是测度空间, $p \in (0, 1)$, 考虑空间 $L^p(X, \mu)$, 带度量
\[
d(f, g) = \int_X|f(x) - g(x)|^p\,\dd\mu(x).
\]
这使得 $L^p(X, \mu)$ 成为拓扑向量空间, 若 $\mu$ 无原子且有限 (例如 $\mu$ 是 $\R^n$ 的有界子集上的 Lebesgue 测度, 或紧群上的 Haar 测度), 则它不局部凸.
\end{example}

\begin{theorem}[Hahn-Banach 定理, 第一几何版本]
\label{thm:B.3.7}
设 $V$ 是拓扑 $\R$-向量空间, $A$ 和 $B$ 是 $V$ 的两个非空凸子集. 假设 $A$ 开且 $A \cap B = \emptyset$.
则存在连续线性泛函 $f \colon V \to \R$ 和 $c \in \R$ 使得对每个 $x \in A$ 和每个 $y \in B$,
\[
f(x) \leq c \leq f(y).
\]
\end{theorem}

我们将用作函数 $p$ 的是所谓开凸集 $C \ni 0$ 的\index{规范函数}规范函数 (gauge).

\begin{lemma}[规范函数]
\label{lem:B.3.8}
设 $C$ 是 $V$ 的非空开凸子集, 假设 $0 \in C$. 定义 $C$ 的规范函数 (gauge) $p \colon V \to \R_{\geq 0}$:
\[
p(v) = \inf\{\alpha > 0 \mid v \in \alpha C\}.
\]
则 $p$ 满足定理 \ref{thm:B.3.1} 的条件 (a) 和 (b), 且还有:
\begin{enumerate}
\item[(c)] 若 $V$ 是赋范向量空间, 则存在 $M \in \R_{\geq 0}$ 使得对每个 $v \in V$,
\[
0 \leq p(v) \leq M\|v\|.
\]
\item[(d)] $C = \{v \in V \mid p(v) < 1\}$.
\end{enumerate}
\end{lemma}

\begin{proof}
$p(\lambda v) = \lambda p(v)$ (对每个 $\lambda \in \R_{>0}$ 和每个 $v \in V$) 这一事实立即由定义得出 (不用 $C$ 的凸性或开性).
证明 (c). 由于 $C$ 开且 $0 \in C$, 存在 $r > 0$ 使得 $C \supset \{v \in V \mid \|v\| < r\}$. 则对每个 $v \in V - \{0\}$, $\frac{r}{\|v\|}v \in C$, 所以 $p(v) \leq \frac{1}{r}\|v\|$.
证明 (d). 设 $v \in C$. 由于 $C$ 开, 存在 $\varepsilon > 0$ 使得 $(1 + \varepsilon)v \in C$. 所以 $p(v) \leq \frac{1}{1+\varepsilon} < 1$. 反之, 设 $v \in V$ 满足 $p(v) < 1$. 则存在 $\alpha \in (0, 1)$ 使得 $v \in \alpha C$, 即 $\frac{1}{\alpha}v \in C$, 然后 $v = \alpha(\frac{1}{\alpha}v) + (1 - \alpha)0 \in C$, 因为 $C$ 凸.
最后证明 $p$ 次可加, 即条件 (b). 设 $v, w \in V$. 设 $\varepsilon > 0$. 由 (b) (和证明的第一个性质), $\frac{1}{p(v)+\varepsilon}v \in C$, $\frac{1}{p(w)+\varepsilon}w \in C$. 由于 $C$ 凸, 这蕴含对每个 $t \in [0, 1]$,
\[
t\frac{1}{p(v)+\varepsilon}v + (1 - t)\frac{1}{p(w)+\varepsilon}w \in C.
\]
取 $t = \frac{p(v)+\varepsilon}{p(v)+p(w)+2\varepsilon}$, 得到
\[
\frac{1}{p(v) + p(w) + 2\varepsilon}(v + w) \in C,
\]
即 $p(v + w) \leq p(v) + p(w) + 2\varepsilon$. 由于 $\varepsilon > 0$ 任意, 这蕴含 $p(v + w) \leq p(v) + p(w)$.
\end{proof}

\begin{lemma}[凸集与线性泛函]
\label{lem:B.3.9}
设 $C \subset V$ 是非空开凸子集, $v_0 \in V - C$.
则存在 $V$ 上的连续线性泛函 $F$ 使得对每个 $v \in C$, $F(v) < F(v_0)$.
\end{lemma}

\begin{proof}
可以假设 $0 \in C$ (通过平移). 设 $p \colon V \to \R_{\geq 0}$ 是 $C$ 的规范函数, 即前一引理定义的函数.
设 $E = \R v_0$, 定义线性泛函 $f \colon E \to \R$: $f(\lambda v_0) = \lambda$, 对每个 $\lambda \in \R$. 证明 $f \leq p$. 若 $\lambda \leq 0$, 则 $f(\lambda v_0) \leq 0 \leq p(\lambda v_0)$. 若 $\lambda > 0$, 则 $f(\lambda v_0) = \lambda \leq p(\lambda v_0)$, 因为 $\frac{1}{\lambda}(\lambda v_0) = v_0 \notin C$.
所以可以应用 Hahn-Banach 定理的解析形式得到线性函数 $F \colon V \to \R$ 使得对每个 $v \in V$ 有 $F(v) \leq p(v)$. 特别地, $F(v_0) = 1$, 且若 $v \in C$, 则 $F(v) \leq p(v) < 1$ (由第一个引理中的 (d)).
最后证明 $F$ 连续. 注意若 $v \in -C$, 则 $-F(v) = F(-v) < 1$. 所以对 $0$ 的开邻域 $U := C \cap (-C)$ 中的每个 $v$, 有 $|F(v)| < 1$. 若 $\varepsilon > 0$, 则 $\varepsilon U$ 是 $V$ 中 $0$ 的开邻域, 且对每个 $v \in \varepsilon U$ 有 $|F(v)| < \varepsilon$. 所以 $F$ 在 $0$ 处连续. 由于 $F$ 线性且平移在 $V$ 上连续, 这蕴含 $F$ 在 $V$ 的每个点连续.
\end{proof}

\begin{proof}[定理的证明]
设 $C = A - B = \{x - y \mid x \in A, y \in B\}$. 则 $C$ 显然凸, $C$ 开因为它等于 $\bigcup_{y \in B}(A - y)$, 且 $0 \notin C$ 因为 $A \cap B = \emptyset$. 所以可以应用上面的第二个引理得到连续线性泛函 $f \colon V \to \R$ 使得对每个 $x \in C$ 有 $f(x) < 0$. 则对每个 $x \in A$ 和每个 $y \in B$, $f(x) < f(y)$. 所以对 $f$ 和 $c = \sup_{x \in A}f(x)$ 结论成立.
\end{proof}

\begin{theorem}[Hahn-Banach 定理, 第二几何版本]
\label{thm:B.3.10}
设 $V$ 是局部凸拓扑 $\R$-向量空间, $A$ 和 $B$ 是 $V$ 的两个非空凸子集. 假设 $A$ 闭, $B$ 紧, 且 $A \cap B = \emptyset$.
则存在连续线性泛函 $f \colon V \to \R$ 和 $c \in \R$ 使得对每个 $x \in A$ 和每个 $y \in B$,
\[
f(x) < c < f(y).
\]
\end{theorem}

\begin{proof}
先找 $V$ 中 $0$ 的凸开邻域 $U$ 使得 $(A + U) \cap (B + U) = \emptyset$. (注意: 这只用 $V$ 局部凸以及 $A$ 闭、$B$ 紧这一事实, 不用 $A$ 和 $B$ 凸.)
对每个 $x \in B$, 选择对称凸开邻域 $U_x$ 使得 $(x + U_x + U_x + U_x) \cap A = \emptyset$; 由于 $U_x$ 对称, 这等价于说 $(x + U_x + U_x) \cap (A + U_x) = \emptyset$. 由于 $B$ 紧, 可以找到 $x_1, \ldots, x_n \in B$ 使得 $B \subset \bigcup_{i=1}^n(x_i + U_{x_i})$. 设 $U = \bigcap_{i=1}^n U_{x_i}$. 则 $U$ 是 $0$ 的凸开邻域, 且 $B + U \subset \bigcup_{i=1}^n(x_i + U_{x_i} + U)$, $A + U \subset \bigcap_{i=1}^n(A + U_{x_i})$, 所以 $(B + U) \cap (A + U) = \emptyset$.
集合 $A + U$ 和 $B + U$ 凸且开, 所以由定理 \ref{thm:B.3.7}, 存在连续线性泛函 $f \colon V \to \R$ 和 $c_0 \in \R$ 使得对每个 $x \in A + U$ 和每个 $y \in B + U$ 有 $f(x) \leq c_0 \leq f(y)$. 由于 $B$ 紧且 $f$ 连续, 存在 $y_0 \in B$ 使得 $f(y_0) = \min_{y \in B}f(y)$. 特别地, $c_0 < \min_{y \in B}f(y)$. 选择 $c \in \R$ 使得 $c_0 < c < \min_{y \in B}f(y)$. 则对每个 $x \in A$ 和每个 $y \in B$ 有 $f(x) < c < f(y)$.
\end{proof}

\section{Banach-Alaoglu 定理}
见 {[}18{]} 第 15.1 节或 {[}20{]} 定理 3.15. 此定理也称为``Alaoglu 定理''.

\begin{theorem}[Banach-Alaoglu 定理]
\label{thm:B.4.1}
设 $V$ 是赋范向量空间. 则 $\Hom(V, \CC)$ 中的闭单位球对弱*拓扑紧 Hausdorff.
\end{theorem}

与下列通常称为``Riesz 引理''或``Riesz 定理''的结果比较 (见 {[}18{]} 第 13.3 节或 {[}20{]} 定理 1.22):

\begin{theorem}[Riesz 定理]
\label{thm:B.4.2}
设 $V$ 是赋范向量空间. 则 $V$ 的闭单位球紧当且仅当 $V$ 有限维.
\end{theorem}

\section{\index{Krein-Milman 定理}Krein-Milman 定理}
见 {[}18{]} 第 14.6 节 (或 {[}20{]} 定理 3.23).

\begin{definition}[端点 (extremal point)]
\label{def:B.5.1}
设 $V$ 是 $\R$-向量空间, $C$ 是 $V$ 的凸子集. 若每当 $x = ty + (1 - t)z$ ($t \in (0, 1)$, $y, z \in C$), 必有 $y = z = x$, 则说 $x \in C$ 端点 (extremal).
\end{definition}

\begin{theorem}[Krein-Milman 定理]
\label{thm:B.5.2}
设 $V$ 是局部凸拓扑 $\R$-向量空间, $K$ 是 $V$ 的非空紧凸子集. 则 $K$ 是其端点集合的凸包的闭包.
\end{theorem}

\begin{lemma}[紧凸集的端点]
\label{lem:B.5.3}
设 $V$ 是局部凸拓扑 $\R$-向量空间, $K$ 是 $V$ 的非空紧凸子集. 则 $K$ 有端点.
\end{lemma}

\begin{proof}
说 $K$ 的子集 $S$ 端 (extremal), 若对每个 $x \in S$, 若 $x = ty + (1 - t)z$, $y, z \in K$, $t \in (0, 1)$, 则必有 $y, z \in S$. (注意点 $x \in K$ 端点当且仅当 $\{x\}$ 端.)
设 $\mathcal{X}$ 是 $K$ 的非空闭端子集的集合, 按反向包含排序. 设 $Y$ 是 $\mathcal{X}$ 的非空全序子集, 证明它有极大元素. 由于 $Y$ 全序, 对所有 $T_1, \ldots, T_n \in Y$, 存在 $i \in \{1, \ldots, n\}$ 使得对所有 $j \in \{1, \ldots, n\}$ 有 $T_i \subset T_j$, 然后 $T_1 \cap \cdots \cap T_n \supset T_i \neq \emptyset$. 由于 $K$ 紧, 这蕴含每个 $S := \bigcap_{T \in Y}T$ 非空. 集合 $S$ 显然闭, 所以若能证明它端, 就完成了. 设 $x \in S$, 写 $x = ty + (1 - t)z$, $y, z \in K$, $t \in (0, 1)$. 对每个 $T \in Y$, 由于 $T$ 端, 必有 $y, z \in T$. 所以 $y, z \in S$, $S$ 确实端.
由 Zorn 引理, 集合 $\mathcal{X}$ 有极大元素, 称之为 $S$. 为完成证明, 只需证明 $S$ 是单点集. 若 $|S| \geq 2$, 设 $x, y \in S$ 满足 $x \neq y$. 由 Hahn-Banach 定理的几何版本 (定理 \ref{thm:B.3.10}), 存在连续线性泛函 $f \colon V \to \R$ 使得 $f(x) < f(y)$. 由于 $S$ 紧, 连续函数 $f$ 在 $S$ 上达到最小值. 设 $m = \min_{z \in S}f(z)$, $S' = \{z \in S \mid f(z) \leq m\}$. 则 $S'$ 闭, 由刚才的观察非空, 且 $S' \neq S$, 因为 $y \notin S'$. 证明 $S'$ 端, 这将给出矛盾 (且蕴含 $S$ 本来是单点集). 设 $z \in S'$, 写 $z = tz' + (1 - t)z''$, $z', z'' \in K$, $t \in (0, 1)$. 由于 $S$ 端, 有 $z', z'' \in S$. 由 $m$ 的定义, $m = f(z) = tf(z') + (1 - t)f(z'') \leq tm + (1 - t)m$, 这迫使 $m = f(z') = f(z'')$, 即 $z', z'' \in S'$.
\end{proof}

\begin{proof}[定理的证明]
设 $L$ 是 $K$ 的端点集合的凸包的闭包. 则 $L$ 凸、闭且包含于 $K$; 特别地, $L$ 也紧. 假设 $L \neq K$, 设 $x \in K \setminus L$. 由 Hahn-Banach 定理的几何版本 (定理 \ref{thm:B.3.10}), 存在连续线性泛函 $f \colon V \to \R$ 使得 $\max_{y \in L}f(y) < f(x)$. 设 $M = \max_{z \in K}f(z)$, $K' = \{z \in K \mid f(z) = M\}$. 则 $K'$ 是 $K$ 的闭凸子集 (所以紧), 且 $K' \cap L = \emptyset$. 由引理, $K'$ 必有端点 $z$, 且容易看到 (如在引理的证明中) $z$ 也是 $K$ 的端点. 但此时 $z$ 应在 $L$ 中, 矛盾.
\end{proof}

\section{Stone-Weierstrass 定理}
见 {[}18{]} 第 12.3 节或 {[}20{]} 定理 5.7 (紧空间的情形).

\begin{theorem}[Stone-Weierstrass 定理]
\label{thm:B.6.1}
设 $X$ 是局部紧 Hausdorff 拓扑空间, $A$ 是 $\mathcal{C}_0(X)$ 的 $\CC$-子代数, 使得:
\begin{enumerate}
\item[(a)] 对每个 $f \in A$, 函数 $x \longmapsto \overline{f(x)}$ 也在 $A$ 中;
\item[(b)] 对所有满足 $x \neq y$ 的 $x, y \in X$, 存在 $f \in A$ 使得 $f(x) \neq f(y)$ (``$A$ 分离 $X$ 的点'');
\item[(c)] 对每个 $x \in X$, 存在 $f \in A$ 使得 $f(x) \neq 0$ (``$A$ 在 $X$ 上任何地方不消失'').
\end{enumerate}
则 $A$ 在 $\mathcal{C}_0(X)$ 中稠密.
\end{theorem}


% 参考文献
% ============ 参考文献 (原书 Bibliography, pp.353-354) ============
% 由 AJbook.cls 的 biblatex 机制处理, sorting=nyt 与原书字母序一致
\cleardoublepage
\phantomsection
\nocite{*}
\printbibliography[title={参考文献}]


% 索引
\printindex

\end{document}
